\newcommand{\DoNotLoadEpstopdf}{}
\documentclass[11pt,reqno]{amsart}
\usepackage[T1]{fontenc}
\usepackage{lmodern,amsmath,amssymb,amsthm,mathtools,microtype,array}
\usepackage[margin=1.1in]{geometry}
\usepackage[hidelinks]{hyperref}
\hypersetup{pdftitle={Collision records, arithmetic, and raw local inversion in a triangular Lorentz gas},pdfauthor={Qian Qi}}
\setlength{\emergencystretch}{2em}
\allowdisplaybreaks[2]
\numberwithin{equation}{section}
\newtheorem{theorem}{Theorem}[section]
\newtheorem{maintheorem}{Theorem}
\renewcommand{\themaintheorem}{\Alph{maintheorem}}
\newtheorem{lemma}[theorem]{Lemma}
\newtheorem{proposition}[theorem]{Proposition}
\newtheorem{corollary}[theorem]{Corollary}
\theoremstyle{definition}\newtheorem{definition}[theorem]{Definition}
\theoremstyle{remark}\newtheorem{remark}[theorem]{Remark}
\newcommand{\R}{\mathbb R}\newcommand{\Z}{\mathbb Z}\newcommand{\T}{\mathbb T}
\newcommand{\dd}{\,\mathrm d}\newcommand{\one}{\mathbf1}
\newcommand{\supp}{\operatorname{supp}}\newcommand{\TV}{\operatorname{TV}}
\title[Collision records and raw local inversion]{Collision records, arithmetic, and raw local inversion in a triangular Lorentz gas}
\author{Qian Qi}
\date{October 7, 2026. A2-DYN, revision 33}
\subjclass[2020]{37D50, 37A50, 60F05, 42A38}
\keywords{Sinai billiards, collision records, local limits, inducing, Fourier inversion, geometric response}
\begin{document}
\raggedbottom
\begin{abstract}
We study the joint lattice displacement, collision count and flight time
of genuine returns to a moving section in a triangular periodic Lorentz
gas, uniformly for radii in $[0.45,0.47]$. Bounded collision compensation
and exact stopping give Gaussian and functional laws, all fixed marked
moments and quantitative convergence of the actual covariances.
Measurable phase holonomy and finite physical covers prove joint
nondegeneracy. Fixed-order spectral jets and damped removal of smoothing
give integrated Gaussian comparisons at a prescribed return count,
with one complex actual-return insertion. A weighted-volume criterion
extends these estimates to a nonrectangular region in all four
coordinate directions, with fixed expansion orders and a retained rate.

For the raw density, finite-count power--logarithm extraction gives an
absolutely invertible residual. Exact edge--residual cancellation makes
the signed raw correction coherent under changes of cutoff. We now
bound the signed integral of that correction over unsmoothed windows,
and obtain their actual Gaussian denominators by order-preserving
Fourier envelopes. A maximal return-window estimate transfers a
three-dimensional denominator to observations at deterministic physical
time under a section-start law. The relative symmetric difference of
physical and return events tends to zero at an explicit rate for the
stated mesoscopic lattice and time windows. A wider collision-clock
band further gives finer unsmoothed windows under equilibrium flow.
The physical-time length bias is treated exactly, and stationary
physical, deterministic-return and last-return events have asymptotically
identical conditional laws. A joint block-frequency estimate and a
rare-event tightness argument identify the common physical path limit
as a Gaussian bridge to the observed endpoint. They also give explicit
denominators and conditional path limits for fixed interior cylinder
observations. A separate geometric extraction now removes a quantified
small-variation measure from every finite itinerary. Its residual has
an explicit second-derivative budget at a linear collision cutoff.
This gives finite-band inversion of the full actual law at fixed
lattice labels and bounded roof intervals, with arbitrarily small
polynomial errors and a roof bandwidth whose logarithm is $O(n\log n)$.
A positive finite-band kernel further approximates the complete mixed
law in total variation and the unnormalized microscopic conditional
measure on the original trajectory space. The latter estimate is uniform
against every bounded measurable path selector. It also holds under
the actual stationary entrance bias and retains singleton lattice labels.
For the original stationary physical observation at a prescribed collision
count, exact integration of the initial flight age gives a different
inversion: both within-flight cell corrections are retained, the physical
time profiles have a uniform second-derivative measure bound, and the
far-frequency error is $O(B^{-1})$. We evaluate its three-dimensional
Gaussian central term and obtain a pointwise Gaussian-plus-complement
identity with polynomial bandwidth. The associated positive event
likelihood is accurate in source total variation for arbitrary bounded
selectors. The original return reconstruction bandwidth may still be
$\exp(O(n\log n))$.
These finite-band representations do not evaluate the complete
middle-frequency integrals.
The pointwise common correction, that complement and the microscopic
weighted denominator and physical-event comparison remain the
requirements of the original raw mixed-density local limit problem.
\end{abstract}
\maketitle

\paragraph{The physical record.}
Let $T_R$ be the collision map, $Y_R^*$ the section constructed below,
and $F_R^*$ its actual first-return map. Write
$J_{n,R}=(K_{n,R},N_{n,R},T_{n,R})$ for the four-coordinate record of
$n$ returns, with the two coordinates of $K_{n,R}$ displayed separately.
Its law is taken under the normalized collision measure $\nu_R^*$.
The following synopsis records the unconditional probabilistic
conclusions. The geometry, normalization, and limiting matrices are
specified intrinsically in the cited theorems.

\begin{maintheorem}[Gaussian laws of the physical record]
\label{thm:intro-gaussian}
For $R\in I=[0.45,0.47]$ set
\[
 c_* = \frac{91}{10000\pi},\qquad
 \bar\tau_R=\frac{\sqrt3/2-\pi R^2}{2R},\qquad
 \bar G_R=(0,0,c_*^{-1},\bar\tau_R/c_*).
\]
There is a continuous positive semidefinite matrix $D_R$ such that,
for every $K<\infty$ and $n\ge2$,
\[
 \sup_{R\in I,\ |t|\le K}
 \left|\int e^{it\cdot(J_{n,R}-n\bar G_R)/\sqrt n}\dd\nu_R^*
              -e^{-t^{\mathsf T}D_Rt/2}\right|
 \le C_Kn^{-1/26}\sqrt{\log(2+n)}.
\]
The polygonal centered return processes converge, uniformly in $R$,
to Brownian motion with covariance $D_R$. The stationary physical
processes of lattice displacement and centered collision count also
satisfy a parameter-uniform functional Gaussian limit.

More precisely, $D_R=c_*^{-1}\Gamma_R$, where $\Gamma_R$ is the
absolutely convergent collision Green--Kubo matrix of
$h_R=f_R-\bar G_R\one_{Y_R^*}$. Its kernel consists exactly of directions
$v$ for which $v\cdot(G_R-\bar G_R)$ is an $L^2$ coboundary for
$F_R^*$. No nondegeneracy assumption is used in these conclusions.
\end{maintheorem}
\begin{proof}
Theorems~\ref{thm:collision-covariance},
\ref{thm:actual-record-gaussian}, \ref{thm:gaussian-kernel-coboundary},
and \ref{thm:actual-return-functional}, together with
Corollary~\ref{cor:physical-functional-gaussian}, prove the assertions.
\end{proof}

\begin{maintheorem}[An integrated central band]
\label{thm:intro-major-arc}
For the same actual return record and the matrix $D_R$ in Theorem A,
\[
 \sup_{R\in I}\int_{|v|\le2n^{1/200}}
 \left|\int e^{iv\cdot(J_{n,R}-n\bar G_R)/\sqrt n}\dd\nu_R^*
                       -e^{-v^{\mathsf T}D_Rv/2}\right|\dd v
       \le Cn^{-3/280}\sqrt{\log(2+n)}.
\]
The variable $v$ is rescaled: the corresponding physical frequencies
have size at most $2n^{-1/2+1/200}$. The same estimate holds for complex
initial insertions whose zero extensions to the collision cylinder
have uniformly bounded supremum and variation norms, with the Gaussian
multiplied by the integral of the insertion.
\end{maintheorem}
\begin{proof}
Apply Theorem~\ref{thm:growing-major-arc}; the original probability is
the admissible insertion $a_R=s_R$.
\end{proof}

\paragraph{Complex insertion convention.}
In Theorem B the insertion is a density with respect to the collision
probability $\nu$. Its amplitude $\alpha_R=\int a_R\dd\nu$ may be
complex; the associated pushforward is a finite complex measure, not
in general a probability. The original section probability corresponds
to $a_R=s_R=c_*^{-1}\one_{Y_R^*}$.

\begin{maintheorem}[A mark at any actual return]
\label{thm:intro-marked-return}
Let $a$ be a complex function on the collision cylinder, zero outside
$Y_R^*$, and put $M_a=\|a\|_\infty$, $V_a=\|a\|_{\mathrm{BV}}$,
and $\alpha_a=\int a\dd\nu$. For $0\le k\le n$ set
\[
 C^{[k],a}_{n,R}(v)=\int_{Y_R^*}a((F_R^*)^k x)
       e^{iv\cdot(J_{n,R}(x)-n\bar G_R)/\sqrt n}\dd\nu(x).
\]
There is a constant $C$, independent of $R,n,k,a$, such that, for $n\ge2$,
\begin{align*}
 \int_{|v|\le2n^{1/200}}
 \left|C^{[k],a}_{n,R}(v)-\alpha_a e^{-v^{\mathsf T}D_Rv/2}\right|\dd v
 \le C\left(M_an^{-3/280}\sqrt{\log(2+n)}+V_an^{-9/175}\right).
\end{align*}
Thus the same central band accommodates initial, terminal, and one
intermediate return-state insertion, including a mark whose index
varies with $n$. The displayed transform is integrated against the
unnormalized restricted collision measure $\nu|_{Y_R^*}$; division by
$c_*$ gives the corresponding transform under the section probability
$\nu_R^*$. No positivity of $a$ or nonsingularity of $D_R$ is assumed.
\end{maintheorem}
\begin{proof}
This is Theorem~\ref{thm:marked-return-band}.
\end{proof}


\begin{maintheorem}[Actual second moments and a marked quadratic law]
\label{thm:intro-actual-moments}
Put $U_{n,R}=J_{n,R}-n\bar G_R$. Uniformly for $R\in I$,
\[
 \int |U_{n,R}|^4\dd\nu_R^*\le Cn^2,\qquad
 \left\|\frac1n\int U_{n,R}\otimes U_{n,R}\dd\nu_R^*-D_R\right\|
       \le Cn^{-1/6}.
\]
The count component satisfies $e_3^{\mathsf T}D_Re_3\ge d_N>0$
uniformly in $R$, where $e_3=(0,0,1,0)$.
For the complex insertion and normalization of Theorem C,
\begin{align*}
 &\left\|\int_{Y_R^*}a((F_R^*)^kx)
       \frac{U_{n,R}(x)\otimes U_{n,R}(x)}n\dd\nu(x)-\alpha_aD_R\right\|\\
 &\hspace{20mm}\le C\left(M_an^{-1/6}+(M_a+V_a)n^{-1}\right),
       \qquad 0\le k\le n.
\end{align*}
The matrix $D_R$ is also the uniform limit of the induced Green--Kubo
formula with Cesaro weights. No induced mixing or absolute convergence
of its correlation series is assumed.
\end{maintheorem}
\begin{proof}
Apply Theorems~\ref{thm:marked-L2-stopping} and
\ref{thm:marked-quadratic-moments}, and
Corollaries~\ref{cor:induced-cesaro-covariance} and
\ref{cor:positive-count-variance}.
\end{proof}

\begin{maintheorem}[Uniform nondegeneracy of the joint record]
\label{thm:intro-joint-nondegeneracy}
The covariance of the same actual four-coordinate return record satisfies
\[
 d_0 I_4\le D_R\le D_0 I_4\qquad(R\in I)
\]
for constants $0<d_0\le D_0<\infty$. Thus every nonzero real joint
direction has positive Gaussian variance. For all sufficiently large
$n$, the normalized actual covariance in Theorem D is at least
$(d_0/2)I_4$, uniformly in $R$. The three-dimensional physical-time
fluctuation covariance is uniformly positive definite as well.
\end{maintheorem}
\begin{proof}
Apply Theorem~\ref{thm:joint-uniform-nondegeneracy} and
Corollary~\ref{cor:actual-uniform-ellipticity}.
\end{proof}
In every subsequent use of the matrix from Theorems A--D, the stronger
uniform positive definiteness of Theorem E is available. The positive
semidefinite formulation of the earlier statements records their weaker
logical input, not a remaining covariance hypothesis.


\begin{maintheorem}[Physical phase arithmetic and quantitative defects]
\label{thm:intro-quantitative-phases}
Let $\omega=(u,s,t)\in\T^2\times\T\times\R$ and
$\phi_{R,\omega}=u\cdot\kappa_R+s+t\tau_R$ on the collision cylinder.
A measurable circle-valued solution of
$q\circ T_R=e^{i\phi_{R,\omega}}q$ exists exactly when $\omega=0$;
in that case it is constant almost everywhere.

For each fixed $R$ and compact nonzero frequency set $\mathcal K$,
there is $c_{R,\mathcal K}>0$ such that, for $H\ge2$,
\[
 \inf_{\omega\in\mathcal K}
 \|q\circ T_R-e^{i\phi_{R,\omega}}q\|_1
                  \ge\frac{c_{R,\mathcal K}}{1+\log H}
 \quad (|q|=1,\ \|q\|_{\mathrm{BV}}\le H).
\]
For complex $f$ with $\|f\|_2=1$ and
$\|f\|_\infty+\|f\|_{\mathrm{BV}}\le H$, the corresponding bound is
\[
 \inf_{\omega\in\mathcal K}
 \|f\circ T_R-e^{i\phi_{R,\omega}}f\|_2
                  \ge\frac{c_{R,\mathcal K}}{(1+\log H)^2}.
\]
No lower bound on $|f|$ is assumed. For fixed $H$ and $\mathcal K$,
the second infimum is also positive uniformly over $R\in I$.
The latter assertion does not specify a uniform logarithmic rate in $H$.
\end{maintheorem}
\begin{proof}
Theorem~\ref{thm:complete-physical-phase} proves the exact assertion.
Theorems~\ref{thm:compact-log-phase-defect},
\ref{thm:compact-vector-defect} and
\ref{thm:uniform-compact-bv-separation} prove the defect bounds.
\end{proof}

\begin{maintheorem}[Near-origin defects for the actual return record]
\label{thm:intro-actual-return-defects}
Set $g_R=G_R-\bar G_R$ and
\[
 \Lambda=1+\log(H/r),\qquad K=\Lambda\log(2+\Lambda),
                 \qquad H\ge2,\quad 0<r<1/2.
\]
There are $a,c,r_*>0$, independent of $R$, such that $r\le r_*$
and $rK\le a$ imply the following conclusions for $|z|=r$.
Every unit-modulus section function $q$ whose zero extension has
variation at most $H$ satisfies
\[
 \|q\circ F_R^*-e^{iz\cdot g_R}q\|_{L^1(\nu_R^*)}\ge cr^2.
\]
Every complex section function $f$ with
$\|f\|_{L^2(\nu_R^*)}=1$ and
$\|f\|_\infty+\|f^0\|_{\mathrm{BV}}\le H$ satisfies
\[
 \|f\circ F_R^*-e^{iz\cdot g_R}f\|_{L^2(\nu_R^*)}
                                    \ge cr^2/K.
\]
The function may vanish. With $H\le H_0n^\zeta$, these bounds apply
uniformly on the physical annulus
$2n^{-99/200}\le|z|\le n^{-2/5}$ for all sufficiently large $n$.
The corresponding rescaled radii are $2n^{1/200}$ and $n^{1/10}$.
\end{maintheorem}
\begin{proof}
Apply Theorems~\ref{thm:actual-return-circle-defect} and
\ref{thm:actual-return-vector-defect}, and
Corollary~\ref{cor:actual-annulus-defect}.
\end{proof}
\paragraph{The spectral parameter in Theorem G.}
The displayed defect in Theorem G has unit spectral parameter for the
centered return multiplier. A constant per return is not the same as a
constant per collision. The next theorem treats both the physical
frequency $z$ and the return spectral phase $\xi$, with
$\lambda=e^{i\xi}$, and strengthens the regularity constraint.

\begin{maintheorem}[A joint peripheral neighborhood]
\label{thm:intro-peripheral-neighborhood}
For $\xi\in[-\pi,\pi]$, set
\[
 \rho=(|z|^2+\xi^2)^{1/2},\quad
 \Lambda=1+\log(H/\rho),\quad K=\Lambda\log(2+\Lambda).
\]
There are $a,b,\rho_0>0$, independent of $R$, such that
$0<\rho\le\rho_0$ and $\rho^2K\le a$ imply
\[
 \|q\circ F_R^*-e^{i(\xi+z\cdot g_R)}q\|_{L^1(\nu_R^*)}
       \ge b\rho^2
\]
for every unit section phase with $\|q^0\|_{\mathrm{BV}}\le H$,
and
\[
 \|f\circ F_R^*-e^{i(\xi+z\cdot g_R)}f\|_{L^2(\nu_R^*)}
       \ge b\rho^2/K
\]
for every normalized complex section function with
$\|f\|_\infty+\|f^0\|_{\mathrm{BV}}\le H$. Zeros are allowed.
These estimates include $z=0$, $\xi\ne0$. They are local in the joint
parameters, not estimates over the full peripheral circle.

For the physical annulus $2n^{-99/200}\le|z|\le n^{-2/5}$,
$|\xi|\le n^{-\gamma}$, and any
$0<\beta<\min\{4/5,2\gamma\}$, they apply to regularity budgets
$H\le H_0\exp(n^\beta)$ for all sufficiently large $n$.
The rescaled annulus is $2n^{1/200}\le|\sqrt n\,z|\le n^{1/10}$.
\end{maintheorem}
\begin{proof}
Apply Theorem~\ref{thm:joint-return-spectral-defect} and
Corollary~\ref{cor:peripheral-annulus-budgets}.
\end{proof}

Theorem~\ref{thm:compressed-peripheral-resolvent} converts these physical
defects into a lower singular-value bound for finite-rank orthogonal
compressions of the actual return operator. The reconstruction has
supremum and variation budget $C/h$ at mesh $h$; no positive lower
modulus is assumed. Proposition~\ref{prop:compressed-finite-time-comparison}
separately bounds the difference between the compressed powers and the
unmodified characteristic function. A mesh-uniform spectral limit and
the complete complementary-frequency integral are not inferred from
that finite-time comparison.


\begin{maintheorem}[Finite-count extraction of the raw density]
\label{thm:intro-finite-count-extraction}
Fix $R\in I$, $n\ge1$ and a collision-count cutoff $L$. For every
bounded finite-record subanalytic weight in
Definition~\ref{def:finite-record-weight}, the restricted actual law
\[
 \mu_{n,R}^{w,L}=(J_{n,R})_*
                    (w\one_{\{N_{n,R}\le L\}}\nu_R^*)
\]
has an exact decomposition $\mu_{n,R}^{w,L}=E^L+Q^L$.
The extracted density is a finite sum of cutoff one-sided terms
$x^\alpha(\log x)^k$, with $\alpha\in\mathbb Q$, $-1<\alpha\le1$,
and $k$ a nonnegative integer. It includes all singular boundary
contributions to the finite-count law. Each lattice component of the
residual density $r^L$ belongs to $W^{2,1}(\R)$. In particular, if
$A_2=\sum_\ell\|\partial_t^2 r^L(\ell,\cdot)\|_1$, then
\[
 \widehat Q^L\in L^1(\T^3\times\R),\qquad
 \int_{\T^3\times\{|b|>B\}}|\widehat Q^L|\dd\omega
                   \le 2(2\pi)^3 A_2/B .
\]
Here $A_2$ is finite at each fixed $(n,L,R,w)$; no uniform long-time
bound for this quantity is part of the assertion.

For the original unweighted law, choose $L_n=\lceil\lambda n\rceil$
with $\lambda>c_*^{-1}+1$. Its density agrees exactly with the restricted
density on every fixed central window for large $n$.
The finite extraction gives the raw inversion identity
\eqref{eq:count-localized-exact-inversion} on that window. The
high-count measure enters only through a low-frequency convolution,
whose normalized central contribution is $O(n^{-P})$ for every $P>0$.
The physical cutoff is $2n^{-99/200}$, or $2n^{1/200}$ after rescaling.
\end{maintheorem}
\begin{proof}
Apply Theorem~\ref{thm:finite-count-raw-extraction},
Lemma~\ref{lem:central-count-support}, and
Theorem~\ref{thm:finite-extraction-central-inversion}.
\end{proof}

Theorem I supplies the complete finite-count residual, rather than only
an initial-coordinate variation bound. Corollary~\ref{cor:finite-extraction-error-budget}
displays the remaining uniform residual-growth, complementary-integral
and local-edge estimates for the original raw density.


\begin{maintheorem}[Exponential reconstruction and multiple observations]
\label{thm:intro-window-reconstruction}
The finite-record variation budget may be taken as $Ce^{CL}$.
Consequently, with $\rho=(|z|^2+\xi^2)^{1/2}$ and
$\Lambda=1+\log(H/\rho)$, the joint return spectral estimates of
Theorem H strengthen to
\[
 \|\mathcal E_{R,z,\xi}q\|_1\ge b\rho^2,\qquad
 \|\mathcal E_{R,z,\xi}f\|_2\ge b\rho^2/\Lambda
\]
whenever $0<\rho\le\rho_0$ and $\rho^2\Lambda\le a$.
Here $q$ is a unit section phase of zero-extension variation at most
$H$, while $f$ is a normalized complex section function with supremum
plus zero-extension variation at most $H$. The constants are uniform
in $R$, and no positive lower modulus of $f$ is required.

There are constants $a_*,\theta_*,\beta_*>0$ for the original family
with the following property. Let $m_n\le a_*\log n$ and choose any
$0\le\ell_1<\cdots<\ell_q\le m_n$. Let $E_j\subset Y_R^*$ be
observation sets whose zero-extended indicators have total variation
at most a fixed power of $\log(2+n)$. If the unchanged event
\[
 A_{n,k,R}=\bigcap_{j=1}^{q}
       \{x:(F_R^*)^{k+\ell_j}x\in E_j\},\qquad
                       0\le k\le n-m_n,
\]
has probability at least $p_0 n^{-\beta_*}$, then, for
$U_{n,R}=J_{n,R}-n\bar G_R$,
\[
 \sup_{R,k}\int_{|v|\le2n^{\theta_*}}
 \left|\mathbb E[e^{iv\cdot U_{n,R}/\sqrt n}\mid A_{n,k,R}]
                          -e^{-v^{\mathsf T}D_Rv/2}\right|\dd v
                           \longrightarrow0.
\]
Its conditional normalized mean tends to zero and its conditional
covariance tends to $D_R$, uniformly. All expectations use $\nu_R^*$.
The physical cutoff is $2n^{-1/2+\theta_*}$. Explicit constants and
more general probability, variation and frequency tradeoffs are given
in \eqref{eq:window-exponent-margins} and
\eqref{eq:nonempty-window-budget}.
\end{maintheorem}
\begin{proof}
Theorems~\ref{thm:exponential-finite-record} and
\ref{thm:single-log-return-defect} prove the first part.
Corollary~\ref{cor:logarithmic-window-conditioning} and
Theorem~\ref{thm:window-weighted-moments}, with
\eqref{eq:nonempty-window-budget}, prove the remaining assertions.
\end{proof}

The single-logarithm estimate is local in the joint physical and return
spectral parameters; its exact lifting identity remains valid at every
return phase. The growing observation window is a new class, not an
assumed BV bound on a long pullback. Its smaller positive central band
is additional to the unchanged $n^{1/200}$ band of Theorem C.
Neither statement replaces the complementary-integral or local-edge
estimates for the raw density. The original multiple-time event is
never exchanged for an event with a collision cutoff.

\begin{maintheorem}[Direct annular averages for the uncompressed record]
\label{thm:intro-direct-annular-averages}
Let $\mathcal U_{R,z}f=e^{-iz\cdot g_R}f\circ F_R^*$ on
$L^2(\nu_R^*)$. For all sufficiently large $n$, all $T\ge n$,
$N\ge0$, $a\in L^2(\nu_R^*)$, and physical frequencies
\[
 2n^{-99/200}\le|z|\le n^{-2/5},
 \qquad 2n^{1/200}\le|\sqrt n\,z|\le n^{1/10},
\]
one has, uniformly in $R$,
\begin{align*}
 \frac1T\sum_{j=0}^{T-1}
 \left|\int a e^{-iz\cdot U_{N+j,R}}\dd\nu_R^*\right|^2
                          &\le C n^{-4/495}\|a\|_2^2,\\
 \sup_{\xi\in\R}
 \left\|\frac1T\sum_{j=0}^{T-1}
         e^{-ij\xi}\mathcal U_{R,z}^{N+j}\one\right\|_2^2
                          &\le C n^{-4/495}.
\end{align*}
For an unchanged event $A$ of probability $p>0$, the average of the
squared conditional characteristic functions is at most
$C p^{-1}n^{-4/495}$. No BV condition on $\one_A$ is needed.

If $\varsigma_{R,ru}$ is the spectral probability of $\one$ for
$\mathcal U_{R,ru}$, where $|u|=1$, its principal angle divided by
$r^2$ converges in law, uniformly in $R,u$, to the Cauchy law of scale
$u^{\mathsf T}D_Ru/2$.
\end{maintheorem}
\begin{proof}
Theorems~\ref{thm:direct-moving-averages} and
\ref{thm:weighted-annular-averages}, with
Corollary~\ref{cor:direct-annular-averages}, prove the averaged bounds.
Theorem~\ref{thm:spectral-Cauchy-limit} proves the spectral conclusion.
\end{proof}

\paragraph{The prescribed-count target.}
Theorem K concerns averages in the return count. The raw local theorem
instead requires, at the prescribed $n$, an estimate of the form
$n^2\int|1-\chi_n|\,|H_{n,R}|\dd\omega=o(1)$ for the appropriate
edge-subtracted residual, together with its local extracted-edge term.
The averaged bound does not imply this estimate: its existing annular
normalization gives only $n^{2/5-2/495}$ after Cauchy--Schwarz and the
raw factor $n^2$.

The common lag scale is $m(|z|)=\lfloor|z|^{-200/99}\rfloor<n/4$.
It requires no spatial mesh, and no BV reconstruction of a time-$n$
vector. The estimates are time averages or cyclic spectral statements,
not decay at each specified return count. The physical record still has
the Gaussian law in Theorem A; the Cauchy law here concerns a different,
explicitly identified spectral probability.
The Cauchy characteristic convergence is for fixed rescaled spectral
test frequency. It supplies no uniform estimate at $t=n|z|^2$
when that quantity grows across the physical annulus.

\begin{maintheorem}[Fixed-return cancellation and a larger central band]
\label{thm:intro-fixed-return-annulus}
For the insertion and measure conventions of Theorem C, let
$\Phi^{[k],a}_{n,R}(z)=\int_{Y_R^*}a((F_R^*)^kx)e^{iz\cdot J_{n,R}(x)}\dd\nu(x)$.
Uniformly for $R\in I$ and $0\le k\le n$,
\begin{align*}
 &n^2\int_{2n^{-99/200}\le|z|\le2n^{-12/25}}
                    |\Phi^{[k],a}_{n,R}(z)|\dd z\\
 &\qquad\le C\left(M_an^{-1/200}\sqrt{\log(2+n)}+V_an^{-13/100}\right).
\end{align*}
This holds at every prescribed sufficiently large return count, not
only on average. Moreover
\begin{align*}
 &\int_{|v|\le2n^{1/50}}
 |C^{[k],a}_{n,R}(v)-\alpha_a e^{-v^{\mathsf T}D_Rv/2}|\dd v\\
 &\qquad\le C\left(M_an^{-1/200}\sqrt{\log(2+n)}+V_an^{-9/175}\right).
\end{align*}
The new physical support radius is $2n^{-12/25}$; the old,
sharper estimate on $2n^{-99/200}$ remains unchanged. The unweighted
law is $a=s_R$. Theorem~\ref{thm:expanded-count-localized-inversion}
places the enlarged cutoff in the exact raw-density decomposition.
\end{maintheorem}
\begin{proof}
Apply Theorem~\ref{thm:fixed-return-inner-annulus} and
Corollary~\ref{cor:fixed-count-expanded-band}.
\end{proof}

\begin{maintheorem}[Damped unsmoothing and a wider fixed-return band]
\label{thm:intro-damped-unsmoothing}
For the insertion, measure and frequency conventions of Theorems C and L,
\begin{align*}
 &n^2\int_{2n^{-99/200}\le|z|\le2n^{-19/42}}
          |\Phi^{[k],a}_{n,R}(z)|\dd z
        \le C\left(M_an^{-5/861}+V_an^{-59/21}\right),\\
 &\int_{|v|\le2n^{1/21}}
   \left|C^{[k],a}_{n,R}(v)-\alpha_a e^{-v^{\mathsf T}D_Rv/2}\right|\dd v
        \le C\left(M_an^{-5/861}+V_an^{-9/175}\right).
\end{align*}
Both estimates hold at every sufficiently large prescribed return count,
uniformly in $R$, $0\le k\le n$ and the stated bounded-BV insertions.
The physical support radius is $2n^{-19/42}$ and its rescaled radius is
$2n^{1/21}$. More generally, a fixed-order construction reaches every
rescaled exponent $1/200<\varepsilon<1/20$ with a positive decay rate.
The prior central and inner-annulus theorems remain valid on their bands.
\end{maintheorem}
\begin{proof}
Theorems~\ref{thm:damped-cubic-unsmoothing} and
\ref{thm:wider-fixed-count-annulus},
Corollary~\ref{cor:wider-marked-central-band}, and
Proposition~\ref{prop:damped-unsmoothing-band-range} prove the assertions.
\end{proof}

The new weighted raw identity in
Theorem~\ref{thm:sharp-weighted-raw-inversion} uses this larger cutoff
and recomputes its local edge correction. Its finite-band residual,
long-time derivative budget and exact-event comparison are not replaced
by the new marked characteristic estimate.


\begin{maintheorem}[Multiscale stopping and a rate-preserving band]
\label{thm:intro-multiscale-stopping}
For the actual return record and insertion convention of Theorems C and D,
let $\mathcal G_d(D)$ denote the degree-$d$ moment tensor of $N(0,D)$.
For every fixed integer $d\ge1$, put $\rho_d=1/2$ for odd $d$ and
$\rho_d=1$ for even $d$. Then
\begin{align*}
 &\left\|\int_{Y_R^*}a((F_R^*)^kx)
       \left(\frac{U_{n,R}(x)}{\sqrt n}\right)^{\otimes d}\dd\nu(x)
                      -\alpha_a\mathcal G_d(D_R)\right\|\\
 &\hspace{15mm}\le C_d\left(M_an^{-1/4}+(M_a+V_a)n^{-\rho_d}\right).
\end{align*}
In particular, the actual normalized covariance and its induced Cesaro
formula converge to $D_R$ at rate $O(n^{-1/4})$. Moreover,
\begin{align*}
 &n^2\int_{2n^{-99/200}\le|z|\le2n^{-633/1400}}
                    |\Phi^{[k],a}_{n,R}(z)|\dd z
      \le C\left(M_an^{-3/280}+V_an^{-983/350}\right),\\
 &\int_{|v|\le2n^{67/1400}}
       |C^{[k],a}_{n,R}(v)-\alpha_a e^{-v^{\mathsf T}D_Rv/2}|\dd v\\
 &\hspace{15mm}\le C\left(M_an^{-3/280}\sqrt{\log(2+n)}
                                   +V_an^{-9/175}\right).
\end{align*}
The estimates are uniform in $R$, the prescribed sufficiently large
return count $n$, and one arbitrary mark $0\le k\le n$. The two support
radii are $2n^{-633/1400}$ physically and $2n^{67/1400}$ after rescaling.
No return-count averaging or normalization by annular volume is used.
The same-event consequence retains the original probability and variation
budget of Theorem C on this larger band.
\end{maintheorem}
\begin{proof}
Apply Theorems~\ref{thm:multiscale-marked-stopping},
\ref{thm:all-marked-gaussian-moments}, and
\ref{thm:rate-preserving-annulus}, together with
Corollary~\ref{cor:rate-preserving-same-event}.
\end{proof}

\begin{maintheorem}[Separate coordinate widths at the prescribed count]
\label{thm:intro-anisotropic-band}
Let
\[
 \mathcal B_n^{\rm roof}
   =[-2n^{1/100},2n^{1/100}]^3
                         \times[-2n^{9/100},2n^{9/100}].
\]
For the same insertion and measure convention as Theorem C,
\begin{align*}
 &n^2\int_{\substack{\sqrt n z\in\mathcal B_n^{\rm roof}\\
                            |z|\ge2n^{-99/200}}}
                  |\Phi^{[k],a}_{n,R}(z)|\dd z
       \le C\left(M_an^{-1/40}+V_an^{-497/25}\right).
\end{align*}
The estimate is uniform in $R$, one prescribed mark $0\le k\le n$,
and all sufficiently large return counts $n$. On the union
$\mathcal U_n=\mathcal B_n^{\rm roof}\cup\{|v|\le2n^{67/1400}\}$,
\begin{align*}
 &\int_{\mathcal U_n}
 |C^{[k],a}_{n,R}(v)-\alpha_a e^{-v^{\mathsf T}D_Rv/2}|\dd v\\
 &\hspace{12mm}\le C\left(M_an^{-3/280}\sqrt{\log(2+n)}
                                      +V_an^{-9/175}\right).
\end{align*}
The new physical widths are $2n^{-49/100}$ in the three discrete
coordinates and $2n^{-41/100}$ in flight time. They extend the proved
region without asserting cancellation on the whole isotropic annulus.
More generally separate rescaled exponents $\theta_i\ge1/200$ are
admissible when $\max_i\theta_i<1/10$ and
$\sum_i\theta_i+\max_i\theta_i<1/4$, with the rates specified in
Theorem~\ref{thm:anisotropic-fixed-count-band}.
\end{maintheorem}
\begin{proof}
Apply Theorem~\ref{thm:high-order-damped-unsmoothing},
Theorem~\ref{thm:anisotropic-fixed-count-band}, and
Corollary~\ref{cor:roof-box-fixed-count}.
\end{proof}

\begin{maintheorem}[A nonrectangular region at the original central rate]
\label{thm:intro-adaptive-cross}
Let $\mathcal H_n$ be the dyadic region defined in
\eqref{eq:adaptive-cross-scales}--\eqref{eq:adaptive-cross-support}.
Its base scale is $2n^{1/200}$, its maximal coordinate scale is
$n^{9/100}$, and its product-times-maximum budget is
$n^{67/280}/[\log(2+n)]^3$. Put
\[
 \mathcal U_n^\sharp=\mathcal H_n\cup\mathcal B_n^{\rm roof}
                         \cup\{|v|\le2n^{67/1400}\}.
\]
For the actual prescribed count, mark and measure convention of
Theorem C,
\begin{align*}
 \int_{\mathcal U_n^\sharp}
 |C^{[k],a}_{n,R}(v)-\alpha_a e^{-v^{\mathsf T}D_Rv/2}|\dd v
 \le C\left(M_an^{-3/280}\sqrt{\log(2+n)}+V_an^{-9/175}\right).
\end{align*}
The physical region is $n^{-1/2}\mathcal U_n^\sharp$. It retains
both the old isotropic ball and the roof box, and reaches physical
scale $n^{-41/100}$ on every coordinate axis. It is not an isotropic
ball of radius $n^{-2/5}$. Constants are uniform in $R,n,k,a$;
the residual order $P=26$ and spectral degree $Q=29$ are fixed.
\end{maintheorem}
\begin{proof}
Use Theorem~\ref{thm:adaptive-cross-fixed-count} and
Corollary~\ref{cor:adaptive-cross-directions-event}.
\end{proof}

\begin{maintheorem}[Coherent transport of the raw correction]
\label{thm:intro-coherent-raw-transport}
For an additionally finite-record admissible marked weight, fix one
exact finite-count extraction $\mu^{a,L}=E^{a,L}+Q^{a,L}$.
For a cutoff $\chi$ with kernel $K_\chi$, put
\[
 \mathcal R_\chi^{a,L}=e^{a,L}-K_\chi*E^{a,L}
            +\mathcal F^{-1}[(1-\chi)\widehat Q^{a,L}].
\]
For any two of the new dyadic, preceding box and preceding isotropic
cutoffs, and $L_n=\lceil\lambda n\rceil$ with
$\lambda>c_*^{-1}+1$, on each fixed central window,
\begin{align*}
 n^2\mathop{\rm ess\,sup}_{\mathcal W_{n,R,A}}
 |\mathcal R_\chi^{a,L_n}-\mathcal R_\psi^{a,L_n}|
 \le C\left(M_an^{-3/280}+V_an^{-983/350}\right)
        +C_{A,\lambda,P_0}M_an^{-P_0}
\end{align*}
for every fixed $P_0>0$. The difference is a bounded continuous
function even when an individual correction has infinite essential
supremum. This compares the entire signed correction, not either of
its two absolute bounds separately.
\end{maintheorem}
\begin{proof}
Proposition~\ref{prop:exact-raw-cutoff-transport} and
Theorem~\ref{thm:coherent-raw-cutoff-rate} prove the assertion.
\end{proof}

\begin{maintheorem}[Unsmoothed windows and the common signed correction]
\label{thm:intro-window-remainder}
Let $U_{n,R}=J_{n,R}-n\bar G_R$, $\rho=67/1400$, and
$B_n(\xi)=\prod_{i=1}^4[\xi_i-n^{-1/25},\xi_i+n^{-1/25}]$.
Uniformly on $|\xi|\le A$ and $R\in I$,
\[
 \nu_R^*\{U_{n,R}/\sqrt n\in B_n(\xi)\}
 =16n^{-4/25}g_{D_R}(\xi)
       \left(1+O_A(n^{-11/1400}\sqrt{\log(2+n)})\right).
\]
The event uses the unmodified record and actual integer intervals.
More generally, a rectangular frequency box contained in the retained
region gives the window estimates in
Theorem~\ref{thm:unsmoothed-window-denominator}, including positive
marked weights with their exact mass.

For a finite-record admissible complex mark, the same physical windows
$\mathcal B=n\bar G_R+\sqrt nB_n(\xi)$ satisfy
\[
 \frac{n^2}{\mathfrak m(\mathcal B)}
 \left|\int_{\mathcal B}\mathcal R_\chi^{a,L_n}\dd\mathfrak m\right|
 \le C_A\left(M_an^{-11/1400}\sqrt{\log(2+n)}+V_an^{-9/175}\right)
\]
for each of the three retained cutoffs and $L_n=\lceil\lambda n\rceil$,
$\lambda>c_*^{-1}+1$. Here $\mathfrak m$ is counting measure times
Lebesgue measure. The absolute value is outside the integral; the
assertion is not an essential-supremum bound.
\end{maintheorem}
\begin{proof}
Apply Corollary~\ref{cor:concrete-four-window} and
Theorem~\ref{thm:averaged-common-raw-remainder}.
\end{proof}

\begin{maintheorem}[Relative windows at deterministic physical time]
\label{thm:intro-relative-physical-window}
Start the orbit at a section collision with law $\nu_R^*$. Let
$V_R^Y(t)=(W_R^Y(t),C_R^Y(t)-t/\bar\tau_R)$ and
$n_R(t)=\lfloor c_*t/\bar\tau_R\rfloor$.
For $B_t=\prod_{i=1}^3[\xi_i-t^{-1/25},\xi_i+t^{-1/25}]$, define
\[
 A_t^{\rm phys}=\{V_R^Y(t)/\sqrt t\in B_t\},\qquad
 A_t^{\rm ret}=\{B_RU_{n_R(t),R}/\sqrt t\in B_t\},
\]
where $B_R(x)=(x_1,x_2,x_3-x_4/\bar\tau_R)$.
Uniformly for $|\xi|\le A$ and $R\in I$, both probabilities equal
\[
 8t^{-3/25}g_{\mathcal V_R}(\xi)
       \left(1+O_A(t^{-11/1400}\sqrt{\log(2+t)})\right),
\]
and
\[
 \frac{\nu_R^*(A_t^{\rm phys}\mathbin\triangle A_t^{\rm ret})}
      {\nu_R^*(A_t^{\rm ret})}
       \le C_A t^{-11/1400}\sqrt{\log(2+t)}.
\]
The conditional initial laws obey the same total variation bound.
These are exact physical window events of half-width $t^{23/50}$,
not singleton lattice events or events under a different initial law.
\end{maintheorem}
\begin{proof}
This is Theorem~\ref{thm:relative-physical-window}.
\end{proof}

\begin{maintheorem}[Stationary physical observations and conditional laws]
\label{thm:intro-stationary-physical-windows}
Represent equilibrium flow by the actual return suspension
\[
 \dd\mathbb P_R(y,s)=\frac{c_*}{\bar\tau_R}\dd\nu_R^*(y)\dd s,
 \qquad 0\le s<\tau_R^*(y).
\]
Let $A_t^{\rm phys}$ be the event that the stationary physical
increment $(W_R(t),C_R(t)-t/\bar\tau_R)/\sqrt t$ lies in the
three-dimensional box of center $\xi$ and half-width $t^{-2/25}$.
For $n_R(t)=\lfloor c_*t/\bar\tau_R\rfloor$, let $A_t^{\rm ret}$
be the event that $B_RU_{n_R(t),R}(y)/\sqrt t$ lies in the same box,
where $B_R(x_1,x_2,x_3,x_4)=(x_1,x_2,x_3-x_4/\bar\tau_R)$.
Then, uniformly in $R$ and bounded $\xi$,
\begin{align*}
 \mathbb P_R(A_t^{\rm phys})
 &=8t^{-6/25}g_{\mathcal V_R}(\xi)
       \left(1+O(t^{-23/2800}\sqrt{\log(2+t)})\right),\\
 \frac{\mathbb P_R(A_t^{\rm phys}\triangle A_t^{\rm ret})}
             {\mathbb P_R(A_t^{\rm ret})}
 &\le C t^{-23/2800}\sqrt{\log(2+t)}.
\end{align*}
The return event has the same denominator asymptotic. The same
comparison holds with the actual last-completed-return event, and
for the conditional laws of any common bounded path statistic.
The three physical coordinate half-widths are $t^{21/50}$.
The return origin $y$ has its equilibrium roof-length bias; it is
not drawn from the unbiased section probability.
\end{maintheorem}
\begin{proof}
Theorems~\ref{thm:wide-collision-band} and
\ref{thm:projected-collision-windows} provide the collision-clock
window estimate. Lemma~\ref{lem:stationary-length-bias} and
Proposition~\ref{prop:stationary-observation-coupling} give the
normalization and coupling. Theorem~\ref{thm:stationary-physical-windows}
proves the displayed assertions and their conditional-law consequence.
\end{proof}

\begin{maintheorem}[Physical paths conditioned on stationary windows]
\label{thm:intro-stationary-window-bridge}
Let $\mathsf X_{t,R}$ be a continuous interpolation of
\[
 s\longmapsto t^{-1/2}
   \bigl(W_R(ts),C_R(ts)-ts/\bar\tau_R\bigr),\qquad 0\le s\le1,
\]
whose supremum distance from the physical path is $O(t^{-1/2})$.
With the stationary initial law and the exact endpoint windows of
Theorem T, uniformly in $R$ and bounded $\xi$,
\[
 \operatorname{Law}(\mathsf X_{t,R}\mid A_t^{\rm phys})
       \ \Longrightarrow\
 \operatorname{Law}\bigl(s\xi+
     \mathcal V_R^{1/2}(W(s)-sW(1))\bigr)
       \quad\hbox{on }C([0,1],\R^3).
\]
Here $W$ is standard three-dimensional Brownian motion. Uniformity
is in bounded-Lipschitz distance. The same limit holds under the
prescribed-return, last-completed-return and collision reference events
on the same stationary space. For fixed interior times and bounded
boxes of nonempty interior, imposing those additional path observations
multiplies the endpoint denominator by their explicitly positive
Gaussian bridge probability; the resulting conditional path limit is
the bridge restricted to those boxes. The physical endpoint half-width
remains $t^{21/50}$.
\end{maintheorem}
\begin{proof}
Theorem~\ref{thm:multiblock-collision-band} and
Proposition~\ref{prop:window-conditional-characteristic} give joint
endpoint-window comparisons uniform in the block boundaries.
Theorem~\ref{thm:collision-window-bridge} proves conditional tightness.
Lemma~\ref{lem:whole-physical-clock-coupling},
Theorem~\ref{thm:stationary-physical-bridge} and
Corollary~\ref{cor:bridge-cylinder-events} give the physical conclusions.
\end{proof}

\paragraph{Microscopic reconstruction and conditional-measure reductions.}
Theorems V and W below are unconditional approximation theorems at the
original lattice resolution. They are separated from the Gaussian
asymptotic conclusions above: the remaining finite integral is not
identified with a Gaussian main term.

\begin{maintheorem}[Quantitative extraction and microscopic inversion]
\label{thm:intro-geometric-raw-reduction}
Use the normalized section probability and the exact finite-count
measure $\mu_{n,R}^{w,L}$. For a fixed product of bounded-BV return-state
marks, let $M,V$ be the budgets in \eqref{eq:geometric-weight-budgets}.
There is an exact decomposition $\mu_{n,R}^{w,L}=E_\varepsilon+Q_\varepsilon$
whose residual density $q_\varepsilon$ satisfies, uniformly in $R,n\le L$
and the mark locations,
\begin{align*}
 \|E_\varepsilon\|_{\TV}
    &\le C\{\varepsilon V+M(L+1)\varepsilon^{1/16}\},\\
 \sum_{\ell\in\Z^3}\|\partial_t^2q_\varepsilon(\ell,\cdot)\|_1
    &\le CM\exp\{C(L+1)\log(C/\varepsilon)\}.
\end{align*}
For $w=1$ both measures are positive. For polynomial $M,V$, any fixed
$P>0$ and any fixed upper bound on $|I|$, a roof cutoff with
$\log B_n=O(n\log n)$ gives
\[
 \sup_{R,\ell,I}n^2\left|
 \mu_{n,R}^w(\{\ell\}\times I)
       -\mathcal I_{n,R,B_n}^{w}(\ell,I)\right|\le C_Pn^{-P}.
\]
Here $\mathcal I$ is the finite Fourier integral of the full, unmodified
actual law in \eqref{eq:actual-microscopic-finite-integral}. No lattice
coordinate is averaged, and $I$ need not grow. This is a reduction of
the microscopic probability, not its Gaussian asymptotic or a pointwise
bound for the removed density.
\end{maintheorem}
\begin{proof}
Theorem~\ref{thm:geometric-extraction-budget} proves the extraction.
Proposition~\ref{prop:geometric-far-roof},
Theorem~\ref{thm:microscopic-finite-band-reduction} and
Corollary~\ref{cor:linear-cutoff-microscopic-inversion} give the inversion.
\end{proof}

\begin{maintheorem}[Positive microscopic reconstruction on the original space]
\label{thm:intro-positive-microscopic-conditioning}
Use either the normalized section law or the stationary return
suspension law \eqref{eq:equilibrium-source-for-inversion}. In the latter
case the record starts at the actual preceding section origin and all
trajectory statistics retain the original elapsed age.
For every $P>0$ there is a roof bandwidth $B_n$ with
$\log B_n=O(n\log n)$ and a nonnegative mixed density $p_{n,R,B_n}$
of mass one such that
\[
 n^2\|p_{n,R}-p_{n,R,B_n}\|_{L^1(\Z^3\times\R)}\le C_Pn^{-P}
\]
uniformly in $R$. The density is the finite integral
\eqref{eq:positive-full-law-density}; no lattice coordinate is smoothed.
For every microscopic event $A=\{J_{n,R}\in\{\ell\}\times I\}$, there
is a likelihood $0\le\lambda_{A,B_n}\le1$ on the same probability
space with
\[
 n^2\sup_{|W|\le1}
 \left|\mathbb E[W\one_A]-\mathbb E[W\lambda_{A,B_n}]\right|
        \le C_Pn^{-P}.
\]
Here $I$ is any bounded interval, including a shrinking interval, and
$W$ is any bounded measurable trajectory statistic. The second
expectation has the finite Fourier formula
\eqref{eq:all-selector-Fourier-formula}. More generally, a polynomial
number of intervals per lattice fibre is permitted. Normalization gives
a conditional total-variation bound after division by the same actual
denominator, as stated in Corollary~\ref{cor:positive-microscopic-posterior}.
\end{maintheorem}
\begin{proof}
Proposition~\ref{prop:stationary-geometric-extraction} supplies the
positive residual for both initial laws. Apply
Theorem~\ref{thm:uniform-positive-microscopic-inversion} and
Corollary~\ref{cor:positive-microscopic-posterior}.
\end{proof}

\begin{maintheorem}[Exact microscopic stationary observations]
\label{thm:intro-exact-physical-endpoint}
Fix the same convex polygonal fundamental cell at both physical
observation times. For a prescribed physical collision count $m\ge1$,
let
\[
 p_{m,R}(k,t)=\mathbb P_R\{W_R(t)=k,\ C_R(t)=m\}.
\]
The initial and terminal within-flight cell corrections enter the exact
three-frequency inverse \eqref{eq:stationary-absolute-inversion}.
Uniformly in $m,R$, its time profiles satisfy
\[
 \sum_k\|D_t^2p_{m,R}(k,\cdot)\|_{\TV}\le C,
 \qquad
 \sum_k\|p_{m,R}(k,\cdot)-p^{\rm sharp}_{m,R,B}(k,\cdot)\|_\infty
                                                      \le C/B.
\]
The derivative is a signed measure. With $\Sigma_R$ and
$\zeta_{m,R}$ defined in
Section~\ref{sec:stationary-gaussian-complement}, the central physical
radius $2m^{-99/200}$, or $2m^{1/200}$ after rescaling, gives
\begin{align*}
 m^{3/2}p_{m,R}(k,t)
 &=\bar\tau_Rg_{\Sigma_R}(\zeta_{m,R}(k,t))
       +m^{3/2}\mathfrak M_{m,R,B}(k,t)\\
 &\quad+O\left(m^{-11/700}\sqrt{\log(2+m)}+m^{3/2}/B\right),
\end{align*}
where $\mathfrak M_{m,R,B}$ is the explicitly specified signed finite
middle-frequency inverse. Thus $B_m=m^{P+3/2}$ gives a pointwise
far-frequency error $O(m^{-P})$ on the normalized physical scale,
without a discarded-density term.

For the same original physical event there is a positive likelihood
$0\le\lambda_{m,k,B}\le1$ whose unnormalized source total-variation
error is at most $C/B$, uniformly against all bounded measurable path
selectors. A relative posterior estimate requires its actual selected
denominator. The microscopic Gaussian denominator requires the signed
middle term to vanish; it is not asserted by the tail bound alone.
\end{maintheorem}
\begin{proof}
Theorems~\ref{thm:stationary-overlap-regularity},
\ref{thm:stationary-polynomial-band},
\ref{thm:physical-time-selector-reconstruction} and
\ref{thm:physical-singleton-reduction} prove the statements.
\end{proof}

\paragraph{From Gaussian laws to raw local inversion.}
The local problem concerns densities with respect to counting measure
in the lattice coordinates and Lebesgue measure in flight time, not
convolution against a fixed test function. Theorem~\ref{thm:intro-gaussian}
therefore does not replace the complementary-frequency and raw residual
estimates in Theorem~\ref{thm:LLT}. The explicit periodic geometry
and edge terms are retained as inputs to that problem. Theorem E proves
nondegeneracy by a different route: the section indicator is removed
before measurable phase rigidity is applied on positive-measure
stable--unstable product sets. No periodic representative is assumed.
The article develops the Gaussian limit, joint nondegeneracy, geometric
estimates and exact local inversion requirements for the same record.

\paragraph{Structure of the proof.}
The joint periodic arithmetic and quantitative coercivity in
Sections~\ref{sec:joint-arithmetic} and~\ref{sec:periodic-coercivity}
concern actual physical orbits. The collision compensation, Gaussian
limits, and functional time changes lead to Theorem A.
Section~\ref{sec:growing-major-arc} then tracks every frequency factor
in the smoothing and stopping errors to prove Theorem B. Its corollary
inserts that proved central band into the exact local decomposition:
positive definiteness is supplied by Theorem E, whereas the
complementary-frequency integral and local edge errors remain distinct
analytical requirements. In particular, an outer-tail
estimate beginning at $n^{-2/5}$ would not cover the intervening annulus.
Appendices~\ref{app:bv-details} and~\ref{app:functional-details} expose
the single-collision variation, smooth-density embeddings,
chronological products, and interpolation estimates. Section~\ref{sec:marked-return-band} recenters the physical orbit at a
chosen return and applies the collision splitting on both sides of a
single smooth multiplier. This proves Theorem C without assuming
mixing of the induced map, and without replacing a randomly evaluated
weight by its deterministic-clock value. Section~\ref{sec:unsmoothed-moments} then proves a finite-product
decoupling estimate by smoothing at the chosen collision gap. This
yields unsmoothed fourth moments, an $L^2$ comparison with the actual
two-sided stopping, and Theorem D.
Sections~\ref{sec:measurable-phase-rigidity} and
\ref{sec:joint-nondegeneracy} use the exact collision action one-form
and conditional pairs to exclude nonzero roof phases. Rotation products
exclude a nonintegral constant phase, and ergodicity of a finite physical
cover excludes a real spatial coboundary. This proves Theorem E for
every joint direction. Sections~\ref{sec:complete-phase-arithmetic} and
\ref{sec:quantitative-phase-defects} then complete the physical circle
arithmetic, quantify the conditional-pair argument, and reconstruct a
bounded-variation phase from an approximate complex vector. This proves
Theorem F on its stated collision-function class.
Sections~\ref{sec:near-origin-defects} and~\ref{sec:actual-return-defects}
then use diffusive damping, untwisted endpoint blocks and an explicitly
regularized actual tower to prove Theorem G. The annulus conclusion is
a function-defect bound; converting it to a complementary transform
estimate requires the stated operator comparison.
Section~\ref{sec:peripheral-return-phases} uses centered endpoint
moments and the exact spectral-phase lift to prove Theorem H.
Section~\ref{sec:compressed-return-resolvents} then supplies an explicit
reconstruction and residual comparison for finite-rank compressions,
together with a finite-time approximation bound.
Section~\ref{sec:finite-count-extraction} integrates the entire finite
physical graph and extracts all second-derivative singularities from its
time-density germs. Section~\ref{sec:count-localized-inversion} uses
the observed count coordinate to obtain the exact central identity in
Theorem I. Section~\ref{sec:exponential-reconstruction} retains the
finite binomial component bound to sharpen the tower and compressed
resolvent budgets. Section~\ref{sec:logarithmic-return-windows} then
approximates a whole multiple-time return window, applies the marked
estimates to that approximation, and removes it in the original
measure. These steps prove Theorem J and its weighted raw inversion
identity. Sections~\ref{sec:direct-orbit-bounds} and
\ref{sec:weighted-spectral-averages} record the incompatible grid
budgets and replace their use, for averaged coefficients, by a finite
Gram-matrix argument on the exact orbit. They prove Theorem K and its
weighted, peripheral and spectral consequences. The raw inversion
routes retain the same mixed-density endpoint.

Sections~\ref{sec:finite-order-damping} and
\ref{sec:fixed-return-annulus} prove Theorem L by fixed-order connected
collision correlations and a finite spectral Taylor polynomial. The
coefficients, but not the analytic radius, are uniform in the smoothing
scale. Real collision damping, unsmoothing, and the exact two-sided
stopping yield a fixed-count inner-annulus integral on the raw scale.
The enlarged-cutoff inversion identity still displays the farther
residual integral, derivative budget, and new local edge correction.
Sections~\ref{sec:damped-unsmoothing} and~\ref{sec:wider-fixed-count-band}
then retain the first three residual-phase corrections in a damped
chronological word, with a fine insertion scale separate from the
coarse twisted-observable scale. A small-mass cumulant bound controls
the fourth-order remainder. Fixed higher moments improve the actual
clock window and give Theorem M. The same unchanged marked event and
the common BV/finite-record subanalytic class pass to the new central
and exact raw identities, with their remaining weighted terms explicit.


Sections~\ref{sec:multiscale-stopping} and~\ref{sec:rate-preserving-band}
replace a single stopping window by disjoint dyadic deviation shells.
A deterministic maximal estimate and the actual clock tails yield a
fourth-root comparison in every fixed finite moment. Anchored mixed
cumulants then identify every fixed marked Gaussian moment. Combining
the same stopping estimate with damped cubic unsmoothing proves Theorem N:
the wider band retains the original central rate. Its weighted raw
identity uses a newly defined kernel and keeps the remaining local edge,
complementary residual, and long-time derivative terms explicit.

Sections~\ref{sec:high-order-damped-unsmoothing} and
\ref{sec:anisotropic-fixed-count-bands} separate the residual order from
the spectral degree, then balance the four coordinate widths with the
raw four-dimensional volume. They prove Theorem O and the matching
kernel identity. Sections~\ref{sec:shape-adaptive-fixed-count} and
\ref{sec:coherent-raw-cutoff} use one fixed-order weighted-volume
estimate to obtain Theorem P, then retain the exact edge/residual
cancellation to prove Theorem Q. Thus enlarging the known Fourier
region does not introduce an unrelated signed raw remainder.
Appendix~\ref{app:dependency-guide} retains the inherited dependency
chain. The new signed comparison does not supply an isolated
absolute edge estimate.

\paragraph{A short verification roadmap.}
Theorems A, D and N supply actual centered moments and deterministic
return-window maximal estimates; Theorem E supplies the same uniformly
positive covariance throughout. The prescribed-count estimates of
Theorems N--P provide the integrated characteristic error, not a
frequency-plane trace. Theorem Q identifies the common signed raw
correction. Section~\ref{sec:unsmoothed-window-remainders} combines
positive interval envelopes with those prescribed-count estimates to
prove Theorem R, keeping lattice endpoints explicit. Section~\ref{sec:relative-physical-windows}
uses a four-dimensional truncation to obtain projected denominators,
then a return-window maximum and an inverse clock to prove Theorem S
at deterministic physical time. The remaining microscopic step is a
bound for the common correction in the required pointwise norm and
for the original exact-label events; it is not a new cutoff transition.

\paragraph{Roadmap for the local and conditional conclusions.}
The collision estimates and actual stopping give the Gaussian,
functional and moment theorems; measurable phase rigidity gives
Theorem E. Finite-count extraction makes the exact raw decomposition
available, while fixed-order damped unsmoothing yields the
prescribed-count domains of Theorems N--P. Theorem Q identifies one
common signed raw correction across those cutoffs. Theorem R controls
its signed window average and exact window denominators; Theorem S
uses these for section-start physical observations. Theorem T instead
uses the wider deterministic-collision band of
Section~\ref{sec:stationary-collision-band}, followed by the exact
roof-length bias, to reach finer stationary physical windows.
Theorem U retains joint path frequencies while integrating the endpoint
frequency. Uniformity at coalescing block boundaries gives conditional
tightness after division by the rare denominator; a whole-clock
comparison then identifies the physical Gaussian bridge and its
interior path-selection probabilities.
The remaining microscopic step is the pointwise common correction
or an equivalent complete raw inversion bound. Neither the collision
band nor a signed average is substituted for that step.

\paragraph{Logical roles of the main results.}
The following diagram separates asymptotic conclusions from exact
reconstruction. Every arrow denotes a proved input implication, not
completion of the raw local limit.
\[
\begin{array}{c}
\text{collision covariance, moments and phase rigidity}\\
\Downarrow\\
\text{actual Gaussian laws and prescribed-count central bands}\\[2mm]
\text{source geometry}\ \Longrightarrow\
\text{positive residual with a long-time derivative budget}\\
\Downarrow\\
\text{finite-band microscopic measures and conditional likelihoods}
\end{array}
\]
\begin{center}
\begin{tabular}{p{0.27\textwidth}p{0.65\textwidth}}
\hline
Part & Proved conclusion and norm\\
\hline
Gaussian and phase theory & Theorems A--H: actual weak and functional
limits, moment identification, uniform nondegeneracy and specified
phase-defect estimates.\\
Raw analysis & Theorems I--R and V: finite extraction, specified
Fourier regions, coherent correction and signed window estimates;
Theorem V is a microscopic finite-band reduction.\\
Physical windows & Theorems S--U: explicit mesoscopic denominators,
relative event comparison and conditional Gaussian bridges.
Theorems T and U retain half-width $t^{21/50}$.\\
Positive reconstruction & Theorem W: global mixed $L^1$ approximation
and a total-variation estimate for microscopic event measures on the
original source, uniform in bounded path selectors.\\
Remaining raw endpoint & The pointwise common correction and full
finite complement, followed by a microscopic Gaussian denominator
and the original physical-event comparison at that scale.\\
\hline
\end{tabular}
\end{center}
Equation~\eqref{eq:geometric-common-raw-ledger} is the detailed pointwise
ledger. Global $L^1$ reconstruction in Theorem W does not replace its
local essential-supremum requirement. The same distinction applies
under the stationary entrance law, whose bias is handled in
Proposition~\ref{prop:stationary-geometric-extraction}.

\paragraph{The original stationary endpoint.}
Theorem X concerns the exact physical singleton, not the preceding-section
return record used by Theorem W. The flight-age overlap supplies a
uniform pointwise tail estimate without a geometric removal. Its
Gaussian central contribution is evaluated directly in three frequency
variables. Section~\ref{sec:stationary-gaussian-complement} gives the
remaining signed middle integral and an exact local equivalence; it
does not identify this physical inverse with the four-dimensional
return-density correction.

\paragraph{Relation to the limit theory of billiards.}
The existence of Gaussian and functional laws for finite-horizon
billiards is not in itself new: tower and invariance-principle methods
already give such laws for collision maps and flows
\cite{Young,MN}. Perturbative collision-space methods give stability of
spectral data under geometric changes \cite{DZ}. Local limit theorems
for cell displacement, including randomly deformed tables, provide a
further comparison \cite{SV,DPZ}. Our problem instead keeps the lattice
displacement, collision count, and flight time of a genuine moving-section
first return in one four-coordinate record. The compensation identity
allows quantitative, radius-uniform central integrals for this
unbounded induced record using bounded collision observables. Its role
here is to connect that central analysis to the arithmetic and exact
edge subtraction of the raw mixed-density problem. None of the cited
cell-displacement theorems is invoked as a theorem for this joint raw
record, and the central-band result is not identified with a complete
local limit theorem.

For the completed parts, the closest comparisons have different
observables or norms. Sz\'asz--Varj\'u \cite{SV} prove a local central
limit theorem for Young systems and apply it to the planar Lorentz
process. Dolgopyat--N\'andori \cite{DN} formulate suspension-flow local
limit conditions and verify them for several classes, including
finite-horizon Sinai billiards. Baladi--Demers--Liverani \cite{BDL}
prove exponential correlation decay for the finite-horizon Sinai flow
through anisotropic distribution spaces. These results establish the
relevant background for Gaussian and local statistics; none is cited
here as a direct theorem for the joint count, displacement and roof
record of this moving-section first return.

The present completed conclusions retain that record through exact
collision compensation, quantify parameter and insertion budgets, and
identify the stationary conditional path at the specified shrinking
window scale. The source-side extraction has a different purpose: it
gives an explicit all-label derivative budget at linear collision
cutoff. The positive reconstruction then controls microscopic event
measures uniformly against arbitrary bounded trajectory statistics
without imposing a multiplier norm on those statistics. Its conclusion
is finite-band reconstruction, not an additional mixing or Gaussian
local-limit theorem. This distinction locates the new argument relative
to the existing local-limit and billiard-flow literature.

\section{The mechanical question and its interfaces}
Let $a=(1,0)$, $b=(1/2,\sqrt3/2)$, $\Lambda=\Z a+\Z b$, and
\[
 I=[9/20,47/100],\qquad Q_R=(\R^2/\Lambda)\setminus\overline B(0,R).
\]
The speed is one and every reflection is specular. Three indices will be kept distinct: a physical collision count $N$, an induced return count $n$, and the number of experimental bits. An induced record has values in $\Z^2\times\Z\times\R$, with counting measure in the first three coordinates and Lebesgue measure in the last. A four-dimensional Lebesgue density would be the wrong object.

The geometric inverse of the preceding article identifies the table from two fixed collision fields \cite{A2G}. It supplies neither independent flights nor an induced transfer-operator estimate. Conversely, the stationary law of the equilibrium billiard is not the unknown preparation law of that inverse experiment. The distinction permits geometric estimates to be used for equilibrium statistics without first estimating the entire preparation law.

The joint arithmetic result is Theorem~\ref{thm:joint-periodic}. It concerns
the complete periodic phase of an explicitly enlarged physical first-return
section. Lemma~\ref{lem:excursion-family} constructs its infinite period
family, and Theorem~\ref{thm:period-excess} proves the strict excess
inequalities. These results go beyond the zero-roof determinant by ruling
out every nonzero roof frequency, without a numerical irrationality claim.
The real-rank and compact-frequency consequences are stated separately
from measurable cohomology regularity and operator contraction.

The original physical record, winding geometry, critical-edge and
conditional inversion arguments remain in the article. Theorem
\ref{thm:local-edge-LLT} refines the inversion requirement by using local
edge errors instead of small total edge mass. Propositions
\ref{prop:exact-induced-means} and \ref{prop:induced-suspension} give exact
normalizations and stationary endpoint control for the enlarged section.
The full raw local limit is not inferred from the arithmetic theorem:
its spectral, residual and weighted-insertion requirements are listed in
the final section. All probability laws here refer to actual specular
orbits, not independent flight surrogates.

The relevant prior local-limit theory includes the Lorentz-process theorem of Sz\'asz--Varj\'u \cite{SV} and the suspension framework of Dolgopyat--N\'andori \cite{DN}. Definition 2.2 and Proposition 6.3 of the latter concern a mixing local limit tested against fixed compactly supported functions under minimality and group assumptions. They do not state absolute integrability of every raw finite-time characteristic function, or the particular geometrically uniform density theorem sought here. Demers--Zhang \cite{DZ} give a perturbation framework for Lorentz-gas maps; its spectral continuity is not automatically twice differentiability in a moving geometric parameter. We use the corrected version of the flow-mixing work of Baladi--Demers--Liverani \cite{BDL}. No novelty is claimed for Fourier inversion, the Morse lemma, or the suspension identity as general methods.


Theorem~\ref{thm:return-stability} compares the actual moving return
events in density $L^1$ and first moment, at every fixed return count.
Corollary~\ref{cor:uniform-stationary-endpoints} then supplies uniform
stationary endpoint control over the whole radius interval. The proof uses
finite-itinerary stability, coarea and exact Kac means, not a hypothetical
common spectral space. These finite-record moduli do not purport to supply
the quantitative long-time estimates in the local-limit interface.

Theorem~\ref{thm:quantitative-excess} strengthens strict increase to
uniform two-sided exponential bounds, obtained from the exact nonlinear
half action. Theorem~\ref{thm:log-period-separation} uses their upper and
lower bounds together: the former selects periods of logarithmic length,
and the latter supplies a polynomial discrepancy rather than mere
nonvanishing. Corollary~\ref{cor:approximate-phase-bound} states precisely
the continuous-phase consequence. It does not assume the measurable
regularity needed for an operator application.

At the level of the clock, Theorems~\ref{thm:uniform-return-fluid} and
\ref{thm:uniform-physical-clock} turn fixed-record continuity and exact
means into a uniform functional first-order law. This includes the maximum
unfinished return over the observation window at scale $t$, whereas the
retained stationary endpoint result is at scale $\sqrt t$ for finitely
many deterministic endpoints. The distinction of scales is essential.
Corollary~\ref{cor:physical-rate-error} gives a geometric error term for
the physical rate without assuming independence of the radius estimate
and the observed orbit. No originality is claimed for the general ergodic,
subadditive or renewal inversion arguments; the point is their verified
application to the actual moving return family.

\section{Physical records and an exact marked identity}
Write an outgoing section point as
\[
 q=R n_\alpha,\qquad v=\sqrt{1-p^2}\,n_\alpha+p\,t_\alpha,
 \qquad (\alpha,p)\in\T\times(-1,1).
\]
The collision map is $T_R$, its next physical flight is $\tau_R$, and
\begin{equation}\label{eq:flux}
 \dd\nu=\frac{\dd\alpha\dd p}{4\pi},\qquad
 \bar\tau_R=\frac{\sqrt3/2-\pi R^2}{2R},\qquad
 \dd\mu_R=\bar\tau_R^{-1}\dd\nu\dd s\quad(0\le s<\tau_R).
\end{equation}
These formulas follow by integrating the boundary flux $\cos\phi\dd s_{\partial Q}\dd\phi$: total section flux is $4\pi R$ and phase volume is $2\pi(\sqrt3/2-\pi R^2)$. Reflection preserves this flux. Thus $T_R$ preserves $\nu$.

Each next disk has a lattice label $d\ne0$. With $D=d-q$ and $L=D\cdot v$, its candidate entry time is
\begin{equation}\label{eq:entry}
 \tau_d=L-\sqrt{R^2-|D|^2+L^2}.
\end{equation}
The physical branch is the positive entry root with the smallest time, not an arbitrarily selected root. Its endpoint normal is $n'=(q+\tau_dv-d)/R$, and its outgoing velocity is $v'=v-2(v\cdot n')n'$. These expressions define the label, point, both velocities, and time on every regular branch. On a compact finite itinerary separated from tangencies and competing roots they depend analytically on $R$ and the initial coordinates, by the implicit function theorem. This is a finite-itinerary assertion, not a common norm for all iterates.

Distinct disks have distance at least $g_R=1-2R\ge3/50$. A straight outgoing flight cannot return to its initial convex disk. In particular there is no finite-time collision accumulation. This family has uniform finite horizon: a rational lattice direction has row spacing at most $\sqrt3/2$, so every line in that direction meets an open disk when $R>\sqrt3/4$; an irrational line is dense on the torus. Compactness then excludes arbitrarily long free segments at the smallest radius. Increasing the radius preserves this upper bound. Tangencies and their finite preimages have section measure zero.

Put $S_j\tau(y)=\sum_{i=0}^{j-1}\tau(T_R^iy)$. For $y$ interpreted as the first future outgoing collision and $u$ as its arrival time, the first $k$ future collisions have times $u+S_j\tau(y)$, $0\le j<k$. Their points, lattice labels, and incoming and outgoing velocities are obtained by iterating \eqref{eq:entry}. The final incomplete flight has duration $t-u-S_{k-1}\tau(y)$.

\begin{proposition}[Marked stationary record]\label{prop:marked}
For every bounded measurable functional $\Psi$ of this record and its terminal flight,
\begin{align}\label{eq:marked}
 \mathbb E_{\mu_R}[\Psi;N_t=k]
 =\frac1{\bar\tau_R}\int\dd\nu(y)\int_0^{\tau(T_R^{-1}y)}
 &\one_{\{u+S_{k-1}\tau(y)\le t<u+S_k\tau(y)\}}\nonumber\\[-2mm]
 &\hspace{-8mm}\cdot\Psi(\mathcal R_{k,t}(y,u))\dd u,\qquad k\ge1.
\end{align}
The initial phase point is $\Phi_R^{-u}y$, so the entire initial residual flight is retained as well. For $t=11/100$ there are at most two future collisions, and \eqref{eq:marked} with $k=1,2$, together with the zero-collision case, gives the full marked law without an independence assumption.
\end{proposition}
\begin{proof}
Under \eqref{eq:flux}, set $y=T_Rx$ and $u=\tau(x)-s$. Invariance of $\nu$ gives the displayed domain for $u$ with Jacobian one. The $k$th and $(k+1)$st future collision times are exactly the two endpoints in the indicator. This proves the identity for arbitrary bounded measurable $\Psi$, first for indicators and then by integration. The event of equality at an endpoint has zero suspension measure. Three collisions require at least two complete intercollision flights after the first event; their total duration is at least $6/50>11/100$.
\end{proof}


\section{Explicit periodic geometry and a genuine inducing set}
Two elementary cycles have zero winding. The normal nearest-neighbor orbit has two flights, total length
\begin{equation}\label{eq:two-three}
 L_2=2(1-2R),\qquad
 L_3=3(1-\sqrt3 R)
\end{equation}
is the length of the inner equilateral three-cycle on centers $0,a,b$. Its contact angles are $\pi/6,5\pi/6,3\pi/2$ and its outgoing $p$ coordinates are all $-1/2$. The reflection cosines are respectively $1$ and $\sqrt3/2$. The equilateral segments lie inside their elementary triangle and outside the three vertex disks; any other lattice center is at distance at least $\sqrt3/2$ from that triangle. Hence all other disks are strictly avoided.

\begin{lemma}[A winding four-cycle]\label{lem:winding}
For every $R\in I$ there is an analytic angle $\theta_R\in(0.92,0.967)$ satisfying
\begin{equation}\label{eq:theta}
 f_R(\theta)=2R\sin\theta-\sin\theta\cos\theta
                       +\frac{\sqrt3}{2}\cos(2\theta)=0.
\end{equation}
The itinerary with centers $0,a,a-b,2a-b,a$ and normals at angles
\[
 -\theta_R,\quad\pi+\theta_R,\quad\theta_R,\quad\pi-\theta_R
\]
is a physical four-collision periodic orbit modulo $\Lambda$, with winding $a$. All its incidence cosines exceed $1/2$ and its clearance from every third disk exceeds $1/200$. Rotation through $\pi/3$ gives another cycle with winding $b$ and the same length.
\end{lemma}
\begin{proof}
Put $c=\cos\theta$, $s=\sin\theta$. Its first four contact points and the translated fifth point are
\begin{align*}
 q_0&=(Rc,-Rs),&q_1&=(1-Rc,-Rs),\\
 q_2&=(1/2+Rc,-\sqrt3/2+Rs),&q_3&=(3/2-Rc,-\sqrt3/2+Rs),\\
 q_4&=q_0+a.
\end{align*}
Two segments are horizontal. Write $d=2Rc-1/2$, $h=\sqrt3/2-2Rs$. The other two are $(d,-h)$ and $(d,h)$. The scalar equation $d\sin2\theta+h\cos2\theta=0$ is precisely \eqref{eq:theta}. Since $d,h>0$, it makes the diagonal direction the reflection of the horizontal direction at the intervening normal. The remaining reflection equations follow by symmetry, with cosine $c$ at every contact. The total length is
\begin{equation}\label{eq:fourlength}
 L_4=2(1-2Rc)+2\sqrt{d^2+h^2}.
\end{equation}

Here are reproducible rational bounds for the inequalities and absence of occlusion. Divide $I$ into 32 closed intervals of width $1/1600$. On each, bisect the endpoint equations \eqref{eq:theta} between $9/10$ and $1$ to width at most $10^{-10}$. The interval derivative
\[
 2R\cos\theta-\cos2\theta-\sqrt3\sin2\theta
\]
is less than $-1/2$. The endpoint signs and positivity of $\partial_R f$ show that the two endpoint root brackets enclose the root for the entire parameter interval. The resulting angles lie in $(0.92,0.967)$; $d,h>0$ and $c>1/2$ on every interval.

For clarity, the arithmetic of this finite certificate uses no floating root solver. Square roots are enclosed by integer square roots at denominator $10^{22}$. For sine and cosine use the Taylor polynomials through degrees 25 and 24 at the rational midpoint, with errors at most $|x|^{26}/26!$ and $|x|^{25}/25!$, respectively; add the interval radius, since both functions are 1-Lipschitz. All interval operations are rational.

For each segment check centers $ia+jb$, $|i|,|j|\le4$, excluding its two endpoints. The distance from the rational midpoint segment is computed by clamping the rational orthogonal projection to $[0,1]$. Subtract the endpoint and center enclosure radii and the largest obstacle radius. Each lower bound exceeds $1/200$. Every contact point has norm less than 3; a center outside the index box has norm at least $5\sqrt3/2$, so it too has distance greater than the obstacle radius from every segment. The complete deterministic calculation is specified in \texttt{tools/certify\_winding.py}; it obtains lower bounds greater than $0.00925$ for clearance and $0.569$ for incidence, and a derivative upper bound below $-0.731$. The signs and strict margins prove the assertions on the full intervals, not only at sampled radii. The implicit function theorem gives analytic dependence and uniqueness in the printed root interval.
\end{proof}

Choose $\epsilon=1/200$. In $(\alpha,p)$ coordinates let $Y_R$ be the union of the open squares of half-side $\epsilon$ centered at
\begin{equation}\label{eq:Y}
 (0,0),\quad(\pi/6,-1/2),\quad
 (-\theta_R,\sin\theta_R),\quad
 (\pi/3-\theta_R,\sin\theta_R).
\end{equation}
Angles are taken modulo $2\pi$. The squares are disjoint and
$\nu(Y_R)=4\epsilon^2/\pi$. Let $r_R$ be the first return time to $Y_R$ and $F_R=T_R^{r_R}$ on its recurrent full-measure domain, with invariant probability $\nu_R^Y=\nu|_{Y_R}/\nu(Y_R)$. Define
\begin{equation}\label{eq:induced}
 \kappa_R=\sum_{j<r_R}\kappa_{\rm coll}\circ T_R^j,
 \qquad \tau_R^Y=\sum_{j<r_R}\tau_R\circ T_R^j.
\end{equation}
Poincar\'e recurrence and invariance give this induced probability; they do not give an exponential return tail, a Markov partition, or a differentiable family of spectral projectors.

\begin{proposition}[Full augmented lattice at zero roof frequency]\label{prop:lattice}
The four cycles above give fixed points of $F_R$ with data $(\kappa_1,\kappa_2,r,\text{return period})$ equal to the rows of
\begin{equation}\label{eq:matrix}
 V=\begin{pmatrix}0&0&2&1\\0&0&3&1\\1&0&4&1\\0&1&4&1\end{pmatrix},
                  \qquad |\det V|=1.
\end{equation}
Consequently a circle-valued coboundary identity
$e^{i(u\cdot\kappa+s r)}=e^{ic}g/(g\circ F_R)$ that can be evaluated at these fixed points forces $u=s=c=0$ modulo $2\pi$.
\end{proposition}
\begin{proof}
Only one contact of each cycle belongs to \eqref{eq:Y}. For the four-cycle the successive $p$ coordinates are $s,s,-s,-s$. The same is true after rotation. The two positive-$p$ unselected contacts are separated from the selected centers; the smallest relevant angular separation is $2\pi/3-2\theta_R>0.16$. All negative-$p$ contacts are separated in $p$ from the selected winding contacts, and also from $0,-1/2$. The remaining contacts of the two- and three-cycles are separated by their displayed angles. Thus the physical first return counts are $2,3,4,4$, not freely assigned labels. The four selected contacts are interior to their squares, so the same finite return branches exist in neighborhoods. Direct subtraction of the first two rows gives determinant of magnitude one. Evaluation first gives $2s=c=3s$, hence $s=c=0$; the last two rows then give $u_1=u_2=0$.
\end{proof}

The qualification about evaluation is essential. A merely measurable transfer function is not automatically defined at a periodic point. A spectral application requires its own regularity or Liv\v{s}ic argument. Moreover, adding a roof twist $b\tau_R^Y$ to this proposition leaves $bL_j$ in every periodic equation. The four rows do not eliminate it. Joint roof nonarithmeticity and returned temporal nonintegrability remain distinct requirements. As one smaller unconditional conclusion, \eqref{eq:two-three} excludes a continuous cohomology $\tau_R=c+g-g\circ T_R$: the two periodic averages are $1-2R$ and $1-\sqrt3R$, which differ by $(2-\sqrt3)R>0$.


\section{A physical period family and the complete periodic annihilator}
\label{sec:joint-arithmetic}
The four records in Proposition~\ref{prop:lattice} settle the integer
coordinates at zero roof frequency. We now add a family of actual periodic
trajectories that settles every roof frequency. A slightly larger section is
specified explicitly; the original section and every assertion about it
remain as stated in the preceding sections.

\subsection{The enlarged section}
Put
\begin{equation}\label{eq:enlarged-section}
 U=\{(\alpha,p):|\alpha-\pi/6|<3/50,\ |p|<3/20\},
 \qquad Y_R^*=Y_R\cup U,\qquad c_*:=\nu(Y_R^*)=\frac{91}{10000\pi}.
\end{equation}
The rectangle $U$ is disjoint from the four squares defining $Y_R$.
At the square with the same angular center the $p$ coordinate is near
$-1/2$; the other three squares are separated in angle or in $p$.
The area of $U$ is $4(3/50)(3/20)$, and the section density is
$1/(4\pi)$, proving the value of $c_*$. Write
$F_R^*=T_R^{r_R^*}$ for the first return to $Y_R^*$, and let
$\kappa_R^*$ and $\tau_R^*$ be the sums of the physical displacement
and flight length before that return. Its invariant probability is
$\nu_R^*=\nu|_{Y_R^*}/c_*$. An assertion about $F_R$ does not become
an assertion about $F_R^*$ merely by changing notation.

Rotate coordinates through $-\pi/6$ and put
\begin{gather*}
 D=\sqrt3,\qquad O=(0,0),\quad C=(D,0),\quad
 A=(D/2,-1/2),\quad B=(D/2,1/2),\\
 d=\frac12-R,\qquad g=D-2R.
\end{gather*}
The rotated lattice consists of $(Dk/2,l/2)$ with integers $k,l$ of
the same parity. The normal orbit between $O$ and $C$ is unobstructed:
the two intermediate centers $A,B$ have distance $1/2$ from its
supporting line, and all other centers are farther away. It has two
physical flights of total length $2g$, zero winding, and exactly one
visit to $Y_R^*$ per period. Its selected state is $(\pi/6,0)$ in the
unrotated coordinates. The four cycles of Proposition~\ref{prop:lattice}
still have exactly one visit per period. Indeed, the equilateral cycle
has $p=-1/2$, the two-cycle has angles $0,\pi$, and the winding cycles
have $|p|=\sin\theta_R>0.79$. None gains a visit to $U$.

\subsection{A convex optical action}
For $-d\le y,z\le0$ define
\begin{align}
 x_R(y)&=\sqrt{R^2-y^2},\qquad u_R(y)=R-x_R(y),\nonumber\\
 e_R(y)&=\sqrt{(D/2-x_R(y))^2+(d+y)^2},\label{eq:excursion-lengths}\\
 \ell_R(y,z)&=\sqrt{(D-x_R(y)-x_R(z))^2+(z-y)^2}.\nonumber
\end{align}
The function $e_R$ is the length from the top contact
$q_A=(D/2,-d)$ of the disk at $A$ to a contact at height $y$ on the
facing arc of the disk at $O$ or $C$. The function $\ell_R$ is a
flight length between the two facing arcs. For $m\ge1$ let
\begin{equation}\label{eq:excursion-action}
 \mathcal A_{m,R}(y_1,\ldots,y_{2m})
 =e_R(y_1)+\sum_{j=1}^{2m-1}\ell_R(y_j,y_{j+1})+e_R(y_{2m}),
 \qquad L_m(R)=\min_{[-d,0]^{2m}}\mathcal A_{m,R}.
\end{equation}

\begin{lemma}[A uniformly regular physical family]\label{lem:excursion-family}
For each $m\ge1$ and $R\in I$, the minimum in
\eqref{eq:excursion-action} is unique. Its coordinates satisfy
\begin{equation}\label{eq:excursion-height}
 -\frac{2d}{5}<y_j<0,\qquad y_j=y_{2m+1-j}.
\end{equation}
The contacts with these heights form a genuine specular periodic orbit
with center word $A,(O,C)^m,A$. Every flight avoids all third disks.
The orbit has zero winding, $2m+1$ physical collisions, and exactly
$m$ returns to $Y_R^*$. The period length $L_m(R)$ is analytic in $R$
for each fixed $m$.
\end{lemma}
\begin{proof}
We supply the variational, mechanical and return checks separately.
On the whole box, $x_R(y)>11/25$, $x_R(y)\le47/100$, and
$1.73<D<1.74$. The positive horizontal quantities
\[
 a(y)=D/2-x_R(y),\qquad h(y,z)=D-x_R(y)-x_R(z)
\]
satisfy $a>79/200$ and $h>79/100$. The second derivative of
$-x_R(y)$ is $R^2/x_R(y)^3>0$. The Hessian formula for a Euclidean
norm shows that the Hessian of $e_R$ or $\ell_R$ is the sum of a
positive semidefinite term and the Hessian of its horizontal
component multiplied by its positive horizontal direction cosine.
The latter cosine exceeds $9/10$ on this box. Thus each segment
length is strictly convex in its variables; in fact the Hessian of
$\mathcal A_{m,R}$ dominates the identity. One may check the last
numerical bound from
$(9/10)(9/20)^2/(47/100)^3>1$. Compactness and strict convexity give
a unique minimizer.

No coordinate equals $-d$. For such a coordinate $y$ and an incident
interior segment, the derivative is
\[
 \partial_y\ell_R(y,z)=\frac{h(y,z)y/x_R(y)+y-z}{\ell_R(y,z)}<0,
\]
since $z\ge y=-d$. For an end segment,
$e_R'(-d)=-a(-d)d/(x_R(-d)e_R(-d))<0$.
Increasing that coordinate therefore decreases the action.
At an end coordinate equal to zero, $e_R'(0)=d/e_R(0)>0$, and
all other incident derivatives are nonnegative. At an interior zero
adjacent to a negative coordinate, its derivative is positive.
If a block of zero coordinates had no such neighbor it would reach
an endpoint. A variation into the box excludes all these cases.
Every coordinate is consequently in $(-d,0)$. Reflection in the
line $x=D/2$ reverses the center word and preserves the action, so
uniqueness gives the symmetry in \eqref{eq:excursion-height}.

Here is a uniform improvement of the height bound. A smallest
coordinate $y=-w$ cannot be interior, because both incident
derivatives there are negative. Take it to be $y_1$, using reversal
if needed, and write $z=y_2$. Stationarity gives
\[
 d-w=\left(a+e_R(y)\frac{h}{\ell_R(y,z)}\right)\frac{w}{x_R(y)}
           +\frac{e_R(y)}{\ell_R(y,z)}(w+z).
\]
The last term is nonnegative since $z\ge-w$. Moreover,
\[
 \frac{a+e_R(y)h/\ell_R(y,z)}{x_R(y)}
 >\frac{(19/10)(79/200)}{47/100}=\frac{1501}{940}>\frac32.
\]
It follows that $w<2d/5$. This bounds every coordinate.

Interior stationarity of the optical length is precisely equality
of tangential incoming and outgoing velocities at each contact.
The outgoing directions point into the exterior normal half-plane.
For example, at a contact on $O$, its outward normal is
$(x_R(y),y)/R$. Its dot product with a direction to $C$ has numerator
at least $x_R(y)h-|y||z-y|>0$; for a direction to $q_A$ both the
horizontal contribution and $y(-d-y)$ are positive. The checks at
$C$ are reflections of these. The symmetry of the two end contacts
implies the reflection law at $q_A$, whose normal is vertical. There
$e_R(y_1)<11/25$ and $d+y_1>3d/5\ge9/500$, so the normal
incidence cosine exceeds $9/220>1/25$. The other contacts are
uniformly transverse by the preceding horizontal-component bounds. Thus the stationary polygon
is specular rather than the straight-through solution of the
stationarity equation.

The long flights lie in $-2d/5<y<0$, so their distance from the disk
at $A$ is at least $3d/5\ge9/500$, and from the disk at $B$ at
least $d\ge3/100$. End flights join $q_A$ to the appropriate facing
arc and leave both endpoint disks in their exterior supporting
half-planes. Their third disk $O$ or $C$ has horizontal separation
at least $D/2>R$, and $B$ is vertically separated. All segments
lie in the rectangle
\[
       11/25<x<D-11/25,\qquad -d\le y\le0.
\]
For the remaining lattice centers, the possibilities $k=0,2$ with
nonzero $l$ have vertical separation at least $1-d\ge0.95$;
$k=1$, $|l|\ge3$ have separation at least $3/2-d$; and
$k\notin\{0,1,2\}$ have horizontal separation exceeding $1.3$.
This enumerates the whole lattice and proves nonocclusion, not just
a check of a finite list of nearby obstacles.

At an outgoing contact on $O$, the normal angle differs from
$\pi/6$ by at most $|y|/x_R(y)<1/22<3/50$. The outgoing velocity
points toward $C$ and its angular deviation from the horizontal is
at most $|z-y|/h<2/79$. Thus
\[
 |p|\le \frac1{22}+\frac2{79}<\frac3{20}.
\]
All $m$ contacts on $O$ belong to $U$. Contacts on $C$ have angles
near $7\pi/6$, and $q_A$ has angle $2\pi/3$; none belongs to
$Y_R^*$. The return counts are $2$ between successive contacts on
$O$, with one count $3$ through $C,A,O$. Their sum is $2m+1$,
and there are $m$ induced returns. The orbit is closed in the plane,
so its winding is zero. All selected contacts and unselected contacts
have strict membership margins, giving actual regular return branches.
Finally the positive Hessian and the interior minimum allow the
analytic implicit function theorem. This proves analyticity for each
fixed $m$; no bound for high parameter derivatives uniform in $m$ is
being inferred.
\end{proof}

\subsection{Strictly increasing bounded excess lengths}
The following argument supplies an infinite sequence of period relations.
It does not infer nonarithmeticity from finitely many decimal lengths.

\begin{theorem}[Strict period excess]\label{thm:period-excess}
For the family of Lemma~\ref{lem:excursion-family}, define
$E_m(R)=L_m(R)-2m g$. Then, for every $R\in I$,
\begin{equation}\label{eq:period-excess}
 0<E_m(R)<E_{m+1}(R)\le \frac{2d^2}{g}<\frac1{100}.
\end{equation}
In particular $E_{m+1}(R)-E_m(R)>0$ and these differences tend to
zero. The convergence of $E_m$ is used pointwise in $R$; the subsequent
compact-frequency assertion supplies its own uniformity argument.
\end{theorem}
\begin{proof}
Suppress $R$ and set $K(y,z)=\ell(y,z)-g$ on $[-d,0]^2$.
The horizontal component of the flight gives
\begin{equation}\label{eq:kernel-coercivity}
 K(y,z)\ge u(y)+u(z),
\end{equation}
with equality exactly when $y=z$. Define continuous value functions
\[
 V_0(y)=u(y),\qquad
 V_{j+1}(y)=\min_{z\in[-d,0]}\{K(y,z)+V_j(z)\}.
\]
They are nonnegative and vanish at zero. For $j\ge1$, choosing
$z=0$ gives $V_j(y)\le K(y,0)=O(y^2)$ near zero; $V_0$ has the
same bound. The implied bound is independent of $j$.
Equation~\eqref{eq:kernel-coercivity} implies $V_1(y)>V_0(y)$ for
$y<0$: equality at a minimizing $z$ would require both $z=0$ and
equality in \eqref{eq:kernel-coercivity}, hence $y=0$.

For any $y<0$, zero is not a minimizer of $K(y,z)+V_j(z)$. Indeed,
\[
 \partial_zK(y,0)=\frac{-y}{\ell(y,0)}>0,
\]
whereas $0\le V_j(z)=O(z^2)$. A sufficiently small negative $z$
therefore reduces the value. If $V_j>V_{j-1}$ on $[-d,0)$, evaluate
at a minimizer $z_*<0$ for $V_{j+1}(y)$. It follows that
\[
 V_{j+1}(y)=K(y,z_*)+V_j(z_*)
   >K(y,z_*)+V_{j-1}(z_*)\ge V_j(y).
\]
Induction proves the strict inequality for every $j$ and $y<0$.

The minimizing path is symmetric, and its middle edge has length
$\ell(y_m,y_m)=g+2u(y_m)$. Splitting the action at that edge gives
the exact identity
\begin{equation}\label{eq:bellman-excess}
 E_m=2\min_{y\in[-d,0]}\{e(y)-g/2+V_{m-1}(y)\}.
\end{equation}
The outer minimizer cannot be zero, since $e'(0)=d/e(0)>0$ and
$V_{m-1}(y)=O(y^2)$. Evaluating the formula for $E_{m+1}$ at its
negative minimizer, and using strict increase of $V_j$ there, proves
$E_{m+1}>E_m$. Since $e(y)\ge D/2-x_R(y)\ge g/2$, equality
$E_m=0$ would force $y=0$ and then $d=0$, which is excluded.
Taking all heights zero gives
\[
 E_m\le2\sqrt{(g/2)^2+d^2}-g
      =\frac{2d^2}{\sqrt{(g/2)^2+d^2}+g/2}\le\frac{2d^2}{g}.
\]
Here $d\le1/20$ and $g>79/100$, proving the last bound in
\eqref{eq:period-excess}. A bounded increasing sequence has summable
positive increments, which must tend to zero.
\end{proof}

\subsection{Joint arithmetic and its exact scope}
A regular periodic return orbit has a total record $(K,N,T,h)$:
winding $K\in\Z^2$, physical collision count $N$, flight length $T$,
and induced period $h$. Define its phase at
$(u,s,\beta,c)\in\T^2\times\T\times\R\times\T$ to be
$u\cdot K+sN+\beta T-ch$, with $\T=\R/(2\pi\Z)$.

\begin{theorem}[Trivial joint periodic annihilator]\label{thm:joint-periodic}
For every $R\in I$, if
\begin{equation}\label{eq:joint-periodic-phase}
 \exp i(u\cdot K+sN+\beta T-ch)=1
\end{equation}
for every regular periodic orbit of $F_R^*$, then
$\beta=0$ and $u=s=c=0$ on their tori. It suffices to test the four
cycles of Proposition~\ref{prop:lattice}, the long normal two-cycle,
and the family of Lemma~\ref{lem:excursion-family}.
\end{theorem}
\begin{proof}
The long normal cycle gives $\exp i(2s+2\beta g-c)=1$.
The $m$th excursion gives
$\exp i((2m+1)s+\beta L_m-mc)=1$. Dividing the latter by the
$m$th power of the former yields
\[
                    \exp i(s+\beta E_m)=1.
\]
Taking consecutive quotients shows that
$\beta(E_{m+1}-E_m)\in2\pi\Z$ for all $m$. If $\beta\ne0$,
Theorem~\ref{thm:period-excess} makes these quantities nonzero and
strictly smaller than $2\pi$ in absolute value for large $m$, a
contradiction. Hence $\beta=0$. The four unchanged first-return
records in \eqref{eq:matrix} now give $u=s=c=0$.
\end{proof}

\begin{corollary}[Finite uniform separation on compact frequency sets]
\label{cor:compact-periodic}
For every $B<\infty$ and $\varepsilon>0$, there are a finite list of
these physical periodic words and $\eta>0$ such that, uniformly in
$R\in I$ and in $|\beta|\le B$ at distance at least $\varepsilon$
from the trivial phase, at least one word has
$|\exp i(u\cdot K+sN+\beta T-ch)-1|\ge\eta$.
\end{corollary}
\begin{proof}
Each word has fixed integer record and continuous length on $I$.
Theorem~\ref{thm:joint-periodic} gives a nonzero discrepancy at every
point of the compact parameter--frequency set in question. The open
sets on which a given word has positive discrepancy cover that set.
Take a finite subcover and minimize the maximum of its discrepancies.
This minimum is positive and supplies $\eta$.
\end{proof}

\begin{corollary}[Full real period rank]\label{cor:real-period-rank}
The real vectors $(K_1,K_2,N,h,T)$ of five fixed cycles span $\R^5$.
The absolute determinant is $2(\sqrt3-1)$, independently of $R$.
Consequently no nonzero real linear combination of the four induced
record coordinates is cohomologous to a constant by a transfer
function for which periodic evaluation is valid.
\end{corollary}
\begin{proof}
Append lengths $L_2,L_3,L_4,L_4$ to the four rows in
\eqref{eq:matrix}, and append $(0,0,2,1,2g)$ as the fifth row.
Subtracting the first row from the fifth leaves only the last entry
$2g-L_2=2(\sqrt3-1)$. The remaining determinant is that of $V$,
of absolute value one. Evaluation of a real cohomology on these
cycles therefore forces all coefficients, including the constant,
to vanish.
\end{proof}

The preceding theorems are statements about real billiard orbits, not
merely formal integer records. They also exclude a circle-valued
coboundary with a continuous transfer function on the regular return
branches: an almost-everywhere identity then extends to each regular
periodic point by continuity and full support of the section measure.
A transfer function assumed only measurable requires a separate
regularity theorem. Similarly, a covariance conclusion requires a
proved variance--coboundary criterion in the relevant function space.
If that criterion and a continuous covariance family are established,
Corollary~\ref{cor:real-period-rank} gives positive definiteness and,
on $I$, a uniform positive lower eigenvalue. Existence and continuity
of that covariance are not consequences of the determinant alone.
Compact periodic separation is not a quantitative operator contraction
or a returned temporal nonintegrability estimate.

\section{Quantitative separation by logarithmic-length periods}
\label{sec:quantitative-periods}
The complete periodic annihilator was determined in
Theorem~\ref{thm:joint-periodic}. We now estimate the period increments
uniformly in both the radius and the number of collisions. This gives an
explicit separation at large roof frequencies, using three actual
periodic words. Throughout this section $g=\sqrt3-2R$, $d=1/2-R$,
and $e$, $\ell$ and $u$ are the functions in+\eqref{eq:excursion-lengths}. The coordinates below are heights, not
arc coordinates.

Define the half action on $[-d,0]^m$ by
\begin{equation}\label{eq:half-action-v4}
 W_{m,R}(y)=e_R(y_1)-g/2+
       \sum_{j=1}^{m-1}\bigl(\ell_R(y_j,y_{j+1})-g\bigr)+u_R(y_m).
\end{equation}
Reflection symmetry and uniqueness of the full minimizer imply
$E_m(R)=2\min W_{m,R}$. Denote its minimizing coordinates by
$y^{(m,R)}_1,\ldots,y^{(m,R)}_m$.

\begin{lemma}[Uniform Hessian and endpoint estimates]
\label{lem:uniform-half-hessian}
On the entire height box, the Hessian $H=\nabla^2 W_{m,R}$ is
tridiagonal and satisfies
\begin{equation}\label{eq:half-hessian-bounds}
 3I_m\le H\le12I_m,\qquad H_{jj}\le10,\qquad
 \frac12\le-H_{j,j+1}\le2\quad(1\le j<m).
\end{equation}
The last condition is empty when $m=1$. For all $1\le j\le m$,
\begin{equation}\label{eq:height-two-sided}
 \frac3{440}\,20^{-(j-1)}
 \le -y^{(m,R)}_j
 \le \frac1{21}\left(\frac7{10}\right)^{j-1}.
\end{equation}
All constants are independent of $m$ and $R\in I$.
\end{lemma}
\begin{proof}
Write $x(y)=\sqrt{R^2-y^2}$, $a=D/2-x(y)$,
$h=D-x(y)-x(z)$, $v=z-y$, with $D=\sqrt3$.
The elementary bounds on the whole box are
\[
 \begin{gathered}
 x>11/25,\quad |x'|\le5/44,\quad79/200<a<43/100,\\
 79/100<h<86/100,\quad e<11/25,\quad\ell<9/10.
 \end{gathered}
\]
Also $u''=R^2/x^3\ge1/R\ge100/47$ and
$u''\le55225/21296<3$. For the norm of a vector-valued function,
the Hessian is its transverse positive semidefinite quadratic form
plus the Hessian of the vector dotted with its unit direction.
Consequently
\[
 e''\ge(a/e)u''>\frac{79}{88}\frac{100}{47}>\frac32,
 \qquad
 \nabla^2\ell\ge(h/\ell)\operatorname{diag}(u''(y),u''(z))
 >\frac32 I_2.
\]
Each variable of $W_m$ occurs in two such terms, with $u''>3/2$
at the terminal endpoint. This proves the lower bound $H\ge3I_m$,
including $m=1$.

The same norm formula gives
\[
 \ell_{yy},\ell_{zz}
 \le\frac{1+(5/44)^2}{79/100}+\frac{55225}{21296}<4,
 \qquad
 e''\le\frac{1+(5/44)^2}{79/200}+\frac{55225}{21296}<6.
\]
An explicit mixed derivative is
\begin{equation}\label{eq:mixed-length-hessian}
 -\ell_{yz}
 =\frac{(h+vy/x(y))(h-vz/x(z))}{\ell^3}.
\end{equation}
Each numerator factor lies between $78/100$ and $87/100$.
Using $79/100<\ell<9/10$ proves $1/2<-\ell_{yz}<2$.
The diagonal bounds for $W_m$ are therefore $10$ at its first
endpoint, $8$ in its interior, and $7$ at its last endpoint; when
$m=1$ the bound is $9$. Absolute row sums are at most $12$.
The symmetric matrix bound $H\le12I_m$ follows, proving
\eqref{eq:half-hessian-bounds}.

At zero the gradient is $a_0 e_1$, where
$a_0=d/e_R(0)$ satisfies $3/44\le a_0\le1/7$.
Integrate the Hessian on the segment from zero to the minimizer and
write the resulting matrix as $A$. Stationarity gives
$Ay^{(m,R)}=-a_0e_1$. The matrix $B=10I_m-A$ is entrywise
nonnegative and tridiagonal. Its eigenvalues lie in $[-2,7]$, so
$\|B/10\|_2\le7/10$. Hence
\[
 A^{-1}=\frac1{10}\sum_{k\ge0}(B/10)^k.
\]
For the $(j,1)$ entry every power before $j-1$ vanishes. The
nearest-neighbor product at power $j-1$ is at least $20^{-(j-1)}$.
Entrywise nonnegativity and the norm bound yield
\[
 \frac1{10}20^{-(j-1)}
 \le (A^{-1})_{j1}
 \le\frac1{10}\sum_{k\ge j-1}(7/10)^k
 =\frac13(7/10)^{j-1}.
\]
Multiplication by the two bounds on $a_0$ proves
\eqref{eq:height-two-sided}. The matrix is an averaged Hessian of
the actual nonlinear action; no quadratic approximation of the
minimizing orbit has been made.
\end{proof}

\begin{theorem}[Two-sided excess increments]
\label{thm:quantitative-excess}
For every $R\in I$ and $m\ge1$,
\begin{equation}\label{eq:quantitative-excess}
 \frac1{10000}\,400^{-m}
 \le E_{m+1}(R)-E_m(R)
 \le\frac1{300}\left(\frac{49}{100}\right)^{m-1}.
\end{equation}
In particular $E_m$ converges uniformly on $I$ to a continuous
function $E_\infty$, with
\begin{equation}\label{eq:uniform-excess-limit}
 0<E_\infty(R)-E_m(R)
 \le\frac1{153}\left(\frac{49}{100}\right)^{m-1}.
\end{equation}
\end{theorem}
\begin{proof}
For a minimizer $y$ of $W_m$, append the coordinate zero. With
$K(y,z)=\ell(y,z)-g$, the resulting change in half action is
\[
 K(y_m,0)-u(y_m)
 =\sqrt{(g+u(y_m))^2+y_m^2}-(g+u(y_m))
 \le\frac{y_m^2}{2g}.
\]
It follows that $E_{m+1}-E_m\le y_m^2/g$. The upper bound in
\eqref{eq:height-two-sided} and $g>79/100$ give the upper
coefficient $100/34839<1/300$ in \eqref{eq:quantitative-excess}.

Conversely, truncate a minimizer $z$ of $W_{m+1}$ after its first
$m$ coordinates. The removed contribution is
\[
 K(z_m,z_{m+1})+u(z_{m+1})-u(z_m)\ge2u(z_{m+1}),
\]
by \eqref{eq:kernel-coercivity}. Thus
\[
 E_{m+1}-E_m\ge4u(z_{m+1})\ge\frac{2z_{m+1}^2}{R}
 \ge\frac9{45496}\,400^{-m}>\frac1{10000}\,400^{-m}.
\]
Here $u(z)=z^2/(R+x(z))\ge z^2/(2R)$ and
\eqref{eq:height-two-sided} were used. Both comparisons concern
minima of the exact actions. Summing the geometric upper bound
proves \eqref{eq:uniform-excess-limit}. Each $E_m$ is continuous
by Lemma~\ref{lem:excursion-family}; the uniform limit is continuous.
\end{proof}

Let $P_0$ denote the long normal orbit and $P_m$ the $m$th excursion.
For a periodic word $P$ write
\[
 Z_P(u,s,\beta,c;R)
 =\exp i\bigl(u\cdot K_P+sN_P+\beta T_P-ch_P\bigr).
\]
Put $a=\log(100/49)$ and, for $b\ge1$, define
\begin{equation}\label{eq:frequency-period-budget}
 m(b)=1+\left\lceil\frac{\log b}{a}\right\rceil,
 \qquad \gamma=\frac{\log400}{a}-1.
\end{equation}

\begin{theorem}[Explicit large-roof phase separation]
\label{thm:log-period-separation}
For $|\beta|=b\ge1$, all $R\in I$, and arbitrary torus
frequencies $u,s,c$, the three true periods $P_0,P_{m(b)},P_{m(b)+1}$
satisfy
\begin{equation}\label{eq:log-period-separation}
 \max_P|Z_P-1|
 \ge\frac{b\,400^{-m(b)}}{15000\pi}
 \ge\frac{b^{-\gamma}}{2400000000\pi}.
\end{equation}
The exponent is $\gamma=\log400/\log(100/49)-1$.
For a whole band $1\le|\beta|\le B$, the fixed three periods
$P_0,P_{m(B)},P_{m(B)+1}$ suffice, with lower bound
$400^{-m(B)}/(15000\pi)$. Their longest physical word has
$2m(B)+3$ collisions. All assertions are uniform in the radius.
\end{theorem}
\begin{proof}
The actual counts and lengths give the exact identity
\begin{equation}\label{eq:three-period-cancellation}
 Z_{P_{m+1}}\overline{Z_{P_m}}\overline{Z_{P_0}}
       =\exp\bigl(i\beta(E_{m+1}-E_m)\bigr).
\end{equation}
The winding, physical count and induced period cancel. With $m=m(b)$,
\eqref{eq:quantitative-excess} implies
\[
 \frac{b\,400^{-m}}{10000}
 \le b(E_{m+1}-E_m)\le\frac1{300}.
\]
For $|v|\le\pi$, $|e^{iv}-1|\ge2|v|/\pi$. The distance
from one of a product of three unit complex numbers is at most
the sum of their individual distances from one. Applying these
facts to \eqref{eq:three-period-cancellation} proves the first
bound in \eqref{eq:log-period-separation}. Since
$m(b)\le2+\log b/a$, the second bound follows. For the band
assertion use $m(B)$; the upper phase bound still holds, while
$b\ge1$ supplies the printed lower bound. The number of collisions
was proved in Lemma~\ref{lem:excursion-family}.
\end{proof}

\begin{corollary}[An approximate-phase lower bound]
\label{cor:approximate-phase-bound}
Let $q:Y_R^*\to\mathbb C$ have modulus one and be continuous in
neighborhoods of all states in the three cycles used above. Suppose
$q$ is measurable elsewhere and let
\[
 D(q)=\mathop{\rm ess\,sup}_{x\in Y_R^*}
 \left|e^{i(u\cdot\kappa_R^*(x)+s r_R^*(x)+\beta\tau_R^*(x))}
                 q(F_R^*x)-e^{ic}q(x)\right|.
\]
For $b=|\beta|\ge1$ and $m=m(b)$,
\begin{equation}\label{eq:approximate-phase-bound}
 D(q)\ge\frac{b\,400^{-m}}{10000\pi(m+1)}.
\end{equation}
\end{corollary}
\begin{proof}
The return branches through these cycles are regular, and their
section measure has full support. Continuity therefore extends the
essential bound to every state of each selected cycle. Multiplication
around a cycle of induced period $h$ telescopes the unit phases and
gives $|Z_P-1|\le hD(q)$. In
\eqref{eq:three-period-cancellation} the sum of the three periods is
$1+m+(m+1)=2(m+1)$. The lower sine bound in the preceding proof
therefore gives
$2b(E_{m+1}-E_m)/\pi\le2(m+1)D(q)$. Apply the lower bound
in \eqref{eq:quantitative-excess}.
\end{proof}

The estimates are quantitative statements about the specified mechanical
return map, and in particular about continuous circle-valued approximate
phases. They do not assert regularity of a merely measurable transfer
function, a lower bound for a spectral eigenfunction on these orbits,
or a transfer-operator resolvent estimate. Those are distinct analytical
steps when using \eqref{eq:approximate-phase-bound} in a high-frequency
argument. The raw characteristic-function tail is not replaced by the
period discrepancy.

\section{Uniform phase separation and positive-measure coercivity}
\label{sec:periodic-coercivity}
We strengthen the frequency estimate by comparing consecutive exact
minimizers. The choice of a period will depend on the frequency and
possibly on the radius. This dependence is essential: a fixed triple
of periods is not asserted to separate the entire unbounded frequency
axis by a positive constant.

Retain $K_R(y,z)=\ell_R(y,z)-g$ and $W_{m,R}$ from
\eqref{eq:half-action-v4}. Write
\[
 a_m(R)=-y_m^{(m,R)}>0,\qquad
 \Delta_m(R)=E_{m+1}(R)-E_m(R)>0.
\]
The symbol $a_m$ denotes a terminal height; it is unrelated to the
logarithmic constant in \eqref{eq:frequency-period-budget}.

\subsection{Comparison of successive exact minima}
\begin{lemma}[The terminal value function]
\label{lem:terminal-response}
For $y\in[-d,0]$ put
\[
 V_R(y)=\min_{z\in[-d,0]}\{K_R(y,z)+u_R(z)\},\qquad
 \delta_R(y)=V_R(y)-u_R(y).
\]
The minimizer $f_R(y)$ is unique. It belongs to $(-d,0)$ for $y<0$
and equals zero for $y=0$. The functions are twice continuously
differentiable, with one-sided derivatives at the endpoints, and
\begin{equation}\label{eq:terminal-response}
 f_R'(y)\ge\frac1{14},\qquad
 |f_R(y)|\ge\frac{|y|}{14},\qquad
 |\delta_R''(y)|\le2,\qquad |\delta_R'(y)|\le2|y|.
\end{equation}
These estimates hold on the entire height interval, uniformly in $R$.
\end{lemma}
\begin{proof}
The second derivative in $z$ is $K_{zz}+u''>0$. At $z=-d$,
\[
 K_z(y,-d)=\frac{h(y,-d)(-d)/x_R(-d)-d-y}{\ell_R(y,-d)}<0,
 \qquad u_R'(-d)<0.
\]
At $z=0$ and $y<0$ the derivative is $-y/\ell_R(y,0)>0$.
Consequently the minimizer is interior and unique. For $y=0$,
\eqref{eq:kernel-coercivity} gives the unique minimizer zero.
The implicit function theorem applies to the stationarity equation,
including its smooth extension near the endpoints, and gives
\[
 f_R'(y)=\frac{-K_{yz}(y,f_R(y))}
                    {K_{zz}(y,f_R(y))+u_R''(f_R(y))}
 \ge\frac{1/2}{4+3}=\frac1{14}.
\]
Here the individual segment bounds established in the proof of
Lemma~\ref{lem:uniform-half-hessian} were used. Integration from $y$
to zero proves the second inequality of \eqref{eq:terminal-response}.

The envelope formula is
\[
 V_R''(y)=K_{yy}(y,f_R(y))-
       \frac{K_{yz}(y,f_R(y))^2}{K_{zz}(y,f_R(y))+u_R''(f_R(y))}.
\]
The Hessian of $K_R(y,z)+u_R(z)$ dominates $(3/2)I_2$.
Its Schur complement therefore is at least $3/2$: evaluate its
quadratic form on $(1,t)$ and minimize over $t$. Also
$V_R''\le K_{yy}\le4$. Since $100/47\le u_R''<3$, subtraction
shows $|\delta_R''|\le2$. The first derivatives of $K_R$ and $u_R$
vanish at the origin, so $\delta_R'(0)=0$. Integration proves the
last inequality. All estimates concern the nonlinear action itself.
\end{proof}

\begin{proposition}[Consecutive terminal heights]
\label{prop:terminal-comparison}
For every $m\ge1$ and $R\in I$,
\begin{equation}\label{eq:terminal-comparison}
 a_{m+1}(R)\ge\frac3{70}a_m(R).
\end{equation}
\end{proposition}
\begin{proof}
Let $y$ minimize $W_{m,R}$ and let $(X,z)$ minimize $W_{m+1,R}$,
where $X$ consists of the first $m$ coordinates. Eliminating the
last coordinate gives exactly
\[
 \min_z W_{m+1,R}(X,z)=W_{m,R}(X)+\delta_R(X_m),
 \qquad z=f_R(X_m).
\]
All minimizing coordinates are interior by
Lemma~\ref{lem:excursion-family}. Stationarity and the fundamental
theorem of calculus thus yield
\[
 A(X-y)=-\delta_R'(X_m)e_m,\qquad
 A=\int_0^1\nabla^2W_{m,R}(y+t(X-y))\dd t\ge3I_m.
\]
The segment stays in the height box. In particular
$0<(A^{-1})_{mm}\le1/3$. Taking the last coordinate and applying
Lemma~\ref{lem:terminal-response}, we obtain
\[
 |y_m|\le |X_m|+\frac13|\delta_R'(X_m)|
             \le\frac53|X_m|.
\]
Finally $|z|=|f_R(X_m)|\ge|X_m|/14$. These two inequalities give
\eqref{eq:terminal-comparison}. No monotonicity of the terminal
heights, and no quadratic replacement of $W_{m,R}$, is assumed.
\end{proof}

\begin{corollary}[A lower ratio for the excess increments]
\label{cor:increment-ratio}
For every $m\ge1$ and $R\in I$,
\begin{equation}\label{eq:increment-ratio}
 \Delta_{m+1}(R)\ge10^{-5}\Delta_m(R).
\end{equation}
\end{corollary}
\begin{proof}
The append and truncate comparisons in the proof of
Theorem~\ref{thm:quantitative-excess} give
\[
 \Delta_m\le\frac{a_m^2}{g},\qquad
 \Delta_{m+1}\ge\frac2R a_{m+2}^2.
\]
Applying \eqref{eq:terminal-comparison} twice, and using
$g>79/100$ and $R\le47/100$, yields
\[
 \Delta_{m+1}\ge\frac{2g}{R}\left(\frac3{70}\right)^4\Delta_m
 >\frac{158}{47}\left(\frac3{70}\right)^4\Delta_m
 >\frac{\Delta_m}{100000}.
\]
This is a comparison between adjacent increments, stronger for the
present purpose than separate exponential upper and lower bounds.
\end{proof}

\subsection{A frequency-dependent triple and a fixed finite menu}
Put
\begin{equation}\label{eq:uniform-phase-constants}
 \eta_0=\frac1{4000000},\qquad
 \sigma_0=\frac1{6000000\pi},\qquad
 M(b)=m(b)+1,
\end{equation}
where $m(b)$ is defined in \eqref{eq:frequency-period-budget}.

\begin{theorem}[Uniform positive phase separation]
\label{thm:uniform-phase-separation}
For every $R\in I$ and $b\ge1$ the integer
\begin{equation}\label{eq:adaptive-period}
 j(R,b)=\min\{j\ge1:b\Delta_j(R)\le1/2\}
\end{equation}
is well defined and satisfies $j(R,b)\le m(b)$. If $|\beta|=b$,
then, for arbitrary $u,s,c$,
\begin{equation}\label{eq:uniform-phase-separation}
 \max_{P\in\{P_0,P_{j(R,b)},P_{j(R,b)+1}\}}
       |Z_P(u,s,\beta,c;R)-1|\ge\sigma_0.
\end{equation}
The longest of these physical words has at most $2m(b)+3$ collisions.
For the entire band $1\le|\beta|\le B$, the fixed menu
\begin{equation}\label{eq:fixed-period-menu}
 \mathcal P_B(R)=\{P_0(R),P_1(R),\ldots,P_{m(B)+1}(R)\}
\end{equation}
satisfies the same lower bound for its maximum, uniformly in $R$.
It contains $m(B)+2$ periods.
\end{theorem}
\begin{proof}
The upper bound of \eqref{eq:quantitative-excess} gives
$b\Delta_{m(b)}\le1/300$. Thus the defining set in
\eqref{eq:adaptive-period} is nonempty. If $j(R,b)=1$, its lower
bound gives $b\Delta_1\ge\Delta_1\ge1/4000000=\eta_0$.
If $j(R,b)>1$, minimality and \eqref{eq:increment-ratio} give
\[
 b\Delta_{j(R,b)}>10^{-5}/2\ge\eta_0.
\]
In either case
\begin{equation}\label{eq:selected-phase-window}
 \eta_0\le b\Delta_{j(R,b)}\le1/2.
\end{equation}
The exact three-period identity \eqref{eq:three-period-cancellation}
therefore has distance at least $2\eta_0/\pi$ from one. The distance
of a product of three unit complex numbers from one is at most the
sum of their three distances. This proves
\eqref{eq:uniform-phase-separation}, since
$2\eta_0/(3\pi)=\sigma_0$. The physical and induced counts are
those of Lemma~\ref{lem:excursion-family}; hence the stated word
length. Finally $m(b)\le m(B)$ on the band, so the selected triple
is contained in \eqref{eq:fixed-period-menu}.

The index selection uses a first crossing, not monotonicity of
$\Delta_j$. It need not be continuous in $R$ or $b$. The menu has
fixed indices and consists of analytically varying true periods,
which is the asserted parameter-uniform conclusion.
\end{proof}

\begin{corollary}[Logarithmic uniform-norm coercivity]
\label{cor:logarithmic-coercivity}
Let $q$ be circle-valued and measurable on $Y_R^*$, and continuous
in neighborhoods of every state of the triple in
\eqref{eq:uniform-phase-separation}. With $D(q)$ as in
Corollary~\ref{cor:approximate-phase-bound}, for $b=|\beta|\ge1$,
\begin{equation}\label{eq:logarithmic-coercivity}
 D(q)\ge\frac1{4000000\pi(j(R,b)+1)}
       \ge\frac1{4000000\pi M(b)}.
\end{equation}
\end{corollary}
\begin{proof}
Continuity and full support extend the essential bound to the
selected regular periodic states. Telescoping around each period
bounds its phase discrepancy by its induced length times $D(q)$.
The three induced lengths sum to $2(j(R,b)+1)$. Apply the product
identity and \eqref{eq:selected-phase-window} to obtain
\[
 2\eta_0/\pi\le2(j(R,b)+1)D(q).
\]
This proves \eqref{eq:logarithmic-coercivity}. The modulus-one
condition is used in telescoping and is not a consequence of an
arbitrary operator normalization.
\end{proof}

\subsection{Regular tubes of the actual return map}
We use the Euclidean metric in local $(\alpha,p)$ coordinates. The
following uniformity concerns the specified periodic family, not all
branches of the induced map.

\begin{lemma}[Uniform return neighborhoods of the period family]
\label{lem:uniform-period-tubes}
There are constants $0<\rho_0<1$, $L\ge2$ and $B_0\ge1$, independent
of $R$ and $m$, with the following properties. For every selected
state $x_P$ of $P_0$ or any $P_m$, the ball $B(x_P,\rho_0)$ lies
in $Y_R^*$ and in one regular first-return branch. On this ball
$r_R^*$ and $\kappa_R^*$ are constant, $r_R^*\in\{2,3\}$, and
\[
 \operatorname{Lip}(F_R^*)\le L,\qquad
 \operatorname{Lip}(\tau_R^*)\le B_0.
\]
If $P$ has induced period $h$ and $0<\rho\le\rho_0L^{-h}$, every
$x\in B(x_P,\rho)$ follows the same $h$ return branches as $x_P$,
and, writing $F=F_R^*$,
\begin{align}\label{eq:period-tube-bounds}
 d(F^k x,F^k x_P)&\le L^k\rho &&(0\le k\le h),\\
 |T_{h,R}(x)-T_{h,R}(x_P)|&\le2B_0L^h\rho,\\
 \nu_R^*(B(x_P,\rho))&=\rho^2/(4c_*).
\end{align}
\end{lemma}
\begin{proof}
We verify that the family stays a uniform distance from every local
singularity relevant to a single return. The estimates in
Lemma~\ref{lem:excursion-family} hold for all its contact heights:
$-2d/5<y<0$, the long-flight third-disk clearances are at least
$9/500$, and the incidence cosine at $A$ exceeds $1/25$.
At the facing contacts on $O,C$ it is bounded below by a positive
constant from the horizontal-component inequalities in that proof.
End flights have a uniform positive distance from every third disk:
the opposite facing center is horizontally separated by $D/2$,
$B$ is vertically separated, and the remaining lattice centers obey
the enumerated uniform separation bounds. All flight lengths are
bounded above and below by positive constants.

Return membership has uniform margins as well. The selected contacts
on $O$ satisfy
\[
 |\alpha-\pi/6|\le1/22<3/50,\qquad
 |p|\le1/22+2/79<3/20.
\]
The unselected contacts on $C$ and $A$ have angles near $7\pi/6$
and equal to $2\pi/3$, respectively, uniformly separated from $U$
and the four squares in $Y_R$. The long normal orbit has the same
strict margins. Each selected-to-selected path consequently contains
two physical flights, or three through $C,A,O$.

For clarity, compactness is applied only to these bounded-length
paths. The parameters $R$ and their contact coordinates range in a
compact box. Every limiting path satisfies the non-grazing,
nonocclusion and return-membership margins just established. In a
neighborhood of such a path the selected entry roots are simple;
the physical collision map and flight time are analytic in the
initial state and parameter. There are only the finitely many
center itineraries just listed, up to lattice translation. A finite
cover of their compact closures therefore gives a common neighborhood
radius and uniform first-derivative bounds for the actual two- or
three-flight first returns. Decrease the radius so its coordinate
balls lie in $Y_R^*$ and contain no change of lattice labels; enlarge
the derivative bounds to $L\ge2$, $B_0\ge1$. This proves the first
part without taking a derivative of an arbitrarily long return.

Starting in a ball of radius at most $\rho_0L^{-h}$, induction gives
the first inequality in \eqref{eq:period-tube-bounds} and preserves
the same return decisions at each step. Summing the roof variations
gives $B_0\rho\sum_{k<h}L^k\le2B_0L^h\rho$. Finally
$\dd\nu_R^*=(4\pi c_*)^{-1}\dd\alpha\dd p$ on the ball, whose
Euclidean area is $\pi\rho^2$. This proves the last equality.
\end{proof}

\begin{theorem}[Positive-measure coercivity for regular phases]
\label{thm:positive-measure-coercivity}
Fix $0<\alpha\le1$, $H\ge1$ and $1\le p<\infty$. Suppose $q$ is
measurable and circle-valued on $Y_R^*$ and has an $\alpha$-H\"older
representative, of constant at most $H$, on the balls of radius
$\rho_0$ around all states of the selected triple. Define
\[
 \mathcal D_q(x)=e^{i(u\cdot\kappa_R^*(x)+s r_R^*(x)+\beta\tau_R^*(x))}
                      q(F_R^*x)-e^{ic}q(x).
\]
For $b=|\beta|\ge1$, put
\begin{equation}\label{eq:coercivity-radius}
 \rho(b,H)=L^{-M(b)}
 \min\left\{\frac{\rho_0}{2},\frac{\sigma_0}{8bB_0},
                       \left(\frac{\sigma_0}{8H}\right)^{1/\alpha}\right\}.
\end{equation}
Then the actual invariant return probability satisfies
\begin{equation}\label{eq:positive-measure-coercivity}
 \|\mathcal D_q\|_{L^p(\nu_R^*)}
 \ge\frac{\sigma_0}{2M(b)}
             \left(\frac{\rho(b,H)^2}{4c_*}\right)^{1/p}.
\end{equation}
All geometric constants are uniform over $R\in I$. No independence
of the returns is assumed.
\end{theorem}
\begin{proof}
Select from the triple a period $P$ with $|Z_P-1|\ge\sigma_0$.
Let its selected state be $x_P$ and its induced period be $h\le M(b)$.
Write $\varphi=u\cdot\kappa_R^*+s r_R^*+\beta\tau_R^*$, $F=F_R^*$
and $S_h\varphi=\sum_{k<h}\varphi\circ F^k$. The lattice and count
terms are constant along the corresponding branches in the tube.
Thus for $x\in B(x_P,\rho(b,H))$, Lemma~\ref{lem:uniform-period-tubes}
gives
\[
 |S_h\varphi(x)-S_h\varphi(x_P)|\le2bB_0L^h\rho(b,H).
\]
At $x_P$, the modulus of
$e^{iS_h\varphi(x_P)}q(F^h x_P)-e^{ich}q(x_P)$ is $|Z_P-1|$.
Since $F^h x_P=x_P$, the H\"older bounds at the two endpoints and
\eqref{eq:coercivity-radius} imply, throughout the smaller ball,
\begin{align*}
 \left|e^{iS_h\varphi(x)}q(F^h x)-e^{ich}q(x)\right|
 &\ge\sigma_0-2bB_0L^h\rho(b,H)-2H(L^h\rho(b,H))^\alpha\\
 &\ge\sigma_0/2.
\end{align*}
Choose the indicated representative on these balls; changes on null
sets do not affect the following integrals. Telescoping $h$ times
bounds the left side by $\sum_{k<h}|\mathcal D_q(F^k x)|$ almost
everywhere. Minkowski's inequality on the smaller ball and invariance
of $\nu_R^*$ give
\[
 \frac{\sigma_0}{2}\nu_R^*(B(x_P,\rho(b,H)))^{1/p}
 \le\sum_{k<h}\|\mathcal D_q\circ F^k\|_{L^p(B(x_P,\rho(b,H)),\nu_R^*)}
 \le h\|\mathcal D_q\|_{L^p(\nu_R^*)}.
\]
Use the exact ball mass and $h\le M(b)$ to conclude. The estimate
uses measure preservation, not disjointness or probabilistic
independence of the iterates of the ball.
\end{proof}

\begin{corollary}[A polynomial budget for the phase regularity]
\label{cor:polynomial-regularity-budget}
Fix $\alpha,p,H_0$ as above and $\zeta\ge0$. If the H\"older constants
satisfy $H\le H_0b^\zeta$, set
\begin{equation}\label{eq:coercivity-exponent}
 \kappa=\frac{\log L}{\log(100/49)}+
                       \max\{1,\zeta/\alpha\}.
\end{equation}
There is $C>0$, depending only on these fixed constants and the
mechanical family, such that
\begin{equation}\label{eq:polynomial-coercivity}
 \|\mathcal D_q\|_{L^p(\nu_R^*)}
             \ge C\frac{b^{-2\kappa/p}}{1+\log b}\qquad(b\ge1).
\end{equation}
\end{corollary}
\begin{proof}
Writing $a=\log(100/49)$, we have $M(b)\le3+\log b/a$ and hence
$L^{-M(b)}\ge L^{-3}b^{-\log L/a}$. The minimum in
\eqref{eq:coercivity-radius} is at least
\[
 \min\left\{\frac{\rho_0}{2},\frac{\sigma_0}{8B_0},
           \left(\frac{\sigma_0}{8H_0}\right)^{1/\alpha}\right\}
                     b^{-\max\{1,\zeta/\alpha\}}.
\]
Substitution into \eqref{eq:positive-measure-coercivity} proves the
claim with a positive constant independent of $R$ and $b$.
\end{proof}

\subsection{The operator interface}
Theorems~\ref{thm:uniform-phase-separation} and
\ref{thm:positive-measure-coercivity} have different hypotheses.
The former is an unconditional estimate for the specified mechanical
periods. The latter is an estimate for a stated regularity class of
circle-valued phases, in the actual invariant measure. In particular
it supplies a positive-measure obstruction rather than a conclusion
obtained by evaluating an unspecified measurable representative at
isolated periodic points.

To apply \eqref{eq:polynomial-coercivity} to a transfer operator,
one must construct from its approximate spectral vectors phases in
that class, bound the H\"older constant, and compare their physical
$L^p$ defect with the operator defect. Such a comparison would
contradict an error smaller than the right side of
\eqref{eq:polynomial-coercivity}. It is not proved merely by writing
an eigenvalue equation. In particular a division by the modulus of a
spectral vector requires a separate lower-modulus estimate.
General Liv\v{s}ic regularity results for Markov systems
\cite{BHN} concern precise dynamical and cohomological hypotheses;
they do not supply a uniform quantitative regularity bound for these
approximate phases without checking those hypotheses and the defect
comparison. No such comparison is assumed here for the full induced
operator. The common operator realization and the all-branch raw
residual estimates in the final section remain the requirements for
the original joint density theorem.

\section{Exact return means and stationary unfinished flights}
\label{sec:exact-return-means}
The enlarged section gives explicit normalization constants without any
local-limit assumption. We use ergodicity of the finite-horizon
periodic dispersing collision map, as supplied by the standard billiard
framework \cite{Young}. No spectral conclusion for the new first-return
map is imported from that fact.

\begin{proposition}[Exact induced means]\label{prop:exact-induced-means}
For the actual return records to $Y_R^*$,
\begin{equation}\label{eq:induced-mean-vector}
 \int(\kappa_{R,1}^*,\kappa_{R,2}^*,r_R^*,\tau_R^*)\dd\nu_R^*
 =\left(0,0,\frac1{c_*},\frac{\bar\tau_R}{c_*}\right),
 \quad
 \bar\tau_R=\frac{\sqrt3}{4R}-\frac{\pi R}{2}.
\end{equation}
The mean roof satisfies
\begin{equation}\label{eq:mean-roof-modulus}
 \left|\frac{\mathrm d}{\mathrm dR}\int\tau_R^*\dd\nu_R^*\right|
 <\frac{15}{4c_*}\qquad(R\in I).
\end{equation}
The equilibrium physical collision rate remains $1/\bar\tau_R$.
\end{proposition}
\begin{proof}
For any integrable physical collision observable $f$, the disjoint
return-tower levels partition the collision section up to null sets.
Invariance, first for nonnegative functions and then for positive and
negative parts, gives the Kac identity
\[
 \int_{Y_R^*}\sum_{j=0}^{r_R^*(x)-1} f(T_R^j x)\dd\nu(x)
                         =\int f\dd\nu.
\]
Taking $f=1$ and dividing by $c_*$ gives the mean return count.
Taking $f=\tau_R$ gives the mean induced roof. The physical
suspension normalization proved earlier gives the displayed value
of $\bar\tau_R$. The physical displacement is integrable by the
finite-horizon bound. Spatial inversion of the entire lattice maps
its value to its negative and preserves $\nu$, so its mean is zero.
This proves \eqref{eq:induced-mean-vector}.

The section volume $c_*$ is independent of $R$. Differentiating the
explicit expression for $\bar\tau_R$ yields
$-\sqrt3/(4R^2)-\pi/2$. The bounds $\sqrt3<7/4$,
$\pi<22/7$ and $R\ge9/20$ show its absolute value is less than
$15/4$. Finally the ratio of mean physical collisions per return to
mean physical time per return is $1/\bar\tau_R$.
\end{proof}

\begin{proposition}[The complete induced suspension]\label{prop:induced-suspension}
The equilibrium flow is represented by the suspension over $F_R^*$
with roof $\tau_R^*$ and probability
\begin{equation}\label{eq:induced-suspension}
 \frac{\nu_R^*(\mathrm dx)\,\mathrm ds}{\int\tau_R^*\dd\nu_R^*},
               \qquad 0\le s<\tau_R^*(x).
\end{equation}
If $A_R$ and $B_R$ denote respectively the age since the last visit
and the time to the next visit to $Y_R^*$ under this stationary law,
then, for $q\ge0$,
\begin{equation}\label{eq:residual-tails}
 \Pr(A_R>q)=\Pr(B_R>q)
 =\frac{\int(\tau_R^*-q)_+\dd\nu_R^*}
              {\int\tau_R^*\dd\nu_R^*}.
\end{equation}
At either one fixed endpoint of a time interval of length $t$,
the incomplete induced flight, its physical collision count and its
lattice displacement divided by $\sqrt t$ converge to zero in
probability as $t\to\infty$, for each fixed $R$.
\end{proposition}
\begin{proof}
In the physical suspension over $T_R$, group the consecutive
single-flight intervals before each first return to $Y_R^*$.
The tower identity and the normalization
$c_*\int\tau_R^*\dd\nu_R^*=\bar\tau_R$ convert the original
suspension measure exactly to \eqref{eq:induced-suspension}.
This rearranges the actual dependent path; it introduces no new
launch law and no independent flight variables. The marked formula
of Proposition~\ref{prop:marked} can consequently be grouped into
these complete return records while keeping both endpoint pieces.

For a given roof $\tau$, the set of $s\in[0,\tau)$ with age $s>q$
has length $(\tau-q)_+$. The set with remaining time $\tau-s>q$
has the same length. Integration proves \eqref{eq:residual-tails}.
Its right side tends to zero by the integrability in
Proposition~\ref{prop:exact-induced-means}. Each endpoint roof is
finite almost surely under the length-biased stationary law;
stationarity gives the same marginal at the other fixed endpoint.
Thus its ratio to $\sqrt t$ tends to zero in probability.

Every physical flight has length at least $1-2R\ge3/50$. The
number of collisions in an incomplete interval of length $s$ is
at most $1+s/(1-2R)$. Its displacement between fundamental cells
is bounded by a constant plus a fixed multiple of $s$, since the
speed is one and the lattice cell has bounded diameter. These
bounds prove the same endpoint statement for the two marked rewards.
\end{proof}

A fixed-endpoint assertion is not a bound on the maximum residual over
all times in $[0,t]$. Nor does pointwise integrability in $R$ imply a
uniform residual estimate over $I$. Such uniform or process-level
claims require uniform integrability or a quantitative tail for the
induced roof. The functional limit assumptions in
Proposition~\ref{prop:clock} therefore remain explicit. The exact
mean and the modulus \eqref{eq:mean-roof-modulus}, by contrast, are
already available for geometric error propagation. They do not
require recovery of the entire preparation distribution in A2-GEOM.

\section{Stability of the actual return records}
\label{sec:return-stability}
We compare the first returns themselves, including their changing domains.
The comparison is made on the common collision space
$M=\T\times(-1,1)$ with the probability $\nu$ in \eqref{eq:flux}.
It is not a comparison between independently chosen inducing schemes.
The squares defining $Y_R$ move analytically with $R$, and $U$ is fixed.
Consequently $Y_R^*$ has constant measure $c_*$, and its boundary is a
finite union of smooth arcs, continuously varying with $R$.

We use only ergodicity, not a parameter-uniform mixing rate, from the
finite-horizon Sinai theory. In the present family the obstacles are
smooth, strictly dispersing and disjoint, and the horizon is finite.
The flow-mixing theorem \cite[Corollary 1.3 and correction note]{BDL}
therefore implies ergodicity of the collision map: the suspension of a
nontrivial invariant collision set would be a nontrivial invariant flow
set. Thus every $Y_R^*$ has full saturation. This explicitly supplies the
full-tower hypothesis in the Kac identity used above.

Write $r=r_R^*$, $F=F_R^*$ and $\tau^*=\tau_R^*$ in sums indexed by
induced returns. For $n\ge1$ define the actual cumulative record
\[
 J_{n,R}=\left(\sum_{j<n}\kappa_R^*\circ F^j,
                   \sum_{j<n}r_R^*\circ F^j,
                   \sum_{j<n}\tau_R^*\circ F^j\right)
       =(K_{n,R},N_{n,R},T_{n,R})
\]
on $Y_R^*$. Its values lie in $\Z^3\times\R$; in particular the last
coordinate is one cumulative time, not $n$ independent roof coordinates.
A tilde denotes extension by zero to $M$. Let $p_{n,R}$ denote the raw
mixed density of $J_{n,R}$ under $\nu_R^*$, whose existence is included
in Theorem~\ref{thm:return-stability} below. Norms of densities use
counting measure times Lebesgue measure, denoted by $\mathfrak m$.

\begin{lemma}[Almost-everywhere stability of a finite return record]
\label{lem:finite-return-stability}
Fix $R_0\in I$ and a finite $n$. Outside a $\nu$-null subset of $M$,
$\widetilde J_{n,R}(x)\to\widetilde J_{n,R_0}(x)$ as $R\to R_0$.
For $x\in Y_{R_0}^*$ outside that null set, the physical word, every
intermediate return decision, $K_{n,R}(x)$ and $N_{n,R}(x)$ are constant
for all sufficiently close $R$. On a neighborhood of such a point the
cumulative time is the physical analytic length of that fixed word.
The neighborhood and parameter tolerance may depend on $x$ and $n$.
\end{lemma}
\begin{proof}
Remove the boundary of $Y_{R_0}^*$, all singular initial states and all
preimages under $T_{R_0}$ of that boundary. These sets are null: the
boundary has zero area, the map preserves $\nu$, and the billiard
singularities and their iterates have zero flux measure. Remove also
the null set on which a required return is infinite. For a remaining
$x\in Y_{R_0}^*$, the first $n$ returns contain finitely many regular
physical collisions. Every tested state lies strictly inside or strictly
outside the return set. The finite-itinerary analyticity following
\eqref{eq:entry}, and the continuous motion of the return boundary,
preserve all these decisions for nearby $(R,x)$. They also preserve
the selected entry roots and lattice labels. The final physical count
and displacement are therefore unchanged, while the sum of flight
lengths varies analytically. A remaining point outside $Y_{R_0}^*$
stays outside $Y_R^*$ for nearby $R$, and its zero extension stays zero.
These statements prove the asserted convergence. No common neighborhood
for infinitely many words is used.
\end{proof}

\begin{theorem}[Finite-record stability in density and first moment]
\label{thm:return-stability}
For every fixed $n\ge1$, the maps
\begin{equation}\label{eq:return-Lone}
 R\longmapsto\widetilde J_{n,R}\in L^1(\nu;\R^4),\qquad
 R\longmapsto p_{n,R}\in L^1(\mathfrak m)
\end{equation}
are continuous. Both maps are uniformly continuous on $I$.
Moreover $\{\widetilde N_{n,R}:R\in I\}$ and
$\{\widetilde T_{n,R}:R\in I\}$ are uniformly integrable.
If $\eta_{n,R}$ is the record law, then
\begin{equation}\label{eq:record-Wone}
 W_1(\eta_{n,R},\eta_{n,S})
 \le c_*^{-1}\|\widetilde J_{n,R}-\widetilde J_{n,S}\|_{L^1(\nu)},
\end{equation}
where $W_1$ uses the Euclidean distance on $\Z^3\times\R$.
In particular, for the original, unmodified characteristic functions,
\begin{align}\label{eq:raw-frequency-comparison}
 |\Phi_{n,R}(\omega)-\Phi_{n,S}(\omega)|
 &\le \|p_{n,R}-p_{n,S}\|_{L^1(\mathfrak m)},\nonumber\\
 |\Phi_{n,R}(\omega)-\Phi_{n,S}(\omega)|
 &\le c_*^{-1}|\omega|\,
       \|\widetilde J_{n,R}-\widetilde J_{n,S}\|_{L^1(\nu)}.
\end{align}
No modulus uniform in $n$, and no frequency-integrated estimate, is
asserted by these inequalities.
\end{theorem}
\begin{proof}
The Kac identity and invariance of $F_R^*$ give, exactly,
\begin{equation}\label{eq:block-Kac-mass}
 \int_M\widetilde N_{n,R}\dd\nu=n,\qquad
 \int_M\widetilde T_{n,R}\dd\nu=n\bar\tau_R.
\end{equation}
For nonnegative functions $h_j\to h$ almost everywhere with
$\int h_j\to\int h$, the identity
\[
 \|h_j-h\|_1=\int h_j+\int h-2\int\min(h_j,h)
\]
and dominated convergence for $\min(h_j,h)\le h$ prove $L^1$
convergence. Apply this to the two nonnegative coordinates, using
Lemma~\ref{lem:finite-return-stability}, \eqref{eq:block-Kac-mass}
and continuity of $\bar\tau_R$. Uniform finite horizon bounds every
single-flight lattice label by a fixed constant. Thus
$|\widetilde K_{n,R}|\le C\widetilde N_{n,R}$. The already established
$L^1$ convergence of the counts makes these bounds uniformly integrable
along each convergent parameter sequence. Almost-everywhere convergence
and truncation then give $L^1$ convergence of the displacement as well.
This proves the first continuity in \eqref{eq:return-Lone}.

A continuous image of the compact interval $I$ is compact in $L^1$.
For completeness, an $L^1$-compact nonnegative family is uniformly
integrable. Cover it by finitely many $L^1$ balls of radius $\epsilon$
with centers $h_1,\ldots,h_q$. For a member $h$ in the ball about $h_i$,
\[
 \int_{\{h>A\}}h
 \le 2\|h-h_i\|_1+
              2\int_{\{h_i>A/2\}}h_i.
\]
Choose $A$ for the finite set of centers and then let $\epsilon$
decrease. This proves the stated uniform integrability without a tail
assumption or an exponential return estimate.

We next prove existence and $L^1$ continuity of the density. Partition
$Y_R^*$, up to null sets, by total physical count and ordered lattice
word. The argument of Corollary~\ref{cor:raw-density} applies to this
actual first-return event as well: on each regular word the cumulative
length has at most one critical point by
Proposition~\ref{prop:general-edge}; elsewhere it is a submersion.
Coarea on a countable exhaustion and then summation give absolute
continuity. This argument recomputes the return event for $Y_R^*$;
it does not transfer a density normalization from $Y_R$.

Fix $R_0$ and $\epsilon>0$. After deleting the null sets just described,
choose finitely many smooth nonnegative cutoffs $\psi_i$ on $M$, with
$\sum_i\psi_i\le1$, whose supports lie compactly in regular
first-$n$-return patches for $R_0$, such that
\[
 c_*^{-1}\int\sum_i\psi_i\dd\nu>1-\epsilon.
\]
Refine the patches so that on each one a coordinate derivative of the
cumulative length is nonzero. Compact support and the finite-return
lemma preserve the same word, all return decisions and this derivative
for $R$ close to $R_0$. On each patch the map consisting of the other
initial coordinate and the cumulative length is a local diffeomorphism.
It and its inverse depend continuously in $C^1$ on $R$ on slightly
larger compact patches. Changing variables and integrating the other
coordinate therefore shows that the pushforward of
$c_*^{-1}\psi_i\nu$ has a density continuous in $L^1$ with $R$.
One can see this directly by extending the transformed smooth weight
by zero off its coordinate patch: the compactly supported cutoff
vanishes near the patch boundary, the Jacobian is bounded away from
zero, and dominated convergence applies on a fixed bounded coordinate
box. The lattice coordinates are constant on the patch.

The remainder is a positive measure. Its mass is
$1-c_*^{-1}\int\sum_i\psi_i\dd\nu<\epsilon$, for all these $R$,
because the cutoff supports stay inside $Y_R^*$ and $\nu(Y_R^*)=c_*$.
Consequently
\[
 \|p_{n,R}-p_{n,R_0}\|_1
 \le\sum_i\|p_{i,R}-p_{i,R_0}\|_1+2\epsilon.
\]
Let $R\to R_0$ and then $\epsilon\to0$. This proves the second
continuity in \eqref{eq:return-Lone}. Compactness of $I$ gives uniform
continuity of both maps.

For a 1-Lipschitz function $f$, subtract the constant $f(0)$. Then
\[
 \int f\dd\eta_{n,R}
 =c_*^{-1}\int_M f(\widetilde J_{n,R})\dd\nu.
\]
The same identity holds for $S$. Subtraction and the Lipschitz bound,
followed by the dual characterization of $W_1$, give
\eqref{eq:record-Wone}. All first moments are finite by
\eqref{eq:block-Kac-mass}. Direct integration against the densities
gives the first bound in \eqref{eq:raw-frequency-comparison}. For the
second, use the zero extensions and
$|e^{ia}-e^{ib}|\le|a-b|$; the off-section constants cancel because
both section measures equal $c_*$. These are estimates of the raw
transforms, with no convolution or change of observable.
\end{proof}

\begin{proposition}[Weighted records and a positive denominator]
\label{prop:weighted-return-stability}
Fix $n$. Suppose measurable insertions $w_R:M\to\R$ satisfy
$|w_R|\le M_0$ and, for every $R_0$, $w_R(x)\to w_{R_0}(x)$ for
$\nu$-almost every $x$ as $R\to R_0$. The $w_R$-weighted record
densities $p^w_{n,R}$ are continuous in $L^1(\mathfrak m)$.
Set
\[
 d_n(\delta)=\sup_{|R-S|\le\delta}\|p_{n,R}-p_{n,S}\|_1,
 \qquad
 d_n^w(\delta)=\sup_{|R-S|\le\delta}\|p^w_{n,R}-p^w_{n,S}\|_1.
\]
These moduli tend to zero. For any Borel observation event $B$ in the
record space, if $\eta_{n,R}(B)\ge b>0$ and
$d_n(\delta)<b$, then, whenever $|R-S|\le\delta$,
\begin{equation}\label{eq:finite-conditioning-stability}
 \left|\frac{\int_Bp^w_{n,R}\dd\mathfrak m}{\eta_{n,R}(B)}
       -\frac{\int_Bp^w_{n,S}\dd\mathfrak m}{\eta_{n,S}(B)}\right|
 \le\frac{d_n^w(\delta)+M_0d_n(\delta)}{b-d_n(\delta)}.
\end{equation}
\end{proposition}
\begin{proof}
Use the compact regular patches in the preceding proof. Replacing
$w_R$ by $w_{R_0}$ costs at most
$c_*^{-1}\int|w_R-w_{R_0}|\dd\nu\to0$, by dominated convergence
and contraction of total variation under pushforward. On each fixed
patch approximate the bounded measurable weight $w_{R_0}$ in $L^1$
by a smooth one. The approximation error in its pushforward is at
most the same $L^1$ error, uniformly in $R$. For the smooth weight,
the coordinate-change proof above applies. The leftover measure has
total variation at most $M_0\epsilon$. This proves weighted density
continuity, and compactness gives the modulus.

Write $a_R=\int_Bp^w_{n,R}\dd\mathfrak m$ and
$b_R=\eta_{n,R}(B)$. Then
$|a_R-a_S|\le d_n^w(\delta)$,
$|b_R-b_S|\le d_n(\delta)$ and $|a_R|\le M_0b_R$.
Since $b_S\ge b-d_n(\delta)>0$, subtraction of the two ratios gives
\eqref{eq:finite-conditioning-stability}.
\end{proof}

The event $B$ may fix the lattice coordinates and restrict the roof
to an interval. The denominator in
\eqref{eq:finite-conditioning-stability} is an actual event probability.
The estimate does not remove its deterioration for small events. In
particular density $L^1$ continuity alone is not the shrinking-interval
LLT of Proposition~\ref{prop:conditioning}; that proposition retains
its stronger pointwise density hypotheses. An error bound
$|\widehat R-R|\le\delta$ can be inserted into the displayed moduli,
but no numerical sampling budget follows from a noneffective modulus.

\begin{corollary}[Uniform stationary endpoint control]
\label{cor:uniform-stationary-endpoints}
Put $H_R=\widetilde T_{1,R}$ and
$b_*=\min_{R\in I}\bar\tau_R>0$. The stationary induced suspensions,
regarded as probabilities on $M\times\R_+$, have densities
\[
 h_R(x,s)=\bar\tau_R^{-1}\one_{\{0\le s<H_R(x)\}}
\]
with respect to $\nu(\dd x)\dd s$, and
\begin{equation}\label{eq:suspension-Lone}
 \|h_R-h_S\|_1\le
 b_*^{-1}\bigl(\|H_R-H_S\|_1+|\bar\tau_R-\bar\tau_S|\bigr).
\end{equation}
Furthermore
\begin{equation}\label{eq:uniform-roof-tail}
 \Psi(q):=b_*^{-1}\sup_{R\in I}
                   \int_M H_R\one_{\{H_R>q\}}\dd\nu
             \longrightarrow0\quad(q\to\infty).
\end{equation}
Under the stationary flow, at either deterministic endpoint of an
interval of length $t$, the incomplete induced time, physical collision
count and lattice displacement are $o_P(\sqrt t)$, uniformly over
$R\in I$. The same is true jointly at any fixed finite number of
deterministic endpoints.
\end{corollary}
\begin{proof}
The normalization in \eqref{eq:induced-suspension} and
$c_*\int\tau_R^*\dd\nu_R^*=\bar\tau_R$ give $h_R$ exactly.
For fixed $x$, the symmetric difference of the intervals
$[0,H_R(x))$ and $[0,H_S(x))$ has length $|H_R(x)-H_S(x)|$.
Integrating, and then comparing the two normalizing constants, proves
\eqref{eq:suspension-Lone}. The uniform integrability in
Theorem~\ref{thm:return-stability} proves
\eqref{eq:uniform-roof-tail}.

The full roof containing a stationary endpoint has tail
\[
 \Pr_R\{\text{containing roof}>q\}
 =\bar\tau_R^{-1}\int_MH_R\one_{\{H_R>q\}}\dd\nu
 \le\Psi(q).
\]
This is the length-biased law; replacing it by
$\nu_R^*(\tau_R^*>q)$ would be incorrect. Both endpoint pieces
are at most this full roof. The collision count of an incomplete
piece of length $s$ is at most $1+(50/3)s$, and its lattice
displacement is bounded by $C(1+s)$, by the physical estimates in
Proposition~\ref{prop:induced-suspension}. Hence there are fixed
constants $C_0,C_1$ such that each marked endpoint norm is bounded
by $C_0+C_1H$ when its containing roof has length $H$. Its probability of exceeding
$\epsilon\sqrt t$ is at most
$\Psi((\epsilon\sqrt t-C_0)/C_1)$ for large $t$, uniformly in $R$.
Stationarity applies at every deterministic endpoint, including an
endpoint depending deterministically on $t$. A union bound proves
the finite-endpoint assertion; independence is not used.
\end{proof}

This removes the additional uniform-integrability premise for fixed
stationary endpoints in this family. It is not a bound on a supremum
over a growing time interval or at an arbitrary stopping time. The
functional clock theorem still needs its printed process hypotheses.
Likewise \eqref{eq:raw-frequency-comparison} is a uniform comparison
of raw characteristic functions, not an integrable tail for any one
of them. The high-frequency and central all-branch remainder problem
is unchanged by an $L^1$ comparison of probability densities.

\section{Exponential tails for the genuine return record}
\label{sec:exponential-returns}
We now use the spectral theory of the collision map, in addition to the
geometric and ergodic inputs used above. The operator in this section is
not the operator of the induced map. A smooth killing weight on the
collision space gives a scalar estimate for avoiding the actual return
section. In particular no multiplication by its discontinuous indicator
on an anisotropic space will be needed.

\subsection{The collision-space input}
We use the strong and weak distribution spaces of Demers--Zhang
\cite[Section 3.3]{DZ}. Their properties needed here are the local uniform
spectral gap and power bounds for unforced dispersing collision maps,
including small changes of the scatterers and their boundary lengths
\cite[Theorems 2.1--2.6 and Remark 2.7(b)]{DZ}. The spectral gap for an
individual unforced Lorentz map is the input from \cite{DZ11} used there.
Here is the precise consequence, together with the checks for this family.

\begin{lemma}[Local uniform collision spaces]
\label{lem:collision-spectral-input}
There is a finite cover of $I$ by parameter neighborhoods $I_j$. On each
$I_j$ there is a Banach space $\mathcal B_j$ of distributions on a common
collision space, a continuous mass functional $\ell_j$, and transfer
operators $\mathcal L_R$, such that
\begin{equation}\label{eq:collision-split}
 \mathcal L_R=\Pi_R+\mathcal N_R,\quad
 \Pi_R h=\ell_j(h)\nu,\quad
 \Pi_R\mathcal N_R=\mathcal N_R\Pi_R=0,\quad
 \|\mathcal N_R^k\|_{\mathcal B_j}\le C\vartheta^k
 \quad(k\ge0),
\end{equation}
where $0<\vartheta<1$ and $C<\infty$ can be chosen for all the finitely
many neighborhoods. The norms of $\nu$, $\ell_j$ and $\mathcal L_R$ are
uniformly bounded. Multiplication by a smooth function $g$ satisfies
\begin{equation}\label{eq:smooth-multiplier}
 \|g h\|_{\mathcal B_j}\le C\|g\|_{C^2}\|h\|_{\mathcal B_j}.
\end{equation}
The transfer convention is $(\mathcal L_R h)(f)=h(f\circ T_R)$.
No assertion about the unbounded induced twists is included in this lemma.
\end{lemma}
\begin{proof}
For the collision-space construction use coordinates $(\alpha,\varphi)$,
where $p=\sin\varphi$. The invariant probability is the same for all
parameters: $\dd\nu=(4\pi)^{-1}\cos\varphi\dd\alpha\dd\varphi$.
We do not conjugate a distribution space by the inverse of
$p=\sin\varphi$ at grazing. At a fixed $R_0$, use the boundary coordinate
$r_0=R_0\alpha$. The parametrization of the nearby circle is then
$r_0\mapsto R n_{r_0/R_0}$. It is a $C^3$-bounded deformation, $C^2$-close
to the original parametrization, and its constant speed $R/R_0$ tends to
one. This is precisely the boundary-length reparametrization permitted in
\cite[Remark 2.7(b)]{DZ}.

The physical gap between different disks is at least $3/50$. The
curvatures belong to $[100/47,20/9]$, and the parametrized boundaries have
uniform $C^3$ bounds. The finite horizon is uniform on $I$: a line with
no disk intersection at the smallest radius would be an infinite free
line in that table, which has no corridor; compactness of position and
direction then bounds the free length, and increasing the radius cannot
increase it. If a rectangular fundamental domain is desired, use the
two-scatterer cover with periods $(1,0)$ and $(0,\sqrt3)$ and disk centers
$0,b$. Its deck translation commutes with the collision map. Restricting
to the deck-invariant distributions gives the original table. Thus the
local billiard estimates are applied without an affine distortion of
specular reflection.

The uniform geometric hypotheses of \cite[Theorem 2.5]{DZ} and the
perturbation hypothesis of its Theorem 2.6 are satisfied. The unforced
collision map at each $R_0$ has the simple eigenvalue one and a spectral
gap. The strong/weak perturbation theorem, specifically the complementary
power bound in \cite[Theorem 2.1(3)]{DZ}, gives \eqref{eq:collision-split}
on a neighborhood of $R_0$. The eigenprojection has the displayed form
because the invariant probability is $\nu$ and mass is preserved. A
finite cover of the compact parameter interval makes the constants
uniform. This uses local common spaces and does not assert that all
arbitrary inducing schemes share a Banach space.

For completeness, smooth multiplication follows directly from the
stable-test and unstable-comparison norms defining these spaces. On a
single admissible stable curve, multiplying its test function by $g$
changes its H\"older norm by at most $C\|g\|_{C^1}$. For two matched
curves at distance $\epsilon$, split the difference of the weighted tests
into the old test difference and the difference of the two restrictions
of $g$. After graph matching, the latter has $C^q$ norm at most
$C\|g\|_{C^2}\epsilon^{1-q}$ by interpolation between its $C^0$ and
$C^1$ bounds. The defining exponents satisfy $\beta<1-q$, so division
by $\epsilon^\beta$ is bounded for $\epsilon\le1$. The stable part of the
norm controls this second term and the unstable part controls the first.
The weak norm has the same single-curve multiplication bound. Density of
the defining smooth distributions extends these bounds to
$\mathcal B_j$, proving \eqref{eq:smooth-multiplier}. The coordinate
changes above and the finite cover have uniformly bounded derivatives.
\end{proof}
\subsection{A smooth avoidance estimate}
Choose once and for all $\chi\in C^\infty(M)$ such that
\begin{equation}\label{eq:fixed-killing-weight}
 0\le\chi\le1,\quad
 \supp\chi\Subset\{ |\alpha|<1/400,\ |p|<1/400\},\quad
 a_\chi:=\int\chi\dd\nu>0.
\end{equation}
The indicated rectangle is inside the square of $Y_R$ centered at
$(0,0)$ for every $R$; hence it is inside $Y_R^*$. Its support is away
from grazing, so it is smooth also in the collision coordinates used in
Lemma~\ref{lem:collision-spectral-input}. This is a proof weight, not an
extra field in the collision record.

\begin{proposition}[Uniform soft-killing estimate]
\label{prop:soft-killing}
There exist $z_0,c,C>0$ such that for every $R\in I$ and $k\ge0$,
\begin{equation}\label{eq:soft-killing}
 \int_M\exp\left(-z_0\sum_{j=0}^{k-1}\chi\circ T_R^j\right)\dd\nu
       \le Ce^{-ck}.
\end{equation}
\end{proposition}
\begin{proof}
On a local collision space set
$\mathcal L_{R,z}=\mathcal L_R M_{e^{-z\chi}}$ for complex $z$.
The multiplier bound proves analyticity in the strong operator norm and,
uniformly in $R$, $\|\mathcal L_{R,z}-\mathcal L_R\|\le C|z|$ near
zero. Equation~\eqref{eq:collision-split} bounds the unperturbed
resolvents uniformly on a circle around one and on a second circle
$|w|=\vartheta_1$ with $\vartheta<\vartheta_1<1$, chosen disjointly.
A uniformly convergent resolvent Neumann series therefore defines a
rank-one analytic projection and a single eigenvalue $\lambda_R(z)$
near one for all $|z|<z_*$. The complementary powers are bounded by
$C\vartheta_1^k$, after decreasing $z_*$ if necessary.

Since $\ell_j(\nu)=1$, differentiation at zero gives
\[
 \lambda_R(0)=1,\qquad
 \lambda_R'(0)=-\ell_j(\mathcal L_R(\chi\nu))=-a_\chi.
\]
The fixed complex disk and the uniform resolvent bounds give a uniform
bound on the second derivative of $\lambda_R$ on a smaller disk, by
Cauchy's formula. Thus
$|\lambda_R(z)-1+a_\chi z|\le C_2|z|^2$ uniformly in $R$.
For real small $z$ the eigenvalue is real, by conjugation and uniqueness,
and positive, by its closeness to one. Choose a fixed $z_0>0$ such that
$C_2z_0\le a_\chi/2$ and $z_0<z_*$. Then
$0<\lambda_R(z_0)\le1-a_\chi z_0/2$. The spectral decomposition gives
$\|\mathcal L_{R,z_0}^k\|\le Ce^{-ck}$. Only finitely many collision
spaces occur, so the same $z_0,c,C$ work on all of $I$.

Iteration of the transfer identity, first for densities and then as
bounded functionals, yields
\[
 \ell_j(\mathcal L_{R,z_0}^k\nu)
 =\int\prod_{j=0}^{k-1}e^{-z_0\chi(T_R^j x)}\dd\nu(x).
\]
The norms of $\ell_j$ and $\nu$ are bounded. This proves the asserted
scalar estimate. In particular the argument does not apply a local
large-deviation formula outside its stated neighborhood of the mean.
\end{proof}

\begin{theorem}[Exponential moments of the actual return block]
\label{thm:exponential-return}
Let
$Q_R=r_R^*+\tau_R^*+|\kappa_R^*|$ on the genuine first-return section.
There are constants $c_1,a_1,C_1>0$, independent of $R\in I$, such that
\begin{equation}\label{eq:exponential-return}
 \nu_R^*(r_R^*>k)\le C_1e^{-c_1k}\quad(k\ge0),\qquad
 \int e^{a_1 Q_R}\dd\nu_R^*\le C_1.
\end{equation}
For $Q_{n,R}=\sum_{j<n}Q_R\circ(F_R^*)^j$ and every $n\ge1$,
\begin{equation}\label{eq:block-exponential-moment}
 \int\exp(a_1 Q_{n,R}/n)\dd\nu_R^*\le C_1,\qquad
 \nu_R^*(N_{n,R}>L)\le C_1e^{-a_1L/n}.
\end{equation}
These are estimates for the original deterministic mechanical records.
\end{theorem}
\begin{proof}
If $x\in Y_R^*$ and $r_R^*(x)>k$, none of $T_Rx,\ldots,T_R^kx$
belongs to $Y_R^*$. Every term of $\sum_{j=1}^k\chi\circ T_R^j$ is
therefore zero. Positivity and invariance of $\nu$ imply
\[
 \nu_R^*(r_R^*>k)
 \le c_*^{-1}\int_M e^{-z_0\sum_{j=1}^k\chi\circ T_R^j}\dd\nu
 \le c_*^{-1}Ce^{-ck}.
\]
This comparison uses neither the regularity of $\one_{Y_R^*}$ as a
multiplier nor an independent-return construction.

Uniform finite horizon and the uniform bound on a single-flight lattice
step give $Q_R\le C_0 r_R^*$ for a fixed $C_0$. Summation of the
exponential tail of the positive integer $r_R^*$ gives a uniform
exponential moment of $Q_R$ for $a_1>0$ sufficiently small. Enlarge
$C_1$ to include the first bound. Finally H\"older's inequality with
$n$ equal exponents and invariance of $F_R^*$ give
\[
 \int\prod_{j<n}\exp\bigl(a_1Q_R\circ(F_R^*)^j/n\bigr)\dd\nu_R^*
 \le\prod_{j<n}\left(\int e^{a_1Q_R\circ(F_R^*)^j}\dd\nu_R^*\right)^{1/n}
 \le C_1.
\]
Since $N_{n,R}\le Q_{n,R}$, exponential Markov inequality proves the
last assertion. No independence is used in this step either.
\end{proof}

\begin{corollary}[All finite moments under parameter variation]
\label{cor:all-moment-continuity}
For fixed $n$ and $1\le p<\infty$, the map
$R\mapsto\widetilde J_{n,R}$ is continuous into $L^p(\nu;\R^4)$.
Every polynomial moment of a fixed finite list of return records is
finite and continuous in $R$. In particular every finite-lag covariance
of the induced record is continuous.
\end{corollary}
\begin{proof}
Almost-everywhere convergence is Lemma~\ref{lem:finite-return-stability}.
Equation~\eqref{eq:block-exponential-moment}, the fixed section mass and
$|J_{n,R}|\le Q_{n,R}$ make all fixed powers uniformly integrable.
Truncation followed by dominated convergence proves $L^p$ convergence
along every convergent parameter sequence. A finite list of records can
be expressed as differences of cumulative records. H\"older's inequality
and the same exponential moment bound give uniform integrability of each
polynomial in that list. Truncation proves its continuity. This proves
finite-lag, not yet infinite Green--Kubo, covariance continuity.
\end{proof}

\section{Exponential control of cumulative collision counts}
\label{sec:cumulative-returns}
The avoidance estimate controls more than a single return. It also
controls the time of the $n$th visit to the section, without separating
the successive return blocks. This observation improves the truncation
budget in \eqref{eq:physical-count-truncation} and supplies exponential
moments on a fixed complex neighborhood for the entire physical record.

Retain $z_0,c,C$ from Proposition~\ref{prop:soft-killing}. Write
$N_{n,R}=\sum_{j<n}r_R^*\circ(F_R^*)^j$, and enlarge a fixed constant
$D\ge1$, if necessary, so that
\begin{equation}\label{eq:cumulative-record-domination}
 |J_{n,R}|_1\le D N_{n,R},\qquad Q_{n,R}\le D N_{n,R}.
\end{equation}
Such a constant exists because each physical flight has uniformly
bounded length and lattice displacement. All probabilities and
expectations in this section are with respect to $\nu_R^*$.

\begin{theorem}[The time of the $n$th actual return]
\label{thm:cumulative-return-tail}
There are constants $A\ge1$, $a,c>0$, independent of $n$, $L$ and $R$,
such that
\begin{equation}\label{eq:cumulative-return-tail}
 \nu_R^*\{N_{n,R}>L\}\le A\exp(an-cL)
       \qquad(n\ge1,\ L\ge0).
\end{equation}
One may take $a=z_0$ and the exponent $c$ of the soft-killing estimate.
Put $v=a/c$. For $0<s<c$ there is a constant $B_s<\infty$, uniform in
$n,R$, for which
\begin{equation}\label{eq:cumulative-exponential-moment}
 \int e^{s(N_{n,R}-vn)_+}\dd\nu_R^*\le B_s,
 \qquad \int e^{sN_{n,R}}\dd\nu_R^*\le B_s e^{svn}.
\end{equation}
Consequently, for every $0<\lambda<c/D$,
\begin{equation}\label{eq:fixed-strip-record-moment}
 \int e^{\lambda Q_{n,R}}\dd\nu_R^*
          \le B_{\lambda D}e^{\lambda Dvn}.
\end{equation}
The constants $B_s$ are locally bounded for $0\le s<c$, with $B_0=1$.
The number $v$ is an upper envelope, not the Kac mean $c_*^{-1}$.
\end{theorem}
\begin{proof}
For an integer $L\ge0$ let
\[
 H_{L,R}(x)=\sum_{j=1}^{L}\one_{Y_R^*}(T_R^j x).
\]
Since $N_{n,R}$ is the physical collision time of the $n$th return,
$N_{n,R}>L$ if and only if $H_{L,R}<n$ on $Y_R^*$. The fixed smooth
weight satisfies $0\le\chi\le\one_{Y_R^*}$. Thus, on this event,
$\sum_{j=1}^{L}\chi(T_R^j x)\le n-1$. In particular, pointwise on
$Y_R^*$,
\[
 \one_{\{N_{n,R}>L\}}
 \le e^{z_0(n-1)}
       \exp\left(-z_0\sum_{j=1}^{L}\chi(T_R^j x)\right).
\]
Integrate on $Y_R^*$, enlarge the domain to $M$, and use invariance to
shift the sum by one collision. Equation~\eqref{eq:soft-killing} gives
\[
 \nu_R^*\{N_{n,R}>L\}\le c_*^{-1}C e^{z_0(n-1)-cL}.
\]
For real $L$ replace it by $\lfloor L\rfloor$ and enlarge the constant
by $e^c$. This proves \eqref{eq:cumulative-return-tail}. It uses a single
orbit's visit count, not a product of probabilities of return blocks.

It follows that, for $u\ge0$,
$\nu_R^*\{N_{n,R}>vn+u\}\le A e^{-cu}$. For a nonnegative random
variable $X$, Tonelli's theorem gives
\[
 \int e^{sX}\dd\nu_R^*
       =1+s\int_0^\infty e^{su}\nu_R^*\{X>u\}\dd u.
\]
Apply this with $X=(N_{n,R}-vn)_+$ to obtain
$B_s=1+As/(c-s)$. Since $N_{n,R}\le vn+X$, the second bound in
\eqref{eq:cumulative-exponential-moment} follows. Finally apply
\eqref{eq:cumulative-record-domination} with $s=\lambda D$.
\end{proof}

We now apply the tail to the actual record measure. For a measurable
insertion $w_{n,R}$ define the finite signed or complex measure
\[
 \mu_{n,R}^{w}=(J_{n,R})_*(w_{n,R}\nu_R^*),\qquad
 \mu_{n,R}^{w,L}=(J_{n,R})_*
       (w_{n,R}\one_{\{N_{n,R}\le L\}}\nu_R^*).
\]
The transform uses the convention $\widehat\mu(\omega)=
\int e^{i\omega\cdot j}\dd\mu(j)$, with
$\omega=(u_1,u_2,s,b)\in\T^3\times\R$.

\begin{proposition}[Weighted finite-band truncation]
\label{prop:linear-count-budget}
Fix a nonnegative integer $d$ and suppose
$|w_{n,R}|\le M(1+Q_{n,R})^d$. For every multi-index $\gamma$ there is a
constant $C_{d,\gamma}$, independent of $n,R,L,w$, such that
\begin{align}\label{eq:weighted-count-truncation}
 \|\mu_{n,R}^{w}-\mu_{n,R}^{w,L}\|_{\TV}
  &\le C_{d,0}M(1+L)^d e^{an-cL},\\
 \sup_{\omega\in\T^3\times\R}
 \left|\partial_\omega^\gamma
   (\widehat\mu_{n,R}^{w}-\widehat\mu_{n,R}^{w,L})(\omega)\right|
  &\le C_{d,\gamma}M(1+L)^{d+|\gamma|}e^{an-cL}.\nonumber
\end{align}
For a bounded insertion ($d=0$), a frequency region of finite volume
$V>0$, and $\varepsilon>0$, the sufficient integer cutoff for integrated
transform error at most $\varepsilon$ is
\begin{equation}\label{eq:linear-count-budget}
 L\ge\left\lceil\frac{a}{c}n+
       \frac1c\max\left\{0,\log\frac{C_{0,0}MV}{\varepsilon}\right\}
        \right\rceil.
\end{equation}
If $V_n\le V_0e^{\kappa n}$ and the required error is $e^{-\delta n}$,
then a cutoff linear in $n$ suffices. This remains true for each fixed
$d,\gamma$, with any fixed positive slack in the exponential budget.
\end{proposition}
\begin{proof}
For an integer $m\ge0$, integration of the tail gives
\begin{align*}
 &\int (1+N_{n,R})^m\one_{\{N_{n,R}>L\}}\dd\nu_R^*\\
 &\quad=(1+L)^m\nu_R^*\{N_{n,R}>L\}
    +m\int_L^\infty(1+t)^{m-1}\nu_R^*\{N_{n,R}>t\}\dd t\\
 &\quad\le C_m(1+L)^m e^{an-cL}.
\end{align*}
The integral term is zero when $m=0$. For $m\ge1$, the last inequality
follows by writing $t=L+u$ and using
$1+L+u\le(1+L)(1+u)$. Pushforward does not increase total variation.
For the derivative estimate differentiate under the integral: its
absolute integrand is at most
$M(1+Q_{n,R})^d|J_{n,R}|_1^{|\gamma|}
\one_{\{N_{n,R}>L\}}$. Equation~\eqref{eq:cumulative-record-domination}
and the preceding tail integral prove both estimates.

Multiplication by $V$ bounds the integral of the transform error. Solving
$C_{0,0}MV e^{an-cL}\le\varepsilon$ gives
\eqref{eq:linear-count-budget}. More generally choose a slope
$\ell>(a+\kappa+\delta)/c$. Then, for $L_n=\lceil\ell n\rceil$,
\[
 V_n(1+L_n)^{d+|\gamma|}e^{an-cL_n}
       \le C(1+n)^{d+|\gamma|}e^{-(\delta+\epsilon_0)n}
\]
for some $\epsilon_0>0$. A fixed additive increase of the cutoffs absorbs
the remaining constant and polynomial, proving the last assertion.
\end{proof}

\begin{remark}
This improves the sufficient cutoff of order
$n\log(V/\varepsilon)$ obtained from
\eqref{eq:physical-count-truncation} to order
$n+\log(V/\varepsilon)$. On an exponentially growing frequency band the
former budget is quadratic in $n$, whereas the new one is linear.
Neither estimate bounds the number of critical words, the variation of
an inverse coarea Jacobian, or a pointwise raw density. In particular the
all-branch derivative sum in Lemma~\ref{lem:raw-residual-sum} remains a
separate requirement. The truncation removes only actual high-count
records in an error estimate; the limiting datum is still the full
unmodified physical record.
\end{remark}

\section{A common realization and the physical renewal identity}
\label{sec:common-renewal}
The preceding moment estimate gives a common realization of the
unbounded record twists on a scale of Lebesgue spaces. It is useful to
construct this realization before seeking an anisotropic spectral
estimate: the collision and induced operators, their damping variables,
and their chronological composition can then be identified without
regularity assumptions on an indicator multiplier. The renewal algebra
is classical; compare \cite{GouezelRenewal}. The assertion here is its
exact realization for the moving physical first-return record.

\subsection{First-return operators on the common collision space}
Use $M=\T\times(-1,1)$ and $\dd\nu=(4\pi)^{-1}\dd\alpha\dd p$.
Let $P_R$ be multiplication by $\one_{Y_R^*}$ and let $E_R=I-P_R$.
Recall locally that
\[
 c_* = \nu(Y_R^*)=\frac{91}{10000\pi},\qquad
 \dd\nu_R^*=c_*^{-1}\one_{Y_R^*}\dd\nu.
\]
On $L^1(M,\nu)$, define
\[
 (\mathcal A_{R,\omega}h)(y)
   =e^{i\omega\cdot f_R(T_R^{-1}y)}h(T_R^{-1}y),\qquad
 f_R=(\kappa_{{\rm coll},1},\kappa_{{\rm coll},2},1,\tau_R),
\]
for real $\omega=(u_1,u_2,s,b)$. Collision singularities are discarded
as null sets. These are isometries. The operators $P_R,E_R$ are
contractions on every $L^p$, although no such assertion about their
action on anisotropic distribution spaces is needed or made here.

\begin{theorem}[Exact damped renewal realization]
\label{thm:common-renewal}
For $m\ge1$ set
\begin{equation}\label{eq:first-return-operator}
 \mathcal R_{m,R,\omega}
 =P_R\mathcal A_{R,\omega}
       (E_R\mathcal A_{R,\omega})^{m-1}P_R,
 \qquad
 \mathcal R_R(z,\omega)=\sum_{m\ge1}z^m\mathcal R_{m,R,\omega}.
\end{equation}
For $|z|<1$ the series converges in operator norm on $L^1(M,\nu)$ and
\begin{equation}\label{eq:induced-damped-bound}
 \|\mathcal R_R(z,\omega)h\|_1\le |z|\|P_Rh\|_1.
\end{equation}
On the closed subspace $P_RL^1$, with identity $P_R$, one has
\begin{equation}\label{eq:physical-renewal-identity}
 P_R(I-z\mathcal A_{R,\omega})^{-1}P_R
       =(I-\mathcal R_R(z,\omega))^{-1}.
\end{equation}
Moreover,
\begin{equation}\label{eq:schur-return-operator}
 \begin{split}
 \mathcal R_R(z,\omega)={}&zP_R\mathcal A P_R\\
 &+z^2P_R\mathcal A E_R
 (I-zE_R\mathcal A E_R)^{-1}
             E_R\mathcal A P_R,
 \end{split}
\end{equation}
where $\mathcal A=\mathcal A_{R,\omega}$ and the middle inverse acts
on $E_R L^1$.

For $|z|\le1$, the strongly defined boundary series represents the
actual first-return twist: for $y\in Y_R^*$ and $x=(F_R^*)^{-1}y$,
\begin{equation}\label{eq:actual-induced-realization}
 (\mathcal R_R(z,\omega)h)(y)
       =z^{r_R^*(x)}e^{i\omega\cdot G_R(x)}h(x).
\end{equation}
It is zero off $Y_R^*$. In particular, for every $n\ge1$,
\begin{equation}\label{eq:actual-record-pairing}
 c_*^{-1}\int_M\mathcal R_R(z,\omega)^n\one_{Y_R^*}\dd\nu
       =\int_{Y_R^*}z^{N_{n,R}}e^{i\omega\cdot J_{n,R}}\dd\nu_R^*.
\end{equation}
\end{theorem}
\begin{proof}
The collision map is invertible and preserves $\nu$. Poincar\'e
recurrence in both time directions makes $F_R^*$ invertible and
measure-preserving on $Y_R^*$, modulo null sets. In the product
\eqref{eq:first-return-operator}, the initial and final states lie in
$Y_R^*$ and the $m-1$ intervening states lie outside it. Thus its initial
domain is exactly $\{r_R^*=m\}$. Its weight is the exponential of the
sum of the $m$ actual collision records, hence $e^{i\omega\cdot G_R}$.
These initial domains partition $Y_R^*$ up to a null set. Their images
under the corresponding return branches also partition $Y_R^*$, by
invertibility of the induced map. Consequently
\[
 \sum_{m\ge1}|z|^m\|\mathcal R_{m,R,\omega}h\|_1
       =\int_{Y_R^*}|z|^{r_R^*}|h|\dd\nu
       \le |z|\|P_Rh\|_1 \quad (|z|\le1).
\]
Each summand has operator norm at most one, so for $|z|<1$ the
operator-norm series converges geometrically. At the boundary the last
display proves strong convergence for each $h$ and gives
\eqref{eq:actual-induced-realization}. It does not assert operator-norm
convergence of the boundary series.

Let $\mathcal U_0=P_R$ and $\mathcal U_m=P_R\mathcal A^mP_R$ for
$m\ge1$. Partition a trajectory with endpoints in $Y_R^*$ according to
its first return time. Chronological composition, with the first segment
on the right, gives
\[
 \mathcal U_m=\sum_{r=1}^m\mathcal U_{m-r}\mathcal R_{r,R,\omega}.
\]
All generating series converge in operator norm for $|z|<1$. Summing
the identity gives $\mathcal U(z)=P_R+\mathcal U(z)\mathcal R_R(z,\omega)$.
The inverse of $I-\mathcal R_R$ exists by
\eqref{eq:induced-damped-bound}, and
$\mathcal U(z)=P_R(I-z\mathcal A)^{-1}P_R$, proving
\eqref{eq:physical-renewal-identity}. Separating $m=1$ from the series
and summing the intervening outside-section iterates proves
\eqref{eq:schur-return-operator}. All these inverses are used only in
$|z|<1$.

Finally iterate \eqref{eq:actual-induced-realization} and change
variables by $(F_R^*)^n$. Its collision counts and record sums are
$N_{n,R}$ and $J_{n,R}$, respectively. Division by $c_*$ yields
\eqref{eq:actual-record-pairing}. No independence of successive
excursions enters the identity.
\end{proof}

\subsection{A fixed complex neighborhood for the unbounded twists}
To cross the damping boundary at $z=1$, use $z=e^w$. Define the operator
directly by the actual $n$th return, rather than by an unproved
continuation of a resolvent. For $y\in Y_R^*$ let $x=(F_R^*)^{-n}y$ and set
\begin{equation}\label{eq:block-scale-operator}
 (\mathcal T_{n,R}(w,\omega)h)(y)
       =e^{wN_{n,R}(x)+i\omega\cdot J_{n,R}(x)}h(x),
\end{equation}
with zero extension to $M\setminus Y_R^*$. Here $w\in\mathbb C$ and
$\omega\in\mathbb C^4$. Holomorphy is stated on this universal cover.
The formulas are invariant under real shifts by $2\pi\mathbb Z^3$
in the first three frequency coordinates. Local complex charts about
real torus frequencies are restrictions of the same formulas; no
additional global complex-torus construction is being asserted. The
symbols $J_{n,R}$ and $G_R$ denote four real coordinates, with the
displacement contributing two coordinates.

\begin{theorem}[Holomorphic realization on a common Lebesgue scale]
\label{thm:common-scale-holomorphy}
Fix $1\le q<p\le\infty$ and put $t=(q^{-1}-p^{-1})^{-1}$. Let
\[
 \sigma(w,\omega)=|\operatorname{Re}w|+D|\operatorname{Im}\omega|_\infty,
 \qquad \Omega_{p,q}=\{(w,\omega):t\sigma(w,\omega)<c\}.
\]
For every $n\ge1$, \eqref{eq:block-scale-operator} defines a
holomorphic map from $\Omega_{p,q}$ to
$\mathcal L(L^p(M,\nu),L^q(M,\nu))$. Its norm satisfies
\begin{equation}\label{eq:scale-operator-bound}
 \|\mathcal T_{n,R}(w,\omega)\|_{p\to q}
       \le c_*^{1/t}B_{t\sigma}^{1/t}e^{\sigma vn}.
\end{equation}
On every compact subset of $\Omega_{p,q}$, all fixed-order derivatives in $(w,\omega)$ have bounds $C_k e^{C_k n}$ uniform in $R$. At every fixed
$(w,\omega)\in\Omega_{p,q}$ and every $h\in L^p(M,\nu)$,
\begin{equation}\label{eq:scale-parameter-continuity}
 R\longmapsto\mathcal T_{n,R}(w,\omega)h
       \quad\hbox{is continuous in }L^q(M,\nu).
\end{equation}
For $\operatorname{Re}w<0$ and real $\omega$, this operator agrees
with $\mathcal R_R(e^w,\omega)^n$ wherever the latter is considered
as an $L^1$ operator. The resulting inserted record pairings extend+holomorphically across $w=0$ on this fixed neighborhood.
\end{theorem}
\begin{proof}
Change variables $y=(F_R^*)^n x$ in the $q$th power of the norm. The
absolute weight is at most $e^{\sigma N_{n,R}(x)}$, by
\eqref{eq:cumulative-record-domination}. H\"older's inequality with
$1/q=1/p+1/t$ and Theorem~\ref{thm:cumulative-return-tail} give
\[
 \|\mathcal T_{n,R}(w,\omega)h\|_q
 \le\left(\int_{Y_R^*}e^{t\sigma N_{n,R}}\dd\nu\right)^{1/t}
       \|h\|_p
 \le c_*^{1/t}B_{t\sigma}^{1/t}e^{\sigma vn}\|h\|_p.
\]
For $\sigma=0$ use $B_0=1$. This proves the operator bound, including
$p=\infty$.

On a compact subset of the stated domain choose a slightly larger
compact neighborhood for which $t\sigma'<c$. A derivative of order
$k$ in $(w,\omega)$ inserts a monomial in $N_{n,R}$ and the four
coordinates of $J_{n,R}$, bounded by $C_kN_{n,R}^k$. Every such
polynomial is absorbed by $e^{(\sigma'-\sigma)N_{n,R}}$. H\"older's
inequality and \eqref{eq:cumulative-exponential-moment} therefore bound
the derivative integrands and the Taylor remainders in operator norm.
The exponential power series on a sufficiently small polydisc converges
in that norm. This proves joint holomorphy, not just weak holomorphy of
individual pairings. The same estimates on the larger neighborhood,
or Cauchy's inequalities there, give the claimed derivative bounds.

We give the parameter-continuity argument because the underlying return
domains move. Fix $R_0$. Remove $\partial Y_{R_0}^*$, the forward and
backward collision singularities, and all finite collision iterates of
$\partial Y_{R_0}^*$. Each is a countable union of null curves, up to
singular endpoints. Backward recurrence holds almost everywhere. At a
remaining point of $Y_{R_0}^*$ the preceding $n$ returns use finitely many
nonsingular physical collisions; none is on a section boundary. The
implicit collision-root argument of
Theorem~\ref{thm:return-stability}, applied backwards, preserves this
finite itinerary for $R$ sufficiently near $R_0$. The backward starting
point and all flight lengths vary continuously; the lattice labels and
collision counts are fixed. At a remaining point outside $Y_{R_0}^*$
the section indicator stays zero. Thus
\eqref{eq:block-scale-operator} converges pointwise almost everywhere
when $h$ is bounded and continuous.

Choose $q'>q$ sufficiently close to $q$ so that $q'<p$ and
$t'\sigma<c$, where $t'=(1/q'-1/p)^{-1}$. The preceding bound in
$L^{q'}$ is uniform in nearby $R$. It gives uniform integrability of the
$q$th powers, so pointwise convergence implies convergence in $L^q$.
For $p<\infty$, approximate an arbitrary $h$ in $L^p$ by bounded smooth
functions and use the uniform $p\to q$ bound. For $p=\infty$, choose a
finite $p_0>q$ so large that $(1/q-1/p_0)^{-1}\sigma<c$. Approximation
in $L^{p_0}$ and its uniform $p_0\to q$ bound give the same conclusion
for $h\in L^\infty$. This avoids assuming norm density of smooth
functions in $L^\infty$.

In the damped region the two operators have the same pointwise formula,
by Theorem~\ref{thm:common-renewal}. For example, for $h\in L^p$ and
$b\in L^{q^\dagger}$, where $q^\dagger$ is the conjugate exponent of $q$,
\[
 c_*^{-1}\int b\,\mathcal T_{n,R}(w,\omega)h\dd\nu
 =\int_{Y_R^*}h(x)b((F_R^*)^nx)
       e^{wN_{n,R}(x)+i\omega\cdot J_{n,R}(x)}\dd\nu_R^*.
\]
This continuous pairing is holomorphic on $\Omega_{p,q}$. It agrees
with the damped renewal pairing there and hence is the asserted
extension across $w=0$.
\end{proof}

\subsection*{Compatibility of the realizations}
For real $\omega$ and $|z|<1$, the constructions fit the diagram
\[
 \begin{array}{ccc}
 \mathcal A_{R,\omega}&\longmapsto&\mathcal R_R(z,\omega)\\
 \Big\downarrow {\scriptstyle P_R(I-z\mathcal A)^{-1}P_R}
 &&\Big\downarrow {\scriptstyle (I-\mathcal R)^{-1}}\\
 \mathcal U_R(z,\omega)&=&\mathcal U_R(z,\omega).
 \end{array}
\]
The upper arrow is the first-return sum
\eqref{eq:schur-return-operator}; the bottom equality is
\eqref{eq:physical-renewal-identity}. On the common domain
$z=e^w$, $\operatorname{Re}w<0$, $\omega\in\mathbb R^4$, the second
compatibility is
\[
 \mathcal R_R(e^w,\omega)^n=\mathcal T_{n,R}(w,\omega),\qquad
 c_*^{-1}\langle b,\mathcal T_{n,R}h\rangle_\nu
 =\int h\,(b\circ(F_R^*)^n)e^{wN_{n,R}+i\omega\cdot J_{n,R}}
                        \dd\nu_R^*.
\]
The first identity is interpreted on $L^p\subset L^1$, with values in
$L^q$, when the scale operator is used; the second has the insertion
exponents of the theorem. Any anisotropic realization must reproduce
these pairings on its admissible test class. Such an additional
realization is not an arrow already established by this diagram.

\begin{remark}[Relation to the spectral and raw-density estimates]
The operator identities concern the genuine four-coordinate record,
and all their spaces live on the same collision probability space.
They specify a common realization with fixed exponential integrability
and strong parameter continuity. They do not assert operator-norm
continuity in $R$, quasi-compactness on $L^p$, or a holomorphic family of
endomorphisms of a single anisotropic space. At real frequencies the
induced $L^1$ operator is an isometry on its section, so its construction
alone gives no spectral gap. Likewise, maps $L^p\to L^q$ with $q<p$
cannot simply be iterated on one space outside the damped region; the
$n$-block operator above was defined directly.

The next analytical step is to establish an anisotropic realization
compatible with \eqref{eq:actual-record-pairing} and
\eqref{eq:physical-renewal-identity}, and to prove its phase
reconstruction and raw-residual bounds. The collision compensation
in Sections~\ref{sec:collision-covariance}--\ref{sec:functional-gaussian}
will establish the actual Gaussian covariance and functional limits
by a different route, without identifying an induced spectral family. The criteria in
Theorem~\ref{thm:phase-resolvent-criterion},
Proposition~\ref{prop:covariance-closure}, and
Lemma~\ref{lem:raw-residual-sum} continue to state these requirements
separately. In particular the common Lebesgue realization is not used
as a substitute for any of those estimates in the raw local limit.
\end{remark}

\section{Common charts for moving return domains}
\label{sec:common-charts}
We give a local change-of-variables estimate and then apply it to the
actual return domains. This makes explicit the domination used in
Theorem~\ref{thm:return-stability}. Critical neighborhoods are discarded
in the initial collision space, not by assuming a uniform density bound
near their images.

\begin{lemma}[A common submersion chart]
\label{lem:common-submersion-chart}
Let $U=U_1\times U_2\Subset\R^2$ be a bounded open rectangle and let
$f(R,u,v)$ be $C^2$ on a neighborhood of $J\times\overline U$, where
$J$ is a compact parameter interval. Suppose
$\partial_v f\ge\eta>0$ there. For $a\in C_c^1(U)$, let $g_R$ be the
density of the pushforward of $a(u,v)\dd u\dd v$ by $f_R$. Then
\begin{equation}\label{eq:chart-Lipschitz}
 \|g_R-g_S\|_{L^1(\R)}\le C_a|R-S|\qquad(R,S\in J).
\end{equation}
The constant depends only on $U$, a common bounded range of $f$, $\eta$,
the $C^2$ bound for $f$, and the $C^1$ bound and support of $a$.
The same assertion holds if $\partial_v f\le-\eta$.
\end{lemma}
\begin{proof}
For fixed $R,u$ the map $v\mapsto f(R,u,v)$ is one-to-one. Write its
inverse as $v_R(u,t)$. The pushforward under $(u,v)\mapsto(u,f_R(u,v))$
has density
\[
 A_R(u,t)=\frac{a(u,v_R(u,t))}{\partial_v f(R,u,v_R(u,t))}
\]
on its image, extended by zero elsewhere. All images lie in one bounded
$(u,t)$ box. Since $a$ vanishes on a fixed neighborhood of $\partial U$,
this zero extension introduces no boundary jump as $R$ varies. Wherever
$a$ or its derivatives are nonzero, the inverse is defined on a common
parameter neighborhood. If $M$ bounds the relevant derivatives of $f$,
implicit differentiation gives
\[
 |\partial_R v_R|\le M/\eta,\qquad
 |\partial_R A_R|
 \le \|a_v\|_\infty M/\eta^2
       +\|a\|_\infty(M/\eta^2+M^2/\eta^3).
\]
At points outside the inverse chart both sides are interpreted by the
zero extension, which vanishes on an open neighborhood of its moving
edge. The same bound therefore gives an $L^1$ majorant on the fixed box.
Integrating in $R$ proves
$\|A_R-A_S\|_1\le C_a|R-S|$. Since
$g_R(t)=\int A_R(u,t)\dd u$, integration in $u$ proves
\eqref{eq:chart-Lipschitz}. Reversing the $v$ coordinate handles the
negative derivative. The result also applies on a constant lattice
component of a mixed density.
\end{proof}

\begin{proposition}[Finite common patches with an explicit remainder]
\label{prop:finite-common-patches}
Fix $n\ge1$, $R_0\in I$, and $\epsilon>0$. There are finitely many
nonnegative smooth cutoffs $\psi_i$ on the common collision space, a
parameter neighborhood $J$ of $R_0$, and $C_{\epsilon,n,R_0}<\infty$
such that $\sum_i\psi_i\le1$, every support lies in one common regular
first-$n$-return patch for all $R\in J$, and
\begin{equation}\label{eq:common-patch-mass}
 c_*^{-1}\int\sum_i\psi_i\dd\nu>1-\epsilon.
\end{equation}
The lattice labels and all intermediate return decisions are fixed on
each patch. If $p^i_{n,R}$ is the pushforward density of
$c_*^{-1}\psi_i\nu$, then
\begin{align}\label{eq:common-patch-bound}
 \sum_i\|p^i_{n,R}-p^i_{n,S}\|_1
       &\le C_{\epsilon,n,R_0}|R-S|,\\
 \|p_{n,R}-p_{n,S}\|_1
       &\le2\epsilon+C_{\epsilon,n,R_0}|R-S|
               \qquad(R,S\in J).\nonumber
\end{align}
Before choosing these patches one may discard $N_{n,R_0}>L$ at cost at
most $C_1e^{-a_1L/n}$, with the constants of
Theorem~\ref{thm:exponential-return}.
\end{proposition}
\begin{proof}
The exponential moment bound proves the last assertion. Choose $L$ so
that this cost is less than $\epsilon/3$. On the remainder remove the
initial boundary of the section, all grazing and competing-root
singularities encountered before collision $L$, and the preimages of the
moving return boundary at $R_0$ through those collisions. These sets have
zero collision measure. On every remaining regular word also remove the
critical points of its cumulative flight length. Proposition
\ref{prop:general-edge} shows that there is at most one such point on a
regular optical word; the submersion complement has full measure.
There are only finitely many possible physical lattice words of length
at most $L$, since each flight label belongs to a fixed finite set.
No estimate for the number of unbounded-length words is used here.

Inner regularity of the collision measure and a countable exhaustion of
the regular submersion patches give a compact subset of mass greater
than $1-\epsilon$ in the normalized section. At every point of this
compact subset, choose a small coordinate rectangle on which one
coordinate derivative of the cumulative time is bounded away from zero.
The finite-itinerary lemma preserves the entire physical word and every
return decision on a slightly larger rectangle for $R$ close to $R_0$.
Select finitely many such rectangles and a subordinate family of smooth
cutoffs with sum at most one, equal to one on the compact subset.
Their supports have positive distance from all the excluded sets at
$R_0$. Finiteness gives a common parameter neighborhood $J$ on which
these distances, return decisions and derivative bounds persist. The
time functions and their first two derivatives are uniformly bounded
on these finitely many compact rectangles. Lemma
\ref{lem:common-submersion-chart} applies with
$a=\psi_i/(4\pi c_*)$ in $(\alpha,p)$ coordinates. Summing its bounds
proves the first line of \eqref{eq:common-patch-bound}.

For every $R\in J$ the remainder of the original probability is a
positive measure. Its mass is exactly
$1-c_*^{-1}\int\sum_i\psi_i\dd\nu<\epsilon$, because the section
measure is the fixed number $c_*$. The raw densities exist by the
submersion/coarea argument of Theorem~\ref{thm:return-stability}.
The two remainder densities have total $L^1$ norm less than
$2\epsilon$. This gives the second line. Nothing has been assumed about
the size of a density at a critical value or a sharp image boundary.
\end{proof}

The constants in \eqref{eq:common-patch-bound} expose the order of
choices: physical-count truncation, deletion of small initial-state
neighborhoods, compact common charts, and finally parameter tolerance.
They are not asserted to remain bounded as $n$ or $\epsilon^{-1}$ grows.
A quantitative global density modulus requires quantitative control of
those additional choices. Bounded measurable insertions are handled by
$L^1$ approximation on the same patches, as in
Proposition~\ref{prop:weighted-return-stability}; no artificial observation
coordinate is introduced.

\section{A uniform functional law for the mechanical clock}
\label{sec:uniform-fluid-clock}
Finite-record continuity also has a long-time consequence that does not
require a mixing rate. We prove a parameter-uniform functional law at the
first-order scale, including the maximum unfinished return over an entire
observation interval. The \(\sqrt t\)-scale process theorem in
Proposition~\ref{prop:clock} remains a separate statement.

Set
\[
 G_R=(\kappa_R^*,r_R^*,\tau_R^*),\qquad
 \bar G_R=(0,0,c_*^{-1},\bar\tau_R/c_*),\qquad J_{0,R}=0.
\]
All probabilities on the return section use $\nu_R^*$, not an
independent sequence of records. Write $\bar r=c_*^{-1}$ and
$\bar\tau_R^*=\bar\tau_R/c_*$. Both $\bar\tau_R^*$ and its reciprocal
are bounded on $I$.

\begin{lemma}[A compact subadditive argument]
\label{lem:compact-subadditive}
Let $a_n:I\to[0,\infty)$ be continuous, $a_0=0$, and
$a_{n+m}(R)\le a_n(R)+a_m(R)$. Suppose
$\inf_{n\ge1}a_n(R)/n=0$ for every $R$. Then
\[
 \lim_{n\to\infty}\sup_{R\in I}\frac{a_n(R)}n=0.
\]
More precisely, for every $\epsilon>0$ there is a finite integer $M$
such that, with $C=\sup_R a_1(R)$,
\begin{equation}\label{eq:subadditive-uniform-budget}
 \sup_R a_n(R)/n\le\epsilon+CM/n\qquad(n\ge1).
\end{equation}
\end{lemma}
\begin{proof}
For each $R$ choose a block length $m_R$ with
$a_{m_R}(R)/m_R<\epsilon$. Continuity preserves this strict inequality
on a neighborhood of $R$. Take a finite subcover and let $M$ be the
largest of its block lengths. For an arbitrary parameter choose one of
the covering neighborhoods, with block length $m\le M$, and write
$n=qm+r$, $0\le r<m$. Subadditivity and $a_r\le r a_1$ give
$a_n\le q\epsilon m+Cr\le\epsilon n+CM$. This proves the
estimate and then the limit. The neighborhoods need not have a common
block length; replacing them by an unproved common mixing time would
not be justified.
\end{proof}

\begin{theorem}[Uniform first-order return process]
\label{thm:uniform-return-fluid}
The actual induced records satisfy
\begin{equation}\label{eq:uniform-return-mean-law}
 \lim_{n\to\infty}\sup_{R\in I}
       \int\left|J_{n,R}/n-\bar G_R\right|\dd\nu_R^*=0.
\end{equation}
For every fixed $A<\infty$ and $\epsilon>0$,
\begin{equation}\label{eq:uniform-return-functional}
 \lim_{n\to\infty}\sup_{R\in I}\nu_R^*\left\{
 \sup_{0\le u\le A}
 \left|n^{-1}J_{\lfloor nu\rfloor,R}-u\bar G_R\right|>\epsilon
 \right\}=0.
\end{equation}
\end{theorem}
\begin{proof}
The ergodic collision map in Section~\ref{sec:return-stability}
induces an ergodic map on its positive-measure return section. Indeed,
the complete tower over an invariant subset of the induced map is an
invariant collision set. The two towers over that subset and its
complement partition the full collision space modulo null sets;
collision ergodicity forces one of the two to have zero measure.
Thus the integrable vector $G_R$ obeys the $L^1$ ergodic theorem for
each fixed $R$, with the exact mean in
Proposition~\ref{prop:exact-induced-means}.

Put $a_n(R)=\int|J_{n,R}-n\bar G_R|\dd\nu_R^*$.
Invariance of $F_R^*$ and the cocycle identity make $a_n$ subadditive.
It is continuous: on the common collision space it equals
\[
 c_*^{-1}\int_M
 \left|\widetilde J_{n,R}-n\bar G_R\one_{Y_R^*}\right|\dd\nu.
\]
Theorem~\ref{thm:return-stability} gives the first term's $L^1$
continuity, the section indicators are continuous in $L^1$ because
their boundaries move continuously and have zero measure, and the
mean is continuous by its explicit formula. The pointwise ergodic
theorem gives $a_n(R)/n\to0$. Lemma~\ref{lem:compact-subadditive}
proves \eqref{eq:uniform-return-mean-law}.

For the functional conclusion, first test a fixed finite mesh of
$[0,A]$. Equation~\eqref{eq:uniform-return-mean-law} and Markov's
inequality give convergence in probability at all its mesh points,
uniformly in $R$. The count and roof coordinates of $J_n$ are
nondecreasing. Between mesh points their normalized deviations are
therefore bounded by the endpoint deviations plus a uniformly bounded
mean times the mesh width, with an $O(n^{-1})$ integer-part error.
For the displacement coordinate, finite horizon gives a constant
$C_\kappa$ such that for every $l\ge j$,
\[
 |K_{l,R}-K_{j,R}|\le C_\kappa(N_{l,R}-N_{j,R}).
\]
Consequently its oscillation on a mesh interval is controlled by the
count increment on that interval. Its mean is zero. On the event
where all mesh errors are small, every coordinate has a uniformly
small error between mesh points as well. Let $n$ grow, then let the
mesh width and the mesh error tolerance decrease. This proves
\eqref{eq:uniform-return-functional} without an independence or
maximal-martingale assumption.
\end{proof}

\begin{lemma}[Sublinear maximum return blocks]
\label{lem:uniform-maximal-block}
Put $Q_R=r_R^*+\tau_R^*+|\kappa_R^*|$ and
\[
 U(L)=\sup_R\int Q_R\one_{\{Q_R>L\}}\dd\nu_R^*.
\]
Then $U(L)\to0$. For $n\ge1$, $A\ge0$, and $\epsilon>0$,
\begin{equation}\label{eq:uniform-maximal-block}
 \sup_R\nu_R^*\left\{
 \max_{0\le j\le\lfloor An\rfloor}Q_R\circ(F_R^*)^j>\epsilon n
 \right\}\le\frac{A+1}{\epsilon}U(\epsilon n).
\end{equation}
Both \eqref{eq:uniform-return-functional} and the convergence to zero
in \eqref{eq:uniform-maximal-block} remain valid when the initial
base point has the stationary length-biased law
\begin{equation}\label{eq:length-biased-base-v4}
 \dd\rho_R=\frac{\tau_R^*}{\bar\tau_R^*}\dd\nu_R^*.
\end{equation}
\end{lemma}
\begin{proof}
Finite horizon bounds $\tau_R^*$ and $|\kappa_R^*|$ by a uniform
constant times $r_R^*$. The uniform integrability in
Theorem~\ref{thm:return-stability}, and the constant section mass,
therefore give $U(L)\to0$. Invariance of $F_R^*$, the union bound
and truncated Markov inequality give
\[
 \nu_R^*\{\max_{j\le\lfloor An\rfloor}Q_R\circ F^j>\epsilon n\}
 \le (\lfloor An\rfloor+1)\nu_R^*(Q_R>\epsilon n)
 \le\frac{A+1}{\epsilon}\int Q_R\one_{Q_R>\epsilon n}\dd\nu_R^*.
\]
No independence of the events is used.

Let $d_*:=\inf_R\bar\tau_R^*>0$. For any event $B_R$ determined
by the forward base orbit and every $L>0$,
\begin{equation}\label{eq:length-bias-transfer-v4}
 \rho_R(B_R)\le d_*^{-1}\left(
 L\nu_R^*(B_R)+\int\tau_R^*\one_{\{\tau_R^*>L\}}\dd\nu_R^*\right).
\end{equation}
If $\sup_R\nu_R^*(B_R)\to0$, first hold $L$ fixed and pass to that
limit, then let $L\to\infty$. Uniform integrability of the roof
makes the right side tend to zero uniformly. Apply this argument to
the exceptional events in the two displayed convergence statements.
This transfer uses a length-biased law, not the unweighted launch law.
\end{proof}

For the stationary physical flow, let $V_R(t)$ count future visits to
$Y_R^*$, let $C_R(t)$ count future physical collisions, and let
$W_R(t)\in\Z^2$ be the net fundamental-cell displacement since time
zero. Changing the choice of a fixed fundamental cell changes the
continuous-position comparison only by a bounded term.

\begin{theorem}[Uniform stationary physical clock]
\label{thm:uniform-physical-clock}
For every $\epsilon>0$, under the actual stationary billiard laws
$\mu_R$, uniformly for $R\in I$,
\begin{equation}\label{eq:uniform-physical-clock}
 \sup_{0\le u\le1}\left|
 \frac1t\bigl(V_R(tu),C_R(tu),W_R(tu)\bigr)
 -\left(\frac{uc_*}{\bar\tau_R},
        \frac{u}{\bar\tau_R},0\right)\right|
 \longrightarrow0
 \quad\hbox{in probability as }t\to\infty.
\end{equation}
The maximum marked unfinished return encountered during $[0,t]$,
including the initial return containing time zero, is $o_P(t)$
uniformly in $R$.
\end{theorem}
\begin{proof}
Represent the starting state by $(x,s)$ in the induced suspension,
$0\le s<\tau_R^*(x)$. The marginal law of $x$ is exactly $\rho_R$
in \eqref{eq:length-biased-base-v4}. Write $T_{j,R}$ for the roof
coordinate of $J_{j,R}$. Count future events in $(0,t]$. With this
convention the following inverse identity also holds at event times:
\[
 V_R(tu)=\max\{j\ge0:T_{j,R}(x)\le s+tu\}.
\]
Choose a fixed $A$ with $A d_*>2$. The functional return law,
transferred by Lemma~\ref{lem:uniform-maximal-block}, shows that
$T_{\lfloor At\rfloor,R}>s+t$ with probability tending to one
uniformly. Here $s/t\to0$ uniformly because $s\le Q_R(x)$ and
the same lemma applies at $j=0$. Integer parts do not affect these
statements; the return-process theorem may be applied with
$n=\lceil t\rceil$ and a slightly larger fixed horizon.

On this event all inverse indices lie below $At$. Uniformly for
$0\le j\le At$ we have
$|T_{j,R}-j\bar\tau_R^*|/t=o_P(1)$ and
$\max_{j\le At}Q_R\circ F^j/t=o_P(1)$. The inequalities
\[
 T_{V_R(tu),R}\le s+tu<T_{V_R(tu)+1,R}
\]
therefore imply
\[
 \sup_{u\le1}\left|
 \bar\tau_R^* V_R(tu)/t-u\right|=o_P(1)
\]
uniformly in $R$. Division by $\bar\tau_R^*\ge d_*$ yields the
first component in \eqref{eq:uniform-physical-clock}.

The completed-return count $N_{V_R(tu),R}$ differs from the actual
physical count by the initial prefix and the final incomplete return.
Their sum is at most a fixed constant plus twice the maximal $Q_R$
over these return indices. The displacement comparison has the same
bound, up to a uniform multiplicative constant, since each physical
lattice step is bounded. These bounds hold simultaneously for all
$u\in[0,1]$. The functional return law at the random inverse indices
now gives
\[
 C_R(tu)/t=\bar r\,V_R(tu)/t+o_P(1),\qquad
 W_R(tu)/t=o_P(1)
\]
with uniform suprema. Since $\bar r/\bar\tau_R^*=1/\bar\tau_R$,
this proves the remaining components. Every unfinished piece in
$[0,t]$ belongs to one of the same return blocks and is bounded by
its marked size, proving the final assertion.
\end{proof}

\begin{corollary}[Geometric error in the physical collision rate]
\label{cor:physical-rate-error}
Let $b_*:=\min_{R\in I}\bar\tau_R$ and put
$L_*=15/(4b_*^2)$. There is a function $\omega_t(\epsilon)\to0$
for every fixed $\epsilon>0$ such that, for any $I$-valued radius
estimate $\widehat R$ on a probability space carrying the stationary
billiard at the true radius $R$, and any $\delta>0$,
\begin{align}\label{eq:physical-rate-error}
 \Pr_R\left\{\sup_{0\le u\le1}
 \left|C_R(tu)/t-u/\bar\tau_{\widehat R}\right|
                       >\epsilon+L_*\delta\right\}
 \le\omega_t(\epsilon)+\Pr_R\{|\widehat R-R|>\delta\}.
\end{align}
No independence of $\widehat R$ and the observed orbit is required.
For the visit rate, replace the deterministic error constant by $c_*L_*$.
\end{corollary}
\begin{proof}
The explicit mean satisfies $|\bar\tau_R'|<15/4$, so
$|(1/\bar\tau_R)'|\le L_*$. On the event
$|\widehat R-R|\le\delta$ the two proposed rates differ by at
most $L_*\delta$, uniformly over $u\in[0,1]$. The triangle
inequality and Theorem~\ref{thm:uniform-physical-clock} prove the
claim with $\omega_t$ equal to the supremum, over $R$, of the true
rate's exceptional probability. Multiplication by $c_*$ gives the
visit-rate statement.
\end{proof}

The supremum statement here is at scale $t$. Uniform integrability
alone does not turn it into an $o_P(\sqrt t)$ bound for a growing
window. Likewise \eqref{eq:subadditive-uniform-budget} and
\eqref{eq:physical-rate-error} do not supply a numerical convergence
rate for $\omega_t$ without quantitative finite-block estimates.
They establish the full first-order clock comparison for this moving
family. The Gaussian covariance and processes require the additional
collision estimates proved in Sections~\ref{sec:collision-covariance}--
\ref{sec:functional-gaussian}; raw density and shrinking-conditioning
conclusions require the further local estimates stated there. The
stationary law throughout is not substituted for an unknown experimental
preparation law.

\section{Bounded collision compensation and its covariance}
\label{sec:collision-covariance}
The discontinuity of the section can be treated on the collision clock
without making its indicator a multiplier on a distribution space. The
argument uses bounded variation in the initial collision coordinates,
smooth approximation, and the collision-space spectral splitting. It
will give a covariance for the actual return record by an exact stopping
identity, rather than by an assumed spectral expansion of an induced
operator.

Write $\eta_R=\one_{Y_R^*}$ and recall
\begin{equation}\label{eq:compensation-normalizations}
 c_* = \frac{91}{10000\pi},\qquad
 \bar\tau_R=\frac{\sqrt3/2-\pi R^2}{2R},\qquad
 \bar G_R=\left(0,0,c_*^{-1},\frac{\bar\tau_R}{c_*}\right).
\end{equation}
Here $\bar G_R=\int G_R\dd\nu_R^*$ is the mean of the actual induced
record, by the Kac identities. On the common collision cylinder put
\begin{equation}\label{eq:bounded-compensation}
 h_R=f_R-\bar G_R\eta_R,
 \qquad s_R=c_*^{-1}\eta_R,
 \qquad S_{m,R}u=\sum_{j=0}^{m-1}u\circ T_R^j.
\end{equation}
Thus $\int h_R\dd\nu=0$ and $s_R\nu=\nu_R^*$, with the latter
measure regarded as a probability on $M$. The four components of $h_R$
are bounded, although those of $G_R$ need not be.

\subsection{Variation of a single physical collision}
For a function on $\T\times(-1,1)$, $\mathrm{BV}$ denotes bounded
variation with respect to the flat coordinates $(\alpha,p)$: both
first distributional derivatives are finite signed measures. All
variation norms below include the $L^1(\nu)$ norm. They are not norms
of the density of a pushed-forward record.

\begin{lemma}[Uniform single-collision variation]\label{lem:physical-bv}
There is a constant $B$ such that
\[
 \|f_R\|_\infty+\|f_R\|_{\mathrm{BV}}
 +\|\eta_R\|_\infty+\|\eta_R\|_{\mathrm{BV}}
 +\|h_R\|_\infty+\|h_R\|_{\mathrm{BV}}
 +\|s_R\|_\infty+\|s_R\|_{\mathrm{BV}}\le B
 \qquad(R\in I).
\]
The norms of vector-valued functions mean the sum of the component norms.
\end{lemma}
\begin{proof}
The uniform horizon bounds all possible next-disk centers in a fixed
finite subset of the triangular lattice. Cover the angular circle by
finitely many bounded rational charts
$t=\tan((\alpha-\alpha_0)/2)$, with uniformly bounded chart derivatives
and inverse derivatives. In such a chart $n_\alpha$ and $t_\alpha$ are
rational functions of $t$. The outgoing velocity
$\sqrt{1-p^2}n_\alpha+pt_\alpha$ is semialgebraic in $(t,p)$. For each
candidate center $d$, the admissibility inequalities and the entry root
\[
 (d-Rn_\alpha)\cdot v
 -\sqrt{R^2-|d-Rn_\alpha|^2+((d-Rn_\alpha)\cdot v)^2}
\]
are semialgebraic in $(R,t,p)$. Taking the least admissible positive
entry root and its lattice label involves only finitely many comparisons.
Choose any fixed tie convention on the singular set. Thus the bounded
single-collision record, extended arbitrarily by zero where no regular
root is defined, is a bounded semialgebraic family on each chart.
Only these one-collision functions, not an arbitrary long itinerary,
are used in this assertion.

We spell out the uniformity implication. A bounded semialgebraic family
$u(\lambda,x)$, with $x$ in a bounded interval, has a uniform finite
number of monotonicity intervals in $x$. To see this, apply differentiable
cell decomposition to its graph with the parameters $\lambda$ ordered
first, and refine by the sign of its fiber derivative where defined.
Uniform finiteness bounds the number of remaining exceptional points
and fiber intervals. On each interval the function is monotone or
constant. These are the cell decomposition and monotonicity results of
\cite[Sections 2.1--2.2 and 6.2]{Coste}. In particular the number does
not depend on the parameter or on the other fixed coordinate of a slice.
If $|u|\le M_0$ and there are at most $b_0$ intervals and exceptional
points, its one-dimensional distributional variation is at most
$C b_0 M_0$, including the jumps.

Apply this statement first with $(R,p)$ as parameters and then with
$(R,t)$ as parameters. Integrating the two slice-variation bounds
against the remaining coordinate proves the bounds for the two
first distributional derivatives. Equivalently, integration by parts
on almost every slice bounds the action on every compactly supported
$C^1$ test function by its supremum norm. A fixed finite chart partition
transfers these bounds to $\alpha$; any jumps at the chart endpoints
have uniformly bounded size and length. This proves the assertion for
$f_R$, including the square-root behavior at grazing.

The section is a fixed finite union of coordinate rectangles with
moving centers. Its perimeter and the number of its sides are uniformly
bounded, so the same assertion for $\eta_R$ follows directly, without
requiring the centers to be semialgebraic in $R$. The explicit means in
\eqref{eq:compensation-normalizations} are bounded on $I$. The remaining
bounds now follow from \eqref{eq:bounded-compensation}.
\end{proof}

\begin{lemma}[Mean-preserving smoothing]\label{lem:bv-smoothing}
There are positive linear operators $\mathsf S_\delta$, $0<\delta<1/2$,
which preserve the integral with respect to $\nu$ and satisfy
\begin{equation}\label{eq:bv-smoothing-bounds}
 \begin{split}
 \|\mathsf S_\delta u\|_\infty&\le\|u\|_\infty,\qquad
 \|u-\mathsf S_\delta u\|_1\le C\delta\|u\|_{\mathrm{BV}},\\
 \|\mathsf S_\delta u\|_{\mathrm{BV}}&\le C\|u\|_{\mathrm{BV}},\qquad
 \|\mathsf S_\delta u\|_{C^2(\alpha,\varphi)}
       \le C\delta^{-2}\|u\|_\infty .
 \end{split}
\end{equation}
The last norm uses $p=\sin\varphi$, as in the collision spaces.
\end{lemma}
\begin{proof}
Reflect $u$ evenly across $p=1$ and $p=-1$ to a function periodic with
period four in $p$, retaining the angular periodicity. Convolve with a
nonnegative, even, smooth product kernel of mass one and scale $\delta$,
then restrict to $-1<p<1$. Reflection symmetry is preserved, so the
integral over this half of the extended cylinder remains the original
integral. The reflected function has variation bounded by a fixed
multiple of the original variation; its traces match at the reflection
boundaries. Translation of a bounded-variation function by a vector $a$
changes its $L^1$ norm by at most $|a|$ times its variation. This follows
first for smooth functions by the fundamental theorem of calculus,
then for distributional derivatives by convolution and passage to the
limit. Integrating this inequality against the kernel gives the $L^1$
error. Convolution contracts variation on the extended cylinder.
The $L^1$ norms of the first two derivatives of the kernel are
$O(\delta^{-1})$ and $O(\delta^{-2})$. These facts prove all the flat
coordinate estimates. Composition with $p=\sin\varphi$ preserves the
stated $C^2$ upper bound, since the first two derivatives of $\sin$
are bounded. No inverse coordinate change at grazing is used.
\end{proof}

\subsection{An exponentially summable collision covariance}
\begin{proposition}[Correlations of bounded-variation observables]
\label{prop:bv-correlation}
For every $B_0<\infty$ there are $C,\gamma>0$, uniform in $R$, such
that centered real functions $a,b$ with
$\|a\|_\infty+\|a\|_{\mathrm{BV}}+
\|b\|_\infty+\|b\|_{\mathrm{BV}}\le B_0$ satisfy
\begin{equation}\label{eq:bv-correlation}
 \left|\int a(b\circ T_R^k)\dd\nu\right|\le Ce^{-\gamma k}
 \qquad(k\ge0).
\end{equation}
If also $\|a\|_1\le C_0\delta$, then for $r=a$ and $b=a$,
\begin{equation}\label{eq:small-bv-variance}
 \int|S_{m,R}r|^2\dd\nu
       \le C m\delta(1+|\log\delta|).
\end{equation}
The second constant may depend on $C_0$ and $B_0$, but not on $m,R$.
\end{proposition}
\begin{proof}
For smooth centered $a,b$, apply Lemma~\ref{lem:collision-spectral-input}
to the distribution $a\nu$ and evaluate the result on $b$ using smooth
multiplication and the mass functional. It gives
\[
 \left|\int a(b\circ T_R^k)\dd\nu\right|
       \le C\vartheta^k\|a\|_{C^2}\|b\|_{C^2}.
\]
Smooth both bounded-variation functions at scale $\epsilon$. Invariance
of $\nu$, boundedness, and \eqref{eq:bv-smoothing-bounds} bound the
change in the correlation by $C\epsilon$. The smooth correlation is
bounded by $C\vartheta^k\epsilon^{-4}$. Choose
$\epsilon=\tfrac14\vartheta^{k/5}$ and enlarge $C$ for $k=0$.
This proves \eqref{eq:bv-correlation}. In particular the constants depend
only on the displayed variation and supremum bounds, not on a modulus
of continuity or on a chosen smoothing scale.

For the last assertion, boundedness gives
$|\int r(r\circ T_R^k)\dd\nu|\le C\delta$ at every lag, while
\eqref{eq:bv-correlation} gives the exponential bound. Hence
\[
 \sum_{k\ge0}\left|\int r(r\circ T_R^k)\dd\nu\right|
       \le C\sum_{k\ge0}\min\{\delta,e^{-\gamma k}\}
       \le C\delta(1+|\log\delta|).
\]
Expanding the square of the stationary sum proves
\eqref{eq:small-bv-variance}.
\end{proof}

Put $h_{R,\delta}=\mathsf S_\delta h_R$. Define, on the collision clock,
\begin{equation}\label{eq:collision-gk}
 \begin{split}
 \Gamma_R={}&\int h_R\otimes h_R\dd\nu\\
 &+\sum_{k=1}^{\infty}\int
 \bigl[h_R\otimes(h_R\circ T_R^k)
       +(h_R\circ T_R^k)\otimes h_R\bigr]\dd\nu .
 \end{split}
\end{equation}
Use the same formula with $h_{R,\delta}$ for $\Gamma_{R,\delta}$.

\begin{theorem}[A continuous collision Green--Kubo matrix]
\label{thm:collision-covariance}
The series \eqref{eq:collision-gk} and its smoothed versions converge
absolutely, with exponential tails uniform in $R$ and $\delta$. The
matrix $\Gamma_R$ is continuous and positive semidefinite on $I$, and
\begin{equation}\label{eq:collision-covariance-limit}
 \Gamma_R=\lim_{m\to\infty}\frac1m
   \int S_{m,R}h_R\otimes S_{m,R}h_R\dd\nu,
 \qquad
 \|\Gamma_{R,\delta}-\Gamma_R\|\le C\delta(1+|\log\delta|).
\end{equation}
The first limit is uniform in $R$.
\end{theorem}
\begin{proof}
Lemma~\ref{lem:physical-bv}, smoothing, and
Proposition~\ref{prop:bv-correlation} give the uniform exponential tails.
For each fixed $k$, the integrand in \eqref{eq:collision-gk} converges
almost everywhere as $R\to R_0$. Indeed, exclude the finitely many
collision singularity preimages and section boundaries encountered up
to time $k$ at $R_0$. On the remaining full-measure set the physical
roots and finite itinerary persist and the moving rectangles have
stable membership. All observables are uniformly bounded. Dominated
convergence proves continuity at each fixed lag; the uniform tail then
proves continuity of the infinite series. This argument establishes
long-time uniformity through the correlation estimate, not through a
fixed-record continuity constant.

Expansion of the square gives the finite covariance with lag weights
$1-k/m$. Absolute summability, uniformly in $R$, proves its limit;
indeed the error is $O(m^{-1})$ by the exponential bound. Positivity
follows from positivity of each finite covariance. At each lag the
change made by smoothing is at most $C\delta$, using invariance and the
$L^1$ approximation. Both the original and the smoothed lag correlations
are at most $Ce^{-\gamma k}$. Summing the minimum of these two bounds
proves the final estimate.
\end{proof}

\begin{remark}\label{rem:collision-gk-scope}
The sum in \eqref{eq:collision-gk} is for the bounded, compensated
collision observable $h_R$ under $T_R$. It is not the Green--Kubo sum
of the unbounded observable $G_R-\bar G_R$ under $F_R^*$. Its relation
to the Gaussian law of the latter will be proved by stopping the actual
collision sums. No bounded-variation estimate for a many-return coarea
density, and no second derivative of that density, follows from the
single-collision variation lemma.
\end{remark}

\section{A uniform Gaussian limit for the actual return record}
\label{sec:stopped-gaussian}
We first obtain a quantitative characteristic-function estimate on the
collision clock. Smooth approximation is used only inside the proof;
the limiting statement concerns the original record and the original
section probability.

\subsection{Smooth perturbation with a shrinking scale}
\begin{lemma}[A quantitative smooth collision expansion]
\label{lem:smooth-collision-expansion}
On each of the collision spaces of
Lemma~\ref{lem:collision-spectral-input}, set
\[
 \mathcal L_{R,\delta,z}
       =\mathcal L_R M_{\exp(i z\cdot h_{R,\delta})},
       \qquad z\in\mathbb C^4.
\]
There are $a,C>0$ and $\rho<1$, independent of $R$ and
$0<\delta<1/2$, such that for $|z|<a\delta^2$ this family has a simple
analytic eigenvalue $\lambda_{R,\delta}(z)$ near one and an analytic
rank-one projection $\Pi_{R,\delta}(z)$. Its complementary powers are
bounded by $C\rho^m$, and, on a smaller such ball,
\begin{equation}\label{eq:smooth-cubic-expansion}
 \begin{split}
 \log\lambda_{R,\delta}(z)
       &=-\tfrac12z^{\mathsf T}\Gamma_{R,\delta}z
                  +O(\delta^{-6}|z|^3),\\
 \|\Pi_{R,\delta}(z)-\Pi_R\|
       &\le C\delta^{-2}|z| .
 \end{split}
\end{equation}
All constants are uniform. The logarithm is the branch vanishing at
$z=0$.
\end{lemma}
\begin{proof}
The smooth multiplier estimate and the smoothing bounds give
$\|\mathcal L_{R,\delta,z}-\mathcal L_R\|
\le C|z|\delta^{-2}$ when $|z|\delta^{-2}$ is small.
The uniform resolvent contours around one and around the complementary
spectrum in the proof of Proposition~\ref{prop:soft-killing} therefore
apply on $|z|<a\delta^2$. The resolvent Neumann series gives the analytic
projection, its difference estimate, and complementary powers on a
fixed circle of radius $\rho<1$. Decrease $a$ so that the eigenvalue
stays in a fixed disk about one not containing zero.

The eigenvalue derivative at zero is zero since $\int h_{R,\delta}
\dd\nu=0$. We identify its second derivative, rather than assuming a
variance formula. Fix a real direction $v$, put $u=v\cdot h_{R,\delta}$,
and normalize a right eigenvector $e(t)$ by $\ell_j(e(t))=1$.
Differentiating the eigenvector equation at zero yields
\[
 e'(0)=i\sum_{k\ge1}\mathcal L_R^k(u\nu),
\]
where the series converges in $\mathcal B_j$ by the spectral splitting.
A second differentiation followed by the mass functional gives
\[
 \lambda''(0)=-\int u^2\dd\nu
       -2\sum_{k\ge1}\int u(u\circ T_R^k)\dd\nu
       =-v^{\mathsf T}\Gamma_{R,\delta}v.
\]
Polarization identifies the full Hessian. Since the first derivative
vanishes, the Hessian of the logarithm is the same. Cauchy's estimates
for the bounded analytic logarithm on a ball of radius $a\delta^2$
bound its third derivatives by $C\delta^{-6}$ on a smaller ball.
Taylor's formula proves \eqref{eq:smooth-cubic-expansion}.
\end{proof}

\begin{proposition}[Uniform collision characteristic estimate]
\label{prop:collision-gaussian}
For every $K<\infty$ there is $C_K$ such that, for $m\ge2$ and for
$\sigma_R=1$ or $\sigma_R=s_R$,
\begin{equation}\label{eq:collision-gaussian}
 \sup_{R\in I,\ |t|\le K}
 \left|\int \sigma_R e^{it\cdot S_{m,R}h_R/\sqrt m}\dd\nu
       -e^{-t^{\mathsf T}\Gamma_Rt/2}\right|
       \le C_K m^{-1/26}\sqrt{\log(2+m)}.
\end{equation}
\end{proposition}
\begin{proof}
Smooth both $h_R$ and the initial density at scale $\delta$. The
smoothed density $\sigma_{R,\delta}=\mathsf S_\delta\sigma_R$ is
nonnegative, has integral one, and satisfies
$\|\sigma_{R,\delta}\nu\|_{\mathcal B_j}\le C\delta^{-2}$.
For $z=t/\sqrt m$, the transfer pairing and the spectral decomposition
of Lemma~\ref{lem:smooth-collision-expansion} give
\[
 \int\sigma_{R,\delta}e^{it\cdot S_{m,R}h_{R,\delta}/\sqrt m}\dd\nu
 =\lambda_{R,\delta}(z)^m
       \ell_j\bigl(\Pi_{R,\delta}(z)\sigma_{R,\delta}\nu\bigr)
       +O(\delta^{-2}\rho^m).
\]
At $z=0$ the amplitude is one; its difference from one is
$O_K(\delta^{-4}m^{-1/2})$. The logarithmic expansion gives, whenever
$K/\sqrt m<a\delta^2$ and $\delta^{-6}m^{-1/2}$ is sufficiently small,
\begin{equation}\label{eq:smoothed-characteristic-error}
 \left|\int\sigma_{R,\delta}
 e^{it\cdot S_{m,R}h_{R,\delta}/\sqrt m}\dd\nu
       -e^{-t^{\mathsf T}\Gamma_{R,\delta}t/2}\right|
 \le C_K\bigl(\delta^{-6}m^{-1/2}
          +\delta^{-4}m^{-1/2}+\delta^{-2}\rho^m\bigr).
\end{equation}
Here positive semidefiniteness of $\Gamma_{R,\delta}$ bounds the
Gaussian factor by one. The estimate for the exponential of the
Taylor remainder follows, for example, from
$|e^w-1|\le |w|e^{|w|}$.

The change of initial density costs at most $C\delta$ in the original
pairing, because the exponential has modulus one. The centered
residual $r_{R,\delta}=h_R-h_{R,\delta}$ has uniformly bounded
supremum and variation norms and $L^1$ norm $O(\delta)$. Therefore
Proposition~\ref{prop:bv-correlation}, component by component, gives
\[
 \int\sigma_{R,\delta}
       |S_{m,R}r_{R,\delta}|^2\dd\nu
       \le C m\delta(1+|\log\delta|),
\]
since the initial density is bounded independently of $\delta$.
The inequality $|e^{ix}-e^{iy}|\le|x-y|$ and Cauchy--Schwarz bound
the change of observable by
$C_K\sqrt{\delta(1+|\log\delta|)}$. Finally,
Theorem~\ref{thm:collision-covariance} bounds the change of the Gaussian
factor by $C_K\delta(1+|\log\delta|)$.

Take $\delta=\tfrac14 m^{-1/13}$. The smoothing error is
$O(m^{-1/26}\sqrt{\log(2+m)})$ and the cubic spectral error is
$O(m^{-1/26})$. The two smallness requirements hold for all sufficiently
large $m$, depending only on $K$. The amplitude and complementary errors
are smaller. Increasing $C_K$ handles the remaining finitely many $m$.
This proves the estimate uniformly over the finite cover of collision
spaces. It does not use an induced spectral gap.
\end{proof}

\subsection{Stopping at the genuine return times}
\begin{lemma}[Exact compensation and a return-clock window]
\label{lem:stopped-compensation}
For $\nu_R^*$-almost every initial state and every $n\ge1$,
\begin{equation}\label{eq:exact-stopped-compensation}
 J_{n,R}-n\bar G_R=S_{N_{n,R},R}h_R.
\end{equation}
There is $C$ such that for $n\ge2$ and $2\le a\le n$,
\begin{equation}\label{eq:return-clock-window}
 \nu_R^*\{|N_{n,R}-n/c_*|>a\}\le Cn/a^2.
\end{equation}
\end{lemma}
\begin{proof}
A trajectory starting in $Y_R^*$ has exactly $n$ visits to this section
among collision times $0,\ldots,N_{n,R}-1$. These are its initial visit
and its first $n-1$ returns. Also
$J_{n,R}=S_{N_{n,R},R}f_R$. Substitution in
\eqref{eq:bounded-compensation} proves the identity with no remainder.

Let $H_{L,R}=\sum_{j=1}^L\eta_R\circ T_R^j$. Its stationary mean is
$c_*L$. Proposition~\ref{prop:bv-correlation} applied to
$\eta_R-c_*$ gives
\[
 \int|H_{L,R}-c_*L|^2\dd\nu\le CL.
\]
The same upper bound, with a changed constant, holds under $\nu_R^*$
because its density relative to $\nu$ is at most $c_*^{-1}$.
The exact equivalences
$N_{n,R}>L\iff H_{L,R}<n$ and
$N_{n,R}\le L\iff H_{L,R}\ge n$ reduce both tails to this estimate.
Choose integers within distance one of $n/c_*+a$ and $n/c_*-a$.
The corresponding deviations from $c_*L$ are at least $c_*a-O(1)$,
and $L\le Cn$. Chebyshev's inequality proves the assertion when $a$
exceeds a fixed constant. For the bounded remaining range, enlarge
$C$ so that the right side is at least one.
\end{proof}

\begin{lemma}[A deterministic-window maximal estimate]
\label{lem:dyadic-collision-maximal}
For $L\ge1$, $r\ge0$, and $R\in I$,
\begin{equation}\label{eq:dyadic-collision-maximal}
 \int\max_{0\le j\le L}
 |S_{r+j,R}h_R-S_{r,R}h_R|^2\dd\nu
       \le CL[\log(2+L)]^2.
\end{equation}
The same estimate holds under $\nu_R^*$ and for backward sums on a
deterministic interval contained in the nonnegative collision times.
\end{lemma}
\begin{proof}
Absolute summability of the correlations gives second moment at most
$C\ell$ for the sum over every interval of $\ell$ collisions, by
stationarity. Enlarge $L$ to a power of two smaller than $2L$.
Every prefix of that interval is a disjoint union of at most one dyadic
interval at each scale. Cauchy--Schwarz bounds the square of its sum by
the number of scales times the sum of the squares of these interval
sums. Bound the latter by the sum over all dyadic intervals at all
scales, independently of the prefix. At each scale, the expectations
sum to at most $CL$, since the intervals partition the full interval.
There are $O(\log(2+L))$ scales, proving the bound. The proof applies
to each component and to intervals in reverse order. A bounded initial
density gives the assertion under $\nu_R^*$ as well.
\end{proof}

\begin{theorem}[Parameter-uniform Gaussian limit of the physical record]
\label{thm:actual-record-gaussian}
Define $D_R=c_*^{-1}\Gamma_R$. It is a continuous positive semidefinite
matrix on $I$. For every $K<\infty$ there is $C_K$ such that
\begin{equation}\label{eq:actual-record-gaussian}
 \sup_{R\in I,\ |t|\le K}
 \left|\int_{Y_R^*}
 e^{it\cdot(J_{n,R}-n\bar G_R)/\sqrt n}\dd\nu_R^*
          -e^{-t^{\mathsf T}D_Rt/2}\right|
       \le C_K n^{-1/26}\sqrt{\log(2+n)}
       \qquad(n\ge2).
\end{equation}
In particular the original, unsmoothed return record has a centered
Gaussian weak limit with covariance $D_R$. The estimate is uniform in
the physical parameter, including along convergent parameter sequences.
\end{theorem}
\begin{proof}
Put $m_n=\lfloor n/c_*\rfloor$ and
$a_n=\lceil n^{3/5}\rceil$. The clock-window lemma, with a harmless
change for rounding, bounds the probability of
$|N_{n,R}-m_n|>a_n$ by $Cn/a_n^2$. On the complementary event, the
difference $S_{N_{n,R},R}h_R-S_{m_n,R}h_R$ is bounded by the larger
of the forward and backward maxima on intervals of length $a_n$ about
$m_n$. The intervals have nonnegative endpoints since $c_*<1$.
The preceding lemma and Cauchy--Schwarz consequently give
\begin{equation}\label{eq:stopped-characteristic-error}
 \begin{split}
 &\left|\int e^{it\cdot S_{N_{n,R},R}h_R/\sqrt n}\dd\nu_R^*
       -\int e^{it\cdot S_{m_n,R}h_R/\sqrt n}\dd\nu_R^*\right|\\
 &\hspace{20mm}\le C_K\left(\frac{n}{a_n^2}
                  +\sqrt{\frac{a_n}{n}}\log(2+a_n)\right)
       \le C_K n^{-1/5}\log(2+n).
 \end{split}
\end{equation}
The exceptional event costs at most twice its probability. No estimate
of an unbounded sum on that event is used.

Apply Proposition~\ref{prop:collision-gaussian} at time $m_n$ and
frequency $t\sqrt{m_n/n}$ with initial density $s_R$. This frequency
stays in a fixed compact set. Since $m_n/n=c_*^{-1}+O(n^{-1})$ and
$\Gamma_R$ is uniformly bounded, its Gaussian differs from the one in
\eqref{eq:actual-record-gaussian} by $O_K(n^{-1})$. The identity
\eqref{eq:exact-stopped-compensation} completes the estimate; the error
in \eqref{eq:stopped-characteristic-error} is smaller than the stated
one. The continuity of $D_R$ follows from
Theorem~\ref{thm:collision-covariance}. The characteristic-function
continuity theorem yields the weak limit, allowing a singular Gaussian.
Along $R_n\to R_0$, the same bound and continuity give the Gaussian
with covariance $D_{R_0}$.
\end{proof}

\subsection{The Gaussian kernel and measurable cohomology}
\begin{theorem}[The kernel is an actual $L^2$ coboundary]
\label{thm:gaussian-kernel-coboundary}
For $R\in I$ and $v\in\R^4$, the following three statements are
equivalent:
\begin{enumerate}
\item $v^{\mathsf T}D_Rv=0$;
\item there is $b\in L^2(M,\nu)$ such that
 $v\cdot h_R=b-b\circ T_R$ almost everywhere;
\item there is $b_*\in L^2(Y_R^*,\nu_R^*)$ such that
 $v\cdot(G_R-\bar G_R)=b_*-b_*\circ F_R^*$ almost everywhere.
\end{enumerate}
If the first condition holds, one may choose $b$ with
$\|b\|_2\le C|v|$, uniformly in $R$, and take $b_*=b|_{Y_R^*}$.
\end{theorem}
\begin{proof}
Put $u=v\cdot h_R$ and $c_k=\int u(u\circ T_R^k)\dd\nu$.
By \eqref{eq:bv-correlation}, $|c_k|\le C|v|^2e^{-\gamma k}$.
If $v^{\mathsf T}D_Rv=0$, then $c_0+2\sum_{k\ge1}c_k=0$.
The exact variance identity gives
\[
 \int|S_{m,R}u|^2\dd\nu
       =-2\sum_{k=1}^{m-1}k c_k-2m\sum_{k\ge m}c_k
       \le C|v|^2 \qquad(m\ge1).
\]
Let $Uu=u\circ T_R$, a unitary operator on $L^2(\nu)$, and form
$b_M=M^{-1}\sum_{m=1}^M S_{m,R}u$. This is a bounded sequence in
$L^2$, with norm at most $C|v|$. Moreover
\[
 (I-U)b_M=u-\frac1M\sum_{m=1}^M U^m u\longrightarrow u
       \quad\hbox{in }L^2.
\]
Here the mean ergodic theorem applies; the collision map is ergodic
and $u$ has mean zero. A weakly convergent subsequence of $b_M$ has
limit $b$, and bounded linearity of $I-U$ shows $(I-U)b=u$.
This proves the second assertion and its norm estimate.

Discard a countable union of collision translates of the null set
where this identity or a representative fails. On the remaining
invariant full-measure set, telescope until the actual first return.
The compensation identity for one block gives
\[
 v\cdot(G_R-\bar G_R)(x)=b(x)-b(F_R^*x).
\]
The restriction belongs to $L^2(\nu_R^*)$ and has norm at most
$c_*^{-1/2}\|b\|_2$. Thus the second assertion implies the third.
Conversely, the third assertion telescopes for the stationary induced
map, so the $L^2(\nu_R^*)$ norm of
$v\cdot(J_{n,R}-n\bar G_R)/\sqrt n$ is at most
$2\|b_*\|_2/\sqrt n$. Its characteristic function tends to one.
Theorem~\ref{thm:actual-record-gaussian} identifies the same limit as
$\exp(-t^2v^{\mathsf T}D_Rv/2)$ for every real $t$, proving the first
assertion. This also completes the equivalence without presupposing
summability of induced correlations.
\end{proof}

\begin{remark}[What the low-frequency theorem supplies]
\label{rem:gaussian-versus-raw}
The matrix $D_R$ is now established by an absolutely convergent
collision Green--Kubo formula and is the covariance of the Gaussian
law of the true induced record. The zero-Gaussian-variance condition
has been identified with a genuine induced $L^2$ coboundary. A
measurable representative of its transfer function has not thereby
been made continuous at a selected periodic orbit. Thus the periodic
rank calculation is not applied to this representative. The separate
phase argument in Theorem~\ref{thm:joint-uniform-nondegeneracy} proves
uniform positive definiteness without periodic evaluation. Neither
convergence of the induced Green--Kubo series nor convergence of its
normalized second moments follows from this stopping argument alone;
the actual second moments and the Cesaro formula are established in
Theorem~\ref{thm:marked-quadratic-moments} and
Corollary~\ref{cor:induced-cesaro-covariance}.

The estimate \eqref{eq:actual-record-gaussian} controls fixed compact
sets of rescaled central frequencies. It does not give the integrated
complementary-frequency bound in the raw local inversion theorem, an
induced eigenvalue expansion, or the complete critical/singular branch
sum. In particular it is not a raw density local limit, a weighted
local-density estimate, or a functional limit under a changed exact
conditioning event. These distinctions retain the original density
problem while supplying its actual Gaussian weak-limit input.
\end{remark}

\section{Functional Gaussian limits and the physical clock}
\label{sec:functional-gaussian}
A central-frequency estimate alone does not give tightness of the record
process. We prove it by smoothing on a slower scale, using fourth moments
for the smooth sums and a maximal second-moment estimate for the residual.
The estimates remain uniform over the moving physical family.

\subsection{Fourth moments and removal of smoothing}
\begin{lemma}[A fourth-moment bound at a specified smoothing scale]
\label{lem:smooth-fourth-moment}
For $m\ge1$, $R\in I$, and $0<\delta<1/2$,
\begin{equation}\label{eq:smooth-fourth-moment}
 \int |S_{m,R}h_{R,\delta}|^4\dd\nu
       \le C\bigl(m^2+(m+1)\delta^{-8}\bigr).
\end{equation}
For $r_{R,\delta}=h_R-h_{R,\delta}$ and $L\ge1$,
\begin{equation}\label{eq:residual-process-maximal}
 \int\max_{0\le j\le L}|S_{j,R}r_{R,\delta}|^2\dd\nu
 \le CL\delta(1+|\log\delta|)[\log(2+L)]^2.
\end{equation}
Both bounds hold, with changed constants, after multiplication by any
nonnegative initial density bounded by a fixed constant.
\end{lemma}
\begin{proof}
For the first assertion consider one real component $u$ of
$h_{R,\delta}$ and the one-variable family in
Lemma~\ref{lem:smooth-collision-expansion}. Write
$\psi(z)=\log\lambda(z)$ and $A(z)=\ell_j(\Pi(z)\nu)$. The spectral
pairing for the stationary characteristic function is
\[
 \int e^{izS_{m,R}u}\dd\nu=e^{m\psi(z)}A(z)+E_m(z).
\]
On a complex disk of radius $a\delta^2$, with $a$ decreased once,
$|E_m(z)|\le C\rho^m$, $|A(z)|\le C$, and $|\psi(z)|\le C$.
Cauchy's inequalities therefore give
$|A^{(j)}(0)|\le C_j\delta^{-2j}$ for $j\le4$,
$|\psi^{(j)}(0)|\le C_j\delta^{-2j}$ for $j=3,4$, and
$|E_m^{(4)}(0)|\le C\delta^{-8}\rho^m$.
In addition $A(0)=1$, $\psi'(0)=0$, and
$|\psi''(0)|\le C$, by the uniformly summable smoothed covariance.
The fourth derivative of the principal term at zero is exactly
\[
 A^{(4)}(0)+6m\psi''(0)A''(0)+4m\psi'''(0)A'(0)
       +m\psi^{(4)}(0)+3m^2\psi''(0)^2.
\]
It is bounded by $C(m^2+(m+1)\delta^{-8})$. Differentiating under the
finite-time integral is legitimate because $u$ is bounded. The fourth
derivative is its fourth moment. Summing the component estimates proves
\eqref{eq:smooth-fourth-moment}.

For the second assertion, \eqref{eq:small-bv-variance} bounds the second
moment of every interval sum of length $\ell$ by
$C\ell\delta(1+|\log\delta|)$. Apply the dyadic proof of
Lemma~\ref{lem:dyadic-collision-maximal} with this constant. The final
assertion follows by domination of the initial density; it does not
assert stationarity under that density.
\end{proof}

For $A_0<\infty$, let $X_{n,R}$ be the polygonal interpolation on
$[0,A_0]$ with values
\[
 X_{n,R}(j/n)=n^{-1/2}S_{j,R}h_R.
\]
Let $\mathcal W_C$ denote Brownian motion starting at zero with covariance
$\mathbb E[\mathcal W_C(s)\mathcal W_C(t)^{\mathsf T}]=(s\wedge t)C$;
$C$ is allowed to be singular. We use the bounded-Lipschitz distance on
probability laws on continuous path space, with its supremum norm.

\begin{theorem}[Uniform collision functional central limit]
\label{thm:collision-functional}
Fix $A_0,B_0<\infty$. Uniformly for $R\in I$ and all nonnegative
probability densities $\sigma_R$ satisfying
$\|\sigma_R\|_\infty+\|\sigma_R\|_{\mathrm{BV}}\le B_0$,
\begin{equation}\label{eq:collision-functional}
 d_{\rm BL}\bigl(\operatorname{Law}_{\sigma_R\nu}(X_{n,R}),
                     \operatorname{Law}(\mathcal W_{\Gamma_R})\bigr)
       \longrightarrow0.
\end{equation}
In particular this holds under $\nu$ and under $\nu_R^*$.
\end{theorem}
\begin{proof}
Set $\delta_n=\tfrac14 n^{-1/32}$ and let $X_{n,R}^{\delta_n}$ be the
polygonal process with $h_R$ replaced by $h_{R,\delta_n}$.
Lemma~\ref{lem:smooth-fourth-moment} implies
\[
 \sup_{R,\sigma_R}\int\|X_{n,R}-X_{n,R}^{\delta_n}\|_\infty^2
             \sigma_R\dd\nu
       \le C_{A_0,B_0}\delta_n[\log(2+n)]^3\longrightarrow0.
\]
Integer parts and interpolation change only the constant. Thus it
suffices to establish the limit for the smoothed process. This scale
is slower than the one optimizing the one-time estimate.

We first prove finite-dimensional convergence. Fix
$0=t_0<t_1<\cdots<t_q\le A_0$ and vectors $v_1,\ldots,v_q$.
Put $m_j=\lfloor nt_j\rfloor-\lfloor nt_{j-1}\rfloor$ and
$z_j=v_j/\sqrt n$. Smooth the initial density at the same scale;
its $L^1$ error is $O(\delta_n)$ and its distribution norm is
$O(\delta_n^{-2})$. The characteristic function of the consecutive
increments, before the harmless polygonal endpoint corrections, is
\[
 \ell_j\bigl(\mathcal L_{R,\delta_n,z_q}^{m_q}\cdots
       \mathcal L_{R,\delta_n,z_1}^{m_1}
                   (\mathsf S_{\delta_n}\sigma_R)\nu\bigr).
\]
The product is chronological, with the earliest segment on the right.
For each fixed $q$, substitute the spectral decompositions of
Lemma~\ref{lem:smooth-collision-expansion}. Every term with a
complementary power tends to zero uniformly, since $m_j$ is a positive
multiple of $n$ and its remaining factors are polynomially bounded in
$\delta_n^{-1}$. The product of principal amplitudes tends to one:
$\Pi(0)=\nu\ell_j$, and its error is
$O_q(\delta_n^{-4}n^{-1/2})$. The cubic error in the logarithm is
$O_q(\delta_n^{-6}n^{-1/2})=o(1)$. The resulting characteristic function
is therefore, uniformly in $R$ and in the initial densities,
\[
 \exp\left(-\frac12\sum_{j=1}^q(t_j-t_{j-1})
            v_j^{\mathsf T}\Gamma_Rv_j\right)+o(1).
\]
Here \eqref{eq:collision-covariance-limit} removes smoothing from the
covariance. This is the joint law of the independent Brownian
increments, not merely a collection of marginal Gaussian laws.

For tightness put $L_n=\lceil\delta_n^{-8}\rceil$ and interpolate the
smoothed sums only at collision indices $0,L_n,2L_n,\ldots$.
Call this coarse polygonal process $Z_{n,R}$. An interval of $m\ge L_n$
collisions has fourth moment at most $Cm^2$, by
\eqref{eq:smooth-fourth-moment} and stationarity of $\nu$. Domination
by $\sigma_R\le B_0$ gives the same bound from every deterministic
starting index under the initial law. It follows at coarse grid points
that the fourth moment of an increment is at most $C|t-s|^2$.
For arbitrary points, split an increment into its two fractional end
segments and the complete grid interval. Within one grid interval,
linear interpolation gives a factor $(|t-s|n/L_n)^4$ multiplying a
fourth moment $O((L_n/n)^2)$, which is bounded by $C|t-s|^2$.
The same estimate follows across adjacent intervals and then across
any number of intervals, with a fixed constant. Hence
\[
 \mathbb E_{\sigma_R\nu}|Z_{n,R}(t)-Z_{n,R}(s)|^4
       \le C_{A_0,B_0}|t-s|^2.
\]
The usual dyadic proof of the continuity and tightness criterion applies
uniformly to these continuous processes. Extend the last grid interval
past $A_0$ and restrict back to $[0,A_0]$ to avoid an endpoint convention.
Finally, boundedness of $h_{R,\delta_n}$ gives deterministically
\[
 \|Z_{n,R}-X_{n,R}^{\delta_n}\|_\infty
       \le C L_n/\sqrt n=O(n^{-1/4})\longrightarrow0.
\]
This proves tightness of the original processes as well.

To obtain the stated uniform distance, suppose it failed and select a
violating sequence of radii and initial densities. By compactness a
subsequence of the radii converges to $R_0$. The uniform tightness and
finite-dimensional limits identify every subsequential path law as
$\mathcal W_{\Gamma_{R_0}}$. Continuity of $\Gamma_R$, including at
singular matrices, gives continuity of the Brownian law; for example
couple by $\Gamma_R^{1/2}$ times one standard Brownian motion. This
contradicts the violating bounded-Lipschitz distance and proves
\eqref{eq:collision-functional}.
\end{proof}

\subsection{The induced record process}
Let $Y_{n,R}$ be the polygonal interpolation on $[0,1]$ of
$n^{-1/2}(J_{j,R}-j\bar G_R)$ at times $j/n$.

\begin{theorem}[Functional limit of the actual return process]
\label{thm:actual-return-functional}
Uniformly for $R\in I$,
\begin{equation}\label{eq:actual-return-functional}
 d_{\rm BL}\bigl(\operatorname{Law}_{\nu_R^*}(Y_{n,R}),
                         \operatorname{Law}(\mathcal W_{D_R})\bigr)
       \longrightarrow0,
       \qquad D_R=c_*^{-1}\Gamma_R.
\end{equation}
\end{theorem}
\begin{proof}
The count coordinate of Theorem~\ref{thm:uniform-return-fluid} gives
\[
 \sup_{0\le s\le1}|N_{\lfloor ns\rfloor,R}/n-s/c_*|
                  \longrightarrow0
\]
in probability, uniformly in $R$. Choose $A_0>c_*^{-1}+1$ and apply
Theorem~\ref{thm:collision-functional} on $[0,A_0]$ under $\nu_R^*$.
Tightness with continuous limits gives a uniform probabilistic modulus
of continuity for $X_{n,R}$. Consequently replacing its random argument
$N_{\lfloor ns\rfloor,R}/n$ by $s/c_*$ changes it by $o_P(1)$ in
supremum norm. The random argument remains in $[0,A_0]$ with probability
tending to one uniformly; independence from $X_{n,R}$ is not required.
The exact identity \eqref{eq:exact-stopped-compensation} identifies this
composition with the step version of the centered induced process.

For its polygonal version, stationarity under $F_R^*$ and the true
one-block tail give, for every $\epsilon>0$,
\[
 \nu_R^*\left\{\max_{j\le n}r_R^*\circ(F_R^*)^j>
                  \epsilon\sqrt n\right\}
       \le C(n+1)e^{-c\epsilon\sqrt n}\longrightarrow0.
\]
Each centered return increment has norm at most $C r_R^*$, by
boundedness of $h_R$ and exact compensation. Thus interpolation changes
the process by $o_P(1)$ uniformly. Brownian time change $s\mapsto s/c_*$
has covariance $c_*^{-1}\Gamma_R$, which proves the assertion by the
same compactness argument as above.
\end{proof}

\subsection{An unconditioned physical-time consequence}
Recall the actual stationary physical collision count $C_R(t)$ and
lattice displacement $W_R(t)$ of
Theorem~\ref{thm:uniform-physical-clock}. Define the linear map
\[
 B_R(x_1,x_2,x_3,x_4)=(x_1,x_2,x_3-x_4/\bar\tau_R),\qquad
 \mathcal V_R=\bar\tau_R^{-1}B_R\Gamma_R B_R^{\mathsf T}.
\]
\begin{corollary}[Physical displacement and collision fluctuations]
\label{cor:physical-functional-gaussian}
Under the stationary laws $\mu_R$, the processes
\begin{equation}\label{eq:physical-functional-gaussian}
 t^{-1/2}\bigl(W_R(ts),C_R(ts)-ts/\bar\tau_R\bigr),\quad 0\le s\le1,
\end{equation}
converge uniformly in $R$ to Brownian motion with covariance $\mathcal V_R$.
The convergence may be taken in the Skorokhod space, or for continuous
interpolations whose uniform difference from the displayed processes
is $O(t^{-1/2})$. The matrices $\mathcal V_R$ are continuous and positive
semidefinite. No conditioning event is changed or imposed in this result.
\end{corollary}
\begin{proof}
In the physical collision suspension, the marginal of the last outgoing
collision state $x$ is $\sigma_R\nu$, where
$\sigma_R=\tau_R/\bar\tau_R$. Its supremum and variation norms are
uniformly bounded by Lemma~\ref{lem:physical-bv} and
$\min_I\bar\tau_R>0$. The elapsed time since that collision is bounded
by the uniform single-flight horizon. Thus
Theorem~\ref{thm:collision-functional} applies with this initial density.

The section compensation disappears under $B_R$:
\[
 B_Rh_R=(\kappa_{{\rm coll},1},\kappa_{{\rm coll},2},
                    1-\tau_R/\bar\tau_R),
 \qquad B_R\bar G_R=0.
\]
The collision clock satisfies
$\sup_{s\le1}|C_R(ts)/t-s/\bar\tau_R|=o_P(1)$ uniformly, by
Theorem~\ref{thm:uniform-physical-clock}. Apply the continuous-limit
random-time argument to $B_RX_{\lfloor t\rfloor,R}$ on a fixed interval
larger than $1/\min_I\bar\tau_R+1$. Rounding $t$ affects neither the
normalization nor the limit.

The initial and final partial physical flights have uniformly bounded
durations. The collision sum of the lattice labels differs from
$W_R(ts)$ by a uniformly bounded cell-position term. Also the sum of
completed flight times differs from $ts$ by a uniformly bounded amount.
Consequently the stopped sum of $B_Rh_R$ differs from the numerator in
\eqref{eq:physical-functional-gaussian} by a uniformly bounded vector,
simultaneously for all $s$. The jumps of the lattice and collision counts
are uniformly bounded; interpolation between collision times, together
with the same cell-position bound, has error $O(t^{-1/2})$. The limiting
Brownian time change is $s/\bar\tau_R$, giving $\mathcal V_R$. Continuity follows
from that of $\Gamma_R$ and the explicit strictly positive mean roof.
\end{proof}

The functional results supply the unconditioned process input for the
original clock analysis. They neither evaluate a measurable transfer
function at the selected periods nor control a shrinking conditioning
event. Uniform positive definiteness is supplied separately by
Corollary~\ref{cor:actual-uniform-ellipticity}; the weighted raw density
estimates remain distinct from these path-space conclusions.

\section{One-step phase coercivity and the operator criterion}
\label{sec:operator-bridge}
The positive-measure estimate in Section~\ref{sec:periodic-coercivity}
uses an entire periodic tube. There is a stronger estimate if one first
locates a large one-step defect on the period, and only then thickens
that single state. It avoids the exponentially shrinking tube of the
whole period. We retain the constants $\rho_0,L,B_0$ of
Lemma~\ref{lem:uniform-period-tubes}, and set $r_0=\rho_0/(2L)$.

\begin{theorem}[One-step thickening of the periodic defect]
\label{thm:one-step-coercivity}
Fix $0<\alpha\le1$ and $1\le p<\infty$. Let $q$ be circle-valued and
measurable, with an $\alpha$-H\"older representative of constant $H\ge1$
on the $\rho_0$ balls around the states of the triple selected in
Theorem~\ref{thm:uniform-phase-separation}. For $b=|\beta|\ge1$ put
\[
 d_b=\frac1{4000000\pi M(b)},\qquad
 H_b=H(1+L^\alpha),\qquad
 r_b=\min\{r_0,d_b/(4bB_0),(d_b/(4H_b))^{1/\alpha}\}.
\]
For the actual one-step defect $\mathcal D_q$ defined in
Theorem~\ref{thm:positive-measure-coercivity},
\begin{equation}\label{eq:one-step-coercivity}
 \|\mathcal D_q\|_{L^p(\nu_R^*)}
       \ge\frac{d_b}{2}\left(\frac{r_b^2}{4c_*}\right)^{1/p}.
\end{equation}
If $H\le H_0b^\zeta$ with fixed $H_0\ge1$ and $\zeta\ge0$, this gives
\begin{equation}\label{eq:sharper-phase-power}
 \|\mathcal D_q\|_{L^p(\nu_R^*)}
 \ge C\,b^{-\lambda}(1+\log b)^{-\eta},\qquad
 \lambda=\frac2p\max\{1,\zeta/\alpha\},\quad
 \eta=1+\frac2{\alpha p}.
\end{equation}
The constants are uniform in the radius and all torus frequencies.
\end{theorem}
\begin{proof}
Telescoping the three exact periods as in
Corollary~\ref{cor:logarithmic-coercivity} shows that at some state $x_0$
of the selected triple the one-step defect has modulus at least $d_b$.
This is initially a statement about the specified continuous
representatives at finitely many states: the sum of the three induced
periods is at most $2M(b)$, while their product discrepancy is at least
$2/(4000000\pi)$.

On $B(x_0,r_0)$ the return has one fixed lattice label and one fixed
physical count, and $F_R^*$ maps that ball into the $\rho_0$ ball about
the successor of $x_0$. For $x$ in that ball,
\[
 |\mathcal D_q(x)-\mathcal D_q(x_0)|
 \le H_b|x-x_0|^\alpha+bB_0|x-x_0|.
\]
Each term on the right is at most $d_b/4$ on the ball of radius
$r_b$. On that ball the defect
is consequently at least $d_b/2$. Its actual invariant measure is
$r_b^2/(4c_*)$. Integration proves \eqref{eq:one-step-coercivity}.
Changes of representatives on null sets do not alter the integral.

Finally $M(b)\le3+\log b/\log(100/49)$ and
$H_b\le C b^\zeta$. The minimum defining $r_b$ is bounded
below by a positive constant times
$b^{-\max\{1,\zeta/\alpha\}}(1+\log b)^{-1/\alpha}$.
Substitution proves \eqref{eq:sharper-phase-power}. Unlike the
whole-period tube estimate, its exponent contains no $\log L$ term.
The regularity premise on $q$ remains essential.
\end{proof}

The next theorem specifies exactly what is required to turn this
physical coercivity into an operator resolvent. The hypothesis is a
reconstruction of a phase from an approximate spectral vector; it is
not silently identified with an eigenvalue equation.

\begin{theorem}[A quantitative phase-reconstruction criterion]
\label{thm:phase-resolvent-criterion}
For each $R,u,s,\beta$, let $\mathcal L_{R,u,s,\beta}$ be a bounded
operator on a Banach space $\mathcal B_R$. Fix the constants
$\alpha,p,H_0,\zeta$ of Theorem~\ref{thm:one-step-coercivity}, and let
$A\ge1$, $v\ge0$, $\theta>0$. Suppose that for every $b=|\beta|\ge1$,
$|z|=1$, and $h\in\mathcal B_R$ with $\|h\|=1$ and
$e=\|(z-\mathcal L_{R,u,s,\beta})h\|\le1$, one can construct a
circle-valued $q$ with the local H\"older bound $H_0b^\zeta$ and
\begin{equation}\label{eq:phase-reconstruction}
 \|\mathcal D_q\|_{L^p(\nu_R^*)}\le A b^v e^\theta,
 \qquad e^{ic}=z.
\end{equation}
Suppose also that $z-\mathcal L_{R,u,s,\beta}$ is Fredholm of index
zero. Then it is invertible and
\begin{equation}\label{eq:phase-resolvent}
 \|(z-\mathcal L_{R,u,s,\beta})^{-1}\|
 \le C b^{(\lambda+v)/\theta}
                (1+\log b)^{\eta/\theta},
\end{equation}
with $\lambda,\eta$ as in \eqref{eq:sharper-phase-power}.
\end{theorem}
\begin{proof}
Combining \eqref{eq:sharper-phase-power} and
\eqref{eq:phase-reconstruction} yields
\[
 e\ge (C_0/A)^{1/\theta}
       b^{-(\lambda+v)/\theta}(1+\log b)^{-\eta/\theta}
\]
when $e\le1$. Decreasing the positive constant to at most one makes
the same lower bound valid when $e>1$. Homogeneity gives a uniform
lower bound for $z-\mathcal L$ on every vector. It is therefore
injective and has closed range. Fredholm index zero makes its cokernel
zero as well, so it is surjective. The lower bound gives
\eqref{eq:phase-resolvent}. The Fredholm assumption cannot be omitted:
an injective operator bounded below need not be onto.
\end{proof}

For the unbounded induced twists, the phase reconstruction
\eqref{eq:phase-reconstruction}, including nonvanishing modulus and
quantitative regularity, is an unverified premise, not a result of
Lemma~\ref{lem:collision-spectral-input}. The latter treats a smooth
collision-space weight near zero. It provides the return-tail theorem
but is not a construction of $\mathcal B_R$ or of the raw high-frequency
operator. The criterion separates the necessary analytical construction
from the now explicit phase lower bound.

\paragraph{Scope of the retained edge calculation.}
The next section uses the original section $Y_R$ in \eqref{eq:Y}, not
$Y_R^*$. Its pointwise physical critical-word calculations are retained,
but an induced count or normalization is not transferred between the
sections without a new return computation.
\section{The exact raw density at a minimal flight length}
Let $L_N(x)=S_N\tau_R(x)$ on the collision section, and put
\[
 g=1-2R,\qquad \zeta_R=\operatorname{arcosh}(1+g/R)>0.
\]
Near the orbit alternating normally between centers $0$ and $a$, choose the initial arc coordinate $y$ in the upward tangent direction and the outgoing tangential velocity $p$. Both vanish on the orbit. Use the same upward signed arc coordinate $y_j$ at its successive endpoints.

\begin{lemma}[Reduced action Hessian]\label{lem:hessian}
The second variation of the length with all contact coordinates free is
\begin{equation}\label{eq:quadratic}
 \frac12\sum_{j=0}^{N-1}
 \left\{\frac{y_j^2+y_{j+1}^2}{R}
             +\frac{(y_{j+1}-y_j)^2}{g}\right\}.
\end{equation}
After solving the internal reflection equations, the endpoint Hessian is
\begin{equation}\label{eq:schur}
 S_N=\begin{pmatrix}A_N&-B_N\\-B_N&A_N\end{pmatrix},\quad
 A_N=\frac{\sinh\zeta_R}{g}\coth(N\zeta_R),\quad
 B_N=\frac{\sinh\zeta_R}{g\sinh(N\zeta_R)}.
\end{equation}
In the physical initial coordinates $(y,p)$, $L_N$ has a nondegenerate minimum $Ng$ and Hessian $H_N$ satisfying
\begin{equation}\label{eq:detH}
             \sqrt{\det H_N}=\sinh(N\zeta_R).
\end{equation}
\end{lemma}
\begin{proof}
At successive facing circles the longitudinal endpoint corrections are $y_j^2/(2R)$ and $y_{j+1}^2/(2R)$. Expanding their distance gives \eqref{eq:quadratic}. Its full Hessian is positive definite, since it is a sum of positive square terms including endpoint squares. The interior stationarity equations are
\[
 y_{j+1}-2(1+g/R)y_j+y_{j-1}=0.
\]
Their solution with given endpoints is
\[
 y_j=\frac{\sinh((N-j)\zeta_R)}{\sinh(N\zeta_R)}y_0
       +\frac{\sinh(j\zeta_R)}{\sinh(N\zeta_R)}y_N.
\]
Substitution into the endpoint gradients yields \eqref{eq:schur}. The nonlinear interior equations are uniquely solvable near zero by their positive definite Hessian. Their solution is the physical reflected path by the first variation and the strict incidence at the normal orbit.

The endpoint first variation gives $p=-\partial_{y_0}L_N$ at the linearized level, so $p=-A_Ny_0+B_Ny_N$. The linear map from $(y_0,p)$ to $(y_0,y_N)$ has determinant $1/B_N$. Thus
\[
 \det H_N=\frac{A_N^2-B_N^2}{B_N^2}
                         =\sinh^2(N\zeta_R).
\]
Positivity is preserved under this nonsingular coordinate change, proving the minimum assertion.
\end{proof}

\begin{theorem}[A physical edge coefficient]\label{thm:edge}
Let $\chi$ be a smooth nonnegative section cutoff supported in the regular alternating itinerary, equal to one near its normal initial point. The pushforward of $\chi\nu$ by $L_N$ has, near $Ng$, the form
\begin{equation}\label{eq:edge}
 f_{N,R}(Ng+s)=\one_{\{s\ge0\}}a_{N,R}(s),\qquad
 a_{N,R}(0)=J_{N,R}:=\frac1{2R\sinh(N\zeta_R)}.
\end{equation}
The function $a_{N,R}$ is smooth on a right neighborhood of zero. For each fixed $N$, the construction and smooth bounds may be made uniform on $I$. No assertion of uniform derivative bounds as $N\to\infty$ is made.
\end{theorem}
\begin{proof}
The section density in initial arc and tangential velocity coordinates is $(4\pi R)^{-1}\dd y\dd p$. The Morse lemma applied to Lemma~\ref{lem:hessian} gives $L_N=Ng+|z|^2/2$. The transformed density $w(z)$ is smooth with
\[
 w(0)=\frac1{4\pi R\sqrt{\det H_N}}.
\]
In polar coordinates its pushforward density is
$\one_{\{s\ge0\}}\int_0^{2\pi}w(\sqrt{2s}\,n_\phi)\dd\phi$ near zero. The angular integral of each odd Taylor term vanishes; Taylor's theorem at arbitrary order gives smoothness in $s\ge0$. Its right value is $2\pi w(0)$, which is \eqref{eq:edge}. Compactness in $R$ and the strictly positive determinant give the fixed-$N$ uniform statement.
\end{proof}

\begin{proposition}[Regular critical words and a uniform edge bound]
\label{prop:general-edge}
Fix a regular physical word with $N$ flights between circular obstacles. Every critical point of its total flight length in the two initial section coordinates is normal at both endpoints. There is at most one such point for the fixed sequence of lattice centers. It is a nondegenerate minimum. If $S=\left(\begin{smallmatrix}A&C\\C&D\end{smallmatrix}\right)$ is its reduced endpoint action Hessian, its section-density jump is
\begin{equation}\label{eq:general-edge}
 J=\frac{|C|}{2R\sqrt{AD-C^2}}.
\end{equation}
There is a fixed explicitly computable constant $C_*$ such that, for every regular word and $N\ge2$,
\begin{equation}\label{eq:edge-decay}
                  0<J\le C_*(47/53)^{N-2}.
\end{equation}
No lower incidence margin for the interior collisions is needed for this coefficient bound. This is an individual-word estimate, not a bound for the sum over words.
\end{proposition}
\begin{proof}
Let $q_j(y_j)$ be arc coordinates and $v_j$ the unit velocity of segment $j$. Its second length variation is
\[
 \frac{c_j}{R}y_j^2+\frac{c_{j+1}}{R}y_{j+1}^2+
 \frac{\{v_j^\perp\cdot(t_{j+1}y_{j+1}-t_jy_j)\}^2}{\ell_j},
\]
where $c_j,c_{j+1}>0$ are its outgoing and incoming incidence cosines. At an interior reflected contact the two cosines agree. Summing gives a positive definite tridiagonal matrix with nonzero off-diagonal entries. Its interior Schur complement $S$ is positive definite and $C\ne0$: eliminating the successive interior variables expresses $C$ as the nonzero product of neighboring entries divided by positive pivots.

The endpoint first variation of the reduced action is $-p_0\dd y_0+p_N\dd y_N$. Moreover $\partial_{p_0}y_N=-1/C\ne0$. Thus the derivative of the physical length in $p_0$ vanishes exactly when $p_N=0$, and its other derivative then vanishes exactly when $p_0=0$. At such a critical point the coordinate change gives $\det H=(AD-C^2)/C^2>0$. The same Morse calculation as in Theorem~\ref{thm:edge} proves \eqref{eq:general-edge}.

There cannot be two such regular points for one word. Regard the polygonal length as a convex function of all contact positions in the product of the closed disks. At a normal endpoint its gradient is $-n$, and at an interior reflected contact it is $-2c_j n_j$. These are strictly inward normals. For a different point $x_j$ of the same strictly convex disk, $n_j\cdot(x_j-q_j)<0$. The convex gradient inequality therefore makes every different tuple have strictly larger action. Hence the normal-to-normal critical tuple is the unique global minimizer of that convex problem. Occlusion by a disk not in the word is excluded separately by physical admissibility, not by this minimization argument.

To prove the uniform bound, set $z_j=c_jy_j$ and choose tangent orientations successively so that the off-diagonal entries become negative. At the normal endpoints $c_0=c_N=1$. Write $t_j=1/\ell_j\le t_*=50/3$. The interior matrix has off-diagonals $-t_j$ and diagonals
\[
         d_j=t_{j-1}+t_j+2/(Rc_j)\ge t_{j-1}+t_j+v_*,
             \qquad v_*=200/47.
\]
Its factorization $D_0(I-K)$ has $K\ge0$ with row sums at most
$q=2t_*/(2t_*+v_*)=47/53$. In the Neumann series, an entry connecting interior indices $1$ and $N-1$ vanishes before power $N-2$. Therefore
\[
 |C|\le t_*^2\frac{q^{N-2}}{v_*(1-q)}.
\]
The reduced quadratic form dominates the two endpoint curvature terms, so $S\ge (100/47)I$. Inserting these estimates in \eqref{eq:general-edge}, and using $R\ge9/20$, gives \eqref{eq:edge-decay} with
\[
 C_* = \frac{47}{90}\frac{t_*^2}{v_*(1-q)}.
\]
Arbitrarily small positive interior incidence only increases the diagonal potential, so it does not weaken the estimate.
\end{proof}

The proposition classifies regular critical contributions and bounds each jump uniformly. Singular itinerary boundaries and the number and arrangement of critical values remain to be controlled. Multiplying \eqref{eq:edge-decay} by an uncontrolled word count would not prove a summable global correction.

\begin{corollary}[Raw finite-record absolute continuity]\label{cor:raw-density}
For every $R\in I$ and $n\ge1$, the actual induced record
$(S_n\kappa_R,S_n r_R,S_n\tau_R^Y)$ under $\nu_R^Y$ has a density with respect to counting measure on $\Z^3$ times Lebesgue measure. For fixed total physical count $N$ and displacement $k$, it has the coarea representation, for almost every $t$,
\begin{equation}\label{eq:raw-coarea}
 p_{n,R}(k,N,t)=\frac1{4\pi\nu(Y_R)}
 \sum_{\mathbf d}\int_{D_{n,\mathbf d}\cap\{L_N=t\}}
              \frac{\dd\mathcal H^1}{|\nabla_{\alpha,p}L_N|}.
\end{equation}
Here $D_{n,\mathbf d}$ is the actual first-return event for the ordered center word $\mathbf d$, total count $N$ and displacement $k$. Bounded measurable insertions have the same formula with the insertion in the numerator. Neither a uniform density bound nor a long-time local limit is asserted by this statement.
\end{corollary}
\begin{proof}
Partition the recurrent regular initial states by total collision count and ordered center word. Uniform finite horizon makes the available lattice steps finite; for fixed $N$ there are finitely many words. On every regular branch the physical length is smooth, and Proposition~\ref{prop:general-edge} makes its critical set empty or a singleton. Away from that null set, the coarea formula applied on a countable exhaustion where $|\nabla L_N|$ is bounded below proves absolute continuity and \eqref{eq:raw-coarea}. Restriction to the measurable first-return event does not alter absolute continuity of the initial measure. The singular initial states have zero flux measure, and summing over the countable possible $N$ gives the full law. The same exhaustion with a bounded insertion proves its signed version.
\end{proof}

For the induced system \eqref{eq:Y}, fix $n$ and restrict its joint law to $\kappa=0$, $r=2n$. Every physical flight is at least $g$. Equality of total length with $2ng$ requires every flight to have length $g$, and hence a nearest-neighbor normal alternating orbit. Of the six possible normal starting states only $(0,0)$ lies in $Y_R$. A neighborhood of this point makes exactly $n$ returns of physical length two each. Compactness, or directly the equality cases of the distance bound, excludes other itineraries from a sufficiently small right neighborhood of $2ng$. Consequently this lattice component of the induced law has right density
\begin{equation}\label{eq:inducededge}
                \frac{J_{2n,R}}{\nu(Y_R)}>0
\end{equation}
at its left endpoint, and zero density on its left.

\begin{corollary}[The boundary term must be retained]\label{cor:nonLone}
For the specified inducing set, the raw joint characteristic function $\Phi_{n,R}$ is not in $L^1(\T^3\times\R)$, for any $n\ge1$. In particular a global integrable bound for this raw function cannot be justified merely by making the normal itinerary exponentially rare.
\end{corollary}
\begin{proof}
If it were integrable, taking its $(0,0,2n)$ lattice Fourier coefficient would give an integrable function of the roof frequency. Its inverse would be a continuous density of that component. A continuous representative cannot agree almost everywhere with zero to the left of $2ng$ and with a function having the positive right limit \eqref{eq:inducededge}. This contradicts Theorem~\ref{thm:edge}.
\end{proof}

The central local-limit problem is not decided by this corollary. The endpoint may be far outside the central window, and $J_{2n,R}$ decays exponentially. What changes is the global inversion method, not the raw observable or the intended central theorem.

\begin{proposition}[Exact edge subtraction]\label{prop:subtract}
Choose a smooth roof cutoff supported in the right neighborhood of Theorem~\ref{thm:edge}, equal to one near its endpoint, and denote the resulting subdensity by $f$. With $J=J_{N,R}$, one has the exact identity
\begin{equation}\label{eq:subtract}
 f(Ng+s)=J e^{-s}\one_{\{s\ge0\}}+q(s),\qquad
 \widehat f(b)=e^{ibNg}\left(\frac{J}{1-ib}+\widehat q(b)\right),
\end{equation}
where $q\in L^1(\R)$, its distributional second derivative is a finite signed measure, and
$|\widehat q(b)|\le C_{N,R}(1+b^2)^{-1}$.
\end{proposition}
\begin{proof}
Writing $f(Ng+s)=\one_{s\ge0}\eta(s)a(s)$, the right value of $\eta a-Je^{-s}$ at zero is zero. Thus $q$ is continuous across zero, piecewise smooth, and exponentially decaying at infinity. Its first derivative may jump at zero, so its second distributional derivative has one finite atom and an integrable density. Distributional integration by parts gives $b^2|\widehat q(b)|\le\|D^2q\|_{\TV}$; the $L^1$ norm gives the bound near zero. Direct integration of the exponential proves \eqref{eq:subtract}.
\end{proof}
This is subtraction and exact addition of an explicit density, not convolution with noise or removal of a physical event. It proves an integrable remainder for the isolated edge. Bounds for all remaining branches and edges are still needed for a long-time result.


\section{A raw residual-inversion theorem}

\paragraph{Frequency convention.}
The rescaled frequency is $v=\sqrt n\,\omega$. When the proved growing
central band is used below, its radii are $2n^{1/200}$ in $v$ and
$2n^{-99/200}$ in $\omega$. General cutoffs in the following statements
retain their stated hypotheses; a Gaussian-tail estimate is not a bound
for the complementary transform of the raw measure.

Let $\eta_{n,R}$ be finite measures on $\Z^3\times\R$ with characteristic functions $\Phi_{n,R}$. Frequencies are $\omega=(u,s,b)\in[-\pi,\pi]^3\times\R$, and our transform convention is $e^{i\omega\cdot z}$. Suppose there are explicit signed edge measures $E_{n,R}$ with densities $e_{n,R}$, and put $H_{n,R}=\Phi_{n,R}-\widehat E_{n,R}$. The following hypotheses are an interface, not consequences of the mechanical constructions above.

\begin{theorem}[Residual criterion for the mixed density LLT]\label{thm:LLT}
Assume the following uniformly for $R\in I$.

\smallskip\noindent
(i) There are $c,C,\delta>0$, $\rho<1$, means $m_R\in\R^4$, and symmetric matrices $cI\le\Sigma_R\le CI$ such that, for $|\omega|\le\delta$,
\[
 \Phi_{n,R}(\omega)=e^{inm_R\cdot\omega}
 \{A_R(\omega)\lambda_R(\omega)^n+O(\rho^n)\},
\]
where $A_R(0)=1$, $|A_R(\omega)-1|\le C|\omega|$,
$\log\lambda_R(\omega)=-\tfrac12\omega^T\Sigma_R\omega+O(|\omega|^3)$, and $|\lambda_R(\omega)|\le e^{-c|\omega|^2}$.

\smallskip\noindent
(ii) $H_{n,R}\in L^1(\T^3\times\R)$ and
\begin{equation}\label{eq:tailcriterion}
 \int_{|\omega|>n^{-2/5}}|H_{n,R}(\omega)|\dd\omega=o(n^{-2}).
\end{equation}

\smallskip\noindent
(iii) $\|E_{n,R}\|_{\TV}=o(n^{-2})$ and, for each fixed $K$,
\[
 n^2\mathop{\rm ess\,sup}_{|(z-nm_R)/\sqrt n|\le K}|e_{n,R}(z)|\longrightarrow0.
\]
Then $\eta_{n,R}$ has a density with respect to counting measure times Lebesgue measure, and on each such central window
\begin{equation}\label{eq:LLT}
 n^2 p_{n,R}(k,m,t)=
 \frac{\exp(-\xi^T\Sigma_R^{-1}\xi/2)}{(2\pi)^2\sqrt{\det\Sigma_R}}+o(1),
 \quad \xi=\frac{(k_1,k_2,m,t)-nm_R}{\sqrt n},
\end{equation}
with uniform essential-supremum error. No smoothing of $\eta_{n,R}$ is used.
\end{theorem}
\begin{proof}
Fourier inversion of $H_{n,R}$ gives a continuous density of the signed measure $\eta_{n,R}-E_{n,R}$. To justify this without assuming a density, take each lattice Fourier coefficient of $H$. It is an integrable roof transform of the corresponding signed component. Convolve that component with a Gaussian approximate identity, use integrability to pass to the inverse-transform limit, and test against compactly supported continuous functions. Uniqueness of finite Fourier transforms identifies each component, hence the full measure. This is a proof of inversion, not a change to the datum. Adding the known density of $E_{n,R}$ gives $p_{n,R}$.

On $|\omega|\le n^{-2/5}$, replacing $H$ by $\Phi$ costs at most a fixed volume times $\|E\|_{\TV}=o(n^{-2})$. Inserting (i), let $v=\sqrt n\,\omega$. The cubic error is uniformly $O(n^{-1/5})$ on this region; the bound on $\lambda$ supplies an integrable Gaussian dominator. Uniform Taylor bounds, truncation first to $|v|\le M$, and then $M\to\infty$, give the inverse Gaussian with error $o(n^{-2})$. The inverse normalization is $(2\pi)^{-4}$ and the Gaussian integral is $(2\pi)^2/\sqrt{\det\Sigma_R}$, giving the constant in \eqref{eq:LLT}. Hypothesis (ii) bounds the remaining frequencies, and (iii) bounds the added density in the central window. All estimates are uniform in $\xi$ in that window.
\end{proof}

For weighted inserted measures, the same proof applies with $A_R(0)=a_R$ and uniformly bounded amplitudes, provided the edge and tail hypotheses are checked for those measures as well. It gives $a_R$ times the Gaussian. It does not assert uniformity for arbitrary $n$-dependent path functionals without such estimates.

\begin{lemma}[An explicit frequency splice]\label{lem:splice}
Suppose a residual transform has a compact minor-arc bound $Ce^{-cn}$ and, for $1\le|b|\le e^{\kappa n}$, a bound $C(1+|b|)^A e^{-cn}$. Beyond $e^{\kappa n}$ suppose it is bounded by
\[
 C_M\rho^n(1+|b|)^{-2}+C_Me^{c_M n}(1+|b|)^{-M}.
\]
If one fixed integer $M>1$ and one $\kappa>0$ satisfy
\begin{equation}\label{eq:splice}
       \frac{c_M}{M-1}<\kappa<\frac{c}{A+1},
\end{equation}
these regions contribute exponentially small $L^1$ mass. Together with the Gaussian major-arc bound between $n^{-2/5}$ and $\delta$, this proves \eqref{eq:tailcriterion}.
\end{lemma}
\begin{proof}
The middle integral is bounded by $C\exp(-(c-(A+1)\kappa)n)$. The two outer integrals are bounded by $C_M\exp((\log\rho-\kappa)n)$ and $C_M\exp(-( (M-1)\kappa-c_M)n)/(M-1)$. The torus has finite volume. The major-arc Gaussian tail is bounded by a polynomial factor times $e^{-c n^{1/5}}$, and the compact remainder is exponential. Strict positivity of the three exponent margins proves the claim.
\end{proof}
The inequality \eqref{eq:splice} must be verified with the actual derivative-growth constants. The phrases ``take $\kappa$ small'' and ``then take $M$ large'' do not imply it: for example, $c_M=2(M-1)$ and $c/(A+1)=1$ leave no possible $\kappa$. This checks a proposed estimate, not the truth or falsity of a central LLT. The residual formulation also requires an independently proved bound for the sum of all extracted edge terms; Proposition~\ref{prop:subtract} treats one physical family only.


\section{Local inversion with nonnegligible remote edge mass}

\paragraph{Frequency convention.}
The rescaled frequency is $v=\sqrt n\,\omega$. When the proved growing
central band is used below, its radii are $2n^{1/200}$ in $v$ and
$2n^{-99/200}$ in $\omega$. General cutoffs in the following statements
retain their stated hypotheses; a Gaussian-tail estimate is not a bound
for the complementary transform of the raw measure.

\label{sec:localized-inversion}
Theorem~\ref{thm:LLT} gives a sufficient global smallness condition on
extracted edges. A central local limit needs less: the extracted measure
must have negligible density and negligible low-frequency convolution
on the central window. Its total variation need not tend to zero.
The following refinement retains the raw datum and the density norm.

Fix an even real function $\chi\in C_c^\infty(\R^4)$ equal to one on
the unit ball and supported in the ball of radius two. Let $a_n\to\infty$
with $a_n=o(n^{1/6})$, put $B_n=a_n/\sqrt n$, and set
$\chi_n(\omega)=\chi(\omega/B_n)$. For all sufficiently large $n$ its
support fits in the fundamental torus interval in the first three
coordinates. With the transform convention already fixed, its inverse
kernel on $\Z^3\times\R$ is
\begin{equation}\label{eq:localized-kernel}
 K_n(k,t)=B_n^4\check\chi(B_nk,B_nt),\qquad
 |K_n(k,t)|\le C_M B_n^4(1+B_n|(k,t)|)^{-M}.
\end{equation}
The check denotes the inverse transform on $\R^4$, including its
$(2\pi)^{-4}$ factor. Convolution here and below uses counting measure
in the discrete coordinates and Lebesgue measure in the roof coordinate.
For fixed $K$, let $W_{n,R,K}=\{z:|z-nm_R|\le K\sqrt n\}$.

\begin{theorem}[A local edge criterion]\label{thm:local-edge-LLT}
Assume the small-frequency hypothesis (i) of Theorem~\ref{thm:LLT}.
Let $E_{n,R}$ be finite signed measures with densities $e_{n,R}$ and
suppose $H_{n,R}=\Phi_{n,R}-\widehat E_{n,R}$ belongs to
$L^1(\T^3\times\R)$. Assume, uniformly in $R$, that
\begin{equation}\label{eq:local-tail}
 \int |1-\chi_n(\omega)|\,|H_{n,R}(\omega)|\dd\omega=o(n^{-2}).
\end{equation}
For every fixed central window suppose also that
\begin{equation}\label{eq:local-edge-errors}
 n^2\mathop{\rm ess\,sup}_{W_{n,R,K}}|e_{n,R}|\longrightarrow0,
 \qquad
 n^2\sup_{W_{n,R,K}}|K_n*E_{n,R}|\longrightarrow0.
\end{equation}
Then the raw density conclusion \eqref{eq:LLT} holds on that window.
There is no small-total-variation hypothesis on $E_{n,R}$.
\end{theorem}
\begin{proof}
As in the proof of Theorem~\ref{thm:LLT}, integrability of $H$ first
identifies a continuous density for $\eta-E$, and adding $e$ gives
an actual mixed density for $\eta$. Splitting this identity with
$\chi_n$ gives the exact equality almost everywhere
\begin{equation}\label{eq:exact-local-decomposition}
 p_{n,R}=K_n*\eta_{n,R}
       +(e_{n,R}-K_n*E_{n,R})
       +\mathcal F^{-1}[(1-\chi_n)H_{n,R}].
\end{equation}
All terms are defined: $E$ is finite, $K_n$ is bounded and integrable,
and the last inverse transform is absolutely convergent.
The last term is at most $(2\pi)^{-4}$ times the integral in
\eqref{eq:local-tail}. The middle term is controlled by
\eqref{eq:local-edge-errors}.

For the first term, its Fourier integral has the factor $\chi_n$.
Insert the assumed small-frequency expansion and use
$v=\sqrt n\,\omega$. The support has $|v|\le2a_n$ and
$nO(|\omega|^3)=O(a_n^3/\sqrt n)=o(1)$. The uniform quadratic
bound on $\lambda_R$ gives a Gaussian dominator, while the amplitude
converges to one. The exponentially small spectral remainder, after
integration and multiplication by $n^2$, is bounded by
$Ca_n^4\rho^n$ and tends to zero. Truncate first at a fixed $|v|$,
then let that truncation grow. This proves uniformly in $R$ and the
central variable that $n^2K_n*\eta_{n,R}$ converges to the Gaussian
in \eqref{eq:LLT}. Combining the three terms proves the assertion.

The kernel is only an analytical device in the exact decomposition
\eqref{eq:exact-local-decomposition}. The two correction terms are
retained and estimated; the result is not a local limit for a
smoothed replacement of the observed law.
\end{proof}

\begin{corollary}[Remote mass need not be small]\label{cor:remote-edges}
Suppose $\|E_{n,R}\|_{\TV}\le C$ and, for the window under
consideration, the support of $E_{n,R}$ is at distance at least
$\varepsilon\sqrt n$ from $W_{n,R,K}$, for a fixed
$\varepsilon>0$. Then \eqref{eq:local-edge-errors} holds.
In particular an extracted density of bounded, nonvanishing mass
can satisfy the local edge hypotheses.
\end{corollary}
\begin{proof}
The density vanishes almost everywhere in the central window. For
$z$ in that window, \eqref{eq:localized-kernel} gives
\[
 n^2|(K_n*E_{n,R})(z)|
 \le CC_M a_n^4(1+\varepsilon a_n)^{-M}\longrightarrow0
\]
by taking $M>4$. The Schwartz bound is available for every fixed $M$
because $\chi$ is smooth. These constants are kernel constants, not
the derivative-growth constants of the billiard in Lemma~\ref{lem:splice}.
\end{proof}

The same conclusion holds for a sum of remote extracted terms whose
total variations are summable with bounded total, and a remaining
nearby extracted density satisfying \eqref{eq:local-edge-errors}.
This follows by linearity and the triangle inequality. Thus individual
edge coefficients do not all have to be $o(n^{-2})$ merely because
one proves a central limit. Nevertheless, one must still construct
the full extraction and prove \eqref{eq:local-tail}. Neither this
corollary nor the single-word estimate in Proposition~\ref{prop:general-edge}
controls the sum of all central critical words or singular boundaries.
A restriction of a density to a sharp window can itself create new
jump terms; it is not automatically an integrable residual.

For weighted measures the same proof gives the corresponding amplitude
times the Gaussian, provided the small-frequency expansion, residual
estimate and local edge errors are verified for that weighted class.
This is the precise additional input needed before applying
Proposition~\ref{prop:conditioning} to a proposed family of insertions.

\section{An integrated growing central band for the physical record}

\paragraph{Frequency convention.}
The rescaled frequency is $v=\sqrt n\,\omega$. When the proved growing
central band is used below, its radii are $2n^{1/200}$ in $v$ and
$2n^{-99/200}$ in $\omega$. General cutoffs in the following statements
retain their stated hypotheses; a Gaussian-tail estimate is not a bound
for the complementary transform of the raw measure.

\label{sec:growing-major-arc}
The estimate on compact rescaled frequency sets can be strengthened to
an integrated estimate on an expanding band. The distinction matters
for local inversion: convergence after integration controls the inverse
transform uniformly in the observation variable. We prove this estimate
for the actual return record, including bounded-variation initial
insertions. No induced spectral realization is used.

An admissible initial insertion is a complex-valued function $a_R$ on
$M$, vanishing almost everywhere outside $Y_R^*$, such that
\[
 \|a_R\|_\infty+\|a_R\|_{\mathrm{BV}}\le B.
\]
Its integral is denoted by $\alpha_R$. The original probability has
$a_R=s_R$ and $\alpha_R=1$. Write
\begin{equation}\label{eq:inserted-central-characteristic}
 C^a_{n,R}(v)=\int_{Y_R^*}a_R(x)
 \exp\left(\frac{i v\cdot(J_{n,R}(x)-n\bar G_R)}{\sqrt n}\right)
 \dd\nu(x),\qquad v\in\R^4.
\end{equation}
Here $v$ is the rescaled frequency; the physical Fourier frequency is
$\omega=v/\sqrt n$. Constants below depend on $B$ but not on the
insertion, the radius, or $n$. No continuity of the insertion in $R$
is required.

\subsection{A frequency-dependent error bound}
\begin{lemma}[Collision estimate with explicit frequency dependence]
\label{lem:explicit-frequency-error}
Put $\ell_\delta=1+|\log\delta|$. There are $c,C>0$ and $\rho<1$
such that, for $m\ge2$, $0<\delta<1/2$, and real $t$ satisfying
\[
 |t|\le c\sqrt m\,\delta^2,\qquad
 \delta^{-6}|t|^3m^{-1/2}\le c,
\]
every complex function $a$ with
$\|a\|_\infty+\|a\|_{\mathrm{BV}}\le B$ satisfies
\begin{align}
 &\left|\int a\exp(it\cdot S_{m,R}h_R/\sqrt m)\dd\nu
       -\left(\int a\dd\nu\right)e^{-t^{\mathsf T}\Gamma_Rt/2}\right|
       \notag\\[-2mm]
 &\quad\le C\left\{
 \delta+|t|\sqrt{\delta\ell_\delta}+|t|^2\delta\ell_\delta
 +\frac{|t|\delta^{-4}+|t|^3\delta^{-6}}{\sqrt m}
 +\delta^{-2}\rho^m\right\}.\label{eq:explicit-frequency-error}
\end{align}
\end{lemma}
\begin{proof}
Set $a_\delta=\mathsf S_\delta a$ and $z=t/\sqrt m$.
The operator in Lemma~\ref{lem:smooth-collision-expansion} gives
\[
 \int a_\delta e^{it\cdot S_{m,R}h_{R,\delta}/\sqrt m}\dd\nu
 =\lambda_{R,\delta}(z)^m
   \ell_j\bigl(\Pi_{R,\delta}(z)a_\delta\nu\bigr)
       +O(\delta^{-2}\rho^m).
\]
The embedding bound in Lemma~\ref{lem:density-norm-details} below
makes the initial-vector factor explicit. The principal amplitude
at zero is $\int a\dd\nu$, and its change is at most
$C|t|\delta^{-4}/\sqrt m$. Moreover
\[
 m\log\lambda_{R,\delta}(t/\sqrt m)
 =-\tfrac12t^{\mathsf T}\Gamma_{R,\delta}t
       +O(\delta^{-6}|t|^3/\sqrt m).
\]
The stated smallness conditions put $z$ in the analytic ball and bound
the last remainder. Positive semidefiniteness of
$\Gamma_{R,\delta}$ therefore bounds the principal power by a fixed
constant. The inequality $|e^w-1|\le |w|e^{|w|}$ proves the two
spectral error terms in \eqref{eq:explicit-frequency-error}.

Mean-preserving smoothing changes the initial insertion by $O(\delta)$
in $L^1$. If $r=h_R-h_{R,\delta}$, then
\[
 \int |a_\delta|\,|S_{m,R}r|^2\dd\nu
       \le C m\delta\ell_\delta,
\]
by \eqref{eq:small-bv-variance} and $|a_\delta|\le B$.
Cauchy--Schwarz bounds the characteristic error from removing this
smoothing by $C|t|\sqrt{\delta\ell_\delta}$. This argument uses
absolute values of the insertion and so applies to signed and complex
insertions as well. Finally, the segment joining the two covariance
matrices consists of positive semidefinite matrices. Differentiation
along that segment and \eqref{eq:collision-covariance-limit} give
\[
 \left|e^{-t^{\mathsf T}\Gamma_{R,\delta}t/2}
             -e^{-t^{\mathsf T}\Gamma_Rt/2}\right|
       \le C|t|^2\delta\ell_\delta.
\]
Combining these estimates proves the lemma.
\end{proof}

\begin{proposition}[Stopped estimate on finite frequency balls]
\label{prop:stopped-band-error}
For $U\ge1$, let
\[
 \mathcal E_n(U)=\sup_{R,a_R}
   \int_{|v|\le2U}
    |C^a_{n,R}(v)-\alpha_Re^{-v^{\mathsf T}D_Rv/2}|\dd v.
\]
If $U\le c\sqrt n\,\delta^2$ and
$\delta^{-6}U^3/\sqrt n\le c$, with $c$ sufficiently small, then
\begin{align}
 \mathcal E_n(U)\le C\bigl\{
 &\delta U^4+\sqrt{\delta\ell_\delta}\,U^5
       +\delta\ell_\delta U^6
       +n^{-1/2}\delta^{-4}U^5+n^{-1/2}\delta^{-6}U^7
       \notag\\
 &+\delta^{-2}\rho_0^nU^4
       +n^{-1/5}U^4+n^{-1/5}\log(2+n)U^5+n^{-1}U^6
                          \bigr\},\label{eq:integrated-band-error}
\end{align}
where $\rho_0<1$. The supremum is over the admissible insertions.
\end{proposition}
\begin{proof}
Put $m_n=\lfloor n/c_*\rfloor$ and $b_n=\lceil n^{3/5}\rceil$.
The exact compensation identity \eqref{eq:exact-stopped-compensation}
permits replacement of the centered record by $S_{N_{n,R},R}h_R$.
The exceptional set $|N_{n,R}-m_n|>b_n$ has $\nu_R^*$-probability
$O(n/b_n^2)$ by Lemma~\ref{lem:stopped-compensation}; rounding changes
only the constant. Its contribution with insertion $a_R$ is bounded
by the same order, since $|a_R|\nu\le Bc_*\nu_R^*$ on the section.
On its complement use the forward and backward deterministic-window
maxima from Lemma~\ref{lem:dyadic-collision-maximal}. The elementary
exponential difference bound and Cauchy--Schwarz give
\begin{align}
 &\left|C^a_{n,R}(v)
       -\int a_R e^{iv\cdot S_{m_n,R}h_R/\sqrt n}\dd\nu\right|
       \notag\\
 &\hspace{12mm}\le C\left\{n^{-1/5}
                  +|v|n^{-1/5}\log(2+n)\right\}.
                 \label{eq:weighted-stopping-error}
\end{align}
No second moment of an unbounded stopped sum on the exceptional set
is used. In particular the estimate is valid for complex insertions.

Apply Lemma~\ref{lem:explicit-frequency-error} at $m_n$ with
$t=v\sqrt{m_n/n}$. These frequencies satisfy its restrictions after
changing the uniform constant $c$. The covariance of the deterministic
sum is $(m_n/n)\Gamma_R$, which differs from $D_R=c_*^{-1}\Gamma_R$
by $O(n^{-1})$. Its Gaussian characteristic function consequently
changes by at most $C|v|^2/n$. Integrating these bounds and
\eqref{eq:weighted-stopping-error} uses only
$\int_{|v|\le2U}|v|^j\dd v=C_jU^{4+j}$, for $j\ge0$.
This proves \eqref{eq:integrated-band-error}. Finitely many small values
of $n$ can be absorbed by increasing the constants wherever the stated
smallness conditions hold.
\end{proof}

\subsection{A polynomially expanding major arc}
\begin{theorem}[Integrated major arc of the actual return record]
\label{thm:growing-major-arc}
Let
\[
 0<\varepsilon<\frac1{134},\qquad
 10\varepsilon<\theta<\frac{1/2-7\varepsilon}{6},\qquad
 r=\min\left\{\frac\theta2-5\varepsilon,
              \frac12-6\theta-7\varepsilon,
              \frac15-5\varepsilon\right\}.
\]
Then $r>0$ and
\begin{equation}\label{eq:general-growing-arc}
 \sup_{R,a_R}\int_{|v|\le2n^\varepsilon}
 |C^a_{n,R}(v)-\alpha_Re^{-v^{\mathsf T}D_Rv/2}|\dd v
       \le C n^{-r}\log(2+n).
\end{equation}
In particular, with $\varepsilon=1/200$ and $\theta=1/14$,
\begin{equation}\label{eq:explicit-growing-arc}
 \sup_{R,a_R}\int_{|v|\le2n^{1/200}}
 |C^a_{n,R}(v)-\alpha_Re^{-v^{\mathsf T}D_Rv/2}|\dd v
       \le C n^{-3/280}\sqrt{\log(2+n)}.
\end{equation}
These estimates hold for $n\ge2$. They require only the positive
semidefinite covariance already established, not positive definiteness.
\end{theorem}
\begin{proof}
Take $U=n^\varepsilon$ and $\delta=\tfrac14n^{-\theta}$ in
Proposition~\ref{prop:stopped-band-error}. Its analytic restrictions hold
for all sufficiently large $n$: in particular
$3\varepsilon+6\theta<1/2$ follows from
$7\varepsilon+6\theta<1/2$, and
$\varepsilon+2\theta<1/2$ follows as well. Constants in the radius
of the analytic ball do not affect these strict inequalities.

The three potentially slow terms in \eqref{eq:integrated-band-error}
are, respectively,
\[
 n^{-(\theta/2-5\varepsilon)}\sqrt{\log(2+n)},\qquad
 n^{-(1/2-6\theta-7\varepsilon)},\qquad
 n^{-(1/5-5\varepsilon)}\log(2+n).
\]
All other polynomial terms decay at least as fast as one of these;
this follows by subtracting their exponents, using $\theta>10\varepsilon$.
For the rounding term also note that
$1-6\varepsilon>1/5-5\varepsilon$ in the indicated range.
The exponentially decaying term is smaller than every fixed negative
power. The interval for $\theta$ is nonempty precisely when
$67\varepsilon<1/2$, proving \eqref{eq:general-growing-arc}.

At the displayed explicit choice the three margins are
\[
 \frac3{280},\qquad \frac{51}{1400},\qquad \frac7{40}.
\]
Only the first is minimal, and that term carries a square root of the
logarithm. The strict excess of the other margins absorbs their
logarithmic factors. Increasing the constant handles the finite initial
range, since both characteristic functions are bounded uniformly by
a constant depending on $B$. This proves \eqref{eq:explicit-growing-arc}.
\end{proof}

\subsection{Insertion into exact raw inversion}
Define the finite complex measure $\mu^a_{n,R}$ by pushing forward
$a_R\nu$ under $J_{n,R}$, and write $\Phi^a_{n,R}$ for its transform.
Take $a_n=n^{1/200}$ in the cutoff and kernel of
\eqref{eq:localized-kernel}, and use $m_R=\bar G_R$ in the central
windows. The preceding theorem supplies a proved replacement for the
spectral small-frequency hypothesis in the following inversion route.

\begin{corollary}[Raw inversion from the proved central band]
\label{cor:raw-from-growing-arc}
Use the uniform bound $D_R\ge d_0 I_4$ of
Theorem~\ref{thm:joint-uniform-nondegeneracy}.
For an admissible family of insertions, let $E^a_{n,R}$ be finite complex
measures with mixed densities $e^a_{n,R}$ and put
$H^a_{n,R}=\Phi^a_{n,R}-\widehat E^a_{n,R}$. Assume that
$H^a_{n,R}\in L^1(\T^3\times\R)$ and, uniformly in the family,
\begin{equation}\label{eq:new-complementary-tail}
 \int |1-\chi_n(\omega)|\,|H^a_{n,R}(\omega)|\dd\omega=o(n^{-2}).
\end{equation}
Assume also the two local edge conditions
\eqref{eq:local-edge-errors} with $e,E$ replaced by $e^a,E^a$.
Then $\mu^a_{n,R}$ has a mixed density $p^a_{n,R}$, and for each fixed
$K$,
\begin{align}
 \mathop{\rm ess\,sup}_{|z-n\bar G_R|\le K\sqrt n}
 \left|n^2p^a_{n,R}(z)
 -\frac{\alpha_R}{(2\pi)^2\sqrt{\det D_R}}
   \exp\left(-\tfrac12\xi^{\mathsf T}D_R^{-1}\xi\right)\right|
 &\longrightarrow0,
 \quad \xi=\frac{z-n\bar G_R}{\sqrt n},
 \label{eq:inserted-raw-conclusion}
\end{align}
uniformly in $R$ and the admissible family. No fixed-neighborhood
spectral expansion of an induced operator is assumed in this corollary.
\end{corollary}
\begin{proof}
Fourier inversion for $H^a$ identifies a continuous density for
$\mu^a-E^a$. Add $e^a$ and use the exact decomposition
\eqref{eq:exact-local-decomposition}, which is linear and remains valid
for complex measures. After $v=\sqrt n\,\omega$, the difference between
$n^2K_n*\mu^a$ and the inverse Gaussian integral with cutoff
$\chi(v/n^{1/200})$ is bounded, independently of $z$, by
$(2\pi)^{-4}\|\chi\|_\infty$ times the integral in
\eqref{eq:explicit-growing-arc}. Positive definiteness bounds the
removed Gaussian tail by
$C\int_{|v|>n^{1/200}}e^{-d_0|v|^2/2}\dd v$, which tends to zero
uniformly. This proves the main term in
\eqref{eq:inserted-raw-conclusion}. The other two terms of the exact
decomposition vanish on the stated scale by the assumed complementary
tail and local edge conditions. These steps prove the raw density
conclusion, not a conclusion only for a convolved observation.
\end{proof}

\begin{remark}[Location and scope of the remaining estimates]
\label{rem:growing-arc-boundary}
The physical radius of the proved central band is
$n^{-1/2+1/200}$. It is smaller than $n^{-2/5}$. Thus
\eqref{eq:new-complementary-tail} is not a consequence of the outer-tail
hypothesis in Theorem~\ref{thm:LLT}: the intervening annulus also needs
an estimate. The present theorem removes the need for an assumed induced
spectral expansion on the central band only. It does not assert a
resolvent estimate on its complement or the full physical branch
decomposition. Uniform positive definiteness is proved independently
in Theorem~\ref{thm:joint-uniform-nondegeneracy}. Both inversion routes are
retained, with their distinct frequency boundaries.

The insertion class is an initial-state class with a uniform norm on
the collision cylinder. Terminal insertions, multiple-time insertions,
and exact-event indicators need not belong to it. In particular
\eqref{eq:explicit-growing-arc} does not justify a change of exact
conditioning event or provide the weighted raw tail required for such
a change. The count truncation and the initial-coordinate variation
bounds continue to concern different objects from the derivative sums
of the raw residual densities.
\end{remark}

\section{Central integrals marked at an arbitrary actual return}
\label{sec:marked-return-band}
A weight at a terminal or intermediate return cannot simply be smoothed
as a function of the initial state: composition with a long first-return
iterate need not preserve a uniform variation bound. Instead we change
the origin of the orbit to the marked return. The record then has a
backward and a forward part, while the weight is a function of the new
origin. Both parts can be estimated with the same forward collision
operator. This argument does not require invariance of the section under
time reversal, or mixing of the induced map.

All measures and normalizations in this section are as follows. The
measure $\nu$ is the probability on the full collision cylinder;
$\nu(Y_R^*)=c_*$, and $\nu_R^*=c_*^{-1}\nu|_{Y_R^*}$. For a complex
function $a$ that is zero almost everywhere outside $Y_R^*$ put
\begin{equation}\label{eq:marked-norms}
 M_a=\|a\|_\infty,\qquad V_a=\|a\|_{\mathrm{BV}},\qquad
 \alpha_a=\int_M a\dd\nu.
\end{equation}
For complex functions the variation norm is the sum of the norms of the
real and imaginary parts. It includes their $L^1$ norms and is taken
in the original collision coordinates, including the section boundary.
No upper bound for $V_a$ independent of $n$ is imposed below.

For $0\le k\le n$ define the finite complex measure and its centered,
rescaled transform by
\begin{align}
 \mu^{[k],a}_{n,R}(B)
   &=\int_{Y_R^*}a((F_R^*)^kx)\one_{\{J_{n,R}(x)\in B\}}\dd\nu(x),
       \label{eq:marked-measure}\\
 C^{[k],a}_{n,R}(v)
   &=\int_{Y_R^*}a((F_R^*)^kx)
          e^{iv\cdot(J_{n,R}(x)-n\bar G_R)/\sqrt n}\dd\nu(x).
       \label{eq:marked-transform}
\end{align}
Its mass is $\alpha_a$, which may be complex. The endpoints $k=0,n$
are the initial and terminal insertions. Taking $a=s_R$ gives the
original probability for every $k$. The physical frequency is always
$\omega=v/\sqrt n$.

\subsection{Changing the origin to the mark}
Let $N^-_{j,R}(y)$ be the collision distance to the $j$th preceding
section visit, so that
\[
 (F_R^*)^{-j}y=T_R^{-N^-_{j,R}(y)}y,\qquad N^-_{0,R}=0.
\]
The forward distance is $N^+_{j,R}=N_{j,R}$, with $N^+_{0,R}=0$.
These distances are finite almost everywhere, by recurrence and
invertibility. For nonnegative integers $r,s$ write
\begin{equation}\label{eq:two-sided-collision-sum}
 \mathsf H_{r,s,R}(y)=\sum_{j=-r}^{s-1}h_R(T_R^jy),
\end{equation}
where a sum over an empty interval is zero.

\begin{lemma}[Exact compensation about a marked return]
\label{lem:marked-compensation}
For almost every $y\in Y_R^*$,
\begin{equation}\label{eq:marked-compensation}
 J_{n,R}((F_R^*)^{-k}y)-n\bar G_R
   =\mathsf H_{N^-_{k,R}(y),N^+_{n-k,R}(y),R}(y).
\end{equation}
Consequently
\begin{equation}\label{eq:marked-recentering}
 C^{[k],a}_{n,R}(v)=\int_{Y_R^*}a(y)
 e^{iv\cdot\mathsf H_{N^-_{k,R}(y),N^+_{n-k,R}(y),R}(y)/\sqrt n}
 \dd\nu(y).
\end{equation}
\end{lemma}
\begin{proof}
The collision interval from $(F_R^*)^{-k}y$ up to, but excluding,
$(F_R^*)^{n-k}y$ has left endpoint $-N^-_{k,R}(y)$ and right endpoint
$N^+_{n-k,R}(y)$ in coordinates based at $y$. The sum of $f_R$ on this
interval is precisely the original four-coordinate record. Its negative
part contains exactly $k$ section visits, including the left endpoint
and excluding zero. Its nonnegative part contains exactly $n-k$ visits,
including zero if that part is nonempty and excluding the right endpoint.
Thus the sum of $\eta_R$ over the entire interval is $n$. Substitute
$h_R=f_R-\bar G_R\eta_R$. This proves the first identity, including
$k=0,n$. The actual first-return map preserves $\nu|_{Y_R^*}$, since it
preserves its normalization $\nu_R^*$. The change of variable
$y=(F_R^*)^kx$ gives the second identity. Only additivity and invariance
are used; the two parts are not asserted independent.
\end{proof}

\begin{lemma}[Forward and backward return-clock windows]
\label{lem:two-sided-clock-window}
For $j\ge0$, $b\ge2$, and $m_j=\lfloor j/c_*\rfloor$, uniformly in $R$,
\begin{equation}\label{eq:two-sided-clock-window}
 \nu_R^*\{|N^\pm_{j,R}-m_j|>b\}\le C\frac{j+b}{b^2}.
\end{equation}
Moreover the collision maximal estimate of
Lemma~\ref{lem:dyadic-collision-maximal} holds for an interval beginning
at any integer collision time, positive or negative, and in either order.
\end{lemma}
\begin{proof}
The case $j=0$ is immediate. For either sign let
$H_L^\pm=\sum_{i=1}^L\eta_R\circ T_R^{\pm i}$. Its mean under $\nu$
is $c_*L$. The variance is at most $CL$ by
Proposition~\ref{prop:bv-correlation}. For the negative sign the scalar
lag correlation is the same as for the positive sign by invariance.
Domination $\nu_R^*\le c_*^{-1}\nu$ gives
\[
 \int|H_L^\pm-c_*L|^2\dd\nu_R^*\le CL.
\]
The equivalences $N_j^\pm>L\iff H_L^\pm<j$ and
$N_j^\pm\le L\iff H_L^\pm\ge j$ hold for the actual visits. Take
integer thresholds within a fixed distance of $m_j+b$ and $m_j-b$.
When the latter threshold is negative, its lower-tail event is empty.
Otherwise the required deviation from $c_*L$ is at least $c_*b-O(1)$,
while $L\le C(j+b)$. Chebyshev's inequality proves the estimate for
$b$ above a fixed constant. For the remaining bounded range enlarge $C$
so that its right-hand side is at least one. In particular no assumption
$b\le j$ is needed.

For the maximal assertion, translate any deterministic integer interval
to a nonnegative one using invariance of $\nu$. Every subinterval still
has second moment bounded by its length, from the absolutely summable
collision correlations. The same dyadic prefix proof applies in the
reverse order. Domination gives its version under $\nu_R^*$ as well.
The assertion concerns fixed intervals, not a conditional distribution
of the past given the future.
\end{proof}

\begin{proposition}[Uniform stopping error about the mark]
\label{prop:marked-stopping}
For all $0\le k\le n$, $n\ge2$, and $v\in\R^4$,
\begin{align}
 &\left|C^{[k],a}_{n,R}(v)
  -\int_M a(y)e^{iv\cdot\mathsf H_{m_k,m_{n-k},R}(y)/\sqrt n}\dd\nu(y)
       \right|\notag\\
 &\hspace{12mm}\le C M_a\left(n^{-1/5}
             +|v|n^{-1/5}\log(2+n)\right).
       \label{eq:marked-stopping-error}
\end{align}
\end{proposition}
\begin{proof}
Use \eqref{eq:marked-recentering} and let
$b_n=\lceil n^{3/5}\rceil$. By Lemma~\ref{lem:two-sided-clock-window},
the union of the two events on which a clock differs from its
corresponding $m_j$ by more than $b_n$ has probability at most
$Cn^{-1/5}$ under $\nu_R^*$. Its contribution to the difference of
characteristic functions is at most $C M_an^{-1/5}$, since
$|a|\nu\le M_ac_*\nu_R^*$ on the section.

On the complement, the difference of the two collision sums is bounded
by the sum of two maxima on deterministic windows of length at most
$2b_n+2$ around their respective endpoints. The preceding maximal
estimate bounds its squared integral under $\nu$ by
$C b_n[\log(2+b_n)]^2$. It remains valid when a window crosses collision
time zero. Cauchy--Schwarz and $|e^{ix}-e^{iy}|\le|x-y|$ give
$C M_a|v|\sqrt{b_n/n}\log(2+b_n)$, which is the second term in
\eqref{eq:marked-stopping-error}. Nothing is estimated in an unbounded
stopped sum on the exceptional event. Finitely many initial values of
$n$ are included by increasing the constant.
\end{proof}

\subsection{A collision multiplier between two spectral powers}
Set $a_\delta=\mathsf S_\delta a$ and $h_{R,\delta}=\mathsf S_\delta h_R$.
The former need not be supported in the section. This is harmless:
\begin{equation}\label{eq:marked-smoothing-norms}
 \|a-a_\delta\|_1\le C\delta V_a,\qquad
 \|a_\delta\|_\infty\le M_a,\qquad
 \|M_{a_\delta}\|_{\mathcal B_j\to\mathcal B_j}
       \le C M_a\delta^{-2}.
\end{equation}
The last bound is Lemma~\ref{lem:density-norm-details}. In particular
its constant depends on the supremum norm, not on $V_a$.

\begin{lemma}[A marked deterministic collision integral]
\label{lem:marked-spectral}
Fix $c_0,C_0>0$. Suppose $c_0n\le r+s\le C_0n$ with $r,s\ge0$,
and let $q\ge1$ be an integer satisfying $2q\le r+s$. There are
$c,C>0$ and $\rho<1$ such that, whenever
\[
 |v|/\sqrt n\le c\delta^2,\qquad
 \delta^{-6}|v|^3/\sqrt n\le c,
\]
the following bound holds, with $\ell_\delta=1+|\log\delta|$:
\begin{align}
 &\left|\int_M a\,e^{iv\cdot\mathsf H_{r,s,R}/\sqrt n}\dd\nu
      -\alpha_a e^{-\frac{r+s}{2n}v^{\mathsf T}\Gamma_Rv}\right|
       \notag\\
 &\quad\le C\delta V_a+C M_a\left\{
 |v|\sqrt{\delta\ell_\delta}+|v|^2\delta\ell_\delta
 +\frac{|v|\delta^{-4}+|v|^3\delta^{-6}}{\sqrt n}
 +\delta^{-2}\rho^q+\frac{|v|q}{\sqrt n}
 +\frac{|v|^2q}{n}\right\}.
       \label{eq:marked-spectral-bound}
\end{align}
All constants are independent of the location of the mark and of $a$.
\end{lemma}
\begin{proof}
First replace $a$ by $a_\delta$. The cost is at most $C\delta V_a$.
Let $u_\delta=h_R-h_{R,\delta}$. Stationarity and
\eqref{eq:small-bv-variance} give
\[
 \int\left|\sum_{i=-r}^{s-1}u_\delta(T_R^iy)\right|^2\dd\nu(y)
       \le C(r+s)\delta\ell_\delta.
\]
The bounded weight $|a_\delta|\le M_a$ and Cauchy--Schwarz therefore
bound the cost of replacing $h_R$ by $h_{R,\delta}$ by
$C M_a|v|\sqrt{\delta\ell_\delta}$.

Put $z=v/\sqrt n$ and $L_z=\mathcal L_{R,\delta,z}$, using one of
the fixed local collision spaces. Its mass functional will be denoted
by $\ell$ to avoid confusion with block indices. Invariance and the
transfer convention give the exact chronological pairing
\begin{equation}\label{eq:marked-chronological-pairing}
 \int_M a_\delta(y)
 e^{iz\cdot\sum_{i=-r}^{s-1}h_{R,\delta}(T_R^iy)}\dd\nu(y)
       =\ell\bigl(L_z^s M_{a_\delta} L_z^r\nu\bigr).
\end{equation}
Indeed, starting at $x=T_R^{-r}y$, the rightmost power records collisions
$0,\ldots,r-1$, the multiplier is evaluated at $T_R^rx=y$, and the
leftmost power records collisions $r,\ldots,r+s-1$. Thus all weights
are evaluated at the prescribed states.

Suppose first that $r,s\ge q$. Write
$L_z^m=\lambda(z)^m\Pi(z)+E_m(z)$ with
$\|E_m(z)\|\le C\rho^m$. The hypotheses and
\eqref{eq:smooth-cubic-expansion} bound each real-frequency principal
power of length at most $r+s$ by a uniform constant. Expanding
\eqref{eq:marked-chronological-pairing}, every complementary term is
bounded by $C M_a\delta^{-2}(\rho^r+\rho^s)$; the term with both
complementary powers is included in this bound. For the principal
amplitude use $\Pi(0)=\nu\ell$ and \eqref{eq:marked-smoothing-norms}:
\begin{align*}
 \ell(\Pi(0)M_{a_\delta}\Pi(0)\nu)&=\alpha_a,\\
 |\ell(\Pi(z)M_{a_\delta}\Pi(z)\nu)-\alpha_a|
       &\le C M_a\delta^{-4}|z|.
\end{align*}
The cubic remainder accumulated over $r+s$ collisions is at most
$C\delta^{-6}|v|^3/\sqrt n$. Since the quadratic form is nonnegative,
split the principal comparison as
$\lambda(z)^{r+s}(A(z)-\alpha_a)
 +\alpha_a(\lambda(z)^{r+s}-e^{-(r+s)z^{\mathsf T}\Gamma_{R,\delta}z/2})$,
where $A(z)$ is its principal amplitude. The first term uses the
amplitude bound above; the second uses $|\alpha_a|\le M_a$ and
$|e^w-1|\le |w|e^{|w|}$. Thus the cubic error is multiplied only
by $C M_a$, not by the multiplier norm. This proves the required spectral bound in the
case of two long blocks.

If $r<q$, delete the first $r$ terms from the real exponent on the left
of \eqref{eq:marked-chronological-pairing}. The cost is at most
$C M_a|v|q/\sqrt n$, by boundedness of $h_{R,\delta}$. The remaining
pairing is $\ell(L_z^s(a_\delta\nu))$. Its principal amplitude differs
from $\alpha_a$ by at most $C M_a\delta^{-4}|z|$, by the smooth-density
embedding, and its complementary term is bounded by
$C M_a\delta^{-2}\rho^s$. Here $s\ge q$. If $s<q$, the remaining
pairing is instead $\ell(M_{a_\delta}L_z^r\nu)$, with exactly the same
bounds. At most one of $r,s$ is less than $q$. Restoring the removed
Gaussian time coefficient changes its value by at most
$C M_a|v|^2q/n$. This handles zero blocks as well; no spectral expansion
of an identity operator is used.

Finally replace $\Gamma_{R,\delta}$ by $\Gamma_R$. The segment joining
these matrices is positive semidefinite, and
\eqref{eq:collision-covariance-limit} bounds the change by
$C M_a|v|^2\delta\ell_\delta$. All the displayed errors combine to
\eqref{eq:marked-spectral-bound}. Only forward collision operators
occurred in the proof.
\end{proof}

\subsection{The growing band and norm growth}
\begin{theorem}[Integrated central band with an arbitrary return mark]
\label{thm:marked-return-band}
There is a constant $C$ such that for every $n\ge2$, $R\in I$,
$0\le k\le n$, and complex $a$ as in \eqref{eq:marked-norms},
\begin{align}
 &\int_{|v|\le2n^{1/200}}
   \left|C^{[k],a}_{n,R}(v)-\alpha_a e^{-v^{\mathsf T}D_Rv/2}\right|\dd v
      \notag\\
 &\hspace{15mm}\le C\left\{
 M_an^{-3/280}\sqrt{\log(2+n)}+V_an^{-9/175}\right\}.
      \label{eq:marked-return-band}
\end{align}
In particular the error tends to zero uniformly for bounded $M_a$ and
$V_a=O(n^\kappa)$ with any $\kappa<9/175$.
\end{theorem}
\begin{proof}
Take $r=m_k$ and $s=m_{n-k}$, so
\[
 (r+s)/n=c_*^{-1}+O(n^{-1}),
\]
with an error independent of $k$. They satisfy the length bounds in
Lemma~\ref{lem:marked-spectral}. Choose
\begin{equation}\label{eq:marked-scales}
 U_n=n^{1/200},\qquad \delta_n=\tfrac14n^{-1/14},\qquad
 q_n=\left\lceil\left(1+\frac4{|\log\rho|}\right)\log(2+n)\right\rceil.
\end{equation}
Then $\rho^{q_n}\le(2+n)^{-4}$ and $2q_n\le r+s$ for sufficiently
large $n$, uniformly in the mark. On $|v|\le2U_n$, both analytic
restrictions hold for large $n$: their positive margins are
$1/2-2/14-1/200$ and $1/2-6/14-3/200$.

Apply the marked spectral estimate and
Proposition~\ref{prop:marked-stopping}. Replacing $(r+s)\Gamma_R/n$
by $D_R$ costs at most $C M_a|v|^2/n$. Integration uses
$\int_{|v|\le2U}|v|^j\dd v=C_jU^{4+j}$. The contribution involving
$V_a$ is only
\[
 C V_a\delta_nU_n^4=C'V_an^{-9/175}.
\]
For clarity, the negative-power margins of the remaining polynomial
terms are displayed below; logarithmic factors are recorded separately.
\begin{center}
\begin{tabular}{lll}
Source & Margin & Logarithmic factor\\ \hline
Observable smoothing & $3/280$ & $\sqrt{\log(2+n)}$\\
Covariance smoothing & $29/700$ & $\log(2+n)$\\
Principal amplitude & $53/280$ & $1$\\
Cubic remainder & $51/1400$ & $1$\\
Clock exception & $9/50$ & $1$\\
Clock window & $7/40$ & $\log(2+n)$\\
Short collision block & $19/40$ & $\log(2+n)$\\
Time coefficient and rounding & $97/100$ & $\log(2+n)$
\end{tabular}
\end{center}
The complementary-power contribution is at most
$C M_an^{2/14+4/200}(2+n)^{-4}$ and is smaller than the displayed
polynomial errors. The first margin in the table is strictly smaller
than all the others, so it absorbs their logarithmic factors.
This proves \eqref{eq:marked-return-band} for large $n$. For the
remaining finite range, both transforms have modulus at most $M_a$,
and the frequency ball has finite volume; enlargement of $C$ proves
the stated estimate for all $n\ge2$. The last assertion follows by
substitution, without changing the smoothing scale with the insertion.
\end{proof}

\begin{corollary}[Conditioning on a state at a marked return]
\label{cor:marked-state-conditioning}
Let $E_{n,R}\subset Y_R^*$ have positive measure and a zero-extended
indicator in $\mathrm{BV}(M)$. Set
$V_{n,R}=\|\one_{E_{n,R}}\|_{\mathrm{BV}}$. Under the section probability,
condition on the exact event
\[
 A_{n,k,R}=\{x:(F_R^*)^kx\in E_{n,R}\}
\]
Uniformly for $0\le k\le n$,
\begin{align}
 &\int_{|v|\le2n^{1/200}}
 \left|\mathbb E_{\nu_R^*}\left[
 e^{iv\cdot(J_{n,R}-n\bar G_R)/\sqrt n}\mid A_{n,k,R}\right]
          -e^{-v^{\mathsf T}D_Rv/2}\right|\dd v\notag\\
 &\quad\le\frac{C}{\nu(E_{n,R})}
  \left(n^{-3/280}\sqrt{\log(2+n)}+V_{n,R}n^{-9/175}\right).
       \label{eq:marked-conditional-band}
\end{align}
For example this tends to zero if, uniformly in $R$,
$\nu_R^*(E_{n,R})\ge c n^{-\beta}$ and $V_{n,R}\le Cn^\kappa$,
where $\beta<3/280$ and $\beta+\kappa<9/175$.
\end{corollary}
\begin{proof}
Invariance gives $\nu_R^*(A_{n,k,R})=\nu(E_{n,R})/c_*$.
The conditional characteristic function is exactly
$C^{[k],\one_{E_{n,R}}}_{n,R}/\nu(E_{n,R})$. Apply
Theorem~\ref{thm:marked-return-band} with $M_a=1$. The two strict
exponent inequalities absorb the logarithm. This argument uses the
same event throughout and makes no comparison with another event.
\end{proof}

\paragraph{Relation to the raw density theorem.}
These estimates control the central integral of the actual complex
measures \eqref{eq:marked-measure}. The exact local decomposition
\eqref{eq:exact-local-decomposition} also applies to these measures by
linearity. More explicitly, use $D_R\ge d_0 I_4$ from
Theorem~\ref{thm:joint-uniform-nondegeneracy}; suppose they admit
extracted measures $E^{[k],a}_{n,R}$ with mixed densities; and suppose
the residual transform $H^{[k],a}_{n,R}$ is in
$L^1(\T^3\times\R)$, satisfies the complementary-tail condition
\eqref{eq:new-complementary-tail}, and satisfies the local edge
conditions \eqref{eq:local-edge-errors}, all uniformly in the chosen
marks and insertions. If the right-hand side of
\eqref{eq:marked-return-band} tends to zero uniformly, the proof of
Corollary~\ref{cor:raw-from-growing-arc} gives its raw density conclusion
for these marked measures. Thus its central hypothesis is now proved
also for terminal and single intermediate state insertions. The other
hypotheses have not been proved by the recentering argument.

The band still has physical radius $n^{-99/200}$, not $n^{-2/5}$.
Multiple separated return-state factors are not one insertion at a
single mark, and composing them with a long induced iterate does not
give the variation norm used here. Likewise an exact final lattice or
flight-time constraint is a condition on the record, not generally a
collision-state set $E_{n,R}$ of the class in
Corollary~\ref{cor:marked-state-conditioning}. Its relative error bound
is consequently not a replacement for the weighted raw local limits
needed for those constraints. Nor does a two-sided orbit identity
provide a pointwise representative of the zero-variance coboundary.
Covariance nondegeneracy is now proved in
Theorem~\ref{thm:joint-uniform-nondegeneracy}. The complete complementary
integral and full critical/singular residual estimates retain their
separate roles.

\section{Unsmoothed moments and the covariance of actual returns}
\label{sec:unsmoothed-moments}
The covariance of a weak Gaussian limit need not be the limit of the
normalized second moments. We establish that identification for the
actual return record, including a bounded-variation insertion at any
return. The argument also gives fourth moments for the unsmoothed
collision sums and improves the two-sided stopping estimate. It uses
only the collision splitting, not decay of correlations for $F_R^*$.

Write $\|u\|_{\mathcal V}=\|u\|_\infty+\|u\|_{\mathrm{BV}}$.
All expectations without a subscript in this section are with respect
to the collision probability $\nu$. Constants are uniform in $R\in I$.
The smoothing operators preserve means and have the bounds in
Lemma~\ref{lem:bv-smoothing}.

\subsection{Decoupling at a deterministic collision gap}
\begin{lemma}[A finite-product decoupling estimate]
\label{lem:bv-finite-product}
For $2\le q\le4$, integers $t_1\le\cdots\le t_q$, and
$1\le p<q$, bounded-variation functions $u_1,\ldots,u_q$ satisfy
\begin{align}
 &\left|\mathbb E\prod_{j=1}^q u_j\circ T_R^{t_j}
 -\left(\mathbb E\prod_{j=1}^p u_j\circ T_R^{t_j}\right)
  \left(\mathbb E\prod_{j=p+1}^q u_j\circ T_R^{t_j}\right)\right|
 \notag\\
 &\hspace{18mm}\le C\exp\{-\gamma(t_{p+1}-t_p)\}
                         \prod_{j=1}^q\|u_j\|_{\mathcal V}.
 \label{eq:bv-finite-product}
\end{align}
The times may be negative or repeated. The constants $C,\gamma>0$
can be chosen for all the indicated values of $q$ and $p$.
\end{lemma}
\begin{proof}
By multilinearity assume $\|u_j\|_{\mathcal V}\le1$; a zero norm
makes the assertion immediate. Replace every $u_j$ by
$u_{j,\epsilon}=\mathsf S_\epsilon u_j$. Telescoping the products,
invariance of $\nu$, and the uniform supremum bounds show that the
change in the left-hand expression is at most $C_q\epsilon$,
independently of all the times. In particular, no variation norm of
$u_j\circ T_R^{t_j}$ is estimated.

Translate $t_1$ to zero. Put $d_j=t_{j+1}-t_j$ and use the transfer
convention of Lemma~\ref{lem:collision-spectral-input}. The smoothed
joint expectation is
\[
 \ell\bigl(M_{u_{q,\epsilon}}\mathcal L_R^{d_{q-1}}
 M_{u_{q-1,\epsilon}}\cdots
 \mathcal L_R^{d_1}M_{u_{1,\epsilon}}\nu\bigr).
\]
Replace the power across the chosen gap by $\Pi_R=\nu\ell$.
The resulting expression is exactly the product of the two smoothed
expectations in \eqref{eq:bv-finite-product}. The difference is bounded
by $C_q\epsilon^{-2q}\vartheta^{d_p}$. Indeed each smooth multiplier
costs at most $C\epsilon^{-2}$, every other power is uniformly bounded,
and the complementary power across the gap costs $C\vartheta^{d_p}$.
For a gap of length zero use
$\mathcal L_R^0-\Pi_R=I-\Pi_R$, with its uniform bound.

Choose $\epsilon=\tfrac14\vartheta^{d_p/(2q+1)}$. Smoothing and
operator errors are at most $C_q\vartheta^{d_p/(2q+1)}$. The finite
parameter cover permits $\gamma=-\log\vartheta/9$ after changing $C$.
Invariance justifies the initial translation also for negative times.
\end{proof}

\begin{proposition}[Fourth moments without smoothing]
\label{prop:unsmoothed-fourth}
Let $u$ be real, centered, and $\|u\|_{\mathcal V}\le B$.
For every deterministic integer interval of length $m\ge1$,
\begin{equation}\label{eq:unsmoothed-fourth}
 \mathbb E\left|\sum_{j=r}^{r+m-1}u\circ T_R^j\right|^4\le C_Bm^2,
 \qquad
 \mathbb E\max_{0\le l\le m}
 \left|\sum_{j=r}^{r+l-1}u\circ T_R^j\right|^4\le C_Bm^2.
\end{equation}
Both estimates hold in the reverse order of the interval and under
any nonnegative initial density bounded by a fixed constant. They hold
componentwise, and hence in Euclidean norm, for $h_R$.
\end{proposition}
\begin{proof}
Put $c_u(a)=\mathbb E[u(u\circ T_R^a)]$. For $a,b,c\ge0$, apply
Lemma~\ref{lem:bv-finite-product} to four copies of $u$ at times
$0,a,a+b,a+b+c$. Splitting across the middle gap bounds
\[
 \left|\mathbb E[u(u\circ T_R^a)(u\circ T_R^{a+b})
                   (u\circ T_R^{a+b+c})]-c_u(a)c_u(c)\right|
       \le C_B e^{-\gamma b}.
\]
Splitting off the first or last centered factor bounds the fourfold
expectation by $C_Be^{-\gamma a}$ and $C_Be^{-\gamma c}$,
respectively. Also $|c_u(a)c_u(c)|\le C_Be^{-\gamma(a+c)}$.
Taking the minimum of these three bounds gives
\begin{align}\label{eq:four-product-gap}
 &\left|\mathbb E[u(u\circ T_R^a)(u\circ T_R^{a+b})
                   (u\circ T_R^{a+b+c})]-c_u(a)c_u(c)\right|
 \le C_B e^{-\gamma\max\{a,b,c\}}.
\end{align}
This includes repeated times.

Expand the fourth power and order its four indices; their multiplicity
is at most $24$. For each initial index the sum of
$e^{-\gamma(a+c)}$ over the two outer gaps is bounded, while the middle
gap has at most $m$ choices. This contributes $O(m^2)$. The sum over
all three nonnegative gaps of $e^{-\gamma\max\{a,b,c\}}$ is finite,
so the remaining contribution is $O(m)$. This proves the first bound.

For the second enlarge $m$ to a power of two $M<2m$. Every prefix is
a disjoint union of at most one dyadic interval of each length.
Consequently its absolute sum is bounded by the sum, over scales, of
the maximum absolute interval sum at that scale. At interval length
$l$, the fourth power of this maximum has expectation at most
$C_B(M/l)l^2=C_BMl$. Minkowski's inequality therefore gives an $L^4$
norm at most
\[
 C_B\sum_{j=0}^{\log_2 M}(M^2 2^{-j})^{1/4}
       \le C_B M^{1/2}.
\]
There is no logarithmic loss. Stationarity translates deterministic
intervals; reversal gives suffixes of the same interval. Domination
by a bounded initial density proves the remaining assertions.
\end{proof}

\subsection{Fourth-moment windows and two-sided stopping}
Retain $N^\pm_{j,R}$, $m_j=\lfloor j/c_*\rfloor$ and
$\mathsf H_{r,s,R}$ from Section~\ref{sec:marked-return-band}. Define
\begin{equation}\label{eq:moment-stopped-pair}
 Z_{n,k,R}=\mathsf H_{N^-_{k,R},N^+_{n-k,R},R},\qquad
 W_{n,k,R}=\mathsf H_{m_k,m_{n-k},R}.
\end{equation}
Thus $Z_{n,k,R}(y)=J_{n,R}((F_R^*)^{-k}y)-n\bar G_R$ almost
everywhere. Zero-length backward and forward intervals are empty.

\begin{lemma}[Fourth-moment return-clock window]
\label{lem:fourth-clock-window}
For $j\ge0$, $b\ge2$, and either sign,
\begin{equation}\label{eq:fourth-clock-window}
 \nu_R^*\{|N^\pm_{j,R}-m_j|>b\}
       \le C\frac{(j+b)^2}{b^4}.
\end{equation}
\end{lemma}
\begin{proof}
The case $j=0$ is immediate. The centered visit count
$\sum_{i=1}^L(\eta_R-c_*)\circ T_R^{\pm i}$ has fourth moment at
most $CL^2$ under $\nu_R^*$, by
Proposition~\ref{prop:unsmoothed-fourth}. The exact visit equivalences
in Lemma~\ref{lem:two-sided-clock-window} turn the upper and lower
return-clock events into deviations of these counts. At integer
thresholds within a bounded distance of $m_j\pm b$, the required
deviation is at least $c_*b-O(1)$ and $L\le C(j+b)$. Markov's
inequality proves \eqref{eq:fourth-clock-window}. A negative lower
threshold contributes no event. Increase $C$ for bounded $b$.
\end{proof}

\begin{theorem}[Moment control of the actual marked stopping]
\label{thm:marked-L2-stopping}
Uniformly for $n\ge2$, $0\le k\le n$, and $R\in I$,
\begin{align}
 \int (|Z_{n,k,R}|^4+|W_{n,k,R}|^4)\dd\nu_R^*&\le Cn^2,
       \label{eq:actual-marked-fourth}\\
 \left\|\frac{Z_{n,k,R}-W_{n,k,R}}{\sqrt n}
                  \right\|_{L^2(\nu_R^*)}&\le Cn^{-1/6}.
       \label{eq:actual-marked-L2}
\end{align}
For every real $v\in\R^4$ and insertion $a$ of
\eqref{eq:marked-norms}, the stopping comparison also satisfies
\begin{equation}\label{eq:improved-marked-stopping}
 \left|C^{[k],a}_{n,R}(v)
       -\int_M a e^{iv\cdot W_{n,k,R}/\sqrt n}\dd\nu\right|
       \le C M_a(1+|v|)n^{-2/9}.
\end{equation}
\end{theorem}
\begin{proof}
The deterministic interval has length at most $Cn$, so its fourth
moment follows from Proposition~\ref{prop:unsmoothed-fourth} and
$\nu_R^*\le c_*^{-1}\nu$. For the random interval put
$Q=N^-_{k,R}+N^+_{n-k,R}$. Invariance of the actual return map gives
$N^-_{j,R}\overset{\rm law}=N^+_{j,R}$ under $\nu_R^*$: use
$N^-_{j,R}((F_R^*)^jx)=N^+_{j,R}(x)$. Hence
Theorem~\ref{thm:cumulative-return-tail} and a union bound imply
\[
 \nu_R^*(Q>L)\le C e^{an-cL/2}.
\]
No independence of the two clocks is used. Fix $A>2a/c+2$.
On $Q\le An$, $|Z_{n,k,R}|$ is bounded by the sum of the backward
and forward collision maxima up to $\lceil An\rceil$. Their fourth
moments are $O(n^2)$. On $Q>An$, boundedness of $h_R$ gives
$|Z_{n,k,R}|\le C Q$. Integrating the displayed exponential tail
bounds that fourth-moment contribution by
$C(1+n)^4e^{-c' n}$ for some $c'>0$. This proves
\eqref{eq:actual-marked-fourth}.

Let $2\le b\le n$ and let $\mathcal E_b$ be the union of the two
clock-window exceptions. Lemma~\ref{lem:fourth-clock-window} gives
$\nu_R^*(\mathcal E_b)\le Cn^2/b^4$. Off this event, the difference
$Z_{n,k,R}-W_{n,k,R}$ is bounded by maxima on two deterministic
collision windows of length $O(b)$, including windows that cross zero.
Their second moments are $O(b)$ by the fourth-moment maximal estimate.
On the exceptional event, Cauchy--Schwarz and
\eqref{eq:actual-marked-fourth} give
\[
 \int |Z_{n,k,R}-W_{n,k,R}|^2\one_{\mathcal E_b}\dd\nu_R^*
       \le Cn\,\nu_R^*(\mathcal E_b)^{1/2}\le Cn^2/b^2.
\]
Thus the normalized squared $L^2$ error is at most
$C(b/n+n/b^2)$. Taking $b=\lceil n^{2/3}\rceil$ proves
\eqref{eq:actual-marked-L2}.

For the characteristic comparison, the exceptional event costs only
$CM_a n^2/b^4$. On its complement, the same deterministic windows and
$|e^{ix}-e^{iy}|\le|x-y|$ give $CM_a|v|\sqrt{b/n}$.
Choose instead $b=\lceil n^{5/9}\rceil$. Both powers are
$n^{-2/9}$. The recentering identity
\eqref{eq:marked-recentering} and $|a|\nu\le c_*M_a\nu_R^*$
complete the proof. All rounding exceptions are absorbed in $C$.
\end{proof}

\subsection{Quadratic insertions and the induced covariance}
The following estimate supplies the moment identification that is not
asserted by the characteristic-function proof in
Remark~\ref{rem:gaussian-versus-raw}.

\begin{lemma}[A centered insertion in a quadratic collision sum]
\label{lem:quadratic-collision-insertion}
Let $a$ be a complex bounded-variation function on $M$, and let
$\alpha=\int a\dd\nu$. For all $r,s\ge0$, with $m=r+s$,
\begin{align}
 \left|\int a\mathsf H_{r,s,R}\dd\nu\right|
       &\le C(\|a\|_\infty+\|a\|_{\mathrm{BV}}),
       \label{eq:linear-insertion-moment}\\
 \left\|\int a\mathsf H_{r,s,R}\otimes\mathsf H_{r,s,R}\dd\nu
                          -\alpha m\Gamma_R\right\|
       &\le C(\|a\|_\infty+\|a\|_{\mathrm{BV}}).
       \label{eq:quadratic-insertion-moment}
\end{align}
The constants do not depend on the location of the origin in the interval.
\end{lemma}
\begin{proof}
Center $a$ by subtracting $\alpha$. Each component of $h_R$ is
centered. The two-factor version of
Lemma~\ref{lem:bv-finite-product} bounds
$\int(a-\alpha)h_R\circ T_R^i\dd\nu$ by
$C\|a\|_{\mathcal V}e^{-\gamma|i|}$. Summation proves the first
assertion.

For components $u,v$ of $h_R$, sort the three times $0,i,j$ and
denote their two consecutive gaps by $d_1,d_2$. All three factors
$a-\alpha,u,v$ are centered. Splitting off the first and the last
factor in \eqref{eq:bv-finite-product} gives
\[
 \left|\int(a-\alpha)(u\circ T_R^i)(v\circ T_R^j)\dd\nu\right|
       \le C\|a\|_{\mathcal V}e^{-\gamma\max\{d_1,d_2\}}.
\]
Because $\max(d_1,d_2)\ge\max(|i|,|j|)/2$, the right side is
summable over all $(i,j)\in\Z^2$. The centered insertion therefore
changes the quadratic moment of any such interval by at most
$C\|a\|_{\mathcal V}$. Finally the exponentially summable collision
covariances give
\[
 \left\|\int\mathsf H_{r,s,R}\otimes\mathsf H_{r,s,R}\dd\nu
                                  -m\Gamma_R\right\|\le C.
\]
This is the finite covariance expansion with the lag weights $m-l$;
the sum of $l$ times the absolute lag covariances is finite. Multiplying
by $\alpha$ proves the second assertion. Real and imaginary parts
handle complex $a$ without any positivity assumption.
\end{proof}

\begin{theorem}[Moments at an arbitrary actual return mark]
\label{thm:marked-quadratic-moments}
Put $U_{n,R}=J_{n,R}-n\bar G_R$. For $a$ as in
\eqref{eq:marked-norms}, uniformly for $n\ge2$ and $0\le k\le n$,
\begin{align}
 \left|\int_{Y_R^*}a((F_R^*)^kx)\frac{U_{n,R}(x)}{\sqrt n}\dd\nu(x)\right|
 &\le C\left(M_an^{-1/6}+(M_a+V_a)n^{-1/2}\right),
       \label{eq:marked-first-moment}\\
 \left\|\int_{Y_R^*}a((F_R^*)^kx)
       \frac{U_{n,R}(x)\otimes U_{n,R}(x)}n\dd\nu(x)-\alpha_a D_R\right\|
 &\le C\left(M_an^{-1/6}+(M_a+V_a)n^{-1}\right).
       \label{eq:marked-second-moment}
\end{align}
These are moments of the finite complex measure of
\eqref{eq:marked-measure}, using $\nu|_{Y_R^*}$, not its normalization.
In particular
\begin{equation}\label{eq:actual-covariance-rate}
 \sup_{R\in I}\left\|\frac1n\int U_{n,R}\otimes U_{n,R}\dd\nu_R^*
                                               -D_R\right\|
       \le C n^{-1/6}.
\end{equation}
\end{theorem}
\begin{proof}
Recenter at $y=(F_R^*)^kx$. The weighted moments then involve
$a(y)Z_{n,k,R}(y)$ and $a(y)Z_{n,k,R}(y)\otimes Z_{n,k,R}(y)$.
Theorem~\ref{thm:marked-L2-stopping} changes the normalized first
moment by at most $CM_an^{-1/6}$. For the second use
\[
 \|z\otimes z-w\otimes w\|\le |z-w|(|z|+|w|)
\]
and Cauchy--Schwarz. Equations~\eqref{eq:actual-marked-fourth} and
\eqref{eq:actual-marked-L2} show that its normalized change is also
at most $CM_an^{-1/6}$. The factor $c_*$ in passing from
$\nu|_{Y_R^*}$ to $\nu_R^*$ is fixed.

Apply Lemma~\ref{lem:quadratic-collision-insertion} to
$W_{n,k,R}=\mathsf H_{m_k,m_{n-k},R}$. Its length satisfies
$m_k+m_{n-k}=n/c_*+O(1)$ uniformly in $k$. Since
$D_R=\Gamma_R/c_*$, division by $\sqrt n$ or $n$ gives
\eqref{eq:marked-first-moment} and \eqref{eq:marked-second-moment}.
For \eqref{eq:actual-covariance-rate} take $a=s_R=c_*^{-1}\eta_R$,
whose mass is one and whose two norms are uniformly bounded. Invariance
of $F_R^*$ and the Kac mean give $\int U_{n,R}\dd\nu_R^*=0$.
Thus the last display is convergence of the actual covariance matrices.
\end{proof}

\begin{corollary}[The induced Green--Kubo formula with Cesaro weights]
\label{cor:induced-cesaro-covariance}
Let $g_R=G_R-\bar G_R$ and
$C^*_{j,R}=\int g_R\otimes(g_R\circ(F_R^*)^j)\dd\nu_R^*$.
Then
\begin{equation}\label{eq:induced-cesaro-covariance}
 \left\|C^*_{0,R}+\sum_{j=1}^{n-1}\left(1-\frac jn\right)
          (C^*_{j,R}+(C^*_{j,R})^{\mathsf T})-D_R\right\|
       \le C n^{-1/6}.
\end{equation}
\end{corollary}
\begin{proof}
The exponential one-return moment makes each term finite. Expand the
second moment of $\sum_{j=0}^{n-1}g_R\circ(F_R^*)^j$ and use stationarity
under $F_R^*$. The left matrix before subtracting $D_R$ is exactly
$n^{-1}\int U_{n,R}\otimes U_{n,R}\dd\nu_R^*$. Apply
\eqref{eq:actual-covariance-rate}.
\end{proof}

\begin{corollary}[Moments under the unchanged marked-state event]
\label{cor:marked-conditional-moments}
Use the event $A_{n,k,R}$ and variation $V_{n,R}$ of
Corollary~\ref{cor:marked-state-conditioning}, and put
$p_{n,R}=\nu_R^*(E_{n,R})>0$. The conditional first moment of
$U_{n,R}/\sqrt n$ is bounded by
\[
 \frac C{p_{n,R}}\left(n^{-1/6}+(1+V_{n,R})n^{-1/2}\right).
\]
Its conditional second moment about the original center $n\bar G_R$
differs from $D_R$ by at most
\[
 \frac C{p_{n,R}}\left(n^{-1/6}+(1+V_{n,R})n^{-1}\right).
\]
Consequently, if $p_{n,R}\ge c n^{-\beta}$ and
$V_{n,R}\le Cn^\kappa$, the latter difference tends to zero when
$\beta<1/6$ and $\beta+\kappa<1$. The conditional mean tends to zero,
and the conditional covariance tends to $D_R$, under the stronger
conditions $\beta<1/6$ and $\beta+\kappa<1/2$.
\end{corollary}
\begin{proof}
Take $a=c_*^{-1}\one_{E_{n,R}}$ in
Theorem~\ref{thm:marked-quadratic-moments}. Its mass is $p_{n,R}$.
Divide the two estimates by this unchanged probability. Subtracting
the tensor square of the conditional mean identifies the conditional
covariance. The strict exponent inequalities prove convergence.
\end{proof}

\paragraph{Role in raw local inversion.}
Theorems~\ref{thm:marked-L2-stopping} and
\ref{thm:marked-quadratic-moments} concern the unsmoothed physical
record, not a replacement process. They identify its covariance by
actual second moments, establish uniform integrability of its squared
normalized size, and strengthen the stopping part of Theorem C.
In particular all the rare marked-state events in
Corollary~\ref{cor:marked-state-conditioning} also satisfy the new
conditional moment conclusions. Cesaro convergence in
\eqref{eq:induced-cesaro-covariance} does not assert absolute summability
of the induced correlations.

Positive definiteness is a distinct assertion. The estimates above give
no positive lower bound for the smallest eigenvalue of $D_R$ and no
periodic-evaluable representative of its zero-variance coboundary.
Theorem~\ref{thm:joint-uniform-nondegeneracy} supplies positive
definiteness by a separate measurable phase argument, not from these
moment bounds. Likewise \eqref{eq:improved-marked-stopping} is not a complementary
frequency estimate. The physical central radius remains
$2n^{-99/200}$, or $2n^{1/200}$ after rescaling. The annulus beyond it,
the complete raw critical and singular branch sum, and the weighted
tails for an exact physical observation retain their roles in
Theorem~\ref{thm:LLT}. No conditioning event has been replaced.

\section{Nondegeneracy of the collision-count component}
\label{sec:count-nondegeneracy}
\begin{corollary}[The count component is uniformly nondegenerate]
\label{cor:positive-count-variance}
Let $e_3=(0,0,1,0)$, the collision-count direction. There is a constant
$d_N>0$ such that
\begin{equation}\label{eq:positive-count-variance}
 e_3^{\mathsf T}D_Re_3\ge d_N\qquad(R\in I).
\end{equation}
In particular the actual normalized variance of $N_{n,R}$ is bounded
below by $d_N/2$ for all sufficiently large $n$, uniformly in $R$.
\end{corollary}
\begin{proof}
First, $T_R$ is mixing for the collision probability. The splitting in
Lemma~\ref{lem:collision-spectral-input} proves mixing for smooth
densities and smooth tests. Approximation in $L^2(\nu)$ and invariance
extend it to arbitrary $L^2$ functions: the approximation error in a
correlation is bounded, uniformly in its lag, by Cauchy--Schwarz.
Explicitly, if $f_\epsilon,g_\epsilon$ approximate $f,g$ in $L^2(\nu)$,
then, uniformly in $m$,
\[
 \left|\int\overline f(g\circ T_R^m)\dd\nu
       -\int\overline{f_\epsilon}(g_\epsilon\circ T_R^m)\dd\nu\right|
 \le\|f-f_\epsilon\|_2\|g\|_2
      +\|f_\epsilon\|_2\|g-g_\epsilon\|_2.
\]
The product of the two means converges under the same approximation.
This implication is used for each fixed $R$, without a rate for the
nonsmooth functions below.

Suppose $e_3^{\mathsf T}D_Re_3=0$. By
Theorem~\ref{thm:gaussian-kernel-coboundary}, there is a real
$b\in L^2(\nu)$ with
\[
 1-c_*^{-1}\eta_R=b-b\circ T_R
 \quad\hbox{almost everywhere}.
\]
Set $q=\exp(2\pi i c_*b)$, so $|q|=1$. Since
$2\pi c_*=91/5000$ and $\eta_R$ is integer-valued,
\[
 q\circ T_R=e^{-2\pi i c_*}q.
\]
Here $0<c_*<1$, so the multiplier $\lambda=e^{-2\pi i c_*}$ is
not one. Invariance gives $\int q\dd\nu=\lambda\int q\dd\nu$,
hence $\int q\dd\nu=0$. Mixing then gives
$\int\overline q(q\circ T_R^m)\dd\nu\to0$. But this integral is
exactly $\lambda^m$, whose modulus is one. This contradiction proves
strict positivity at every $R$. The continuity of $D_R$ and compactness
of $I$ give the uniform constant in \eqref{eq:positive-count-variance}.
Finally \eqref{eq:actual-covariance-rate} identifies the actual count
variance and gives the asserted lower bound for large $n$.

The argument uses the exact integer-valued section indicator and the
nonintegral section mass. It neither chooses a pointwise representative
at a periodic orbit nor, by itself, proves positive definiteness in
every joint direction. The latter conclusion is proved separately in
Theorem~\ref{thm:joint-uniform-nondegeneracy}.
\end{proof}


\section{Measurable collision phases and symplectic area}
\label{sec:measurable-phase-rigidity}
The measurable coboundary in Theorem~\ref{thm:gaussian-kernel-coboundary}
need not have a value at a prescribed periodic orbit. We shall instead
use conditional pairs on a positive-measure hyperbolic product set.
The discontinuous section contribution will be removed algebraically
before this argument is applied. Throughout this section $R$ is fixed;
none of the local product constants is asserted uniform in $R$.

Write $\kappa_R$ for the one-collision lattice label, in the basis
$a=(1,0)$, $b=(1/2,\sqrt3/2)$. Thus
$f_R=(\kappa_{R,1},\kappa_{R,2},1,\tau_R)$. All identities between
measurable functions are understood almost everywhere with respect to
$\nu$. Let $\mathbb S^1=\{z\in\mathbb C:|z|=1\}$.

\subsection{A qualitative local product input}
\begin{lemma}[Regular product sets for the collision map]
\label{lem:regular-collision-product}
There is a compact positive-$\nu$-measure set $\mathcal P$ in a regular
collision coordinate rectangle, away from grazing, with the following
properties. It has product coordinates
$\Psi:K^u\times K^s\longrightarrow\mathcal P$. Write a product point as $\Psi(u,s)$. Fixing $u$ places points on one
local homogeneous stable curve; fixing $s$ places them on one local
homogeneous unstable curve. Thus $K^u$ labels stable leaves and $K^s$
labels unstable leaves. These curves are $C^1$, have transverse tangent cones,
and each stable curve meets each unstable curve exactly once.

The probability $m=\nu(\mathcal P)^{-1}\nu|_{\mathcal P}$ is equivalent,
in these coordinates, to a product $m^u\otimes m^s$ of nonatomic
transversal probabilities. A stable arc joining two points on the same
curve remains a connected regular arc under every forward collision
iterate, and its length tends to zero. The corresponding statement
holds for unstable arcs under backward iterates. In particular the
successive lattice labels agree along the respective arcs.

For every sufficiently small $\varepsilon>0$, the product coordinates
can be restricted to positive-measure subsets so that every four-corner
stable--unstable quadrilateral lies in a coordinate box of diameter
less than $\varepsilon$. Almost every such quadrilateral with distinct
coordinates is a simple, nondegenerate Jordan domain.
\end{lemma}
\begin{proof}
We use the qualitative billiard product construction in
\cite[Sections 8.2--8.3]{Young}. Section 8.2B gives transverse stable
and unstable cones. Homogeneous local curves and their regular iterates
are described in Section 8.2D and constructed in Section 8.3,
Sublemma 1. The product set chosen there has positive invariant area;
this is stated at the end of its construction, before the verification
of (P3)--(P5). The absolute-continuity formula (P5)(b), verified there,
has positive finite Jacobian. Apply the same local statement in the
reverse time direction. Since the invariant collision measure is
smooth and positive away from grazing, disintegration and the two
holonomy formulas identify its null sets on the product rectangle with
those of the product of the transversal measures. These measures are
nonatomic, being equivalent to restricted arc length on transversals.
Restriction to compact positive-measure subsets gives the stated form.
No assertion that $\mathcal P$ contains an open set is needed.

For completeness, contraction here means contraction of the actual
curves, not just their endpoints in a symbolic metric. The billiard
estimates of \cite[Section 8.2C]{Young} give exponential contraction of
the stable arcs in its adapted metric and
$\operatorname{length}(\gamma)\le C\sqrt{\operatorname{length}_{\rm ad}(\gamma)}$
for a stable or unstable curve. Thus, for fixed $R$ and the chosen product set, the ordinary lengths
of these iterated arcs are bounded by $C_{\mathcal P}\theta_{\mathcal P}^n$
for some $C_{\mathcal P}<\infty$ and $0<\theta_{\mathcal P}<1$.
In particular they tend to zero. No uniform choice in $R$ is made. Passing from $(\alpha,\varphi)$ to $(\alpha,p)$ uses only
$p=\sin\varphi$, whose derivative is bounded; no inverse change at
grazing is made. Homogeneous local curves do not cross the relevant
finite-iterate collision singularities. The label $\kappa_R$ is
locally constant on each such physical branch, which proves the label
assertion.

The hypotheses of that construction hold here: the disks are disjoint
with separation at least $3/50$, their boundaries are smooth with
strictly positive curvature, the free domain is connected, and the
horizon is finite. These are properties of the physical Euclidean
table. The same local construction works on its flat triangular torus;
alternatively the rectangular two-scatterer cover in
Lemma~\ref{lem:collision-spectral-input} gives the local product set in
a covering chart. There is no affine change of the reflection law.

Finally choose a product support point and restrict each transversal
to a sufficiently short positive-measure interval around it. Continuity
of the local product map and the transverse graph cones put the
connecting arcs into a small coordinate box. Unique intersections of
the graph curves make the four distinct corners bound a simple Jordan
domain. Nonatomicity excludes repeated transversal coordinates almost
everywhere. This also proves the last assertion for product sets with
gaps.
\end{proof}


\paragraph{Conditional pairs in product coordinates.}
To fix the measure conventions, write $\mu^u=m^u$, $\mu^s=m^s$,
$\mu=\mu^u\otimes\mu^s$ and
\[
 \dd m_{\mathcal P}=\rho(u,s)\dd\mu^u(u)\dd\mu^s(s),\qquad
 m_{\mathcal P}=\nu(\mathcal P)^{-1}\nu|_{\mathcal P}.
\]
Here $0<\rho<\infty$ almost everywhere and $\int\rho\dd\mu=1$.
Put $\rho_u(u)=\int\rho(u,s)\dd\mu^s(s)$ and
$\rho_s(s)=\int\rho(u,s)\dd\mu^u(u)$. The stable and unstable pair
probabilities are explicitly
\begin{align}\label{eq:explicit-conditional-pair-densities}
 \dd\mathfrak m_s(u,s_0,s_1)
 &=\frac{\rho(u,s_0)\rho(u,s_1)}{\rho_u(u)}
           \dd\mu^u(u)\dd\mu^s(s_0)\dd\mu^s(s_1),\\
 \dd\mathfrak m_u(u_0,u_1,s)
 &=\frac{\rho(u_0,s)\rho(u_1,s)}{\rho_s(s)}
           \dd\mu^u(u_0)\dd\mu^u(u_1)\dd\mu^s(s).\notag
\end{align}
The denominators are positive and finite almost everywhere. Integration
of either unused coordinate shows that both marginals are
$m_{\mathcal P}$. The displayed positive finite densities also give
exactly the triple-product measure equivalences used below, without
a boundedness assertion for $\rho$ or its reciprocal.

\subsection{The physical action one-form}
On the collision cylinder define
\begin{equation}\label{eq:collision-action-form}
 \vartheta_R=Rp\,\dd\alpha,\qquad
 \dd\vartheta_R=R\,\dd p\wedge\dd\alpha.
\end{equation}
The one-form is globally defined on the cylinder and bounded in both
the flat and the original collision coordinates.

\begin{lemma}[Exact action on a regular branch]
\label{lem:exact-collision-action}
On every regular collision branch,
\begin{equation}\label{eq:exact-collision-action}
 \dd\tau_R=T_R^*\vartheta_R-\vartheta_R.
\end{equation}
Consequently, if an arc $\gamma$ from $x$ to $y$ stays on regular
branches through time $n$, then
\begin{equation}\label{eq:action-telescope-arc}
 S_{n,R}\tau_R(y)-S_{n,R}\tau_R(x)
       =\int_{T_R^n\gamma}\vartheta_R-\int_\gamma\vartheta_R.
\end{equation}
For backward regular arcs the corresponding identity is
\begin{equation}\label{eq:backward-action-arc}
 S_{n,R}\tau_R(T_R^{-n}y)-S_{n,R}\tau_R(T_R^{-n}x)
       =\int_\gamma\vartheta_R-\int_{T_R^{-n}\gamma}\vartheta_R.
\end{equation}
\end{lemma}
\begin{proof}
Lift the flight to its actual initial disk and the next disk with fixed
label $d$. Its endpoints are $q=Rn_\alpha$ and
$q'=d+Rn_{\alpha'}$. If $v$ is its unit velocity, differentiation of
its Euclidean length gives
$\dd\tau_R=v\cdot\dd q'-v\cdot\dd q$.
The first tangential component is $p$, and specular reflection preserves
the tangential component at arrival, which is $p'$ in the outgoing
coordinates of $T_Rx$. Hence
$\dd\tau_R=Rp'\dd\alpha'-Rp\dd\alpha$. The lattice label is constant
on the branch and contributes no derivative. This proves
\eqref{eq:exact-collision-action}. Summing its pullbacks and integrating
over the arc proves \eqref{eq:action-telescope-arc}; applying that
identity to $T_R^{-n}\gamma$ proves \eqref{eq:backward-action-arc}.
\end{proof}

\begin{proposition}[A measurable physical phase has zero roof frequency]
\label{prop:measurable-roof-rigidity}
Suppose $q:M\to\mathbb S^1$ is measurable and, for constants
$u\in\R^2$ and $s,t\in\R$,
\begin{equation}\label{eq:physical-measurable-phase}
 q\circ T_R=\exp\{i(u\cdot\kappa_R+s+t\tau_R)\}\,q.
\end{equation}
Then $t=0$.
\end{proposition}
\begin{proof}
Discard an invariant null set so that every finite iterate of the phase
identity holds at the remaining points. Use the product set and the
probability $m$ of Lemma~\ref{lem:regular-collision-product}.
Disintegrate $m$ over its stable curves and, conditionally on the curve,
draw two independent points $x,y$. Denote this pair probability by
$\mathfrak m_s$. Both of its marginals are exactly $m$; in particular
they are bounded by $\nu(\mathcal P)^{-1}\nu$.

The following argument justifies endpoint agreement for a merely
measurable $q$. Let $q_\epsilon$ be a smooth approximation in $L^1(\nu)$,
with supremum norm at most one. Invariance of $\nu$ and the marginal
bound give, for every $n$,
\begin{align}\label{eq:conditional-pair-approximation}
 &\int |q(T_R^ny)-q(T_R^nx)|\dd\mathfrak m_s(x,y)\notag\\
 &\quad\le \frac{2}{\nu(\mathcal P)}\|q-q_\epsilon\|_{L^1(\nu)}
       +\int|q_\epsilon(T_R^ny)-q_\epsilon(T_R^nx)|\dd\mathfrak m_s(x,y).
\end{align}
For fixed $\epsilon$, the last term tends to zero by contraction and
bounded convergence. Then let $\epsilon$ tend to zero. Thus the left
side tends to zero. This uses bounded marginals, not invariance of the
pair probability or a claim about every stable curve.

Let $\gamma_s(x,y)$ be the stable arc oriented from $x$ to $y$.
Its forward labels agree at all times. Iteration of
\eqref{eq:physical-measurable-phase} and
\eqref{eq:action-telescope-arc} gives
\[
 \frac{q(T_R^ny)}{q(T_R^nx)}
 =\exp\!\left(it\left[\int_{T_R^n\gamma_s(x,y)}\vartheta_R
                         -\int_{\gamma_s(x,y)}\vartheta_R\right]\right)
                         \frac{q(y)}{q(x)}.
\]
The first arc integral tends to zero, because $\vartheta_R$ is bounded
and the arc length tends to zero. Since $|q|=1$, the left side minus one
has $L^1(\mathfrak m_s)$ norm tending to zero by
\eqref{eq:conditional-pair-approximation}. The right side has the
pointwise limit just displayed. Fatou's lemma therefore gives
\begin{equation}\label{eq:measurable-stable-holonomy}
 \frac{q(y)}{q(x)}
       =\exp\!\left(it\int_{\gamma_s(x,y)}\vartheta_R\right)
       \quad\text{for }\mathfrak m_s\text{-almost every }(x,y).
\end{equation}
Disintegrating instead over unstable curves and using backward iterates
in \eqref{eq:conditional-pair-approximation} gives, by
\eqref{eq:backward-action-arc},
\begin{equation}\label{eq:measurable-unstable-holonomy}
 \frac{q(y)}{q(x)}
       =\exp\!\left(it\int_{\gamma_u(x,y)}\vartheta_R\right)
       \quad\text{for almost every unstable pair}.
\end{equation}
The sign is the same in the two formulas: the backward action in
\eqref{eq:backward-action-arc} has limiting value
$\int_{\gamma_u(x,y)}\vartheta_R$.

We now use the measure-class part of the local product input. The stable
and unstable pair probabilities are equivalent, in product coordinates,
to the corresponding three-coordinate product measures. Fubini's theorem
therefore permits \eqref{eq:measurable-stable-holonomy} and
\eqref{eq:measurable-unstable-holonomy} to be multiplied around almost
every four-corner product quadrilateral $Q$. The oriented corner order is
\begin{gather*}
 x_{00}=\Psi(u_0,s_0)\ \longrightarrow\
 x_{01}=\Psi(u_0,s_1)\ \longrightarrow\ x_{11}=\Psi(u_1,s_1),\\
 x_{11}\ \longrightarrow\ x_{10}=\Psi(u_1,s_0)
            \ \longrightarrow\ x_{00}.
\end{gather*}
The edges are stable, unstable, stable with reversed $s$-orientation,
and unstable with reversed $u$-orientation. This specifies the boundary
orientation used in Stokes' formula; its sign need not be fixed when
only the absolute area is used. Their left sides telescope, so
\begin{equation}\label{eq:measurable-area-obstruction}
 \exp\!\left(it\oint_{\partial Q}\vartheta_R\right)=1.
\end{equation}
This assertion is on a positive-measure family of quadrilaterals; it does
not assign values to $q$ on a preselected closed orbit.

Suppose $t\ne0$. Restrict to a positive-measure product subrectangle
small enough that every such quadrilateral has coordinate area less
than $2\pi/(R|t|)$. Almost every quadrilateral is a nondegenerate Jordan
domain. Stokes' formula and \eqref{eq:collision-action-form} give
\[
 0<\left|t\oint_{\partial Q}\vartheta_R\right|
       =|t|R\operatorname{area}_{\alpha,p}(Q)<2\pi,
\]
contradicting \eqref{eq:measurable-area-obstruction}. Hence $t=0$.
\end{proof}

\subsection{Rotations and the constant collision phase}
Let $S_j(\alpha,p)=(\alpha+j\pi/3,p)$, $j\in\Z$. In lattice
coordinates rotation by $\pi/3$ is
\begin{equation}\label{eq:triangular-rotation-matrix}
 A=\begin{pmatrix}0&-1\\1&1\end{pmatrix},\qquad
 A^3=-I_2,\qquad I_2+A^2+A^4=0.
\end{equation}
The actual physical symmetries give
\begin{equation}\label{eq:collision-rotation-identities}
 T_RS_j=S_jT_R,\qquad
 \kappa_R\circ S_j=A^j\kappa_R,\qquad
 \tau_R\circ S_j=\tau_R,\qquad (S_j)_*\nu=\nu.
\end{equation}
They follow by rotating both endpoints, the velocity and the selected
next disk. The first admissible physical root remains the first root.
No invariance of $Y_R^*$ under these rotations is asserted.

\begin{theorem}[Rigidity of the roof and constant parts of a phase]
\label{thm:physical-phase-rigidity}
If a measurable circle-valued $q$ satisfies
\eqref{eq:physical-measurable-phase}, then
\begin{equation}\label{eq:physical-phase-conclusion}
 t=0,\qquad s\in2\pi\Z.
\end{equation}
\end{theorem}
\begin{proof}
The proposition gives $t=0$. Define
\[
 Q_2=q\,(q\circ S_3),\qquad
 Q_3=q\,(q\circ S_2)(q\circ S_4).
\]
By \eqref{eq:triangular-rotation-matrix} and
\eqref{eq:collision-rotation-identities}, the lattice phases cancel
exactly and
\[
 Q_2\circ T_R=e^{2is}Q_2,\qquad Q_3\circ T_R=e^{3is}Q_3.
\]
Both functions have modulus one. The mixing collision map has no
nontrivial circle eigenvalue, by the $L^2$ argument in the proof of
Corollary~\ref{cor:positive-count-variance}. Therefore
$e^{2is}=e^{3is}=1$, and their quotient gives $e^{is}=1$.
\end{proof}

The theorem is qualitative. It excludes exact measurable physical
phases with nonzero roof frequency or a nonintegral constant phase.
It is not a norm estimate for approximate spectral vectors, and supplies
no decay rate for a high-frequency collision or induced resolvent.

\section{Uniform nondegeneracy of the joint return record}
\label{sec:joint-nondegeneracy}
We now combine measurable phase rigidity with the exact integer-valued
section compensation. The proof treats every joint direction of the
same four-coordinate record. It does not pass from the scalar count
variance to a determinant by an inequality.

\subsection{A finite-cover obstruction for displacement}
\begin{lemma}[No real spatial collision coboundary]
\label{lem:no-spatial-coboundary}
If $w\in\R^2$ and $b\in L^2(\nu;\R)$ satisfy
\begin{equation}\label{eq:spatial-coboundary}
 w\cdot\kappa_R=b-b\circ T_R,
\end{equation}
then $w=0$.
\end{lemma}
\begin{proof}
Regard $w$ as a row vector and suppose $w\ne0$.
Composition with $S_1$, using \eqref{eq:collision-rotation-identities},
gives $(wA)\cdot\kappa_R=b\circ S_1-(b\circ S_1)\circ T_R$.
The matrix with rows $w=(w_1,w_2)$ and
$wA=(w_2,-w_1+w_2)$ has determinant
\begin{equation}\label{eq:spatial-rotation-determinant}
 -(w_1^2-w_1w_2+w_2^2),
 \qquad w_1^2-w_1w_2+w_2^2\ge\tfrac12|w|^2>0.
\end{equation}
Solving these two real identities gives a vector-valued
$H\in L^2(\nu;\R^2)$ with $\kappa_R=H-H\circ T_R$.

Consider the physical billiard on the flat torus $\R^2/(2\Lambda)$.
There are four equal scatterers, labeled by
$\ell\in\Lambda/(2\Lambda)\simeq(\Z/2\Z)^2$. In collision
coordinates its map and invariant probability are exactly
\begin{equation}\label{eq:four-sheet-collision}
 \widetilde T_R(x,\ell)
       =(T_Rx,\ell+\kappa_R(x)\bmod2),\qquad
 \widetilde\nu=\nu\otimes\operatorname{unif}((\Z/2\Z)^2).
\end{equation}
This map is ergodic. The planar lift has precisely the original disk
array, so a free flight in the cover has the original uniform length
bound. The enlarged torus contains finitely many pairwise disjoint
embedded disks; their complement is connected. Thus the enlarged table
has smooth strictly convex scatterers, connected free space and finite
horizon; the finite-horizon dispersing-billiard theorem
\cite[Section 8.1, Theorem 6]{Young} applies. To work only with a
rectangular fundamental domain, take the physical cover with periods
$(2,0)$ and $(0,2\sqrt3)$. It has eight scatterers and satisfies the
same hypotheses. Its quotient by the deck group is precisely the
four-scatterer table above, so its ergodicity implies that of
\eqref{eq:four-sheet-collision}. Neither argument changes Euclidean
length or specular reflection.

However,
\[
 F(x,\ell)=\exp\{i\pi(H_1(x)+\ell_1)\}
\]
is invariant under \eqref{eq:four-sheet-collision}: the integer
increment $\kappa_{R,1}$ cancels $H_1\circ T_R-H_1$, and reduction
of $\ell_1$ modulo two does not alter the exponential. The function
has modulus one and integral zero by averaging over $\ell_1$.
It cannot be almost everywhere constant, contradicting ergodicity.
This proves the lemma. No value of $b$ or $H$ at a periodic point is
used.
\end{proof}

\subsection{Removing the section contribution}
\begin{theorem}[Uniform positive definiteness]
\label{thm:joint-uniform-nondegeneracy}
There exist constants $0<d_0\le D_0<\infty$ such that
\begin{equation}\label{eq:joint-uniform-nondegeneracy}
 d_0I_4\le D_R\le D_0I_4\qquad(R\in I).
\end{equation}
Equivalently the actual $L^2$ coboundaries in
Theorem~\ref{thm:gaussian-kernel-coboundary} have only the zero joint
direction. The discontinuity of $\eta_R=\one_{Y_R^*}$ is allowed.
\end{theorem}
\begin{proof}
Fix $R$ and suppose $v^{\mathsf T}D_Rv=0$, where $v=(w,a,b)$ with
$w\in\R^2$. The collision-clock coboundary characterization gives a
real $B\in L^2(\nu)$ such that
\begin{equation}\label{eq:joint-zero-coboundary}
 w\cdot\kappa_R+a+b\tau_R-\zeta\eta_R=B-B\circ T_R,
 \qquad \zeta=\frac{a+b\bar\tau_R}{c_*}.
\end{equation}
All terms are the physical single-collision quantities. We separate
whether the coefficient of the section indicator vanishes.

If $\zeta\ne0$, set $q=\exp(2\pi i B/\zeta)$. Since $\eta_R$ is
integer-valued, exponentiation removes it exactly and yields
\[
 q\circ T_R
 =\exp\!\left(-\frac{2\pi i}{\zeta}
          (w\cdot\kappa_R+a+b\tau_R)\right)q.
\]
Theorem~\ref{thm:physical-phase-rigidity} implies $b=0$ and
$-2\pi a/\zeta\in2\pi\Z$. But now
$\zeta=a/c_*\ne0$, so $a/\zeta=c_*$. Since $0<c_*<1$, this is
impossible. Notice that the rotations are applied only after the
section indicator has disappeared. Their use requires no symmetry of
the actual return section.

If $\zeta=0$, set $q=\exp(iB)$ in
\eqref{eq:joint-zero-coboundary}. Proposition~\ref{prop:measurable-roof-rigidity}
implies $b=0$. The definition of $\zeta$ then gives $a=0$.
The remaining identity is the real spatial coboundary
$w\cdot\kappa_R=B-B\circ T_R$, and
Lemma~\ref{lem:no-spatial-coboundary} gives $w=0$.

Thus $v=0$ in all cases. The symmetric matrix $D_R$ is positive definite
for each $R$. Theorem~\ref{thm:collision-covariance} gives its continuity
on the compact interval $I$. The smallest eigenvalue is consequently
bounded below by a positive number, and the largest is bounded above.
This proves \eqref{eq:joint-uniform-nondegeneracy}. Uniformity is obtained
at this last step, not by assuming uniform constants for measurable
phase regularity or for the product set.
\end{proof}

\begin{remark}[The role of the section]
The algebra in \eqref{eq:joint-zero-coboundary} uses only that the
section indicator is integer-valued and that its mean lies strictly
between zero and one. The two rotation products force an integral
constant phase without an arithmetic assumption on this mean. The
specific section geometry is used earlier to establish the bounded
collision covariance and its $L^2$ kernel characterization. Thus the
present proof neither evaluates a discontinuous section weight on
stable arcs nor replaces it by a smoother section.
\end{remark}

\begin{corollary}[Actual covariances and the physical-time projection]
\label{cor:actual-uniform-ellipticity}
For all sufficiently large $n$, uniformly in $R$,
\begin{equation}\label{eq:finite-n-joint-ellipticity}
 \frac1n\int U_{n,R}\otimes U_{n,R}\dd\nu_R^*
       \ge\frac{d_0}{2}I_4.
\end{equation}
The physical-time fluctuation covariance $\mathcal V_R$ in
Corollary~\ref{cor:physical-functional-gaussian} is also uniformly
positive definite on its three-dimensional space.
\end{corollary}
\begin{proof}
Equation~\eqref{eq:actual-covariance-rate} and
\eqref{eq:joint-uniform-nondegeneracy} give
\eqref{eq:finite-n-joint-ellipticity}. For the physical covariance recall
$\Gamma_R=c_*D_R$ and
$\mathcal V_R=\bar\tau_R^{-1}B_R\Gamma_RB_R^{\mathsf T}$, where
$B_R(x_1,x_2,x_3,x_4)=(x_1,x_2,x_3-x_4/\bar\tau_R)$.
For $z\in\R^3$, $|B_R^{\mathsf T}z|\ge|z|$. Hence
$z^{\mathsf T}\mathcal V_Rz\ge
 c_*d_0(\max_I\bar\tau_R)^{-1}|z|^2$. All constants are positive.
\end{proof}

\subsection{The Gaussian part of exact local inversion}
\begin{corollary}[Uniform Gaussian inversion and the remaining terms]
\label{cor:elliptic-central-inversion}
The centered Gaussian density associated with the actual return record is
\begin{equation}\label{eq:uniform-joint-gaussian-density}
 g_R(\xi)=\frac{1}{(2\pi)^2\sqrt{\det D_R}}
              \exp\!\left(-\tfrac12\xi^{\mathsf T}D_R^{-1}\xi\right).
\end{equation}
It is defined for every $R\in I$. There are $C,c>0$ such that
\begin{equation}\label{eq:elliptic-gaussian-tail}
 \sup_R\int_{|v|>n^{1/200}}
          e^{-v^{\mathsf T}D_Rv/2}\dd v
       \le C\exp(-c n^{1/100}).
\end{equation}
In Corollary~\ref{cor:raw-from-growing-arc}, uniform positive definiteness
is therefore a proved conclusion for this physical family, not an
additional hypothesis. Under its stated residual integrability,
complementary-transform estimate and local-edge conditions, the raw
mixed-density conclusion follows with the density
\eqref{eq:uniform-joint-gaussian-density}.
\end{corollary}
\begin{proof}
Positive definiteness gives the Gaussian inversion formula. In four
dimensions, integration in polar coordinates gives, for $A>0$,
\[
 \int_{|v|>A}e^{-d_0|v|^2/2}\dd v
   =2\pi^2\left(\frac{A^2}{d_0}+\frac{2}{d_0^2}\right)
                         e^{-d_0A^2/2}.
\]
Absorb the polynomial into a weaker exponential and put $A=n^{1/200}$.
The last assertion follows from the exact decomposition and the proved
central comparison in Corollary~\ref{cor:raw-from-growing-arc}.
\end{proof}

The two Fourier tails in this last argument concern different objects.
Equation~\eqref{eq:elliptic-gaussian-tail} estimates the limiting Gaussian.
It does not estimate the complementary transform of the physical raw
record or of its edge-subtracted residual. The proved central band is
still $|v|\le2n^{1/200}$, equivalently
$|\omega|\le2n^{-99/200}$. The annulus beyond it, the full critical and
singular edge extraction, and the weighted tails and relative event
comparison for exact physical conditioning remain the analytical tasks
in the raw mixed-density theorem. The covariance obstruction has been
removed without changing that theorem's record or target measure.

\section{Complete physical phase arithmetic}
\label{sec:complete-phase-arithmetic}
The roof and constant parts of a physical collision phase have already
been determined in Theorem~\ref{thm:physical-phase-rigidity}. We now
determine its two spatial parts. The argument uses an atom of the phase
on a positive-measure product set, not a value at a periodic point.
It applies to arbitrary real spatial frequencies, including irrational
ones.

Write
\[
 \Omega=\T^2\times\T\times\R,\qquad
 \phi_{R,\omega}=u\cdot\kappa_R+s+t\tau_R,
 \quad \omega=(u,s,t).
\]
The three torus coordinates have period $2\pi$. The origin of $\Omega$
therefore means $u\in2\pi\Z^2$, $s\in2\pi\Z$, $t=0$.

\begin{lemma}[Separation of a proper integer subgroup]
\label{lem:integer-subgroup-separation}
If $L$ is a proper subgroup of $\Z^d$, then there are a prime $p$ and a
nonzero $a\in(\Z/p\Z)^d$ with $a\cdot L=0$ modulo $p$.
\end{lemma}
\begin{proof}
An integer subgroup has an integer basis. If its rank is less than $d$,
choose a nonzero integer row annihilating it and a prime not dividing
all coordinates of that row. If it has full rank, its index is the
absolute determinant of a basis matrix. This integer exceeds one.
Modulo a prime divisor of the index, the basis matrix is singular and
has a nonzero left null vector. These two constructions prove the
claim without a rationality assumption on the character that might
have defined $L$.
\end{proof}

\begin{theorem}[Complete measurable physical phase rigidity]
\label{thm:complete-physical-phase}
For every $R\in I$, a measurable $q:M\to\mathbb S^1$ satisfies
\begin{equation}\label{eq:complete-physical-phase}
 q\circ T_R=e^{i\phi_{R,\omega}}q
\end{equation}
if and only if $\omega=0$ in $\Omega$ and $q$ is constant almost
everywhere.
\end{theorem}
\begin{proof}
Theorem~\ref{thm:physical-phase-rigidity} gives $t=0$ and
$s\in2\pi\Z$. Both holonomy formulas
\eqref{eq:measurable-stable-holonomy} and
\eqref{eq:measurable-unstable-holonomy} consequently have right side
one. In the product coordinates of
Lemma~\ref{lem:regular-collision-product}, the first says that $q$
is almost everywhere a function only of the unstable transversal
coordinate, and the second that it is almost everywhere a function
only of the stable transversal coordinate. Product measure equivalence
and Fubini imply that $q=c$ almost everywhere on $\mathcal P$, for
one $c\in\mathbb S^1$.

The saturation $\bigcup_{j\ge0}T_R^j\mathcal P$ has full measure by
ergodicity and recurrence. Iterating \eqref{eq:complete-physical-phase}
from $\mathcal P$ therefore shows that
\begin{equation}\label{eq:countable-phase-range}
 q(x)/c\in\{e^{iu\cdot k}:k\in\Z^2\}
       \quad\hbox{for almost every }x.
\end{equation}
Let $L=\{k\in\Z^2:u\cdot k\in2\pi\Z\}$. The map
$[k]\mapsto e^{iu\cdot k}$ identifies the countable group $\Z^2/L$
with the displayed range. Its inverse on that countable range is
measurable. Thus \eqref{eq:countable-phase-range} defines a measurable
$h:M\to\Z^2/L$ with
\begin{equation}\label{eq:countable-phase-cocycle}
 h(T_Rx)=h(x)+[\kappa_R(x)].
\end{equation}
This is a countable-valued lift, not a real logarithm with a chosen
branch. No integrability of the lift is required.

If $L\ne\Z^2$, apply Lemma~\ref{lem:integer-subgroup-separation}.
The resulting nonzero row $a$ defines a homomorphism
$\chi:\Z^2/L\to\Z/p\Z$. On the actual physical billiard modulo
$p\Lambda$ the collision map is
\[
 \widetilde T_R(x,\ell)
 =(T_Rx,\ell+\kappa_R(x)\bmod p),\qquad
 \ell\in(\Z/p\Z)^2.
\]
It preserves $\nu$ times the uniform sheet probability and is ergodic.
Indeed its lift to the plane has exactly the original disk array, so
its free-flight bound is unchanged. Its finitely many disjoint circular
scatterers have positive curvature, and the complement of these disks
in the torus is connected. Thus the finite-horizon billiard theorem
\cite[Section 8.1]{Young} applies. Equivalently use the rectangular
Euclidean cover with periods $(p,0)$ and $(0,p\sqrt3)$, with $2p^2$
scatterers, and then its quotient. No affine deformation or independent
sheet dynamics is introduced.

Equation~\eqref{eq:countable-phase-cocycle} makes
\[
 \exp\!\left(\frac{2\pi i}{p}
             [a\cdot\ell-\chi(h(x))]\right)
\]
invariant under this collision map. The brackets in this formula are
interpreted modulo $p$. The function has modulus one and mean zero
upon averaging over $\ell$, since $a\ne0$. This contradicts ergodicity.
Hence $L=\Z^2$, or $u\in2\pi\Z^2$. The original phase equation now
says that $q$ is invariant, so it is constant. The converse is immediate.
\end{proof}

\subsection{Finite positive-measure return witnesses}
This consequence will replace a nonquantitative appeal to absence of
exact eigenfunctions in the next section. The auxiliary set below is
not the return section $Y_R^*$ of the physical record.

\begin{lemma}[A finite generating family of product-set returns]
\label{lem:finite-return-witnesses}
Fix $R$ and a positive-measure regular product set $\mathcal P$.
Let $r_{\mathcal P}$ be its actual first-return time under $T_R$.
There are finitely many pairs $(n_i,k_i)\in\Z_{>0}\times\Z^2$ such that
\begin{equation}\label{eq:positive-return-witness}
 E_i=\{x\in\mathcal P:r_{\mathcal P}(x)=n_i,
                  S_{n_i,R}\kappa_R(x)=k_i\},\qquad \nu(E_i)>0,
\end{equation}
and the vectors $g_i=(k_{i,1},k_{i,2},n_i)$ generate $\Z^3$ over $\Z$.
In particular, integers $b_{ji}$ can be chosen with
$e_j=\sum_i b_{ji}g_i$ for $1\le j\le3$.
\end{lemma}
\begin{proof}
Let $L$ be the subgroup generated by all positive-measure cells in
\eqref{eq:positive-return-witness}. The cells form a countable partition
of $\mathcal P$ modulo a null set. If $L$ were proper,
Lemma~\ref{lem:integer-subgroup-separation} would give a nontrivial
finite character $\theta=(u,s)\in\T^3$ annihilating every cell record.
For $\phi=u\cdot\kappa_R+s$ this means
$e^{iS_{r_{\mathcal P},R}\phi}=1$ on $\mathcal P$ almost everywhere.

The first-return tower over $\mathcal P$ represents almost every point
of the invertible ergodic collision system uniquely as $T_R^j y$,
where $y\in\mathcal P$ and $0\le j<r_{\mathcal P}(y)$.
Set $q(T_R^j y)=e^{iS_{j,R}\phi(y)}$ on these measurable tower levels.
The preceding return identity verifies the phase equation also at
the top of each level. This is an exact measurable physical phase at
$(u,s,0)\ne0$, contradicting
Theorem~\ref{thm:complete-physical-phase}. Thus $L=\Z^3$.
Writing the three standard basis vectors as integer combinations uses
only finitely many cell records, which gives the asserted family.
The witnesses have positive measure; they are not isolated orbits.
\end{proof}

\section{Quantitative defects of regular collision phases}
\label{sec:quantitative-phase-defects}
Absence of exact phases does not itself control approximate vectors.
We give a quantitative version on the space of bounded functions of
bounded variation. The phase is reconstructed from a complex vector;
no pointwise lower bound for its modulus is assumed. The logarithmic
bounds in this section are at a fixed radius and on a fixed compact
frequency set. A separate compactness argument gives a parameter-uniform
bound when the regularity budget is fixed.

For $\omega\in\Omega$ put
\begin{equation}\label{eq:collision-phase-defects}
 \begin{split}
 d_{R,\omega}(q)&=\|q\circ T_R-e^{i\phi_{R,\omega}}q\|_{L^1(\nu)},
                      \qquad (|q|=1),\\
 \mathcal E_{R,\omega}(f)&=\|f\circ T_R-e^{i\phi_{R,\omega}}f\|_{L^2(\nu)}.
 \end{split}
\end{equation}
The variation is always in the original collision coordinates. In
particular it is neither an anisotropic distribution norm nor a norm
of a many-return pushforward density.

\subsection{Finite-time measurable holonomy with a defect}
Use the product coordinates and the pair probabilities
\eqref{eq:explicit-conditional-pair-densities}. For a stable or unstable
oriented arc from $x$ to $y$ define
\[
 H_{s/u}(x,y)=
 \left|\frac{q(y)}{q(x)}-
            \exp\!\left(it\int_{\gamma_{s/u}(x,y)}\vartheta_R\right)\right|.
\]
These are bounded measurable functions of pairs. Null sets are taken
with respect to the indicated pair probabilities.

\begin{lemma}[A finite-time holonomy defect]
\label{lem:finite-holonomy-defect}
Fix a regular product set $\mathcal P$ at radius $R$. There are
$C_{\mathcal P}<\infty$ and $0<\theta_{\mathcal P}<1$ such that,
for $m\ge1$, $0<\epsilon<1/2$, and $q:M\to\mathbb S^1$ with
$\|q\|_{\mathrm{BV}}\le H$, $H\ge2$,
\begin{equation}\label{eq:finite-holonomy-defect}
 \int H_{s/u}\dd\mathfrak m_{s/u}
 \le C_{\mathcal P}\left[
 m d_{R,\omega}(q)+H\epsilon+
                  (\epsilon^{-1}+|t|)\theta_{\mathcal P}^{m}\right].
\end{equation}
The same expression controls both signs of time. The constants do not
depend on $u,s,t,H,m,\epsilon$.
\end{lemma}
\begin{proof}
Iterating the approximate phase equation along one orbit introduces
an error at most the sum of its $m$ one-step absolute defects. The
stable-pair version therefore has errors at its two endpoints. Both
pair marginals equal $m_{\mathcal P}=\nu(\mathcal P)^{-1}\nu|_{\mathcal P}$.
Invariance of $\nu$ bounds the integrated sum of endpoint defects by
$2m\nu(\mathcal P)^{-1}d_{R,\omega}(q)$.

Let $q_\epsilon=\mathsf S_\epsilon q$. Positivity of smoothing gives
$|q_\epsilon|\le1$, and
$\|q-q_\epsilon\|_1\le C H\epsilon$,
$\|q_\epsilon\|_{C^1(\alpha,p)}\le C\epsilon^{-1}$.
The two endpoint approximation errors are bounded at every iterate
by $2\nu(\mathcal P)^{-1}\|q-q_\epsilon\|_1$.
The actual stable arc lengths are at most
$C_{\mathcal P}\theta_{\mathcal P}^m$ after $m$ forward iterates,
by the contraction estimate in the proof of
Lemma~\ref{lem:regular-collision-product}. Thus the smoothed endpoint
difference costs $C_{\mathcal P}\epsilon^{-1}\theta_{\mathcal P}^m$.
The terminal action integral in \eqref{eq:action-telescope-arc} costs
at most $C_{\mathcal P}|t|\theta_{\mathcal P}^m$.
The lattice and constant phases cancel along the stable pair exactly.
These estimates prove \eqref{eq:finite-holonomy-defect} for stable arcs.

For an unstable pair telescope backwards. Its defect is the sum of the
one-step defects at $T_R^{-j}x,T_R^{-j}y$, $1\le j\le m$; the same
marginal domination applies by invertibility and invariance. Use
\eqref{eq:backward-action-arc} and backward contraction. The holonomy
sign is again positive, as in
\eqref{eq:measurable-unstable-holonomy}. This proves the other case.
No invariant pair measure and no independent orbit assumption is used.
\end{proof}

We make explicit how a measure-class assertion can be used
quantitatively without replacing it by an unproved bounded-density
assertion. Write $\mu=\mu^u\otimes\mu^s$ and
$\dd m_{\mathcal P}=\rho\dd\mu$ in the product coordinates.
Let $\lambda_s=\mu^u\otimes\mu^s\otimes\mu^s$ and
$\lambda_u=\mu^u\otimes\mu^u\otimes\mu^s$ be the reference triple
probabilities, and set
\[
 w_s=\frac{\dd\lambda_s}{\dd\mathfrak m_s},\qquad
 w_u=\frac{\dd\lambda_u}{\dd\mathfrak m_u}.
\]
These derivatives are positive and finite almost everywhere. For
$0\le F\le2$ and either pair type,
\begin{equation}\label{eq:pair-density-truncation}
 \int F\dd\lambda_{s/u}
 \le K\int F\dd\mathfrak m_{s/u}
            +2\lambda_{s/u}\{w_{s/u}>K\}.
\end{equation}
The last probability tends to zero as $K\to\infty$. Its value is part
of the constant below; no uniform pointwise bound on $\rho$ is needed.

\begin{lemma}[An almost constant phase on the product set]
\label{lem:almost-constant-product-phase}
For every $\eta>0$ there is a finite $C_{\mathcal P,\eta}$ such that
one can choose $c\in\mathbb S^1$ with
\begin{align}\label{eq:almost-constant-product-phase}
 \int_{\mathcal P}|q-c|\dd\nu
 \le\eta+C_{\mathcal P,\eta}\bigl[
 |t|+m d_{R,\omega}(q)+H\epsilon
              +(\epsilon^{-1}+|t|)\theta_{\mathcal P}^m\bigr].
\end{align}
This holds for the same $q,H,m,\epsilon$ as in the preceding lemma.
\end{lemma}
\begin{proof}
The action integrals on the original arcs are uniformly bounded on
$\mathcal P$. Hence the stable and unstable pair averages of
$|q(y)-q(x)|$ are bounded by the right side of
\eqref{eq:finite-holonomy-defect} plus $C_{\mathcal P}|t|$.

Here are the density choices needed to pass to two arbitrary product
points. Choose $J$ so that
$\int_{\{\rho>J\}}\rho\dd\mu$ is less than a prescribed small number.
Choose $K$ so that both tail probabilities in
\eqref{eq:pair-density-truncation} are less than that number divided
by $J$. Joining $(u,s)$ to $(u,s')$ and then to $(u',s')$, Fubini
bounds the $\mu\otimes\mu$ average of $|q(x)-q(y)|$ by the sum of
the two reference triple averages. Thus there is a value $c=q(y)$
for which the $\mu$ average of $|q-c|$ is at most that sum. Finally
\[
 \int|q-c|\dd m_{\mathcal P}
 \le J\int|q-c|\dd\mu+2\int_{\{\rho>J\}}\rho\dd\mu.
\]
Take the prescribed small numbers sufficiently small in terms of
$\eta$ and $\nu(\mathcal P)$, and apply
\eqref{eq:pair-density-truncation}. The remaining multiplier is a
finite function of $J,K,C_{\mathcal P},\nu(\mathcal P)$.
This proves \eqref{eq:almost-constant-product-phase}.
\end{proof}

\subsection{A quantitative roof obstruction}
\begin{proposition}[A logarithmic defect bound on a roof band]
\label{prop:logarithmic-roof-defect}
Fix $R$ and $0<t_0\le T<\infty$. There is $c_{R,t_0,T}>0$ such that
\begin{equation}\label{eq:logarithmic-roof-defect}
 d_{R,(u,s,t)}(q)\ge\frac{c_{R,t_0,T}}{1+\log H}
\end{equation}
whenever $t_0\le|t|\le T$, $|q|=1$, $H\ge2$, and
$\|q\|_{\mathrm{BV}}\le H$. The estimate is independent of $u,s$.
\end{proposition}
\begin{proof}
Restrict a regular product set so that every oriented four-corner
quadrilateral has action $\mathcal A=\oint\vartheta_R$ satisfying
$|\mathcal A|\le\pi/T$. By Stokes' formula and nonatomicity,
$|\mathcal A|>0$ for almost every quadruple under
$\lambda_4=(\mu^u)^2\otimes(\mu^s)^2$.
Choose $a_0>0$ such that
\[
 p_0:=\lambda_4\{|\mathcal A|\ge a_0\}>0.
\]
For this set and the entire roof band,
$|e^{it\mathcal A}-1|\ge2t_0a_0/\pi$.

Each of the four edge marginals of $\lambda_4$ is one of the reference
triple measures. Choose $K$ so that the sum of their probabilities
of $w_{s/u}>K$ is less than $p_0/2$. Intersect the set
$\{|\mathcal A|\ge a_0\}$ with all four complementary conditions;
call the resulting set $\mathcal Q$. It has measure at least $p_0/2$,
and each of its unnormalized edge marginals is bounded by
$K\mathfrak m_s$ or $K\mathfrak m_u$.

Multiply the four oriented approximate holonomies. The product of
$q$-ratios is exactly one, and the error in a product of unit complex
numbers is at most the sum of its four errors. Integration over
$\mathcal Q$ and Lemma~\ref{lem:finite-holonomy-defect} give
\begin{equation}\label{eq:quantitative-area-defect}
 \frac{p_0t_0a_0}{\pi}
 \le4K C_{\mathcal P}\bigl[
 m d_{R,\omega}(q)+H\epsilon
                  +(\epsilon^{-1}+T)\theta_{\mathcal P}^m\bigr].
\end{equation}
All constants have now been selected from the product set, its density
tails and its action distribution, independently of $q$ and $H$.
Choose $\epsilon=c_1/H$, with a fixed sufficiently small $c_1>0$.
Then choose $m=\lceil c_2+c_3\log H\rceil$ so that the final term in
brackets is also less than one quarter of the required lower bound.
The positive contraction exponent makes this possible. Equation
\eqref{eq:quantitative-area-defect} leaves $m d_{R,\omega}(q)\ge c_4>0$.
This proves \eqref{eq:logarithmic-roof-defect}.
\end{proof}

\subsection{All compact nonzero physical frequencies}
\begin{theorem}[A logarithmic phase defect on a compact band]
\label{thm:compact-log-phase-defect}
For a fixed $R\in I$ and a compact set
$\mathcal K\subset\Omega\setminus\{0\}$ there is $c_{R,\mathcal K}>0$
such that, for every $\omega\in\mathcal K$,
\begin{equation}\label{eq:compact-log-phase-defect}
 |q|=1,\quad \|q\|_{\mathrm{BV}}\le H,\quad H\ge2
 \quad\Longrightarrow\quad
 d_{R,\omega}(q)\ge\frac{c_{R,\mathcal K}}{1+\log H}.
\end{equation}
The set includes compact nonzero spatial frequencies with $s=t=0$.
\end{theorem}
\begin{proof}
Set
\[
 \rho(u,s)=|e^{iu_1}-1|+|e^{iu_2}-1|+|e^{is}-1|,
 \qquad \rho_*:=\min_{\mathcal K}(|t|+\rho(u,s))>0.
\]
Let $T=1+\max_{\mathcal K}|t|$. Fix a product set $\mathcal P$ and
the finite witnesses $(n_i,k_i),E_i,b_{ji}$ in
Lemma~\ref{lem:finite-return-witnesses}. Write $p_i=\nu(E_i)>0$ and
$d=d_{R,\omega}(q)$. For a constant $c\in\mathbb S^1$, iteration
of the approximate phase equation on $E_i$ gives
\begin{align}\label{eq:witness-phase-defect}
 p_i|e^{i(u\cdot k_i+s n_i)}-1|
 \le 2\int_{\mathcal P}|q-c|\dd\nu+n_i d
                     +p_i|t|n_i\|\tau_R\|_\infty.
\end{align}
Indeed $T_R^{n_i}E_i\subset\mathcal P$, its contribution to the
endpoint error is bounded by the integral over $\mathcal P$ by
invariance, and the $n_i$ one-step defects integrate to at most $n_i d$.
The roof sum is bounded by $n_i\|\tau_R\|_\infty$.

The integer identities $e_j=\sum_i b_{ji}g_i$ and
$|z^b-1|\le |b||z-1|$ for $|z|=1$, $b\in\Z$, show that finite
constants $A,B,C$ satisfy
\begin{equation}\label{eq:witness-resonance-control}
 \rho(u,s)\le A\int_{\mathcal P}|q-c|\dd\nu+B d+C|t|.
\end{equation}
For example one may take
$A=2\sum_{j,i}|b_{ji}|/p_i$,
$B=\sum_{j,i}|b_{ji}|n_i/p_i$ and
$C=\|\tau_R\|_\infty\sum_{j,i}|b_{ji}|n_i$.
These constants include the actual positive masses of the witnesses.

Choose $\eta>0$ with $A\eta<\rho_*/16$, and use
Lemma~\ref{lem:almost-constant-product-phase} with this $\eta$.
Choose $t_0\in(0,\rho_*/4)$ so small that all terms proportional
to $|t|$ in \eqref{eq:witness-resonance-control}, after substitution
of that lemma, are less than $\rho_*/16$ when $|t|\le t_0$.
Choose $\epsilon=c_5/H$, then $m=\lceil c_6+c_7\log H\rceil$,
so that the two remaining approximation terms contribute less than
$\rho_*/8$. For $|t|\le t_0$, the left side is at least
$3\rho_*/4$. The inequalities thus force $m d\ge c_8>0$.
For $|t|\ge t_0$ use Proposition~\ref{prop:logarithmic-roof-defect}
on $[t_0,T]$. Taking the smaller of the two constants proves the
assertion. All selections precede the choice of $q,H$.
\end{proof}

\subsection{Reconstructing a phase without a lower-modulus assumption}
\begin{lemma}[A bounded-variation circle truncation]
\label{lem:bv-circle-truncation}
For every complex $f\in\mathrm{BV}(M)$ there is a measurable
$q:M\to\mathbb S^1$ such that
\begin{equation}\label{eq:bv-circle-truncation}
 \|q\|_{\mathrm{BV}}\le C(1+\|f\|_{\mathrm{BV}}),\qquad
 \|q-f\|_2\le3\||f|-1\|_2
\end{equation}
whenever the latter norm is finite.
\end{lemma}
\begin{proof}
Put $a=|f|$. The BV chain inequality gives $\TV(a)\le C\TV(f)$.
The coarea formula gives a level $\tau\in[1/4,1/2]$ for which
$E=\{a>\tau\}$ has perimeter at most $4\TV(a)$.
These chain, product and coarea statements are for first distributional
derivatives in the collision coordinates; see \cite[Chapter 3]{AFP}.
Let $N_\tau(z)=z/\max\{|z|,\tau\}$ and define
$q=N_\tau(f)\one_E+\one_{M\setminus E}$.
The map $N_\tau$ is uniformly Lipschitz for the selected levels and
has modulus at most one. The BV product inequality, including the
jump on the finite-perimeter boundary of $E$, proves the first bound.
On $E$, $|q-f|=|1-a|$. On its complement $a\le1/2$ and
$|q-f|\le1+a\le3(1-a)$. This proves the second bound.
In particular no value of $f/|f|$ is used at a zero of $f$.
\end{proof}

\begin{theorem}[A logarithmic obstruction for complex approximate vectors]
\label{thm:compact-vector-defect}
Fix $R$ and a compact $\mathcal K\subset\Omega\setminus\{0\}$.
There is $c'_{R,\mathcal K}>0$ such that
\begin{equation}\label{eq:compact-vector-defect}
 \|f\|_2=1,\quad \|f\|_\infty+\|f\|_{\mathrm{BV}}\le H,\quad H\ge2
 \quad\Longrightarrow\quad
 \inf_{\omega\in\mathcal K}\mathcal E_{R,\omega}(f)
             \ge\frac{c'_{R,\mathcal K}}{(1+\log H)^2}.
\end{equation}
The bound applies to complex $f$, with no restriction on its zero set.
\end{theorem}
\begin{proof}
Write $\varepsilon=\mathcal E_{R,\omega}(f)$, $a=|f|$,
$\bar a=\int a\dd\nu$ and $x=\|a-\bar a\|_2$.
The reverse triangle inequality gives
$\|a\circ T_R-a\|_2\le\varepsilon$, hence
$\|a\circ T_R^m-a\|_2\le m\varepsilon$.
The BV correlation bound of Proposition~\ref{prop:bv-correlation},
rescaled by homogeneity, gives constants $C_0,\gamma>0$ with
\begin{equation}\label{eq:approximate-modulus-variance}
 x^2\le C_0 H^2e^{-\gamma m}+x m\varepsilon.
\end{equation}
Indeed add and subtract the lag-$m$ covariance of the centered $a$ and
use Cauchy--Schwarz on the remaining term. The constants in this
collision estimate can be taken uniform in $R$.

Given $0<\eta<1$, choose
$m\le C[1+\log H+|\log\eta|]$ with
$C_0H^2e^{-\gamma m}\le\eta^2$.
Equation~\eqref{eq:approximate-modulus-variance} implies
$x\le\eta+m\varepsilon$. Since $\|a\|_2=1$ and $\bar a\ge0$,
\[
 \|a-1\|_2^2=2(1-\bar a)
       \le2(1-\bar a^2)=2x^2.
\]
Lemma~\ref{lem:bv-circle-truncation} constructs $q$ with
$\|q\|_{\mathrm{BV}}\le C(1+H)$ and
$\|q-f\|_2\le3\sqrt2(\eta+m\varepsilon)$. Invariance and the
unit modulus of the multiplier yield
\begin{equation}\label{eq:phase-reconstruction-defect}
 d_{R,\omega}(q)
 \le\mathcal E_{R,\omega}(q)
 \le\varepsilon+6\sqrt2(\eta+m\varepsilon).
\end{equation}
This is an explicit reconstruction and defect comparison.

Let $L=1+\log H$. Theorem~\ref{thm:compact-log-phase-defect}, with
the bound $C(1+H)$ for $q$, makes the left side at least $c_9/L$.
Choose $\eta=c_{10}/L$ so that $6\sqrt2\eta<c_9/(2L)$.
The selected $m$ is at most $C_{R,\mathcal K}L$, since
$\log L\le L$. Equation~\eqref{eq:phase-reconstruction-defect}
then gives $\varepsilon\ge c'_{R,\mathcal K}/L^2$.
No lower-modulus estimate for $f$ was assumed; it was replaced by
\eqref{eq:approximate-modulus-variance} and the controlled truncation.
\end{proof}

\begin{corollary}[Polynomial regularity budgets on a fixed band]
\label{cor:polynomial-bv-vector-budget}
For fixed $R,\mathcal K,H_0,\zeta$ with $H_0\ge2$, $\zeta\ge0$,
normalized vectors with
$\|f_N\|_\infty+\|f_N\|_{\mathrm{BV}}\le H_0N^\zeta$
satisfy $\inf_{\omega\in\mathcal K}\mathcal E_{R,\omega}(f_N)
\ge c/(1+\log(2+N))^2$. In particular no such vectors have defect
$O(N^{-a})$ for any fixed $a>0$.
\end{corollary}
\begin{proof}
Substitute the regularity budget into
\eqref{eq:compact-vector-defect}. A positive power of $N$ dominates
$(1+\log(2+N))^2$.
\end{proof}

\subsection{Uniform compact separation across the physical family}
\begin{theorem}[Parameter-uniform separation at bounded regularity]
\label{thm:uniform-compact-bv-separation}
For every $H<\infty$ and compact
$\mathcal K\subset\Omega\setminus\{0\}$ there is
$c_{H,\mathcal K}>0$ such that
\begin{equation}\label{eq:uniform-compact-bv-separation}
 \inf_{\substack{R\in I,\ \omega\in\mathcal K,\ \|f\|_2=1\\
                \|f\|_\infty+\|f\|_{\mathrm{BV}}\le H}}
          \mathcal E_{R,\omega}(f)\ge c_{H,\mathcal K}.
\end{equation}
An empty admissible class causes no difficulty. The theorem asserts
no uniform logarithmic dependence of this constant on $H$.
\end{theorem}
\begin{proof}
Suppose a sequence made the infimum zero. Pass to subsequences with
$R_j\to R$, $\omega_j\to\omega\in\mathcal K$ and $f_j\to f$ in
$L^1(\nu)$. The last subsequence exists by BV compactness. In the present
coordinates it also follows directly from
Lemma~\ref{lem:bv-smoothing}: the approximation is uniform in $L^1$,
and at every fixed smoothing scale the bounded smoothed functions
form a precompact family on the reflected compact cylinder.
The supremum bound upgrades this convergence to $L^2$, preserves
$\|f\|_2=1$, and gives $|f|\le H$ almost everywhere.

Outside the finite one-collision singularity set at $R$,
$T_{R_j}x\to T_Rx$, $\kappa_{R_j}(x)\to\kappa_R(x)$ and
$\tau_{R_j}(x)\to\tau_R(x)$, by persistence of the regular physical
root. Thus $e^{i\phi_{R_j,\omega_j}}\to e^{i\phi_{R,\omega}}$
almost everywhere and in every finite $L^p$ norm. For a smooth $g$,
invariance of the common probability gives
\begin{align*}
 \|f_j\circ T_{R_j}-f\circ T_R\|_2
 \le\|f_j-f\|_2+2\|f-g\|_2
                      +\|g\circ T_{R_j}-g\circ T_R\|_2.
\end{align*}
Take $j\to\infty$, then $g\to f$ in $L^2$. Dominated convergence
controls the last term. The multiplier term also converges in $L^2$.
The limiting defect is therefore zero.

The identity for $f$ makes $|f|$ invariant. Ergodicity and $\|f\|_2=1$
give $|f|=1$ almost everywhere. Theorem~\ref{thm:complete-physical-phase}
now forces $\omega=0$, contradicting $\mathcal K\subset\Omega\setminus\{0\}$.
This proves the positive infimum, without claiming continuity of
product sets or their density-tail constants in $R$.
\end{proof}

\paragraph{Relation to the complementary integral.}
Theorems~\ref{thm:compact-vector-defect} and
\ref{thm:uniform-compact-bv-separation} are estimates on actual
collision functions, with their stated norms. They cannot be applied
to an anisotropic distribution until a map to such functions, its
normalization, regularity and defect comparison have been proved.
They also do not assert that these BV classes are invariant under
the twisted collision operator. Absence of exact phases alone has not
been substituted for either assertion.

The logarithmic estimate allows a growing regularity budget on a fixed
compact nonzero band at each fixed $R$. The parameter-uniform theorem
instead fixes that budget. No estimate of $c_{R,\mathcal K}$ is supplied
when the roof band grows or $\mathcal K$ approaches the origin.
Consequently the physical annulus beginning at
$2n^{-99/200}$, equivalently the rescaled boundary $2n^{1/200}$,
is not declared controlled by these theorems. The full complementary
integral, the complete critical and singular raw residual sum and the
weighted exact-event comparison retain their distinct roles in
Theorem~\ref{thm:LLT}. The present estimates give a quantitative
phase and function-vector step toward that same inversion problem.

\section{Uniform defect bounds at frequencies approaching zero}
\label{sec:near-origin-defects}
The compact bands of Section~\ref{sec:quantitative-phase-defects}
exclude the origin by a fixed amount. We now allow the physical
frequency to approach zero and the variation budget to grow. The
observable is the bounded compensation $h_R$, including its original
section discontinuity. Uniform ellipticity, rather than compact-band
phase rigidity, supplies the obstruction.

For $z\in\R^4$, $\sigma\in\R$, and a collision function $f$, put
\begin{equation}\label{eq:compensated-defect-definition}
 \mathcal D_{R,z,\sigma}f
       =f\circ T_R-e^{i(\sigma+z\cdot h_R)}f.
\end{equation}
The constant phase $\sigma$ is arbitrary. All norms in this section
are with respect to $\nu$ on the collision cylinder. Constants are
uniform in $R$; none depends on $\sigma$.

\subsection{A twisted middle block with untwisted ends}
\begin{lemma}[Damping on the diffusive collision scale]
\label{lem:diffusive-middle-damping}
There exist $r_0>0$, $A>1$, and $C<\infty$ with the following
properties. If $0<r=|z|\le r_0$, set
\[
 \delta=\tfrac14r^{1/12},\qquad m=\lceil A/r^2\rceil,\qquad
 P_{R,z}=\mathcal L_R M_{\exp(-iz\cdot h_{R,\delta})}.
\]
On every relevant local collision space,
\begin{equation}\label{eq:diffusive-middle-damping}
 \sup_{j\ge0}\|P_{R,z}^{j}\|\le C,
       \qquad |\ell(P_{R,z}^{m}\nu)|\le\tfrac18.
\end{equation}
Moreover
\begin{equation}\label{eq:middle-unsmoothing-error}
 r\|S_{m,R}(h_R-h_{R,\delta})\|_2
       \le C r^{1/24}\sqrt{1+|\log r|}.
\end{equation}
\end{lemma}
\begin{proof}
By Theorem~\ref{thm:joint-uniform-nondegeneracy},
$\Gamma_R=c_*D_R\ge c_*d_0I_4$. The covariance smoothing estimate
makes $\Gamma_{R,\delta}\ge(c_*d_0/2)I_4$ for small $r$.
The spectral expansion of Lemma~\ref{lem:smooth-collision-expansion}
applies since $r\delta^{-2}=O(r^{5/6})$. Its cubic error is
$O(r^3\delta^{-6})=O(r^{5/2})$. Reducing $r_0$ gives
\[
 \mathop{\rm Re}\log\lambda_{R,\delta}(-z)\le-c_1r^2
\]
for some uniform $c_1>0$. The projections and complementary powers
are uniformly bounded, proving the power bound. The stationary
principal amplitude is $1+O(r\delta^{-2})$. Choose $A$ large enough
that $2e^{-c_1A}<1/16$, and then reduce $r_0$ so that the
complementary term at $m$ is less than $1/16$. This proves the
second bound. Finally Proposition~\ref{prop:bv-correlation} gives
$\|S_m(h-h_\delta)\|_2^2\le Cm\delta(1+|\log\delta|)$.
Since $mr^2\le A+1$, it proves
\eqref{eq:middle-unsmoothing-error}. No multiplier bound for the
unsmoothed $h_R$ is asserted.
\end{proof}

\begin{theorem}[Quadratic coercivity for collision phases]
\label{thm:near-origin-circle-defect}
There exist $a_0,c_0,r_0>0$ such that, if $H\ge2$,
$0<r=|z|\le r_0$, and $r(1+\log H)\le a_0$, then
\begin{equation}\label{eq:near-origin-circle-defect}
 |q|=1,\quad \|q\|_{\mathrm{BV}}\le H
 \quad\Longrightarrow\quad
 \inf_{R\in I,\ \sigma\in\R}
        \|\mathcal D_{R,z,\sigma}q\|_1\ge c_0r^2.
\end{equation}
Here the infimum means the same uniform bound for every admissible
$q$ at each radius; $q$ itself may depend on $R,z,\sigma$.
\end{theorem}
\begin{proof}
Fix the radius and write $d=\|\mathcal D_{R,z,\sigma}q\|_1$.
Let $q_\epsilon=\mathsf S_\epsilon q$, where
$\epsilon=b_0/H$ and $b_0>0$ is a sufficiently small fixed number.
Then $|q_\epsilon|\le1$, $\|q-q_\epsilon\|_1\le Cb_0$, and
$\|q_\epsilon\|_{C^2}\le C\epsilon^{-2}$.
Choose integers
\[
 l=\lceil b_1+b_2\log H\rceil,\qquad N=m+2l,
\]
where $m$ is as in Lemma~\ref{lem:diffusive-middle-damping}.
The constants $b_1,b_2$ will be fixed before $H,r,q$ are chosen.

Iteration of the defect and invariance give
\begin{equation}\label{eq:near-origin-lower-pairing}
 \left|\int\overline q\,(q\circ T_R^N)
              e^{-iz\cdot S_{N,R}h_R}\dd\nu\right|\ge1-Nd.
\end{equation}
Indeed the corresponding unit-modulus integrand differs in $L^1$
from $e^{iN\sigma}$ by at most $Nd$.
Remove the phase from the first and last $l$ collisions. Since $h_R$
is bounded, this changes the pairing by at most $2lr\|h_R\|_\infty$.
Replacing its two endpoint factors by $q_\epsilon$ costs at most
$2\|q-q_\epsilon\|_1$. Smoothing only the middle collision observable
costs at most \eqref{eq:middle-unsmoothing-error}, by stationarity
and the unit supremum bounds on the endpoint factors.

The resulting pairing is exactly
\begin{equation}\label{eq:untwisted-end-word}
 \ell\left(M_{q_\epsilon}\mathcal L_R^l
                 P_{R,z}^m\mathcal L_R^l(\overline{q_\epsilon}\nu)\right).
\end{equation}
Replace both untwisted end powers by $\Pi_R=\nu\ell$.
The total error is bounded by
$C\epsilon^{-4}\vartheta^l$: both smooth endpoint factors cost
$C\epsilon^{-2}$, all other untwisted powers are bounded, and the
middle power is bounded by \eqref{eq:diffusive-middle-damping}.
The remaining expression is
\[
 \left|\int q_\epsilon\dd\nu\right|^2
                          \ell(P_{R,z}^m\nu),
\]
whose modulus is at most $1/8$.

Choose $b_0$ to make the endpoint smoothing error less than $1/16$;
then choose $b_1,b_2$ to make $C\epsilon^{-4}\vartheta^l<1/16$
for every $H\ge2$. Reducing $r_0$ makes the middle unsmoothing error
less than $1/16$. Finally choose $a_0$ so that
$r(1+\log H)\le a_0$ makes $2lr\|h_R\|_\infty<1/16$.
Thus the modulus in \eqref{eq:near-origin-lower-pairing} is at most
$1/2$. Since $N\le C/r^2$ under the same constraint, $Nd\ge1/2$
proves the assertion. Notice that the end blocks are untwisted only
inside this comparison; the original record is unchanged.
\end{proof}

\subsection{Complex vectors and their modulus}
\begin{theorem}[A uniform small-frequency function bound]
\label{thm:near-origin-vector-defect}
There exist $a_1,c_1,r_1>0$ such that, with
\[
 \Lambda(H,r)=1+\log(H/r),\qquad H\ge2,\quad 0<r\le r_1,
\]
$r\Lambda(H,r)\le a_1$ implies
\begin{equation}\label{eq:near-origin-vector-defect}
 \|f\|_2=1,\quad\|f\|_\infty+\|f\|_{\mathrm{BV}}\le H
 \quad\Longrightarrow\quad
 \|\mathcal D_{R,z,\sigma}f\|_2
                    \ge\frac{c_1r^2}{\Lambda(H,r)}
       \quad (|z|=r).
\end{equation}
The assertion is uniform in $R$ and $\sigma$, and allows zeros of $f$.
\end{theorem}
\begin{proof}
Put $e=\|\mathcal D_{R,z,\sigma}f\|_2$, $a=|f|$ and
$x=\|a-\int a\dd\nu\|_2$. The modulus estimate
\eqref{eq:approximate-modulus-variance} uses only the unit modulus of
the twist and the collision BV correlation estimate. It therefore
applies here with constants uniform in $R$:
\[
 x^2\le C H^2e^{-\gamma j}+xje.
\]
Given $\eta=b r^2$ with fixed small $b>0$, choose
$j\le C\Lambda(H,r)$ so that the first term is at most $\eta^2$.
Then $x\le\eta+je$. The circle truncation of
Lemma~\ref{lem:bv-circle-truncation} gives a phase $q$ with
\[
 \|q\|_{\mathrm{BV}}\le C(1+H),\qquad
 \|q-f\|_2\le3\sqrt2(\eta+je).
\]
Consequently
\[
 \|\mathcal D_{R,z,\sigma}q\|_1
       \le e+6\sqrt2(\eta+je).
\]
After decreasing $a_1,r_1$, the hypotheses of
Theorem~\ref{thm:near-origin-circle-defect} hold for the budget
$C(1+H)$. Choose $b$ so that the $\eta$ term is less than
$c_0r^2/2$. The remaining inequality gives
$e\ge c r^2/(1+j)\ge c_1r^2/\Lambda(H,r)$.
This proves the claim without dividing by $f$ on its zero set.
\end{proof}

\paragraph{The frequency being estimated.}
The twist in \eqref{eq:compensated-defect-definition} contains
$z\cdot h_R$, not just the physical one-collision displacement and
flight time. Its sum over a genuine return is exactly
$z\cdot(G_R-\bar G_R)$. This exact relation is used next to obtain
bounds for the actual induced record. The preceding theorems bound
defects of functions; neither is an integrated Fourier-tail estimate.

\section{Regularizing the actual tower and excluding return defects}
\label{sec:actual-return-defects}
A function of an actual return state lifts to the collision tower with
an exactly localized defect. That lift need not have bounded variation.
We construct a regular approximation with a controlled cost before
using Section~\ref{sec:near-origin-defects}. The construction concerns
functions and first distributional derivatives, not anisotropic
vectors or inverse-coarea densities.

\subsection{Finite collision records in the initial coordinates}
\begin{lemma}[A quantitative finite-record variation bound]
\label{lem:finite-record-variation}
There is $C\ge1$ such that, for every $L\ge1$ and $R\in I$, the
following bounded functions have variation norm at most
\begin{equation}\label{eq:finite-record-variation-budget}
 B_L=\exp\{C L\log(2+L)\}:
\end{equation}
all components of
$\chi\circ T_R^j$, $h_R\circ T_R^j$, and $\eta_R\circ T_R^j$,
$|j|\le L$, where $\chi(\alpha,p)=(\cos\alpha,\sin\alpha,p)$.
Their supremum norms are bounded independently of $L$. Null singular
sets are assigned any fixed bounded semialgebraic convention.
Moreover, if $g$ is smooth on the reflected collision cylinder, then
\begin{equation}\label{eq:finite-record-smooth-composition}
 \|g\circ T_R^j\|_{\mathrm{BV}}
              \le C\|g\|_{C^1(\alpha,p)}B_L\qquad(|j|\le L).
\end{equation}
The constant in $B_L$ is independent of the coefficients specifying
the moving section rectangles.
\end{lemma}
\begin{proof}
We supply the finite-complexity argument, since a one-collision BV
bound does not alone imply this lemma. A fixed finite set of disk
labels contains every next collision, by the uniform horizon. Use
normal vectors $n_j\in\mathbb S^1$, outgoing velocities
$v_j\in\mathbb S^1$, and flight lengths $\tau_j$ as auxiliary
variables. For a fixed label word $d_j$, each step is given by
\begin{equation}\label{eq:polynomial-flight-graph}
 Rn_j+\tau_jv_j-d_j=Rn_{j+1},\qquad
 v_{j+1}=v_j-2(v_j\cdot n_{j+1})n_{j+1},\qquad
 p_j=v_j\cdot Jn_j,
\end{equation}
where $J$ rotates vectors by $\pi/2$. Require $\tau_j>0$,
$v_j\cdot n_j>0$, and $v_j\cdot n_{j+1}<0$.
Require also that the open flight not enter any other candidate disk.
These last conditions are finite Boolean combinations of polynomial
sign conditions of bounded degree: for each center $d$, minimize
$|Rn_j+t v_j-d|^2$ over $0\le t\le\tau_j$.
The minimizing value occurs at an endpoint or at
$t=(d-Rn_j)\cdot v_j$; splitting into these three cases gives the
required inequalities without a quantified time variable. Endpoints
and tangencies are included in the declared singular convention.
Thus the graph records the actual first collision, not merely some
root of the flight equation.

On each bounded rational angular chart, the initial normal and outgoing
velocity also have such a graph, using a positive square-root variable
for $\sqrt{1-p^2}$. A word of length at most $2L+2$ therefore uses
at most $C_2L$ variables and $C_2L$ polynomials, of degree bounded by
a fixed $d$. All these variables can be confined to a bounded ball.
The same representation gives negative iterates by fixing the terminal
rather than the initial state; the regular billiard map is invertible.
Every lifted fiber over a regular initial state is unique.

Section membership is equally finite. Each angular interval has length
less than $\pi$ and is described by two oriented half-plane tests on
$n_j$, together with the two inequalities for $p_j$. The moving angular
centers enter only as real coefficients. Taking the finite union of
rectangles gives $\eta_R$; adjoining its two possible values also gives
the graph of $h_R=f_R-\bar G_R\eta_R$. All degree and variable counts
remain as above. There are at most $C_3^L$ choices of label word and
of these finite membership alternatives.

Here is the topological bound we use. The finite sign-condition bound
recorded in \cite[Section 3.2, equation (3.3)]{BPR}, with the zeroth
Betti number and ambient dimension $k$, bounds the total number of
connected components by
\[
 d(2d-1)^{k-1}\sum_{i=0}^{k}\binom{s}{i}4^i
                       \le (C(1+s)d)^{Ck}.
\]
The bound applies uniformly in the polynomial coefficients and to any
Boolean union of the sign conditions. We use this finite inequality,
not an asymptotic formula whose dimension is held fixed.
Fix the second collision coordinate and a real level $a$. Add the
condition that a chosen scalar output $u$ exceed $a$. The lifted
semialgebraic set has the preceding complexity. Its projection onto
the remaining coordinate has at most as many connected components:
every connected component projects to an interval or a point. Summing
over the at most $C_3^L$ alternatives bounds the number of components
of the slice superlevel set by $\exp\{C L\log(2+L)\}$.
The same reasoning applies with the two input coordinates interchanged.

If $|u|\le M_0$, the perimeter of a one-dimensional superlevel set
with at most $b$ components is at most $2b$ in the interior. The
one-dimensional coarea formula therefore bounds the variation of
that slice by $4M_0b$. Integrating both slice bounds gives the two
first distributional derivative bounds, including jumps. The fixed
angular chart cover has bounded derivatives and inverse derivatives;
its seams add only uniformly bounded boundary terms. This proves
\eqref{eq:finite-record-variation-budget}, after increasing $C$.
No inverse derivative near grazing has been bounded pointwise.

Finally extend $g$ from the embedded cylinder
$\chi(M)\subset\R^3$ to a neighborhood, with Lipschitz norm at most
$C\|g\|_{C^1(\alpha,p)}$, by a fixed tubular chart and cutoff. The
BV chain inequality applied to the three components of
$\chi\circ T_R^j$ proves \eqref{eq:finite-record-smooth-composition}.
This does not assert a BV composition bound for an arbitrary nonsmooth
$g$; only the smoothed functions below require this estimate.
\end{proof}

\subsection{Exact lifting and its regular approximation}
Write $Y=Y_R^*$, $F=F_R^*$, $r_R^*$ for its actual return time and
$g_R=G_R-\bar G_R$. For almost every collision state let its age since
the most recent section visit be $j$. The corresponding disjoint levels
are
\begin{equation}\label{eq:actual-age-levels}
 D_{j,R}=T_R^j\{y\in Y:r_R^*(y)>j\},\qquad j\ge0.
\end{equation}
They partition a full-measure set, by recurrence and invertibility.
For a measurable section function $f$, define its exact phase lift by
\begin{equation}\label{eq:exact-phase-tower-lift}
 (\mathsf L_{R,z}f)(T_R^jy)
       =e^{iz\cdot S_{j,R}h_R(y)}f(y),\qquad 0\le j<r_R^*(y).
\end{equation}
This is a function lift, not an induced transfer operator.

\begin{lemma}[The defect is supported at the true tower tops]
\label{lem:exact-tower-defect}
For $1\le p<\infty$, whenever the terms are finite,
\begin{align}
 \|\mathsf L_{R,z}f\|_{L^p(\nu)}^p
       &=c_*\int_Y r_R^*|f|^p\dd\nu_R^*,
                      \label{eq:tower-lift-norm}\\
 \|\mathcal D_{R,z,0}(\mathsf L_{R,z}f)\|_{L^p(\nu)}^p
       &=c_*\int_Y|f\circ F-e^{iz\cdot g_R}f|^p\dd\nu_R^*.
                      \label{eq:tower-exact-defect}
\end{align}
The analogous supremum bound is
$\|\mathsf L_{R,z}f\|_\infty\le\|f\|_\infty$.
\end{lemma}
\begin{proof}
Invariance and the disjoint levels give, for every nonnegative $b$,
\[
 \int_M b\dd\nu
   =c_*\int_Y\sum_{j=0}^{r_R^*-1}b(T_R^jy)\dd\nu_R^*(y).
\]
Apply this first to the modulus of the lift. At every nontop level
its next value is exactly $e^{iz\cdot h_R}$ times its current value.
At the last level the defect is
\[
 f(Fy)-e^{iz\cdot S_{r_R^*,R}h_R(y)}f(y)
               =f(Fy)-e^{iz\cdot g_R(y)}f(y),
\]
using the one-return compensation identity. This proves
\eqref{eq:tower-exact-defect}. Only one top occurs per return; its
measure is integrated against $c_*\nu_R^*$, not against a
return-time-biased probability.
\end{proof}

For a section function $f$ let $f^0$ be its zero extension to $M$,
$M_f=\|f\|_\infty$ and $V_f=\|f^0\|_{\mathrm{BV}}$.
Set $f_\epsilon=\mathsf S_\epsilon f^0$ and define
\begin{equation}\label{eq:finite-age-approximation}
 Q_{L,\epsilon}(x)=\sum_{j=0}^{L-1}\one_{D_{j,R}}(x)
 e^{iz\cdot\sum_{i=1}^j h_R(T_R^{-i}x)}
                              f_\epsilon(T_R^{-j}x).
\end{equation}
The sum for $j=0$ is empty. The approximant is zero on ages at least
$L$; the original lift is not truncated in any theorem statement.

\begin{proposition}[A uniform cost for approximating the lift]
\label{prop:tower-bv-approximation}
For $L\ge2$ and $0<\epsilon<1/2$, uniformly in $R,z,f$,
\begin{align}
 \|Q_{L,\epsilon}\|_\infty&\le M_f,
                                      \label{eq:tower-sup-bound}\\
 \|\mathsf L_{R,z}f-Q_{L,\epsilon}\|_1
        &\le C(L\epsilon V_f+M_fe^{-cL}),
                                      \label{eq:tower-L1-approximation}\\
 \|\mathsf L_{R,z}f-Q_{L,\epsilon}\|_2^2
        &\le C(M_fL\epsilon V_f+M_f^2e^{-cL}),
                                      \label{eq:tower-L2-approximation}\\
 \|Q_{L,\epsilon}\|_{\mathrm{BV}}
        &\le M_f\exp\{C L\log(2+L)\}(1+\epsilon^{-1}+|z|).
                                      \label{eq:tower-BV-budget}
\end{align}
\end{proposition}
\begin{proof}
The level indicators are disjoint, proving the supremum bound.
On the first $L$ levels, invariance bounds the total $L^1$ change
of the initial function by
$L\|f^0-f_\epsilon\|_1\le CL\epsilon V_f$.
The remaining age probability is exactly
\[
 \nu\{\text{age}\ge L\}
       =c_*\int_Y(r_R^*-L)_+\dd\nu_R^*\le Ce^{-cL},
\]
by Theorem~\ref{thm:exponential-return}. This proves the $L^1$
bound. Both functions have supremum at most $M_f$; multiplying the
$L^1$ bound by $2M_f$ gives the $L^2$ bound.

For variation, the exact level indicator is
\[
 \one_{D_{j,R}}=(\eta_R\circ T_R^{-j})
                    \prod_{i=0}^{j-1}(1-\eta_R\circ T_R^{-i}).
\]
The finite-record lemma and the BV product inequality bound its
variation by $C(j+1)B_L$. The same lemma bounds the phase sum
variation by $CjB_L$. The exponential is a Lipschitz function of this
real vector sum, with Lipschitz constant at most $|z|$.
Moreover $\|f_\epsilon\|_{C^1}\le C M_f\epsilon^{-1}$, so
\eqref{eq:finite-record-smooth-composition} applies to its pullback.
Apply the BV product inequality to each summand and then sum over
$j<L$. Absorb the resulting fixed powers of $L$ into $B_L$ by
increasing its constant. This proves \eqref{eq:tower-BV-budget}.
Neither the levels nor their images have been declared smooth; their
indicator jumps are included in the BV estimates.
\end{proof}

\subsection{Defects for the actual unbounded return record}
For $H\ge2$, $0<r<1/2$ define
\begin{equation}\label{eq:return-log-budget}
 \Lambda=1+\log(H/r),\qquad \mathcal K(H,r)=\Lambda\log(2+\Lambda).
\end{equation}
The letter $\mathcal K$ here denotes a scalar budget, not a compact
frequency set as in Section~\ref{sec:quantitative-phase-defects}.

\begin{theorem}[Uniform return-phase coercivity near the origin]
\label{thm:actual-return-circle-defect}
There exist $a_2,c_2,r_2>0$ such that, for $H\ge2$,
$0<r=|z|\le r_2$, and $r\mathcal K(H,r)\le a_2$,
\begin{equation}\label{eq:actual-return-circle-defect}
 |q|=1\text{ on }Y_R^*,\quad\|q^0\|_{\mathrm{BV}}\le H
 \quad\Longrightarrow\quad
 \|q\circ F_R^*-e^{iz\cdot(G_R-\bar G_R)}q\|_{L^1(\nu_R^*)}
                         \ge c_2r^2.
\end{equation}
All constants are independent of $R$, and no induced mixing hypothesis
is used.
\end{theorem}
\begin{proof}
The exact lift $Q=\mathsf L_{R,z}q$ has modulus one. Given
$\eta=b r^2$ with $b>0$ fixed and small, choose
$L\le C(1+|\log\eta|)$ and
$\epsilon=\eta/(C L H)$ so that
\eqref{eq:tower-L1-approximation} gives
$\|Q-Q_{L,\epsilon}\|_1\le\eta$.
The circle truncation in Lemma~\ref{lem:bv-circle-truncation}, applied
to $Q_{L,\epsilon}$, gives a phase $\widetilde q$ satisfying
\[
 \|\widetilde q-Q\|_1\le4\eta,\qquad
 1+\log(2+\|\widetilde q\|_{\mathrm{BV}})
                                     \le C\mathcal K(H,r).
\]
For the first bound, use the pointwise inequality from that lemma,
$|\widetilde q-Q_{L,\epsilon}|\le3\big||Q_{L,\epsilon}|-1\big|
\le3|Q_{L,\epsilon}-Q|$. The second follows from
\eqref{eq:tower-BV-budget}; the constants include the chosen $b$.

After decreasing $a_2,r_2$, Theorem~\ref{thm:near-origin-circle-defect}
applies to $\widetilde q$. Changing a function by $4\eta$ in $L^1$
changes its collision defect by at most $8\eta$.
Equation~\eqref{eq:tower-exact-defect} therefore gives
\[
 c_0r^2\le c_*\|q\circ F_R^*-e^{iz\cdot g_R}q\|_1+8br^2.
\]
Choose $b<c_0/16$. This proves the assertion with a uniform $c_2>0$.
In particular the regularity of the exact lift was not assumed.
\end{proof}

\begin{theorem}[Uniform coercivity for actual return functions]
\label{thm:actual-return-vector-defect}
There exist $a_3,c_3,r_3>0$ such that
$0<r=|z|\le r_3$, $H\ge2$, and $r\mathcal K(H,r)\le a_3$ imply
\begin{align}
 &\|f\|_{L^2(\nu_R^*)}=1,\quad M_f+V_f\le H
 \quad\Longrightarrow\notag\\
 &\qquad
 \|f\circ F_R^*-e^{iz\cdot(G_R-\bar G_R)}f\|_{L^2(\nu_R^*)}
                     \ge\frac{c_3r^2}{\mathcal K(H,r)}.
                         \label{eq:actual-return-vector-defect}
\end{align}
The section function may be complex and may vanish. The conclusion
is uniform in the radius and in the direction of $z$.
\end{theorem}
\begin{proof}
Let $Q=\mathsf L_{R,z}f$ and $S=\|Q\|_2$. By
\eqref{eq:tower-lift-norm}, $\sqrt{c_*}\le S\le H$; the upper bound
also follows from $\|Q\|_\infty\le H$ and $\nu(M)=1$.
If $e$ denotes the defect on the left of
\eqref{eq:actual-return-vector-defect}, then
\begin{equation}\label{eq:normalized-tower-defect}
 \|\mathcal D_{R,z,0}(Q/S)\|_2=\sqrt{c_*}\,e/S\le e.
\end{equation}

Choose an accuracy
$\eta=b r^2/\mathcal K(H,r)$, where $0<b<1$ will be fixed.
The bounds in Proposition~\ref{prop:tower-bv-approximation} allow
$L\le C[\Lambda+|\log b|]$ and
$\epsilon\ge c\eta^2/(H^2L)$ with
$\|Q-Q_{L,\epsilon}\|_2\le\sqrt{c_*}\eta$.
For example choose $L$ to make the exponential term at most
$c_*\eta^2/2$, then choose $\epsilon$ equal to a sufficiently small
fixed multiple of $\eta^2/(H^2L)$. Both choices are uniform in $R$.
Normalize $Q_{L,\epsilon}/S$ to an $L^2$ unit function $v$.
For $\eta<1/4$ this is well defined and
\[
 \|v-Q/S\|_2\le2\eta,\qquad
 1+\log(2+\|v\|_\infty+\|v\|_{\mathrm{BV}})+|\log r|
       \le C(1+|\log b|)^2\mathcal K(H,r).
\]
The last bound is \eqref{eq:tower-BV-budget}, using the lower bound
on $S$, and $\log\mathcal K(H,r)\le C\Lambda$.
The displayed dependence on $b$ prevents a circular choice of accuracy.

Choose $b$ so small that
$4b<c_1/[2C(1+|\log b|)^2]$; such a choice exists because
$b(1+|\log b|)^2\to0$. Fix this $b$ once and for all. Decrease
$a_3,r_3$ so that Theorem~\ref{thm:near-origin-vector-defect} applies
to $v$ under the preceding budget. That theorem and
\eqref{eq:normalized-tower-defect} give
\[
 e+4\eta\ge\|\mathcal D_{R,z,0}v\|_2
   \ge\frac{c_1r^2}{C(1+|\log b|)^2\mathcal K(H,r)}.
\]
The selected accuracy absorbs $4\eta$ and proves
\eqref{eq:actual-return-vector-defect}. No lower bound for $|f|$
and no mixing rate for $F_R^*$ enter the proof.
\end{proof}

\begin{corollary}[The intervening annulus with polynomial budgets]
\label{cor:actual-annulus-defect}
Fix $H_0\ge2$ and $\zeta\ge0$. For all sufficiently large $n$,
uniformly in $R$ and in physical frequencies
\begin{equation}\label{eq:actual-annulus-frequencies}
 2n^{-99/200}\le|z|\le n^{-2/5},\qquad
        2n^{1/200}\le|\sqrt n\,z|\le n^{1/10},
\end{equation}
unit section phases with zero-extension variation at most $H_0n^\zeta$
have defect at least $c|z|^2$. Normalized complex section functions
with $M_f+V_f\le H_0n^\zeta$ have defect at least
\begin{equation}\label{eq:actual-annulus-vector-rate}
 \frac{c|z|^2}{\log(2+n)\log\log(e^e+n)}.
\end{equation}
In particular this last lower bound is at least
$c n^{-99/100}/[\log(2+n)\log\log(e^e+n)]$ on the whole annulus.
\end{corollary}
\begin{proof}
On this annulus $\Lambda(H_0n^\zeta,|z|)\le C\log(2+n)$, so
$\mathcal K\le C\log(2+n)\log\log(e^e+n)$.
The largest possible value of $|z|\mathcal K$ is bounded by
$C n^{-2/5}\log(2+n)\log\log(e^e+n)$, which tends to zero.
Thus both preceding theorems apply. Squaring the lower frequency
endpoint gives the final power. The interval is nonempty for all
sufficiently large $n$, since its endpoint ratio is
$n^{19/200}/2\to\infty$.
\end{proof}

\paragraph{What the annulus estimate supplies.}
The norm in \eqref{eq:actual-annulus-vector-rate} is a one-step
cohomological defect for the \emph{actual unbounded return record},
with the normalization and first-variation budget stated above.
The lifted defect and the regularization error have both been proved,
rather than assumed to come from an induced spectral construction.
No conclusion is made here for arbitrary eigenvalues of the induced
operator: the multiplier in the theorem is precisely
$e^{iz\cdot(G_R-\bar G_R)}$.

This estimate is a quantitative input on the intervening annulus, not
the complementary Fourier integral itself. To deduce decay of powers
from it on a distribution space still requires a construction of that
space, a function reconstruction with the displayed budgets, and a
bound relating the operator defect to the present physical $L^2$
defect. Nor does the finite-record BV lemma estimate any second
derivative of a raw pushforward density. The growing roof region,
the complete critical/singular extraction, and the weighted exact-event
comparison remain distinct requirements of Theorem~\ref{thm:LLT}.
All conditioning indicators in the inherited central and moment
theorems are unchanged.

\section{Peripheral return phases and a quadratic regularity budget}
\label{sec:peripheral-return-phases}
The spectral phase of a return operator is not a constant phase per
collision. We keep these parameters separate. The physical Fourier
variable is $z\in\R^4$, the return spectral phase is
$\xi\in[-\pi,\pi]$, and the corresponding unit spectral parameter is
$\lambda=e^{i\xi}$. For a section function define
\begin{equation}\label{eq:peripheral-return-defect}
 \mathcal E_{R,z,\xi}f
   =f\circ F_R^*-e^{i(\xi+z\cdot g_R)}f,
       \qquad g_R=G_R-\bar G_R.
\end{equation}
Theorem~\ref{thm:intro-actual-return-defects} concerns $\xi=0$.
We prove a joint small-parameter estimate, including $z=0$, and an
exact lifting identity for every $\xi$. We also improve the regularity
constraint by using the centered second moment when removing the
short collision end blocks.

\subsection{Collision end blocks in the stationary second moment}
\begin{lemma}[Quadratic endpoint budget]
\label{lem:quadratic-endpoint-budget}
There are $a,b,r_0>0$, uniform in $R$, such that
\begin{equation}\label{eq:quadratic-endpoint-budget}
 0<r=|w|\le r_0,\quad H\ge2,\quad r^2(1+\log H)\le a
 \quad\Longrightarrow\quad
 \|\mathcal D_{R,w,\sigma}q\|_1\ge br^2
\end{equation}
for every $\sigma\in\R$ and every collision phase $q$ with
$|q|=1$ and $\|q\|_{\mathrm{BV}}\le H$.
\end{lemma}
\begin{proof}
Use the diffusive middle block of
Lemma~\ref{lem:diffusive-middle-damping}, of length
$m=\lceil A/r^2\rceil$, and the endpoint smoothing in the proof of
Theorem~\ref{thm:near-origin-circle-defect}. Thus
$\epsilon=b_0/H$, $\ell=\lceil b_1+b_2\log H\rceil$, and
$N=m+2\ell$. The lower pairing is at least $1-Nd$, where
$d=\|\mathcal D_{R,w,\sigma}q\|_1$.

The improvement concerns the removal of the phase on the first and
last $\ell$ collisions. Since $h_R$ is centered and its collision
covariances are absolutely summable, uniformly in $R$,
\[
 \|S_{\ell,R}h_R\|_2\le C\sqrt\ell.
\]
All endpoint factors have modulus at most one. Invariance and
$|e^{iu}-e^{iv}|\le|u-v|$ therefore bound the change in the pairing by
\begin{equation}\label{eq:quadratic-endpoint-error}
 |w|\int\left|S_{\ell,R}h_R+
       (S_{\ell,R}h_R)\circ T_R^{m+\ell}\right|\dd\nu
       \le Cr\sqrt\ell.
\end{equation}
No independence of the two end blocks is used. This replaces the
supremum estimate $C r\ell$ in the earlier proof.

After endpoint smoothing and middle-block smoothing, the same
chronological word \eqref{eq:untwisted-end-word} remains. Its two
untwisted powers may be replaced by $\Pi_R$ at cost
$C\epsilon^{-4}\vartheta^\ell$. Its principal pairing has modulus at
most $1/8$. Choose $b_0,b_1,b_2$ exactly as in the earlier proof so
that the two endpoint smoothing and mixing errors total less than
$1/8$. Reduce $r_0$ so that the middle unsmoothing error is less than
$1/16$. Finally choose $a$ so that \eqref{eq:quadratic-endpoint-error}
is less than $1/16$. The original pairing then has modulus at most
$1/2$. The new constraint also gives $N\le C/r^2$, whence
$Nd\ge1/2$ proves \eqref{eq:quadratic-endpoint-budget}.
\end{proof}

\begin{lemma}[Joint small collision frequency and constant phase]
\label{lem:joint-collision-phase-defect}
Let $s=\operatorname{dist}(\sigma,2\pi\Z)$ and
$t=(|w|^2+s^2)^{1/2}$. There are $a,b,t_0>0$ such that
\begin{equation}\label{eq:joint-collision-circle}
 0<t\le t_0,\quad t^2(1+\log H)\le a
 \quad\Longrightarrow\quad
 \|\mathcal D_{R,w,\sigma}q\|_1\ge bt^2
\end{equation}
for every unit-modulus $q$ of variation at most $H\ge2$.
All constants are uniform in $R$.
\end{lemma}
\begin{proof}
Write $r=|w|$ and choose a fixed large number $B$.
If $s\le Br^2$, then $r>0$ and
$t^2\le(1+B^2t_0^2)r^2$. Lemma~\ref{lem:quadratic-endpoint-budget}
proves the claim, after adjusting the constants.

Suppose $s>Br^2$. Centering $q$ in the collision correlation bound
and normalizing its BV norm give
\[
 \left|\int\overline q(q\circ T_R^j)\dd\nu
                   -\left|\int q\dd\nu\right|^2\right|
       \le C H^2e^{-\gamma j}.
\]
Choose $\ell=\lceil b_1+b_2\log H\rceil$ so that this bound is at
most $1/16$ for every $j\ge\ell$. Reduce $t_0$ below $1/4$.
There is an integer
\[
 \ell\le j\le\ell+\lceil2\pi/s\rceil+1
       \quad\hbox{such that}\quad \cos(j\sigma)\le-3/4.
\]
Indeed consecutive phases, in the direction of the representative
of $\sigma$ of absolute value $s$, traverse a full circle in this
many steps; a phase within $s$ of $\pi$ has the required cosine.
This observation includes a negative representative.

The exact iteration of the phase defect implies
\[
 \left|\int\overline q(q\circ T_R^j)\dd\nu
       -e^{ij\sigma}\int e^{iw\cdot S_{j,R}h_R}\dd\nu\right|
       \le j d,\qquad d=\|\mathcal D_{R,w,\sigma}q\|_1.
\]
The centered second moment bounds the difference of the last integral
from one by $C r\sqrt j$. The selected time satisfies
\[
 r^2j\le C(r^2(1+\log H)+r^2/s)\le C(a+B^{-1}).
\]
First choose $B$ large and then $a$ small, so that
$C r\sqrt j\le1/8$. The real part of the first correlation is at
least $-1/16$, whereas that of the second term is at most
$-3/4+1/8$. Hence $jd\ge1/2$.
Finally
\[
 jt^2\le C\bigl(t^2(1+\log H)+r^2/s+s\bigr)
                  \le C(a+B^{-1}+t_0),
\]
so $d\ge bt^2$. This argument also covers $w=0$, $s>0$; thus no
pure constant-phase direction has been lost.
\end{proof}

\begin{proposition}[Joint collision defects for complex functions]
\label{prop:joint-collision-vector-defect}
With $t$ as in Lemma~\ref{lem:joint-collision-phase-defect}, put
$L(H,t)=1+\log(H/t)$. There are uniform $a,b,t_0>0$ such that
\begin{equation}\label{eq:joint-collision-vector}
 0<t\le t_0,\quad t^2L(H,t)\le a,
 \quad \|f\|_2=1,\quad
 \|f\|_\infty+\|f\|_{\mathrm{BV}}\le H
 \quad\Longrightarrow\quad
 \|\mathcal D_{R,w,\sigma}f\|_2\ge\frac{bt^2}{L(H,t)}.
\end{equation}
The function $f$ may vanish.
\end{proposition}
\begin{proof}
Put $e=\|\mathcal D_{R,w,\sigma}f\|_2$ and
$x=\||f|-\int|f|\dd\nu\|_2$. The modulus-variance estimate
\eqref{eq:approximate-modulus-variance} gives
$x^2\le C H^2e^{-\gamma j}+xje$.
For $\eta=\varepsilon_0t^2$, choose $j\le C L(H,t)$ so that the
first term is at most $\eta^2$. Then $x\le\eta+je$.
Lemma~\ref{lem:bv-circle-truncation} provides a unit phase $q$ with
\[
 \|q\|_{\mathrm{BV}}\le C(1+H),\qquad
 \|q-f\|_2\le3\sqrt2(\eta+je).
\]
Its $L^1$ collision defect is at most $e+6\sqrt2(\eta+je)$.
After reducing $a,t_0$, Lemma~\ref{lem:joint-collision-phase-defect}
applies to this phase. Choose $\varepsilon_0$ small enough to absorb
the $\eta$ term. The resulting inequality proves
\eqref{eq:joint-collision-vector}. This is a controlled modulus
argument, not division by an unspecified nonvanishing representative.
\end{proof}

\subsection{The exact spectral phase on the actual tower}
Recall $\eta_R=\one_{Y_R^*}$ and $e_3=(0,0,1,0)$.
For every real $\xi$ define
\begin{equation}\label{eq:peripheral-lift-definition}
 b_{R,\xi}=\eta_R+e^{i\xi}(1-\eta_R),\qquad
 \mathsf L_{R,z,\xi}f=b_{R,\xi}\mathsf L_{R,z}f,
 \qquad \sigma=c_*\xi,\quad w=z-\sigma e_3.
\end{equation}
The multiplier $b_{R,\xi}$ has modulus one and uniformly bounded
variation, independently of $\xi$. In particular its introduction
does not entail a variation estimate for a long return iterate.

\begin{lemma}[Exact peripheral lift and its two normalizations]
\label{lem:exact-peripheral-lift}
For every real $\xi$, every $z\in\R^4$, and $1\le p<\infty$,
whenever the terms are finite,
\begin{align}
 \|\mathsf L_{R,z,\xi}f\|_{L^p(\nu)}^p
     &=c_*\int_{Y_R^*}r_R^*|f|^p\dd\nu_R^*,
                      \label{eq:peripheral-lift-norm}\\
 \|\mathcal D_{R,w,\sigma}(\mathsf L_{R,z,\xi}f)\|_{L^p(\nu)}^p
     &=c_*\|\mathcal E_{R,z,\xi}f\|_{L^p(\nu_R^*)}^p.
                      \label{eq:peripheral-lift-defect}
\end{align}
The defect vanishes at every nontop level of the genuine return tower.
For $Q_{L,\epsilon}$ in \eqref{eq:finite-age-approximation}, the
approximant $b_{R,\xi}Q_{L,\epsilon}$ has the same $L^1$, $L^2$ and
supremum error bounds as in Proposition~\ref{prop:tower-bv-approximation};
its BV bound is at most a fixed multiple of
\eqref{eq:tower-BV-budget}.
\end{lemma}
\begin{proof}
The third collision compensation coordinate is
$h_{R,3}=1-c_*^{-1}\eta_R$. Consequently the real identity
\begin{equation}\label{eq:peripheral-frequency-shift}
 \sigma+w\cdot h_R=z\cdot h_R+\xi\eta_R
\end{equation}
holds, before taking exponentials. On a tower starting at $y\in Y_R^*$,
\[
 \sum_{i=0}^{j-1}\eta_R(T_R^iy)
       =\begin{cases}0,&j=0,\\1,&1\le j\le r_R^*(y).\end{cases}
\]
Thus the lift in \eqref{eq:peripheral-lift-definition} accumulates
exactly the phase in \eqref{eq:peripheral-frequency-shift}. At a
nontop level its collision equation is exact. At the last level its
defect is
\[
 f(F_R^*y)-e^{i(\xi+z\cdot g_R(y))}f(y).
\]
This includes $r_R^*(y)=1$: the initial level is then itself the top.
The disjoint-level integration formula in
Lemma~\ref{lem:exact-tower-defect} gives both identities. The norm
has the height weight; the top defect does not.

Multiplication by $b_{R,\xi}$ preserves pointwise absolute errors.
The BV product inequality gives
$\|b_{R,\xi}Q_{L,\epsilon}\|_{\mathrm{BV}}
 \le C\|Q_{L,\epsilon}\|_{\mathrm{BV}}+CM_f$,
which proves the last assertion. No BV regularity is assumed for the
full exact lift. The construction is periodic in $\xi$, although the
real coordinates $(w,\sigma)$ used to express it change with its lift
from the circle.
\end{proof}

\subsection{A neighborhood of the induced peripheral parameter}
For $H\ge2$, $0<\rho<1/2$, write
\begin{equation}\label{eq:joint-return-budget}
 \Lambda(H,\rho)=1+\log(H/\rho),\qquad
 K(H,\rho)=\Lambda(H,\rho)\log(2+\Lambda(H,\rho)).
\end{equation}

\begin{theorem}[Joint physical and return spectral phases]
\label{thm:joint-return-spectral-defect}
There are uniform constants $a,b,\rho_0>0$ such that, with
\begin{equation}\label{eq:joint-return-domain}
 \xi\in[-\pi,\pi],\qquad
 \rho=(|z|^2+\xi^2)^{1/2},\qquad
 0<\rho\le\rho_0,\qquad \rho^2K(H,\rho)\le a,
\end{equation}
the following estimates hold:
\begin{align}
 |q|=1\text{ on }Y_R^*,\quad \|q^0\|_{\mathrm{BV}}\le H
 &\quad\Longrightarrow\quad
 \|\mathcal E_{R,z,\xi}q\|_{L^1(\nu_R^*)}\ge b\rho^2,
                             \label{eq:joint-return-circle}\\
 \|f\|_{L^2(\nu_R^*)}=1,\quad
 M_f+V_f\le H
 &\quad\Longrightarrow\quad
 \|\mathcal E_{R,z,\xi}f\|_{L^2(\nu_R^*)}
                                  \ge\frac{b\rho^2}{K(H,\rho)}.
                             \label{eq:joint-return-vector}
\end{align}
Zeros of $f$ are allowed. The case $\xi=0$ improves the sufficient
budget of Theorem~\ref{thm:intro-actual-return-defects} from
$|z|K(H,|z|)\le a$ to $|z|^2K(H,|z|)\le a$.
\end{theorem}
\begin{proof}
Put $(w,\sigma)=(z-c_*\xi e_3,c_*\xi)$. For small $\rho$,
$\operatorname{dist}(\sigma,2\pi\Z)=|\sigma|$. The linear map
$(z,\xi)\mapsto(w,\sigma)$ is invertible, with fixed coefficients.
Therefore there are $c_1,C_1>0$ such that
\begin{equation}\label{eq:joint-frequency-comparability}
 c_1\rho\le t=(|w|^2+\sigma^2)^{1/2}\le C_1\rho.
\end{equation}
In particular the line $w=0$ causes no loss: its nonzero constant
collision phase is covered by
Lemma~\ref{lem:joint-collision-phase-defect}.

For a unit section phase, let $Q=\mathsf L_{R,z,\xi}q$.
It has modulus one. Choose accuracy $\eta=\varepsilon_0\rho^2$.
The approximation in Lemma~\ref{lem:exact-peripheral-lift} and
circle truncation give a phase $\widetilde q$ with
\[
 \|\widetilde q-Q\|_1\le4\eta,\qquad
 1+\log(2+\|\widetilde q\|_{\mathrm{BV}})
                           \le C_{\varepsilon_0}K(H,\rho).
\]
Explicitly one may take $L\le C(1+|\log\eta|)$ and
$\epsilon=\eta/(CLH)$. Apply
Lemma~\ref{lem:joint-collision-phase-defect} at $(w,\sigma)$ and
use \eqref{eq:peripheral-lift-defect}. Changing a function by
$4\eta$ changes its $L^1$ collision defect by at most $8\eta$.
Choose $\varepsilon_0$ to absorb this error relative to $b_0c_1^2\rho^2$,
then decrease $a,\rho_0$ so that the required quadratic budget holds.
This proves \eqref{eq:joint-return-circle}.

For a normalized complex section function let
$Q=\mathsf L_{R,z,\xi}f$ and $S=\|Q\|_2$. The exact norm identity
gives $\sqrt{c_*}\le S\le H$. If
$e=\|\mathcal E_{R,z,\xi}f\|_2$, then
\[
 \|\mathcal D_{R,w,\sigma}(Q/S)\|_2
                        =\sqrt{c_*}\,e/S\le e.
\]
Choose now $\eta=\varepsilon_0\rho^2/K(H,\rho)$.
Take $L\le C(1+\log(H/\eta))$ and
$\epsilon=\eta^2/(CH^2L)$ in the $L^2$ approximation and normalize
the resulting approximant. It gives a collision function $v$ with
\begin{align*}
 \|v\|_2&=1,\qquad \|v-Q/S\|_2\le2\eta,\\
 1+\log(2+\|v\|_\infty+\|v\|_{\mathrm{BV}})+|\log t|
       &\le C(1+|\log\varepsilon_0|)^2K(H,\rho).
\end{align*}
These inequalities follow directly by taking logarithms in
\eqref{eq:tower-BV-budget};
$\log K(H,\rho)\le C\Lambda(H,\rho)$ and
\eqref{eq:joint-frequency-comparability} are used here.
They display the dependence on the accuracy choice.
Proposition~\ref{prop:joint-collision-vector-defect} yields
\[
 e+4\eta\ge
 \frac{b_0c_1^2\rho^2}
      {C(1+|\log\varepsilon_0|)^2K(H,\rho)}.
\]
Fix $\varepsilon_0$ so small that the accuracy term is absorbable;
this is possible since
$\varepsilon_0(1+|\log\varepsilon_0|)^2\to0$.
Only after that choice reduce $a,\rho_0$ to satisfy the collision
quadratic budget. This proves \eqref{eq:joint-return-vector} with
no circular choice of a regularity constant.
\end{proof}

\begin{corollary}[Peripheral annuli with enlarged regularity classes]
\label{cor:peripheral-annulus-budgets}
Consider the physical annulus and its rescaling
\begin{equation}\label{eq:peripheral-annulus-scales}
 2n^{-99/200}\le|z|\le n^{-2/5},\qquad
 2n^{1/200}\le|\sqrt n\,z|\le n^{1/10}.
\end{equation}
For $H\le H_0n^\zeta$, there is $a_0>0$ such that the preceding
theorem holds throughout this annulus and
\[
 |\xi|\le\frac{a_0}{\sqrt{\log(2+n)\log\log(e^e+n)}}
\]
for all sufficiently large $n$. Its complex-function lower bound is
at least $b(|z|^2+\xi^2)/[\log(2+n)\log\log(e^e+n)]$.

More generally, fix $\gamma>0$ and
$0<\beta<\min\{4/5,2\gamma\}$. For
$H\le H_0\exp(n^\beta)$ and $|\xi|\le n^{-\gamma}$, the theorem
holds uniformly on \eqref{eq:peripheral-annulus-scales}, with
complex-function lower bound at least
$b(|z|^2+\xi^2)/(n^\beta\log(2+n))$.
\end{corollary}
\begin{proof}
The lower annulus endpoint gives $\rho\ge2n^{-99/200}$.
For polynomial budgets this implies
$K(H,\rho)\le C\log(2+n)\log\log(e^e+n)$.
Choose $a_0$ small enough that the $\xi^2K$ term meets
\eqref{eq:joint-return-domain}; the $|z|^2K$ term tends to zero.
For the exponential budget, $K(H,\rho)\le Cn^\beta\log(2+n)$,
whereas $\rho^2\le n^{-4/5}+n^{-2\gamma}$.
The strict inequalities on $\beta$ make their product tend to zero.
The lower bounds follow from \eqref{eq:joint-return-vector}.
\end{proof}

\paragraph{Spectral scope.}
The lift in Lemma~\ref{lem:exact-peripheral-lift} treats every return
phase exactly. The quantitative conclusion, in contrast, is on the
joint region \eqref{eq:joint-return-domain}; it is not a bound on the
entire unit circle at a growing regularity budget. The pure spectral
direction $z=0$, $\xi\ne0$ and the shifted-frequency cancellation line
are included. Theorem~\ref{thm:joint-return-spectral-defect} remains a
physical function theorem. The next section gives a precise
operator-defect comparison for finite-rank compressions, without
identifying them with an anisotropic realization of the unmodified
return dynamics.

\section{Function reconstruction for compressed actual return operators}
\label{sec:compressed-return-resolvents}
A distributional spectral vector does not automatically belong to a
bounded-variation class. We give a setting where the reconstruction
and the defect comparison can both be proved: finite-rank orthogonal
compressions of the actual return operator. The multiplier uses the
full unbounded return record. Only the function space is compressed.
We keep the mesh dependence explicit and prove a separate finite-time
comparison with the uncompressed characteristic function.

\subsection{An explicit regular function class}
For a sufficiently small mesh size $h>0$, subdivide the five disjoint
coordinate rectangles making up $Y_R^*$ into rectangles with side
lengths between $h/2$ and $h$. Equal subdivisions of each side have
this property. Angular coordinates are taken in a chart through each
rectangle; no artificial small cell is created at the angular seam.
Denote the partition by $\mathcal P_{R,h}$, its piecewise constant
space by $V_{R,h}\subset L^2(\nu_R^*)$, and its conditional expectation
by $P_{R,h}$. The side lengths of the original rectangles are fixed;
only their locations move with $R$.

\begin{lemma}[Uniform reconstruction on the section partition]
\label{lem:section-grid-reconstruction}
Uniformly in $R$ and $h$,
\begin{equation}\label{eq:grid-reconstruction-budget}
 \dim V_{R,h}\le C h^{-2},\qquad
 \|f\|_\infty+\|f^0\|_{\mathrm{BV}}
                         \le C h^{-1}\|f\|_{L^2(\nu_R^*)}
                    \quad(f\in V_{R,h}).
\end{equation}
The map $P_{R,h}$ is an orthogonal projection, preserves constants,
and contracts both $L^2$ and $L^\infty$. For every bounded section
function $v$ with zero extension of bounded variation,
\begin{align}
 \|(I-P_{R,h})v\|_{L^1(\nu_R^*)}
                 &\le Ch\|v^0\|_{\mathrm{BV}},
                       \label{eq:grid-L1-approximation}\\
 \|(I-P_{R,h})v\|_{L^2(\nu_R^*)}^2
                 &\le Ch\|v\|_\infty\|v^0\|_{\mathrm{BV}}.
                       \label{eq:grid-L2-approximation}
\end{align}
The zero extension in \eqref{eq:grid-reconstruction-budget} includes
all jumps at the outer section boundary.
\end{lemma}
\begin{proof}
The normalized collision density on the section is the constant
$(4\pi c_*)^{-1}$ in $(\alpha,p)$ coordinates. Each cell has mass
between $c h^2$ and $C h^2$, and there are at most $C h^{-2}$ cells.
If $f$ has value $a_B$ on cell $B$ and norm one, then
$\sum_B|a_B|^2\le C h^{-2}$ and $\max_B|a_B|\le C h^{-1}$.
Every cell side has length at most $h$. Summing internal jumps and
outer traces bounds the first variation by
\[
 C h\sum_B|a_B|
 \le C h(\#\mathcal P_{R,h})^{1/2}
                       \left(\sum_B|a_B|^2\right)^{1/2}
 \le C h^{-1}.
\]
Real and imaginary parts give the stated complex BV norm. The $L^1$
part of the norm is uniformly bounded, proving
\eqref{eq:grid-reconstruction-budget}.

For a smooth $v$, compare its value to its cell average by integrating
$|v(x)-v(y)|$ over pairs in that cell. Change the two coordinates one
at a time and integrate the fundamental theorem of calculus. Since
the aspect ratios are bounded, the result on a cell is at most $Ch$
times the integral of its first derivatives. BV approximation, or
the same one-dimensional slicing argument for distributional
variation, gives \eqref{eq:grid-L1-approximation}. Summation over
cells and the fixed normalization leave the constant uniform.
Finally $|(I-P)v|\le2\|v\|_\infty$; multiplication of the $L^1$
bound gives \eqref{eq:grid-L2-approximation}. The projection and
contraction properties follow directly from conditional expectation.
\end{proof}

Define the actual twisted Koopman operator and its compression by
\begin{equation}\label{eq:actual-compressed-operator}
 \mathcal U_{R,z}f=e^{-iz\cdot g_R}f\circ F_R^*,\qquad
 A_{R,h,z}=P_{R,h}\mathcal U_{R,z}|_{V_{R,h}}.
\end{equation}
The first operator is unitary on $L^2(\nu_R^*)$: $F_R^*$ is invertible
and measure preserving, and the real-frequency multiplier has modulus
one. Thus $\|A_{R,h,z}\|\le1$. In the orthonormal cell basis
$e_B=\one_B/\sqrt{\nu_R^*(B)}$, its exact matrix entries are
\begin{equation}\label{eq:actual-compressed-matrix}
 (A_{R,h,z})_{BC}=
 \frac{\displaystyle\int_{B\cap(F_R^*)^{-1}C}
                         e^{-iz\cdot g_R}\dd\nu_R^*}
      {\sqrt{\nu_R^*(B)\nu_R^*(C)}}.
\end{equation}
These integrals retain every return length; they are not truncated
flight records or matrices from an independent stochastic model.

\begin{lemma}[Compression residual versus physical defect]
\label{lem:compression-defect-comparison}
Let $U$ be an isometry of a Hilbert space, $P$ an orthogonal projection,
and $A=PU|_{\operatorname{ran}P}$. If $\|f\|=1$, $Pf=f$,
$|\lambda|=1$, and $\varepsilon=\|(A-\lambda)f\|\le1$, then
\begin{equation}\label{eq:compression-defect-comparison}
 \|(U-\lambda)f\|^2
   =\|(A-\lambda)f\|^2+1-\|Af\|^2
   \le3\varepsilon.
\end{equation}
In particular, a physical lower bound
$\|(U-\lambda)f\|\ge d$ for all normalized $f\in\operatorname{ran}P$,
with $0<d\le1$, gives
\begin{equation}\label{eq:compression-lower-singular}
 \|(A-\lambda)f\|\ge(d^2/3)\|f\|.
\end{equation}
\end{lemma}
\begin{proof}
The orthogonal decomposition of $(U-\lambda)f$ has components
$(A-\lambda)f$ and $(I-P)Uf$. The latter has squared norm
$1-\|Af\|^2$. Also $\|Af\|\ge1-\varepsilon$, so
$1-\|Af\|^2\le2\varepsilon-\varepsilon^2$ when
$\varepsilon\le1$. This proves the displayed, slightly weaker,
constant $3$. If $\varepsilon>1$, the lower bound in
\eqref{eq:compression-lower-singular} is immediate. Homogeneity
finishes the proof. The argument controls arbitrary approximate
vectors, not only exact eigenvectors, and makes no normality assumption
on the compression.
\end{proof}

\begin{theorem}[Quantitative peripheral resolvents of the compression]
\label{thm:compressed-peripheral-resolvent}
Choose $H_h=C_0h^{-1}\ge2$, with $C_0$ large enough for
Lemma~\ref{lem:section-grid-reconstruction}. Set
\[
 \rho=(|z|^2+\xi^2)^{1/2},\quad K_h=K(H_h,\rho),\quad
 \delta_h=b_0\rho^4/K_h^2.
\]
There are uniform $a,\rho_0,b_0>0$ such that, whenever
\begin{equation}\label{eq:compressed-peripheral-domain}
 \xi\in[-\pi,\pi],\quad0<\rho\le\rho_0,\quad \rho^2K_h\le a,
\end{equation}
the finite matrix $e^{i\xi}-A_{R,h,z}$ is invertible and
\begin{equation}\label{eq:compressed-resolvent}
 \|(e^{i\xi}-A_{R,h,z})^{-1}\|_{L^2\to L^2}
                                     \le\delta_h^{-1}.
\end{equation}
For $u\ge1$ the same bound holds at $ue^{i\xi}$. It holds with
constant $2\delta_h^{-1}$ for $1-\delta_h/2\le u\le1$.
An eigenvalue $\lambda=ue^{i\xi}$ satisfying the angular and physical
conditions in \eqref{eq:compressed-peripheral-domain} obeys
\begin{equation}\label{eq:compressed-eigenvalue-gap}
 1-u\ge\delta_h.
\end{equation}
The adjoint compression has the corresponding resolvent bound at
$e^{-i\xi}$.
\end{theorem}
\begin{proof}
Every normalized vector of $V_{R,h}$ satisfies the actual section
function budget $H_h$. Also
\[
 \|\mathcal U_{R,z}f-e^{i\xi}f\|_2
                         =\|\mathcal E_{R,z,\xi}f\|_2.
\]
Theorem~\ref{thm:joint-return-spectral-defect} therefore supplies
$d=b\rho^2/K_h$, reduced if necessary so that $d\le1$.
Apply Lemma~\ref{lem:compression-defect-comparison} and choose
$b_0=b^2/3$. Finite dimensionality turns the lower singular-value
bound into invertibility and \eqref{eq:compressed-resolvent}.

For a normalized $f$ and $u\ge1$,
\begin{align*}
 &\|(ue^{i\xi}-A)f\|^2-\|(e^{i\xi}-A)f\|^2\\
 &\qquad=(u-1)\bigl(u+1-2\operatorname{Re}
                         (e^{-i\xi}\langle f,Af\rangle)\bigr)\ge0,
\end{align*}
where $\langle f,g\rangle=\int\overline f g\dd\nu_R^*$ and
$|\langle f,Af\rangle|\le1$.
For the inner radial interval use the reverse triangle inequality;
the lower singular value is at least $\delta_h-(1-u)\ge\delta_h/2$.
If $Af=ue^{i\xi}f$, the singular-value estimate at $e^{i\xi}$ gives
$1-u\ge\delta_h$. Finally take the adjoint of the inverse to obtain
the last assertion. All statements are resolvent bounds for the
specified compression, not a spectral assertion about $\mathcal U_{R,z}$
on the full $L^2$ space.
\end{proof}

\paragraph{Compatibility of scales.}
The finite-time estimate below and the local resolvent estimate have
separate domains. Proposition~\ref{prop:grid-scale-separation} shows
explicitly that even their improved versions cannot be concatenated
on the target annulus using a mesh that makes the displayed time-$n$
grid error vanish. Sections~\ref{sec:direct-orbit-bounds} and
\ref{sec:weighted-spectral-averages} give a direct alternative for
averaged physical coefficients, not an unstated annular power theorem.

\subsection{Finite-time consistency with the unmodified record}
To relate this operator to the raw problem, we record an error bound
which is independent of a spectral convergence assertion. Use
$U_{j,R}=J_{j,R}-j\bar G_R$ and $N_{0,R}=0$.

\begin{lemma}[Variation after a collision-count cutoff]
\label{lem:truncated-return-phase-variation}
For integers $1\le j\le n\le L$ the zero extension of
\[
 v_{j,L}(x)=\one_{\{N_{j,R}(x)\le L\}}
                                  e^{-iz\cdot U_{j,R}(x)},\quad x\in Y_R^*,
\]
satisfies, uniformly in $R$,
\begin{equation}\label{eq:truncated-return-phase-variation}
 \|v_{j,L}\|_\infty\le1,\qquad
 \|v_{j,L}^0\|_{\mathrm{BV}}
                     \le(1+|z|)\exp\{C L\log(2+L)\}.
\end{equation}
\end{lemma}
\begin{proof}
Split according to the collision time $m=N_{j,R}$, $j\le m\le L$.
Starting in $Y_R^*$, this event says that
$\eta_R\circ T_R^m=1$ and exactly $j-1$ of the indicators
$\eta_R\circ T_R^i$, $1\le i<m$, equal one.
Write it as a disjoint union of the at most $2^{m-1}$ binary
membership patterns, multiplied by the initial section indicator.
Each pattern is a product of at most $m+1$ factors from
$\eta_R\circ T_R^i$ and $1-\eta_R\circ T_R^i$.
On this event exact compensation gives $U_{j,R}=S_{m,R}h_R$.

Lemma~\ref{lem:finite-record-variation} and the BV product inequality
bound the variation of each pattern by $C(m+1)B_L$. The phase
$e^{-iz\cdot S_mh_R}$ has supremum one and variation at most
$C(1+|z|m)B_L$, by the BV chain inequality. Sum their product bounds
over $m\le L$ and all patterns. The factors $2^L$ and the polynomial
powers of $L$ are absorbed by increasing the constant in $B_L$.
All section and itinerary jumps, including the initial boundary, have
been counted. No variation bound for an inverse-coarea density is
used.
\end{proof}

\begin{proposition}[An exact telescoping comparison of powers]
\label{prop:compressed-finite-time-comparison}
There are $C,C_1,a,c>0$ such that for $n\ge2$, $L\ge n$,
\begin{align}
 &\left|\int_{Y_R^*}e^{-iz\cdot U_{n,R}}\dd\nu_R^*
             -\langle1,A_{R,h,z}^n1\rangle\right|\notag\\
 &\quad\le Cn\left[
       e^{(an-cL)/2}
       +\sqrt{h(1+|z|)\exp\{C_1L\log(2+L)\}}
                    \right].
                 \label{eq:compressed-finite-time-error}
\end{align}
In particular the matrix pairing approximates the actual characteristic
function at every fixed return time when the cutoff and mesh are
chosen so that the right side tends to zero.
\end{proposition}
\begin{proof}
Abbreviate $\mathcal U_{R,z}$, $P_{R,h}$ and $A_{R,h,z}$ by $U,P,A$.
For $f\in V_{R,h}$ the exact identity
\begin{equation}\label{eq:compression-telescoping}
 A^nf-PU^nf=\sum_{j=0}^{n-1}A^{n-1-j}PU(P-I)U^jf
\end{equation}
follows by telescoping. All operators on the right are contractions.
For $f=1$, $U^j1=e^{-iz\cdot U_{j,R}}$ under the section law.
The cumulative return tail gives
\[
 \|U^j1-v_{j,L}\|_2
       =\nu_R^*(N_{j,R}>L)^{1/2}\le C e^{(an-cL)/2}.
\]
For $j=0$, $P1=1$ and no error occurs. For $1\le j<n$,
Lemma~\ref{lem:section-grid-reconstruction} applied to
Lemma~\ref{lem:truncated-return-phase-variation} gives
\[
 \|(I-P)U^j1\|_2\le C e^{(an-cL)/2}
       +C\sqrt{h(1+|z|)\exp\{C_1L\log(2+L)\}}.
\]
Take the norm in \eqref{eq:compression-telescoping} and pair with the
unit vector $1$. Since $P1=1$, the pairing of $PU^n1$ is the
unmodified characteristic function. This proves the estimate.
\end{proof}

\paragraph{The remaining power and inversion estimates.}
The operator \eqref{eq:actual-compressed-operator} supplies an explicit
regular reconstruction class, a physical defect comparison, and a
local peripheral resolvent. These are proved for every mesh meeting
\eqref{eq:compressed-peripheral-domain}. Passing to the uncompressed
raw record additionally requires a mesh and time scale satisfying
\eqref{eq:compressed-finite-time-error}, control of the other spectral
arcs, and a quantitative conversion of the resulting contour estimates
into powers. The finite-time consistency statement alone supplies
none of these spectral limits. The full $L^2$ operator is unitary;
its contraction properties must not be confused with the spectral
behavior sought on an anisotropic space.

The physical annulus remains \eqref{eq:peripheral-annulus-scales}:
$2n^{-99/200}\le|z|\le n^{-2/5}$, or
$2n^{1/200}\le|\sqrt n\,z|\le n^{1/10}$ after rescaling.
The compact nonzero and growing roof-frequency regions, the far roof
tail, and the global critical/singular raw density extraction remain
the other terms of Theorem~\ref{thm:LLT}. The finite-count cutoff is
used only in the proved error comparison; the matrix entries and the
characteristic function in that comparison retain the full actual
return record. No conditioning event or denominator is replaced.

\section{Complete edge extraction at a finite physical count}
\label{sec:finite-count-extraction}
The individual critical-word calculation does not describe singular
itinerary boundaries. We now include all such boundaries in a finite-count
extraction. The argument integrates the full finite collision graph,
rather than extending an inverse-coarea Jacobian through grazing.
It gives an integrable Fourier remainder for every fixed count cutoff.
The constants of this remainder retain their dependence on the cutoff;
they will not be treated as uniform long-time estimates.

Throughout this section $R\in I$, $n\ge1$, and $L\ge1$ are fixed.
The section is the actual $Y_R^*$, the measure is its probability
$\nu_R^*=c_*^{-1}\nu|_{Y_R^*}$, and $J_{n,R}$ is unchanged. Write
$\ell=(k,m)\in\Z^2\times\mathbb N$ for displacement and physical
collision count. Define
\begin{equation}\label{eq:finite-count-measure}
 \mu_{n,R}^{w,L}=(J_{n,R})_*
       \bigl(w\one_{\{N_{n,R}\le L\}}\nu_R^*\bigr).
\end{equation}
The weight is not normalized. In particular $w=1$ gives the restricted
original law, not the law conditioned on the cutoff event.

\subsection{The whole finite collision graph}
\begin{definition}[Finite-record subanalytic weights]
\label{def:finite-record-weight}
A bounded complex weight is admissible at $(n,L,R)$ if its real and
imaginary parts, on every finite collision chart with $N_{n,R}\le L$,
are globally subanalytic functions of the initial coordinates. Values
on null collision singularities may be fixed arbitrarily. Denote its
supremum bound by $M_w$. This class includes the weight one.
\end{definition}
Here globally subanalytic has its usual real-geometric meaning; it is
not a regularity assertion about every bounded-variation function.
Finite Boolean combinations of subanalytic conditions on a finite list
of actual return states and partial records up to the $n$th return
belong to this class.
Indeed, after fixing the successive return counts, each state is a
specified finite collision iterate. There are only finitely many such
choices below $L$. Products of bounded subanalytic initial, intermediate
and terminal weights are allowed for the same reason. No assertion
about their variation uniformly in the number of observations is made.

\begin{lemma}[Finite graph and actual mixed density]
\label{lem:finite-graph-density}
The measure \eqref{eq:finite-count-measure} has a compactly supported
mixed density $p_{n,R}^{w,L}(\ell,t)$. It has only finitely many nonzero
lattice components. For every fixed $\ell$ and admissible $w$, the real
and imaginary parts of this density have constructible representatives
in the sense of Cluckers--Miller: finite sums of products of globally
subanalytic functions and logarithms of positive globally subanalytic
functions. Each representative is analytic off a finite set.
\end{lemma}
\begin{proof}
Use finitely many bounded rational angular charts and disjoint measurable
chart pieces. Uniform finite horizon gives a finite set of possible
next-disk labels. The finite physical graph in the proof of
Lemma~\ref{lem:finite-record-variation} specifies flight equations, reflection,
positive entry roots and the exclusion of an earlier collision with any
candidate disk. For a fixed word of at most $L$ flights, the graph is
semialgebraic in its initial and auxiliary variables. Projection gives
semialgebraic finite collision records. Return membership adds the exact
rectangle tests at every collision, not just the terminal test. Angular
intervals become finite rational-chart intervals; at fixed $R$ their
endpoints are fixed real coefficients. Thus the events of exactly $n$
returns, total count $m\le L$, and displacement $k$ are finite unions of
semialgebraic chart pieces. The collision probability in a rational
angular coordinate has a bounded semialgebraic density. In particular,
the construction includes the limiting grazing, competing-root and
return-membership boundaries through the domains of integration.

On a regular word, the total flight time is analytic.
Proposition~\ref{prop:general-edge} identifies its critical points as
normal-to-normal nondegenerate minima; in particular its gradient is
nonzero almost everywhere. That proposition is a physical collision-word
statement and does not depend on which measurable section is used to
restrict the initial state. Apply coarea on an exhaustion where the
gradient is bounded below, multiply by the bounded weight, and then sum
the finitely many words. This gives the density, with the normalization
$1/(4\pi c_*)$ in $(\alpha,p)$ coordinates. Singular initial states and
chart boundaries have zero initial measure. Restricting to $Y_R^*$
does not import the old section's return counts. The time support lies
in $[0,L\tau_{\max}]$; the count and displacement labels also range over
finite sets.

We identify the density without bounding an inverse Jacobian near a
singularity. On each chart let $b(x)$ be the product of its measure
density and $w$, and let $\Omega_\ell$ be its exact finite-record event.
For real $t$, the cumulative component is a sum of integrals
\begin{equation}\label{eq:constructible-cumulative}
 P_\ell(t)=\sum_{\rm charts}\int
   \one_{\Omega_\ell}(x)b(x)\one_{\{\tau_m(x)<t\}}\dd x .
\end{equation}
Extend the chart integrands by zero. They are globally subanalytic and
integrable for every $t$. The stability-under-integration theorem
\cite[Theorem 1.3]{CMInt} makes $P_\ell$ constructible. The absolute
continuity just proved identifies its almost-everywhere derivative with
$p_{n,R}^{w,L}(\ell,\cdot)$.

For completeness, only one-variable differentiation is required here.
Write a constructible representation as
$\sum_i a_i(t)\prod_j\log b_{ij}(t)$, with $b_{ij}>0$.
A common finite subanalytic partition makes all generators analytic on
each open interval. Their derivatives are subanalytic there, and the
product rule, with $(\log b)'=b'/b$, gives a constructible derivative.
Define its value arbitrarily at the finitely many endpoints. This is a
constructible representative of the density. The same argument treats
real and imaginary parts and proves the last assertion.
\end{proof}

\begin{remark}[Relation to the retained physical edge calculation]
\label{rem:finite-count-inherited-edge}
The regular critical points and their exponentially small individual
jumps were already classified in Proposition~\ref{prop:general-edge}.
They are not new consequences of the present lemma. The added conclusion
is constructible structure for the \emph{entire} finite-count component,
including all boundaries of its exact integration domain. A singular
boundary may produce a fractional power or a logarithm instead of a
finite jump. The extraction below allows all of these forms.
\end{remark}

\subsection{One-sided power--logarithm jets}
\begin{lemma}[A finite jet removes every second-derivative singularity]
\label{lem:power-log-jet}
Let $f\in L^1(\R;\mathbb C)$ have compact support and constructible
real and imaginary parts. There is a finite set $S\subset\R$ such that
$f$ is analytic off $S$. For every $s\in S$ and sign
$\epsilon\in\{-1,1\}$, on some interval $0<x<d_s$ one has
\begin{equation}\label{eq:power-log-expansion}
 f(s+\epsilon x)=\sum_{k=0}^{K_s}(\log x)^k
                         \sum_{j=j_s}^{\infty}c_{s,\epsilon,j,k}x^{j/q_s},
 \qquad q_s\in\mathbb N.
\end{equation}
The inner series are convergent Laurent series in $x^{1/q_s}$; their
negative parts are finite. After collecting equal exponents, every
nonzero coefficient with $j/q_s\le-1$ vanishes. Put
\begin{equation}\label{eq:power-log-jet}
 J_{s,\epsilon}(x)=
    \sum_{-1<j/q_s\le1}\sum_{k=0}^{K_s}
                           c_{s,\epsilon,j,k}x^{j/q_s}(\log x)^k.
\end{equation}
Then the difference in \eqref{eq:power-log-expansion} after subtracting
this finite jet has, for $a=0,1,2$, a bound
\begin{equation}\label{eq:power-log-remainder}
 \left|\frac{d^a}{dx^a}
       \bigl(f(s+\epsilon x)-J_{s,\epsilon}(x)\bigr)\right|
 \le C_{s,\epsilon}x^{\beta_s-a}(1+|\log x|^{K_s}),
 \qquad \beta_s>1,
\end{equation}
after decreasing $d_s$. An identically zero remainder satisfies the
same assertion with any $\beta_s>1$.
\end{lemma}
\begin{proof}
The convergence used here is part of the analytic-unit representation,
not a conclusion obtained by differentiating an asymptotic expansion.
In \cite[Definitions 2.2--2.3]{CMInt}, a strong unit is $U\circ\phi$,
where $\phi$ is a bounded rational-monomial map and $U$ is analytic
on an open neighborhood of its image closure. At a one-variable
endpoint moved to zero, each nonconstant bounded monomial has a
nonnegative exponent. Clearing denominators turns $\phi$ into
nonnegative integral powers of $x^{1/q}$, so Taylor expansion of $U$
converges near that endpoint. A nonzero prepared center instead gives
an ordinary analytic germ there. The rational prefactor permits only
a finite Laurent principal part. For a positive generator the unit
is positive at zero, and its logarithm has a convergent Taylor series.
These observations also follow directly from
\cite[Theorem 2.4]{CMInt} applied to the subanalytic generators.

The one-variable specialization of constructible preparation
\cite[Theorem 3.11]{CMInt} writes a germ as a finite sum of rational
powers, nonnegative integral logarithmic powers, and subanalytic units.
At a one-sided endpoint, a strong subanalytic unit is analytic in
finitely many bounded rational monomials. Choose a common denominator
for these monomials and the rational prefactors. Expanding the analytic
units gives convergent Laurent series in the corresponding root of $x$.
Equivalently, one may first take the convergent one-variable Puiseux
germs of the subanalytic generators: if a positive generator is
$x^r u(x^{1/q})$ with $u(0)>0$, its logarithm is
$r\log x+\log u(x^{1/q})$. This gives exactly
\eqref{eq:power-log-expansion}, not merely an undifferentiated asymptotic
formula. Finitely many generators and partition endpoints suffice for
all of $f$. Include the endpoints of its support in $S$.

For the integrability assertion, take the smallest exponent having a
nonzero coefficient polynomial in $\log x$. Its highest nonzero log
power dominates all the lower log powers; every strictly higher
rational exponent is smaller by a positive power of $x$ times a fixed
logarithmic power. If that smallest exponent were at most $-1$, the
absolute value would not be integrable at zero. This contradicts
$f\in L^1$. Repeating the argument, or simply using minimality, excludes
all such exponents. Equal-exponent cancellations have already been
performed. The argument works also for complex coefficients.

Only finitely many terms have exponent at most one. After they are
removed, each of the finitely many log coefficients is either zero or
$x^{\beta}a(x^{1/q_s})$ with $\beta>1$ and $a$ analytic at zero.
Differentiate this representation twice. Bounded derivatives of $a$ on
a smaller root interval give \eqref{eq:power-log-remainder}, with the
least of the remaining exponents as $\beta_s$. In particular, termwise
differentiation is justified by the convergent analytic representation.
\end{proof}

Choose disjoint neighborhoods $(s-d_s,s+d_s)$, with $0<d_s<1$.
Let $\chi_s(x)$ be a $C^2$ piecewise-polynomial cutoff on $[0,\infty)$,
equal to one for $x\le d_s/2$ and zero for $x\ge d_s$. Its first two
derivatives vanish at the transition endpoints. Define
\begin{equation}\label{eq:full-finite-jet-density}
 e_f(t)=\sum_{s\in S}\sum_{\epsilon=\pm1}
  \one_{\{\epsilon(t-s)>0\}}\chi_s(\epsilon(t-s))
                         J_{s,\epsilon}(\epsilon(t-s)),
 \qquad r_f=f-e_f.
\end{equation}
The cutoffs need not be analytic; their piecewise-polynomial choice
keeps them subanalytic and is sufficient for two weak derivatives.

\begin{proposition}[Exact subtraction and its derivative budget]
\label{prop:finite-jet-subtraction}
One has $e_f\in L^1(\R)$ and $r_f\in W^{2,1}(\R)$, both with compact
support. The representative of $r_f$ is $C^1$ and has value and first
derivative zero at every point of $S$.

The budget for its second derivative can be recorded without suppressing
its endpoint constants. Put
\[
 \mathcal J(a,k,d)=\int_0^d x^{a-1}|\log x|^k\dd x
 =d^a\sum_{j=0}^k\frac{k!}{(k-j)!}
                      \frac{(-\log d)^{k-j}}{a^{j+1}}
 \quad(a>0,\ 0<d<1).
\]
On $0<x<d_s/2$, the contribution to $\|r_f''\|_1$ is bounded by
\begin{equation}\label{eq:explicit-jet-budget}
 C_{s,\epsilon}\left(
   \mathcal J(\beta_s-1,0,d_s/2)
  +\mathcal J(\beta_s-1,K_s,d_s/2)\right).
\end{equation}
On each transition interval the contribution is bounded by the integral
of $|f''|$ there plus
\[
 \sum_{j,k}|c_{s,\epsilon,j,k}|C_{j/q_s,k}\,
                d_s^{j/q_s-1}(1+|\log d_s|^k).
\]
On the remaining compact regular intervals it is the integral of
$|f''|$. Thus all constants in this finite budget are attached to actual
germs and actual regular intervals.
\end{proposition}
\begin{proof}
Every extracted exponent exceeds $-1$, so $e_f$ is integrable.
Near $s$, \eqref{eq:power-log-remainder} gives zero traces for both
$r_f$ and its first derivative on both sides, and gives an integrable
second derivative since $\beta_s-2>-1$. Away from these points all
functions are $C^2$, with compact support. The matching traces show by
integration by parts that no point mass or derivative of a point mass
remains in the second distributional derivative. This proves
$r_f\in W^{2,1}$ and the asserted representative.

Integrating \eqref{eq:power-log-remainder} for $a=2$ gives
\eqref{eq:explicit-jet-budget}. The identity for $\mathcal J$ follows
by repeated integration by parts, with the boundary term at zero equal
to zero. On $[d_s/2,d_s]$, two differentiations of
$\chi_s(x)x^\alpha(\log x)^k$ and the scaled cutoff bounds give the
stated transition estimate. The remaining regular intervals are
compact subsets of analytic pieces, so their displayed integrals are
finite. These are bounds for the cancellation remainder, not separate
integrals of the divergent second derivatives at $s$.
\end{proof}

\begin{remark}[Why constants and slopes are extracted]
\label{rem:constant-slope-extraction}
Even when the classical second derivative of a one-sided constant or
linear term vanishes, its zero extension has a distributional boundary
term. The constant produces a derivative of a point mass, and the slope
produces a point mass. Including both exponents $0$ and $1$ in
\eqref{eq:power-log-jet} gives an $L^1$ second derivative, stronger than
merely a finite-measure second derivative. Terms such as
$x^{1/2}$, $\log x$, or $x\log x$ must likewise be extracted. The
selection rule concerns the prepared raw density, not the variation of
an initial-coordinate observable.
\end{remark}

\subsection{All finite-count components at once}
\begin{theorem}[Complete finite-count raw extraction]
\label{thm:finite-count-raw-extraction}
For each fixed $n,L,R$ and every admissible weight, there is an exact
identity of finite complex measures
\begin{equation}\label{eq:finite-count-exact-extraction}
 \mu_{n,R}^{w,L}=E_{n,R}^{w,L}+Q_{n,R}^{w,L},
\end{equation}
where $E$ has the finite power--logarithm density
\eqref{eq:full-finite-jet-density} in each nonzero lattice component,
and the density $r_{n,R}^{w,L}(\ell,\cdot)$ of $Q$ belongs to
$W^{2,1}(\R)$. In particular
\begin{equation}\label{eq:finite-count-residual-norms}
 A_j(n,L,R,w)=\sum_\ell\|\partial_t^j
                              r_{n,R}^{w,L}(\ell,\cdot)\|_1<\infty
 \quad(j=0,2).
\end{equation}
Its mixed Fourier transform satisfies
\begin{align}
 \int_{\T^3\times\R}|\widehat Q|\dd\omega
 &\le4(2\pi)^3\sqrt{A_0A_2},
       \label{eq:finite-count-Fourier-integrable}\\
 \int_{\T^3\times\{|b|>B\}}|\widehat Q|\dd\omega
 &\le\frac{2(2\pi)^3}{B}A_2\qquad(B>0).
       \label{eq:finite-count-far-tail}
\end{align}
The construction includes every regular critical, singular itinerary
and image-boundary contribution to the restricted measure. It does
not assume that its only singularities are the retained Morse jumps.
\end{theorem}
\begin{proof}
Apply Lemma~\ref{lem:finite-graph-density} and
Proposition~\ref{prop:finite-jet-subtraction} to each of the finitely
many lattice components. Alternatively apply them to every chart-word
subdensity and then sum; that version records every word's derivative
budget separately. In either construction all actual finite graph
pieces are integrated in \eqref{eq:constructible-cumulative}. Boundary
points themselves have zero initial measure, but their effects on
one-sided density limits are included in these integrals. The prepared
germs of the whole density include those effects. There is no assumption
that an inverse coarea map extends smoothly across such a boundary.

For a component $r_\ell$, distributional integration by parts gives
\[
 |\widehat r_\ell(b)|\le
       \min\{\|r_\ell\|_1,|b|^{-2}\|r_\ell''\|_1\}.
\]
The integral on the real line of $\min\{a,d/b^2\}$ is
$4\sqrt{ad}$ for $a,d>0$. Sum over $\ell$, apply Cauchy--Schwarz,
and integrate the torus of volume $(2\pi)^3$ to obtain
\eqref{eq:finite-count-Fourier-integrable}. A zero norm is handled by
continuity, or by observing that a compactly supported affine $L^1$
function is zero. Integration only over $|b|>B$ gives the other bound.
\end{proof}

The constants in \eqref{eq:finite-count-residual-norms} are finite for
each fixed packet. Neither the preparation theorem nor compactness of
$I$ by itself bounds them uniformly in $R$ as $n,L$ grow. Small cutoff
widths, approaching singular values, jet coefficients and the number of
pieces all occur in Proposition~\ref{prop:finite-jet-subtraction}.
Their total contribution must be estimated before using
\eqref{eq:finite-count-far-tail} in a long-time frequency splice.

\section{Count-localized inversion for the unmodified law}
\label{sec:count-localized-inversion}
The collision count is one of the observed lattice coordinates. Thus a
finite physical-count restriction is exact, not approximate, at every
lattice point whose count lies below the cutoff. This permits the
finite extraction of Section~\ref{sec:finite-count-extraction} to enter
an identity for the original raw density on a central window. The
high-count measure remains in that identity through an explicitly
controlled low-frequency convolution.

Write $\mu_{n,R}=(J_{n,R})_*\nu_R^*$, and let $p_{n,R}$ denote its
mixed density. The density exists by the same countable-word coarea
argument used in Lemma~\ref{lem:finite-graph-density}; the nonnegative
word masses sum to one. For a cutoff $L$, put
\[
 \mu^L_{n,R}=(J_{n,R})_*(\one_{\{N_{n,R}\le L\}}\nu_R^*),
 \qquad \mathsf T^L_{n,R}=\mu_{n,R}-\mu^L_{n,R}.
\]
In this section $E^L,Q^L$ are the exact measures constructed in
Theorem~\ref{thm:finite-count-raw-extraction} for $w=1$, with densities
$e^L,r^L$. We keep the dependence of these choices on $n,L,R$.

\subsection{An exact support property}
\begin{lemma}[Central labels require only a linear count cutoff]
\label{lem:central-count-support}
For $\ell=(k,m)$ with $m\le L$,
\begin{equation}\label{eq:central-count-support}
 p_{n,R}(\ell,t)=p^L_{n,R}(\ell,t)
               =e^L_{n,R}(\ell,t)+r^L_{n,R}(\ell,t)
 \quad\hbox{for almost every }t.
\end{equation}
For $A<\infty$ define the central window
\[
 \mathcal W_{n,R,A}=
 \{(\ell,t)\in\Z^3\times\R:
                  |(\ell,t)-n\bar G_R|\le A\sqrt n\}.
\]
If $\lambda>c_*^{-1}+1$ and $L_n=\lceil\lambda n\rceil$, then
\eqref{eq:central-count-support} holds throughout this window for all
$n\ge\max\{1,A^2\}$, in the mixed almost-everywhere sense.
Moreover, with the constants of
Theorem~\ref{thm:cumulative-return-tail},
\begin{equation}\label{eq:count-tail-for-inversion}
 \|\mathsf T^{L}_{n,R}\|_{\mathrm{TV}}
           \le C e^{an-cL},\qquad
 \sup_\omega|\widehat{\mathsf T^{L}_{n,R}}(\omega)|
           \le C e^{an-cL}.
\end{equation}
\end{lemma}
\begin{proof}
The restriction $N_{n,R}\le L$ is exactly the restriction that the
third lattice coordinate is at most $L$. The measures therefore agree
on each such lattice fiber, and their densities agree almost everywhere
there. The finite extraction supplies the second equality.
On the displayed window the count satisfies
$m\le n/c_*+A\sqrt n\le (c_*^{-1}+1)n$ for $n\ge A^2$.
This proves the asserted coverage. The tail has nonnegative mass equal
to $\nu_R^*(N_{n,R}>L)$; use the cumulative-return tail and then the
basic Fourier bound by total variation. No pointwise density bound for
the high-count measure is deduced from its small mass.
\end{proof}

\subsection{Four terms in an exact inversion}
Choose the cutoff $\chi$ in \eqref{eq:localized-kernel} also to satisfy
$0\le\chi\le1$. Use the proved growing band, namely
\[
 a_n=n^{1/200},\qquad B_n=n^{-99/200},\qquad
 \chi_n(\omega)=\chi(\omega/B_n).
\]
The support has physical radius $2n^{-99/200}$ and rescaled radius
$2n^{1/200}$. The kernel $K_n$ is the mixed inverse transform of
$\chi_n$, with the $(2\pi)^{-4}$ convention of
\eqref{eq:localized-kernel}.

\begin{theorem}[An exact finite extraction on a raw central window]
\label{thm:finite-extraction-central-inversion}
Take $L_n$ as in Lemma~\ref{lem:central-count-support}. On
$\mathcal W_{n,R,A}$, for all sufficiently large $n$,
\begin{align}
 p_{n,R}={}&K_n*\mu_{n,R}
       +(e^{L_n}_{n,R}-K_n*E^{L_n}_{n,R})\notag\\
 &+\mathcal F^{-1}
       [(1-\chi_n)\widehat Q^{L_n}_{n,R}]
       -K_n*\mathsf T^{L_n}_{n,R}.
 \label{eq:count-localized-exact-inversion}
\end{align}
The equality holds almost everywhere. The residual inverse transform
is absolutely convergent, by the proved finite extraction, and
\begin{equation}\label{eq:count-convolution-loss}
 n^2\|K_n*\mathsf T^{L}_{n,R}\|_\infty
          \le C n^{1/50}e^{an-cL}.
\end{equation}
For $L_n=\lceil\lambda n\rceil$ with $\lambda>c_*^{-1}+1$, one also
has, for every $P>0$,
\begin{equation}\label{eq:central-count-Schwartz-tail}
 n^2\sup_{\mathcal W_{n,R,A}}
        |K_n*\mathsf T^{L_n}_{n,R}|\le C_{A,\lambda,P}n^{-P}.
\end{equation}
This last estimate does not require $c\lambda>a$. If that stronger
inequality is imposed, \eqref{eq:count-convolution-loss} additionally
gives exponential smallness on the entire mixed space.
\end{theorem}
\begin{proof}
Theorem~\ref{thm:finite-count-raw-extraction} gives
$\widehat Q^L\in L^1(\T^3\times\R)$, so its inverse is the continuous
representative of the residual density. Splitting that inverse with
$\chi_n$ gives, on the full mixed space,
\[
 p^L=K_n*\mu^L+(e^L-K_n*E^L)
              +\mathcal F^{-1}[(1-\chi_n)\widehat Q^L].
\]
Now substitute $\mu^L=\mu-\mathsf T^L$ and use the exact fiber
agreement of Lemma~\ref{lem:central-count-support} on the window.
Every convolution is defined because its measure is finite and $K_n$
is bounded. This proves \eqref{eq:count-localized-exact-inversion}.
Finally $\|K_n\|_\infty\le CB_n^4$ and
$n^2B_n^4=n^{1/50}$. Combine this with
\eqref{eq:count-tail-for-inversion} to obtain
\eqref{eq:count-convolution-loss}. The power $n^{1/50}$ is retained;
the cutoff error is not declared to be a pointwise tail-density bound.

For the sharper central assertion, the count distance between a point
of the central window and the support of $\mathsf T^{L_n}$ is at least
$(\lambda-c_*^{-1}-1)n$ when $n\ge A^2$. The Schwartz bound in
\eqref{eq:localized-kernel}, and $\|\mathsf T^{L_n}\|_{\mathrm{TV}}\le1$,
therefore give, for every integer $M\ge1$,
\[
 n^2\sup_{\mathcal W_{n,R,A}}|K_n*\mathsf T^{L_n}|
       \le C_{A,\lambda,M} n^{1/50-101M/200}.
\]
Choose $M$ large enough to obtain \eqref{eq:central-count-Schwartz-tail}.
This is decay of the proof kernel across an exactly separated count
support, not decay of the original physical Fourier transform.
\end{proof}

Let
\[
 \mathfrak g_R(y)=\frac{1}{(2\pi)^2\sqrt{\det D_R}}
                    \exp\!\left(-\tfrac12 y^{\mathsf T}D_R^{-1}y\right),
 \qquad y=\frac{(\ell,t)-n\bar G_R}{\sqrt n}.
\]
Uniform positive definiteness was proved in
Theorem~\ref{thm:joint-uniform-nondegeneracy}. Define the actual local
edge correction and finite-band residual integral by
\begin{align*}
 \mathcal E_{n,R,A}^{L}
   &=\mathop{\rm ess\,sup}_{\mathcal W_{n,R,A}}
                   |e^L_{n,R}-K_n*E^L_{n,R}|,\\
 \mathcal C_{n,R}^{L}(B)
   &=\int_{\T^3\times[-B,B]}
               |1-\chi_n(\omega)|\,|\widehat Q^L_{n,R}(\omega)|\dd\omega.
\end{align*}
The first quantity may be infinite until the local density germs are
estimated. Its definition preserves possible cancellation in the exact
edge correction; separate bounds for its two terms also suffice.

\begin{corollary}[A quantitative raw error with the finite branch sum exposed]
\label{cor:finite-extraction-error-budget}
Let $\lambda>c_*^{-1}+1$ and $L_n=\lceil\lambda n\rceil$.
For $P>0$, $B\ge1$, sufficiently large $n$, and every $R\in I$,
the following extended-real inequality holds:
\begin{align}
 &\mathop{\rm ess\,sup}_{\mathcal W_{n,R,A}}
       |n^2p_{n,R}(\ell,t)-\mathfrak g_R(y)|\notag\\
 &\quad\le C n^{-3/280}\sqrt{\log(2+n)}
              +Ce^{-c_0n^{1/100}}
              +C_{A,\lambda,P}n^{-P}\notag\\
 &\qquad\quad+n^2\mathcal E_{n,R,A}^{L_n}
       +(2\pi)^{-4}n^2\mathcal C_{n,R}^{L_n}(B)
       +\frac{n^2}{\pi B}A_2(n,L_n,R,1).
 \label{eq:finite-extraction-error-budget}
\end{align}
Here $C,c_0$ are uniform in the radius; the displayed $A_2$ and edge
terms retain their actual parameter dependence. Thus the raw central
local limit follows if, for some $B=B_n'\ge1$, the last three terms
tend to zero uniformly in $R$ for each fixed $A$.
\end{corollary}
\begin{proof}
In the Fourier integral for $K_n*\mu_{n,R}$, rescale
$v=\sqrt n\,\omega$ and remove the deterministic mean phase.
Theorem~\ref{thm:growing-major-arc} bounds its difference from the
cutoff Gaussian integral by
$C n^{-3/280}\sqrt{\log(2+n)}$ after multiplication by $n^2$.
The complement of the cutoff Gaussian has $|v|\ge n^{1/200}$ and is
bounded by $Ce^{-c_0 n^{1/100}}$ using uniform ellipticity.
These estimates are uniform in the evaluation point, since its Fourier
factor has modulus one.

Apply \eqref{eq:count-localized-exact-inversion}. Its edge term gives
$n^2\mathcal E$, and its high-count convolution gives the third term
in \eqref{eq:finite-extraction-error-budget} by
\eqref{eq:central-count-Schwartz-tail}. Split the absolutely
integrable residual at $|b|=B$. The part inside has precisely the stated
finite-band bound. For $B\ge1$ and sufficiently large $n$, $\chi_n=0$
on $|b|>B$. Equation~\eqref{eq:finite-count-far-tail} bounds this part
by $2(2\pi)^3A_2/B$ before the inverse-transform factor. Since
$(2\pi)^{-4}2(2\pi)^3=1/\pi$, the final constant is as displayed.
This proves the inequality and its sufficient convergence criterion.
\end{proof}

\subsection{Finite bands and weighted observations}
\begin{proposition}[The exact cutoff and a finite-band error budget]
\label{prop:cutoff-finite-band-budget}
For any frequency set $\Omega_n$ of finite volume $V_n$,
\[
 \int_{\Omega_n}|\widehat\mu_{n,R}-\widehat\mu^{L_n}_{n,R}|\dd\omega
       \le C V_n e^{an-cL_n}.
\]
If $V_n\le C e^{\kappa n}$, the choice
\[
 \lambda>\max\{c_*^{-1}+1,(a+\kappa+\delta)/c\},
 \qquad L_n=\lceil\lambda n\rceil,
\]
gives an $O(e^{-\delta n})$ bound and still agrees exactly with the
original raw density on every fixed central window for large $n$.
For a bounded complex weight $w$, the same two statements hold with
an extra factor $\|w\|_\infty$ in the tail estimate. When $w$ is
admissible in Definition~\ref{def:finite-record-weight}, the extraction
and the identity \eqref{eq:count-localized-exact-inversion} also hold
for that weighted measure.
\end{proposition}
\begin{proof}
Integrate \eqref{eq:count-tail-for-inversion} over $\Omega_n$.
The strict inequality in the slope leaves a positive exponential slack.
The support argument is independent of the weight, and the total
variation of the omitted weighted measure is at most its supremum
norm times the omitted probability. Apply
Theorem~\ref{thm:finite-count-raw-extraction} to the weighted finite
measure and repeat the algebra of the inversion identity.
\end{proof}

The last proposition is compatible with the stronger polynomial-weight
truncation estimates in Proposition~\ref{prop:linear-count-budget}.
It is stated here for bounded weights because that is the class used in
the complete finite extraction. For a single initial or actual-return
mark in the inherited BV class, the corresponding central Gaussian
term has the already proved marked estimate, whenever the weight also
satisfies the finite-record subanalytic condition. A finite multiple-time
subanalytic indicator has the structural extraction, but is not thereby
proved to have that marked central estimate. In every case the weight
and its probability denominator remain unchanged.

\paragraph{The remaining long-time estimates.}
The structural finite extraction, its $W^{2,1}$ residual, the exact
central support agreement and the rapidly decaying count correction
are now conclusions. The quantities $\mathcal C_{n,R}^{L_n}(B)$,
$A_2(n,L_n,R,1)$ and $\mathcal E_{n,R,A}^{L_n}$ are the specific remaining
estimates in \eqref{eq:finite-extraction-error-budget}. The first still
contains the intervening annulus, compact nonzero torus frequencies and
the roof frequencies below $B$; its full complement is not bounded by a
one-step function defect. The second has the actual finite-count jet
budget of Proposition~\ref{prop:finite-jet-subtraction}, not a presumed
polynomial growth rate. The third requires local control of singular
values and coefficients on the original central windows. The same
quantitative issues, with the actual indicators, remain in the exact
physical conditioning application. The identity retains the original
raw density throughout; no smoothing or conditioning replaces it.

\section{Exponential finite-record complexity and reconstruction}
\label{sec:exponential-reconstruction}
The finite graph used in Lemma~\ref{lem:finite-record-variation}
has both its number of variables and its number of polynomial tests
linear in the collision length. Retaining the binomial form of the
component bound, rather than its fixed-dimension asymptotic estimate,
gives an exponential variation bound. This removes the additional
logarithm from the regularized tower budget.

\subsection{The finite binomial bound}
\begin{lemma}[Linear-size sign descriptions]
\label{lem:linear-sign-complexity}
Fix $a_0,a_1,d\ge1$. Suppose a family of lifted graphs of length $L$
uses at most $a_0(L+1)$ real variables and $a_1(L+1)$ polynomial
sign tests of degrees at most $d$, and has at most $a_0^{L+1}$
discrete alternatives. The number of connected components of any
Boolean combination of those sign conditions, summed over the
alternatives, is at most $C e^{CL}$. The same bound holds after
coordinate projection and after fixing a bounded number of coordinates
or adjoining a bounded number of fixed-degree tests. The constants
depend on $a_0,a_1,d$, not on the real coefficients.
\end{lemma}
\begin{proof}
Use the finite bound in \cite[Section 3.2, (3.3)]{BPR}, with zeroth
Betti number and ambient variety the full space. If $s$ is the number
of tests and $k\ge1$ the number of variables, it bounds the sum of
the numbers of connected components of all realizable sign conditions
by
\begin{equation}\label{eq:finite-binomial-exponential}
 d(2d-1)^{k-1}\sum_{j=0}^{k}\binom{s}{j}4^j
 \le d(2d-1)^{k-1}5^s.
\end{equation}
Here $\binom{s}{j}=0$ for $j>s$. The right side is exponential in
$k+s$, so it is $Ce^{CL}$ in the stated range. A Boolean combination
is a union of realizable sign conditions. Each of its connected
components contains a connected component of one of those conditions;
thus its number is no larger than their total. Summing over the
exponentially many discrete alternatives preserves an exponential bound.

The image of a connected set under a coordinate projection is connected.
A union of $m$ connected sets has at most $m$ connected components.
This proves the projection assertion without using quantitative
elimination of quantifiers. Fixed coordinates and additional tests
only change the linear bounds by constants. A bounding-ball inequality
can likewise be adjoined. In zero dimension the assertion is immediate.
We use the finite inequality in \eqref{eq:finite-binomial-exponential},
not an $s\to\infty$ expansion with constants depending on $k$.
\end{proof}

\begin{theorem}[Exponential variation of physical finite records]
\label{thm:exponential-finite-record}
There is $C_{\mathrm g}\ge1$, uniform in $R\in I$, such that all the
conclusions of Lemma~\ref{lem:finite-record-variation} hold with
\begin{equation}\label{eq:exponential-record-budget}
 B_L=C_{\mathrm g}e^{C_{\mathrm g}L}
\end{equation}
in place of $\exp\{CL\log(2+L)\}$. In particular, for $|j|\le L$,
\[
 \|\chi\circ T_R^j\|_{\mathrm{BV}}
 +\|h_R\circ T_R^j\|_{\mathrm{BV}}
 +\|\eta_R\circ T_R^j\|_{\mathrm{BV}}\le B_L,
 \qquad
 \|g\circ T_R^j\|_{\mathrm{BV}}\le C\|g\|_{C^1}B_L.
\]
The same exponential bound holds for every binary section-membership
word of length at most $2L+2$, extended by zero off its domain.
The constants are independent of the moving rectangle coefficients.
\end{theorem}
\begin{proof}
Retain the first-admissible-flight graph
\eqref{eq:polynomial-flight-graph}, including its no-earlier-hit
conditions, incidence signs and section-membership tests. Its degree
is bounded independently of $L$, whereas its variable and test counts
are $O(L)$. The number of label and membership alternatives is
exponential in $L$. Fixing one initial coordinate and a superlevel of
a bounded record component adds only a bounded number of variables
and tests. Lemma~\ref{lem:linear-sign-complexity} therefore bounds
the number of interval components of every projected slice superlevel
by $Ce^{CL}$. A finite angular atlas and its seam contributions change
only the constants.

For a bounded scalar function $u$ on an interval, with $|u|\le M$,
one-dimensional coarea gives
\[
 |Du|\le\int_{-M}^{M}\operatorname{Per}\{u>t\}\dd t.
\]
A superlevel with at most $b$ interval components has perimeter at most
$2b$, with the endpoint convention incorporated by zero extension.
Thus the slice variation is at most $4Mb$. Integrating in the other
initial coordinate proves both first distributional derivative bounds.
The statement for binary words follows in exactly the same way by
adjoining all their membership tests to the finite graph; it does not
require differentiating the word's boundary pointwise.

For smooth $g$, cover the angular circle by the same fixed regular
charts and apply the BV chain and product inequalities to the bounded
coordinate record $\chi\circ T_R^j$. The trace terms at the finitely
many chart seams have the same component bound. This gives the stated
$C^1$ composition estimate. Negative iterates use the unique regular
lifted fiber with the terminal state fixed, as in the original graph.
All coefficient dependence enters the sign conditions, to which
\eqref{eq:finite-binomial-exponential} applies uniformly.
\end{proof}

\subsection{The regularized tower with one logarithm}
Use $\mathsf L_{R,z,\xi}$ and $\mathcal E_{R,z,\xi}$ from
Section~\ref{sec:peripheral-return-phases}. Set
$M_f=\|f\|_\infty$ and $V_f=\|f^0\|_{\mathrm{BV}}$ for a section
function. The constants $a_{\mathrm t},c_{\mathrm t}>0$ below are chosen
so that the cumulative-return tail reads
\begin{equation}\label{eq:path-cumulative-tail-convention}
 \nu_R^*\{N_{m,R}>L\}\le C e^{a_{\mathrm t}m-c_{\mathrm t}L}.
\end{equation}
The one-return tail may have different positive constants.

\begin{proposition}[Exponential regularization budget]
\label{prop:exponential-tower-regularization}
For $L\ge1$ and $0<\epsilon<1/2$, the finite-age approximant of
Lemma~\ref{lem:exact-peripheral-lift}, denoted here by $Q$, satisfies
\begin{align}
 \|Q-\mathsf L_{R,z,\xi}f\|_1
   &\le C\{L\epsilon V_f+M_f e^{-cL}\},
       \label{eq:exponential-tower-error}\\
 \|Q\|_\infty&\le M_f,\qquad
 \|Q\|_{\mathrm{BV}}
   \le C M_f e^{CL}(1+\epsilon^{-1}+|z|).
       \label{eq:exponential-tower-BV}
\end{align}
Its squared $L^2$ error is bounded by $2M_f$ times the right side
of \eqref{eq:exponential-tower-error}. These assertions hold for every
real return phase $\xi$.
\end{proposition}
\begin{proof}
The age truncation, initial smoothing, exact top convention and
multiplication by $b_{R,\xi}$ are unchanged. In particular their
$L^1$ errors are precisely those of
Proposition~\ref{prop:tower-bv-approximation} and
Lemma~\ref{lem:exact-peripheral-lift}. For age $j<L$, the level
indicator is a binary backward membership word. Apply
Theorem~\ref{thm:exponential-finite-record} to that word, the smoothed
initial function evaluated at $T_R^{-j}$, and the accumulated phase.
The product and chain inequalities bound the variation by
$CM_f e^{Cj}(1+\epsilon^{-1}+j|z|)$. Summing the ages and absorbing
polynomial factors in $e^{CL}$ gives \eqref{eq:exponential-tower-BV}.
The age levels are disjoint, so the supremum bound is $M_f$, not $LM_f$.
Multiplication by $b_{R,\xi}$ costs only its uniform BV norm. Finally
both the lift and its approximant have modulus at most $M_f$, giving
the $L^2$ assertion. No regularity of the full lift is presumed.
\end{proof}

\begin{theorem}[A single-logarithm joint return estimate]
\label{thm:single-log-return-defect}
Put
\[
 \rho=(|z|^2+\xi^2)^{1/2},\qquad
 \Lambda=1+\log(H/\rho),\qquad H\ge2,
       \quad \xi\in[-\pi,\pi].
\]
There are uniform $a,b,\rho_0>0$ such that $0<\rho\le\rho_0$
and $\rho^2\Lambda\le a$ imply
\begin{align}
 |q|=1,\quad \|q^0\|_{\mathrm{BV}}\le H
 &\quad\Longrightarrow\quad
 \|\mathcal E_{R,z,\xi}q\|_{L^1(\nu_R^*)}\ge b\rho^2,
       \label{eq:single-log-circle}\\
 \|f\|_{L^2(\nu_R^*)}=1,\quad M_f+V_f\le H
 &\quad\Longrightarrow\quad
 \|\mathcal E_{R,z,\xi}f\|_{L^2(\nu_R^*)}
                    \ge b\rho^2/\Lambda.
       \label{eq:single-log-vector}
\end{align}
In the second assertion $f$ may vanish. The domain contains the pure
spectral direction and the shifted-frequency cancellation line.
\end{theorem}
\begin{proof}
Write $\sigma=c_*\xi$ and $w=z-\sigma e_3$. On a sufficiently small
fixed joint ball,
\[
 c_1\rho\le t=(|w|^2+\operatorname{dist}(\sigma,2\pi\Z)^2)^{1/2}
                         \le c_2\rho.
\]
This follows from the invertible linear frequency change in
\eqref{eq:peripheral-frequency-shift}; it does not divide by $|w|$.

For a unit section phase the exact lift has modulus one on the full
collision cylinder. Choose an $L^1$ error $\eta=\epsilon_0\rho^2$,
with $\epsilon_0>0$ fixed later. In
Proposition~\ref{prop:exponential-tower-regularization}, take
$L\le C(1+|\log\eta|)$ large enough for the tail and
$\epsilon=\eta/(CLH)$. Circle truncation, with the same coarea threshold
as in the proof of Theorem~\ref{thm:joint-return-spectral-defect},
gives a unit collision phase $v$ with
\[
 \|v-\mathsf L_{R,z,\xi}q\|_1\le C\eta,\qquad
 1+\log(2+\|v\|_{\mathrm{BV}})
               \le C(1+|\log\epsilon_0|)\Lambda.
\]
Indeed the logarithm in \eqref{eq:exponential-tower-BV} is at most
$CL+\log H+\log(1+\epsilon^{-1})+C$, and
$\log L\le C(1+|\log\eta|)$. Apply
Lemma~\ref{lem:joint-collision-phase-defect}, and use the exact
$L^1$ defect identity \eqref{eq:peripheral-lift-defect}. The change of
defect costs at most twice the approximation error. Choose
$\epsilon_0$ small enough to absorb it and then reduce the joint
smallness constants. This proves \eqref{eq:single-log-circle}.

For a normalized complex section function, set
$Q_*=\mathsf L_{R,z,\xi}f$ and $S=\|Q_*\|_2$. The exact norm
identity gives $\sqrt{c_*}\le S\le H$. If
$e=\|\mathcal E_{R,z,\xi}f\|_2$, the collision defect of $Q_*/S$
is at most $e$. Choose
\[
 \eta=\epsilon_0\rho^2/\Lambda,\qquad
 L\le C(1+\log(H/\eta)),\qquad
 \epsilon=\eta^2/(CH^2L),
\]
with the fixed constants large enough to make the normalized $L^2$
error at most $\eta$. Normalizing the approximant gives $\|v\|_2=1$
and
\begin{align*}
 \|v-Q_*/S\|_2&\le2\eta,\\
 1+\log(2+\|v\|_\infty+\|v\|_{\mathrm{BV}})+|\log t|
       &\le C(1+|\log\epsilon_0|)\Lambda.
\end{align*}
The last inequality follows from
$\log(H/\eta)=\log H+2\log(1/\rho)+\log\Lambda+
\log(1/\epsilon_0)$ and $\log\Lambda\le\Lambda$.
Proposition~\ref{prop:joint-collision-vector-defect} consequently gives
\[
 e+4\eta\ge\frac{b_0c_1^2\rho^2}
                         {C(1+|\log\epsilon_0|)\Lambda}.
\]
Since $\epsilon_0(1+|\log\epsilon_0|)\to0$, choose $\epsilon_0$
first so that the approximation term is absorbable. Only then reduce
$a,\rho_0$ to meet the collision budget. This proves
\eqref{eq:single-log-vector} with no circular accuracy choice.
\end{proof}

\begin{corollary}[Peripheral annuli and compressed resolvents]
\label{cor:single-log-compressed}
For $H\le H_0n^\zeta$, the preceding theorem applies on
\[
 2n^{-99/200}\le|z|\le n^{-2/5},\qquad
 |\xi|\le a_0/\sqrt{\log(2+n)}
\]
for sufficiently small fixed $a_0>0$ and sufficiently large $n$.
Its complex-function lower bound is
$c(|z|^2+\xi^2)/\log(2+n)$. The rescaled physical radii are still
$2n^{1/200}$ and $n^{1/10}$.

For the exact section-grid compression $A_{R,h,z}$ in
Section~\ref{sec:compressed-return-resolvents}, let
$H_h=C/h$ and $\Lambda_h=1+\log(H_h/\rho)$. On the joint region
$0<\rho\le\rho_0$, $\rho^2\Lambda_h\le a$,
\begin{equation}\label{eq:single-log-compressed-resolvent}
 \|(e^{i\xi}-A_{R,h,z})^{-1}\|
                          \le C\Lambda_h^2/\rho^4.
\end{equation}
The same bound, with twice the constant, holds when the resolvent
parameter lies within $c\rho^4/\Lambda_h^2$ of $e^{i\xi}$.
\end{corollary}
\begin{proof}
On the annulus $\Lambda\le C\log(2+n)$, and
$|z|^2\Lambda\to0$ uniformly. Choosing $a_0$ small controls the
spectral contribution to the budget. For a normalized grid vector,
Lemma~\ref{lem:section-grid-reconstruction} gives the budget $H_h$.
If its compressed residual is $\varepsilon\le1$,
Lemma~\ref{lem:compression-defect-comparison} bounds the squared
physical defect by $3\varepsilon$. Theorem~\ref{thm:single-log-return-defect}
therefore gives $\varepsilon\ge c\rho^4/\Lambda_h^2$. Residuals
larger than one satisfy the same lower bound after changing $c$.
This is a smallest-singular-value estimate, hence implies
\eqref{eq:single-log-compressed-resolvent}. Moving the scalar resolvent
parameter by at most half that lower bound preserves half the singular
value by the triangle inequality. No normality assumption is used.
\end{proof}

\begin{proposition}[Improved finite-time consistency]
\label{prop:exponential-compression-consistency}
The finite-time comparison of
Proposition~\ref{prop:compressed-finite-time-comparison} improves to
\begin{equation}\label{eq:exponential-compression-consistency}
 \left|\int_{Y_R^*}e^{-iz\cdot U_{n,R}}\dd\nu_R^*
                       -\langle1,A_{R,h,z}^{\,n}1\rangle\right|
 \le Cn\left\{e^{(a_{\mathrm t}n-c_{\mathrm t}L)/2}
          +\sqrt{h(1+|z|)e^{CL}}\right\},\qquad L\ge n.
\end{equation}
The Fourier sign and Hilbert pairing are those of that proposition;
$U_{n,R}=J_{n,R}-n\bar G_R$ is the actual centered record.
\end{proposition}
\begin{proof}
For each $j\le n$, restrict the actual $j$-return phase to
$\{N_{j,R}\le L\}$. Resolve this event by its binary section word
up to the last collision, and use
Theorem~\ref{thm:exponential-finite-record}. Products of word indicators
and the phase chain rule give variation at most
$C(1+|z|)e^{CL}$; the number of words and all factors polynomial in
$L$ are absorbed in the exponential. The discarded part has $L^2$
norm at most $Ce^{(a_{\mathrm t}n-c_{\mathrm t}L)/2}$.
Cellwise BV approximation on the fitted grid gives the second term
inside braces. Substitute these two estimates in the exact projection
power telescoping identity used in the cited proposition and sum
its $n$ terms. Neither the original law nor the matrix entries are
truncated in \eqref{eq:exponential-compression-consistency}.
\end{proof}

These estimates improve the reconstruction and consistency constants.
They do not assert a mesh-uniform power bound for the uncompressed
operator: the consistency scale and the local resolvent scale must
still be reconciled, and the remaining peripheral arcs must be treated.
First variation of a collision record is also distinct from the
second derivative sum of the raw densities in Theorem I.

\section{Several observations in a logarithmic return window}
\label{sec:logarithmic-return-windows}
A product of observations at different actual returns is not covered
by assuming a variation bound for one of its factors. We prove that
such a product has a controlled approximation when its return window
grows logarithmically. The final event and the record remain unchanged;
the collision cutoff occurs only in the approximation used in the proof.

Let $0\le\ell_1<\cdots<\ell_q\le m$ and let $u_1,\ldots,u_q$
be complex section functions with $\|u_j\|_\infty\le1$ and
zero extensions $u_j^0\in\mathrm{BV}(M)$. Set
\begin{equation}\label{eq:window-weight-definition}
 W_R(y)=\prod_{j=1}^{q}u_j((F_R^*)^{\ell_j}y),\qquad
 V=1+\sum_{j=1}^{q}\|u_j^0\|_{\mathrm{BV}},\qquad
 \alpha_R=\int_{Y_R^*}W_R\dd\nu_R^*.
\end{equation}
All the data may depend on $n,R$. Empty products are one. No
factorization of $\alpha_R$ into individual observation masses is
asserted. For $0\le k\le n-m$ put
\begin{equation}\label{eq:window-transform}
 \mathcal C_{n,k,R}^{W}(v)=
 \int_{Y_R^*}W_R((F_R^*)^kx)
 e^{iv\cdot U_{n,R}(x)/\sqrt n}\dd\nu_R^*(x),
 \qquad U_{n,R}=J_{n,R}-n\bar G_R.
\end{equation}
This convention uses the normalized section probability. Thus all the
observations occur at their actual return states between returns $k$
and $k+m$, including a terminal window when $k=n-m$.

\subsection{A controlled approximation to the whole window}
\begin{lemma}[Truncated and smoothed window weights]
\label{lem:window-weight-approximation}
There is $C_{\mathrm w}\ge1$, uniform in $R,m,q,L$, such that for
integers $L\ge\max\{m,1\}$ and $0<\epsilon<1/2$ there is a section
function $W_{L,\epsilon}$ with
\begin{align}
 \|W_{L,\epsilon}\|_\infty&\le1,\qquad
 \|(W_{L,\epsilon})^0\|_{\mathrm{BV}}
       \le C_{\mathrm w}(1+q\epsilon^{-1})e^{C_{\mathrm w}L},
       \label{eq:window-weight-BV}\\
 \|W_R-W_{L,\epsilon}\|_{L^1(\nu_R^*)}
       &\le C\{\epsilon V+e^{a_{\mathrm t}m-c_{\mathrm t}L}\}.
       \label{eq:window-weight-L1}
\end{align}
The tail constants are those of
\eqref{eq:path-cumulative-tail-convention}.
\end{lemma}
\begin{proof}
Let $v_j=\mathsf S_\epsilon u_j^0$ and define, on the actual section,
\[
 W_{L,\epsilon}(y)=\one_{\{N_{m,R}(y)\le L\}}
                \prod_{j=1}^{q}v_j((F_R^*)^{\ell_j}y).
\]
Extend this function by zero to $M$. Its supremum is at most one.
Telescoping products and invariance of the actual first-return map give
\[
 \int_{Y_R^*}\left|
  \prod_j u_j\circ(F_R^*)^{\ell_j}
  -\prod_j v_j\circ(F_R^*)^{\ell_j}\right|\dd\nu_R^*
 \le \sum_j\|u_j-v_j\|_{L^1(\nu_R^*)}\le C\epsilon V.
\]
This use of invariance is valid before imposing the count cutoff.
The cutoff changes the product by at most its event probability, which
is bounded by \eqref{eq:path-cumulative-tail-convention}. This proves
\eqref{eq:window-weight-L1}.

To prove the variation assertion, resolve $\{N_{m,R}\le L\}$ by
its first $L$ binary section-membership symbols. There are at most
$2^L$ possibilities, with fixed initial membership equal to one. On
each such word the collision indices of the first $m$ actual returns
are fixed integers between zero and $L$. Hence the product above
is a product of $q$ smooth functions $v_j\circ T_R^{t_j}$, multiplied
by the word indicator and the initial section indicator. By
Theorem~\ref{thm:exponential-finite-record} and
$\|v_j\|_{C^1}\le C\epsilon^{-1}$, the BV product inequality bounds
its variation by $C(1+q\epsilon^{-1})e^{CL}$ on each word.
Summing the words and enlarging the exponential constant gives
\eqref{eq:window-weight-BV}. All internal word jumps and the initial
section trace are included. The summands have disjoint supports, which
is why their number does not enter the supremum bound. Null singular
sets are assigned the same convention as in the finite-record theorem.
\end{proof}

\subsection{A variable central band}
\begin{lemma}[The marked estimate on a specified subband]
\label{lem:variable-marked-band}
For $0\le\theta\le1/200$, define
\[
 s_0(\theta)=1/28-5\theta,\qquad
 s_1(\theta)=1/14-4\theta.
\]
For a bounded complex section function $a$, with zero-extension
variation $V_a$, supremum $M_a$, and mass
$\alpha_a=\int a\dd\nu_R^*$, the marked transform under $\nu_R^*$
satisfies
\begin{align}
 &\int_{|v|\le2n^\theta}
 \left|\int a((F_R^*)^kx)e^{iv\cdot U_{n,R}/\sqrt n}\dd\nu_R^*
                 -\alpha_a e^{-v^{\mathsf T}D_Rv/2}\right|\dd v
 \notag\\
 &\hspace{15mm}\le C\{M_a n^{-s_0(\theta)}\sqrt{\log(2+n)}
                           +V_a n^{-s_1(\theta)}\}.
       \label{eq:variable-marked-band}
\end{align}
The constants may be chosen on the full displayed interval of $\theta$,
and uniformly in $0\le k\le n$. The physical radius is
$2n^{-1/2+\theta}$, not $2n^\theta$.
\end{lemma}
\begin{proof}
Use the insertion $c_*^{-1}a^0$ in
Lemma~\ref{lem:marked-spectral} and retain
$\delta=\tfrac14n^{-1/14}$ and a logarithmic spectral block length.
Replace only the rescaled ball by $|v|\le2n^\theta$.
The observable-smoothing term integrates to
$C M_a n^{5\theta-1/28}\sqrt{\log(2+n)}$ and the insertion-smoothing
term to $CV_a n^{4\theta-1/14}$.
The remaining polynomial margins are
\[
 1/14-6\theta,\quad 3/14-5\theta,\quad 1/14-7\theta,
 \quad 1/5-4\theta,\quad 1/5-5\theta,
 \quad 1/2-5\theta,\quad 1-6\theta.
\]
They are strictly larger than $s_0(\theta)$ throughout
$0\le\theta\le1/200$. Their logarithmic factors are therefore
absorbed uniformly; the complementary power is smaller still. The two
analytic smallness conditions hold on the largest ball and hence on
all the smaller balls. This proves \eqref{eq:variable-marked-band}.
At $\theta=1/200$ it gives exactly the exponents $3/280$ and $9/175$
of Theorem C. Thus no inherited central estimate is replaced.
\end{proof}

\begin{theorem}[Central integrals with a full return window]
\label{thm:window-central-integral}
For the weight \eqref{eq:window-weight-definition}, $n\ge2$,
$0\le k\le n-m$, $L\ge\max\{m,1\}$, $0<\epsilon<1/2$, and
$0\le\theta\le1/200$,
\begin{align}
 &\int_{|v|\le2n^\theta}
 \left|\mathcal C_{n,k,R}^{W}(v)
                       -\alpha_R e^{-v^{\mathsf T}D_Rv/2}\right|\dd v
 \notag\\
 &\quad\le C\left\{
 n^{-s_0(\theta)}\sqrt{\log(2+n)}
 +(1+q\epsilon^{-1})e^{C_{\mathrm w}L}n^{-s_1(\theta)}
 +n^{4\theta}\bigl(\epsilon V+
                         e^{a_{\mathrm t}m-c_{\mathrm t}L}\bigr)
                      \right\}.
       \label{eq:window-central-integral}
\end{align}
The constant is independent of the location and the length of the
window and of its observation functions. The mass $\alpha_R$ and
the transform on the left use the full original window, with no
collision cutoff and no smoothing of the observations.
\end{theorem}
\begin{proof}
Apply Lemma~\ref{lem:variable-marked-band} to $W_{L,\epsilon}$ at
return $k$. Lemma~\ref{lem:window-weight-approximation} controls its
supremum and variation. Replacing that weight by $W_R$ changes the
transform pointwise by at most the $L^1$ error in
\eqref{eq:window-weight-L1}; invariance allows its evaluation at
$(F_R^*)^k x$. Integration over the four-dimensional ball costs
$Cn^{4\theta}$ times that error. The difference of the two masses
is bounded by the same error. Uniform positive definiteness gives
\[
 \int_{\R^4}e^{-v^{\mathsf T}D_Rv/2}\dd v\le C,
\]
so the Gaussian mass correction is no larger. This proves the result.
At no step is a randomly evaluated observation replaced by the same
observation at a deterministic collision time.
\end{proof}

\subsection{An unchanged multiple-time event}
\begin{corollary}[Logarithmic windows and rare-event budgets]
\label{cor:logarithmic-window-conditioning}
Suppose, for fixed constants $a,b,d>0$, $\kappa\ge0$, and $r\ge0$,
\[
 m\le a\log n,
 \qquad V\le C n^\kappa[\log(2+n)]^r,
 \qquad b>a.
\]
Take $L=\lceil b\log n\rceil$ and $\epsilon=n^{-d}/4$.
For $0<\theta\le1/200$ put
\begin{align}
 \sigma_1&=1/28-5\theta,&
 \sigma_2&=1/14-4\theta-C_{\mathrm w}b-d,\notag\\
 \sigma_3&=d-\kappa-4\theta,&
 \sigma_4&=c_{\mathrm t}b-a_{\mathrm t}a-4\theta.
       \label{eq:window-exponent-margins}
\end{align}
If all four numbers are positive, the integral in
\eqref{eq:window-central-integral} tends to zero at a power rate,
uniformly in $R,k$ and the permitted data.

For indicators $u_j=\one_{E_j}$, let
\begin{equation}\label{eq:exact-window-event}
 A_{n,k,R}=\bigcap_{j=1}^{q}
       \{x:(F_R^*)^{k+\ell_j}x\in E_j\},\qquad
 p_{n,R}=\nu_R^*(A_{n,k,R}).
\end{equation}
This probability is independent of $k$. If
$p_{n,R}\ge p_0n^{-\beta}>0$ and
$\beta<\min_i\sigma_i$, then
\begin{equation}\label{eq:conditioned-window-band}
 \int_{|v|\le2n^\theta}
 \left|\mathbb E_{\nu_R^*}
   [e^{iv\cdot U_{n,R}/\sqrt n}\mid A_{n,k,R}]
                      -e^{-v^{\mathsf T}D_Rv/2}\right|\dd v
                         \longrightarrow0
\end{equation}
uniformly. The event in this conclusion is exactly
\eqref{eq:exact-window-event}, without an added restriction on
$N_{m,R}$.
\end{corollary}
\begin{proof}
Here $q\le m+1=O(\log n)$. Substitution in
\eqref{eq:window-central-integral} gives its four polynomial margins
\eqref{eq:window-exponent-margins}, multiplied only by fixed powers
of $\log(2+n)$. Strict positivity absorbs those powers. For indicator
weights, invariance gives $\alpha_R=p_{n,R}$, not a product of the
individual probabilities. Divide the unconditional estimate by this
same probability. The strict inequalities with $\beta$ again absorb
the logarithmic factors. No induced mixing hypothesis is used.
\end{proof}

The admissible region is nonempty without a relation between the
complexity and return-tail constants. To exhibit one choice, enlarge
$C_{\mathrm w}$ to at least one and put
\begin{equation}\label{eq:nonempty-window-budget}
 \begin{split}
 d&=1/56,\qquad b=1/(56C_{\mathrm w}),\qquad
 a=\min\{b/2,c_{\mathrm t}b/(4a_{\mathrm t})\},\\
 \theta&=\min\{1/448,c_{\mathrm t}b/32\},\qquad
 \beta=\min\{1/448,c_{\mathrm t}b/16\},\qquad \kappa=0.
 \end{split}
\end{equation}
For every fixed $r$, these constants satisfy
$0<\theta<1/200$ and $0<\beta<\min_i\sigma_i$.
Indeed $C_{\mathrm w}b=d=1/56$ and
$a_{\mathrm t}a\le c_{\mathrm t}b/4$; hence the first three margins
are at least $11/448$, $12/448$, and $4/448$, respectively, and the
fourth is at least $5c_{\mathrm t}b/8$. This gives logarithmically
many observations with bounded individual variation and polynomially
rare exact events. The original $n^{1/200}$ band for the original
single-mark class remains Theorem C; \eqref{eq:window-exponent-margins}
states the cost of the enlarged class, rather than assigning it that
band without proof.

\subsection{Actual moments under the same event}
Set
\[
 \mathcal B=C(1+q\epsilon^{-1})e^{C_{\mathrm w}L},\qquad
 \mathcal E=\min\{1,\epsilon V+e^{a_{\mathrm t}m-c_{\mathrm t}L}\}.
\]
\begin{theorem}[Window-weighted actual moments]
\label{thm:window-weighted-moments}
With the same data, uniformly in $R,k$,
\begin{align}
 \left|\int W_R\circ(F_R^*)^k\,\frac{U_{n,R}}{\sqrt n}\dd\nu_R^*\right|
  &\le C\{n^{-1/6}+\mathcal B n^{-1/2}+\mathcal E^{1/2}\},
       \label{eq:window-first-moment}\\
 \left\|\int W_R\circ(F_R^*)^k\,
           \frac{U_{n,R}\otimes U_{n,R}}n\dd\nu_R^*
                           -\alpha_R D_R\right\|
  &\le C\{n^{-1/6}+\mathcal B n^{-1}+\mathcal E^{1/2}\}.
       \label{eq:window-second-moment}
\end{align}
For \eqref{eq:exact-window-event} with $p_{n,R}\ge p_0n^{-\beta}$,
the conditional mean tends to zero and the conditional covariance
tends to $D_R$ whenever
\begin{equation}\label{eq:window-moment-margins}
 \beta<\min\left\{\frac16,\frac12-C_{\mathrm w}b-d,
             \frac{d-\kappa}{2},
             \frac{c_{\mathrm t}b-a_{\mathrm t}a}{2}\right\}.
\end{equation}
The explicit choice \eqref{eq:nonempty-window-budget} satisfies these
inequalities as well.
\end{theorem}
\begin{proof}
Use $c_*^{-1}(W_{L,\epsilon})^0$ in
Theorem~\ref{thm:marked-quadratic-moments}. This gives the first two
terms of each bound with the approximant's exact mass.
The difference between the weights has modulus at most two, so
\eqref{eq:window-weight-L1} and invariance imply
\[
 \| (W_R-W_{L,\epsilon})\circ(F_R^*)^k\|_2
                               \le C\mathcal E^{1/2}.
\]
The actual second and fourth moments satisfy
$\|U_{n,R}/\sqrt n\|_2\le C$ and
$\|\,|U_{n,R}|^2/n\|_2\le C$ by Theorem D.
Cauchy--Schwarz therefore controls the two changes of weighted moment
by $C\mathcal E^{1/2}$. The change of mass is at most $C\mathcal E$
and $D_R$ is uniformly bounded. This proves
\eqref{eq:window-first-moment}--\eqref{eq:window-second-moment}.

For indicators, divide by the unchanged $p_{n,R}$. Under the
logarithmic choices, the four margins in
\eqref{eq:window-moment-margins} give convergence of the first moment
and of the second moment about the original center. Subtract the tensor
square of the conditional first moment to obtain the covariance.
For \eqref{eq:nonempty-window-budget}, these margins are at least
$1/6$, $13/28$, $1/112$, and $3c_{\mathrm t}b/8$, all strictly larger
than the displayed $\beta$.
\end{proof}

\subsection{The weighted raw inversion interface}
The new central estimate concerns the original window-weighted
measure
\[
 \mu^W_{n,R}=(J_{n,R})_*
                ((W_R\circ(F_R^*)^k)\nu_R^*).
\]
If the observation functions also belong to the bounded finite-record
subanalytic class of Theorem~\ref{thm:finite-count-raw-extraction},
that theorem supplies a finite-count extraction for this same measure.
This is an additional regularity requirement: arbitrary BV observations
are not thereby called subanalytic.

Here is the exact error decomposition which identifies the remaining
weighted estimates. Fix $0<\theta\le1/200$, set
$r_n=n^{-1/2+\theta}$, and choose a smooth cutoff $0\le\chi\le1$, equal
to one on the unit ball and zero outside the ball of radius two. Let
$\chi_{n,\theta}(\omega)=\chi(\omega/r_n)$ on a fixed neighborhood
of zero in $\T^3\times\R$, and let $K_{n,\theta}$ be its mixed
inverse Fourier kernel. Write $\mu^{W,L}=E^{W,L}+Q^{W,L}$ for the
finite extraction, $e^{W,L}$ for its extracted density, and
$T^{W,L}=\mu^W-\mu^{W,L}$. For a central point $y$ whose collision
count does not exceed $L$, the same algebra as in
\eqref{eq:count-localized-exact-inversion} gives
\begin{align}
 p^W(y)={}&(K_{n,\theta}*\mu^W)(y)
       +e^{W,L}(y)-(K_{n,\theta}*E^{W,L})(y)\notag\\
 &+\mathcal F^{-1}[(1-\chi_{n,\theta})\widehat Q^{W,L}](y)
       -(K_{n,\theta}*T^{W,L})(y).
       \label{eq:window-raw-exact-identity}
\end{align}
The inverse transform of the residual is legitimate by Theorem I.
The identity holds for the raw density representatives wherever those
are defined, and hence almost everywhere.

If $L=\lceil\lambda n\rceil$, $\lambda>c_*^{-1}+1$, every fixed
central window has the required count support for large $n$.
Since $|W_R|\le1$, the omitted high-count measure has total variation
at most one. The Poisson--Schwartz kernel estimate in
Section~\ref{sec:count-localized-inversion}, with $n^{1/200}$ replaced
by $n^\theta$, consequently gives
\[
 n^2\sup_{y\text{ in the central window}}
        |(K_{n,\theta}*T^{W,L})(y)|=O(n^{-P})
                    \quad\hbox{for every }P>0.
\]
Indeed its count labels are a distance proportional to $n$ away,
$r_n n=n^{1/2+\theta}$, and the normalized kernel prefactor is
$n^2r_n^4=n^{4\theta}$; any fixed polynomial loss is absorbed by a
sufficient Schwartz order.

Theorem~\ref{thm:window-central-integral} bounds the normalized error
in the first term of \eqref{eq:window-raw-exact-identity} by its
right-hand side, plus $C e^{-c n^{2\theta}}$ from the Gaussian tail.
The other terms require the actual weighted local-edge correction,
the complementary residual integral, and the growth of the finite
second-derivative sum. For an indicator window these requirements
must hold after division by the same $p_{n,R}$.
Thus the central and moment inputs now include a growing multiple-time
class. No weighted tail or exact physical-event replacement is inferred
from them; the latter still concerns a different event from
\eqref{eq:exact-window-event}.

\section{Direct estimates on finite pieces of the actual orbit}
\label{sec:direct-orbit-bounds}
A small residual for a compressed vector and a small error in a
finite-time compression are different requirements. We first record
their scales. We then obtain an estimate on the uncompressed dynamics
by a different argument: short-lag Gaussian estimates bound a finite
Gram matrix of exact orbit vectors. The resulting estimate holds on
every translate of that orbit segment. It involves no reconstruction
of a long-time vector on a spatial grid.

Throughout this and the next section the Hilbert space is
$\mathcal H_R=L^2(Y_R^*,\nu_R^*)$, with
$\langle f,g\rangle_R=\int\overline f g\dd\nu_R^*$.
Write
\begin{equation}\label{eq:direct-orbit-convention}
 \mathcal U_z=\mathcal U_{R,z},\qquad e=\one,\qquad
 c_{j,R}(z)=\langle e,\mathcal U_z^j e\rangle_R
       =\int e^{-iz\cdot U_{j,R}}\dd\nu_R^*,\quad j\ge0.
\end{equation}
Here $U_{j,R}=J_{j,R}-j\bar G_R$ is the centered record and
$\mathcal U_z f=e^{-iz\cdot g_R}f\circ F_R^*$ is the unitary
operator of \eqref{eq:actual-compressed-operator}. For negative $j$
we use its inverse powers; in particular
$c_{-j,R}(z)=\overline{c_{j,R}(z)}$. The symbol $z\in\R^4$ denotes
physical frequency throughout.

\subsection{The two grid scales}
\begin{proposition}[Separation of the displayed grid budgets]
\label{prop:grid-scale-separation}
Let $C_{\mathrm g}>0$ be the exponential variation constant in
\eqref{eq:exponential-compression-consistency}. Suppose $L_n\ge n$
and the grid contribution in that bound tends to zero, that is,
\begin{equation}\label{eq:grid-term-vanishing}
 n\sqrt{h_n(1+|z_n|)e^{C_{\mathrm g}L_n}}\longrightarrow0.
\end{equation}
If $|z_n|\ge2n^{-99/200}$ and
$\rho_n=(|z_n|^2+\xi_n^2)^{1/2}\le\rho_0<1$, then
\[
 \rho_n^2\{1+\log(C_0/(h_n\rho_n))\}\longrightarrow\infty.
\]
Thus these choices do not satisfy the local resolvent domain in
Corollary~\ref{cor:single-log-compressed}.
\end{proposition}
\begin{proof}
Squaring \eqref{eq:grid-term-vanishing} and taking logarithms gives
\[
 \log(1/h_n)-C_{\mathrm g}L_n-2\log n
                 -\log(1+|z_n|)\longrightarrow+\infty.
\]
Consequently $1+\log(C_0/(h_n\rho_n))\ge c n$ for all sufficiently
large $n$, with $c>0$. Since $\rho_n^2\ge4n^{-198/200}$, the product
is at least $4c n^{1/100}$. This proves the assertion.
\end{proof}

This statement concerns simultaneous use of the two displayed bounds,
not the size of the actual approximation error for every mesh. Neither
finite-dimensional theorem is contradicted. The direct argument below
makes no mesh choice and proves a different long-time conclusion:
average square cancellation, rather than decay at every individual
return count. The fixed-return raw inversion problem is unchanged.

\subsection{A uniform estimate on the short lags}
Put
\begin{equation}\label{eq:direct-exponent-definitions}
 \gamma=\frac1{28}-\frac1{200}=\frac{43}{1400},\qquad
 \alpha=\frac{200}{99},\qquad \kappa=\alpha-2=\frac2{99}.
\end{equation}

\begin{lemma}[Pointwise Gaussian comparison on the growing band]
\label{lem:pointwise-lag-Gaussian}
There is $C$ such that for $j\ge2$ and $|z|\le2j^{-99/200}$,
\begin{equation}\label{eq:pointwise-lag-Gaussian}
 \left|c_{j,R}(z)-e^{-jz^{\mathsf T}D_Rz/2}\right|
             \le Cj^{-\gamma}\sqrt{\log(2+j)}
\end{equation}
uniformly in $R$. This is a pointwise estimate, not an inference from
an integrated characteristic-function bound.
\end{lemma}
\begin{proof}
Set $v=-\sqrt j\,z$ and $\delta=j^{-1/14}/4$. Apply
Lemma~\ref{lem:explicit-frequency-error} at collision time
$\lfloor j/c_*\rfloor$, with initial density $s_R$, and use the
pointwise stopping estimate \eqref{eq:weighted-stopping-error}.
The rescaled frequency obeys $|v|\le2j^{1/200}$. Before integration,
the seven polynomial margins from smoothing, covariance replacement,
amplitude, cubic remainder, stopping and integer rounding are bounded
below by the entries in
\[
 \frac1{14},\quad\frac1{28}-\frac1{200},\quad
 \frac1{14}-\frac2{200},\quad\frac3{14}-\frac1{200},\quad
 \frac1{14}-\frac3{200},\quad\frac15-\frac1{200},\quad
 1-\frac2{200}.
\]
The second entry is the unique smallest one. It carries
$\sqrt{\log(2+j)}$; all other logarithmic factors are absorbed by
the strict larger margins. The complementary spectral power is
exponentially small. The analytic conditions hold because
$6/14+3/200<1/2$ and $2/14+1/200<1/2$. Enlarging the constant treats
the finite initial range of $j$, since the two transforms are bounded
by one. This proves \eqref{eq:pointwise-lag-Gaussian} directly from the
pointwise collision estimate.
\end{proof}

For $0<r=|z|\le r_0$, where $r_0$ is sufficiently small, choose
\begin{equation}\label{eq:direct-memory-scale}
 m=m(r)=\lfloor r^{-\alpha}\rfloor\ge2.
\end{equation}
Only lags shorter than $m$ will be estimated by Gaussian comparison.

\begin{proposition}[Summable short-lag budget]
\label{prop:short-lag-row-budget}
Uniformly in $R$ and $0<|z|=r\le r_0$,
\begin{equation}\label{eq:short-lag-row-budget}
 \mathcal B_m(R,z):=1+2\sum_{j=1}^{m-1}|c_{j,R}(z)|
  \le C\{r^{-2}+m^{1-\gamma}\sqrt{\log(2+m)}\},
 \qquad \frac{\mathcal B_m(R,z)}m\le Cr^\kappa.
\end{equation}
\end{proposition}
\begin{proof}
For every $2\le j<m$, the choice of $m$ gives
$rj^{99/200}\le1$, hence the physical frequency is in the band of
Lemma~\ref{lem:pointwise-lag-Gaussian}. Uniform ellipticity gives
$e^{-jz^{\mathsf T}D_Rz/2}\le e^{-d_0jr^2/2}$.
Summing the geometric series costs $Cr^{-2}$, and summing the error
costs $Cm^{1-\gamma}\sqrt{\log(2+m)}$. The lag one and diagonal
terms are absorbed in the first term. For small $r$ one has
$m\ge r^{-\alpha}/2$, so the normalized terms are at most
\[
 Cr^{\alpha-2},\qquad
 Cr^{\alpha\gamma}\sqrt{1+|\log r|}.
\]
Here $\alpha\gamma=43/693$ and $\kappa=14/693$; their difference is
$29/693>0$. The logarithm is therefore absorbed uniformly in
$Cr^\kappa$. No estimate on $c_j$ for $j\ge m$ is used.
\end{proof}

\subsection{Finite orbit Gram matrices}
\begin{lemma}[Analysis and synthesis of an exact orbit segment]
\label{lem:finite-orbit-frame}
Let $U$ be unitary on a complex Hilbert space, $\|e\|=1$, and put
$B_m=1+2\sum_{j=1}^{m-1}|\langle e,U^je\rangle|$.
For every integer $N$, coefficients $a_0,\ldots,a_{m-1}$ and vector
$b$,
\begin{align}
 \left\|\sum_{j=0}^{m-1}a_jU^{N+j}e\right\|^2
                  &\le B_m\sum_{j=0}^{m-1}|a_j|^2,
                         \label{eq:orbit-synthesis}\\
 \sum_{j=0}^{m-1}|\langle b,U^{N+j}e\rangle|^2
                  &\le B_m\|b\|^2.
                         \label{eq:orbit-analysis}
\end{align}
Both assertions remain true after multiplying each orbit vector by an
arbitrary scalar of modulus one.
\end{lemma}
\begin{proof}
The Gram entry at $(i,j)$ is $\langle e,U^{j-i}e\rangle$; it is
independent of $N$. Every absolute row sum is at most $B_m$.
For completeness, expansion of the squared norm and
$2|a_i||a_j|\le|a_i|^2+|a_j|^2$ bound that quadratic form by
$B_m\sum_j|a_j|^2$. This proves \eqref{eq:orbit-synthesis}.
The synthesis operator from $\mathbb C^m$ to the Hilbert space has
norm at most $\sqrt{B_m}$; its adjoint has the same norm and is the
analysis map in \eqref{eq:orbit-analysis}. Multiplication by unit
scalars conjugates the Gram matrix by a diagonal unitary. This proves
the final assertion without assuming that the vectors are orthogonal
or linearly independent.
\end{proof}

\begin{theorem}[Uncompressed moving averages]
\label{thm:direct-moving-averages}
Let $m=m(r)$, $0<|z|=r\le r_0$, and $T\ge m$ be an integer.
For all integers $N$, every $b\in\mathcal H_R$ and every
$\xi\in\R$,
\begin{align}
 \frac1T\sum_{j=0}^{T-1}
     |\langle b,\mathcal U_z^{N+j}e\rangle_R|^2
                         &\le Cr^\kappa\|b\|_2^2,
                            \label{eq:direct-moving-analysis}\\
 \left\|\frac1T\sum_{j=0}^{T-1}
      e^{-ij\xi}\mathcal U_z^{N+j}e\right\|_2^2
                         &\le Cr^\kappa.
                            \label{eq:direct-moving-synthesis}
\end{align}
The constants are independent of $N,T,R,\xi$ and $b$. In particular
these estimates apply to the original characteristic functions at
arbitrarily late return counts, not to powers of a compressed operator.
\end{theorem}
\begin{proof}
Divide the $T$ indices into complete blocks of length $m$ and one
possibly shorter block. The estimate for a shorter block follows from
the same Gram argument with its row sum bounded by $\mathcal B_m$.
There are at most $T/m+1\le2T/m$ blocks. Summing
\eqref{eq:orbit-analysis} proves
\eqref{eq:direct-moving-analysis} with constant
$2\mathcal B_m/m$.

For synthesis, a complete block of $m$ unit-modulus coefficients has
norm at most $\sqrt{m\mathcal B_m}$, and a shorter block has no
larger norm. The triangle inequality across the at most $2T/m$
blocks, followed by division by $T$, gives squared norm at most
$4\mathcal B_m/m$. Proposition~\ref{prop:short-lag-row-budget}
finishes both bounds. Modulation by $e^{-ij\xi}$ does not change any
absolute Gram entry. Thus there is no restriction on the peripheral
angle $\xi$ in this theorem.
\end{proof}

\begin{corollary}[One feasible scale on the full small annulus]
\label{cor:direct-annular-averages}
For sufficiently large $n$ put
\begin{equation}\label{eq:direct-annulus}
 \mathcal A_n=\{z\in\R^4:2n^{-99/200}\le|z|\le n^{-2/5}\}.
\end{equation}
Its rescaled radii are $2n^{1/200}$ and $n^{1/10}$.
For every $z\in\mathcal A_n$ one has
\begin{equation}\label{eq:direct-annular-scale}
 m(|z|)\le2^{-200/99}n<n/4,
 \qquad |z|^\kappa\le n^{-4/495}.
\end{equation}
Consequently both conclusions of
Theorem~\ref{thm:direct-moving-averages} hold for $T\ge n$ with
$Cr^\kappa$ replaced by $Cn^{-4/495}$, uniformly throughout the
annulus and at every peripheral angle.
\end{corollary}
\begin{proof}
The first inequality follows by raising the lower physical radius to
power $-200/99$; the exponent product is exactly one. The second follows
from the upper radius and $(2/5)(2/99)=4/495$. Thus a single explicit
choice of the lag window is compatible with every averaging time
$T\ge n$. No spatial mesh or long-orbit BV budget enters this choice.
\end{proof}

The use of the word ``average'' in these statements is essential.
They bound a sum of squared matrix coefficients or the norm of a
Cesaro average of the actual orbit. They do not bound
$|c_{n,R}(z)|$ at every specified $n$, and they do not assert that
$\|\mathcal U_z^n\|$ decays; that norm equals one. The next section
records the weighted and peripheral consequences at precisely this
level.

\section{Weighted annular averages and the peripheral spectral measure}
\label{sec:weighted-spectral-averages}
The direct orbit estimate has two consequences which require no
variation bound for the final test vector: cancellation under an
unchanged event, and control of the spectral measure of the constant
vector on every peripheral arc. We also identify the small-frequency
limit of that spectral measure. These are statements about the exact
unitary return dynamics, at the level of averages or of its specified
cyclic vector. They are not operator-norm power bounds.

\subsection{Weights and unchanged events}
For a complex $a\in\mathcal H_R$ define
\begin{equation}\label{eq:averaged-weighted-characteristic}
 C^a_{j,R}(z)=\int a(x)e^{-iz\cdot U_{j,R}(x)}\dd\nu_R^*(x),
       \qquad j\ge0.
\end{equation}
The weight is not normalized. Let $\mathcal A_n$ be the physical
annulus in \eqref{eq:direct-annulus}.

\begin{theorem}[Mean-square cancellation with arbitrary square-integrable weights]
\label{thm:weighted-annular-averages}
For $n$ sufficiently large, $T\ge n$, $N\ge0$, and every
$z\in\mathcal A_n$,
\begin{equation}\label{eq:weighted-annular-averages}
 \frac1T\sum_{j=0}^{T-1}|C^a_{N+j,R}(z)|^2
                    \le Cn^{-4/495}\|a\|_2^2.
\end{equation}
If $A$ is any measurable event with $p=\nu_R^*(A)>0$, then
\begin{equation}\label{eq:unchanged-event-average}
 \frac1T\sum_{j=0}^{T-1}
 \left|\mathbb E_{\nu_R^*}
     [e^{-iz\cdot U_{N+j,R}}\mid A]\right|^2
                    \le \frac C p n^{-4/495}.
\end{equation}
The same event $A$ occurs at every averaging index. The constants are
uniform in $R,N,T,z,a,A$ subject to these conditions.
\end{theorem}
\begin{proof}
Take $b=\overline a$ in
\eqref{eq:direct-moving-analysis} and use
Corollary~\ref{cor:direct-annular-averages}. For the conditional
statement take $a=\one_A/p$, whose squared $L^2$ norm is exactly
$1/p$. This proves both estimates with their actual normalizations.
No smoothing, count restriction, subanalytic assumption or factorization
of the event probability is used.
\end{proof}

For example, if $p\ge p_0 n^{-\beta}$ with $\beta<4/495$, the
average in \eqref{eq:unchanged-event-average} tends to zero uniformly.
More quantitatively, for every $\eta>0$, the fraction of indices
$j<T$ at which the displayed conditional transform exceeds $\eta$
in modulus is at most
\begin{equation}\label{eq:exceptional-return-fraction}
 C\eta^{-2}p^{-1}n^{-4/495}.
\end{equation}
This follows by applying the elementary counting inequality to the
sum of squares. The exceptional set may depend on the frequency and
the event. No claim is made that one specified return count avoids it.

\begin{corollary}[A fixed actual mark and a moving terminal mark]
\label{cor:averaged-actual-marks}
The bound \eqref{eq:weighted-annular-averages} also holds for
\[
 \int a((F_R^*)^k x)e^{-iz\cdot U_{N+j,R}(x)}\dd\nu_R^*(x)
\]
when $0\le k\le N$ is fixed throughout the averaging window. It holds
as well when the weight is $a((F_R^*)^{N+j}x)$, evaluated at the actual
terminal return at each index.
\end{corollary}
\begin{proof}
Write $\psi_l=\mathcal U_z^l e$, including negative integers.
The weighted-composition identity for $\mathcal U_z^{-k}$ is
\[
 \psi_{l-k}(y)=\psi_{-k}(y)\psi_l((F_R^*)^{-k}y).
\]
After changing variables $y=(F_R^*)^kx$, the fixed-mark integral is
$\langle\overline a\,\psi_{-k},\psi_{N+j-k}\rangle_R$.
Its test vector has norm $\|a\|_2$ and is independent of $j$.
Apply \eqref{eq:direct-moving-analysis}.

For the terminal mark, change variables $y=(F_R^*)^{N+j}x$.
The phase becomes $\overline{\psi_{-(N+j)}(y)}$. The conjugate of
the resulting integral is $\langle a,\psi_{-(N+j)}\rangle_R$.
Apply the same theorem to the inverse unitary, whose absolute lag
correlations are unchanged. These changes of variables use the actual
return map and do not replace a random observation by a deterministic
collision-clock observation.
\end{proof}

\begin{corollary}[A normalized annular integral]
\label{cor:normalized-annular-integral}
For the original law, and uniformly in $R,N\ge0,T\ge n$,
\begin{equation}\label{eq:normalized-annular-integral}
 \frac1{|\mathcal A_n|}\int_{\mathcal A_n}
     \frac1T\sum_{j=0}^{T-1}|c_{N+j,R}(z)|^2\dd z
                                      \le Cn^{-4/495}.
\end{equation}
The same conclusion holds for an event as in
\eqref{eq:unchanged-event-average}, with the factor $p^{-1}$.
\end{corollary}
\begin{proof}
Integrate the uniform bounds in
Theorem~\ref{thm:weighted-annular-averages}, and divide by the volume
of this nonempty annulus. Finite sums and bounded integrands justify
interchanging the integrals.
\end{proof}

The normalization in \eqref{eq:normalized-annular-integral} is stated
explicitly. The raw local theorem instead requires an \emph{unnormalized}
complementary integral at a specified return count, with the four-dimensional
factor $n^2$. These requirements are not interchangeable. Indeed
$|\mathcal A_n|\asymp n^{-8/5}$; Cauchy--Schwarz applied to the new
mean-square bound would give at best
$n^{2/5-2/495}$ after that raw normalization. It therefore supplies
no vanishing raw error by itself. Likewise a fixed event across an
averaging window is different from a separately chosen exact physical
observation event for every member of that window.

\subsection{Every peripheral arc for the constant vector}
Let $\varsigma_{R,z}$ be the probability on $[-\pi,\pi)$ associated
by the unitary spectral theorem with $e=\one$ and $\mathcal U_z$.
Thus
\begin{equation}\label{eq:cyclic-spectral-measure}
 c_{j,R}(z)=\int e^{ij\theta}\dd\varsigma_{R,z}(\theta),
                         \qquad j\in\Z.
\end{equation}
This is a spectral measure of the return operator, not the time density
of the physical record. The distinction is relevant to both of the
following statements.

\begin{theorem}[Small peripheral arcs and exterior cyclic resolvents]
\label{thm:cyclic-peripheral-bounds}
For $0<|z|=r\le r_0$, $m=\lfloor r^{-200/99}\rfloor$, and all
$\xi\in\R$,
\begin{align}
 \varsigma_{R,z}\{\theta:
   \operatorname{dist}(\theta-\xi,2\pi\Z)\le m^{-1}\}
                                      &\le Cr^{2/99},
                              \label{eq:cyclic-arc-bound}\\
 \left\|(1-s)(I-se^{-i\xi}\mathcal U_z)^{-1}e\right\|_2^2
                                      &\le Cr^{2/99}
                              \label{eq:cyclic-Abel-bound}
\end{align}
whenever $0<s<1$ and $(1-s)m\le1$. All constants are uniform in $R$.
The inverse in \eqref{eq:cyclic-Abel-bound} is an ordinary exterior
unitary resolvent applied to one specified vector. No inverse of
$e^{i\xi}-\mathcal U_z$ is asserted on the unit circle.
\end{theorem}
\begin{proof}
The squared norm of
$m^{-1}\sum_{j<m}e^{-ij\xi}\mathcal U_z^je$ is the integral against
$\varsigma_{R,z}$ of
\[
 \left|\frac1m\sum_{j=0}^{m-1}e^{ij(\theta-\xi)}\right|^2.
\]
On the arc of radius $1/m$ this kernel is bounded below by a positive
absolute constant. For example, use
$|\sin(m t/2)|/(m|\sin(t/2)|)\ge2/\pi$ for $|t|\le1/m$,
with its continuous value at zero. Equation~\eqref{eq:orbit-synthesis}
and Proposition~\ref{prop:short-lag-row-budget} bound the total integral
by $Cr^{2/99}$. This proves \eqref{eq:cyclic-arc-bound} simultaneously
for all centers, with no choice of peripheral representatives.

For the second assertion expand the resolvent in its norm-convergent
Neumann series. In its block of indices $km,\ldots,km+m-1$,
\eqref{eq:orbit-synthesis} gives the norm bound
\[
 s^{km}\sqrt{\mathcal B_m\sum_{j=0}^{m-1}s^{2j}}
                         \le s^{km}\sqrt{m\mathcal B_m}.
\]
Sum the blocks and multiply by $1-s$. Since
$s^m\le e^{-m(1-s)}$ and
$1-e^{-x}\ge(1-e^{-1})x$ for $0\le x\le1$, the result is at most
$C\sqrt{\mathcal B_m/m}$. Squaring and applying
\eqref{eq:short-lag-row-budget} proves \eqref{eq:cyclic-Abel-bound}.
The proof controls no operator norm on arbitrary input vectors.
\end{proof}

\subsection{A Cauchy scaling law for the exact spectral measure}
For a unit direction $u\in\R^4$ put
\[
 d_{R,u}=u^{\mathsf T}D_Ru,\qquad a_{R,u}=d_{R,u}/2.
\]
Uniform ellipticity bounds $a_{R,u}$ above and below by positive
constants. Let $\mathsf C_a$ be the Cauchy probability on $\R$ with
density $a/[\pi(a^2+x^2)]$ and characteristic function $e^{-a|t|}$.
The principal angle in \eqref{eq:cyclic-spectral-measure} is used
in the rescaling below.

\begin{theorem}[Spectral Cauchy limit at physical frequency zero]
\label{thm:spectral-Cauchy-limit}
Let $\widetilde\varsigma_{R,r,u}$ be the pushforward of
$\varsigma_{R,ru}$ under $\theta\mapsto\theta/r^2$.
Then, as $r\downarrow0$,
\begin{equation}\label{eq:spectral-Cauchy-limit}
 \sup_{R\in I,\ |u|=1}
 d_{\rm BL}\bigl(\widetilde\varsigma_{R,r,u},
                         \mathsf C_{a_{R,u}}\bigr)\longrightarrow0.
\end{equation}
This describes a spectral probability. It does not claim that the
physical return record has a Cauchy limit; its Gaussian limit is
unchanged.
\end{theorem}
\begin{proof}
The one-return second moment is uniformly bounded by
Theorem~\ref{thm:exponential-return}. Using $1-\cos x\le x^2/2$,
\[
 1-\operatorname{Re}c_{1,R}(ru)
      \le Cr^2.
\]
For $|\theta|\le\pi$ one has
$\theta^2\le(\pi^2/2)(1-\cos\theta)$. Therefore
\begin{equation}\label{eq:principal-angle-rounding}
 \int|\theta|\dd\varsigma_{R,ru}(\theta)\le Cr.
\end{equation}
Fix $t>0$ and put $j=\lfloor t/r^2\rfloor$. The difference between
the characteristic function of $\widetilde\varsigma_{R,r,u}$ at $t$
and $c_{j,R}(ru)$ is at most $Cr$, by
\eqref{eq:principal-angle-rounding} and $|t/r^2-j|<1$.
For sufficiently small $r$, $j\ge2$, and Theorem A applied at
rescaled frequency $-r\sqrt j\,u$ gives
\[
 c_{j,R}(ru)=\exp(-jr^2d_{R,u}/2)+
               O_t\bigl(j^{-1/26}\sqrt{\log(2+j)}\bigr).
\]
The rescaled frequencies remain in a fixed ball for fixed $t$.
The main term tends uniformly to $e^{-td_{R,u}/2}$. Negative $t$
uses $c_{-j}=\overline{c_j}$, and $t=0$ is exact. This proves
pointwise characteristic convergence, uniformly in the parameters,
to the asserted Cauchy family.

To justify the uniform bounded-Lipschitz conclusion, take any sequence
$r_l\downarrow0$, $R_l\in I$, $|u_l|=1$. Compactness gives a
subsequence $R_l\to R_*$, $u_l\to u_*$. Continuity of $D_R$ and the
preceding characteristic estimates identify its limit as
$\mathsf C_{a_{R_*,u_*}}$ by the characteristic-function continuity
theorem. The target Cauchy laws converge as well: multiply one standard
Cauchy random variable by $a_{R_l,u_l}$. Thus a sequence violating
\eqref{eq:spectral-Cauchy-limit} cannot exist. No absolute continuity
of $\varsigma_{R,z}$ for a fixed nonzero $z$ is presumed or deduced.
\end{proof}

\paragraph{Place in the raw inversion problem.}
Theorems~\ref{thm:direct-moving-averages},
\ref{thm:weighted-annular-averages} and
\ref{thm:cyclic-peripheral-bounds} bypass the incompatible grid scales
for averaged physical coefficients and for a specified cyclic
resolvent vector. They cover every peripheral angle and allow
square-integrable weights without a BV reconstruction hypothesis.
Theorem~\ref{thm:spectral-Cauchy-limit} identifies the concentration
scale of the same spectral probability. None of these statements
replaces the fixed-return complementary integral in
Theorem~\ref{thm:finite-extraction-central-inversion}, the long-time growth of
$A_2(n,L_n,R,w)$, or the local extracted-edge correction. The exact
physical-event replacement remains a different assertion from the
unchanged-event average proved above.

\section{Finite-order collision cumulants and spectral damping}
\label{sec:finite-order-damping}
The short-lag argument of Section~\ref{sec:direct-orbit-bounds} does
not estimate a prescribed later coefficient. We now return to the
collision spectral splitting and prove stronger finite-order bounds
before stopping. The point is that its derivatives at zero can be
bounded independently of the smoothing scale, although the analytic
neighborhood still shrinks with that scale. A Taylor polynomial of
fixed, sufficiently high degree then gives real-frequency damping on
a larger band. No analyticity of the unsmoothed twist is asserted.

All expectations in this section use the stationary collision
probability $\nu$. Put $\|u\|_{\mathcal V}=\|u\|_\infty+
\|u\|_{\mathrm{BV}}$. Every order below is fixed before the smoothing
scale, collision length, and radius are varied. Constants may depend
on that order. We require no factorial bound uniform in the order.

\subsection{Connected correlations at any fixed order}
\begin{lemma}[Finite-product decoupling at arbitrary fixed order]
\label{lem:all-finite-product}
For each integer $q\ge2$ there are $C_q,\gamma_q>0$, uniform in $R$,
such that, for $t_1\le\cdots\le t_q$ and $1\le p<q$,
\begin{align}
 &\left|\int\prod_{j=1}^q u_j\circ T_R^{t_j}\dd\nu
 -\left(\int\prod_{j\le p}u_j\circ T_R^{t_j}\dd\nu\right)
  \left(\int\prod_{j>p}u_j\circ T_R^{t_j}\dd\nu\right)\right|
 \notag\\
 &\hspace{12mm}\le C_qe^{-\gamma_q(t_{p+1}-t_p)}
                         \prod_{j=1}^q\|u_j\|_{\mathcal V}.
 \label{eq:all-finite-product}
\end{align}
Negative and repeated times are allowed. The constants can be chosen
simultaneously for orders at most any specified integer $Q$.
\end{lemma}
\begin{proof}
The proof of Lemma~\ref{lem:bv-finite-product} works at each fixed
order, and we give its scale explicitly. By multilinearity normalize
all nonzero norms to one. Smooth every factor at scale $\epsilon$.
Telescoping the products and using invariance bounds the total error
by $C_q\epsilon$. Translate $t_1$ to zero, write the expectation as
the chronological word of smooth multipliers and collision powers,
and replace the power across the selected gap by $\Pi_R=\nu\ell$.
The new word is exactly the product of the two block expectations.
The error is at most $C_q\epsilon^{-2q}\vartheta^{t_{p+1}-t_p}$:
there are $q$ smooth multipliers, all other powers are uniformly
bounded, and the selected complementary power decays exponentially.
For a zero gap use $I-\Pi_R$ with its uniform norm. Taking
$\epsilon=\vartheta^{(t_{p+1}-t_p)/(2q+1)}/4$ proves the assertion.
A finite minimum of the resulting exponents treats orders at most $Q$.
No variation of a long pullback is estimated.
\end{proof}

For bounded real random variables, define the joint cumulant by
\begin{equation}\label{eq:finite-order-cumulant-definition}
 \operatorname{cum}(X_1,\ldots,X_q)
 =\sum_{\pi\in\mathcal P_q}(-1)^{|\pi|-1}(|\pi|-1)!
                    \prod_{B\in\pi}\mathbb E\prod_{j\in B}X_j,
\end{equation}
where $\mathcal P_q$ is the set of partitions of $\{1,\ldots,q\}$.
This is the mixed derivative at zero of the logarithm of the joint
moment generating function. In particular it is multilinear and
vanishes when the variables split into two nonempty independent
families. The latter assertion follows because the joint generating
function factors and its logarithm is a sum of two functions in
separate sets of variables. This elementary finite-order use of
cumulants is distinct from an all-orders normal-approximation theorem;
see also \cite{DJS} for the general framework.

\begin{proposition}[Absolute summability of fixed-order connected correlations]
\label{prop:connected-collision-summability}
Let $u$ be real and $\|u\|_{\mathcal V}\le B$. For each $q\ge2$,
\begin{equation}\label{eq:connected-collision-summability}
 \sum_{\mathbf k\in\Z^{q-1}}
 (1+\operatorname{diam}\{0,k_2,\ldots,k_q\})
 \left|\operatorname{cum}(u,u\circ T_R^{k_2},\ldots,
                                      u\circ T_R^{k_q})\right|
 \le C_{q,B}.
\end{equation}
Consequently the absolutely convergent sum without the diameter
factor, denoted by $\mathfrak c_q(R,u)$, satisfies
\begin{equation}\label{eq:cumulant-linear-length}
 \left|\operatorname{cum}_q(S_{m,R}u)-m\mathfrak c_q(R,u)\right|
                \le C_{q,B}\qquad(m\ge1),
 \qquad |\mathfrak c_q(R,u)|\le C_{q,B}.
\end{equation}
Both statements are uniform in $R$ and in the displayed class of $u$.
\end{proposition}
\begin{proof}
Order any $q$ times and split them at a largest consecutive gap $d$.
Compare their joint law with a law in which the entire left block and
the entire right block are independent, but the two block marginals
are unchanged. For every subset meeting both blocks, its moment
changes by at most $C_{q,B}e^{-\gamma_qd}$ by
Lemma~\ref{lem:all-finite-product}; the gap in that subset is at least
$d$. All other subset moments are unchanged. The finite partition
polynomial \eqref{eq:finite-order-cumulant-definition}, whose factors
are bounded, therefore changes by at most the same order. Its value
under the independent-block law is zero. Thus the original joint
cumulant is bounded by $C_{q,B}e^{-\gamma_qd}$. If all times coincide,
the bounded partition formula gives this estimate with $d=0$.

For anchored times $0,k_2,\ldots,k_q$, their diameter $D$ obeys
$d\ge D/(q-1)$. There are at most $(2D+1)^{q-1}$ tuples with all
coordinates of absolute value at most $D$. Exponential decay therefore
sums with the additional factor $1+D$, proving
\eqref{eq:connected-collision-summability}.

By multilinearity and stationarity, the exact finite sum is
\[
 \operatorname{cum}_q(S_{m,R}u)
 =\sum_{\mathbf k\in\Z^{q-1}}(m-D(\mathbf k))_+
   \operatorname{cum}(u,u\circ T_R^{k_2},\ldots,u\circ T_R^{k_q}).
\]
Indeed, after fixing the first index, the remaining indices have these
relative offsets; exactly $(m-D)_+$ translates fit in the interval of
$m$ collisions. The difference between this coefficient and $m$ has
absolute value at most $D$. Applying the first assertion proves
\eqref{eq:cumulant-linear-length}. This also treats repeated indices
without an additional combinatorial multiplicity.
\end{proof}

\subsection{Uniform derivatives of the smoothed eigenvalue}
For a real unit vector $\xi\in\R^4$ let
\[
 u_{R,\delta,\xi}=\xi\cdot h_{R,\delta},\qquad
 \psi_{R,\delta,\xi}(t)=\log\lambda_{R,\delta}(t\xi),
\]
using the branch of the logarithm fixed in
Lemma~\ref{lem:smooth-collision-expansion}. The supremum and BV norms
of $u_{R,\delta,\xi}$ are bounded uniformly, and its mean is zero.

\begin{theorem}[A smoothing-independent finite spectral jet]
\label{thm:finite-spectral-jet}
For each fixed $Q\ge3$ there is $C_Q$ such that
\begin{align}
 \psi'(0)&=0,\qquad
 \psi''(0)=-\xi^{\mathsf T}\Gamma_{R,\delta}\xi,\notag\\
 \psi^{(q)}(0)&=i^q\mathfrak c_q(R,u_{R,\delta,\xi}),\qquad
 |\psi^{(q)}(0)|\le C_Q\quad(2\le q\le Q).
 \label{eq:finite-spectral-jet}
\end{align}
In particular, on $|t|\le a\delta^2$, after decreasing $a>0$,
\begin{equation}\label{eq:finite-spectral-Taylor}
 \psi(t)=-\tfrac12t^2\xi^{\mathsf T}\Gamma_{R,\delta}\xi
           +\sum_{q=3}^{Q}\frac{\psi^{(q)}(0)}{q!}t^q
           +O_Q(\delta^{-2(Q+1)}|t|^{Q+1}).
\end{equation}
Every constant is uniform in $R,\delta,\xi$. The Taylor remainder still
uses the shrinking complex analytic radius; it is not scale independent.
\end{theorem}
\begin{proof}
Fix $R,\delta,\xi$ temporarily. The stationary spectral pairing is
\[
 F_m(t):=\int e^{itS_{m,R}u_{R,\delta,\xi}}\dd\nu
               =e^{m\psi(t)}A(t)+E_m(t),\qquad A(t)=\ell(\Pi(t\xi)\nu).
\]
Here $A(0)=1$, and the complementary term has norm at most $C\rho^m$
on the given complex neighborhood. On a smaller neighborhood,
$|A(t)|\ge1/2$ and $|e^{\psi(t)}|\ge\rho_1>\rho$. Hence the relative
error is analytic and bounded by $C(\rho/\rho_1)^m$. For large $m$,
choose its logarithm near one. Cauchy's inequalities on a still smaller
neighborhood give, for each fixed $q$,
\[
 \frac1m(\log F_m)^{(q)}(0)\longrightarrow\psi^{(q)}(0).
\]
The derivatives of $\log A$ divided by $m$ vanish. All these
identification steps are made at fixed $\delta$; their constants need
not be uniform as $\delta\downarrow0$.

Boundedness of the finite collision sum and the cumulant definition
give $(\log F_m)^{(q)}(0)=i^q\operatorname{cum}_q(S_m u)$.
Proposition~\ref{prop:connected-collision-summability} identifies the
limit and bounds it independently of $\delta,R,\xi$.
The first derivative vanishes by centering; the second is the collision
Green--Kubo variance, also identified in
Lemma~\ref{lem:smooth-collision-expansion}.

Finally the analytic logarithm is uniformly bounded on a complex disk
of radius a fixed multiple of $\delta^2$, as in the proof of that
lemma. Cauchy's estimate for the Taylor tail on a concentric smaller
disk is $C_Q\delta^{-2(Q+1)}|t|^{Q+1}$. This proves
\eqref{eq:finite-spectral-Taylor}. Only the finitely many coefficients,
not the analytic disk or an infinite cumulant series, have gained
smoothing-independent bounds.
\end{proof}

\begin{corollary}[Real-frequency collision damping]
\label{cor:finite-jet-damping}
For each fixed $Q\ge3$ there are $c,C,a,\delta_0>0$ such that
$0<\delta\le\delta_0$ and
\begin{equation}\label{eq:finite-jet-damping-domain}
 |z|\le a\delta^2,\qquad
 \delta^{-2(Q+1)}|z|^{Q-1}\le a
\end{equation}
imply, on every local collision space,
\begin{equation}\label{eq:finite-jet-damping}
 |\lambda_{R,\delta}(z)|\le e^{-c|z|^2},\qquad
 \|\mathcal L_{R,\delta,z}^{m}\|\le Ce^{-cm|z|^2}
                \quad(m\ge0,\ z\in\R^4).
\end{equation}
\end{corollary}
\begin{proof}
Theorem~\ref{thm:joint-uniform-nondegeneracy} and
\eqref{eq:collision-covariance-limit} give
$\Gamma_{R,\delta}\ge c_1I_4$ for all sufficiently small $\delta$,
uniformly in $R$. Write $z=t\xi$. The finite sum of degree at least
three in \eqref{eq:finite-spectral-Taylor} is bounded by $C_Q|t|^3$
for $|t|\le1$. Its remainder is at most $C_Qa|t|^2$ under
\eqref{eq:finite-jet-damping-domain}. Decreasing $a$ and $\delta_0$
therefore gives $\operatorname{Re}\psi(t)\le-c_2|t|^2$.
The bounded spectral projection and the complementary power estimate
then give
$\|\mathcal L_z^m\|\le C(e^{-c_2m|z|^2}+\rho^m)$.
Decrease the constant in the exponent so that $\rho^m$ is dominated
by the same bound on the stated small real-frequency region. This
also holds for $m=0$ after increasing $C$.
\end{proof}

The estimate is a power bound for the \emph{smoothed collision}
operator on its previously constructed local distribution space.
It is not an operator-norm decay assertion for the unitary twisted
return Koopman operator. The next section removes smoothing and
stops at the actual prescribed return, keeping those errors explicit.

\section{A fixed-return inner annulus on the raw inversion scale}
\label{sec:fixed-return-annulus}
We apply the finite spectral jet before any averaging in the return
count. The result estimates the original characteristic function at
each specified count, also with an insertion at any actual return.
The new band extends beyond the previously proved central ball.
The farther complement and the extracted-edge estimates remain
separate parts of the same raw inversion problem.

Retain $M_a,V_a,\alpha_a$ and $C^{[k],a}_{n,R}$ from
\eqref{eq:marked-norms} and \eqref{eq:marked-transform}. Thus $a$ is
zero outside the actual section and the transform uses
$\nu|_{Y_R^*}$, not the normalized section probability.
The choice $a=s_R=c_*^{-1}\one_{Y_R^*}$ recovers that probability.
Put
\[
 \Phi^{[k],a}_{n,R}(z)=\int_{Y_R^*}a((F_R^*)^kx)e^{iz\cdot J_{n,R}(x)}\dd\nu(x).
\]
Its absolute value equals that of $C^{[k],a}_{n,R}(\sqrt n\,z)$.
The insertion may depend on $n,R,k$, but is fixed as frequency varies.

\subsection{A deterministic marked estimate with damping}
\begin{lemma}[Damping before the two-sided stopping]
\label{lem:fixed-count-damped-comparison}
Fix $Q\ge3$. Suppose $\delta\le\delta_0$ and every frequency under
consideration obeys \eqref{eq:finite-jet-damping-domain} with
$z=v/\sqrt n$. Uniformly in $0\le k\le n$,
\begin{align}
 |C^{[k],a}_{n,R}(v)|
 \le {}& CM_a\delta^{-2}e^{-c|v|^2}+C\delta V_a\notag\\
 &+CM_a\left\{|v|\sqrt{\delta(1+|\log\delta|)}
              +(1+|v|)n^{-2/9}\right\}.
 \label{eq:fixed-count-damped-comparison}
\end{align}
\end{lemma}
\begin{proof}
Use the exact recentering at the marked return and then
\eqref{eq:improved-marked-stopping}. This replaces the actual centered
record by $\mathsf H_{m_k,m_{n-k},R}$, at cost
$CM_a(1+|v|)n^{-2/9}$. Its deterministic collision length is
$m_k+m_{n-k}=n/c_*+O(1)$, uniformly in the mark.

Smooth the insertion at scale $\delta$, at $L^1$ cost $C\delta V_a$.
The mean-zero residual $h_R-h_{R,\delta}$ has interval variance at
most $Cn\delta(1+|\log\delta|)$ by
\eqref{eq:small-bv-variance}. Stationarity applies to the whole
interval, including its negative part. Since $|a_\delta|\le M_a$,
Cauchy--Schwarz bounds the change of observable by
$CM_a|v|\sqrt{\delta(1+|\log\delta|)}$.

The remaining integral is exactly
\[
 \ell\left(\mathcal L_{R,\delta,z}^{m_{n-k}}
          M_{a_\delta}\mathcal L_{R,\delta,z}^{m_k}\nu\right)
\]
by \eqref{eq:marked-chronological-pairing}. Each power now has the
real-frequency damping of Corollary~\ref{cor:finite-jet-damping}.
The multiplier norm is at most $CM_a\delta^{-2}$. Their product is
therefore bounded by
$CM_a\delta^{-2}e^{-c(m_k+m_{n-k})|z|^2}$, which is the first term
of \eqref{eq:fixed-count-damped-comparison}. The power bound includes
length zero, so the initial and terminal marks require no exceptional
argument. No return-count average, induced spectral gap, or spatial
compression has been used.
\end{proof}

\subsection{A feasible fixed-count scale}
\begin{theorem}[Raw-scale cancellation on a fixed-return inner annulus]
\label{thm:fixed-return-inner-annulus}
For all sufficiently large $n$, let
\begin{equation}\label{eq:new-fixed-inner-annulus}
 \mathcal A_n^{\rm in}
 =\{z\in\R^4:2n^{-99/200}\le |z|\le2n^{-12/25}\}.
\end{equation}
Its rescaled radii are $2n^{1/200}$ and $2n^{1/50}$. Uniformly in
$R$, in $0\le k\le n$, and in the stated insertions,
\begin{equation}\label{eq:fixed-inner-annulus-bound}
 n^2\int_{\mathcal A_n^{\rm in}}
       |\Phi^{[k],a}_{n,R}(z)|\dd z
 \le C\left(M_an^{-1/200}\sqrt{\log(2+n)}
                         +V_an^{-13/100}\right).
\end{equation}
In particular the left side tends to zero for the original probability
at each prescribed $n$, uniformly in the physical radius.
\end{theorem}
\begin{proof}
Choose once and for all
\begin{equation}\label{eq:fixed-count-feasible-scales}
 Q=19,\qquad \delta=\tfrac14n^{-21/100},
 \qquad |v|\le2n^{1/50},\qquad z=v/\sqrt n.
\end{equation}
The two damping restrictions hold for large $n$. Indeed their powers
are
\[
 |z|\delta^{-2}\le Cn^{-3/50},\qquad
 |z|^{18}\delta^{-40}\le Cn^{-6/25}.
\]
The constants, including those from $\delta= n^{-21/100}/4$, are
fixed before $n$ tends to infinity. The smoothed covariance is
uniformly positive definite at this scale.

Rescaling $v=\sqrt n\,z$ cancels the raw factor $n^2$ exactly.
Integrate \eqref{eq:fixed-count-damped-comparison} between the stated
rescaled radii. The Gaussian damping term is bounded by
$CM_an^{21/50}e^{-c n^{1/100}}$ and hence by every fixed negative
power of $n$ times $M_a$. The three other integrated losses are
\[
 CV_a n^{-21/100+4/50},\qquad
 CM_a n^{-21/200+5/50}\sqrt{\log(2+n)},\qquad
 CM_a n^{-2/9+5/50}.
\]
Their decay margins are respectively $13/100$, $1/200$, and
$11/90$. This proves \eqref{eq:fixed-inner-annulus-bound}.
The four-dimensional volume factor is included in each term. Taking
$a=s_R$ gives uniformly bounded $M_a,V_a$ and unit mass.
\end{proof}

\begin{corollary}[A larger integrated band at the prescribed count]
\label{cor:fixed-count-expanded-band}
Under the same conventions,
\begin{equation}\label{eq:fixed-count-expanded-band}
 \int_{|v|\le2n^{1/50}}
 \left|C^{[k],a}_{n,R}(v)-\alpha_a e^{-v^{\mathsf T}D_Rv/2}\right|\dd v
 \le C\left(M_an^{-1/200}\sqrt{\log(2+n)}+V_an^{-9/175}\right).
\end{equation}
This is an additional wider-band estimate. The sharper rate
$n^{-3/280}$ on the original $2n^{1/200}$ band in Theorem C is
retained unchanged.
\end{corollary}
\begin{proof}
On the original central ball apply
Theorem~\ref{thm:marked-return-band}. On its complement up to the new
radius use Theorem~\ref{thm:fixed-return-inner-annulus} and the
uniformly elliptic Gaussian tail. Here $|\alpha_a|\le M_a$,
$3/280>1/200$, and $13/100>9/175$.
\end{proof}

\begin{proposition}[The range of the finite-order construction]
\label{prop:finite-order-band-range}
Let
\[
 \frac1{200}<\varepsilon<\frac1{42},\qquad
 10\varepsilon<\theta<\frac14-\frac\varepsilon2,
\]
and choose a fixed integer $Q\ge3$ with
\[
 (Q-1)(1/2-\varepsilon)>2(Q+1)\theta.
\]
The inner-annulus assertion holds with outer physical radius
$2n^{-1/2+\varepsilon}$ and bound
\begin{equation}\label{eq:general-fixed-inner-annulus}
 C\left\{V_a n^{-(\theta-4\varepsilon)}
    +M_a n^{-(\theta/2-5\varepsilon)}\sqrt{\log(2+n)}
    +M_a n^{-(2/9-5\varepsilon)}\right\}.
\end{equation}
The rescaled outer radius is $2n^\varepsilon$. Every exponent in this
bound is positive, and all choices are independent of $n,R,k,a$.
\end{proposition}
\begin{proof}
Set $\delta=n^{-\theta}/4$. The strict upper bound on $\theta$
puts $z$ inside the analytic disk; the inequality for $Q$ is exactly
the second damping condition, with a strictly decaying power of $n$.
Such a fixed $Q$ exists because $1/2-\varepsilon>2\theta$.
Integration proceeds as above. The interval for $\theta$ is nonempty
because $\varepsilon<1/42$, and $\theta>10\varepsilon$ makes the
first two margins positive. Also $2/9-5\varepsilon>0$ in this range.
This is a sufficient range for the present proof, not a barrier for
other methods or for the original raw theorem.
\end{proof}

\subsection{Unchanged marked events and raw inversion}
\begin{corollary}[The same marked event on the wider band]
\label{cor:expanded-same-event-band}
Let $E_{n,R}\subset Y_R^*$ have zero-extended indicator of variation
$V_{n,R}$ and probability $p_{n,R}=\nu_R^*(E_{n,R})>0$.
For the unchanged event
$A_{n,k,R}=\{x:(F_R^*)^kx\in E_{n,R}\}$,
\begin{align}
 &\int_{|v|\le2n^{1/50}}
 \left|\mathbb E_{\nu_R^*}
 [e^{iv\cdot U_{n,R}/\sqrt n}\mid A_{n,k,R}]
                         -e^{-v^{\mathsf T}D_Rv/2}\right|\dd v\notag\\
 &\hspace{14mm}\le \frac{C}{p_{n,R}}
       \left(n^{-1/200}\sqrt{\log(2+n)}+V_{n,R}n^{-9/175}\right).
 \label{eq:expanded-same-event-band}
\end{align}
Thus convergence follows, for example, if
$p_{n,R}\ge c n^{-\beta}$ and $V_{n,R}\le Cn^\kappa$, where
$\beta<1/200$ and $\beta+\kappa<9/175$.
\end{corollary}
\begin{proof}
Take $a=c_*^{-1}\one_{E_{n,R}}$ in
Corollary~\ref{cor:fixed-count-expanded-band}. Its mass is $p_{n,R}$;
invariance gives the same probability to $A_{n,k,R}$. Divide the
numerator estimate by precisely this probability. No event is replaced.
The factor $p^{-1}$ here belongs to a fixed-count absolute-error
estimate; it is not the squared-$L^2$ normalization of the averaged
theorem. Its allowed events may vary with $n$ subject to these budgets.
\end{proof}

For the original unweighted law, use the finite extraction and the
notation of Section~\ref{sec:count-localized-inversion}. To distinguish
the new cutoff from the inherited one, put
\begin{equation}\label{eq:expanded-inversion-cutoff}
 \widetilde B_n=n^{-12/25},\quad
 \widetilde\chi_n(\omega)=\chi(\omega/\widetilde B_n),\quad
 \widetilde K_n=\mathcal F^{-1}\widetilde\chi_n.
\end{equation}
The support has physical radius $2n^{-12/25}$ and rescaled radius
$2n^{1/50}$. Define $\widetilde{\mathcal E}_{n,R,A}^{L}$ and
$\widetilde{\mathcal C}_{n,R}^{L}(B)$ by the definitions preceding
Corollary~\ref{cor:finite-extraction-error-budget}, replacing $K_n$
and $\chi_n$ by $\widetilde K_n$ and $\widetilde\chi_n$.
In particular the local edge correction is recomputed with the new
kernel; it is not transferred from the old kernel without an estimate.

\begin{theorem}[Fixed-count raw inversion with the enlarged cutoff]
\label{thm:expanded-count-localized-inversion}
Let $L_n=\lceil\lambda n\rceil$, $\lambda>c_*^{-1}+1$.
On each central window $\mathcal W_{n,R,A}$, the exact identity
\begin{align}
 p_{n,R}={}&\widetilde K_n*\mu_{n,R}
 +(e^{L_n}_{n,R}-\widetilde K_n*E^{L_n}_{n,R})\notag\\
 &+\mathcal F^{-1}[(1-\widetilde\chi_n)\widehat Q^{L_n}_{n,R}]
                     -\widetilde K_n*\mathsf T^{L_n}_{n,R}
 \label{eq:expanded-count-localized-inversion}
\end{align}
holds almost everywhere for large $n$. For every $P>0$ and $B\ge1$,
\begin{align}
 &\mathop{\rm ess\,sup}_{\mathcal W_{n,R,A}}
          |n^2p_{n,R}(\ell,t)-\mathfrak g_R(y)|\notag\\
 &\quad\le Cn^{-1/200}\sqrt{\log(2+n)}+Ce^{-c n^{1/25}}
                       +C_{A,\lambda,P}n^{-P}\notag\\
 &\qquad+n^2\widetilde{\mathcal E}_{n,R,A}^{L_n}
       +(2\pi)^{-4}n^2\widetilde{\mathcal C}_{n,R}^{L_n}(B)
       +\frac{n^2}{\pi B}A_2(n,L_n,R,1).
 \label{eq:expanded-raw-error-budget}
\end{align}
As before, this is an extended-real inequality until the displayed
local-edge quantity is bounded. The remaining three raw terms retain
their full radius and count dependence.
\end{theorem}
\begin{proof}
The support identity and absolutely convergent finite residual inverse
used in Theorem~\ref{thm:finite-extraction-central-inversion} do not
depend on the particular cutoff. Splitting with $\widetilde\chi_n$
gives \eqref{eq:expanded-count-localized-inversion} exactly.
Corollary~\ref{cor:fixed-count-expanded-band}, with $a=s_R$, bounds
the normalized inverse of the cutoff full law by the stated polynomial
error. The omitted Gaussian frequencies have $|v|\ge n^{1/50}$,
so their integral is at most $Ce^{-c n^{1/25}}$.

The count distance from the central window to the high-count support
is at least $(\lambda-c_*^{-1}-1)n$. The Schwartz estimate for the
new kernel and the mass bound $\|\mathsf T^{L_n}\|_{\mathrm{TV}}\le1$
give, for every integer $M\ge1$,
\[
 n^2\sup_{\mathcal W_{n,R,A}}|\widetilde K_n*\mathsf T^{L_n}|
                \le C_{A,\lambda,M}n^{2/25-13M/25}.
\]
Choose $M$ large enough for the desired $P$. The local extracted term
is precisely the new $\widetilde{\mathcal E}$, not the old edge
correction. Split the residual inverse at roof frequency $B$ and use
\eqref{eq:finite-count-far-tail}; its inverse-transform constant is
again $1/\pi$. This proves \eqref{eq:expanded-raw-error-budget}.
\end{proof}

\paragraph{The remaining fixed-count region.}
The estimate \eqref{eq:fixed-inner-annulus-bound} is for the actual
full-law transform at the specified count, with the raw $n^2$
normalization and no division by annular volume. It is not deduced
from Theorem K, cyclic arc mass, or the fixed-test-frequency spectral
Cauchy limit. The part from the enlarged cutoff toward $n^{-2/5}$,
compact nonzero lattice frequencies, and growing roof frequencies
still require fixed-count control. The far-roof derivative budget and
the complete local extracted-edge correction in
\eqref{eq:expanded-raw-error-budget} are unchanged analytical tasks.
A first-variation estimate for collision observables does not bound
those inverse-coarea derivatives. Likewise, the same-event corollary
does not compare a completed-return event with an exact physical-time
or lattice observation. These distinctions leave the original joint
raw-density endpoint in place while moving its proved fixed-count
central boundary.

\section{Damped removal of smoothing and higher-moment stopping}
\label{sec:damped-unsmoothing}
The first-order comparison in
Lemma~\ref{lem:fixed-count-damped-comparison} takes the absolute value
of the entire change of observable before using spectral damping.
We instead retain a finite Taylor polynomial of that change inside
the chronological collision word. Every term of this polynomial
still has damping over the full deterministic collision interval.
Only its fourth-order remainder is estimated without damping.
A second, finer smoothing scale is used for the finitely many
insertions; it does not shrink the analytic domain of the twisted
collision powers. We also improve the actual stopping estimate by
using a fixed higher moment of the visit count.

All constants below are uniform in $R$. Their dependence on a fixed
moment order is allowed. We neither let that order grow with $n$
nor use an all-orders analytic expansion. Write
$\|u\|_{\mathcal V}=\|u\|_\infty+\|u\|_{\mathrm{BV}}$ and
$L_\delta=1+|\log\delta|$. Expectations without a subscript are with
respect to the stationary collision probability $\nu$.

\subsection{Cumulants of a small bounded-variation residual}
\begin{lemma}[Small-mass connected correlations]
\label{lem:small-mass-cumulants}
Fix $q\ge2$ and $B<\infty$. If $u$ is real,
$\|u\|_{\mathcal V}\le B$, and $\|u\|_1\le B\delta$ for
$0<\delta<1/2$, then
\begin{align}
 &\sum_{\mathbf j\in\Z^{q-1}}
 \left|\operatorname{cum}(u,u\circ T_R^{j_2},\ldots,
                                      u\circ T_R^{j_q})\right|
       \le C_{q,B}\delta L_\delta^{q-1},
       \label{eq:small-mass-connected}\\
 &|\operatorname{cum}_q(S_{m,R}u)|
       \le C_{q,B}m\delta L_\delta^{q-1}\qquad(m\ge1).
       \label{eq:small-mass-cumulant-length}
\end{align}
If in addition $\int u\dd\nu=0$, then
\begin{equation}\label{eq:small-residual-fourth}
 \mathbb E|S_{m,R}u|^4
       \le C_B\bigl(m^2\delta^2L_\delta^2
                          +m\delta L_\delta^3\bigr).
\end{equation}
The same fourth-moment bound holds on every deterministic integer
interval, also with its orientation reversed.
\end{lemma}
\begin{proof}
In each term of the partition formula
\eqref{eq:finite-order-cumulant-definition}, the block containing the
first factor has moment at most $B^{q-1}\|u\|_1$ in absolute value;
all other block moments are bounded. Thus the joint cumulant is at
most $C_{q,B}\delta$, independently of the times. The independent-block
argument in Proposition~\ref{prop:connected-collision-summability}
also bounds it by $C_{q,B}e^{-\gamma D}$, where
$D=\operatorname{diam}\{0,j_2,\ldots,j_q\}$ and $\gamma>0$
depends on the fixed order.

There are at most $(2J+1)^{q-1}$ anchored tuples with $D\le J$.
Choose $J$ to be a sufficiently large fixed multiple of $L_\delta$.
Inside this set use the bound $C\delta$; outside it sum the
exponential bound, absorbing the fixed-degree lattice counting
polynomial in half of the exponential. This gives
$C_{q,B}\delta L_\delta^{q-1}$. The exact anchored multiplicity
$(m-D)_+$ is at most $m$, proving
\eqref{eq:small-mass-cumulant-length} without a boundary error
independent of $\delta$.

For a centered real variable $X$, the finite moment--cumulant identity
is $\mathbb E X^4=\operatorname{cum}_4(X)+
3\operatorname{cum}_2(X)^2$. It follows either by expanding the
exponential of its cumulant generating polynomial through degree
four, or directly by the partition formula. Apply
\eqref{eq:small-mass-cumulant-length} at orders two and four.
Stationarity translates an arbitrary integer interval, and reversal
does not change its sum.
\end{proof}

\begin{lemma}[Every fixed even collision moment and its maximum]
\label{lem:fixed-even-maximal}
Fix an integer $p\ge2$. For a real centered $u$ with
$\|u\|_{\mathcal V}\le B$ and $m\ge1$,
\begin{equation}\label{eq:fixed-even-maximal}
 \mathbb E\left|\sum_{j=r}^{r+m-1}u\circ T_R^j\right|^{2p}
 +\mathbb E\max_{0\le l\le m}
       \left|\sum_{j=r}^{r+l-1}u\circ T_R^j\right|^{2p}
       \le C_{p,B}m^p.
\end{equation}
The location $r\in\Z$ is arbitrary. The estimate holds in reverse
order, for the vector $h_R$, and under a nonnegative initial density
bounded by a fixed constant.
\end{lemma}
\begin{proof}
By Proposition~\ref{prop:connected-collision-summability}, every
cumulant of $S_m u$ of order $2\le j\le2p$ is $O_{p,B}(m)$.
The coefficient identity
\[
 \mathbb E X^{2p}=\sum_{\pi\in\mathcal P_{2p}}
                         \prod_{B\in\pi}\operatorname{cum}_{|B|}(X)
\]
is obtained by exponentiating the generating polynomial through that
fixed degree. For centered $X$, partitions with singleton blocks
vanish. Every remaining partition has at most $p$ blocks. The
moment of $S_m u$ is consequently at most $C_{p,B}m^p$.
No convergence of an infinite cumulant series is used.

Enlarge $m$ to a power of two $M<2m$. Every prefix is a union of at
most one aligned dyadic interval at each scale. At interval length
$l=M2^{-j}$ the $2p$th moment of the maximum interval sum is at most
$C_{p,B}(M/l)l^p$. Minkowski's inequality over scales therefore gives
\[
 \left\|\max_{0\le l\le M}|S_lu|\right\|_{2p}
 \le C_{p,B}\sum_{j=0}^{\log_2 M}
                 (M^p2^{-j(p-1)})^{1/(2p)}
 \le C_{p,B}\sqrt M.
\]
The geometric series converges because $p>1$. Translation and
reversal use invariance, not a changed initial law. A bounded initial
density is handled by domination. Summing the four component bounds
proves the vector assertion.
\end{proof}

\subsection{Higher-moment windows for the genuine return clock}
Retain $N^\pm_{j,R}$, $m_j=\lfloor j/c_*\rfloor$ and
$\mathsf H_{r,s,R}$ from Section~\ref{sec:marked-return-band}.
As there, the zero-return clocks and empty sums are zero.

\begin{proposition}[Higher-moment marked stopping]
\label{prop:higher-moment-marked-stopping}
For each fixed integer $p\ge2$, all $j\ge0$ and $b\ge2$,
\begin{equation}\label{eq:higher-moment-clock}
 \nu_R^*\{|N^\pm_{j,R}-m_j|>b\}
             \le C_p\frac{(j+b)^p}{b^{2p}}.
\end{equation}
Put $s_p=p/(4p+1)$. For every $n\ge2$, $0\le k\le n$ and
$v\in\R^4$, the actual marked transform obeys
\begin{align}
 &\left|C^{[k],a}_{n,R}(v)
 -\int_M a(y)e^{iv\cdot\mathsf H_{m_k,m_{n-k},R}(y)/\sqrt n}
                       \dd\nu(y)\right|\notag\\
 &\hspace{18mm}\le C_pM_a(1+|v|)n^{-s_p}.
 \label{eq:higher-moment-marked-stopping}
\end{align}
No bound on $V_a$ is needed for this stopping comparison.
\end{proposition}
\begin{proof}
The centered visit count
$\sum_{i=1}^L(\eta_R-c_*)\circ T_R^{\pm i}$ has $2p$th moment at
most $C_pL^p$ under $\nu_R^*$, by
Lemma~\ref{lem:fixed-even-maximal} and $\nu_R^*\le c_*^{-1}\nu$.
The exact visit equivalences in
Lemma~\ref{lem:two-sided-clock-window}, followed by Markov's
inequality at integer thresholds within a bounded distance of
$m_j\pm b$, give \eqref{eq:higher-moment-clock}. The required
deviation is at least $c_*b-O(1)$ and $L\le C(j+b)$.
A negative lower threshold yields an empty event; $j=0$ is immediate.
Enlarge the constant for bounded $b$.

Recenter at the actual marked return using
\eqref{eq:marked-recentering}. For $2\le b\le n$, the union of the
two clock exceptions has probability at most $C_pn^p/b^{2p}$.
Off this union, the difference from the deterministic interval is
bounded by maxima on two fixed collision windows of length $O(b)$.
Lemma~\ref{lem:fixed-even-maximal} bounds their second moments by
$C_pb$, including windows crossing collision time zero. On the
exceptional set use only the bound two for a difference of phases.
Consequently the error is at most
\[
 C_pM_a\left(\frac{n^p}{b^{2p}}+|v|\sqrt{\frac bn}\right).
\]
Taking $b=\lceil n^{(2p+1)/(4p+1)}\rceil$ gives
\eqref{eq:higher-moment-marked-stopping}. Finitely many small $n$
are absorbed in the constant. The two clocks are not assumed
independent, and the weight is never evaluated at a deterministic
surrogate for the marked return.
\end{proof}

\subsection{A cubic correction retained inside the damped word}
Let $r,s\ge0$, $m=r+s\ge1$, and define
\[
 I_{r,s,R}^a(z)=\int_M a(y)
                         e^{iz\cdot\mathsf H_{r,s,R}(y)}\dd\nu(y).
\]
The insertion need not be positive. The next result uses the coarse
scale $\delta$ only for $h_{R,\delta}$ and the fine scale $\epsilon$
only for the insertion and a finite number of residual factors.

\begin{theorem}[Damped finite-order removal of smoothing]
\label{thm:damped-cubic-unsmoothing}
Fix $Q\ge3$, and suppose $\delta\le\delta_0$ and real $z$ satisfy
\eqref{eq:finite-jet-damping-domain}. For every $0<\epsilon<1/2$,
\begin{align}
 |I_{r,s,R}^a(z)|\le{}&
 CM_a\epsilon^{-8}(1+m|z|)^3e^{-cm|z|^2}+C\epsilon V_a\notag\\
 &+CM_a\epsilon(1+m|z|)^3\notag\\
 &+CM_a|z|^4\bigl(m^2\delta^2L_\delta^2
                              +m\delta L_\delta^3\bigr).
 \label{eq:damped-cubic-unsmoothing}
\end{align}
The constants are independent of the relative location of the mark
and of the fine scale. No analytic perturbation in the fine-scale
observable is asserted.
\end{theorem}
\begin{proof}
The assertion at $z=0$ follows by increasing $C$, so take $z\ne0$.
Put $\xi=z/|z|$, $u=\xi\cdot(h_R-h_{R,\delta})$,
$a_\epsilon=\mathsf S_\epsilon a$ and
$u_\epsilon=\mathsf S_\epsilon u$. Mean-preserving smoothing gives
\[
 \int u\dd\nu=0,\quad \|u\|_{\mathcal V}\le C,\quad
 \|u\|_1\le C\delta,\quad
 \|u-u_\epsilon\|_1\le C\epsilon.
\]
Also $\|a-a_\epsilon\|_1\le C\epsilon V_a$ and
$\|a_\epsilon\|_\infty\le M_a$. Replace the insertion first,
while the full exponential still has modulus one. This costs exactly
the second term of \eqref{eq:damped-cubic-unsmoothing}.

Translate the interval to collision times $0,\ldots,m-1$. Its mark
is now at time $r\in[0,m]$. With $X=|z|S_m u$, use the real-variable
Taylor formula
\begin{equation}\label{eq:cubic-phase-remainder}
 e^{iX}=\sum_{j=0}^3\frac{(iX)^j}{j!}+\mathcal R_4(X),
 \qquad |\mathcal R_4(X)|\le |X|^4/24.
\end{equation}
Multiplication by the coarse phase
$e^{iz\cdot S_m h_{R,\delta}}$ does not change this bound.
Lemma~\ref{lem:small-mass-cumulants} therefore bounds the remainder
integral by the last term of \eqref{eq:damped-cubic-unsmoothing}.
This estimate is for the original residual $u$, not for a residual
with a presumed independent law.

For $1\le j\le3$, expand $(S_m u)^j$ as a sum over its $m^j$
ordered index tuples, allowing repetitions. Within each product,
replace each factor $u\circ T_R^t$ by $u_\epsilon\circ T_R^t$.
The other factors and the coarse phase are bounded. Telescoping and
invariance bound its integrated error by $C_jM_a\epsilon$.
After multiplying by $|z|^j/j!$ and summing, this costs at most
$CM_a\epsilon(1+m|z|)^3$. No long pullback BV norm is used.

It remains to bound each corrected product. Write
$L_z=\mathcal L_{R,\delta,z}$. Sort its $j$ residual insertion times
together with the mark, retaining repeated times, as
$0\le t_1\le\cdots\le t_d\le m$, where $d=j+1\le4$.
Exactly one $b_i$ is $a_\epsilon$; the others are $u_\epsilon$.
The chronological word is
\begin{equation}\label{eq:corrected-chronological-word}
 \ell\!\left(
 L_z^{m-t_d}M_{b_d}L_z^{t_d-t_{d-1}}M_{b_{d-1}}
 \cdots L_z^{t_2-t_1}M_{b_1}L_z^{t_1}\nu\right).
\end{equation}
For $d=1$ this means
$\ell(L_z^{m-t_1}M_{b_1}L_z^{t_1}\nu)$.
The nonnegative power lengths sum to $m$. Smooth multiplication costs
at most $CM_a\epsilon^{-2}$ for the mark and $C\epsilon^{-2}$
for each residual. Corollary~\ref{cor:finite-jet-damping} thus bounds
\eqref{eq:corrected-chronological-word} by
$C_jM_a\epsilon^{-2(j+1)}e^{-cm|z|^2}$. Summing the finitely
many degrees and all their index tuples proves the first term of
\eqref{eq:damped-cubic-unsmoothing}.

A repeated time is a zero-length power, and a mark at $0$ or $m$
is an endpoint multiplier; both are covered by the same power bound.
There is never a perturbation of $L_z$ by $u_\epsilon$. In particular,
the fine scale does not enter the damping-domain restrictions.
\end{proof}

\paragraph{Why the new estimate has a different scale balance.}
The undamped observable loss in the earlier comparison was proportional
to $|z|\sqrt{m\delta L_\delta}$. In
\eqref{eq:damped-cubic-unsmoothing}, the lower-degree changes remain
inside a word damped for all $m$ collisions; its undamped residual
begins with $|z|^4m^2\delta^2L_\delta^2$. The polynomial cost of
the much finer insertions is harmless on an annulus with a growing
rescaled inner radius. This argument concerns the characteristic
function itself at a deterministic collision length. The passage to a
specified actual return count uses
\eqref{eq:higher-moment-marked-stopping}, not an average over counts.

\section{A wider prescribed-count band and its raw error budget}
\label{sec:wider-fixed-count-band}
We combine the damped removal of smoothing with higher-moment stopping.
The resulting estimate reaches beyond the sufficient exponent range
in Proposition~\ref{prop:finite-order-band-range}. Both the raw
normalization and the single actual-return insertion are retained.
The unchanged-event consequence and a weighted exact inversion below
use the same measures, with no normalization of a truncated law.

Use $M_a,V_a,\alpha_a,C^{[k],a}_{n,R}$ and
$\Phi^{[k],a}_{n,R}$ as in Section~\ref{sec:fixed-return-annulus}.
The insertion is fixed while frequency varies. It is integrated
against $\nu|_{Y_R^*}$; taking $a=s_R$ gives the original probability.

\subsection{An explicit band beyond the first-order range}
\begin{theorem}[A wider raw-scale annulus at every prescribed return count]
\label{thm:wider-fixed-count-annulus}
For every sufficiently large $n$, put
\begin{equation}\label{eq:wider-fixed-count-annulus}
 \mathcal A_n^{\sharp}
 =\{z\in\R^4:2n^{-99/200}\le |z|\le2n^{-19/42}\}.
\end{equation}
The corresponding rescaled radii are $2n^{1/200}$ and $2n^{1/21}$.
Uniformly in $R$, $0\le k\le n$, and the stated complex insertions,
\begin{equation}\label{eq:wider-fixed-count-bound}
 n^2\int_{\mathcal A_n^{\sharp}}
           |\Phi^{[k],a}_{n,R}(z)|\dd z
       \le C\left(M_an^{-5/861}+V_an^{-59/21}\right).
\end{equation}
There is no average over return counts and no division by annular volume.
\end{theorem}
\begin{proof}
Fix the orders and scales
\begin{equation}\label{eq:wider-fixed-scales}
 Q=19,\quad p=10,\quad
 \delta=\tfrac14n^{-1/5},\quad
 \epsilon=\tfrac14n^{-3},\quad |v|\le2n^{1/21},\quad z=v/\sqrt n.
\end{equation}
Here $Q$ is the degree of the inherited spectral jet, $2p=20$ is
only the moment order used for the clock, and the polynomial in the
residual phase has degree three. These three orders have different roles.
The two analytic damping quantities satisfy
\[
 |z|\delta^{-2}\le Cn^{-11/210},\qquad
 |z|^{18}\delta^{-40}\le Cn^{-1/7}.
\]
Thus Corollary~\ref{cor:finite-jet-damping} applies for large $n$,
uniformly in the radius. Increasing the spectral degree alone would
not have improved the first-order unsmoothing loss.

Apply Proposition~\ref{prop:higher-moment-marked-stopping} with
$p=10$. The deterministic interval has length
$m=m_k+m_{n-k}=n/c_*+O(1)$, so $m\asymp n$ uniformly in the mark.
Use Theorem~\ref{thm:damped-cubic-unsmoothing} on that interval.
Rescale $v=\sqrt n\,z$; in four dimensions the Jacobian $n^{-2}$
cancels the factor $n^2$ on the left of
\eqref{eq:wider-fixed-count-bound}.

The first term of \eqref{eq:damped-cubic-unsmoothing} has integral
at most $CM_an^{26}e^{-c n^{1/100}}$, after decreasing $c$.
Indeed, $\epsilon^{-8}=O(n^{24})$ and
$(1+m|z|)^3\le Cn^{3/2}(1+|v|)^3$, while the lower radius is
$2n^{1/200}$. This term is bounded by every fixed inverse power
of $n$ times $M_a$. The insertion and finite-product replacement
errors have integrated bounds
\[
 CV_a n^{-3+4/21}=CV_a n^{-59/21},\qquad
 CM_a n^{-3+3/2+7/21}=CM_a n^{-7/6}.
\]
The fourth-order residual terms have bounds
\[
 CM_a n^{-2/5+8/21}[\log(2+n)]^2
   +CM_a n^{-6/5+8/21}[\log(2+n)]^3.
\]
Their power margins are $2/105$ and $86/105$. Finally the integrated
stopping error is at most
\[
 CM_a n^{-10/41+5/21}=CM_a n^{-5/861};
\]
the term without $|v|$ is smaller. Since $2/105>5/861$, the logarithms
in the fourth-order residual terms are absorbed by the strict power
slack. This proves \eqref{eq:wider-fixed-count-bound} with the
separate variation loss as stated. All orders and constants are fixed
before $n,R,k,a$ are varied.
\end{proof}

\begin{corollary}[Integrated Gaussian comparison on the wider band]
\label{cor:wider-marked-central-band}
Under the same conventions,
\begin{align}
 &\int_{|v|\le2n^{1/21}}
 \left|C^{[k],a}_{n,R}(v)-\alpha_a e^{-v^{\mathsf T}D_Rv/2}\right|\dd v
 \notag\\
 &\hspace{18mm}\le C\left(M_an^{-5/861}+V_an^{-9/175}\right).
 \label{eq:wider-marked-central-band}
\end{align}
The physical support radius is $2n^{-19/42}$. The sharper inherited
rate on $|v|\le2n^{1/200}$ remains a separate valid estimate.
\end{corollary}
\begin{proof}
On the original central ball apply
Theorem~\ref{thm:marked-return-band}. On its complement up to the
new radius apply Theorem~\ref{thm:wider-fixed-count-annulus} and the
uniformly elliptic Gaussian tail. Use $|\alpha_a|\le M_a$,
$3/280>5/861$, and $59/21>9/175$.
\end{proof}

\begin{proposition}[A fixed-order family of wider bands]
\label{prop:damped-unsmoothing-band-range}
Fix $1/200<\varepsilon<1/20$. Choose
\begin{equation}\label{eq:damped-unsmoothing-parameter-range}
 4\varepsilon<\theta<\frac14-\frac\varepsilon2,
 \qquad \frac{p}{4p+1}>5\varepsilon,
 \qquad (Q-1)(1/2-\varepsilon)>2(Q+1)\theta,
\end{equation}
with fixed integers $p\ge2$, $Q\ge3$. For every
\[
 0<\eta<\min\{2\theta-8\varepsilon,\ p/(4p+1)-5\varepsilon\},
\]
the annulus from $2n^{-99/200}$ to $2n^{-1/2+\varepsilon}$ satisfies
\begin{equation}\label{eq:damped-unsmoothing-general-band}
 n^2\int |\Phi^{[k],a}_{n,R}(z)|\dd z
       \le C_{\varepsilon,\theta,p,Q,\eta}
                    \left(M_an^{-\eta}+V_an^{-(3-4\varepsilon)}\right).
\end{equation}
Its rescaled outer radius is $2n^\varepsilon$. Every choice in
\eqref{eq:damped-unsmoothing-parameter-range} is feasible, and all
orders remain fixed as $n$ tends to infinity.
\end{proposition}
\begin{proof}
Choose $\delta=n^{-\theta}/4$ and $\epsilon=n^{-3}/4$.
The strict upper bound on $\theta$ leaves a nonzero analytic-disk
margin, and a sufficiently large fixed $Q$ gives the second damping
margin. The interval for $\theta$ is nonempty for
$\varepsilon<1/18$, hence for the stated range. Since
$p/(4p+1)\uparrow1/4$, a fixed $p$ with the required stopping margin
exists whenever $\varepsilon<1/20$.

The two principal integrated losses are
$n^{-2\theta+8\varepsilon}(\log n)^2$ and
$n^{-p/(4p+1)+5\varepsilon}$. The lower residual term has margin
$1+\theta-8\varepsilon>2\theta-8\varepsilon$, because
$\theta<1$. The fine-scale replacement has margin
$3/2-7\varepsilon$, and the insertion error has margin
$3-4\varepsilon$. Both exceed the required positive loss margins
where needed. Damping absorbs all fixed polynomial prefactors.
Absorb the logarithms by the strict choice of $\eta$ and repeat the
raw rescaling. No constant is asserted uniform as an order grows.
\end{proof}

\paragraph{Comparison of frequency scales.}
The earlier sufficient range was $\varepsilon<1/42$ (approximately
$0.02381$); the new one is $\varepsilon<1/20=0.05$.
The explicit choice $1/21$ is strictly inside the new range.
Neither sufficient range reaches $1/10$, the exponent corresponding
to physical radius $n^{-2/5}$. The remaining interval beyond the new
cutoff is not included by a change of notation. No one of these
sufficient ranges is an impossibility assertion for a different method.

\subsection{The unchanged event and a weighted raw identity}
\begin{corollary}[An unchanged marked-state event on the new band]
\label{cor:wider-same-event-band}
Let $E_{n,R}\subset Y_R^*$ have zero-extended indicator of variation
$V_{n,R}$ and probability $p_{n,R}>0$. For the same event
$A_{n,k,R}=\{x:(F_R^*)^kx\in E_{n,R}\}$ in numerator and denominator,
\begin{align}
 &\int_{|v|\le2n^{1/21}}
 \left|\mathbb E_{\nu_R^*}
       [e^{iv\cdot U_{n,R}/\sqrt n}\mid A_{n,k,R}]
                         -e^{-v^{\mathsf T}D_Rv/2}\right|\dd v\notag\\
 &\hspace{18mm}\le\frac C{p_{n,R}}
                    \left(n^{-5/861}+V_{n,R}n^{-9/175}\right).
 \label{eq:wider-same-event-band}
\end{align}
In particular, $p_{n,R}\ge cn^{-\beta}$ and
$V_{n,R}\le Cn^\kappa$ suffice for convergence if
$\beta<5/861$ and $\beta+\kappa<9/175$.
\end{corollary}
\begin{proof}
Set $a=c_*^{-1}\one_{E_{n,R}}$ in
Corollary~\ref{cor:wider-marked-central-band}. Its insertion mass is
the exact probability $p_{n,R}$, which is also the probability of
$A_{n,k,R}$ by invariance. Divide by that same number. This proves
the claim without a replacement of the conditioning event.
\end{proof}

Put
\begin{equation}\label{eq:sharp-inversion-cutoff}
 B_n^\sharp=n^{-19/42},\quad
 \chi_n^\sharp(\omega)=\chi(\omega/B_n^\sharp),\quad
 K_n^\sharp=\mathcal F^{-1}\chi_n^\sharp.
\end{equation}
The support radii are $2n^{-19/42}$ physically and $2n^{1/21}$
after rescaling. For a marked insertion define
\[
 w_{n,k,R}(x)=c_*a((F_R^*)^kx).
\]
Then $(J_{n,R})_*(w_{n,k,R}\nu_R^*)$ is exactly
$\mu_{n,R}^{[k],a}$ of \eqref{eq:marked-measure}, with mass
$\alpha_a$, not a renormalized probability.

For the next theorem require, in addition to the BV assumptions,
that this weight be admissible in
Definition~\ref{def:finite-record-weight} at the count cutoffs in use.
For example a bounded insertion subanalytic on each collision chart
has this property, by the finite collision graph and exact return
membership. This additional requirement is not imposed on the preceding
Fourier theorems. In particular arbitrary BV and finite-record
subanalytic classes are not being identified.

Let $\mu^{a,L}$ be this measure restricted to $N_{n,R}\le L$,
$\mathsf T^{a,L}=\mu^{[k],a}-\mu^{a,L}$, and let $E^{a,L},Q^{a,L}$
be its exact finite extraction with densities $e^{a,L},r^{a,L}$.
Suppress $n,k,R$ only in these formulas. Define anew
\begin{align}
 \mathcal E^{\sharp,a,L}_{n,k,R,A}
 &=\mathop{\rm ess\,sup}_{\mathcal W_{n,R,A}}
            |e^{a,L}-K_n^\sharp*E^{a,L}|,\notag\\
 \mathcal C^{\sharp,a,L}_{n,k,R}(B)
 &=\int_{\T^3\times[-B,B]}
          |1-\chi_n^\sharp(\omega)|\,|\widehat Q^{a,L}(\omega)|\dd\omega.
 \label{eq:sharp-weighted-raw-terms}
\end{align}
The edge correction is kernel-dependent, and may be infinite before
its density germs are estimated.

\begin{theorem}[Exact weighted raw error at the new cutoff]
\label{thm:sharp-weighted-raw-inversion}
Let $L_n=\lceil\lambda n\rceil$, $\lambda>c_*^{-1}+1$.
For every admissible marked insertion as above, the mixed density
$p_{n,R}^{[k],a}$ obeys, almost everywhere on $\mathcal W_{n,R,A}$,
\begin{align}
 p_{n,R}^{[k],a}={}&K_n^\sharp*\mu_{n,R}^{[k],a}
      +(e^{a,L_n}-K_n^\sharp*E^{a,L_n})\notag\\
 &+\mathcal F^{-1}[(1-\chi_n^\sharp)\widehat Q^{a,L_n}]
                -K_n^\sharp*\mathsf T^{a,L_n}.
 \label{eq:sharp-weighted-raw-inversion}
\end{align}
Writing $y=((\ell,t)-n\bar G_R)/\sqrt n$, for $B\ge1$ and $P>0$,
\begin{align}
 &\mathop{\rm ess\,sup}_{\mathcal W_{n,R,A}}
       |n^2p_{n,R}^{[k],a}(\ell,t)-\alpha_a\mathfrak g_R(y)|\notag\\
 &\quad\le C(M_an^{-5/861}+V_an^{-9/175})
           +CM_ae^{-c n^{2/21}}+C_{A,\lambda,P}M_an^{-P}\notag\\
 &\qquad+n^2\mathcal E^{\sharp,a,L_n}_{n,k,R,A}
       +(2\pi)^{-4}n^2\mathcal C^{\sharp,a,L_n}_{n,k,R}(B)
       +\frac{n^2}{\pi B}A_2(n,L_n,R,w_{n,k,R}).
 \label{eq:sharp-weighted-raw-budget}
\end{align}
The inequality is interpreted in the extended reals until the local
edge correction is bounded. The original unweighted law is the case
$a=s_R$, for which $w_{n,k,R}=1$.
\end{theorem}
\begin{proof}
The full marked measure is dominated in total variation by
$c_*M_a\mu_{n,R}$, so it has a mixed density. Its restriction to a
count fiber at most $L_n$ agrees exactly with that of the full
measure. Apply Theorem~\ref{thm:finite-count-raw-extraction} to the
admissible $w_{n,k,R}$; its residual Fourier inverse is absolutely
convergent. Splitting this inverse by $\chi_n^\sharp$ and substituting
$\mu^{a,L_n}=\mu^{[k],a}-\mathsf T^{a,L_n}$ proves
\eqref{eq:sharp-weighted-raw-inversion}, with the same exact support
argument as Lemma~\ref{lem:central-count-support}.

Corollary~\ref{cor:wider-marked-central-band}, after rescaling in
the cutoff inverse transform, bounds the first term by the first
error in \eqref{eq:sharp-weighted-raw-budget}. The missing Gaussian
frequencies satisfy $|v|\ge n^{1/21}$ and contribute at most
$CM_ae^{-c n^{2/21}}$. For every integer $J\ge1$, count-coordinate
separation and the Schwartz bound on the new kernel give
\[
 n^2\sup_{\mathcal W_{n,R,A}}|K_n^\sharp*\mathsf T^{a,L_n}|
          \le C_{A,\lambda,J}M_a n^{4/21-23J/42}.
\]
Here $n^2(B_n^\sharp)^4=n^{4/21}$ and
$nB_n^\sharp=n^{23/42}$; the mass bound is
$\|\mathsf T^{a,L_n}\|_{\mathrm{TV}}\le c_*M_a$.
Choose $J$ for the desired $P$. Finally split the residual inverse
at roof frequency $B$. The interior integral is exactly
$\mathcal C^{\sharp,a,L_n}_{n,k,R}(B)$; the exterior estimate in
\eqref{eq:finite-count-far-tail}, multiplied by $(2\pi)^{-4}$,
gives $A_2/(\pi B)$. This proves the inequality with every weighted
raw term retained.
\end{proof}

\paragraph{The remaining raw estimates.}
The new contribution is a fixed-count full-law estimate over a wider
portion of the complement, not a function-defect or averaged-return
substitute. The finite spectral jet, quartic residual estimate and
higher visit moments concern bounded collision observables. They do
not estimate $A_2$ or the local extracted density germs at a linear
count cutoff. Likewise, the finite-band residual integral beyond the
new cutoff still includes the farther small annulus, compact nonzero
and peripheral return frequencies, and growing roof frequencies.
These three raw terms must be controlled together in a compatible
splice. The weighted identity specifies the common BV and subanalytic
class for which both the new central theorem and the complete finite
extraction apply. It does not place an arbitrary multiple-return or
exact physical-time event in that class, nor compare two rare events.

\section{Multiscale stopping and all fixed marked moments}
\label{sec:multiscale-stopping}
A single clock window assigns the same exceptional-set cost to every
larger deviation. Summing instead over the actual deviation scales
recovers a fourth-root comparison. This estimate is uniform in the
location of a prescribed return mark; it does not take a maximum over
all marks inside the expectation. The stronger comparison will also
remove the loss of central rate on the wider Fourier band.

Retain the notation of \eqref{eq:moment-stopped-pair}:
\[
 Z_{n,k,R}=\mathsf H_{N^-_{k,R},N^+_{n-k,R},R},\qquad
 W_{n,k,R}=\mathsf H_{m_k,m_{n-k},R},\qquad
 m_j=\lfloor j/c_*\rfloor.
\]
In particular, recentering the orbit at its actual marked return gives
$Z_{n,k,R}(y)=U_{n,R}((F_R^*)^{-k}y)$, where
$U_{n,R}=J_{n,R}-n\bar G_R$. All norms in the stopping argument are
under the normalized section probability $\nu_R^*$. Constants may
depend on a fixed moment order, but not on $R,n,k$.

\subsection{Deviation shells rather than one exceptional set}
\begin{lemma}[Two bounds for the total clock deviation]
\label{lem:total-clock-deviation}
Set
\[
 D_{n,k,R}=|N^-_{k,R}-m_k|+|N^+_{n-k,R}-m_{n-k}|.
\]
For every fixed integer $p\ge2$ there are constants $C_p,a_0,c_0>0$
such that, for $n\ge2$ and $b\ge2$,
\begin{align}
 \nu_R^*(D_{n,k,R}>b)&\le C_p\frac{(n+b)^p}{b^{2p}},
       \label{eq:total-clock-polynomial}\\
 \nu_R^*(D_{n,k,R}>b)&\le C_pe^{a_0n-c_0b}.
       \label{eq:total-clock-exponential}
\end{align}
The second bound may exceed one. It will only be used beyond a fixed
multiple of $n$.
\end{lemma}
\begin{proof}
A deviation greater than $b$ requires one of the two clocks to deviate
by more than $b/2$. Equation~\eqref{eq:higher-moment-clock} and a
union bound prove \eqref{eq:total-clock-polynomial}, with a change of
constant for bounded $b$. This argument includes $k=0,n$; a zero-return
clock is identically zero.

Put $Q=N^-_{k,R}+N^+_{n-k,R}$. Under $\nu_R^*$, a backward clock of
$j$ returns has the same law as a forward clock of $j$ returns, since
$N^-_{j,R}((F_R^*)^jx)=N^+_{j,R}(x)$ and the actual return map preserves
that probability. The cumulative-return estimate therefore gives
\[
 \nu_R^*(Q>L)\le C e^{an-cL/2}\qquad(L\ge0)
\]
by a union bound, with no independence of the two clocks. Since
$D_{n,k,R}\le Q+m_k+m_{n-k}\le Q+c_*^{-1}n$, translating the threshold
proves \eqref{eq:total-clock-exponential}. For negative translated
thresholds, increase $a_0$ and $C_p$ so that the displayed right side
is at least one.
\end{proof}

\begin{theorem}[Fourth-root stopping in every fixed finite moment]
\label{thm:multiscale-marked-stopping}
For every fixed real $q\ge1$ there is $C_q<\infty$ such that
\begin{align}
 \|Z_{n,k,R}\|_q+\|W_{n,k,R}\|_q&\le C_q\sqrt n,
       \label{eq:all-stopped-moment-bound}\\
 \|Z_{n,k,R}-W_{n,k,R}\|_q&\le C_q n^{1/4}
       \label{eq:quarter-stopping}
\end{align}
for all $n\ge2$, $0\le k\le n$, and $R\in I$. Consequently, for every
bounded insertion $a$ supported on the section and every real $v$,
\begin{align}
 &\left|C^{[k],a}_{n,R}(v)
   -\int_M a(y)e^{iv\cdot W_{n,k,R}(y)/\sqrt n}\dd\nu(y)\right|
       \le CM_a|v|n^{-1/4}.
       \label{eq:quarter-characteristic-stopping}
\end{align}
For $q=2$, the proof of \eqref{eq:quarter-stopping} uses only the
fourth collision moments, the fourth-moment clock window, and the
cumulative-return tail.
\end{theorem}
\begin{proof}
Choose a fixed integer $p\ge2$ such that $2p>3q/2$. The deterministic
interval has length at most $c_*^{-1}n$. Its $q$th moment is bounded
by Lemma~\ref{lem:fixed-even-maximal} at a sufficiently large fixed
even order, followed by domination of the initial density.
For $Z$, split at $Q\le An$, with $Q$ as in the preceding proof and
$A$ a sufficiently large fixed constant. On this event, $|Z|$ is at
most the sum of two deterministic collision maxima of length
$\lceil An\rceil$. Their $q$th moments are $O_q(n^{q/2})$. On the
complement, boundedness of $h_R$ gives $|Z|\le C Q$. Integration of
the cumulative exponential tail bounds the remaining $q$th moment by
$C_q(1+n)^q e^{-c'n}$. This proves
\eqref{eq:all-stopped-moment-bound}.

We give the full multiscale argument for the difference. For each
$L\ge1$, let $B_L$ be the sum of the four prefix or suffix maxima
on the deterministic length-$\lceil L\rceil$ windows adjacent to
the two endpoints $-m_k$ and $m_{n-k}$. All these windows are permitted
to cross collision time zero. They are fixed intervals, not intervals
conditioned on the clocks. By Lemma~\ref{lem:fixed-even-maximal},
\begin{equation}\label{eq:deterministic-shell-maximal}
 \int B_L^{2p}\dd\nu_R^*\le C_p L^p.
\end{equation}
On $D_{n,k,R}\le L$ one has the pathwise bound
$|Z-W|\le B_L$. Let $b_0=\lceil\sqrt n\rceil$ and
$L_j=2^j b_0$ for $j\ge1$. Partition the sample space into
$\{D\le b_0\}$ and the disjoint sets
$\{L_j/2<D\le L_j\}$, $j\ge1$. The clocks are finite almost
everywhere, so these sets exhaust the section up to a null set.
The first set contributes at most $C_q b_0^{q/2}$ to
$\int|Z-W|^q\dd\nu_R^*$.

On every other shell, H\"older's inequality gives
\begin{equation}\label{eq:dyadic-shell-holder}
 \int_{\{L_j/2<D\le L_j\}}|Z-W|^q\dd\nu_R^*
 \le C_{p,q}L_j^{q/2}
       \nu_R^*(D>L_j/2)^{1-q/(2p)}.
\end{equation}
The exponent is positive. Fix $A$ large enough that the exponential
bound \eqref{eq:total-clock-exponential} at threshold $L_j/2$ is
at most $Ce^{-c_0L_j/4}$ whenever $L_j>An$.
For the shells $L_j\le An$, use
\eqref{eq:total-clock-polynomial} in
\eqref{eq:dyadic-shell-holder}. Their total is at most
\begin{align}
 C_{p,q,A}\sum_{L_j\le An} n^{p-q/2}L_j^{3q/2-2p}
 &\le C_{p,q,A}n^{q/4}
             \sum_{j\ge1}2^{j(3q/2-2p)}
 \le C_q n^{q/4}.
 \label{eq:dyadic-clock-summation}
\end{align}
Here $b_0\asymp\sqrt n$ and $3q/2-2p<0$. For the remaining shells,
\eqref{eq:dyadic-shell-holder} is bounded by
$C_{p,q}L_j^{q/2}e^{-c_qL_j}$, which is summable and contributes no
more than $C_qn^{q/4}$. Taking the $q$th root proves
\eqref{eq:quarter-stopping}. There is no logarithmic count of shells:
the shell contributions form a convergent geometric series. For $q=2$
one can take $p=2$, in which case the moderate-shell sum is
$C\sum n/L_j\le C\sqrt n$.

Finally, recenter at the actual mark and use
$|e^{ix}-e^{iy}|\le|x-y|$. Since $a$ vanishes outside the section,
$|a|\nu\le c_*M_a\nu_R^*$ there. Equation~\eqref{eq:quarter-stopping}
with $q=1$ proves \eqref{eq:quarter-characteristic-stopping}; its right
side is exactly zero at $v=0$. No martingale representation, conditional
independence, or mixing of the induced map is used.
\end{proof}

\subsection{One anchored factor in an arbitrary fixed moment}
For an integer $d\ge1$ and a symmetric positive semidefinite matrix
$D$, write
\[
 \mathcal G_d(D)=\mathbb E\,\mathcal Z_D^{\otimes d},\qquad
 \mathcal Z_D\sim N(0,D),\qquad
 \rho_d=\begin{cases}1/2,&d\text{ odd},\\1,&d\text{ even}.\end{cases}
\]
The tensor is zero for odd $d$. Its norm below may be any fixed
Euclidean tensor norm.

\begin{lemma}[A fixed marked collision moment]
\label{lem:all-anchored-collision-moments}
Let $a$ be complex and bounded-BV on the collision cylinder, with
$\alpha=\int a\dd\nu$. For a deterministic two-sided interval of
length $m=r+s\ge1$,
\begin{align}
 \left\|\int_M a\,\mathsf H_{r,s,R}^{\otimes d}\dd\nu
                  -\alpha m^{d/2}\mathcal G_d(\Gamma_R)\right\|
 \le C_d(\|a\|_\infty+\|a\|_{\mathrm{BV}})m^{d/2-\rho_d}.
 \label{eq:all-anchored-collision-moments}
\end{align}
The constant is independent of the position of zero in the interval.
\end{lemma}
\begin{proof}
First take a real unit direction $\xi$ and put $u=\xi\cdot h_R$.
Consider a joint cumulant containing $a-\alpha$ at time zero and
$j\ge1$ copies of $u$ at integer times $i_1,\ldots,i_j$.
The independent-block comparison in
Proposition~\ref{prop:connected-collision-summability} applies also
to these different observables: the finite-order decoupling estimate
is multilinear, and the partition polynomial is a finite sum of
products of subset moments. After normalization of $a-\alpha$, its
absolute value is bounded by
\[
 C_j(\|a\|_\infty+\|a\|_{\mathrm{BV}})
    e^{-c_j\operatorname{diam}\{0,i_1,\ldots,i_j\}}.
\]
It is summable over all $\Z^j$. Thus the cumulant containing the
single centered mark and $j$ copies of the sum of $u$ on any interval
is bounded by $C_j(\|a\|_\infty+\|a\|_{\mathrm{BV}})$, independently
of its length or position. Complex marks are handled by real and
imaginary parts.

Expand the moment with one centered mark by the finite
moment--cumulant partition identity. The block containing that mark
must contain at least one copy of the sum. Every other nonzero block
contains at least two copies, because $u$ is centered. At most
$\lfloor(d-1)/2\rfloor$ such blocks remain. Each pure sum cumulant
is $O_d(m)$ by the inherited connected-correlation theorem. Hence
\[
 \left|\int(a-\alpha)(\xi\cdot\mathsf H_{r,s,R})^d\dd\nu\right|
 \le C_d(\|a\|_\infty+\|a\|_{\mathrm{BV}})
                         m^{\lfloor(d-1)/2\rfloor}.
\]
For the unweighted moment, all cumulants of order at least two are
$O_d(m)$ and the second is $m\xi^{\mathsf T}\Gamma_R\xi+O(1)$.
If $d$ is even, the pair partitions give the Gaussian term; replacing
one covariance by its bounded correction, or using a nonpair
partition, lowers the power of $m$ by at least one. If $d$ is odd,
there are at most $(d-1)/2$ nonzero blocks and the Gaussian moment
is zero. For $d=1$ the unweighted moment is exactly zero.
These observations prove \eqref{eq:all-anchored-collision-moments}
against every $\xi^{\otimes d}$. Polarization in the fixed dimension
four gives the symmetric tensor assertion. All cumulant orders used
are fixed, at most $d+1$.
\end{proof}

\begin{theorem}[All fixed Gaussian moments with an actual return mark]
\label{thm:all-marked-gaussian-moments}
For every fixed integer $d\ge1$, all $n\ge2$ and $0\le k\le n$,
\begin{align}
 &\left\|\int_{Y_R^*}a((F_R^*)^kx)
        \left(\frac{U_{n,R}(x)}{\sqrt n}\right)^{\otimes d}\dd\nu(x)
                   -\alpha_a\mathcal G_d(D_R)\right\|\notag\\
 &\hspace{18mm}\le C_d\left(M_an^{-1/4}
                           +(M_a+V_a)n^{-\rho_d}\right).
 \label{eq:all-marked-gaussian-moments}
\end{align}
The convention is the restricted collision measure, with
$\alpha_a=\int a\dd\nu$; no positivity of $a$ is required.
In particular the actual normalized covariance and its induced
Cesaro Green--Kubo expression converge to $D_R$ at rate $O(n^{-1/4})$.
\end{theorem}
\begin{proof}
Recenter at $y=(F_R^*)^kx$. The original numerator is
\[
 \int_{Y_R^*}a(y)
       \left(\frac{Z_{n,k,R}(y)}{\sqrt n}\right)^{\otimes d}\dd\nu(y).
\]
For $d\ge2$, the tensor telescoping inequality
\[
 \|x^{\otimes d}-y^{\otimes d}\|
       \le C_d|x-y|(|x|+|y|)^{d-1}
\]
and H\"older's inequality with exponents $d$ and $d/(d-1)$ show
that its change on replacing $Z$ by $W$ is at most $C_dM_an^{-1/4}$,
by \eqref{eq:all-stopped-moment-bound} and \eqref{eq:quarter-stopping}.
For $d=1$ use the first moment version directly. Restriction from
$\nu$ to the section changes only the fixed factor $c_*$.

Apply Lemma~\ref{lem:all-anchored-collision-moments} to $W$ with
$m=m_k+m_{n-k}=n/c_*+O(1)$. Division by $n^{d/2}$ gives the second
error in \eqref{eq:all-marked-gaussian-moments}. The Gaussian term
is $\alpha_a\mathcal G_d((m/n)\Gamma_R)$, which differs from the
stated term by $O_d(M_a/n)$ for even $d$ and is zero for odd $d$.
The covariance matrices are uniformly bounded. This proves the theorem.

For the unweighted assertion take $a=s_R$, so $\alpha_a=1$ and both
insertion norms are uniformly bounded. The Kac mean makes the actual
first centered moment zero. At $d=2$ the theorem is therefore a
covariance estimate. Expanding that exact covariance under stationarity
of $F_R^*$ gives the Cesaro expression in
\eqref{eq:induced-cesaro-covariance}, with the improved error.
Absolute summability of induced correlations is not asserted.
\end{proof}

\begin{corollary}[Polynomial moments under the unchanged event]
\label{cor:all-same-event-moments}
Let $E_{n,R}\subset Y_R^*$ have probability $p_{n,R}>0$ and
zero-extended indicator variation $V_{n,R}$. For
$A_{n,k,R}=\{x:(F_R^*)^kx\in E_{n,R}\}$ and every fixed $d\ge1$,
\begin{align}
 &\left\|\mathbb E_{\nu_R^*}
   \left[\left(U_{n,R}/\sqrt n\right)^{\otimes d}\mid A_{n,k,R}\right]
                       -\mathcal G_d(D_R)\right\|\notag\\
 &\hspace{18mm}\le\frac{C_d}{p_{n,R}}
         \left(n^{-1/4}+(1+V_{n,R})n^{-\rho_d}\right).
 \label{eq:all-same-event-moments}
\end{align}
For $p_{n,R}\ge cn^{-\beta}$ and $V_{n,R}\le Cn^\kappa$,
$\beta<1/4$ and $\beta+\kappa<\rho_d$ suffice. Every fixed
polynomial moment and the conditional covariance converge under
$\beta<1/4$ and $\beta+\kappa<1/2$.
\end{corollary}
\begin{proof}
Take $a=c_*^{-1}\one_{E_{n,R}}$ in
Theorem~\ref{thm:all-marked-gaussian-moments}. Its mass equals
$p_{n,R}$, also the probability of $A_{n,k,R}$ by invariance.
Divide by this unchanged denominator. The conditional first moment
tends to zero under the stronger stated conditions, so subtracting
its tensor square from the conditional second moment proves the
covariance assertion. A fixed polynomial is a finite linear
combination of the displayed tensor coordinates.
\end{proof}

\paragraph{Scope of the moment improvement.}
The fourth-root rate is obtained for the actual unbounded record at
one prescribed count and mark. It is not a maximal-in-$k$ theorem,
a statement about growing moment orders, or an event-replacement
estimate. In particular it neither assumes induced mixing nor treats
a first-variation bound as a bound for inverse-coarea second derivatives.

\section{A wider fixed-count band with the original central rate}
\label{sec:rate-preserving-band}
The multiscale stopping estimate can be combined with damped cubic
unsmoothing without altering the actual return mark. We choose the
outer radius so that the integrated stopping error has the same power
as the original marked central estimate. Thus the wider band no longer
requires a slower probability-decay budget for a single marked event.

Set
\begin{equation}\label{eq:rate-preserving-scales}
 \varepsilon_\natural=\frac{67}{1400},\qquad
 V_n^\natural=n^{67/1400},\qquad
 B_n^\natural=n^{-633/1400}.
\end{equation}
The smooth Fourier support radii are $2V_n^\natural$ after rescaling
and $2B_n^\natural$ in physical frequency. The preceding radius
$n^{1/21}$ is strictly smaller than $V_n^\natural$. All statements
below are at the prescribed return count $n$, not averaged in $n$.

\begin{theorem}[Rate-preserving prescribed-count annulus]
\label{thm:rate-preserving-annulus}
Uniformly in $R\in I$, $0\le k\le n$ and the marked bounded-BV
insertion $a$, for all sufficiently large $n$,
\begin{align}
 &n^2\int_{2n^{-99/200}\le|z|\le2B_n^\natural}
                     |\Phi^{[k],a}_{n,R}(z)|\dd z
       \le C\left(M_an^{-3/280}+V_an^{-983/350}\right),
       \label{eq:rate-preserving-annulus}\\
 &\int_{|v|\le2V_n^\natural}
       |C^{[k],a}_{n,R}(v)-\alpha_a e^{-v^{\mathsf T}D_Rv/2}|\dd v
       \le C\left(M_an^{-3/280}\sqrt{\log(2+n)}
                           +V_an^{-9/175}\right).
       \label{eq:rate-preserving-central}
\end{align}
The original probability corresponds to $a=s_R$. The transform with a
general insertion is still integrated against $\nu|_{Y_R^*}$, with
mass $\alpha_a=\int a\dd\nu$.
\end{theorem}
\begin{proof}
Recenter at the actual mark. Equation~\eqref{eq:quarter-characteristic-stopping}
compares the original transform with
$I^a_{m_k,m_{n-k},R}(v/\sqrt n)$ at cost $CM_a|v|n^{-1/4}$.
The deterministic length satisfies $m_k+m_{n-k}=n/c_*+O(1)$.
Apply Theorem~\ref{thm:damped-cubic-unsmoothing} with
\[
 Q=19,\qquad \delta=\tfrac14 n^{-1/5},\qquad
 \epsilon=\tfrac14 n^{-3},\qquad |v|\le2n^{67/1400}.
\]
The coarse scale determines the collision perturbation disk, while
only the insertions use the fine scale. For $z=v/\sqrt n$, the two
required quantities satisfy
\begin{align}
 |z|\delta^{-2}&=O(n^{-73/1400}),\notag\\
 |z|^{18}\delta^{-40}&=O(n^{-97/700}).
 \label{eq:rate-preserving-damping-domain}
\end{align}
Hence the damping theorem applies for all sufficiently large $n$.
On the annulus $2n^{1/200}\le|v|\le2V_n^\natural$, the damped
polynomial term is smaller than every fixed inverse power of $n$:
its prefactor is polynomial and its exponential is at most
$e^{-c n^{1/100}}$.

We display every nondamping integrated exponent. Four-dimensional
volume and the largest frequency factor give respectively
\begin{align}
 \epsilon V_a(V_n^\natural)^4
      &\le C V_an^{-983/350},\notag\\
 M_a\epsilon(1+\sqrt n V_n^\natural)^3(V_n^\natural)^4
      &\le C M_an^{-233/200},\notag\\
 M_a\delta^2L_\delta^2(V_n^\natural)^8
      &\le C M_an^{-3/175}[\log(2+n)]^2,\notag\\
 M_an^{-1}\delta L_\delta^3(V_n^\natural)^8
      &\le C M_an^{-143/175}[\log(2+n)]^3,\notag\\
 M_an^{-1/4}(V_n^\natural)^5
      &= M_an^{-3/280}.
 \label{eq:rate-preserving-exponents}
\end{align}
The quartic residual exponent exceeds $3/280$ by $9/1400$, so its
logarithms are absorbed by this strict power margin. Every other
supremum-norm error also decays faster than the last term. Since
$v=\sqrt n z$ has Jacobian $\dd v=n^2\dd z$ in four dimensions,
these estimates prove \eqref{eq:rate-preserving-annulus} with the
raw factor $n^2$ and no division by annular volume.

On $|v|\le2n^{1/200}$, use the sharper inherited
Theorem~\ref{thm:marked-return-band}. On the rest of the new ball,
use \eqref{eq:rate-preserving-annulus} and the uniformly elliptic
Gaussian tail. Since $983/350>9/175$, their sum gives
\eqref{eq:rate-preserving-central}. This does not replace the old
small-ball theorem by a weaker one.
\end{proof}

\begin{remark}[The choice of radius]
The equality
\[
 \frac14-5\varepsilon_\natural=\frac3{280}
\]
selects the radius for which this stopping estimate retains the original
central power. It is a balance of the displayed error bounds, not an
assertion of optimality or an obstruction to another method. With a
smaller positive rate, the preceding construction still reaches every
fixed $1/200<\varepsilon<1/20$. Neither statement covers the rescaled
endpoint $n^{1/10}$, corresponding to physical frequency $n^{-2/5}$.
\end{remark}

\begin{corollary}[The original rare-event budget on the wider band]
\label{cor:rate-preserving-same-event}
For the unchanged marked-state event $A_{n,k,R}$, probability $p_{n,R}$,
and indicator variation $V_{n,R}$ of
Corollary~\ref{cor:all-same-event-moments},
\begin{align}
 &\int_{|v|\le2V_n^\natural}
 \left|\mathbb E_{\nu_R^*}
       [e^{iv\cdot U_{n,R}/\sqrt n}\mid A_{n,k,R}]
                               -e^{-v^{\mathsf T}D_Rv/2}\right|\dd v\notag\\
 &\hspace{18mm}\le\frac C{p_{n,R}}
        \left(n^{-3/280}\sqrt{\log(2+n)}+V_{n,R}n^{-9/175}\right).
 \label{eq:rate-preserving-same-event}
\end{align}
Thus $p_{n,R}\ge cn^{-\beta}$ and $V_{n,R}\le Cn^\kappa$ suffice
when $\beta<3/280$ and $\beta+\kappa<9/175$. Under these conditions,
every fixed conditional polynomial moment also has the Gaussian limit.
\end{corollary}
\begin{proof}
Use $a=c_*^{-1}\one_{E_{n,R}}$ in
\eqref{eq:rate-preserving-central}. Its mass is the exact denominator
$p_{n,R}$. Division proves the bound without changing the event.
The stated inequalities also imply the moment conditions in
Corollary~\ref{cor:all-same-event-moments}.
\end{proof}

\subsection{The raw identity with the new kernel}
Define
\[
 \chi_n^\natural(\omega)=\chi(\omega/B_n^\natural),\qquad
 K_n^\natural=\mathcal F^{-1}\chi_n^\natural.
\]
For a marked insertion additionally satisfying the finite-record
admissibility in Definition~\ref{def:finite-record-weight}, use the
same weight $w_{n,k,R}=c_*a\circ(F_R^*)^k$ and the same exact finite
extraction $\mu^{a,L}=E^{a,L}+Q^{a,L}$ as in
Theorem~\ref{thm:sharp-weighted-raw-inversion}. Set, with the new
kernel throughout,
\begin{align}
 \mathcal E^{\natural,a,L}_{n,k,R,A}
 &=\mathop{\rm ess\,sup}_{\mathcal W_{n,R,A}}
                 |e^{a,L}-K_n^\natural*E^{a,L}|,\notag\\
 \mathcal C^{\natural,a,L}_{n,k,R}(B)
 &=\int_{\T^3\times[-B,B]}
       |1-\chi_n^\natural(\omega)|\,|\widehat Q^{a,L}(\omega)|\dd\omega.
 \label{eq:rate-preserving-raw-terms}
\end{align}
In particular the edge correction is not copied from a different
cutoff. It may still be infinite before its germs are estimated.

\begin{corollary}[Exact weighted raw budget with the retained central rate]
\label{cor:rate-preserving-raw-budget}
Let $L_n=\lceil\lambda n\rceil$ with $\lambda>c_*^{-1}+1$.
For $B\ge1$, $P>0$, and $y=((\ell,t)-n\bar G_R)/\sqrt n$,
\begin{align}
 &\mathop{\rm ess\,sup}_{\mathcal W_{n,R,A}}
        |n^2p_{n,R}^{[k],a}(\ell,t)-\alpha_a\mathfrak g_R(y)|\notag\\
 &\quad\le C\left(M_an^{-3/280}\sqrt{\log(2+n)}+V_an^{-9/175}\right)
      +CM_ae^{-c n^{67/700}}+C_{A,\lambda,P}M_an^{-P}\notag\\
 &\qquad+n^2\mathcal E^{\natural,a,L_n}_{n,k,R,A}
     +(2\pi)^{-4}n^2\mathcal C^{\natural,a,L_n}_{n,k,R}(B)
     +\frac{n^2}{\pi B}A_2(n,L_n,R,w_{n,k,R}).
 \label{eq:rate-preserving-raw-budget}
\end{align}
The inequality has the extended-real interpretation when needed.
Both cutoff radii are as in \eqref{eq:rate-preserving-scales}.
\end{corollary}
\begin{proof}
On the central count labels the high-count tail has no density
contribution. The same exact support and absolute-residual-inversion
argument gives
\begin{align*}
 p_{n,R}^{[k],a}={}&K_n^\natural*\mu_{n,R}^{[k],a}
       +(e^{a,L_n}-K_n^\natural*E^{a,L_n})\\
 &+\mathcal F^{-1}[(1-\chi_n^\natural)\widehat Q^{a,L_n}]
       -K_n^\natural*\mathsf T^{a,L_n}.
\end{align*}
Use \eqref{eq:rate-preserving-central} for the first term. The missing
Gaussian frequencies have $|v|\ge V_n^\natural$ and give the stated
exponential tail. Count separation is unchanged, but the kernel scale
is new. Its Schwartz estimate is
\[
 n^2\sup_{\mathcal W_{n,R,A}}|K_n^\natural*\mathsf T^{a,L_n}|
       \le C_{A,\lambda,J}M_a
                     n^{67/350-767J/1400}.
\]
Indeed $n^2(B_n^\natural)^4=n^{67/350}$ and
$nB_n^\natural=n^{767/1400}$. Choose $J$ for the prescribed power
$P$. Splitting the residual inverse at roof frequency $B$ gives
exactly the last two terms, by the inherited absolute second-derivative
far-tail estimate. The local edge term stays in its combined form.
\end{proof}

\paragraph{What remains beyond this band.}
The new fixed-count estimate improves the physical transform, rather
than an average or a compressed operator. The outer region beginning
at $2n^{-633/1400}$, compact nonzero and peripheral return frequencies,
and growing roof frequencies still enter the complementary residual
integral. The long-time bound for $A_2$ at $L_n\asymp n$ and the local
edge correction remain distinct terms in
\eqref{eq:rate-preserving-raw-budget}. A relative comparison with an
exact physical-time event is not performed. The new all-moment and
same-event conclusions preserve the actual indicator and its exact
probability; they are not a substitute for those raw estimates.

\section{Damped removal of smoothing at arbitrary fixed order}
\label{sec:high-order-damped-unsmoothing}
The cubic expansion is sufficient for the isotropic band of
Section~\ref{sec:rate-preserving-band}. Different widths in the four
physical frequency coordinates require a separate balance. We first
prove a finite-order version in which the residual degree, the spectral
Taylor degree, and the two smoothing scales are independent. The
expansion order is fixed before the collision length varies. No estimate
uniform in that order, and no analytic unsmoothed twist, is used.

\subsection{Small-mass moments and heterogeneous cumulants}
Write $L_\delta=1+|\log\delta|$ and use the stationary probability
$\nu$ in this subsection.

\begin{lemma}[Every fixed even small-mass moment]
\label{lem:all-small-mass-moments}
Fix an integer $P\ge1$ and $B<\infty$. Suppose that $u$ is real,
centered, $\|u\|_{\mathcal V}\le B$, and
$\|u\|_1\le B\delta$, with $0<\delta<1/2$. Then
\begin{equation}\label{eq:all-small-mass-moments}
 \int|S_{m,R}u|^{2P}\dd\nu
 \le C_{P,B}\sum_{b=1}^{P}(m\delta)^bL_\delta^{2P-b}
 \qquad(m\ge1).
\end{equation}
The estimate holds on any translated or reversed deterministic
interval. Its constant is uniform in $R,u,\delta,m$ within the
stated class, but may depend on $P$.
\end{lemma}
\begin{proof}
Lemma~\ref{lem:small-mass-cumulants} gives, simultaneously for the
finitely many orders $2\le j\le2P$,
\[
 |\operatorname{cum}_j(S_{m,R}u)|
          \le C_{P,B}m\delta L_\delta^{j-1}.
\]
Expand the $2P$th moment by the finite partition formula. Centering
removes singleton blocks. A remaining partition with $b$ blocks has
$1\le b\le P$, and its absolute contribution is at most
\[
 C_{P,B}(m\delta)^b
       L_\delta^{\sum_{A\in\pi}(|A|-1)}
 =C_{P,B}(m\delta)^bL_\delta^{2P-b}.
\]
There are finitely many such partitions at the fixed order. Absorb
their numbers in $C_{P,B}$ and sum over $b$. The even moment is the
absolute even moment. Invariance translates an interval, and reversal
only changes the order of summation. This is an exact finite
moment--cumulant calculation, not a truncation of a convergent
infinite cumulant series.
\end{proof}

\begin{lemma}[One heterogeneous cumulant comparison]
\label{lem:heterogeneous-cumulant-comparison}
For each fixed $q\ge2$, and bounded-BV collision observables
$u_1,\ldots,u_q$, the joint cumulant at times $t_1,\ldots,t_q$
satisfies
\begin{equation}\label{eq:heterogeneous-cumulant-comparison}
 |\operatorname{cum}(u_1\circ T_R^{t_1},\ldots,u_q\circ T_R^{t_q})|
 \le C_q e^{-c_q\operatorname{diam}\{t_1,\ldots,t_q\}}
                     \prod_{j=1}^q\|u_j\|_{\mathcal V}.
\end{equation}
The same bound with a fixed factor $1+\operatorname{diam}$ is
summable over all times relative to one anchored time.
\end{lemma}
\begin{proof}
Normalize the nonzero norms to one and order the times, preserving
the identity of each observable. Split at a largest consecutive gap
$d$. Under a comparison law, keep the entire left and right marginal
laws unchanged but make them independent. If $m_A$ denotes the moment
of the variables in a nonempty subset $A$, its change satisfies
$|m_A-m_A^\circ|\le C_qe^{-\gamma_qd}$ whenever $A$ meets both
sides, by Lemma~\ref{lem:all-finite-product}; otherwise its change is
zero. For a partition $\pi=\{A_1,\ldots,A_b\}$, write explicitly
\[
 \prod_{j=1}^b m_{A_j}-\prod_{j=1}^b m_{A_j}^\circ
 =\sum_{i=1}^b(m_{A_i}-m_{A_i}^\circ)
     \prod_{j<i}m_{A_j}\prod_{j>i}m_{A_j}^\circ.
\]
Every factor is bounded. Substitute this identity into
\eqref{eq:finite-order-cumulant-definition}. The comparison cumulant
vanishes by independent-block cancellation, so the original one is
$O_q(e^{-\gamma_qd})$. The diameter is at most $(q-1)d$; coincident
times are treated by the bounded partition formula. Restoring the
norms proves the bound and exponential summability.

For example, take four centered, possibly different variables
$A,B,C,D$ in chronological order, with the chosen split between
$B$ and $C$. The fourth cumulant is
\[
 \mathbb E(ABCD)-\mathbb E(AB)\mathbb E(CD)
 -\mathbb E(AC)\mathbb E(BD)-\mathbb E(AD)\mathbb E(BC).
\]
Under the comparison law, the first two terms agree and the last
two vanish. Each difference is controlled by a cross-gap subset
moment, with the unchanged same-side moments bounded. This includes
$A=a-\alpha$ and three collision factors, or a split on either side
of that mark. No identical-distribution assumption is involved.
For the moment theorem of Section~\ref{sec:multiscale-stopping},
the mark block still consumes at least one collision factor; the
remaining nonzero blocks consume at least two each. Thus the parity
rates there remain $\rho_d=1/2$ for odd $d$ and $\rho_d=1$ for even
$d$. The constants in this comparison are fixed-order constants.
\end{proof}

\subsection{A corrected chronological word of fixed arbitrary degree}
Use $I^a_{r,s,R}(z)$ from Section~\ref{sec:damped-unsmoothing},
with $m=r+s\ge1$. Its mark is a bounded-BV function on the collision
cylinder; it need not be a probability density.

\begin{theorem}[High-order damped unsmoothing]
\label{thm:high-order-damped-unsmoothing}
Fix integers $P\ge1$ and $Q\ge3$. Suppose $0<\delta\le\delta_0(Q)$
and real $z$ satisfy \eqref{eq:finite-jet-damping-domain} for the
spectral degree $Q$. For any fine scale $0<\epsilon<1/2$,
\begin{align}
 |I^a_{r,s,R}(z)|\le{}&
 C_{P,Q}M_a\epsilon^{-4P}(1+m|z|)^{2P-1}e^{-cm|z|^2}
       +C\epsilon V_a\notag\\
 &+C_PM_a\epsilon(1+m|z|)^{2P-1}\notag\\
 &+C_PM_a|z|^{2P}
             \sum_{b=1}^{P}(m\delta)^bL_\delta^{2P-b}.
 \label{eq:high-order-damped-unsmoothing}
\end{align}
The constants are independent of the mark location $0\le r\le m$,
$R,m,\delta,\epsilon,z$ in this domain. The residual Taylor degree
is $2P-1$, not $Q$. The fine scale enters only the multipliers and
not the analytic damping domain.
\end{theorem}
\begin{proof}
For $z\ne0$ put $u=(z/|z|)\cdot(h_R-h_{R,\delta})$. It is real,
centered, uniformly bounded-BV and has $L^1$ norm at most $C\delta$.
The case $z=0$ follows by increasing the constant. First replace
$a$ by $a_\epsilon=\mathsf S_\epsilon a$, while the full phase has
modulus one. This costs $C\epsilon V_a$. Translate the interval so
that it begins at zero and the mark is at time $r$. The exact
real-variable Taylor formula gives
\[
 e^{iX}=\sum_{j=0}^{2P-1}\frac{(iX)^j}{j!}+\mathcal R_{2P}(X),
 \qquad |\mathcal R_{2P}(X)|\le\frac{|X|^{2P}}{(2P)!},
 \qquad X=|z|S_m u.
\]
The coarse phase has modulus one, and
Lemma~\ref{lem:all-small-mass-moments} bounds the remainder integral
by the last line of \eqref{eq:high-order-damped-unsmoothing}.
In particular the term with $b=1$ is retained even when $m\delta$
is small; only the term with $b=P$ would not be a uniform bound.

For degree $j\le2P-1$, expand the $j$th power of the residual sum
as the sum over its $m^j$ time tuples, including repetitions. In each
product replace every $u$ by $u_\epsilon=\mathsf S_\epsilon u$.
Telescoping, invariance and the uniform supremum bounds give an
integrated cost at most $C_jM_a\epsilon$ per tuple. Summation over
the fixed degrees gives the third term of
\eqref{eq:high-order-damped-unsmoothing}.

Sort the residual times together with the mark as
$0\le t_1\le\cdots\le t_d\le m$, $d=j+1\le2P$. The corrected
integral is exactly
\[
 \ell\bigl(L_z^{m-t_d}M_{b_d}L_z^{t_d-t_{d-1}}
  \cdots M_{b_2}L_z^{t_2-t_1}M_{b_1}L_z^{t_1}\nu\bigr),
 \qquad L_z=\mathcal L_{R,\delta,z}.
\]
One multiplier is $a_\epsilon$ and the other $j$ multipliers are
$u_\epsilon$. Each costs at most its fixed supremum bound times
$C\epsilon^{-2}$. The power lengths sum to exactly $m$, including
zero-length powers for repeated times and endpoint marks.
Corollary~\ref{cor:finite-jet-damping} bounds the word by
$C_{j,Q}M_a\epsilon^{-2(j+1)}e^{-cm|z|^2}$. The product of its
finitely many power-bound constants is absorbed in $C_{j,Q}$;
no count-dependent number of factors is present. Summing the tuples
and degrees proves the first term of
\eqref{eq:high-order-damped-unsmoothing}. This proves the assertion
without perturbing the operator by the fine-scale residual.
\end{proof}

\begin{remark}[Fixed orders and the retained cubic theorem]
At $P=2$, the last sum is
$m^2\delta^2L_\delta^2+m\delta L_\delta^3$, and the first
multiplier loss is $\epsilon^{-8}$. Thus
Theorem~\ref{thm:damped-cubic-unsmoothing} is recovered exactly.
Higher fixed $P$ gives a different frequency budget; it does not
assert that $P$ or $Q$ may grow with $n$. The original cubic theorem
and the sharper small-ball estimate remain separate statements.
\end{remark}

\section{Anisotropic Fourier bands at the prescribed return count}
\label{sec:anisotropic-fixed-count-bands}
The four coordinates of the raw record have different roles in the
final inversion: three are discrete and the fourth is flight time.
The same prescribed-count argument permits different Fourier widths
in these coordinates. We prove a family of such estimates, retaining
the four-dimensional Lebesgue factor in frequency and the original
return count and mark. It gives new portions of the farther annulus;
no average over counts, directions, or frequency volume is taken.

\subsection{A feasible region of four width exponents}
Fix $\boldsymbol\theta=(\theta_1,\theta_2,\theta_3,\theta_4)$ and set
\begin{equation}\label{eq:anisotropic-widths}
 \theta_* =\max_i\theta_i,\quad \Theta=\sum_{i=1}^4\theta_i,
 \quad B_{n,i}=n^{-1/2+\theta_i},\quad
 \mathcal B_n=\prod_{i=1}^4[-2n^{\theta_i},2n^{\theta_i}].
\end{equation}
Thus the physical box is $n^{-1/2}\mathcal B_n$ with half-widths
$2B_{n,i}$. The coordinates are ordered as
$(K_1,K_2,N,T)$ and their dual frequencies as $(u_1,u_2,s,b)$.
For all large $n$ the first three physical intervals lie in the
chosen local chart of $\T^3$.

\begin{theorem}[Prescribed-count cancellation with separate widths]
\label{thm:anisotropic-fixed-count-band}
Suppose
\begin{equation}\label{eq:anisotropic-feasible-region}
 \min_i\theta_i\ge\frac1{200},\qquad
 \theta_*<\frac1{10},\qquad
 \Theta+\theta_*<\frac14.
\end{equation}
Fix $0<\eta\le1/4-\Theta-\theta_*$. There is a number $\sigma>0$
such that, uniformly in $R\in I$, $0\le k\le n$ and bounded-BV
section insertions $a$, for all sufficiently large $n$,
\begin{align}
 &n^2\int_{\substack{z\in n^{-1/2}\mathcal B_n\\
                         |z|\ge2n^{-99/200}}}
             |\Phi^{[k],a}_{n,R}(z)|\dd z
       \le C\left(M_an^{-\eta}+V_an^{-\sigma}\right).
 \label{eq:anisotropic-fixed-count-band}
\end{align}
One may choose $\sigma$ larger than any prescribed fixed number by
reducing only the fine multiplier scale. If
$\eta\ge3/280$ and $\sigma\ge9/175$, then
\begin{align}
 &\int_{\mathcal B_n}
 |C^{[k],a}_{n,R}(v)-\alpha_a e^{-v^{\mathsf T}D_Rv/2}|\dd v
 \le C\left(M_an^{-3/280}\sqrt{\log(2+n)}+V_an^{-9/175}\right).
 \label{eq:anisotropic-central-band}
\end{align}
The measure convention is $\nu|_{Y_R^*}$ and
$\alpha_a=\int a\dd\nu$. The original probability is $a=s_R$.
All exponent choices and expansion orders are fixed independently of
$n,R,k,a$.
\end{theorem}
\begin{proof}
Here are choices with strict analytic and residual margins. Select
\begin{equation}\label{eq:anisotropic-order-choices}
 2\theta_*<\alpha<\frac14-\frac{\theta_*}{2},\quad
 P(\alpha-2\theta_*)>\Theta+\eta,\quad
 (Q-1)(1/2-\theta_*)>2\alpha(Q+1),
\end{equation}
where $P\ge1$ and $Q\ge3$ are integers. The interval for $\alpha$
is nonempty because $\theta_*<1/10$. Once $\alpha$ is chosen,
choose $P$ and $Q$ sufficiently large; both displayed strict
inequalities have positive coefficients at large order. Choose a
fixed $\beta$ so that
\begin{equation}\label{eq:anisotropic-fine-choice}
 \beta-\Theta-(2P-1)(1/2+\theta_*)>\eta.
\end{equation}
Set $\delta=n^{-\alpha}/4$ and $\epsilon=n^{-\beta}/4$.
There is no relation $Q=2P-1$.

Recenter the orbit at its prescribed actual-return mark and use
\eqref{eq:quarter-characteristic-stopping}. Its integrated cost is
\begin{equation}\label{eq:anisotropic-stopping-integral}
 CM_an^{-1/4}\int_{\mathcal B_n}|v|\dd v
       \le CM_an^{-1/4+\Theta+\theta_*}\le CM_an^{-\eta}.
\end{equation}
The deterministic interval has length $m=n/c_*+O(1)$ uniformly in
$k$. For $z=v/\sqrt n$ in the box, $|z|\le4n^{-1/2+\theta_*}$.
The two coarse damping quantities are therefore bounded by constants
times
\[
 n^{-(1/2-\theta_*-2\alpha)},\qquad
 n^{-[(Q-1)(1/2-\theta_*)-2\alpha(Q+1)]}.
\]
Both powers decay by \eqref{eq:anisotropic-order-choices}.
Apply Theorem~\ref{thm:high-order-damped-unsmoothing}.
On $|v|\ge2n^{1/200}$ its first term, even after integration, is
smaller than every fixed inverse power of $n$: the prefactor is a
fixed polynomial, and $m|z|^2\ge c n^{1/100}$.

The fine insertion error integrates to $CV_an^{-(\beta-\Theta)}$.
The fine residual replacement is at most
\[
 CM_an^{-[\beta-\Theta-(2P-1)(1/2+\theta_*)]}.
\]
For each of the $P$ residual moment terms, the integrated bound is
\begin{align}
 &CM_an^{\Theta+2P\theta_*-P+(1-\alpha)b}
                       [\log(2+n)]^{2P-b},\qquad 1\le b\le P.
 \label{eq:anisotropic-residual-powers}
\end{align}
Since $\alpha<1$, the largest exponent occurs at $b=P$ and is
$\Theta-P(\alpha-2\theta_*)<-\eta$. Each logarithmic power is
fixed and absorbed by the strict margin. This treats also the
connected terms with fewer than $P$ blocks; they have additional
powers of $n^{-(1-\alpha)}$. Put $\sigma=\beta-\Theta$.
The change of variables satisfies $\dd v=n^2\dd z$, exactly, so the
bounds prove \eqref{eq:anisotropic-fixed-count-band} with the raw
factor $n^2$. Increasing $\beta$ increases $\sigma$ without changing
the coarse damping domain or the polynomial-versus-exponential argument.

The ball $|v|\le2n^{1/200}$ is contained in $\mathcal B_n$.
Use Theorem~\ref{thm:marked-return-band} there. On the rest use
\eqref{eq:anisotropic-fixed-count-band} and the uniformly elliptic
Gaussian tail. Under the stated inequalities on $\eta,\sigma$,
this gives \eqref{eq:anisotropic-central-band}.
\end{proof}

\begin{remark}[Quantifiers and frequency geometry]
The inequalities in \eqref{eq:anisotropic-feasible-region} describe a
sufficient region for this construction, not an optimality theorem.
They allow a single width arbitrarily close to rescaled exponent
$1/10$, while the other widths remain smaller and the sum budget is
respected. They do not give the full isotropic ball with that exponent.
All constants and the threshold in $n$ may depend on the fixed vector
$\boldsymbol\theta$ and the fixed orders. No limit of growing Taylor
orders is taken. The scalar covariance and induced Cesaro conclusions
of earlier sections are unchanged.
\end{remark}

\subsection{A concrete extension in the flight-time frequency}
Choose, in the physical coordinate order above,
\begin{equation}\label{eq:roof-box-concrete-scales}
 \boldsymbol\theta=\left(\frac1{100},\frac1{100},
                              \frac1{100},\frac9{100}\right),
 \quad \alpha=\frac{19}{100},\quad P=16,\quad Q=29,\quad\beta=20.
\end{equation}
Write $\mathcal B_n^{\rm roof}$ for the resulting rescaled box.
Its physical support is
\begin{equation}\label{eq:roof-box-physical-support}
 |u_1|,|u_2|,|s|\le2n^{-49/100},\qquad |b|\le2n^{-41/100}.
\end{equation}
The rescaled half-widths are $2n^{1/100}$ and $2n^{9/100}$.

\begin{corollary}[A fixed-count roof-direction extension]
\label{cor:roof-box-fixed-count}
For the same radius, count, mark and insertion uniformity,
\begin{align}
 &n^2\int_{\substack{z\in n^{-1/2}\mathcal B_n^{\rm roof}\\
                            |z|\ge2n^{-99/200}}}
         |\Phi^{[k],a}_{n,R}(z)|\dd z
   \le C\left(M_an^{-1/40}+V_an^{-497/25}\right).
 \label{eq:roof-box-fixed-count}
\end{align}
The marked central comparison on $\mathcal B_n^{\rm roof}$ has the
original error on the right of \eqref{eq:anisotropic-central-band}.
The same central error, with a changed constant, holds on the union
\begin{equation}\label{eq:roof-box-union}
 \mathcal U_n=\mathcal B_n^{\rm roof}
                    \cup\{v:|v|\le2n^{67/1400}\}.
\end{equation}
Thus the old isotropic band is retained and genuinely new frequencies
beyond it are covered.
\end{corollary}
\begin{proof}
Here $\Theta=3/25$ and $\theta_*=9/100$. The distinct error
budgets in Theorem~\ref{thm:anisotropic-fixed-count-band} are
\begin{align}
 1/2-\theta_*-2\alpha&=3/100,&
 (Q-1)(1/2-\theta_*)-2\alpha(Q+1)&=2/25,\notag\\
 \beta-\Theta&=497/25,&
 \beta-\Theta-(2P-1)(1/2+\theta_*)&=159/100,\notag\\
 P(\alpha-2\theta_*)-\Theta&=1/25,&
 1/4-\Theta-\theta_*&=1/25.
 \label{eq:roof-box-exponent-ledger}
\end{align}
The residual with $b=P$ has the factor
$n^{-1/25}[\log(2+n)]^{16}$; every other residual term has a
strictly faster power, with its own fixed logarithm. They are all
$O(n^{-1/40})$. The residual Taylor degree is $31$ and the spectral
degree is $29$; the two numbers refer to different expansions.
The theorem proves \eqref{eq:roof-box-fixed-count}, and
$1/40>3/280$ gives the central estimate. For the union, add the
nonnegative integrals over the two sets and use
Theorem~\ref{thm:rate-preserving-annulus} for the isotropic ball.
The new roof exponent satisfies $9/100>67/1400$, so the difference
of the two regions is nonempty for all sufficiently large $n$.
\end{proof}

\begin{corollary}[The unchanged marked-state event on the union]
\label{cor:anisotropic-same-event}
Let $A_{n,k,R}$ be the same single marked-state event as in
Corollary~\ref{cor:rate-preserving-same-event}, with probability
$p_{n,R}>0$ and zero-extension variation $V_{n,R}$. Then
\begin{align}
 &\int_{\mathcal U_n}
 \left|\mathbb E_{\nu_R^*}
  [e^{iv\cdot U_{n,R}/\sqrt n}\mid A_{n,k,R}]
                        -e^{-v^{\mathsf T}D_Rv/2}\right|\dd v\notag\\
 &\hspace{15mm}\le\frac C{p_{n,R}}
       \left(n^{-3/280}\sqrt{\log(2+n)}+V_{n,R}n^{-9/175}\right).
 \label{eq:anisotropic-same-event}
\end{align}
The sufficient conditions $p_{n,R}\ge cn^{-\beta_0}$,
$V_{n,R}\le Cn^{\kappa_0}$ remain
$\beta_0<3/280$ and $\beta_0+\kappa_0<9/175$.
\end{corollary}
\begin{proof}
Insert $a=c_*^{-1}\one_{E_{n,R}}$, whose mass is exactly $p_{n,R}$,
and divide the union estimate by that same probability. The event,
its lattice coordinates, and its denominator are not replaced. The
all-moment statements retain their separate conditions with
$\rho_d=1/2$ in odd degree and $\rho_d=1$ in even degree.
\end{proof}

\subsection{Raw inversion with a coordinatewise kernel}
Fix a real even $\chi_0\in C_c^\infty(\R)$ with $0\le\chi_0\le1$,
$\chi_0=1$ on $[-1,1]$ and support in $[-2,2]$. With the widths in
\eqref{eq:roof-box-concrete-scales}, set
\begin{equation}\label{eq:roof-box-kernel}
 \chi_n^{\rm box}(\omega)=\prod_{i=1}^4\chi_0(\omega_i/B_{n,i}),
 \qquad K_n^{\rm box}=\mathcal F^{-1}\chi_n^{\rm box}.
\end{equation}
The inverse is on $\Z^3\times\R$ with the inherited Fourier
normalization. This kernel uses the box alone; the transform estimate
on the larger union does not require silently identifying two kernels.

For a marked insertion which also satisfies the finite-record
admissibility of Definition~\ref{def:finite-record-weight}, put
$w_{n,k,R}=c_*a\circ(F_R^*)^k$ and use the exact extraction
$\mu^{a,L}=E^{a,L}+Q^{a,L}$. Define the active-kernel terms by
\begin{align}
 \mathcal E^{{\rm box},a,L}_{n,k,R,A}
 &=\mathop{\rm ess\,sup}_{\mathcal W_{n,R,A}}
                     |e^{a,L}-K_n^{\rm box}*E^{a,L}|,\notag\\
 \mathcal C^{{\rm box},a,L}_{n,k,R}(B)
 &=\int_{\T^3\times[-B,B]}|1-\chi_n^{\rm box}(\omega)|
                              |\widehat Q^{a,L}(\omega)|\dd\omega.
 \label{eq:roof-box-raw-terms}
\end{align}
The local edge quantity is allowed to be infinite before its germs
are bounded. Bounded variation alone does not imply finite-record
admissibility or a uniform density-derivative estimate.

\begin{corollary}[Exact marked raw budget at separate coordinate scales]
\label{cor:roof-box-raw-budget}
Let $L_n=\lceil\lambda n\rceil$ with $\lambda>c_*^{-1}+1$.
For $B\ge1$, $P_0>0$, and
$y=((\ell,t)-n\bar G_R)/\sqrt n$, one has
\begin{align}
 &\mathop{\rm ess\,sup}_{\mathcal W_{n,R,A}}
      |n^2p_{n,R}^{[k],a}(\ell,t)-\alpha_a\mathfrak g_R(y)|\notag\\
 &\quad\le C\left(M_an^{-3/280}\sqrt{\log(2+n)}+V_an^{-9/175}\right)
      +CM_ae^{-cn^{1/50}}+C_{A,\lambda,P_0}M_an^{-P_0}\notag\\
 &\qquad+n^2\mathcal E^{{\rm box},a,L_n}_{n,k,R,A}
   +(2\pi)^{-4}n^2\mathcal C^{{\rm box},a,L_n}_{n,k,R}(B)
   +\frac{n^2}{\pi B}A_2(n,L_n,R,w_{n,k,R}).
 \label{eq:roof-box-raw-budget}
\end{align}
This is an extended-real inequality where necessary, with physical
cutoffs \eqref{eq:roof-box-physical-support} and rescaled cutoffs
$2n^{1/100}$, $2n^{9/100}$.
\end{corollary}
\begin{proof}
For all sufficiently large $n$, the lattice-frequency support lies
strictly inside the fundamental torus cube. Hence Fourier inversion
and one-dimensional rescaling give the product-kernel estimate
\begin{equation}\label{eq:box-kernel-Schwartz}
 |K_n^{\rm box}(x)|\le C_J\prod_{i=1}^4 B_{n,i}
                  \prod_{i=1}^4(1+B_{n,i}|x_i|)^{-J}
 \qquad(x\in\Z^3\times\R).
\end{equation}
It follows directly from the Schwartz decay of the inverse of $\chi_0$;
no lattice coordinate is approximated by a continuous variable.

Let $\mathsf T^{a,L_n}$ be the high-count part of the marked measure.
On the actual central count labels it has no density contribution,
and the count coordinate of its support is separated from these
labels by at least $c_{A,\lambda}n$ for large $n$. Its total variation
is at most $c_*M_a$. Since the third coordinate is the collision
count, \eqref{eq:box-kernel-Schwartz} gives
\begin{align}
 n^2\sup_{\mathcal W_{n,R,A}}
             |K_n^{\rm box}*\mathsf T^{a,L_n}|
 &\le C_{A,\lambda,J}M_a n^{3/25-51J/100}.
 \label{eq:box-count-separation}
\end{align}
Here $n^2\prod_i B_{n,i}=n^{3/25}$ and
$nB_{n,3}=n^{51/100}$. Select the fixed Schwartz order $J$ so that
the bound is $O(M_an^{-P_0})$.

The exact finite-count decomposition and support identity now give
\begin{align*}
 p_{n,R}^{[k],a}={}&K_n^{\rm box}*\mu_{n,R}^{[k],a}
 +(e^{a,L_n}-K_n^{\rm box}*E^{a,L_n})\\
 &+\mathcal F^{-1}[(1-\chi_n^{\rm box})\widehat Q^{a,L_n}]
       -K_n^{\rm box}*\mathsf T^{a,L_n}.
\end{align*}
For the first term use the central comparison on the box, whose
cutoff is between zero and one. Missing Gaussian frequencies have
at least one $|v_i|\ge n^{\theta_i}$ and hence
$|v|\ge n^{1/100}$. Uniform ellipticity bounds their integral by
$Ce^{-cn^{1/50}}$. The finite-band residual inverse and its far-roof
part are bounded exactly as in
Corollary~\ref{cor:rate-preserving-raw-budget}, giving the last two
terms of \eqref{eq:roof-box-raw-budget}. The local combined edge
correction is the one in \eqref{eq:roof-box-raw-terms}, not a correction
for either preceding isotropic kernel. This proves the inequality.
\end{proof}

\paragraph{Position in the full raw argument.}
The new theorem enlarges the proved prescribed-count domain in a
specific, four-dimensionally integrated way, including flight-time
frequencies of size $n^{-41/100}$. The previously proved isotropic
ball is preserved by \eqref{eq:roof-box-union}. Neither set contains
the full ball of physical radius $n^{-2/5}$: the transverse widths,
compact nonzero torus frequencies, full peripheral regimes and growing
roof ranges require additional estimates. The terms involving
$\mathcal E^{\rm box}$, $\mathcal C^{\rm box}$ and the true long-time
$A_2$ in \eqref{eq:roof-box-raw-budget} are not bounded by the high-order
collision moment lemma. Weighted denominator asymptotics and relative
physical-event replacement remain separate steps in the same raw
mixed-density problem.

\section{Shape-adaptive Fourier cancellation at a prescribed count}
\label{sec:shape-adaptive-fixed-count}
The preceding box theorem assigns one fixed exponent to each coordinate.
Here the frequency region may depend on $n$ without being a box. A
weighted-volume bound replaces the sum of its four width exponents.
We then construct a single nonrectangular region reaching all four
coordinate directions. The logarithmic factor in its definition pays
for the entire dyadic sum; it does not enter the final central rate.

All integrals in this section are in the rescaled variable $v=\sqrt n z$.
The coordinate order remains $(K_1,K_2,N,T)$, with dual physical
coordinates $(u_1,u_2,s,b)$. Put
\begin{equation}\label{eq:shape-fixed-constants}
 \tau_0=\frac1{200},\qquad \theta_0=\frac9{100},\qquad
 s_0=\frac{67}{280},\qquad d_0'=s_0-\tau_0=\frac{41}{175},\qquad
 \eta_0'=\frac14-s_0=\frac3{280}.
\end{equation}
The symbols $d_0'$ and $\eta_0'$ in this section are exponents, not
covariance lower bounds or phase-separation constants.

\subsection{A measurable-region estimate with fixed expansion orders}
\begin{theorem}[A weighted-volume criterion]
\label{thm:shape-volume-criterion}
Fix $C_0<\infty$. For each $n$, let $\Omega_n\subset\R^4$ be any
measurable set such that
\begin{equation}\label{eq:shape-volume-hypotheses}
 \sup_{v\in\Omega_n}|v|\le C_0 n^{\theta_0},\qquad
 |\Omega_n|\le C_0 n^{d_0'},\qquad
 \int_{\Omega_n}|v|\dd v\le C_0 n^{s_0}.
\end{equation}
The supremum may be interpreted essentially. Uniformly in these sets,
$R\in I$, the prescribed count $n$, the mark $0\le k\le n$, and
bounded-BV section insertions $a$, for all sufficiently large $n$,
\begin{align}
 &n^2\int_{\substack{z\in n^{-1/2}\Omega_n\\
                             |z|\ge2n^{-99/200}}}
           |\Phi^{[k],a}_{n,R}(z)|\dd z
 \le C_{C_0}\left(M_a n^{-3/280}+V_a n^{-5559/175}\right).
 \label{eq:shape-volume-cancellation}
\end{align}
No regularity of $\partial\Omega_n$, and no fixed choice of width
exponents for $\Omega_n$, is required.
\end{theorem}
\begin{proof}
Use the exact marked recentering and
\eqref{eq:quarter-characteristic-stopping}. Its integral is at most
$CM_a n^{-1/4}\int_{\Omega_n}|v|\dd v
\le CM_a n^{-3/280}$. The deterministic collision length is
$m=m_k+m_{n-k}=n/c_*+O(1)$ uniformly in $k$.

Apply Theorem~\ref{thm:high-order-damped-unsmoothing} with the fixed choices
\begin{equation}\label{eq:shape-expansion-orders}
 P=26,\qquad Q=29,\qquad
 \delta=\tfrac14n^{-19/100},\qquad
 \epsilon=\tfrac14n^{-32}.
\end{equation}
The residual Taylor degree is $51$ and the spectral degree is $29$.
For $z=v/\sqrt n$ the two analytic margins are $3/100$ and $2/25$,
exactly as for the roof box. The harmless fixed constant $C_0$ is
absorbed by increasing the threshold in $n$. On
$|v|\ge2n^{1/200}$, the damped term is at most a fixed polynomial in
$n$ times $e^{-cn^{1/100}}$, including its $52$ multiplier factors.
Its integral is smaller than every fixed inverse power of $n$.

The fine insertion loss has exponent
$32-d_0'=5559/175$. The fine residual replacement has margin
\begin{equation}\label{eq:shape-fine-margin}
 32-d_0'-51(1/2+\theta_0)=\frac{1173}{700}.
\end{equation}
For the residual term with $b$ nonsingleton blocks, $1\le b\le26$,
\eqref{eq:high-order-damped-unsmoothing} and the volume bound give
\[
 C M_a n^{d_0'+52\theta_0-26+(81/100)b}
                         [\log(2+n)]^{52-b}.
\]
The slowest term has $b=26$ and decay exponent
\begin{equation}\label{eq:shape-residual-margin}
 26(19/100-2\theta_0)-d_0'=\frac9{350}
                 =\frac3{280}+\frac3{200}.
\end{equation}
Every smaller block count improves the power by $81/100$. Thus all
fixed logarithms are absorbed by the strict margin $3/200$.
The $b=1$ contribution is included in this calculation. Finally
$\dd v=n^2\dd z$ and the centering phase has modulus one. This proves
\eqref{eq:shape-volume-cancellation} with constants independent of the
shape of $\Omega_n$. Neither expansion order grows with $n$.
\end{proof}

\subsection{A summable family of coordinate scales}
For $\mathbf j=(j_1,j_2,j_3,j_4)\in\Z_{\ge0}^4$, write
$S(\mathbf j)=\sum_i j_i$ and $J(\mathbf j)=\max_i j_i$.
The following elementary counting fact will be used both for frequency
volume and for the physical convolution kernel.

\begin{lemma}[The weighted dyadic sum]
\label{lem:weighted-dyadic-cross}
For each integer $K\ge0$,
\begin{equation}\label{eq:weighted-dyadic-cross}
 \sum_{S(\mathbf j)+J(\mathbf j)\le K}
                       2^{S(\mathbf j)+J(\mathbf j)}
             \le C2^K(K+1)^3.
\end{equation}
Consequently, if $a>0$, $T\ge a^5$ and $U\ge a$, the finite set
\[
 \mathcal D(a,T,U)=\{\mathbf j:
     a^5 2^{S(\mathbf j)+J(\mathbf j)}\le T,
     \ a2^{J(\mathbf j)}\le U\}
\]
satisfies, with $\lambda_i=a2^{j_i}$ and
$\lambda_*=\max_i\lambda_i$,
\begin{align}
 \sum_{\mathbf j\in\mathcal D}\prod_i\lambda_i\lambda_*
    &\le CT\,[1+\log(T/a^5)]^3,\notag\\
 \sum_{\mathbf j\in\mathcal D}\prod_i\lambda_i
    &\le C(T/a)\,[1+\log(T/a^5)]^3.
 \label{eq:dyadic-cross-scale-sums}
\end{align}
\end{lemma}
\begin{proof}
Fix $t=S+J$ and choose one coordinate attaining the maximum $J$.
For this choice the other three coordinates sum to $t-2J$.
Ignoring their upper bound by $J$ only increases their number, which
is at most $(t+2)^2$. There are at most $t+1$ choices for $J$ and
four choices of its attaining coordinate. Thus the number of tuples
with $S+J=t$ is at most $C(t+1)^3$. Overcounting tied maxima is
harmless. Sum $2^tC(t+1)^3$ for $0\le t\le K$ and use
$\sum_{t=0}^K2^t\le2^{K+1}$. This proves the first assertion.
Take $K=\lfloor\log_2(T/a^5)\rfloor$ and multiply by $a^5$.
The bound on $U$ only removes terms. The second scale sum follows
from $\lambda_*\ge a$.
\end{proof}

Set
\begin{equation}\label{eq:adaptive-cross-scales}
 a_n=2n^{1/200},\quad U_n=n^{9/100},\quad
 T_n=\frac{n^{67/280}}{[\log(2+n)]^3},\quad
 \mathcal D_n=\mathcal D(a_n,T_n,U_n).
\end{equation}
All assertions concerning these sets are for sufficiently large $n$,
so $T_n\ge a_n^5$ and $U_n\ge a_n$. Define the closed union of boxes
\begin{equation}\label{eq:adaptive-cross-support}
 \mathcal H_n=\bigcup_{\mathbf j\in\mathcal D_n}
                       \prod_{i=1}^4[-2\lambda_i,2\lambda_i],
 \qquad \lambda_i=a_n2^{j_i}.
\end{equation}
This is a nonrectangular frequency region. A useful description uses
$x_i(v)=\max\{a_n,|v_i|\}$ and $x_*(v)=\max_i x_i(v)$.
It contains
\begin{equation}\label{eq:adaptive-cross-inner}
 \left\{v:x_*(v)\le U_n/2,\quad
                  x_*(v)\prod_i x_i(v)\le T_n/32\right\},
\end{equation}
and is contained in the same product region with thresholds $2U_n$
and $32T_n$. These constants also allow a smooth cutoff below.

\begin{lemma}[Size and interior of the nonrectangular region]
\label{lem:adaptive-cross-volume}
The set $\mathcal H_n$ satisfies \eqref{eq:shape-volume-hypotheses}
with one fixed constant. It contains $\{|v|\le2n^{1/200}\}$.
The inclusions just stated hold for all sufficiently large $n$.
\end{lemma}
\begin{proof}
Every box has volume $256\prod_i\lambda_i$ and
$|v|\le4\lambda_*$ on it. The two scale sums in
\eqref{eq:dyadic-cross-scale-sums} bound its union volume and first
radial moment. Since $1+\log(T_n/a_n^5)\le C\log(2+n)$, the
factor $[\log(2+n)]^{-3}$ in $T_n$ cancels this entire counting
loss. This gives $|\mathcal H_n|\le Cn^{41/175}$ and
$\int_{\mathcal H_n}|v|\dd v\le Cn^{67/280}$. Its radius is at
most $4U_n$. The zero tuple supplies the central ball.

For the inner inclusion, select scales with
$x_i\le\lambda_i\le2x_i$, taking $j_i=0$ when $x_i=a_n$.
They obey $\lambda_*\prod_i\lambda_i\le32x_*\prod_i x_i\le T_n$
and $\lambda_*\le2x_*\le U_n$. This gives an admissible box
containing $v$. Conversely, membership in a box implies
$x_i\le2\lambda_i$, because $\lambda_i\ge a_n$, and hence the
stated outer bounds.
\end{proof}

\begin{theorem}[One retained-rate region in all four directions]
\label{thm:adaptive-cross-fixed-count}
The physical region is $n^{-1/2}\mathcal H_n$. On its part outside
$|z|\le2n^{-99/200}$,
\begin{align}
 n^2\int |\Phi^{[k],a}_{n,R}(z)|\dd z
       &\le C\left(M_an^{-3/280}+V_an^{-5559/175}\right).
 \label{eq:adaptive-cross-fixed-count}
\end{align}
Let
\begin{equation}\label{eq:adaptive-cross-retained-union}
 \mathcal U_n^\sharp=\mathcal H_n\cup\mathcal B_n^{\rm roof}
                         \cup\{|v|\le2n^{67/1400}\}.
\end{equation}
Then
\begin{align}
 &\int_{\mathcal U_n^\sharp}
 |C^{[k],a}_{n,R}(v)-\alpha_a e^{-v^{\mathsf T}D_Rv/2}|\dd v
 \le C\left(M_an^{-3/280}\sqrt{\log(2+n)}+V_an^{-9/175}\right).
 \label{eq:adaptive-cross-central}
\end{align}
Both estimates are uniform in $R,n,k,a$ in the original marked class.
Every previously proved isotropic and roof-box frequency is retained.
\end{theorem}
\begin{proof}
Theorem~\ref{thm:shape-volume-criterion} and
Lemma~\ref{lem:adaptive-cross-volume} give the first assertion. On
the central ball use Theorem~\ref{thm:marked-return-band}. On its
complement the uniformly positive covariance gives an exponentially
small Gaussian integral, while the first assertion gives the physical
transform integral. Finally add the inherited roof-box and isotropic
bounds. Their overlap can only increase the upper bound. No division
by the volume of any region is made.
\end{proof}

\begin{corollary}[New mixed directions and the unchanged event]
\label{cor:adaptive-cross-directions-event}
For a fixed exponent vector $\boldsymbol\rho$ with $\rho_i\ge0$, set
$\widetilde\rho_i=\max\{1/200,\rho_i\}$ and
$\widetilde\rho_* =\max_i\widetilde\rho_i$. If
\begin{equation}\label{eq:adaptive-cross-exponent-interior}
 \widetilde\rho_*<9/100,\qquad
 \sum_i\widetilde\rho_i+\widetilde\rho_*<67/280,
\end{equation}
then every fixed multiple of the box $|v_i|\le n^{\rho_i}$ is
contained in $\mathcal H_n$ for all sufficiently large $n$.
Also $|v_i|\le U_n/4$ on any one coordinate axis is contained in
$\mathcal H_n$. Thus each physical coordinate has a new axial reach
of order $n^{-41/100}$, not only the roof coordinate.

For the unchanged marked-state event with probability $p_{n,R}>0$
and variation $V_{n,R}$, the right side of
\eqref{eq:adaptive-cross-central} becomes
\[
 \frac C{p_{n,R}}\left(n^{-3/280}\sqrt{\log(2+n)}
                                    +V_{n,R}n^{-9/175}\right).
\]
Its probability and variation conditions are unchanged.
\end{corollary}
\begin{proof}
Apply the inner description \eqref{eq:adaptive-cross-inner}. The
strict power margins absorb both fixed constants and
$[\log(2+n)]^3$. On an axis the product exponent is
$2(9/100)+3/200=39/200<67/280$, so the same argument works at
the asserted constant multiple of $U_n$. As a mixed example,
$\boldsymbol\rho=(1/25,1/25,1/25,1/20)$ has sum plus maximum
$11/50<67/280$ and extends beyond the old isotropic exponent.
The event statement follows by taking $a=c_*^{-1}\one_E$ and
dividing by the exact probability, as in
Corollary~\ref{cor:anisotropic-same-event}.
\end{proof}

\paragraph{Extent of the result.}
This is not the isotropic ball of rescaled radius $n^{1/10}$.
For example, a vector with four coordinates of order $n^{3/50}$
is outside the product budget and outside both retained regions for
large $n$. The expansion orders are fixed. The new result controls a
single explicit nonrectangular domain, rather than a collection of
estimates whose constants might depend on $n$-varying exponents.
The next section connects this domain to the already used raw
extraction without imposing a new, unrelated edge hypothesis.

\section{A coherent raw correction under changes of cutoff}
\label{sec:coherent-raw-cutoff}
Increasing the low-frequency region changes both the extracted-edge
convolution and the residual inverse transform. Estimating their
changes separately loses an exact cancellation. We retain that
cancellation and prove a quantitative comparison with the existing
box and isotropic cutoffs, for the same finite extraction and the
same physical weight. This closes the change-of-cutoff contribution;
it does not assume smallness of the remaining common raw correction.

\subsection{A smooth cutoff for the nonrectangular region}
Choose the function $\chi_0$ in Section~\ref{sec:anisotropic-fixed-count-bands}
also nonincreasing on $[0,\infty)$. Define
\[
 \varphi_0(t)=\chi_0(t),\qquad
 \varphi_j(t)=\chi_0(t/2^j)-\chi_0(t/2^{j-1})\quad(j\ge1).
\]
These functions are nonnegative and sum to one on $\R$. With
$\mathcal D_n,a_n$ from \eqref{eq:adaptive-cross-scales}, put
\begin{equation}\label{eq:cross-smooth-cutoff}
 \Xi_n(v)=\sum_{\mathbf j\in\mathcal D_n}
                 \prod_{i=1}^4\varphi_{j_i}(v_i/a_n),\qquad
 \chi_n^\sharp(\omega)=\Xi_n(\sqrt n\,\omega),\qquad
 K_n^\sharp=\mathcal F^{-1}\chi_n^\sharp.
\end{equation}
For large $n$, the first three physical supports are strictly inside
the chosen torus chart; the expression is extended by zero there.
Its largest coordinate support is $2n^{-41/100}$ and its central
plateau contains the physical ball $|\omega|\le2n^{-99/200}$.
The corresponding rescaled scales are $2n^{9/100}$ and $2n^{1/200}$.

\begin{lemma}[Cutoff plateau and count-coordinate kernel tail]
\label{lem:cross-kernel-tail}
The cutoff satisfies $0\le\Xi_n\le1$, is supported in
$\mathcal H_n$, and is one on the central ball and on the inner
product region \eqref{eq:adaptive-cross-inner}. For every fixed
integer $J\ge0$,
\begin{equation}\label{eq:cross-kernel-Schwartz}
 |K_n^\sharp(x)|\le C_J\sum_{\mathbf j\in\mathcal D_n}
     \prod_{i=1}^4 B_{\mathbf j,i}
          (1+B_{\mathbf j,i}|x_i|)^{-J},\qquad
 B_{\mathbf j,i}=\lambda_i/\sqrt n,
 \quad x\in\Z^3\times\R.
\end{equation}
In particular, for $L_n=\lceil\lambda n\rceil$ with
$\lambda>c_*^{-1}+1$, and the high-count part
$\mathsf T^{a,L_n}$ of the original marked law,
\begin{align}
 n^2\sup_{\mathcal W_{n,R,A}}
        |K_n^\sharp*\mathsf T^{a,L_n}|
 \le C_{A,\lambda,J}M_a n^{41/175-101J/200}.
 \label{eq:cross-count-tail}
\end{align}
The third coordinate, not the largest roof scale, supplies this decay.
\end{lemma}
\begin{proof}
The nonnegative full product partition sums to one, so a subsum lies
between zero and one. Each summand is supported in its box from
\eqref{eq:adaptive-cross-support}. If a summand is nonzero at $v$,
then $\lambda_i\le2\max\{a_n,|v_i|\}$ for every $i$; this is
immediate for $j_i=0$, and follows from
$|v_i|>\lambda_i/2$ for $j_i>0$. At a point of the inner product
region all nonzero summands therefore belong to $\mathcal D_n$,
which proves the plateau. On the central ball only the zero tuple
can contribute and its value is one.

For $j\ge1$, write $\varphi_j(v/a_n)$ as
$\chi_0(v/\lambda_j)-\chi_0(2v/\lambda_j)$. Thus the inverse of
every factor is a dilation of one of two fixed Schwartz functions.
The mixed inverse transform, with the three integer output coordinates
left discrete, gives \eqref{eq:cross-kernel-Schwartz}. Its constants
are uniform over all dyadic indices. The sum is finite at every $n$.

The total variation of the high-count marked measure is at most
$c_*M_a$. On a central window its third-coordinate support is
separated by at least $c_{A,\lambda}n$, by
Lemma~\ref{lem:central-count-support}. Since every
$B_{\mathbf j,3}\ge a_n/\sqrt n=2n^{-99/200}$, each count factor
has decay at most $C n^{-101J/200}$. Multiply the remaining scale
sum by $n^2$ and use \eqref{eq:dyadic-cross-scale-sums}; the result is
$Cn^{41/175}$. This proves \eqref{eq:cross-count-tail}. Arbitrary
inverse powers follow by choosing a fixed $J$ large enough.
\end{proof}

\subsection{The exact signed correction and its cancellation}
Require now the additional finite-record admissibility of the weight
$w_{n,k,R}=c_*a\circ(F_R^*)^k$ from
Definition~\ref{def:finite-record-weight}. Fix one extraction
\[
 \mu^{a,L}=E^{a,L}+Q^{a,L},\qquad
 p^{a,L}=e^{a,L}+r^{a,L},\qquad
 \widehat Q^{a,L}\in L^1(\T^3\times\R).
\]
All cutoffs in this subsection use precisely these same measures.
For a real smooth compactly supported frequency cutoff $\chi$,
with kernel $K_\chi=\mathcal F^{-1}\chi$, define
\begin{equation}\label{eq:coherent-raw-correction}
 \mathcal R_\chi^{a,L}
 =e^{a,L}-K_\chi*E^{a,L}
       +\mathcal F^{-1}[(1-\chi)\widehat Q^{a,L}].
\end{equation}
This is a signed or complex function defined almost everywhere; it
is not the sum of the absolute edge and residual bounds. Its essential
supremum may be infinite. The residual inverse is absolutely convergent,
and the other convolutions are well defined for finite measures.

\begin{proposition}[Exact transport of a raw correction]
\label{prop:exact-raw-cutoff-transport}
For two such cutoffs $\chi,\psi$, almost everywhere on the mixed
space,
\begin{align}
 \mathcal R_\chi^{a,L}&=p^{a,L}-K_\chi*\mu^{a,L},\notag\\
 \mathcal R_\chi^{a,L}-\mathcal R_\psi^{a,L}
             &=(K_\psi-K_\chi)*\mu^{a,L}.
 \label{eq:exact-raw-cutoff-transport}
\end{align}
In particular their difference has a bounded continuous representative,
and
\begin{equation}\label{eq:raw-cutoff-global-comparison}
 \|\mathcal R_\chi^{a,L}-\mathcal R_\psi^{a,L}\|_\infty
 \le(2\pi)^{-4}\int |\chi-\psi|\,|\widehat\mu^{a,L}|\dd\omega.
\end{equation}
The identities hold for any exact finite extraction with the stated
integrable residual transform. They introduce no estimate for $A_2$.
\end{proposition}
\begin{proof}
Absolute residual inversion gives
$\mathcal F^{-1}[(1-\chi)\widehat Q]=r-K_\chi*Q$.
Substitute this in \eqref{eq:coherent-raw-correction} and use
$E+Q=\mu^{a,L}$. Subtract the two identities. In frequency form the
edge change is $(\psi-\chi)\widehat E$ and the residual change is
$(\psi-\chi)\widehat Q$; their sum is
$(\psi-\chi)\widehat\mu^{a,L}$. This is the cancellation that is
lost by taking absolute values of the two changes first. The last
integrand is integrable, since the cutoff difference has compact
support and the measure transform is bounded. Mixed Fourier inversion
therefore proves the norm bound and continuity. The equalities are
almost-everywhere equalities of actual densities, not a subtraction
of two extended-real essential suprema.
\end{proof}

\subsection{Comparison with the already used cutoffs}
Retain $\chi_n^{\rm box},K_n^{\rm box}$ from
\eqref{eq:roof-box-kernel} and
$\chi_n^\natural,K_n^\natural$ from
Section~\ref{sec:rate-preserving-band}. Each is one on the original
central ball for sufficiently large $n$. Define
\begin{equation}\label{eq:raw-cutoff-rate}
 \Delta_n(a)=M_an^{-3/280}+V_an^{-983/350}.
\end{equation}
This exponent for the variation is the slowest of the three
complementary-band estimates, not the weaker central-ball exponent.

\begin{theorem}[A vanishing complete change-of-cutoff contribution]
\label{thm:coherent-raw-cutoff-rate}
For $\chi,\psi$ any two of
$\chi_n^\sharp,\chi_n^{\rm box},\chi_n^\natural$,
$L_n=\lceil\lambda n\rceil$, $\lambda>c_*^{-1}+1$, and fixed
$A,P_0<\infty$ with $P_0>0$,
\begin{align}
 n^2\mathop{\rm ess\,sup}_{\mathcal W_{n,R,A}}
   |\mathcal R_\chi^{a,L_n}-\mathcal R_\psi^{a,L_n}|
       \le C\Delta_n(a)+C_{A,\lambda,P_0}M_an^{-P_0}.
 \label{eq:coherent-raw-cutoff-rate}
\end{align}
The constant is uniform in $R,n,k,a$ in the common weight class.
If in addition $c\lambda>a$ for the constants in
\eqref{eq:count-tail-for-inversion}, the comparison is global, with
an exponentially small tail in place of the last term.
\end{theorem}
\begin{proof}
The cutoff difference vanishes on $|z|\le2n^{-99/200}$ and is
supported in $n^{-1/2}\mathcal U_n^\sharp$. It has modulus at most
one. Equations~\eqref{eq:adaptive-cross-fixed-count},
\eqref{eq:roof-box-fixed-count}, and
\eqref{eq:rate-preserving-annulus} therefore imply
\[
 (2\pi)^{-4}n^2\int |\chi-\psi|\,
                         |\widehat\mu^{[k],a}_{n,R}|\dd z
                       \le C\Delta_n(a).
\]
Here $1/40>3/280$, and the two other variation exponents are larger
than $983/350$. This is an estimate for the full law at the prescribed
return count, not for an averaged coefficient.

In \eqref{eq:exact-raw-cutoff-transport}, write
$\mu^{a,L_n}=\mu^{[k],a}_{n,R}-\mathsf T^{a,L_n}$. The full-law
part has the global bound just proved. The remaining convolution
is bounded on the central window by the two relevant count-separation
bounds: \eqref{eq:cross-count-tail}, \eqref{eq:box-count-separation},
or the isotropic bound in the proof of
Corollary~\ref{cor:rate-preserving-raw-budget}. Choose their fixed
Schwartz orders for $P_0$. This proves the first assertion without
assuming $c\lambda>a$. With that stronger condition the tail total
variation is exponentially small, and all three kernel suprema are
polynomially bounded. Their products are exponentially small globally.
\end{proof}

\begin{corollary}[One common local raw remainder]
\label{cor:one-common-raw-remainder}
Let $\mathcal A_n$ be a class of common admissible weights and marks
for which $\sup_{a\in\mathcal A_n}\Delta_n(a)\to0$ and $M_a$ has
at most fixed polynomial growth. Then
\[
 \sup_{R,k,a\in\mathcal A_n}n^2
        \|\mathcal R_\chi^{a,L_n}\|_{L^\infty(\mathcal W_{n,R,A})}
                   \longrightarrow0
\]
holds for one of the three cutoffs if and only if it holds for all
three. Finite values of the norms likewise transfer, and the difference
of any two finite norms is at most the bound in
\eqref{eq:coherent-raw-cutoff-rate}. For a fixed bounded-BV weight
class the rate is $O(n^{-3/280})$.
\end{corollary}
\begin{proof}
Apply the triangle inequality to the almost-everywhere identity in
Theorem~\ref{thm:coherent-raw-cutoff-rate} in both directions, choosing
$P_0$ larger than the polynomial growth exponent of $M_a$. If one norm
is finite, the bounded difference makes the other finite. No difference
of infinite norms is formed. The uniform convergence statement follows.
\end{proof}

The corollary does not transfer smallness of the isolated edge term or
of the absolute complementary integral. Cancellation between them is
retained in $\mathcal R$. In particular a small value of their signed
sum cannot be used to bound either of the two absolute quantities.
This distinction permits one fixed raw problem to be used throughout
the expansion of the known Fourier region.

\subsection{The exact raw budget at the new cutoff}
For the new kernel define, immediately before use,
\begin{align}
 \mathcal E^{\sharp,a,L}_{n,k,R,A}
 &=\mathop{\rm ess\,sup}_{\mathcal W_{n,R,A}}
                         |e^{a,L}-K_n^\sharp*E^{a,L}|,\notag\\
 \mathcal C^{\sharp,a,L}_{n,k,R}(B)
 &=\int_{\T^3\times[-B,B]}
           |1-\chi_n^\sharp|\,|\widehat Q^{a,L}|\dd\omega.
 \label{eq:cross-raw-terms}
\end{align}
The first quantity is allowed to be infinite. The exact identity on
the central count labels is
\begin{equation}\label{eq:cross-raw-exact-identity}
 p_{n,R}^{[k],a}=K_n^\sharp*\mu_{n,R}^{[k],a}
                  +\mathcal R_{\chi_n^\sharp}^{a,L_n}
                  -K_n^\sharp*\mathsf T^{a,L_n}.
\end{equation}
It follows from central count support and the first identity in
Proposition~\ref{prop:exact-raw-cutoff-transport}.

\begin{corollary}[Raw inversion with the coherent cutoff]
\label{cor:cross-raw-budget}
For $B\ge1$, $P_0>0$, and
$y=((\ell,t)-n\bar G_R)/\sqrt n$,
\begin{align}
 &\mathop{\rm ess\,sup}_{\mathcal W_{n,R,A}}
 |n^2p_{n,R}^{[k],a}(\ell,t)-\alpha_a\mathfrak g_R(y)|\notag\\
 &\quad\le C\left(M_an^{-3/280}\sqrt{\log(2+n)}
                              +V_an^{-9/175}\right)
        +CM_ae^{-cn^{1/100}}+C_{A,\lambda,P_0}M_an^{-P_0}\notag\\
 &\qquad+n^2\mathcal E^{\sharp,a,L_n}_{n,k,R,A}
       +(2\pi)^{-4}n^2\mathcal C^{\sharp,a,L_n}_{n,k,R}(B)
       +\frac{n^2}{\pi B}A_2(n,L_n,R,w_{n,k,R}).
 \label{eq:cross-raw-budget}
\end{align}
The inequality has its extended-real meaning if necessary.
\end{corollary}
\begin{proof}
Apply \eqref{eq:adaptive-cross-central} to the first term of
\eqref{eq:cross-raw-exact-identity}, using $0\le\Xi_n\le1$.
The omitted Gaussian frequencies are outside the plateau containing
$|v|\le2n^{1/200}$; uniform ellipticity gives the exponential term.
Use \eqref{eq:cross-count-tail} for the high-count convolution.
Bound the two parts of the coherent correction separately only now:
its edge term is the one in \eqref{eq:cross-raw-terms}. Splitting the
absolutely convergent residual inverse at $|b|=B$ and using
\eqref{eq:finite-count-far-tail} gives the last two terms. This keeps
the true derivative sum at $L_n\asymp n$ unchanged.
\end{proof}

\paragraph{The remaining absolute and conditional estimates.}
Theorem~\ref{thm:coherent-raw-cutoff-rate} is an unconditional estimate
for the entire signed change of the raw correction, and
Corollary~\ref{cor:one-common-raw-remainder} makes that comparison
independent of which of the three cutoffs is used. It supplies neither
an absolute value for the common correction nor a bound on the genuine
long-time $A_2$. Frequencies outside the retained union, the common
local raw remainder, weighted raw denominators and relative physical
event replacement still require estimates for the same original law.
The last comparison concerns the same indicator on both sides. No
physical observation event has been substituted.

\section{Unsmoothed windows and an averaged common raw remainder}
\label{sec:unsmoothed-window-remainders}
The transport identity compares two raw corrections but does not bound
one of them. Positivity gives a different estimate: it bounds the
integral of the common correction over an actual observation window.
The window is not convolved with noise. Its side lengths must exceed
the resolution of a proved Fourier region. We keep this restriction
explicit, in particular in the three integer coordinates.

Put $X_{n,R}=U_{n,R}/\sqrt n$, $U_{n,R}=J_{n,R}-n\bar G_R$, and set
\begin{equation}\label{eq:window-Fourier-error}
 e=\frac3{280},\quad v_0=\frac9{175},\qquad
 \mathcal E_n(a)=M_an^{-e}\sqrt{\log(2+n)}+V_an^{-v_0}.
\end{equation}
The centered pushforward of the marked measure in
\eqref{eq:marked-measure} is denoted by $\lambda_{n,R}^{[k],a}$.
Its mass is $\alpha_a$. Theorem~\ref{thm:adaptive-cross-fixed-count}
gives the integrated characteristic error $C\mathcal E_n(a)$ on
$\mathcal U_n^\sharp$. Frequencies here are rescaled; a frequency $v$
corresponds to the physical frequency $v/\sqrt n$.

\subsection{Order-preserving interval approximation}
We give the approximation argument, including its values at endpoints.
These values matter because three coordinates of the record are discrete.
The construction is the classical Beurling--Selberg interval construction;
see \cite[Introduction]{CVExt}. Only the elementary facts proved below
are used, not an extremality assertion.

\begin{lemma}[Bandlimited envelopes, including endpoints]
\label{lem:window-interval-envelopes}
For an interval $I$ of length $l>0$ and an angular bandwidth $b>0$,
there are real continuous integrable functions $L_I,U_I$ with Fourier
support in $[-b,b]$ such that
\begin{equation}\label{eq:interval-envelopes}
 L_I\le\one_I\le U_I,\quad 0\le U_I\le3,\quad
 \int U_I=l+\frac{2\pi}{b},\quad
 \int L_I=l-\frac{2\pi}{b}.
\end{equation}
Either convention at either endpoint is allowed in the same inequalities.
The difference $E_I=U_I-L_I$ is nonnegative and has integral $4\pi/b$.
The Fourier convention in this statement is
$\widehat f(v)=\int e^{-ivx}f(x)\dd x$.
\end{lemma}
\begin{proof}
Use the continuous extensions of
\[
 K(x)=\left(\frac{\sin\pi x}{\pi x}\right)^2,\qquad
 H(x)=\sum_{m\in\Z}\operatorname{sgn}(m)K(x-m)+2xK(x),
 \quad \operatorname{sgn}(0)=0.
\]
The identity $\sum_{m\in\Z}K(x-m)=1$ follows from the partial-fraction
identity for $\pi^2/\sin^2\pi x$; the limiting integer cases agree.
For $x>0$ not an integer,
\[
 1-H(x)=K(x)+2\left(\frac{\sin\pi x}{\pi}\right)^2
                  \left(\sum_{j\ge1}\frac1{(x+j)^2}-\frac1x\right).
\]
The decreasing-integrand comparisons give
$1/(x+1)\le\sum_{j\ge1}(x+j)^{-2}\le1/x$.
Since $1/(x+1)\ge1/x-1/x^2$, it follows that
$|1-H(x)|\le K(x)$. Oddness proves
$|\operatorname{sgn}(x)-H(x)|\le K(x)$ everywhere.
Moreover $H-\operatorname{sgn}$ is an odd integrable function.

In the Fourier convention $e^{-2\pi i\xi x}$, $K$ has transform
$(1-|\xi|)_+$: write the sinc as the inverse transform of an interval
and take its square. Each translate of $K$, and $xK$ as a tempered
distribution, has support in $[-1,1]$ after Fourier transformation.
The partial sums defining $H$ converge as tempered distributions,
because their absolute values are bounded by $\sum_m K(x-m)=1$.
Thus $H$ has the same distributional Fourier support.

Write the interval endpoints as $c<d$, and put $\delta=b/(2\pi)$.
Set
\[
 U_I(x)=\tfrac12\{H(\delta(x-c))-H(\delta(x-d))
                       +K(\delta(x-c))+K(\delta(x-d))\},
\]
and change the two plus signs before $K$ to minus signs to obtain $L_I$.
The sign comparison proves the bounds off the endpoints. At $x=c$,
$U_I(c)=(H(\delta(d-c))+1+K(\delta(d-c)))/2\ge1$ and
$L_I(c)=(H(\delta(d-c))-1-K(\delta(d-c)))/2\le0$;
the other endpoint is the same. This proves both endpoint conventions.
The pointwise errors are bounded by the sum of the two $K$ terms,
so $0\le U_I\le3$. The odd integrable error of $H$ has integral
zero, and $\int K=1$. Hence the stated integrals follow. The
support statement follows by translation and dilation of the
compactly supported distributions above; the final functions are
integrable, so their Fourier transforms are continuous. In particular
Fourier inversion and integration against a finite measure are valid.
\end{proof}

For $B=\prod_{i=1}^d I_i$, where $d=3$ or $4$, write
$|B|=\prod_i l_i$ and $\mathcal Q(\mathbf b)=\prod_i[-b_i,b_i]$.
If $b_il_i\ge1$, tensorization gives envelopes with support in
$\mathcal Q(\mathbf b)$ and
\begin{equation}\label{eq:box-envelope-norms}
 L_B\le\one_B\le U_B,\quad
 \|L_B\|_1+\|U_B\|_1\le C_d|B|,\quad
 \int(U_B-\one_B)+\int(\one_B-L_B)
               \le C_d|B|\sum_i(b_il_i)^{-1}.
\end{equation}
Here the minorant is not the product of the minorants, which can
have the wrong sign. Precisely, set
\begin{equation}\label{eq:box-lower-envelope}
 U_B=\prod_i U_{I_i},\qquad
 L_B=U_B-\sum_i E_{I_i}\prod_{j\ne i}U_{I_j}.
\end{equation}
To verify the lower bound, telescope
$\prod_i U_{I_i}-\prod_i\one_{I_i}$ and use
$0\le U_{I_i}-\one_{I_i}\le E_{I_i}$ and $\one_{I_j}\le U_{I_j}$.
Integration of these nonnegative error bounds proves
\eqref{eq:box-envelope-norms}. Products here involve different variables,
so they do not increase a coordinate bandwidth.

\subsection{An actual box denominator, without a smoothing event}
Let $\mathbf h=(h_1,\ldots,h_4)$, $0<h_i\le1$, and define
\begin{equation}\label{eq:physical-window-box}
 B(\xi,\mathbf h)=\prod_i[\xi_i-h_i,\xi_i+h_i],\quad
 V(\mathbf h)=\prod_i2h_i,\quad
 \mathcal B_{n,R}(\xi,\mathbf h)=n\bar G_R+\sqrt n B(\xi,\mathbf h).
\end{equation}
In the last expression the first three coordinates are intersected
with $\Z^3$. Write $\mathfrak m$ for counting measure times Lebesgue
measure on the mixed record space, and put
\begin{equation}\label{eq:window-mixed-Gaussian}
 \mathfrak G_{n,R}(\mathcal B)=n^{-2}\int_{\mathcal B}
       g_{D_R}\left(\frac{x-n\bar G_R}{\sqrt n}\right)\dd\mathfrak m(x).
\end{equation}
The Gaussian density has the standard probability normalization.

\begin{theorem}[Unsmoothed window probabilities at a prescribed count]
\label{thm:unsmoothed-window-denominator}
Fix $A<\infty$. Suppose $|\xi|\le A$, $b_ih_i\ge1$,
$0<b_i\le\sqrt n/2$, and
$\mathcal Q(\mathbf b)\subset\mathcal U_n^\sharp$.
Then, uniformly in $R,n,k,a$ and these boxes,
\begin{align}
 \left|\mu_{n,R}^{[k],a}(\mathcal B)
                      -\alpha_a\mathfrak G_{n,R}(\mathcal B)\right|
 \le C_A V(\mathbf h)
       \left(\mathcal E_n(a)+M_a\sum_i(b_ih_i)^{-1}\right).
 \label{eq:unsmoothed-window-denominator}
\end{align}
For $a\ge0$, the factor $M_a$ in the second term may be replaced by
$\alpha_a$. If $\alpha_a>0$, $a\ge0$, and
\begin{equation}\label{eq:window-relative-budget}
 \epsilon_n(a,\mathbf b,\mathbf h)
       =\mathcal E_n(a)/\alpha_a+\sum_i(b_ih_i)^{-1}\longrightarrow0,
\end{equation}
then the ratio of the probability to its displayed Gaussian mass
converges to one uniformly on these central windows.
\end{theorem}
\begin{proof}
First suppose $a\ge0$, so $\lambda_{n,R}^{[k],a}$ is a positive
measure. Integrate the two envelopes against it. Their transforms
are supported in $\mathcal Q(\mathbf b)$ and bounded by $C V(\mathbf h)$,
by \eqref{eq:box-envelope-norms}. Fourier inversion and the inherited
integrated characteristic estimate give
\[
 \left|\int F\dd\lambda_{n,R}^{[k],a}
               -\alpha_a\int Fg_{D_R}\dd x\right|
       \le C V(\mathbf h)\mathcal E_n(a),\qquad F=L_B,U_B.
\]
This integrates an entire test function against the actual discrete
coordinates; no density with respect to four-dimensional Lebesgue
measure is assigned to that law. Uniform ellipticity bounds
$\|g_{D_R}\|_\infty$. The envelope excess then proves the same
estimate with $\int_Bg_{D_R}$ as Gaussian mass, with an additional
$C\alpha_a V(\mathbf h)\sum_i(b_ih_i)^{-1}$.

We now compare that continuous Gaussian mass with the mixed mass
in \eqref{eq:window-mixed-Gaussian}. The grid spacing is $n^{-1/2}$
in precisely the first three standardized coordinates. In one
coordinate, assign to each included grid point its interval of that
length. The symmetric difference from the target interval has length
at most $2n^{-1/2}$. On full cells use the uniform first-derivative
bound of the Gaussian. Iterating in the three grid coordinates gives
\begin{equation}\label{eq:window-grid-error}
 \left|\mathfrak G_{n,R}(\mathcal B)-\int_Bg_{D_R}\dd x\right|
 \le C_A V(\mathbf h)\sum_{i=1}^3(\sqrt n h_i)^{-1}.
\end{equation}
The conditions imply $\sqrt n h_i\ge2$, so all intermediate grid
volumes are comparable to the interval lengths. Since $b_i\le\sqrt n/2$,
the grid error is absorbed by the stated budget. This explicitly pays
for endpoints and the lattice convention.

For complex $a$, decompose its real and imaginary parts into positive
and negative parts. Their supremum and variation norms have sum at
most a fixed multiple of $M_a+V_a$, by contraction of BV under positive
part. Their masses have sum at most $C M_a$. Add the positive-measure
estimates to obtain \eqref{eq:unsmoothed-window-denominator}.
For the last assertion, uniform ellipticity and $|\xi|\le A$, $h_i\le1$
give $\mathfrak G_{n,R}(\mathcal B)\ge c_A V(\mathbf h)$ for all
sufficiently large $n$ under the vanishing budget. Divide the positive
estimate by that lower bound times $\alpha_a$.
\end{proof}

\begin{corollary}[A concrete rare-window law]
\label{cor:concrete-four-window}
Take $a=s_R$, so $\alpha_a=1$ and the original section probability
is recovered. Set $\rho=67/1400$ and $h_i=n^{-1/25}$.
Uniformly for $|\xi|\le A$,
\begin{equation}\label{eq:concrete-four-window}
 \nu_R^*\{X_{n,R}\in B(\xi,\mathbf h)\}
  =16n^{-4/25}g_{D_R}(\xi)
       \left(1+O_A(n^{-11/1400}\sqrt{\log(2+n)})\right).
\end{equation}
The event uses the original record, integer intervals in its first
three coordinates, and a time interval in its fourth. Each physical
half-width is $n^{23/50}$. No convolution is made in its definition.
\end{corollary}
\begin{proof}
Choose $b_i=n^\rho/4$. The product frequency box lies in the retained
isotropic ball of radius $2n^\rho$. Here
$\rho-1/25=11/1400>0$, $e=15/1400$, and $v_0>e$.
Theorem~\ref{thm:unsmoothed-window-denominator} and the Gaussian
first-derivative bound give the conclusion, since replacing the
Gaussian mass by its value at the center times volume costs relative
$O(n^{-1/25})$. The grid error is smaller. The Fourier supports are
rescaled $b_i$ and physical $b_i/\sqrt n=n^{-633/1400}/4$.
\end{proof}

\subsection{A size bound for the common signed correction}
\begin{theorem}[The common correction averaged over an actual window]
\label{thm:averaged-common-raw-remainder}
In addition to the hypotheses of
Theorem~\ref{thm:unsmoothed-window-denominator}, suppose the original
marked weight is finite-record admissible, and take the same finite
extraction and $L_n=\lceil\lambda n\rceil$, $\lambda>c_*^{-1}+1$,
as in Section~\ref{sec:coherent-raw-cutoff}. For any of the three
cutoffs $\chi_n^\sharp,\chi_n^{\rm box},\chi_n^\natural$, and any
fixed $P>0$,
\begin{align}
 \frac{n^2}{\mathfrak m(\mathcal B)}
 \left|\int_{\mathcal B}\mathcal R_\chi^{a,L_n}\dd\mathfrak m\right|
 \le C_A\left(\mathcal E_n(a)+M_a\sum_i(b_ih_i)^{-1}
                                      +M_an^{-P}\right).
 \label{eq:averaged-common-raw-remainder}
\end{align}
Thus for the concrete windows in \eqref{eq:concrete-four-window} the
common correction itself, rather than only a change between two
cutoffs, has a vanishing normalized signed window integral.
The absolute value is outside the integral.
\end{theorem}
\begin{proof}
Every cutoff is between zero and one, equals one on the original
central ball, and is supported in the retained Fourier union.
Mixed inversion and the integrated characteristic estimate imply
\begin{equation}\label{eq:window-smoothed-density}
 \sup_x\left|n^2(K_\chi*\mu_{n,R}^{[k],a})(x)
          -\alpha_ag_{D_R}((x-n\bar G_R)/\sqrt n)\right|
       \le C\mathcal E_n(a)+C_P M_an^{-P}.
\end{equation}
The omitted Gaussian frequencies are bounded by uniform ellipticity;
the first three physical supports are inside one torus chart.
On the central box, $N\le L_n$ for all sufficiently large $n$, so
$\mu^{a,L_n}(\mathcal B)=\mu_{n,R}^{[k],a}(\mathcal B)$ exactly.
The exact identity \eqref{eq:exact-raw-cutoff-transport} gives
\[
 \int_{\mathcal B}\mathcal R_\chi^{a,L_n}
  =\mu_{n,R}^{[k],a}(\mathcal B)
     -\int_{\mathcal B}K_\chi*\mu_{n,R}^{[k],a}
     +\int_{\mathcal B}K_\chi*\mathsf T^{a,L_n}.
\]
All integrals use $\mathfrak m$. Apply the preceding theorem to the
first term, \eqref{eq:window-smoothed-density} to the second, and
count-coordinate kernel separation to the last. The third-coordinate
support gap is of order $n$, and each of the three kernels gives
any fixed inverse power on that gap. Finally
$\mathfrak m(\mathcal B)\asymp n^2V(\mathbf h)$ by the same grid
count used in \eqref{eq:window-grid-error}. This proves the estimate.
No derivative norm of the extracted density was used.
\end{proof}

\paragraph{Resolution and the raw endpoint.}
The signed estimate \eqref{eq:averaged-common-raw-remainder} is a bound
for one common correction tested on unsmoothed windows, not a bound
for its essential supremum or for its absolute integral. The resolution
condition $b_ih_i\to\infty$ cannot hold for a singleton lattice
coordinate when $b_i=o(\sqrt n)$. It therefore does not establish a
fixed-label raw denominator or the microscopic local limit theorem.
The complete complementary integral, the common correction in the
required pointwise norm, and the long-time raw preparation budget
remain the requirements for that same endpoint. Positivity has here
supplied actual probability and signed-average estimates, not an
unproved inference from cutoff transport.

\section{Relative comparison at an exact physical observation time}
\label{sec:relative-physical-windows}
The previous section gives a denominator for an unchanged, unsmoothed
window. We now compare two different events relative to such a
denominator. The time of physical observation is deterministic. The
initial law in this section is $\nu_R^*$ at a section collision, not
the stationary suspension law. The window widths will be specified;
we do not replace a singleton lattice event by a larger one in the
raw theorem.

\subsection{A three-dimensional denominator from the four-dimensional band}
Put
\begin{equation}\label{eq:physical-window-exponents}
 \rho=\frac{67}{1400},\qquad \sigma=\frac1{25},\qquad
 \epsilon_n=n^{-11/1400}\sqrt{\log(2+n)}.
\end{equation}
For a positive definite $d\times d$ matrix $C$, $g_C$ denotes its
normalized $d$-dimensional Gaussian density.

\begin{lemma}[Projection of the unsmoothed window theorem]
\label{lem:projected-window-denominator}
Let $P_{n,R}:\R^4\to\R^3$ be any linear map which extends to an
invertible $4\times4$ matrix $S_{n,R}$ with
$\|S_{n,R}\|+\|S_{n,R}^{-1}\|\le C_0$. Under $\nu_R^*$, put
$Y=P_{n,R}U_{n,R}/\sqrt n$ and
$\Sigma_{n,R}=P_{n,R}D_RP_{n,R}^{\mathsf T}$.
For $|\xi|\le A$ and half-widths satisfying
$c_0n^{-\sigma}\le h_i\le C_0n^{-\sigma}$, $1\le i\le3$,
\begin{equation}\label{eq:projected-window-denominator}
 \nu_R^*\{Y\in B_3(\xi,\mathbf h)\}
  =\int_{B_3(\xi,\mathbf h)}g_{\Sigma_{n,R}}(y)\dd y
                +O_{A,c_0,C_0}\left(\prod_{i=1}^3(2h_i)\epsilon_n\right).
\end{equation}
The event is the actual projection of the record. No restriction of an
$L^1(\R^4)$ Fourier estimate to a measure-zero frequency plane is used.
\end{lemma}
\begin{proof}
Write $Y'=S_{n,R}U_{n,R}/\sqrt n$ and truncate its fourth coordinate
to $[-T,T]$, where $T=n^{1/400}$. The characteristic error for $Y'$
on the cube $|v_i|\le b$, $b=c(C_0)n^\rho$, has integral at most
$C\mathcal E_n(s_R)$. Indeed $S_{n,R}^{\mathsf T}$ maps this cube
into the retained isotropic ball $|v|\le2n^\rho$, and its determinant
and inverse determinant are bounded. This is a four-dimensional
change of variables, not a trace estimate.

Apply the product envelopes \eqref{eq:box-lower-envelope} to
$B_3(\xi,\mathbf h)\times[-T,T]$. Their $L^1$ norms are
$O(TV_3)$, where $V_3=\prod_{i=1}^3(2h_i)$. Fourier comparison thus
costs $CT V_3\mathcal E_n(s_R)$. We keep a sharper bound for the
Gaussian envelope error. Uniform ellipticity of $S_{n,R}D_RS_{n,R}^{\mathsf T}$
gives a joint density bounded by $Ce^{-c|y|^2}$. For a term containing
an error factor $E_{I_i}$ with $i\le3$, integrate that factor and
the other two output factors in $L^1$; integrate the fourth majorant
against $e^{-cy_4^2}$, using $0\le U_{[-T,T]}\le3$. This is at most
$C V_3/(bh_i)$, independently of $T$. The fourth-coordinate error
term has integral at most $CV_3/b$. The same bounds apply to both
envelopes by \eqref{eq:box-lower-envelope}.

The fourth coordinate of $Y'$ has bounded sixty-fourth moment,
uniformly in $n,R,S_{n,R}$, by
\eqref{eq:all-stopped-moment-bound}. Removing its truncation costs
at most $C T^{-64}$. For the Gaussian, the corresponding cost on
the output box is at most $C V_3 e^{-cT^2}$. Consequently the
absolute probability error divided by $V_3$ is at most
\begin{equation}\label{eq:projection-error-budget}
 C\left[T\mathcal E_n(s_R)+\sum_{i=1}^3(bh_i)^{-1}+b^{-1}
                     +T^{-64}/V_3+e^{-cT^2}\right].
\end{equation}
The four relevant power margins are
\[
 e-\frac1{400}=\frac{23}{2800},\qquad
 \rho-\sigma=\frac{11}{1400},\qquad
 \frac{64}{400}-3\sigma=\frac1{25},\qquad \rho>\frac{11}{1400}.
\]
Also $v_0-1/400>23/2800$, and the norms of $s_R$ are uniform.
Each term in \eqref{eq:projection-error-budget} is therefore
$O(\epsilon_n)$. This proves the lemma, including the discrete endpoint
conventions inherited from the envelopes.
\end{proof}

\subsection{Maximal actual-return increments and inverse clocks}
\begin{lemma}[A deterministic return-window maximal estimate]
\label{lem:actual-return-window-maximal}
For every fixed real $p>2$, every integer $q\ge1$, and $n\ge q$,
\begin{equation}\label{eq:actual-return-window-maximal}
 \int\max_{|j|\le q}|U_{n+j,R}-U_{n,R}|^p\dd\nu_R^*
                       \le C_p q^{p/2}.
\end{equation}
\end{lemma}
\begin{proof}
Every centered return increment over a deterministic interval of
length $l$ has $p$th moment at most $C_p l^{p/2}$, by invariance of
$F_R^*$ and \eqref{eq:all-stopped-moment-bound}. Enlarge $q$ to a
power of two $Q<2q$. At scale $l$, the maximum over its $Q/l$
aligned intervals has $L^p$ norm at most
$C_p Q^{1/p}l^{1/2-1/p}$. A prefix is a sum of at most one interval
per scale. Minkowski's inequality and $p>2$ give a convergent geometric
sum bounded by $C_p\sqrt Q$. Translation by the deterministic index
$n$ is justified by return-map invariance. The backward part uses
suffixes of the length-$q$ interval ending at $n$, and has the same
bound. Adding the two parts proves the assertion. The maximum is over
a short deterministic return window, not over all marks in a
collision Fourier theorem.
\end{proof}

Start the physical orbit at $x\in Y_R^*$, with physical time zero at
that outgoing collision. Let $W_R^Y(t)$ be its lattice displacement
and $C_R^Y(t)$ its completed collision count, using the physical
convention of Section~\ref{sec:uniform-fluid-clock}. Put
\begin{equation}\label{eq:section-start-physical-vector}
 V_R^Y(t)=\left(W_R^Y(t),C_R^Y(t)-t/\bar\tau_R\right),\quad
 B_R(x_1,x_2,x_3,x_4)=(x_1,x_2,x_3-x_4/\bar\tau_R).
\end{equation}
The chosen displacement convention can differ from the sum of completed
flight labels by one uniformly bounded cell-position term; that term
is retained in the estimates below. Define the actual completed-return
clock and its deterministic mean index by
\begin{equation}\label{eq:physical-return-index}
 M_R(t)=\max\{m\ge0:T_{m,R}\le t\},\qquad
 n_R(t)=\left\lfloor c_*t/\bar\tau_R\right\rfloor.
\end{equation}
The positive lower free-flight length makes the maximum finite and
bounds it deterministically by $C(1+t)$.

\begin{proposition}[A coupling at the deterministic physical time]
\label{prop:physical-window-coupling}
For every fixed $P>0$, there is $C_P$ such that, uniformly in $R$ and
all sufficiently large $t$,
\begin{equation}\label{eq:physical-window-coupling}
 \nu_R^*\left\{
   |V_R^Y(t)-B_RU_{n_R(t),R}|>C_P t^{3/8}
          \right\}\le C_Pt^{-P}.
\end{equation}
The same bound holds for the difference
$B_RU_{M_R(t),R}-B_RU_{n_R(t),R}$.
\end{proposition}
\begin{proof}
Put $n=n_R(t)$, $q=\lceil t^{5/8}\rceil$; for large $t$, $q<n$
and $n\asymp t$ uniformly. The return flight-time mean is
$\bar\tau_R/c_*$. The equivalences for the increasing times $T_m$
turn $|M_R(t)-n|>q$ into a deviation of $T_{n+q}$ or $T_{n-q}$
from its mean of size at least $c q$. The rounding error is uniformly
bounded. For every fixed $p$, \eqref{eq:all-stopped-moment-bound}
and Markov's inequality therefore give
\[
 \nu_R^*\{|M_R(t)-n|>q\}\le C_p t^{p/2}q^{-p}
                                      \le C_p t^{-p/8}.
\]
Off this event, \eqref{eq:actual-return-window-maximal} bounds the
probability that $|U_{M_R(t)}-U_n|>t^{3/8}$ by
$C_p q^{p/2}t^{-3p/8}\le C_p t^{-p/16}$. The operator norm of
$B_R$ is uniform, and $B_R\bar G_R=0$.

At the last completed section return, the vector is
$B_RJ_{M_R(t)}=B_RU_{M_R(t)}$. Between this return and time $t$,
the change in displacement, collision count, and elapsed time is
bounded by $C r_R^*\circ(F_R^*)^{M_R(t)}+C$. Uniformly bounded
single flights and labels justify this assertion even when the next
return has not yet occurred. Since $M_R(t)\le C(1+t)$,
invariance and the exponential one-return tail give
\[
 \nu_R^*\left\{\max_{0\le j\le C(1+t)}r_R^*\circ(F_R^*)^j
                                               >t^{3/8}\right\}
       \le C(1+t)e^{-ct^{3/8}}.
\]
This is a union bound, not a conditional independence statement.
Choose a fixed $p>\max\{2,16P\}$ and combine the three events.
The count of moment factors is fixed for each $P$. This proves both
assertions with the actual physical observation time left unchanged.
\end{proof}

\subsection{A relative event theorem with its actual denominator}
Recall the uniformly positive physical covariance
\[
 \mathcal V_R=\bar\tau_R^{-1}B_R\Gamma_RB_R^{\mathsf T}
             =(c_*/\bar\tau_R)B_RD_RB_R^{\mathsf T}.
\]
For $h=t^{-1/25}$ and $|\xi|\le A$, let $B_t=B_3(\xi,(h,h,h))$.
On the same initial probability space $(Y_R^*,\nu_R^*)$ define
\begin{align}
 A_t^{\rm ret}&=\{t^{-1/2}B_RU_{n_R(t),R}\in B_t\},\notag\\
 A_t^{\rm last}&=\{t^{-1/2}B_RU_{M_R(t),R}\in B_t\},\notag\\
 A_t^{\rm phys}&=\{t^{-1/2}V_R^Y(t)\in B_t\}.
 \label{eq:three-physical-window-events}
\end{align}
All three are exact events. The third is measured at deterministic
physical time $t$; it uses integer lattice and collision-count
intervals of physical half-width $t^{23/50}$.

\begin{theorem}[An unsmoothed physical denominator and relative replacement]
\label{thm:relative-physical-window}
Uniformly for $R\in I$, $|\xi|\le A$, and large $t$, all three events
in \eqref{eq:three-physical-window-events} have probability
\begin{equation}\label{eq:physical-window-denominator}
 8t^{-3/25}g_{\mathcal V_R}(\xi)
       \left(1+O_A(t^{-11/1400}\sqrt{\log(2+t)})\right).
\end{equation}
For either $A_t=A_t^{\rm last}$ or $A_t^{\rm phys}$,
\begin{equation}\label{eq:relative-physical-window}
 \frac{\nu_R^*(A_t\mathbin\triangle A_t^{\rm ret})}
                    {\nu_R^*(A_t^{\rm ret})}
       \le C_A t^{-11/1400}\sqrt{\log(2+t)}.
\end{equation}
The total variation distance, defined as the supremum over measurable
sets, between the two conditional initial laws is bounded by the same
right-hand side with a changed constant. Consequently the comparison
holds for every common bounded measurable path statistic, with a
factor twice its supremum norm for the difference of expectations.
\end{theorem}
\begin{proof}
Set $n=n_R(t)$ and take
$P_{n,R}=\sqrt{n/t}\,B_R$ in
Lemma~\ref{lem:projected-window-denominator}. Adjoin the fourth row
$(0,0,0,1)$ to obtain $S_{n,R}$. The ratio $n/t$ is bounded above
and away from zero uniformly, as are $\bar\tau_R$ and its reciprocal;
thus the matrix and its inverse satisfy the lemma's bounds.
Its projected covariance is $(n/t)B_RD_RB_R^{\mathsf T}$, which
differs from $\mathcal V_R$ by $O(t^{-1})$. Uniform ellipticity and
the Gaussian formula make the density error $O_A(t^{-1})$ on central
boxes. Replacing the box Gaussian mass by its center value costs
relative $O_A(h)$. The lemma proves
\eqref{eq:physical-window-denominator} for $A_t^{\rm ret}$,
including a lower bound $c_A t^{-3/25}$.

For the event comparison, use
Proposition~\ref{prop:physical-window-coupling} with $P=2$.
Outside a set of probability $O(t^{-2})$, the standardized vector
changes by at most $d_t=Ct^{-1/8}$. Hence the symmetric difference
is contained in that exceptional set and the event that the return
vector lies in $B_t^+\setminus B_t^-$, where the half-widths are
$h+d_t$ and $h-d_t$. For large $t$ they lie between $h/2$ and $2h$.
Apply the projected-window lemma separately to these two boxes and
subtract their probabilities. The Gaussian boundary volume is at
most $C_A h^3 d_t/h$, with
\[
 d_t/h=O(t^{-17/200}).
\]
The two Fourier-envelope errors have total at most
$C_A h^3 t^{-11/1400}\sqrt{\log(2+t)}$. Thus the symmetric-difference
probability divided by the proved denominator is at most
\[
 C_A\left[t^{-17/200}+t^{-11/1400}\sqrt{\log(2+t)}
                                      +t^{-2+3/25}\right].
\]
The first and third terms have strictly faster power decay, so this proves
\eqref{eq:relative-physical-window}. This argument subtracts masses
of enlarged and contracted boxes; it does not require a separate
Fourier estimate for a thinner unresolved boundary slab.
The difference of event probabilities is bounded by their symmetric
difference, proving \eqref{eq:physical-window-denominator} for the
other events as well.

Finally, if $p=\nu_R^*(A_t^{\rm ret})$ and the symmetric-difference
mass is $\delta\le p/2$, then $\nu_R^*(A_t)\ge p/2$.
Writing the two normalized indicator densities shows directly that
the conditional total variation distance is at most $2\delta/p$.
Integration of any common bounded statistic proves the final assertion.
No event was silently substituted in a numerator or a denominator.
\end{proof}

\paragraph{The still required microscopic comparison.}
Theorem~\ref{thm:relative-physical-window} proves a relative comparison
for different events and a nonzero denominator at deterministic
physical time, but at the displayed growing physical window widths
and under a section-start law. The original raw application also
requires fixed lattice labels, its own exact time constraint, and
its stated initial law. Those are not the events of
\eqref{eq:three-physical-window-events}. In particular neither this
theorem nor the signed average estimate in the preceding section
supplies the full pointwise common remainder, the long-time density
jet budget, or the full weighted fixed-count Fourier complement.
The same raw mixed-density target and all its existing mechanisms
are retained. The new results identify precisely a range in which
actual denominators and relative event errors can already be proved,
without presupposing the microscopic theorem.

\section{A wider collision band and projected window probabilities}
\label{sec:stationary-collision-band}
A prescribed collision interval does not incur a return-clock stopping
error. We exploit this distinction before passing to physical time.
It permits a full isotropic collision band with rescaled radius
$m^{9/100}$, without asserting such a band at the prescribed return
count. The initial collision density may be nonstationary, provided
that its supremum and initial-coordinate variation are bounded.
This includes the length-biased collision marginal of equilibrium flow.

Write $S_{m,R}h_R=\sum_{j=0}^{m-1}h_R\circ T_R^j$ and put
\[
 e_0=\frac3{280},\qquad
 \mathfrak e_m=m^{-e_0}\sqrt{\log(2+m)}.
\]
All initial densities in this section are relative to the collision
probability $\nu$, not to $\nu_R^*$.

\subsection{Eliminating the stopping loss on a fixed collision interval}
\begin{theorem}[An isotropic collision band with bounded-BV initial data]
\label{thm:wide-collision-band}
Fix $B,C_0<\infty$. For every complex collision function $a$ with
$\|a\|_\infty+\|a\|_{\mathrm{BV}}\le B$, and
$\alpha_a=\int a\dd\nu$, one has
\begin{align}
 &\int_{|v|\le C_0m^{9/100}}
 \left|\int a e^{iv\cdot S_{m,R}h_R/\sqrt m}\dd\nu
                   -\alpha_a e^{-v^{\mathsf T}\Gamma_Rv/2}\right|\dd v
 \le C_{B,C_0}\mathfrak e_m .
 \label{eq:wide-collision-band}
\end{align}
The bound is uniform in $R\in I$ and $m\ge2$. The physical frequencies
have radius $C_0m^{-41/100}$. There is no stopping time in the transform.
\end{theorem}
\begin{proof}
On $|v|\le2m^{1/200}$, integrate
Lemma~\ref{lem:explicit-frequency-error} with
$\delta=m^{-1/14}/4$. The six resulting terms, with
$U=m^{1/200}$ and $L_\delta=1+|\log\delta|$, are
\[
 C_B\left(\delta U^4+\sqrt{\delta L_\delta}\,U^5
 +\delta L_\delta U^6+m^{-1/2}\delta^{-4}U^5
 +m^{-1/2}\delta^{-6}U^7+\delta^{-2}\rho^mU^4\right).
\]
The slowest exponent is $1/28-5/200=3/280$. The cubic exponent
$1/2-6/14-7/200=51/1400$ and all other powers have a strict
margin over $3/280$. Thus this integral is $O_B(\mathfrak e_m)$.
No return-clock comparison is used in this calculation.

For the annulus $2m^{1/200}\le|v|\le C_0m^{9/100}$ use
Theorem~\ref{thm:high-order-damped-unsmoothing}, with a mark at the
initial collision, and the fixed choices
\begin{equation}\label{eq:wide-collision-orders}
 P=40,\quad Q=29,\quad
 \delta=m^{-19/100}/4,\quad \epsilon=m^{-50}/4.
\end{equation}
The residual Taylor degree is $79$; the spectral Taylor degree is
$29$. Neither varies with $m$. For $z=v/\sqrt m$, the two analytic
margins in \eqref{eq:finite-jet-damping-domain} are
\[
 \frac12-\frac9{100}-2\frac{19}{100}=\frac3{100},\qquad
 28\left(\frac12-\frac9{100}\right)
              -2\frac{19}{100}\,30=\frac2{25}.
\]
Fixed multiplicative constants $C_0$ are absorbed by taking $m$
sufficiently large. The volume is $O_{C_0}(m^{9/25})$.
The fine insertion and fine residual errors have decay exponents
\[
 50-\frac9{25}=\frac{1241}{25},\qquad
 50-\frac9{25}-79\left(\frac12+\frac9{100}\right)
                                      =\frac{303}{100}.
\]
For the residual moment term indexed by $1\le b\le40$, the power
before taking its negative decay sign is
\begin{equation}\label{eq:wide-collision-residual-powers}
 \frac9{25}-40+80\frac9{100}+\frac{81}{100}b.
\end{equation}
Its largest value, at $b=40$, is $-1/25$. Every smaller $b$ improves
this by $81/100$ per removed block. In particular the connected
$b=1$ term is included. The fixed logarithmic powers are absorbed
by $1/25-3/280=41/1400>0$.

The damped word has a fixed polynomial prefactor, possibly large,
and $m|z|^2\ge4m^{1/100}$. It is smaller than every fixed inverse
power after integration. The omitted Gaussian contribution has the
same property by the already proved uniform ellipticity of
$\Gamma_R=c_*D_R$. This proves the annular estimate and then
\eqref{eq:wide-collision-band}. Increasing the constant covers the
finitely many remaining $m$. The argument enlarges a collision-clock
band, not the known prescribed-return-count Fourier region.
\end{proof}

\subsection{Projection without taking a trace of a Fourier estimate}
For a positive definite matrix $C$, write $g_C$ for its normalized
Gaussian density. Put
\begin{equation}\label{eq:stationary-window-exponents}
 \sigma_1=\frac2{25},\qquad
 \kappa_1=\frac{23}{2800},\qquad
 \varepsilon_m=m^{-\kappa_1}\sqrt{\log(2+m)}.
\end{equation}

\begin{theorem}[Projected collision windows under bounded initial densities]
\label{thm:projected-collision-windows}
Let $\sigma_R\ge0$ satisfy $\int\sigma_R\dd\nu=1$ and
$\|\sigma_R\|_\infty+\|\sigma_R\|_{\mathrm{BV}}\le B$.
Let $P_{m,R}:\R^4\to\R^3$ extend to a matrix $\mathsf A_{m,R}$ with
$\|\mathsf A_{m,R}\|+\|\mathsf A_{m,R}^{-1}\|\le C_0$. Set
\[
 Y=\frac1{\sqrt m}P_{m,R}\sum_{j=0}^{m-1}h_R\circ T_R^j.
\]
For $|\xi|\le A$ and $c_0m^{-\sigma_1}\le h_i\le C_0m^{-\sigma_1}$,
\begin{align}
 \int\sigma_R\one_{\{Y\in B_3(\xi,\mathbf h)\}}\dd\nu
 =\int_{B_3(\xi,\mathbf h)}
       g_{P_{m,R}\Gamma_RP_{m,R}^{\mathsf T}}(y)\dd y
       +O_{A,B,c_0,C_0}\left(V_3\varepsilon_m\right),
 \label{eq:projected-collision-windows}
\end{align}
where $V_3=\prod_{i=1}^3(2h_i)$. Endpoints in the exact event may
be included or excluded. No smoothing is added to that event.
\end{theorem}
\begin{proof}
Write $\mathsf A=\mathsf A_{m,R}$. Use the full
four-dimensional variable $Y'=\mathsf A(\sum_{j<m}h_R\circ T_R^j)/\sqrt m$.
Truncate its fourth coordinate to $[-T,T]$, with $T=m^{1/400}$.
Choose the same bandwidth $b=c(C_0)m^{9/100}$ in all four output
coordinates. The linear change of frequency $v\mapsto\mathsf A^{\mathsf T}v$
maps this cube into the band of Theorem~\ref{thm:wide-collision-band},
with bounded determinant factors. The characteristic error integrated
on the cube is $O(\mathfrak e_m)$.

Apply the endpoint-safe product envelopes
\eqref{eq:box-lower-envelope} to
$B_3(\xi,\mathbf h)\times[-T,T]$. Their $L^1$ norms are $O(TV_3)$,
so Fourier comparison costs $CTV_3\mathfrak e_m$. To estimate the
Gaussian envelope error, use a joint density bounded by
$C\exp(-c|y|^2)$, uniformly in all parameters. In an error term for
one of the first three coordinates, integrate that error factor and
the other two output majorants in $L^1$, but integrate the fourth
majorant against the Gaussian decay. Since it is at most three,
this contributes $CV_3/(bh_i)$, without a factor $T$. The error
factor for the fourth coordinate contributes $CV_3/b$.

Lemma~\ref{lem:fixed-even-maximal} at order $128$, followed by
$\sigma_R\le B$, gives a uniform $128$th moment for $Y'_4$.
Its truncation error is at most $CT^{-128}$. The Gaussian tail
costs at most $CV_3e^{-cT^2}$. Thus the absolute error divided by
$V_3$ is bounded by
\begin{equation}\label{eq:stationary-projection-budget}
 C\left(T\mathfrak e_m+\sum_{i=1}^3(bh_i)^{-1}+b^{-1}
                  +T^{-128}/V_3+e^{-cT^2}\right).
\end{equation}
The relevant exponents are
\[
 \frac3{280}-\frac1{400}=\frac{23}{2800},\quad
 \frac9{100}-\frac2{25}=\frac1{100},\quad
 \frac{128}{400}-3\frac2{25}=\frac2{25}.
\]
They give \eqref{eq:projected-collision-windows}. Fourier comparison
has used a genuine four-dimensional integral and then removed a
coordinate by a moment bound; it has not restricted an $L^1(\R^4)$
estimate to a plane. Positivity and the endpoint envelope inequalities
apply directly even when output coordinates have atoms.
\end{proof}

\section{Equilibrium physical windows and relative conditional laws}
\label{sec:stationary-physical-conditioning}
A typical physical observation does not start at a section collision.
The section state seen from equilibrium is biased by the duration
of its return excursion. We retain this bias exactly. A wider
collision-clock Fourier band supplies the denominator, while the
return-clock moment estimates compare the same stationary physical
event with its completed-return descriptions. This avoids assuming
bounded variation of a length-biased entrance distribution.

\subsection{Two descriptions of one equilibrium probability}
Write $\Theta_R(y)=\tau_R^*(y)$, $r_R^*(y)=r(y)$, and
$\bar\Theta_R=\int\Theta_R\dd\nu_R^*=\bar\tau_R/c_*$. By
Proposition~\ref{prop:induced-suspension}, equilibrium flow is represented on
\begin{equation}\label{eq:stationary-return-suspension}
 \mathcal X_R=\{(y,s):y\in Y_R^*,\ 0\le s<\Theta_R(y)\},\qquad
 \dd\mathbb P_R(y,s)=\bar\Theta_R^{-1}\dd\nu_R^*(y)\dd s.
\end{equation}
The physical observation starts at the flow state reached after time
$s$ from $y$. For $0\le j\le r(y)$ put
$\tau_{j,R}^{\rm sum}(y)=\sum_{i<j}\tau_R(T_R^iy)$.
There is a unique $0\le j<r(y)$ such that
$\tau_{j,R}^{\rm sum}\le s<\tau_{j+1,R}^{\rm sum}$, apart from
null endpoint conventions. Define the current outgoing collision and
its elapsed flight age by
\[
 x=T_R^jy,\qquad u=s-\tau_{j,R}^{\rm sum}(y).
\]

\begin{lemma}[The two length biases and exceptional-set transfer]
\label{lem:stationary-length-bias}
Under $\mathbb P_R$, the marginal law of $y$ is
$(\Theta_R/\bar\Theta_R)\nu_R^*$, whereas the law of $(x,u)$ is
\begin{equation}\label{eq:collision-suspension-marginal}
 \bar\tau_R^{-1}\one_{\{0\le u<\tau_R(x)\}}\dd\nu(x)\dd u.
\end{equation}
In particular the outgoing collision marginal is
$\sigma_R\nu$, where $\sigma_R=\tau_R/\bar\tau_R$ is a probability
density with uniformly bounded supremum and BV norms.
There are uniform $C,c>0$ such that, for measurable $E\subset Y_R^*$
and $L\ge0$,
\begin{equation}\label{eq:stationary-bias-bounds}
 \mathbb P_R\{y\in E\}\le C\nu_R^*(E)^{1/2},\qquad
 \mathbb P_R\{\Theta_R(y)>L\}\le Ce^{-cL}.
\end{equation}
\end{lemma}
\begin{proof}
Integrating $s$ in \eqref{eq:stationary-return-suspension} gives the
first marginal. For an integrable test $F$ of $(x,u)$, partition each
excursion into its actual collision flights. The tower identity gives
\begin{align*}
 \int F(x,u)\dd\mathbb P_R
 &=\frac1{\bar\Theta_R}\int_{Y_R^*}\sum_{j<r(y)}
       \int_0^{\tau_R(T_R^jy)}F(T_R^jy,u)\dd u\dd\nu_R^*(y)\\
 &=\frac1{c_*\bar\Theta_R}\int_M\int_0^{\tau_R(x)}F(x,u)\dd u\dd\nu(x).
\end{align*}
Since $c_*\bar\Theta_R=\bar\tau_R$, this proves the second marginal.
Lemma~\ref{lem:physical-bv} and the positive lower bound for
$\bar\tau_R$ give the density bounds. No assertion that
$\Theta_R/\bar\Theta_R$ is bounded-BV is needed or made.

The return roof has a uniform exponential moment by
Theorem~\ref{thm:exponential-return}. In particular
$\|\Theta_R\|_{L^2(\nu_R^*)}$ is uniform. Cauchy--Schwarz proves
the first estimate in \eqref{eq:stationary-bias-bounds}. Integrating
the exponential roof tail after multiplication by $\Theta_R$, and
reducing the tail exponent, proves the second. The positive normalizing
mean is bounded away from zero. This treats the physical-time length
bias, not merely the collision-count bias by $r(y)$.
\end{proof}

\subsection{Coupling collision, return and physical observations}
Let $V_R(t)=(W_R(t),C_R(t)-t/\bar\tau_R)$ be the actual physical
increment during the interval of length $t$ starting at $(y,s)$.
The lattice convention is the one used throughout the article; any
bounded cell-position term is included in the following error bounds.
Retain $B_R$ from \eqref{eq:section-start-physical-vector}. Set
\begin{align}
 m_R(t)&=\lfloor t/\bar\tau_R\rfloor,&
 n_R(t)&=\lfloor c_*t/\bar\tau_R\rfloor,\notag\\
 L_R(t;y,s)&=\max\{l\ge0:T_{l,R}(y)\le t+s\},\notag\\
 Z_t^{\rm coll}&=t^{-1/2}B_R\sum_{i<m_R(t)}h_R(T_R^ix),&
 Z_t^{\rm ret}&=t^{-1/2}B_RU_{n_R(t),R}(y),\notag\\
 Z_t^{\rm last}&=t^{-1/2}B_RU_{L_R(t;y,s),R}(y),&
 Z_t^{\rm phys}&=t^{-1/2}V_R(t).
 \label{eq:stationary-four-observations}
\end{align}
These are random variables on the same probability space
$(\mathcal X_R,\mathbb P_R)$. The return reference starts at the
actual section origin $y$ preceding the stationary observation, not
at a newly sampled independent section state.

\begin{proposition}[Uniform stationary observation coupling]
\label{prop:stationary-observation-coupling}
For each fixed $P>0$ there is $C_P$ such that, for all large $t$,
\begin{equation}\label{eq:stationary-observation-coupling}
 \mathbb P_R\left\{\max_{j\in\{{\rm ret},{\rm last},{\rm phys}\}}
 |Z_t^j-Z_t^{\rm coll}|>C_Pt^{-1/8}\right\}\le C_Pt^{-P}
 \qquad(R\in I).
\end{equation}
\end{proposition}
\begin{proof}
First compare the collision reference with the physical vector.
The age $u$ is bounded by the uniform one-flight horizon. If
$\ell(t)$ is the completed-flight index measured from $x$, then
its threshold differs from $t$ by at most that horizon. The flight
roof is bounded-BV and has mean $\bar\tau_R$. Fixed collision
moments and the monotone inverse-clock equivalences therefore give,
with $m=m_R(t)$, $q=\lceil t^{5/8}\rceil$, and each fixed $p$,
\[
 \mathbb P_R\{|\ell(t)-m|>q\}\le C_p t^{p/2}q^{-p}
                                            \le C_pt^{-p/8}.
\]
Indeed the outgoing collision density is bounded by
Lemma~\ref{lem:stationary-length-bias}, so stationary collision
moment estimates apply by domination; independence of $u$ and $x$
is not used. Off this event the centered collision-sum difference
lies in the deterministic length-$q$ windows adjacent to $m$.
Lemma~\ref{lem:fixed-even-maximal}, at a sufficiently large fixed
even order, bounds the probability that its size exceeds $t^{3/8}$
by $C_p q^{p/2}t^{-3p/8}\le C_pt^{-p/16}$.
The initial and final partial flights contribute $O(1)$ after applying
$B_R$, because $B_Rh_R=(\kappa_R,1-\tau_R/\bar\tau_R)$.
Choose $p>16P$. This proves the physical-to-collision comparison.

For the return reference, apply
Proposition~\ref{prop:physical-window-coupling} under $\nu_R^*$,
with its exceptional probability exponent chosen larger than $2P$.
The first estimate in \eqref{eq:stationary-bias-bounds} transfers
this exceptional set to $\mathbb P_R$ with probability $O(t^{-P})$.
The actual interval from time $s$ to $s+t$ differs from the
section-start interval from zero to $t$ by two end intervals of
length at most $s\le\Theta_R(y)$. The positive minimum free-flight
length and bounded one-flight displacement give, deterministically,
\[
 |V_R(t)-V_R^Y(t;y)|\le C(1+\Theta_R(y)).
\]
This includes collision-count centering and cell-position conventions.
The exponential tail in \eqref{eq:stationary-bias-bounds} makes this
error at most $Ct^{3/8}$ outside a smaller exceptional set. Thus the
return reference and the physical vector obey the required comparison.

For the last-return reference, its difference from the physical vector
is bounded before normalization by
\[
 C\bigl(1+\Theta_R(y)+r_R^*((F_R^*)^{L_R(t;y,s)}y)\bigr).
\]
The initial partial return contributes the roof term and the final
partial return contributes the last term. On $\Theta_R\le t$,
the positive minimum flight length gives $L_R(t;y,s)\le C(1+t)$.
Under $\nu_R^*$ the union of the events
$r_R^*\circ(F_R^*)^j>t^{3/8}$ over these deterministic indices has
probability at most $C(1+t)e^{-ct^{3/8}}$, by return invariance.
The first bound in \eqref{eq:stationary-bias-bounds} transfers this
union to $\mathbb P_R$; the event $\Theta_R>t$ has exponential
probability as well. This proves the last-return comparison.
A triangle inequality and union bound now give
\eqref{eq:stationary-observation-coupling} for all three variables
against the common collision reference. All moment orders are fixed
for a chosen $P$.
\end{proof}

\subsection{The stationary denominator and relative replacement}
Let $h=t^{-2/25}$ and $B_t=B_3(\xi,(h,h,h))$, $|\xi|\le A$.
For each $j\in\{{\rm coll},{\rm ret},{\rm last},{\rm phys}\}$ define
\begin{equation}\label{eq:stationary-four-events}
 A_t^j=\{Z_t^j\in B_t\}.
\end{equation}
The physical event is an observation at the original deterministic
time $t$, with exact integer displacement and collision-count
intervals of physical half-width $t^{21/50}$.

\begin{theorem}[Stationary physical windows and relative conditional laws]
\label{thm:stationary-physical-windows}
Uniformly for $R\in I$, $|\xi|\le A$, and large $t$, each event in
\eqref{eq:stationary-four-events} satisfies
\begin{equation}\label{eq:stationary-window-denominator}
 \mathbb P_R(A_t^j)=8t^{-6/25}g_{\mathcal V_R}(\xi)
            \left(1+O_A(t^{-23/2800}\sqrt{\log(2+t)})\right).
\end{equation}
For $j={\rm coll},{\rm last},{\rm phys}$,
\begin{equation}\label{eq:stationary-relative-event-bound}
 \frac{\mathbb P_R(A_t^j\mathbin\triangle A_t^{\rm ret})}
                 {\mathbb P_R(A_t^{\rm ret})}
 \le C_A t^{-23/2800}\sqrt{\log(2+t)}.
\end{equation}
The conditional laws $\mathbb P_R(\,\cdot\mid A_t^j)$ and
$\mathbb P_R(\,\cdot\mid A_t^{\rm ret})$ on the full stationary
initial-state space have total variation distance bounded by the
same right side with a changed constant. Thus the comparison holds
for every common bounded measurable statistic of the entire orbit.
\end{theorem}
\begin{proof}
Take $m=m_R(t)$ and the initial density
$\sigma_R=\tau_R/\bar\tau_R$. Apply
Theorem~\ref{thm:projected-collision-windows} with
$P_{m,R}=\sqrt{m/t}\,B_R$, augmenting by the fourth row $(0,0,0,1)$.
The ratios $m/t$, $\bar\tau_R$ and their reciprocals have uniform
bounds for large $t$, so the augmenting matrix and its inverse do
also. The covariance is
\[
 (m/t)B_R\Gamma_RB_R^{\mathsf T}
                         =\mathcal V_R+O(t^{-1}).
\]
The window half-width $t^{-2/25}$ is comparable to $m^{-2/25}$.
Uniform ellipticity controls the change in the Gaussian density from
this $O(t^{-1})$ perturbation. Replacing its integral on the window
by $8h^3$ times its central value costs relative $O_A(h)$.
This proves \eqref{eq:stationary-window-denominator} for
$A_t^{\rm coll}$ and gives a lower bound $c_Ah^3$.

Choose $P=2$ in
Proposition~\ref{prop:stationary-observation-coupling} and put
$d_t=C_2t^{-1/8}$. If $A_t^j$ and $A_t^{\rm coll}$ differ outside
the exceptional set, then $Z_t^{\rm coll}$ belongs to the shell
between boxes of half-widths $h+d_t$ and $h-d_t$ about $\xi$.
For large $t$ those widths are between $h/2$ and $2h$.
The projected collision-window theorem applies to both boxes with
the same constants. Subtract their probabilities before taking
absolute values. The Gaussian shell mass is at most $C_Ah^2d_t$;
the total envelope comparison error is at most $C_Ah^3\varepsilon_t$.
Consequently
\begin{equation}\label{eq:stationary-shell-budget}
 \frac{\mathbb P_R(A_t^j\mathbin\triangle A_t^{\rm coll})}
                         {\mathbb P_R(A_t^{\rm coll})}
 \le C_A\left(t^{-9/200}+\varepsilon_t+t^{-2+6/25}\right)
 \le C_A\varepsilon_t.
\end{equation}
This also proves the denominator asymptotic for each $j$.
The symmetric-difference triangle inequality and the resulting
comparison of the denominators give
\eqref{eq:stationary-relative-event-bound}. No Fourier estimate of
an unresolved thin boundary slab is assumed.

For measurable events $A,B$ with $p=\mathbb P_R(A)>0$ and
$\delta=\mathbb P_R(A\triangle B)\le p/2$, direct subtraction of
$\one_A/p$ and $\one_B/\mathbb P_R(B)$ gives
$d_{\rm TV}(\mathbb P_R(\cdot\mid A),\mathbb P_R(\cdot\mid B))
\le2\delta/p$. Apply this to the proved event comparison. The
expectation difference for a common bounded statistic is at most
twice its supremum norm times that distance. This is a comparison
of conditional measures, not a claim of a conditional invariance
principle or of independence from the event.
\end{proof}

\begin{corollary}[Additional bounded path selection]
\label{cor:stationary-weighted-selection}
Let $w_t\ge0$ be any measurable function on $\mathcal X_R$, possibly
depending on the whole path, with $\|w_t\|_\infty\le B$ and
\[
 \int_{A_t^{\rm ret}}w_t\dd\mathbb P_R
                   \ge a_t\mathbb P_R(A_t^{\rm ret}),\qquad a_t>0.
\]
If $B\varepsilon_t/a_t\to0$, the conditional laws with densities
proportional to $w_t\one_{A_t^j}$ and $w_t\one_{A_t^{\rm ret}}$
have total variation distance $O_A(B\varepsilon_t/a_t)$.
\end{corollary}
\begin{proof}
Their unnormalized difference has $L^1$ norm at most
$B\mathbb P_R(A_t^j\triangle A_t^{\rm ret})$.
Divide by the stated lower bound for the weighted return denominator,
then use the same normalization inequality as in the preceding proof.
The lower bound involving $a_t$ is an explicit hypothesis, not a
weighted raw denominator theorem inferred from the unweighted one.
\end{proof}

\paragraph{Relation to the original raw problem.}
Theorems~\ref{thm:wide-collision-band} and
\ref{thm:stationary-physical-windows} supply a collision-clock
Fourier estimate and a relative physical-event theorem under the
actual equilibrium law. They are not statements under the
unbiased section-start law. The physical half-width $t^{21/50}$
is smaller than the section-start half-width $t^{23/50}$ above,
but still grows: it is not a singleton lattice label. The probability
$t^{-6/25}$ and the relative error in
\eqref{eq:stationary-relative-event-bound} refer to these exact
windows and to no substituted event. The bounded path-selection
corollary states its own additional denominator requirement.

The full prescribed-return-count complement, the pointwise common
signed correction, the long-time inverse-coarea derivative budget,
and the original microscopic conditioning theorem remain distinct
requirements of the same raw mixed-density problem. In particular
removing the collision stopping loss has not removed the return-count
stopping loss from Theorems P and Q. The common raw correction is still
bounded only in the signed window norm proved in
Theorem~\ref{thm:averaged-common-raw-remainder}, unless its stronger
pointwise norm is established separately.

\section{Joint collision blocks under a shrinking endpoint window}
\label{sec:multiblock-bridge-windows}
The comparison of two conditioned initial measures does not identify
their common limit. We now retain finitely many path frequencies while
integrating the endpoint frequency. The estimates are uniform in the
positions of the block boundaries, including coincident boundaries.
This uniformity will be used for conditional tightness, not only for
finite-dimensional convergence.

Write $S_mh=S_{m,R}h_R$, $e_0=3/280$, and
$\mathfrak e_m=m^{-e_0}\sqrt{\log(2+m)}$, as in
Section~\ref{sec:stationary-collision-band}. All collision expectations
in this section are under $\sigma\nu$, where
\begin{equation}\label{eq:bridge-initial-class}
 \sigma\ge0,\qquad \int\sigma\dd\nu=1,\qquad
 \|\sigma\|_\infty+\|\sigma\|_{\mathrm{BV}}\le B.
\end{equation}
The constants may depend on the fixed number of blocks, but not on
any of their lengths.

\subsection{A uniform joint Fourier estimate}
Fix an integer $q\ge1$ and integers
$0=j_0\le j_1\le\cdots\le j_q=m$. Put
$m_l=j_l-j_{l-1}$, $d_l=m_l/m$, and
$H_l=(S_{j_l}h-S_{j_{l-1}}h)/\sqrt m$.
For $u_1,\ldots,u_q\in\R^4$ define
\begin{align}
 \mathcal C_m(v;\mathbf u)
 &=\int\sigma\exp\left(i\sum_{l=1}^q(v+u_l)\cdot H_l\right)\dd\nu,
 \notag\\
 \mathcal G_m(v;\mathbf u)
 &=\exp\left(-\frac12\sum_{l=1}^q
               d_l(v+u_l)^{\mathsf T}\Gamma_R(v+u_l)\right).
 \label{eq:multiblock-characteristic}
\end{align}
An interval of length zero contributes zero to both expressions.

\begin{theorem}[The endpoint band with joint block frequencies]
\label{thm:multiblock-collision-band}
For fixed $q,B,C_0<\infty$, uniformly in $R$, the subdivision, and
$\max_l|u_l|\le m^{1/200}$,
\begin{equation}\label{eq:multiblock-collision-band}
 \int_{|v|\le C_0m^{9/100}}
   |\mathcal C_m(v;\mathbf u)-\mathcal G_m(v;\mathbf u)|\dd v
          \le C_{q,B,C_0}\mathfrak e_m.
\end{equation}
Here $v$ and $u_l$ are rescaled frequencies. Each actual collision
frequency is $(v+u_l)/\sqrt m$.
\end{theorem}
\begin{proof}
On $|v|\le2m^{1/200}$ take $\delta=m^{-1/14}/4$ and smooth
$\sigma$ and $h_R$. Put $z_l=(v+u_l)/\sqrt m$. In chronological
order the smoothed pairing is
\[
 \ell\left(L_{z_q}^{m_q}\cdots L_{z_1}^{m_1}
                         (\mathsf S_\delta\sigma)\nu\right),
 \qquad L_z=\mathcal L_{R,\delta,z}.
\]
Use the decomposition in Lemma~\ref{lem:smooth-collision-expansion}.
For a zero power its complementary part means $I-\Pi(z)$.
The complementary bound $C\rho^l$ and its annihilation of $\Pi(z)$
then hold also at $l=0$.

We explain why a short block is harmless. In the expansion of this
fixed product, the term with only complementary factors is at most
$C_{q,B}\delta^{-2}\rho^m$. Every other nonprincipal term has two
adjacent factors of different types. At such a boundary, for example,
\[
 N(z)^{l}\Pi(w)=N(z)^{l}(\Pi(w)-\Pi(z)).
\]
The reversed product is handled in the same way. Analytic projection
bounds give $\|\Pi(w)-\Pi(z)\|\le C\delta^{-2}|w-z|$ on this
ball. All remaining powers are uniformly bounded: their eigenvalue
factors have modulus at most
$\exp(Cm\delta^{-6}\max_l|z_l|^3)\le C$, and the complementary
powers have their uniform bounds. Thus the sum of the mixed terms,
including the initial density norm, is at most
$C_{q,B}\delta^{-4}(|v|+U)/\sqrt m$, where $U=\max_l|u_l|$.
This argument uses no positive lower bound for $m_l/m$.

The product of principal factors has amplitude
$1+O_{q,B}(\delta^{-4}(|v|+U)/\sqrt m)$. The eigenvalue product
has the Gaussian logarithm with $\Gamma_{R,\delta}$, up to a remainder
bounded by $C_q\delta^{-6}m^{-1/2}(|v|+U)^3$.
Changing the density costs $C_B\delta$. For the observable change,
the small-mass variance on each deterministic interval and Minkowski's
inequality give
\[
 \int\left|\sum_l z_l\cdot
       (S_{j_l}-S_{j_{l-1}})(h_R-h_{R,\delta})\right|
                       (\mathsf S_\delta\sigma)\dd\nu
 \le C_{q,B}(|v|+U)\sqrt{\delta(1+|\log\delta|)}.
\]
Changing the covariance costs
$C_q\delta(1+|\log\delta|)(|v|+U)^2$. Since
$U\le m^{1/200}$, integration of these bounds has exactly the six
power budgets in the proof of Theorem~\ref{thm:wide-collision-band}.
The slowest is $1/28-5/200=3/280$. This proves the central estimate,
uniformly in the permitted growing block frequencies.

On $2m^{1/200}\le|v|\le C_0m^{9/100}$ use the fixed orders and
scales
\[
 P=40,\quad Q=29,\quad \delta=m^{-19/100}/4,
                   \quad\epsilon=m^{-50}/4.
\]
Here $|v+u_l|\ge|v|/2$ and $|v+u_l|\le3|v|/2$. We give the
multiblock version of the argument in
Theorem~\ref{thm:high-order-damped-unsmoothing} to account for the
additional frequencies. Let $r_\delta=h_R-h_{R,\delta}$ and
$Z=\max_l|z_l|$. For $Z>0$, write the residual phase as
\[
 X=\sum_l z_l\cdot(S_{j_l}-S_{j_{l-1}})r_\delta.
\]
The all-small-mass moment bound, on each block, followed by
Minkowski's inequality gives
\[
 \int|X|^{2P}\dd\nu
       \le C_{q,P}Z^{2P}\sum_{b=1}^P(m\delta)^b
                                (1+|\log\delta|)^{2P-b}.
\]
Domination gives the same bound with the smoothed initial density.
Expand $e^{iX}$ through degree $79$. Smooth its residual factors
and the initial density at scale $\epsilon$. Each corrected monomial
is a chronological word with at most $80$ smooth multipliers and the
$q-1$ additional block boundaries. The power lengths in block $l$
sum to $m_l$, so its damping is
\[
 \exp\left(-c\sum_l m_l|z_l|^2\right)
                                      \le e^{-c'|v|^2}.
\]
The number of factors is bounded in terms of $q$ and the fixed order,
not in terms of $m$. The fine insertion, fine residual, and residual
moment errors are consequently bounded exactly as in
\eqref{eq:high-order-damped-unsmoothing}, with $|z|$ replaced by $Z$
and a changed fixed constant. The two analytic margins are $3/100$
and $2/25$. The integrated fine errors have decay exponents
$1241/25$ and $303/100$. The largest residual-moment power is
$-1/25$, attained at $b=40$; every smaller block number improves it
by $81/100$. Thus the margin $41/1400$ pays all fixed logarithms.
The Gaussian in \eqref{eq:multiblock-characteristic} has the same
annular damping. Combining the two regions proves the theorem for
large $m$; enlarging the constant covers the finitely many remaining $m$.
\end{proof}

\subsection{A shrinking endpoint with oscillatory path weights}
Let $P_m:\R^4\to\R^3$ be the first three rows of a matrix
$\mathsf A_m$ satisfying
$\|\mathsf A_m\|+\|\mathsf A_m^{-1}\|\le C_0$. Put
\[
 C_m=P_m\Gamma_RP_m^{\mathsf T},\qquad
 Y_m=P_mS_mh/\sqrt m,\qquad
 E_m=\{Y_m\in B_3(\xi,\mathbf h)\},
\]
where $|\xi|\le A$ and
$c_0m^{-2/25}\le h_i\le C_0m^{-2/25}$. These assumptions make
$C_m$ uniformly positive definite. Write $V_3=\prod_i2h_i$ and
\[
 \varepsilon_m=m^{-23/2800}\sqrt{\log(2+m)}.
\]
For $\theta_1,\ldots,\theta_q\in\R^3$ define
\begin{align}
 \mathcal B_m(\boldsymbol\theta;y)
 =\exp\biggl(&i\sum_l d_l\theta_l\cdot y
 -\frac12\sum_{l,r}(d_l\mathbf1_{l=r}-d_ld_r)
                         \theta_l^{\mathsf T}C_m\theta_r\biggr).
 \label{eq:bridge-increment-characteristic}
\end{align}
This is the characteristic function of Gaussian increments of lengths
$d_l$, conditioned on their sum being $y$.

\begin{proposition}[Uniform conditional block characteristic functions]
\label{prop:window-conditional-characteristic}
There is $c(C_0)>0$ such that
$\max_l|\theta_l|\le c(C_0)m^{1/200}$ implies
\begin{align}
 &\left|\mathbb E_{\sigma\nu}
  \left[\left.\exp\left(i\sum_l\theta_l\cdot P_mH_l\right)
                          \right|E_m\right]
           -\mathcal B_m(\boldsymbol\theta;\xi)\right|
       \le C_{q,A,B,c_0,C_0}\varepsilon_m .
 \label{eq:window-conditional-characteristic}
\end{align}
The bound is uniform over all subdivisions, including zero-length
blocks. It concerns the original endpoint indicator, with either
choice of inclusion of its discrete endpoints.
\end{proposition}
\begin{proof}
First retain the Gaussian integral over the endpoint box rather than
its center. Augment $Y_m$ by the fourth row of $\mathsf A_m$ and
truncate that fourth coordinate to $[-T,T]$, $T=m^{1/400}$.
Use the four-dimensional endpoint envelopes of
\eqref{eq:box-lower-envelope} at bandwidth $b=c_1(C_0)m^{9/100}$.
After the change of frequency $v\mapsto\mathsf A_m^{\mathsf T}v$,
the oscillatory comparison is
Theorem~\ref{thm:multiblock-collision-band}, with
$u_l=P_m^{\mathsf T}\theta_l$. Its error for either envelope is
$C T V_3\mathfrak e_m$.

Positivity cannot sandwich an oscillatory numerator. Instead let
$W=\exp(i\sum_l\theta_l\cdot P_mH_l)$, so $|W|=1$, and use
\[
 |\mathbb E[W(\one_{B_3\times[-T,T]}-U)]|
                  \le\mathbb E(U-\one_{B_3\times[-T,T]})
                  \le\mathbb E(U-L).
\]
The latter is bounded by the unweighted Fourier comparison plus the
Gaussian envelope excess. The Gaussian weighted endpoint measure
also has total variation bounded by its unweighted Gaussian marginal,
because its conditional characteristic function has modulus at most
one. Thus the same replacement is valid on the Gaussian side.
Integrating the fourth Gaussian coordinate against its decay, as in
\eqref{eq:stationary-projection-budget}, gives an error
$C V_3[\sum_i(bh_i)^{-1}+b^{-1}]$, without a factor $T$.
The $128$th collision moment bounds the omitted actual tail by
$CT^{-128}$; $|W|=1$ makes this estimate unchanged. The Gaussian
tail is at most $C V_3e^{-cT^2}$.

It follows that the numerator differs from
\[
 \int_{B_3(\xi,\mathbf h)}g_{C_m}(y)
                         \mathcal B_m(\boldsymbol\theta;y)\dd y
\]
by at most $C V_3\varepsilon_m$. The three exponent margins are
$23/2800$, $1/100$ and $2/25$, exactly as in the unweighted projection.
The Gaussian formula follows by subtracting $d_l$ times the total
from independent centered increments of covariance $d_l C_m$.
These differences have covariance
$(d_l\mathbf1_{l=r}-d_ld_r)C_m$ and are independent of the total.

Theorem~\ref{thm:projected-collision-windows} gives
$\mathbb P_{\sigma\nu}(E_m)=V_3g_{C_m}(\xi)(1+O(\varepsilon_m))$,
including a lower bound $c_A V_3$. On the box the change in
\eqref{eq:bridge-increment-characteristic} from its value at $\xi$
is at most $C m^{-2/25}\max_l|\theta_l|$; the nonoscillatory
Gaussian density changes by $O_A(m^{-2/25})$. The first error is
$O(m^{-3/40})$, hence smaller than $\varepsilon_m$ even for the
permitted growing frequencies. Divide by the exact denominator.
This proves the proposition, without taking a trace of a
four-dimensional $L^1$ Fourier bound.
\end{proof}

\paragraph{Uniformity needed below.}
The number of blocks in these statements is fixed. Their boundaries
may nevertheless depend on $m$ and may approach each other or an
endpoint. The permitted growth of the test frequencies and this
boundary uniformity are both needed to prove tightness after division
by a probability of order $m^{-6/25}$. Finite-dimensional convergence
alone would not justify such a division in a path-space assertion.

\section{Conditioned physical paths and their Gaussian bridges}
\label{sec:stationary-window-bridges}
We identify the common conditional law of
Theorem~\ref{thm:stationary-physical-windows}. The conditioning remains
an exact observation at deterministic physical time, with the same
shrinking standardized endpoint box. The result is a functional limit,
not merely a comparison of two conditional measures. Its proof pays
for the rare denominator in the tightness and clock estimates.

For a positive definite $3\times3$ matrix $C$ and $\xi\in\R^3$, let
\begin{equation}\label{eq:Gaussian-bridge-definition}
 \mathbb B_{C,\xi}(s)=s\xi+C^{1/2}(W(s)-sW(1)),\qquad 0\le s\le1,
\end{equation}
where $W$ is standard three-dimensional Brownian motion. This law on
$\mathcal C=C([0,1],\R^3)$ has mean $s\xi$ and covariance
$(\min\{s,t\}-st)C$. We use the supremum norm on $\mathcal C$
and the bounded-Lipschitz distance on probability laws.

\subsection{Tightness after conditioning on the endpoint}
Under the hypotheses of
Proposition~\ref{prop:window-conditional-characteristic}, let $X_m$
be the polygonal interpolation with
\[
 X_m(j/m)=P_m S_jh/\sqrt m,\qquad 0\le j\le m.
\]
Retain $E_m$, $C_m$ and $\varepsilon_m$ from that proposition.

\begin{lemma}[A trigonometric tail test]
\label{lem:bridge-trigonometric-tail}
For each positive integer $p$ there are $c_p,C_p>0$ such that
\[
 c_p\one_{\{|x|\ge1\}}
 \le F_p(x):=\int_0^1(1-\cos(ux))^p\dd u
 \le C_p|x|^{2p}.
\]
The integrand is a trigonometric polynomial of degree $p$ in $ux$.
\end{lemma}
\begin{proof}
The upper bound follows from $1-\cos v\le v^2/2$.
The function $F_p$ is continuous, even, and strictly positive away
from zero. On $1\le|x|\le4\pi$ it has a positive minimum.
For $|x|\ge4\pi$, write it as the average of the nonnegative,
$2\pi$-periodic function $(1-\cos v)^p$ over an interval of length
$|x|$. The complete periods occupy at least half that length and
have a positive mean. This gives the lower bound. Expanding the
power gives the last assertion.
\end{proof}

\begin{theorem}[A collision bridge under an unsmoothed endpoint window]
\label{thm:collision-window-bridge}
Uniformly in $R$, the initial densities \eqref{eq:bridge-initial-class},
the matrices and windows in
Proposition~\ref{prop:window-conditional-characteristic},
\begin{equation}\label{eq:collision-window-bridge}
 d_{\rm BL}\left(\operatorname{Law}_{\sigma\nu}(X_m\mid E_m),
                         \operatorname{Law}(\mathbb B_{C_m,\xi})\right)
                                      \longrightarrow0.
\end{equation}
No independence of the actual collision increments is assumed.
\end{theorem}
\begin{proof}
For a fixed finite list of times, use
\eqref{eq:window-conditional-characteristic} with bounded test
frequencies and the corresponding consecutive blocks. Rounding times
to collision indices and polygonal interpolation change a fixed
finite list by $O(m^{-1/2})$, since $h_R$ and $P_m$ are bounded.
The displayed Gaussian increment formula proves the finite-dimensional
bridge limit, uniformly in the parameters.

For tightness take any collision-grid increment of length
$\Delta=(k-j)/m$. Apply Lemma~\ref{lem:bridge-trigonometric-tail}
with $p=4$ to each of its three coordinates divided by $a/\sqrt3$.
The characteristic functions involved have three consecutive blocks
and test frequencies at most $4\sqrt3/a$. Therefore, whenever
$4\sqrt3/a\le c(C_0)m^{1/200}$,
\begin{equation}\label{eq:conditional-increment-tail}
 \mathbb P_{\sigma\nu}\{|X_m(k/m)-X_m(j/m)|>a\mid E_m\}
          \le C_A a^{-8}\Delta^4+C_A\varepsilon_m.
\end{equation}
Indeed the corresponding Gaussian bridge increment has mean
$\Delta\xi$, covariance $\Delta(1-\Delta)C_m$, and eighth moment
at most $C_A\Delta^4$. The finite trigonometric coefficient sum is
fixed, so its Fourier comparison costs $C_A\varepsilon_m$,
not a negative power of the endpoint probability.

Set
\begin{equation}\label{eq:bridge-tightness-scales}
 \eta=\frac1{200},\qquad \alpha=\frac18,\qquad
 J_m=\lfloor\eta\log_2m\rfloor.
\end{equation}
Use the nested grids with vertices
$\lfloor ml2^{-j}\rfloor/m$, $0\le l\le2^j$, for
$j\le J_m$. Each adjacent increment has length at most $2^{1-j}$.
Fix $\epsilon>0$ and assign threshold
$a_j=\epsilon2^{-\alpha j}$. For all sufficiently large $m$, all
frequencies in \eqref{eq:conditional-increment-tail} are permitted,
because their largest order is
$C_\epsilon m^{\eta\alpha}=C_\epsilon m^{1/1600}$,
strictly below $m^{1/200}$. A union bound over levels
$j_0\le j\le J_m$ and their adjacent increments gives
\begin{equation}\label{eq:bridge-coarse-grid-bound}
 C_A\epsilon^{-8}\sum_{j\ge j_0}2^{-2j}
       +C_A2^{J_m}\varepsilon_m
 \le C_A\epsilon^{-8}2^{-2j_0}
       +C_A m^{-9/2800}\sqrt{\log(2+m)}.
\end{equation}
Here $4(1-2\alpha)-1=2$ and $23/2800-1/200=9/2800$.
On the complementary event, dyadic chaining bounds the modulus of
the polygon through the finest grid vertices by a constant times
$\epsilon2^{-\alpha j_0}$ on intervals of length $2^{-j_0}$:
pass from each vertex to its ancestor at level $j_0$ and sum at most
two increments per finer level. Adjacent ancestors add at most a
fixed number of level-$j_0$ increments. Rounding the nested grids
only changes these constants.

It remains to control the motion inside the finest cells. This is
where ordinary tightness without a rare-event budget would be
insufficient. Let $\widetilde X_m$ be the polygon through those
vertices. Lemma~\ref{lem:fixed-even-maximal}, at the fixed order
$128$, gives on each cell, under the bounded density $\sigma$,
\[
 \mathbb E\sup_{s\text{ in the cell}}
          |X_m(s)-\widetilde X_m(s)|^{128}
                                     \le C2^{-64J_m}.
\]
The bound follows from deterministic-window partial sums and their
endpoint increment; the largest cell has $O(m2^{-J_m})$ collisions.
There are $2^{J_m}$ cells. Since
$\mathbb P_{\sigma\nu}(E_m)\ge c_A m^{-6/25}$,
\begin{equation}\label{eq:bridge-fine-grid-bound}
 \mathbb P_{\sigma\nu}
 \{\|X_m-\widetilde X_m\|_\infty>\epsilon\mid E_m\}
 \le C_A\epsilon^{-128}m^{6/25}2^{-63J_m}
 \le C_A\epsilon^{-128}m^{-3/40}.
\end{equation}
The strict exponent is $63/200-6/25=3/40$.
First let $m\to\infty$ in \eqref{eq:bridge-coarse-grid-bound}
and \eqref{eq:bridge-fine-grid-bound}, and then let $j_0\to\infty$.
They prove conditional tightness in $\mathcal C$.

Finally suppose uniformity in \eqref{eq:collision-window-bridge}
failed. Along a violating sequence select a subsequence for which
$C_m$ and $\xi$ converge; these lie in compact sets of uniformly
positive matrices and bounded vectors. Tightness and the uniform
finite-dimensional estimate identify every subsequential law as the
bridge with those limiting parameters. The comparison bridge laws
also converge: couple them by the same $W$ in
\eqref{eq:Gaussian-bridge-definition} and use continuity of the
positive matrix square root. This contradicts the violating
bounded-Lipschitz distance and completes the proof.
\end{proof}

\subsection{The whole physical clock on the same stationary space}
Use the stationary suspension $(\mathcal X_R,\mathbb P_R)$ and the
four exact endpoint events $A_t^j$ of
\eqref{eq:stationary-four-events}, with half-width $t^{-2/25}$.
Let $\mathsf X_{t,R}$ be the continuous polygon through the values
\[
 t^{-1/2}V_R(k),\qquad k=0,1,\ldots,\lfloor t\rfloor,
 \quad\hbox{and }t^{-1/2}V_R(t),
\]
at times $k/t$ and $1$, where
$V_R(u)=(W_R(u),C_R(u)-u/\bar\tau_R)$.
The last point is included only once when $t$ is an integer.
Its distance from the original physical step path is $O(t^{-1/2})$:
the minimum flight length is positive and a unit physical time interval
has bounded displacement and bounded collision count.

Let $m=m_R(t)$ and let $\mathsf Y_{t,R}$ be the collision polygon
with values $t^{-1/2}B_RS_jh_R(x)$ at times $j/m$.
The state $x$ is the actual outgoing collision preceding the stationary
observation, not an independently chosen initial state.

\begin{lemma}[A uniform pathwise physical-to-collision comparison]
\label{lem:whole-physical-clock-coupling}
For every fixed $P>0$,
\begin{equation}\label{eq:whole-physical-clock-coupling}
 \mathbb P_R\{\|\mathsf X_{t,R}-\mathsf Y_{t,R}\|_\infty
                                  >C_Pt^{-1/8}\}\le C_Pt^{-P}.
\end{equation}
Consequently the same exceptional probability under any of the four
endpoint conditionings is at most $C_{A,P}t^{-P+6/25}$.
\end{lemma}
\begin{proof}
Use Proposition~\ref{prop:stationary-observation-coupling} at each
integer observation length $k\ge t^{1/4}$ and at $t$. In unnormalized
variables its physical-to-collision error is $C_Q k^{3/8}$ outside
probability $C_Q k^{-Q}$. Summing over at most $t+2$ observation
lengths gives $C_Qt^{1-Q/4}$. Take $Q>4(P+1)$.
Every retained error is at most $C_Qt^{3/8}$.

For $k<t^{1/4}$ both the physical vector and its deterministic collision
reference have size at most $C(1+k)$: use bounded speed and cell-position
terms, the positive minimum flight length, and boundedness of $h_R$.
Their difference is therefore at most $Ct^{1/4}$, which is smaller
than the same threshold. At the physical grid point $k/t$ the index
$\lfloor mk/t\rfloor$ differs from $m_R(k)$ by at most two.
This changes the collision sum by a bounded amount. Between adjacent
physical grid points, the physical interpolation and the collision
polygon each change by a bounded amount before normalization.
Division by $\sqrt t$ proves \eqref{eq:whole-physical-clock-coupling}.
No union over a continuum of clocks has been taken. The lower bounds
in \eqref{eq:stationary-window-denominator} prove the last assertion.
\end{proof}

\begin{theorem}[Stationary physical bridges in exact endpoint windows]
\label{thm:stationary-physical-bridge}
For every $A<\infty$, uniformly for $R\in I$, $|\xi|\le A$, and
$j\in\{{\rm coll},{\rm ret},{\rm last},{\rm phys}\}$,
\begin{equation}\label{eq:stationary-physical-bridge}
 d_{\rm BL}\left(\operatorname{Law}_{\mathbb P_R}
          (\mathsf X_{t,R}\mid A_t^j),
               \operatorname{Law}(\mathbb B_{\mathcal V_R,\xi})\right)
                                      \longrightarrow0.
\end{equation}
In particular the actual physical observation conditions the entire
physical displacement and collision-fluctuation path to a Gaussian
bridge with terminal value $\xi$.
\end{theorem}
\begin{proof}
For $A_t^{\rm coll}$ apply
Theorem~\ref{thm:collision-window-bridge} under the exact collision
marginal $\sigma_R\nu$, $\sigma_R=\tau_R/\bar\tau_R$, with
$P_m=\sqrt{m/t}\,B_R$. This matrix has the uniformly bounded
augmentation used in Theorem~\ref{thm:stationary-physical-windows}.
Its covariance is
$C_m=(m/t)B_R\Gamma_RB_R^{\mathsf T}=\mathcal V_R+O(t^{-1})$,
and its endpoint window widths are comparable to $m^{-2/25}$.
The resulting polygon is exactly $\mathsf Y_{t,R}$.
Lemma~\ref{lem:whole-physical-clock-coupling}, with $P=2$, replaces
it by $\mathsf X_{t,R}$ under this conditioning, since the conditional
exceptional probability is $O(t^{-2+6/25})$. Continuity of Gaussian
bridge laws in their positive covariance gives the stated limit.

Theorem~\ref{thm:stationary-physical-windows} and the triangle
inequality bound the other three conditional initial laws against the
collision-conditioned law in total variation by
\[
 C_A t^{-23/2800}\sqrt{\log(2+t)}.
\]
Pushforward by the same physical path cannot increase this distance. This proves all four assertions. The return-conditioned
statistic remains the actual physical path on the original stationary
suspension; neither its initial section point nor its elapsed age is
resampled.
\end{proof}

\subsection{Actual denominators for a class of path selectors}
\begin{corollary}[Path-weighted window denominators]
\label{cor:bridge-path-weighted-denominator}
Let $F:\mathcal C\to[0,\infty)$ be fixed, bounded and continuous.
Uniformly on the parameter sets of the preceding theorem,
\begin{align}
 \mathbb E_R[F(\mathsf X_{t,R})\one_{A_t^j}]
  =8t^{-6/25}g_{\mathcal V_R}(\xi)
        \left(\mathbb E F(\mathbb B_{\mathcal V_R,\xi})+o(1)\right).
 \label{eq:bridge-path-weighted-denominator}
\end{align}
Whenever the bridge expectation has a positive uniform lower bound,
the conditional path law with weight $F(\mathsf X_{t,R})\one_{A_t^j}$
converges to the bridge law weighted by $F$. The relative comparison
between any two of these weighted endpoint events retains the rate
$O_A(t^{-23/2800}\sqrt{\log(2+t)})$, with a constant depending on
the stated lower bound and $\|F\|_\infty$.
\end{corollary}
\begin{proof}
Uniform weak convergence and uniform tightness give convergence of
expectations of every fixed bounded continuous function: on a common
compact path set approximate it uniformly by bounded Lipschitz
functions, and use boundedness outside that set. Multiply this
conclusion by \eqref{eq:stationary-window-denominator}.
For the reweighted law apply it also to $FG$ for a bounded continuous
path test $G$, and divide by the positive denominator.
For the comparison, the difference of the two unnormalized densities
is supported on the original symmetric difference and is at most
$\|F\|_\infty$ there. Normalization costs at most twice its integral
divided by either weighted mass. The relative event theorem and
\eqref{eq:bridge-path-weighted-denominator} give the claimed rate.
\end{proof}

\begin{corollary}[A concrete nonsingle-time indicator class]
\label{cor:bridge-cylinder-events}
Fix $0<s_1<\cdots<s_r<1$ and bounded boxes
$D_1,\ldots,D_r\subset\R^3$ with nonempty interiors. Define the
exact path event
\[
 D_t=\bigcap_{l=1}^r\{t^{-1/2}V_R(ts_l)\in D_l\}.
\]
Then, uniformly for $R\in I$, $|\xi|\le A$, and all four endpoint
events,
\begin{equation}\label{eq:bridge-cylinder-denominator}
 \mathbb P_R(D_t\cap A_t^j)
 =8t^{-6/25}g_{\mathcal V_R}(\xi)
       \mathfrak b_R(\xi;\mathbf s,\mathbf D)(1+o(1)),
\end{equation}
where
\[
 \mathfrak b_R(\xi;\mathbf s,\mathbf D)
  =\int_{\prod_lD_l}
      g_{\Sigma_R}\bigl((y_l-s_l\xi)_{l=1}^r\bigr)
                           \prod_l\dd y_l,\quad
 (\Sigma_R)_{lk}=(\min\{s_l,s_k\}-s_ls_k)\mathcal V_R.
\]
The factor $\mathfrak b_R$ has a positive uniform lower bound.
The conditional physical path converges to the bridge restricted to
these cylinder boxes. The conditional initial laws for
$D_t\cap A_t^j$ and $D_t\cap A_t^{j'}$ are within
$O_A(t^{-23/2800}\sqrt{\log(2+t)})$ in total variation.
\end{corollary}
\begin{proof}
For distinct interior times the scalar bridge covariance matrix is
positive definite. This follows, for example, by expressing its values
as independent Gaussian increments conditioned on their total: no
nonzero linear combination of the interior partial sums is a multiple
of that total. Uniform ellipticity of $\mathcal V_R$ gives positive
definiteness of $\Sigma_R$. The Gaussian density is uniformly positive
on any fixed compact subbox of $\prod_l\operatorname{int}D_l$, for
$R\in I$ and $|\xi|\le A$. This proves the lower bound.
The cylinder boundary has zero measure under every limiting bridge.
The actual observations differ from the interpolated values by
$O(t^{-1/2})$ uniformly. Enlarging and contracting the fixed cylinder
boxes therefore shows that the path convergence applies to the exact
physical indicator and to that indicator times a bounded continuous
path test. Uniformity follows by
the same compact-parameter subsequence argument; every subsequential
Gaussian has this boundary-null property. These observations prove
the denominator and the restricted bridge limit. Finally
$(D_t\cap A_t^j)\mathbin\triangle(D_t\cap A_t^{j'})
\subset A_t^j\mathbin\triangle A_t^{j'}$; divide the inherited
relative bound by the newly proved positive cylinder factor.
\end{proof}

\paragraph{Relation to the raw conditioning problem.}
The result supplies an actual conditional functional limit and explicit
weighted denominators for fixed interior cylinder observations, without
assuming their long pullbacks have controlled variation. The endpoint
width remains $t^{21/50}$ in physical coordinates. Fixed lattice labels,
microscopic time windows, the pointwise common signed raw correction,
and the complete fixed-return Fourier complement still require their
own estimates. No restriction on those original target events has been
silently imposed: the theorem above concerns the displayed window class
and leaves the original raw mixed-density endpoint intact.

\section{A quantitative geometric extraction on every finite itinerary}
\label{sec:geometric-raw-extraction}
The power--logarithm extraction gives a residual at each finite packet,
but does not bound its derivative norm as the collision cutoff grows.
We construct a different exact extraction. Its removed measure has a
quantified total variation, and its residual has a uniform, explicit
second-derivative budget. The construction takes place before coarea,
on the physical collision cylinder. In particular it neither estimates
individual prepared jets nor assumes separated critical values.

Write $L_m=S_{m,R}\tau_R$, and use $(\alpha,p)$ throughout for derivatives
and area. All constants in this section are uniform in $R\in I$.
Put
\[
 R_-=9/20,\qquad R_+=47/100,\qquad \ell_0=3/50,
 \qquad H_0=1+R_+/\ell_0=53/6.
\]
The integer $m$ is a number of physical collisions, not returns.

\subsection{A direction transverse to every regular roof level}
Lemma~\ref{lem:self-contained-collision-differential} gives a
self-contained derivation of the coordinate differential and action
used below, with both signs and the determinant in invariant coordinates.
\begin{lemma}[A uniform nonstationary direction]\label{lem:roof-transverse-direction}
On a regular $m$-collision branch, $m\ge1$, write
\[
 D T_R^m=\begin{pmatrix}A_m&B_m\\ C_m&D_m\end{pmatrix}.
\]
Then $B_m\ne0$ and $0<A_m/B_m\le H_0$. The total roof satisfies
\begin{equation}\label{eq:uniform-roof-gradient}
 |\nabla L_m|\ge \frac{R_-}{\sqrt{1+H_0^2}}|p|.
\end{equation}
On $p\ne0$ the vector field
\begin{equation}\label{eq:roof-transverse-field}
 X_m=-\frac1{Rp}\left(\partial_\alpha-
                      \frac{A_m}{B_m}\partial_p\right)
 \quad\hbox{satisfies}\quad X_m L_m=1.
\end{equation}
No incidence bound at the interior collisions is required for these
three assertions.
\end{lemma}
\begin{proof}
Let $c=\sqrt{1-p^2}$, $c'=\sqrt{1-(p')^2}$ and $v=\tau_R/R$ for
one regular flight. In the coordinates used here its differential is
\begin{equation}\label{eq:positive-collision-matrix}
 D T_R=-\begin{pmatrix}
 (v+c)/c'&v/(cc')\\
 v+c+c'&(v+c')/c
 \end{pmatrix}=:-P.
\end{equation}
One obtains this formula by differentiating the straight flight and
its reflection: in arc length and outgoing angle the four entries,
before multiplication by $-1/c'$, are
$\tau_R/R+c$, $\tau_R$, $\tau_R/R^2+(c+c')/R$, and
$\tau_R/R+c'$. Change variables by $r=R\alpha$, $p=\sin\varphi$.
This is also the fixed-table specialization of the differential in
\cite[Section 3.1]{SYZ}. Its determinant in $(\alpha,p)$ is one.

All entries of $P$ are positive. The ratios of the first to the second
entry in either row are
\[
 c+\frac{c^2}{v},\qquad c+\frac{c^2}{v+c'},
\]
respectively, and are at most $H_0$. The first row of the positive
product $P_{m-1}\cdots P_0$ is a positive linear combination of the
two rows of $P_0$. Its entry ratio is consequently a weighted average
of these two ratios, with positive weights. This proves the assertion
about $A_m/B_m$, including $m=1$.

The exact action identity of Lemma~\ref{lem:exact-collision-action}
gives
\[
 \partial_\alpha L_m=R(p_m A_m-p),\qquad
 \partial_p L_m=Rp_m B_m.
\]
Subtracting $A_m/B_m$ times the second identity from the first gives
$-Rp$. Cauchy--Schwarz proves \eqref{eq:uniform-roof-gradient}; the
same calculation proves \eqref{eq:roof-transverse-field}.
\end{proof}

\subsection{A one-collision regularity cutoff}
\begin{lemma}[Polynomial sublevel sets in a rectangle]
\label{lem:elementary-polynomial-sublevel}
Fix $D\ge1$ and $c_0>0$. A real polynomial $P$ in two variables, of
degree at most $D$ in each variable and with maximum coefficient
modulus at least $c_0$, satisfies
\begin{equation}\label{eq:elementary-polynomial-sublevel}
 |\{(x,y)\in[-1,1]^2:|P(x,y)|\le h\}|
                  \le C_{D,c_0}h^{1/(2D)},\qquad 0<h<1.
\end{equation}
An upper bound on the other coefficients is not needed.
\end{lemma}
\begin{proof}
For a one-variable polynomial of degree at most $D$, a sublevel set of
length $s>0$ contains $D+1$ points whose pairwise distances are at least
$c_Ds$. For example choose successive quantiles of its restricted
Lebesgue measure, discarding a small part at the endpoints. Lagrange
interpolation on these points bounds every coefficient by $C_D h s^{-D}$
when the polynomial is at most $h$ there. It follows that a coefficient
of modulus at least $b>0$ gives the one-variable bound
$C_D(h/b)^{1/D}$; when this exceeds two the bound is automatic.

Write $P(x,y)=\sum_j a_j(y)x^j$ and choose $a_j$ with a coefficient
of modulus at least $c_0$. The set where $|a_j(y)|\le h^{1/2}$ has
length at most $C_{D,c_0}h^{1/(2D)}$. On its complement, apply the
one-variable bound to $P(\cdot,y)$ with $b=h^{1/2}$. Integrating over
$y$ proves \eqref{eq:elementary-polynomial-sublevel}. Constant
polynomials cause no difficulty in either step.
\end{proof}

Fix an even $\theta\in C^\infty(\R;[0,1])$, zero on $[-1,1]$ and one
outside $(-2,2)$. A finite, parameter-independent set $\mathcal D$ of
nonzero lattice centers contains every possible next collision center,
by the uniform finite horizon. For $d\in\mathcal D$ put
\begin{equation}\label{eq:all-candidate-discriminants}
 \Delta_{d,R}(\alpha,p)
   =R^2-\bigl((d-Rn_\alpha)\mathbin{\times}
          (\sqrt{1-p^2}n_\alpha+pt_\alpha)\bigr)^2.
\end{equation}
The cross denotes the scalar planar cross product. We use all candidate
discriminants, including candidates that are not selected by the actual
flight. This prevents a discontinuity from being missed on the side
where a tangent disk ceases to be the first collision.
Lemma~\ref{lem:complete-first-hit-guard} specifies the finite candidate
set, four bounded charts and uniform coefficient minimum, and proves
that every label transition is covered.

\begin{lemma}[Uniform regularity cutoffs]\label{lem:uniform-collision-cutoff}
There are $\varepsilon_0>0$ and smooth functions
$\chi_{R,\varepsilon}:M\to[0,1]$, $0<\varepsilon\le\varepsilon_0$,
vanishing on neighborhoods of the one-collision singularities, grazing
endpoints and every boundary of $Y_R^*$, such that, with $\kappa_0=1/16$,
\begin{equation}\label{eq:collision-cutoff-measure}
 \nu\{\chi_{R,\varepsilon}\ne1\}\le C\varepsilon^{\kappa_0}.
\end{equation}
On the support of this cutoff and its first three derivatives, every
selected one-collision branch and roof, as well as the cutoff itself,
have $C^3(\alpha,p)$ bounds at most $(C/\varepsilon)^C$.
The actual section membership and lattice label are locally constant
where needed. Values at the excluded singularities are fixed to zero.
\end{lemma}
\begin{proof}
Let $\mathcal A_R$ contain the angular endpoints and $\mathcal B_R$ the
$p$-endpoints of all the finitely many rectangles defining $Y_R^*$.
Set $c(p)=\sqrt{1-p^2}$ and take the product
\begin{equation}\label{eq:explicit-collision-cutoff}
 \theta(c(p)/\varepsilon)
 \prod_{d\in\mathcal D}\theta(\Delta_{d,R}/\varepsilon)
 \prod_{a\in\mathcal A_R}\theta(\sin(\alpha-a)/\varepsilon)
 \prod_{b\in\mathcal B_R}\theta((p-b)/\varepsilon).
\end{equation}
The extra angular zeros at antipodal points only enlarge the removed
set by finitely many strips. All factors are globally defined on the
angular circle. The first factor makes the product identically zero
near $p=\pm1$, so its zero extension is smooth at these endpoints.

Here is an explicit uniform sublevel estimate for each discriminant.
On a bounded rational chart for $\alpha$, write $t$ for its half-angle
coordinate, and put $s=\tan(\varphi/2)$, so that $|s|\le1$,
$c=(1-s^2)/(1+s^2)$ and $p=2s/(1+s^2)$. The cross product in
\eqref{eq:all-candidate-discriminants} has denominator
$(1+t^2)(1+s^2)$ and numerator of total degree at most four. Thus
$\Delta_{d,R}$ has a bounded positive denominator and a numerator of
total degree at most eight. That numerator is never the zero
polynomial: at any fixed initial boundary point, $|d-Rn_\alpha|>R$,
and its cross product with a varying unit direction cannot have
constant square $R^2$ on an interval of directions. Compactness of $I$
and finiteness of $\mathcal D$ and the charts give a positive lower
bound on the maximum coefficient modulus. Apply
Lemma~\ref{lem:elementary-polynomial-sublevel} with $D=8$.
The coordinate Jacobian of $\nu$ in these bounded $(t,s)$ charts is
bounded above. This proves the bound $C\varepsilon^{1/16}$ for the
discriminant factors. The remaining factors remove finitely many
strips of measure $O(\varepsilon)$, or $O(\varepsilon^2)$ at grazing.
A union bound proves \eqref{eq:collision-cutoff-measure}.

Every possible change of first-hit label occurs at a candidate
tangency: a positive entry root cannot cross zero because distinct
disks are separated, and two distinct disks cannot have the same
entry point at the same positive time. Those tangencies are among the
zero sets in \eqref{eq:all-candidate-discriminants}. On a selected
branch with $\chi_{R,\varepsilon}>0$, its entry root is
\[
 (d-Rn_\alpha)\cdot v-\sqrt{\Delta_{d,R}},
 \qquad \Delta_{d,R}>\varepsilon,
\]
with $c(p)>\varepsilon$. Differentiating this formula and the Euclidean
reflection formula through order three gives a fixed polynomial loss
in $\varepsilon^{-1}$. The outgoing point lies on a circle of radius
at least $R_-$, so local angular coordinates add only fixed losses.
Changing from $\varphi$ to $p$ costs further fixed powers of
$c(p)^{-1}$. Differentiating the factors of
\eqref{eq:explicit-collision-cutoff} has the same property. The estimate
also holds on their derivative supports, because a smooth factor which
vanishes on an interval is flat at its endpoints. Section membership
can change only on the explicitly removed rectangle boundaries.
\end{proof}

For $L\ge1$ define, on regular orbits and then by zero extension,
\begin{equation}\label{eq:whole-itinerary-cutoff}
 \Gamma_{L,R,\varepsilon}(x)
  =\theta(p(x)/\varepsilon)
                    \prod_{j=0}^{L}\chi_{R,\varepsilon}(T_R^jx).
\end{equation}
The first factor removes a thin normal strip at the initial state.
This one strip, rather than neighborhoods chosen separately at every
critical value, suffices for \eqref{eq:uniform-roof-gradient}.

\begin{lemma}[Whole-itinerary bounds]\label{lem:whole-itinerary-cutoff}
The function in \eqref{eq:whole-itinerary-cutoff} is globally $C^3$
after zero extension across its excluded boundaries, and
\begin{align}
 \int(1-\Gamma_{L,R,\varepsilon})\dd\nu_R^*
          &\le C(L+1)\varepsilon^{1/16},
             \label{eq:whole-itinerary-mass}\\
 \|\Gamma_{L,R,\varepsilon}\|_{C^3}
          &\le \exp\{C(L+1)\log(C/\varepsilon)\}.
             \label{eq:whole-itinerary-derivatives}
\end{align}
On its derivative supports, for $1\le m\le L$, the field $X_m$ has
$C^2$ norm bounded by the right side of
\eqref{eq:whole-itinerary-derivatives}, after changing $C$.
\end{lemma}
\begin{proof}
Use $1-\prod b_j\le\sum(1-b_j)$ for $0\le b_j\le1$, invariance of
$\nu$, and $\nu_R^*\le c_*^{-1}\nu$. The initial normal strip has
measure $O(\varepsilon)$, and each of the other $L+1$ terms is bounded
by \eqref{eq:collision-cutoff-measure}. This proves
\eqref{eq:whole-itinerary-mass} without counting itinerary cells.

On a nonzero derivative support all earlier branches are regular and
all the one-step bounds of Lemma~\ref{lem:uniform-collision-cutoff}
apply. The chain rule through fixed order three, iterated $j$ times,
gives at most $(C/\varepsilon)^{C(j+1)}$. The product rule for $L+2$
factors contributes at most a fixed power of $L+2$, absorbed by
$\exp(C(L+1))$. At any boundary of a finite itinerary, an earlier
cutoff factor vanishes on a neighborhood on both sides. At a section
boundary the corresponding factor at that time does the same.
Consequently these piecewise derivatives extend by zero and give the
global estimate.

Every entry of the positive matrix in
\eqref{eq:positive-collision-matrix} is at least
$v_0=\ell_0/R_+=6/47$. In particular $|B_m|\ge v_0^m$. The quotient
rule, the $C^3$ bounds for $T_R^m$, and this lower bound give a
$C^2$ bound for $A_m/B_m$ of the same exponential form. On the relevant
supports $|p|\ge\varepsilon$. Formula~\eqref{eq:roof-transverse-field}
therefore gives the assertion about $X_m$.
\end{proof}

\subsection{An exact residual for bounded-variation return marks}
Fix $q\ge0$, indices $0\le k_1,\ldots,k_q\le n$, and complex section
functions $a_1,\ldots,a_q$ with bounded supremum and BV zero extensions
$a_i^0$ to $M$. The number $q$ is fixed; repeated marks are allowed.
Put
\begin{equation}\label{eq:geometric-weight-budgets}
 \begin{split}
 w(x)&=\prod_{i=1}^q a_i((F_R^*)^{k_i}x),\\
 M&=\prod_{i=1}^q\|a_i\|_\infty,\qquad
 V=\sum_{i=1}^q\|a_i^0\|_{\mathrm{BV}}
                         \prod_{j\ne i}\|a_j\|_\infty.
 \end{split}
\end{equation}
For $q=0$ these conventions give $w=1$, $M=1$ and $V=0$.
All measures below use the normalized section law $\nu_R^*$ and retain
the exact weight $w$. For the earlier unnormalized single-mark convention,
choose $w=c_*a\circ(F_R^*)^k$.

\begin{theorem}[All-itinerary geometric extraction]
\label{thm:geometric-extraction-budget}
For $1\le n\le L$ and $0<\varepsilon\le\varepsilon_0$, there is an exact
decomposition of the actual finite-count measure
\begin{equation}\label{eq:geometric-exact-decomposition}
 \mu_{n,R}^{w,L}=E_{n,R,\varepsilon}^{w,L}
                         +Q_{n,R,\varepsilon}^{w,L}
\end{equation}
with compactly supported mixed densities $e_\varepsilon$ and
$q_\varepsilon$, respectively. The residual satisfies
\begin{align}
 \|Q_{n,R,\varepsilon}^{w,L}\|_{\TV}&\le M,\notag\\
 \|E_{n,R,\varepsilon}^{w,L}\|_{\TV}
     &\le C_q\{\varepsilon V+M(L+1)\varepsilon^{1/16}\},
                   \label{eq:geometric-extraction-mass}\\
 \mathcal A_{2,\varepsilon}(n,L,R,w)
  :=\sum_{\ell\in\Z^3}\|\partial_t^2q_\varepsilon(\ell,\cdot)\|_1
     &\le C_q M\exp\{C_q(L+1)\log(C_q/\varepsilon)\}.
                   \label{eq:geometric-derivative-budget}
\end{align}
Each residual component belongs to $W^{2,1}(\R)$. The constants do not
depend on $n$, the marks or the radius. When $w=1$, both terms of
\eqref{eq:geometric-exact-decomposition} are positive measures.
No subanalytic assumption on the marked functions is imposed.
\end{theorem}
\begin{proof}
Set $a_{i,\varepsilon}=\mathsf S_{\varepsilon/8}a_i^0$, using the
mean-preserving smoothing of Lemma~\ref{lem:bv-smoothing}, and let
$w_\varepsilon$ be the corresponding product at the same actual marks.
Differentiating the same smoothing kernel three times gives
$\|a_{i,\varepsilon}\|_{C^3}\le C\varepsilon^{-3}\|a_i\|_\infty$;
composition with $p=\sin\varphi$ is not needed for this flat-coordinate bound.
Define
\begin{equation}\label{eq:geometric-residual-definition}
 Q_{n,R,\varepsilon}^{w,L}=(J_{n,R})_*
 \bigl(\one_{\{N_{n,R}\le L\}}\Gamma_{L,R,\varepsilon}
                                     w_\varepsilon\nu_R^*\bigr),
 \qquad E_{n,R,\varepsilon}^{w,L}=\mu_{n,R}^{w,L}-Q_{n,R,\varepsilon}^{w,L}.
\end{equation}
An empty product is not smoothed. Supremum contraction gives
$|w_\varepsilon|\le M$. Telescoping the products and invariance of
$F_R^*$ give
\[
 \int|w-w_\varepsilon|\dd\nu_R^*
 \le c_*^{-1}\sum_i\|a_i^0-a_{i,\varepsilon}\|_{L^1(\nu)}
                        \prod_{j\ne i}\|a_j\|_\infty
 \le C_q\varepsilon V.
\]
Combining this with \eqref{eq:whole-itinerary-mass} proves the two
variation bounds. When $w=1$ the residual source weight is exactly
$\Gamma_{L,R,\varepsilon}\one_{\{N_n\le L\}}$, which lies between zero
and the source weight of $\mu^{1,L}_{n,R}$. This proves positivity.

It remains to establish the derivative estimate including the sum over
all components. Partition the initial section by its finite regular
collision itinerary through $L$, all its section decisions, and the
value $m=N_{n,R}\le L$. There are finitely many pieces, but no estimate
of their number will be used. On one such piece $D$, every mark is
$T_R^{j_i}x$ for a fixed $j_i\le m$, the lattice label
$\ell=(\sum_{j<m}\kappa_R\circ T_R^j,m)$ is constant, and the time
coordinate is $L_m(x)$. The density on the initial cylinder is
\[
 b_D=(4\pi c_*)^{-1}\one_D\Gamma_{L,R,\varepsilon}w_\varepsilon.
\]
Lemma~\ref{lem:source-partition-zero-extension} supplies the formal
partition and the earliest-forbidden-time argument. Its zero extension
across the boundaries of $D$ is $C^2$: all changes
of itinerary or section decision have already been removed on both
sides by \eqref{eq:whole-itinerary-cutoff}. Fixed-order chain and product
rules, smoothing in $C^3$, and Lemma~\ref{lem:whole-itinerary-cutoff}
give on $D$
\begin{equation}\label{eq:geometric-divergence-bound}
 |b_D|+|\operatorname{div}(b_DX_m)|+
 |\operatorname{div}(X_m\operatorname{div}(b_DX_m))|
       \le C_qM\exp\{C_q(L+1)\log(C_q/\varepsilon)\}.
\end{equation}
All vector-field products are extended by zero outside their regular
pieces; the orbit cutoffs vanish on neighborhoods of every excluded
boundary. A further smooth cutoff in $p$ can remove the factor $1/p$
away from the supports. No extension of the billiard map through a
singularity is used.
Angular integration is periodic, and all source functions vanish near
$p=\pm1$. Hence there are no coordinate boundary terms.

For a test function $\phi\in C_c^\infty(\R)$,
$X_m(\phi\circ L_m)=\phi'\circ L_m$. Integration by parts therefore
identifies the first and second distributional derivatives of the
pushforward density as the pushforwards of
\begin{equation}\label{eq:geometric-density-derivatives}
 \operatorname{div}(b_DX_m),\qquad
 \operatorname{div}(X_m\operatorname{div}(b_DX_m)),
\end{equation}
respectively. The signs are positive in both distributional identities.
On these supports $L_m$ is a submersion by
\eqref{eq:uniform-roof-gradient}, so coarea gives $L^1$ densities for
all three pushforwards. Their norms are at most the integrals of the
absolute source functions.

The interiors of the pieces $D$ are disjoint. Summing these source
integrals, even after summing over all output labels, pays only the
fixed area of the collision cylinder. It does not multiply
\eqref{eq:geometric-divergence-bound} by an itinerary count. This proves
\eqref{eq:geometric-derivative-budget} and the $W^{2,1}$ assertion.
The actual finite-count measure has a density for bounded measurable
weights by coarea on its regular branches; alternatively
\eqref{eq:uniform-roof-gradient} excludes a zero gradient outside the
null initial line $p=0$. Subtraction gives the density of $E$.
All time supports lie in $[0,L\tau_{\max}]$, and all lattice supports
are finite.
\end{proof}

\begin{remark}[What is extracted]\label{rem:geometric-extraction-scope}
The measure $E$ includes the discarded parts of every regular critical
word, every grazing and first-hit transition, every section-decision
boundary, and all their iterated preimages. Any singular output image
created by those parts remains in $E$. There is no unlisted branch in
\eqref{eq:geometric-exact-decomposition}. This extraction is not the
power--logarithm jet extraction of Section~\ref{sec:finite-count-extraction};
its new derivative bound does not bound the old jet residual. In
particular, small total variation of $E$ does not imply a small
essential supremum of $e_\varepsilon$.
\end{remark}

\section{Finite-band inversion at fixed lattice labels}
\label{sec:microscopic-finite-band}
The geometric extraction has two uses at the original microscopic
resolution. First, it supplies a uniform far-roof estimate for a new
exact residual at a linear collision cutoff. Second, it reduces exact
fixed-label interval probabilities to a finite Fourier integral of the
\emph{full actual law}, with an arbitrarily small polynomial error.
Neither conclusion evaluates that finite integral asymptotically.

Let $\mu_{n,R}^w=(J_{n,R})_*(w\nu_R^*)$ and retain the weights
\eqref{eq:geometric-weight-budgets}. The uncentered transform is
\[
 \widehat\mu_{n,R}^w(u,b)=\int e^{i(u\cdot\ell+bt)}\dd\mu_{n,R}^w(\ell,t),
 \qquad (u,b)\in[-\pi,\pi]^3\times\R.
\]
Thus $u_3$ is the collision-count frequency and $b$ is roof frequency.
For a time interval $I=[s,s+h]$, $h\ge0$, put
\begin{equation}\label{eq:microscopic-interval-transform}
 H_I(b)=\int_Ie^{-ibt}\dd t,
 \qquad H_I(0)=h,
 \qquad |H_I(b)|\le\min\{h,2/|b|\}.
\end{equation}
All endpoint conventions give the same probability, because the roof
components of the law are absolutely continuous.

\begin{proposition}[Far-roof bounds for the geometric residual]
\label{prop:geometric-far-roof}
For the exact residual of Theorem~\ref{thm:geometric-extraction-budget},
write $\mathcal A=\mathcal A_{2,\varepsilon}(n,L,R,w)$. For $B>0$,
\begin{align}
 \frac1{(2\pi)^4}\int_{[-\pi,\pi]^3}\int_{|b|>B}
                         |\widehat Q_\varepsilon(u,b)|\dd b\dd u
       &\le\frac{\mathcal A}{\pi B},
                 \label{eq:geometric-far-density}\\
 \frac1{(2\pi)^4}\int_{[-\pi,\pi]^3}\int_{|b|>B}
                |H_I(b)\widehat Q_\varepsilon(u,b)|\dd b\dd u
       &\le\frac{\mathcal A}{\pi B^2}.
                 \label{eq:geometric-far-window}
\end{align}
These estimates include the sum over all lattice labels in $\mathcal A$.
\end{proposition}
\begin{proof}
Distributional integration by parts twice on each compactly supported
$W^{2,1}$ component gives
\[
 |\widehat Q_\varepsilon(u,b)|\le\mathcal A/|b|^2\qquad(b\ne0).
\]
For $|b|\le1$ its transform is bounded by $M$, and for $|b|>1$
it is bounded by $\mathcal A/|b|^2$. Thus the transform is integrable
on the full product of the frequency torus and roof line, and absolute
inversion is legitimate. The torus volume is $(2\pi)^3$. Integrating the displayed
bound over both roof tails gives \eqref{eq:geometric-far-density}.
Multiply it by $2/|b|$ before integrating to obtain
\eqref{eq:geometric-far-window}. The angular Fourier convention is
responsible for the factors $\pi^{-1}$.
\end{proof}

Define the finite integral of the actual law by
\begin{equation}\label{eq:actual-microscopic-finite-integral}
 \mathcal I_{n,R,B}^{w}(\ell,I)
 =\frac1{(2\pi)^4}\int_{[-\pi,\pi]^3}\int_{-B}^{B}
    e^{-iu\cdot\ell}H_I(b)\widehat\mu_{n,R}^w(u,b)\dd b\dd u.
\end{equation}
It is well defined for every finite $B$. There is no cutoff or
conditioning inside $\mu_{n,R}^w$ in this definition.

\begin{theorem}[Microscopic finite-band reduction]
\label{thm:microscopic-finite-band-reduction}
There are $a,c>0$, as in Theorem~\ref{thm:cumulative-return-tail}, such
that, for $n\le L$, $0<\varepsilon\le\varepsilon_0$, $B\ge1$, every
$\ell\in\Z^3$ and every finite interval $I$ of length $h$,
\begin{align}
 &\left|\mu_{n,R}^w(\{\ell\}\times I)
                    -\mathcal I_{n,R,B}^{w}(\ell,I)\right|\notag\\
 &\quad\le C_q\left\{M e^{an-cL}+\varepsilon V
                    +M(L+1)\varepsilon^{1/16}\right\}
                         [1+\log(1+Bh)]\notag\\
 &\qquad+\frac{C_qM\exp\{C_q(L+1)\log(C_q/\varepsilon)\}}{\pi B^2}.
                 \label{eq:microscopic-finite-band-error}
\end{align}
In particular no lattice coordinate is averaged, and the time interval
need not grow with $n$.
\end{theorem}
\begin{proof}
Let $Q=Q_{n,R,\varepsilon}^{w,L}$. The cumulative-return tail and
\eqref{eq:geometric-extraction-mass} give
\begin{equation}\label{eq:full-law-geometric-TV}
 \|\mu_{n,R}^w-Q\|_{\TV}
 \le \delta:=C_q\{M e^{an-cL}+\varepsilon V
                         +M(L+1)\varepsilon^{1/16}\}.
\end{equation}
The direct change of the interval probability costs at most $\delta$.
The corresponding change inside the finite integral is at most
\[
 \frac{\delta}{2\pi}\int_{-B}^B\min\{h,2/|b|\}\dd b
       \le C\delta[1+\log(1+Bh)].
\]
This logarithmic cost is essential. Bounding $H_I$ by its supremum on
the whole band would pay a factor of $B$ and lose the conclusion below.
The remaining difference, for $Q$, is its absolutely invertible roof
tail, bounded by \eqref{eq:geometric-far-window}. Add the three errors
and use \eqref{eq:geometric-derivative-budget}.
\end{proof}

\begin{corollary}[A uniform long-time budget]
\label{cor:linear-cutoff-microscopic-inversion}
Fix $q$, $P>0$, $H<\infty$, and nonnegative exponents $\beta,\gamma$.
Suppose
\[
 M\le M_0 n^\beta,\qquad V\le V_0 n^\gamma.
\]
There are a slope $\lambda>a/c$, a fixed $d>0$, and explicit choices
$L_n=\lceil\lambda n\rceil$, $\varepsilon_n=\varepsilon_0 n^{-d}$,
and $B_n\ge1$ with
\begin{equation}\label{eq:long-time-bandwidth-budget}
 \log B_n\le C n\log(2+n)
\end{equation}
such that, uniformly in $R$, the marks, all $\ell\in\Z^3$ and all
intervals of length at most $H$,
\begin{equation}\label{eq:microscopic-uniform-inversion}
 n^2\left|\mu_{n,R}^w(\{\ell\}\times I)
                    -\mathcal I_{n,R,B_n}^{w}(\ell,I)\right|
       \le C_P n^{-P}.
\end{equation}
For the same choices the total variation in
\eqref{eq:full-law-geometric-TV} is $O(n^{-P-4})$, and the residual
has a second-derivative bound $\mathcal A_{2,\varepsilon_n}\le A_n$
with $\log A_n\le C n\log(2+n)$. Its pointwise far-roof inverse,
multiplied by $n^2$, is at most $C_Pn^{-P-2}$.
\end{corollary}
\begin{proof}
Choose $\lambda>\max\{1,a/c\}+1$ and
\[
 d>\max\{16(P+5+\beta),\ P+4+\gamma\}.
\]
The two geometric errors in \eqref{eq:full-law-geometric-TV} are then
$O(n^{-P-4})$. The truncated high-count error is exponentially small
and hence satisfies the same bound, with a changed constant. Set
\begin{equation}\label{eq:explicit-long-time-bandwidth}
 \begin{split}
 A_n&=1+C_qM_0 n^\beta
           \exp\{C_q(L_n+1)\log(C_q/\varepsilon_n)\},\\
 B_n&=n^{P+4}A_n.
 \end{split}
\end{equation}
These are numerical expressions in the constants already proved; no
unknown value of a density derivative occurs in the choice of $B_n$.
Both logarithms satisfy \eqref{eq:long-time-bandwidth-budget}.
For $h\le H$, the logarithmic factor in
\eqref{eq:microscopic-finite-band-error} is at most $C n\log(2+n)$.
Its contribution after multiplication by $n^2$ is
$O(n^{-P-1}\log(2+n))$, bounded by $C_Pn^{-P}$. The last term is
smaller. Finally \eqref{eq:geometric-far-density} and
$\mathcal A_{2,\varepsilon_n}\le A_n$ give the claimed pointwise tail.
\end{proof}

\begin{remark}[Microscopic resolution and the remaining finite integral]
\label{rem:microscopic-reduction-not-LLT}
The three lattice coordinates in \eqref{eq:microscopic-uniform-inversion}
are singletons. Its time interval may have any fixed positive length,
or may shrink; the stated error is absolute, not relative to an
unproved denominator. Both the observable and its law are the original
ones. Thus the reduction is not a mesoscopic replacement of the raw
event. The price is a roof bandwidth whose logarithm can grow as
$n\log n$. The preceding collision and prescribed-return central
bands do not reach this bandwidth. No estimate of the finite integral
\eqref{eq:actual-microscopic-finite-integral} on its remaining complement
is supplied by this corollary.
\end{remark}

\subsection{The common pointwise correction after the new extraction}
The new extraction also makes one precise alternative to the old jet
budget available in the pointwise inversion identity. Let $\chi$ be one
of the retained smooth cutoffs and $K_\chi$ its mixed inverse kernel.
For the same weight and finite count packet, let
\[
 \mathcal R_\chi^{w,L}=p_{n,R}^{w,L}-K_\chi*\mu_{n,R}^{w,L}.
\]
This is the intrinsic signed correction in
\eqref{eq:coherent-raw-correction}; the definition does not
depend on the chosen extraction. Substituting
\eqref{eq:geometric-exact-decomposition} gives the almost-everywhere
identity
\begin{equation}\label{eq:geometric-common-raw-ledger}
 \begin{split}
 \mathcal R_\chi^{w,L}
 ={}& e_\varepsilon-K_\chi*E_\varepsilon\\
 &+\frac1{(2\pi)^4}\int_{[-\pi,\pi]^3}\int_{-B}^B
 e^{-i(u\cdot\ell+bt)}(1-\chi(u,b))
                               \widehat Q_\varepsilon(u,b)\dd b\dd u
   +\mathcal T_{\varepsilon,B}(\ell,t),
 \end{split}
\end{equation}
provided $B$ contains the roof support of $\chi$. Here
\begin{equation}\label{eq:geometric-ledger-controlled-terms}
 \|K_\chi*E_\varepsilon\|_\infty
      \le\|K_\chi\|_\infty\|E_\varepsilon\|_{\TV},
 \qquad
 \|\mathcal T_{\varepsilon,B}\|_\infty
      \le\frac{\mathcal A_{2,\varepsilon}}{\pi B}.
\end{equation}
The identities are formed before taking any essential supremum.
They do not subtract extended-real norm values.

For clarity, the terms in the final raw formula have the following
status in this extraction.
\begin{center}
\begin{tabular}{p{0.29\textwidth}p{0.63\textwidth}}
\hline
Term & Estimate and scope\\
\hline
Known central convolution & The inherited central and fixed-count bands
are unchanged, with their original weight classes.\\
$K_\chi*E_\varepsilon$ & The total-variation estimate above and the known
kernel supremum give a quantitative pointwise bound.\\
$\mathcal T_{\varepsilon,B}$ & The explicit derivative budget gives a
pointwise far-roof bound, including at linear $L$.\\
$e_\varepsilon$ & Its $L^1$ norm is small; no small essential supremum
is proved. It is positive for the original unweighted law.\\
Finite complementary integral & Balanced, compact nonzero, peripheral
and growing-roof regimes outside the proved domains still require a
compatible bound.\\
\hline
\end{tabular}
\end{center}

In particular, the geometrically removed density can be tall on a very
small set. The elementary estimate
\[
 m\{|e_\varepsilon|>\eta\}
        \le \|E_\varepsilon\|_{\TV}/\eta,\qquad \eta>0,
\]
uses the mixed counting--Lebesgue measure $m$, and cannot be upgraded to
an essential-supremum estimate. The two uncontrolled terms in
\eqref{eq:geometric-common-raw-ledger} may still need to be estimated
together with their sign. The coherent transport identities remain
available for that purpose.

For a single marked BV weight the geometric extraction and the inherited
Fourier estimates now apply to the same law without a further
subanalytic-weight assumption for the extraction. For a fixed product
of several return-state marks, the new extraction and
\eqref{eq:microscopic-uniform-inversion} hold, but the inherited
single-mark central theorem is not thereby a multiple-mark theorem.
The stationary path cylinders of Theorem U and the original microscopic
physical-time event are different weight classes. Their raw denominator
and relative microscopic event replacement are not consequences of
\eqref{eq:microscopic-uniform-inversion}.

\section{Source boundaries and the stationary geometric extraction}
\label{sec:source-boundary-calculus}
We give the coordinate and boundary arguments for the geometric
extraction in one place. We then establish the corresponding positive
extraction under the stationary suspension law. The latter is not the
section law: its entrance marginal has an unbounded roof bias.

\subsection{The differential and the action in invariant coordinates}
\begin{lemma}[Collision differential and action]
\label{lem:self-contained-collision-differential}
On a regular flight from $Rn_\alpha$ to $d+Rn_{\alpha'}$, set
$c=\sqrt{1-p^2}$, $c'=\sqrt{1-(p')^2}$ and $v_0=\tau_R/R$.
In the coordinates $(\alpha,p)$ the derivative is
\[
 DT_R=-\begin{pmatrix}
 (v_0+c)/c'&v_0/(cc')\\
 v_0+c+c'&(v_0+c')/c
 \end{pmatrix},\qquad \det DT_R=1,
\]
and
\begin{equation}\label{eq:self-contained-action-differential}
 d\tau_R=Rp'\,d\alpha'-Rp\,d\alpha.
\end{equation}
These formulas use the outgoing value of $p'$ after reflection.
\end{lemma}
\begin{proof}
Write $\mathcal J$ for counterclockwise rotation by $\pi/2$, and
$V=cn_\alpha+pt_\alpha$ for the unit flight velocity. If $\beta$ is
its direction angle on a local chart, then
\[
 d\beta=d\alpha+dp/c,\qquad dV=\mathcal JV\,d\beta.
\]
At arrival, $V=-c'n_{\alpha'}+p't_{\alpha'}$: reflection preserves
the tangential component and reverses the normal component. Thus
$d\beta=d\alpha'-dp'/c'$. Differentiating the flight equation and
taking its scalar product with $\mathcal JV$ gives
\[
 -Rc'\,d\alpha'=Rc\,d\alpha+\tau_R\,d\beta,
\]
since $\mathcal JV\cdot t_\alpha=c$ and
$\mathcal JV\cdot t_{\alpha'}=-c'$. Substitution gives the first
row of the displayed matrix. The identity
$dp'=c'(d\alpha'-d\beta)$ gives its second row. Its determinant is
\[
 \frac{(v_0+c)(v_0+c')-v_0(v_0+c+c')}{cc'}=1.
\]
Taking instead the scalar product of the differentiated flight equation
with $V$ gives \eqref{eq:self-contained-action-differential}.
The lattice center is constant on this branch. Summing the action
formula gives the two derivatives of $L_m$ used in
Lemma~\ref{lem:roof-transverse-direction}. This derivation also
specifies the sign and coordinate conversion in the fixed-table
specialization of \cite[Section 3.1]{SYZ}.
\end{proof}

\subsection{A complete first-hit guard}
Let $\tau_{\max}$ be a uniform horizon from the earlier geometric
estimates, and choose once and for all
\begin{equation}\label{eq:explicit-finite-center-set}
 \mathcal D=\{d\in\Lambda\setminus\{0\}:
                        |d|\le\tau_{\max}+2R_++1\}.
\end{equation}
Here $\Lambda$ is the physical triangular lattice, not a rescaled
collision coordinate lattice.

\begin{lemma}[First-hit transitions and uniform chart coefficients]
\label{lem:complete-first-hit-guard}
The set \eqref{eq:explicit-finite-center-set} is finite and contains
every next collision center for all $R\in I$. Away from initial grazing
and the zero sets of all $\Delta_{d,R}$, the first-hit label is locally
constant and the selected positive entry root is smooth. In four bounded
angular half-angle charts, the numerator of each discriminant has degree
at most eight, its positive denominator lies in $[1,16]$, and the maximum
modulus of its coefficients is bounded below uniformly in $R$, the
chart and $d\in\mathcal D$.
\end{lemma}
\begin{proof}
A flight of length at most $\tau_{\max}$ from one boundary point to the
next gives $|d|\le\tau_{\max}+2R_+$, proving the first assertion.
For a candidate let $a=d-Rn_\alpha$, $A=a\cdot V$ and
$\Delta=R^2-|a|^2+A^2$. Since $|d|\ge1$ and $R\le47/100$,
\[
 |a|\ge1-R>R,\qquad A^2-\Delta=|a|^2-R^2>0.
\]
On $\Delta>0$ the entry root is $A-\sqrt\Delta$. It is positive
exactly when $A>0$. Its sign cannot change while $\Delta>0$, and a
positive root cannot pass through zero. Two distinct positive entry
roots cannot be equal: the corresponding point on the ray would lie
on two disjoint closed disks. Consequently the ordering of all admissible
entry roots persists locally until a candidate discriminant vanishes.
The margin in \eqref{eq:explicit-finite-center-set} prevents a center
from entering through the edge of the candidate list. The outgoing
half-ray cannot reenter the initial convex disk. This accounts for all
first-hit transitions, including the side of a tangency on which the
candidate is no longer hit.

Take chart centers $\alpha_j=j\pi/2$, $0\le j\le3$, and use
$\alpha=\alpha_j+2\arctan t$ on $|t|<1$; their open images cover the
angular circle. Bounds may be taken on the closed chart intervals.
Let $s=\tan(\varphi/2)\in[-1,1]$, where $p=\sin\varphi$.
Then
\[
 p=\frac{2s}{1+s^2},\qquad c=\frac{1-s^2}{1+s^2},\qquad
 a\mathbin{\times}V
  =c(d\mathbin{\times}n_\alpha)
                  +p(d\mathbin{\times}t_\alpha-R).
\]
The last expression has denominator $(1+t^2)(1+s^2)$ and a
numerator of total degree at most four. Thus the discriminant has
denominator $(1+t^2)^2(1+s^2)^2\in[1,16]$ and a numerator of total
degree at most eight. The coordinate Jacobian of $\nu$ is bounded
above by $1/\pi$, since $|d\alpha/dt|\le2$ and $|dp/ds|\le2$.

For each fixed $d,j,R$, the numerator is not identically zero. Fix an
initial point in the chart. Its vector $a$ has length greater than $R$;
the square of its cross product with directions in an open angular
interval cannot be the constant $R^2$. The maximum modulus of the
finitely many polynomial coefficients is a continuous, strictly positive
function of $R$. Its minimum on $I$ is positive. Taking the minimum
over the finite candidate and chart sets proves the uniform coefficient
bound. This justifies the uniform use of
Lemma~\ref{lem:elementary-polynomial-sublevel}, including near tangency.
\end{proof}

\begin{lemma}[Partition and zero extension]
\label{lem:source-partition-zero-extension}
Fix $n\le L$ and $\varepsilon>0$. Up to a $\nu$-null set, the source
$Y_R^*\cap\{N_{n,R}\le L\}$ is the union of finitely many disjoint open
pieces $D$ relative to the collision cylinder. On each piece the center
itinerary, section decisions, mark collision indices, output label and
$m=N_{n,R}$ are fixed. The source functions $b_D$ in the proof of
Theorem~\ref{thm:geometric-extraction-budget}, and the vector-field
products used there, extend by zero with the regularity needed for both
integrations by parts. In particular there are no distributions
supported on itinerary or section boundaries in those identities.
\end{lemma}
\begin{proof}
At each collision record the first-hit center in the finite set
\eqref{eq:explicit-finite-center-set} and the section membership decision.
There are finitely many such records through $L$ collisions. Remove
all encountered tangencies, grazing points and section boundaries.
Their union is null: the one-step discriminant zero sets are null by
Lemma~\ref{lem:complete-first-hit-guard} and the polynomial sublevel
bound, section sides are null, and the collision map preserves $\nu$.
The regular decision sets are open. They need not be connected for the
argument; all their components with one record are treated together.
On each such set every return index and output label is fixed. No
artificial angular seam is introduced in this partition.

Here is the boundary extension explicitly. At a boundary point choose
the first singular collision or the first section decision at which the
record can change. All preceding regular branches persist on a
neighborhood. At a first-hit tangency the input discriminant of the
preceding flight is zero. At an initial grazing point the initial guard
is zero; grazing produced by a flight is already its tangency. At a
section decision the corresponding rectangle-side guard is zero.
In each case, at the earliest such input time $j$, the previous maps
are continuous and $\chi_{R,\varepsilon}\circ T_R^j$ is identically
zero on a two-sided neighborhood of the boundary point. The factors
at later, possibly undefined, collision states need not be evaluated:
the product is defined to be zero throughout that neighborhood.

At a boundary of a cutoff transition inside a regular decision set,
ordinary smoothness and the flat endpoint derivatives of $\theta$
give the same matching of derivatives. The initial factor
$\theta(p/\varepsilon)$ removes the pole of $X_m$. Hence $b_D$ is
$C^2$ after zero extension, $b_DX_m$ is $C^2$, and
$X_m\operatorname{div}(b_DX_m)$ is $C^1$. Their derivatives have the
bounds of \eqref{eq:geometric-divergence-bound}. For fixed $L$ and
$\varepsilon$ their supports are separated from the excluded boundaries;
no uniform lower bound for that separation is used. The cylinder is
periodic in $\alpha$, and the supports avoid $p=\pm1$.
The distributional integration-by-parts formulas therefore have no
boundary terms. Since the pieces have disjoint interiors, their
absolute source integrals sum over a subset of the initial cylinder,
not over multiple copies of it.
\end{proof}

\subsection{The preceding section under the equilibrium law}
Write $r=r_R^*$ and $\tau^*=\tau_R^*$. On the return suspension put
\begin{equation}\label{eq:equilibrium-source-for-inversion}
 \Omega_R=\{(y,s):y\in Y_R^*,\ 0\le s<\tau^*(y)\},\qquad
 d\mathbb P_R^{\rm eq}=\frac{c_*}{\bar\tau_R}\,d\nu_R^*(y)\,ds.
\end{equation}
Its entrance marginal is $\rho_R\nu_R^*$, where
$\rho_R=c_*\tau^*/\bar\tau_R$. This is the stationary normalization
of Lemma~\ref{lem:stationary-length-bias}. The same $y$ is used for
$J_{n,R}(y)$ and for every physical trajectory observable on $\Omega_R$.

\begin{proposition}[Positive extraction with the actual entrance bias]
\label{prop:stationary-geometric-extraction}
Let $\mu_{n,R}^{\rm eq}=(J_{n,R})_*(\rho_R\nu_R^*)$.
There is a positive measure $Q\le\mu_{n,R}^{\rm eq}$ with a
compactly supported mixed density $q$ whose roof components belong to
$W^{2,1}(\R)$, and constants independent of $R,n,L,\varepsilon$ such that
\begin{align}
 \delta_{\rm eq}:=\|\mu_{n,R}^{\rm eq}-Q\|_{\TV}
 &\le C\left(e^{(an-cL)/2}
                +(L+1)^{1/2}\varepsilon^{1/32}\right),
                     \label{eq:stationary-geometric-mass}\\
 \sum_{\ell\in\Z^3}\|\partial_t^2q(\ell,\cdot)\|_1
 &\le C\exp\{C(L+1)\log(C/\varepsilon)\}=:A_2.
                     \label{eq:stationary-geometric-A2}
\end{align}
In particular $\|Q\|_{\TV}\le1$. For the section probability the same
conclusions hold with the stronger removed-mass estimate
$C(e^{an-cL}+(L+1)\varepsilon^{1/16})$.
\end{proposition}
\begin{proof}
Define
\[
 Q=(J_{n,R})_*\bigl(\rho_R\one_{\{N_{n,R}\le L\}}
                         \Gamma_{L,R,\varepsilon}\nu_R^*\bigr).
\]
All factors are nonnegative and the last two are at most one. The
one-return exponential moment and the positive lower bound for
$\bar\tau_R$ give $\sup_R\|\rho_R\|_{L^2(\nu_R^*)}<\infty$.
Cauchy--Schwarz, the cumulative return tail and
\eqref{eq:whole-itinerary-mass} consequently give
\[
 \int\rho_R\one_{\{N_n>L\}}\,d\nu_R^*
       \le C e^{(an-cL)/2},\qquad
 \int\rho_R(1-\Gamma_{L,R,\varepsilon})\,d\nu_R^*
       \le C\bigl((L+1)\varepsilon^{1/16}\bigr)^{1/2}.
\]
These imply \eqref{eq:stationary-geometric-mass}.

On a source piece of Lemma~\ref{lem:source-partition-zero-extension},
its first return collision time $j_1$ is fixed and at most $L$.
There $\tau^*=\sum_{j<j_1}\tau_R\circ T_R^j$.
The cutoff estimates and chain rule bound this function and its first
three derivatives, on the required supports, by
$C\exp\{C(L+1)\log(C/\varepsilon)\}$. Multiplying $b_D$ by
$\rho_R$ therefore preserves the divergence budget after increasing
$C$. The earliest-decision argument gives the same zero extensions.
The two source integrations by parts and disjoint-source summation
prove \eqref{eq:stationary-geometric-A2}. The argument does not assume
that the untruncated entrance bias is BV. The section-law assertion
is the unweighted case of Theorem~\ref{thm:geometric-extraction-budget}
and its cumulative-tail extension.
\end{proof}

\section{Positive finite-band inversion and microscopic conditioning}
\label{sec:positive-microscopic-conditioning}
A sharp truncation of a Fourier integral need not define a positive
measure. For conditioning it is useful to retain positivity before
normalization. We construct a finite-band likelihood on the original
probability space and control it uniformly against every bounded
trajectory statistic. Only the unweighted law is regularized; the
trajectory statistic is never differentiated or smoothed.

\subsection{A positive kernel and a source-to-density estimate}
Use the real-line transform $\widehat f(b)=\int e^{ibt}f(t)\,dt$.
Define, with the continuous value at zero,
\begin{equation}\label{eq:positive-bandlimited-kernel}
 k(t)=\frac3{8\pi}\left(\frac{\sin(t/4)}{t/4}\right)^4,
 \qquad k_B(t)=B k(Bt),\qquad B>0.
\end{equation}
\begin{lemma}[Kernel normalization and moments]
\label{lem:positive-kernel-moments}
The kernel $k$ is even and nonnegative, $\int k=1$, and
\begin{equation}\label{eq:positive-kernel-transform}
 \widehat k(b)=\eta(b)=
 \begin{cases}
 1-6b^2+6|b|^3,& |b|\le1/2,\\
 2(1-|b|)^3,&1/2\le|b|\le1,\\
 0,&|b|\ge1.
 \end{cases}
\end{equation}
Moreover $\int t k(t)\,dt=0$, $\int t^2k(t)\,dt=12$, and
$\kappa_1:=\int|t|k(t)\,dt\le\sqrt{12}$.
\end{lemma}
\begin{proof}
The identity
$\sin(t/4)/(t/4)=2\int_{-1/4}^{1/4}e^{iut}\,du$
expresses the fourth power by the fourfold convolution of the interval
indicator. Evaluating that elementary convolution gives
\eqref{eq:positive-kernel-transform}, including the normalization
$\eta(0)=1$. The kernel decays as $O((1+|t|)^{-4})$ and hence has a
finite second moment. Differentiating its transform twice at zero
therefore gives $\int t^2k(t)\,dt=-\eta''(0)=12$.
Evenness gives the first signed moment; Cauchy--Schwarz gives the
absolute first-moment bound.
\end{proof}

\begin{lemma}[An all-label interpolation inequality]
\label{lem:all-label-interpolation}
Let $q_\ell\in W^{2,1}(\R)$ have compact support, and suppose
$A_j=\sum_\ell\|q_\ell^{(j)}\|_1<\infty$ for $j=0,2$.
Then
\begin{equation}\label{eq:all-label-interpolation}
 A_1\le2\sqrt{A_0A_2},\qquad
 \sum_\ell\|q_\ell\|_\infty\le A_1/2,
 \qquad
 \sum_\ell\|q_\ell-k_B*q_\ell\|_1\le6A_2/B^2.
\end{equation}
The same statements hold for real or complex functions.
\end{lemma}
\begin{proof}
Taylor's formula in $W^{2,1}$, summed over labels, gives for every $h>0$
\[
 A_1\le 2A_0/h+hA_2/2.
\]
Taking $h=2\sqrt{A_0/A_2}$ proves the first inequality when both
quantities are positive. If either vanishes the claim follows by
letting $h$ tend to zero or infinity. Each component tends to zero at
both ends of the line. Integrating its derivative from either end
bounds its supremum by half its total variation, proving the second
inequality. Finally
\[
 \|q_\ell(\cdot-y)-q_\ell+yq'_\ell\|_1
                       \le y^2\|q''_\ell\|_1/2.
\]
Integrate against $k_B(y)$, use its zero signed first moment and second
moment $12/B^2$, and sum. All sums and integrals are justified by the
nonnegative bounds just displayed.
\end{proof}

Here and below $\|\cdot\|_{\TV}$ is the total variation norm of a signed
measure, equivalently the supremum of its integrals against measurable
complex functions of modulus at most one. Thus it is the $L^1$ norm
for densities, without a factor of one half.

\begin{theorem}[Positive spectral approximation of a mixed law]
\label{thm:positive-spectral-approximation}
Let $Z=(\mathsf L,\mathsf T)$ on a probability space
$(\Omega,\mathbb P)$ take values in $\Z^3\times\R$. Suppose its law
$\mu$ admits a positive submeasure $Q\le\mu$ with compactly supported
roof densities $q_\ell\in W^{2,1}(\R)$. Write
\[
 \delta=\|\mu-Q\|_{\TV},\qquad
 A_2=\sum_\ell\|q''_\ell\|_1<\infty.
\]
Then
\begin{equation}\label{eq:positive-full-law-density}
 p_B(\ell,t)=\frac1{(2\pi)^4}
  \int_{[-\pi,\pi]^3}\int_{-B}^B
 e^{-i(u\cdot\ell+bt)}\eta(b/B)\widehat\mu(u,b)\,db\,du
\end{equation}
is a nonnegative density of total mass one. It smooths only the roof
coordinate and preserves all three lattice coordinates exactly. The
corresponding probability $\mu_B$ satisfies
\begin{equation}\label{eq:positive-full-law-TV}
 \|\mu-\mu_B\|_{\TV}\le2\delta+6A_2/B^2.
\end{equation}
If $\mu$ has density $p$, this is the global mixed $L^1$ estimate for
$p-p_B$, including the sum over all labels.
\end{theorem}
\begin{proof}
Fourier inversion for $k_B$ and torus orthogonality give
\[
 p_B(\ell,t)
   =\mathbb E\bigl[\one_{\{\mathsf L=\ell\}}
                                      k_B(t-\mathsf T)\bigr].
\]
The frequency integral is absolutely convergent, since the band is
finite and $|\widehat\mu|\le1$. This identity proves positivity,
normalization and exact preservation of the lattice marginal.
Write $\mu=Q+E$ with $E\ge0$ and $\|E\|_{\TV}=\delta$.
Roof convolution is a contraction on finite measures, so
$\|E-k_B*E\|_{\TV}\le2\delta$. The last estimate of
Lemma~\ref{lem:all-label-interpolation} bounds the remaining difference
$Q-k_B*Q$, proving \eqref{eq:positive-full-law-TV}.
\end{proof}

\subsection{A microscopic event as a finite-band likelihood}
Consider an event $A=\{Z\in\mathcal A\}$ with fibres
$\mathcal A_\ell\subset\R$, each a union of at most $J$ bounded
intervals. Empty fibres are permitted. Endpoint conventions do not
matter for the billiard laws considered below. Define on the original
probability space
\begin{equation}\label{eq:positive-event-likelihood}
 \lambda_{A,B}
       =(k_B*\one_{\mathcal A_{\mathsf L}})(\mathsf T).
\end{equation}
In particular $0\le\lambda_{A,B}\le1$. This is a weight on the same
initial trajectory, not a newly sampled section point or a change to
its physical evolution.

\begin{theorem}[Uniform weighted inversion on the original space]
\label{thm:all-selector-finite-band}
Under the hypotheses of Theorem~\ref{thm:positive-spectral-approximation},
let
\begin{equation}\label{eq:positive-event-error-budget}
 E_B(J)=\delta+2\kappa_1J\sqrt{A_2}/B.
\end{equation}
Then
\begin{equation}\label{eq:positive-likelihood-L1}
 \|\one_A-\lambda_{A,B}\|_{L^1(\mathbb P)}\le E_B(J).
\end{equation}
For every bounded measurable trajectory statistic $W$, including a
statistic depending on $n$ or $B$,
\begin{equation}\label{eq:all-selector-error}
 \left|\mathbb E[W\one_A]-\mathbb E[W\lambda_{A,B}]\right|
                               \le\|W\|_\infty E_B(J).
\end{equation}
No BV, continuity or finite-record hypothesis on $W$ is required.
For $\mathcal A=\{\ell\}\times I$, the second expectation is exactly
\begin{equation}\label{eq:all-selector-Fourier-formula}
 \frac1{(2\pi)^4}\int_{[-\pi,\pi]^3}\int_{-B}^B
 e^{-iu\cdot\ell}H_I(b)\eta(b/B)
          \mathbb E\bigl[W e^{i(u\cdot\mathsf L+b\mathsf T)}\bigr]
                                                        \,db\,du.
\end{equation}
The bound is independent of the length of $I$. Finite unions of fibres
have the corresponding finite sum of these formulas.
\end{theorem}
\begin{proof}
For a set $I$ that is a union of at most $J$ intervals,
\[
 \int_\R|\one_I(t)-\one_I(t-y)|\,dt\le2J|y|.
\]
Positivity and unit mass of $k_B$ imply pointwise
\[
 |\one_I(t)-(k_B*\one_I)(t)|
    \le\int k_B(y)|\one_I(t)-\one_I(t-y)|\,dy.
\]
On the positive discarded measure $E=\mu-Q$ the integrand is at most
one, and so costs at most $\delta$. On $Q$, bound each component by
$\|q_\ell\|_\infty$ and apply the preceding translation estimate.
Summation and Lemma~\ref{lem:all-label-interpolation} give
\[
 \int|\one_{\mathcal A}-k_B*\one_{\mathcal A}|\,dQ
   \le\frac{2\kappa_1J}{B}\sum_\ell\|q_\ell\|_\infty
   \le\frac{\kappa_1J A_1}{B}
   \le\frac{2\kappa_1J\sqrt{A_2}}B,
\]
since $A_0=\|Q\|_{\TV}\le1$. This proves
\eqref{eq:positive-likelihood-L1}. It also handles countably many
nonempty fibres by summation of nonnegative bounds.

Multiplication by $W$ proves \eqref{eq:all-selector-error}; taking the
supremum over $|W|\le1$ gives precisely the total variation norm of
$\one_A\mathbb P-\lambda_{A,B}\mathbb P$.
Fourier inversion for $k_B$, followed by integration over the bounded
interval $I$ and torus orthogonality, proves
\eqref{eq:all-selector-Fourier-formula}. Fubini applies because the
frequency band and interval are finite and $W$ is bounded.
No derivative of the weighted transform is used. In particular the
proof does not claim a global $L^1$ convolution approximation uniform
in arbitrary oscillating weights $W$; it proves the stated interval-event
operator estimate instead.
\end{proof}

\begin{corollary}[Relative conditioning and selected denominators]
\label{cor:positive-microscopic-posterior}
Let $p_A=\mathbb P(A)$ and $p_{A,B}=\mathbb E\lambda_{A,B}$.
If $p_A>E_B(J)$, the probability
$\Pi_{A,B}=p_{A,B}^{-1}\lambda_{A,B}\mathbb P$ is defined and
\begin{equation}\label{eq:positive-posterior-TV}
 \|\mathbb P(\cdot\mid A)-\Pi_{A,B}\|_{\TV}
                            \le 2E_B(J)/p_A.
\end{equation}
Alternatively, $p_{A,B}>E_B(J)$ proves the rigorous denominator bound
$p_A\ge p_{A,B}-E_B(J)>0$ and the upper bound
$2E_B(J)/(p_{A,B}-E_B(J))$ in \eqref{eq:positive-posterior-TV}.

For every measurable selector $D\subset\Omega$, the same statements
hold with $A\cap D$ and the likelihood $\one_D\lambda_{A,B}$,
using the identical error budget $E_B(J)$. More generally they hold
for a nonnegative path weight bounded by one, with its corresponding
weighted denominator. No lower bound for that denominator is assumed
without being stated.
\end{corollary}
\begin{proof}
Equation~\eqref{eq:positive-likelihood-L1} bounds both the variation
of the unnormalized measures and the difference of their masses.
For two nonnegative densities $f,g$ of masses $a,b>0$,
\[
 \|f/a-g/b\|_1\le\|f-g\|_1/a+|a-b|/a.
\]
Apply this with $f=\one_A$ and $g=\lambda_{A,B}$. Positivity of the
needed masses follows from their difference being at most $E_B(J)$.
Multiplication by $\one_D$, or by a nonnegative weight at most one,
can only decrease the unnormalized $L^1$ error. This proves every
assertion with no regularity assumption on the selector.
\end{proof}

\subsection{Uniform long-time budgets for the actual billiard laws}
Let $\mathbb P_R^{\rm sec}=\nu_R^*$ and let $\mathbb P_R^{\rm eq}$
be \eqref{eq:equilibrium-source-for-inversion}. In the first case set
$Z=J_{n,R}(y)$; in the second use the same record from the actual preceding
section origin $y$, retaining the elapsed age $s$ for all path selectors.
Denote the two output densities by $p_{n,R}^{\mathrm j}$,
$\mathrm j\in\{\mathrm{sec},\mathrm{eq}\}$.

\begin{theorem}[Microscopic positive inversion with arbitrary selectors]
\label{thm:uniform-positive-microscopic-inversion}
Fix $P>0$, $J_0<\infty$ and $\kappa\ge0$. There are choices
$L_n=\lceil\lambda n\rceil$, $\varepsilon_n=\varepsilon_0n^{-d}$
and $B_n\ge1$, uniform in $R$, for which
\begin{equation}\label{eq:positive-inversion-bandwidth}
 \log B_n\le C n\log(2+n),
\end{equation}
and the following assertions hold for both initial laws. Let
$p_{n,R,B_n}^{\mathrm j}$ be the positive finite-band density
\eqref{eq:positive-full-law-density}. Then
\begin{align}
 n^2\|p_{n,R}^{\mathrm j}-p_{n,R,B_n}^{\mathrm j}\|_{L^1(\Z^3\times\R)}
        &\le C_Pn^{-P},\label{eq:uniform-positive-density}\\
 n^2\sup_{\substack{\mathcal A:\ J\le J_0n^\kappa\\ |W|\le1}}
 \left|\mathbb E_R^{\mathrm j}[W\one_{\{J_{n,R}\in\mathcal A\}}]
       -\mathbb E_R^{\mathrm j}[W\lambda_{A,B_n}]\right|
        &\le C_Pn^{-P}.\label{eq:uniform-positive-all-selectors}
\end{align}
The supremum is over the interval-fibre events defined above and all
bounded measurable statistics on the respective original probability
space. In particular singleton lattice labels and arbitrarily shrinking
roof intervals are allowed. Their Fourier representation is
\eqref{eq:all-selector-Fourier-formula}, with $\widehat\mu^{W}$ computed
under the same initial law.
\end{theorem}
\begin{proof}
Choose $\lambda>\max\{1,a/c\}+1$ and $d>32(P+5)$.
Proposition~\ref{prop:stationary-geometric-extraction} gives, for both
laws, positive extractions with
\[
 \delta_n\le C_Pn^{-P-4},\qquad
 A_{2,n}\le A_n:=1+C\exp\{C(L_n+1)\log(C/\varepsilon_n)\}.
\]
Indeed the slower stationary geometric mass is bounded by
$C n^{1/2-d/32}$, and its high-count term is exponentially small.
Increase the common $C$ to cover both derivative budgets and all
initial values of $n\ge2$. Set
\begin{equation}\label{eq:positive-explicit-bandwidth}
 B_n=n^{P+\kappa+4}\sqrt{A_n}.
\end{equation}
This proves \eqref{eq:positive-inversion-bandwidth}. Also
\[
 2\delta_n+6A_n/B_n^2\le C_Pn^{-P-4},\qquad
 \delta_n+2\kappa_1J_0n^\kappa\sqrt{A_n}/B_n
                                    \le C_Pn^{-P-4}.
\]
Apply Theorems~\ref{thm:positive-spectral-approximation} and
\ref{thm:all-selector-finite-band}, then multiply by $n^2$.
The stronger intermediate powers absorb constants and rounding.
All bandwidths here are finite; their logarithms, not the bandwidths
themselves, have order $n\log n$.
\end{proof}

\paragraph{Location of the remaining local theorem.}
The positive likelihood gives a finite-band representation of microscopic
conditioning on the original probability space, including arbitrary
bounded path selectors and the true stationary entrance bias.
It does not evaluate that finite integral by a Gaussian. Its relative
conclusions use the denominator bounds stated in
Corollary~\ref{cor:positive-microscopic-posterior}. Likewise
\eqref{eq:uniform-positive-density} is a global $L^1$ approximation,
not an essential-supremum estimate on a prescribed central window.
The full microscopic conclusion remains governed by
\eqref{eq:geometric-common-raw-ledger}: the finite complement and the
pointwise common correction must still be estimated, followed by the
original physical-event comparison at that denominator scale.

\section{Exact stationary endpoint inversion by flight-age integration}
\label{sec:exact-stationary-endpoint}
The return record from the preceding section and the observation at a
fixed physical time are different microscopic events. In this section
we invert the latter directly. Integrating the actual initial flight
age produces an overlap of two intervals. This overlap has two
distributional derivatives with a uniform measure bound, independently
of the number of intervening collisions. The initial and terminal
cell corrections are retained exactly, not bounded and discarded.

Throughout this section $m\ge1$ is a prescribed \emph{physical collision
count}, not a return count. Put
\[
 S_{m,R}=\sum_{j=0}^{m-1}\tau_R\circ T_R^j,\qquad
 K^{\rm c}_{m,R}=\sum_{j=0}^{m-1}\kappa_R\circ T_R^j,
 \qquad 0<\tau_-\le\tau_R\le\tau_+.
\]
Here $\kappa_R$ is the physical next-disk label and $\tau_-,\tau_+$
are uniform finite-horizon bounds. The equilibrium probability in
current-collision coordinates is, by
Lemma~\ref{lem:stationary-length-bias},
\begin{equation}\label{eq:stationary-endpoint-source}
 d\mathbb P_R(x,a)=\bar\tau_R^{-1}d\nu(x)\,da,
                    \qquad 0\le a<\tau_R(x).
\end{equation}
This is the same stationary trajectory as in the return-suspension
coordinates. No new initial point is sampled.

\subsection{Cell labels during a physical flight}
Fix the fundamental-cell convention used for $W_R(t)$. We take a
bounded convex polygonal fundamental cell $\mathcal P$; in particular
one may use a fixed half-open lattice parallelogram. Let
$\mathcal L:\R^2\to\R^2$ send the standard integer basis to the
triangular lattice basis. For the lifted flight beginning at the disk
with center zero define
\begin{equation}\label{eq:within-flight-cell}
 d_R(x,a)=d\quad\Longleftrightarrow\quad
 Rn_\alpha+a v(x)\in\mathcal P+\mathcal Ld,
                \qquad d\in\Z^2,
\end{equation}
with the fixed half-open boundary convention. The same cell convention
is used at the two observation times.

\begin{lemma}[A uniform finite flight partition]
\label{lem:finite-flight-cell-partition}
There is a fixed finite set $\mathcal D\subset\Z^2$ such that, for
almost every regular $x$ and every $R$, the sets
\[
 I_{R,d}(x)=\{a\in[0,\tau_R(x)):d_R(x,a)=d\},\qquad d\in\mathcal D,
\]
are intervals or empty and partition the flight up to endpoints.
Write $J_* = |\mathcal D|$. Both $J_*$ and $\max_{d\in\mathcal D}|d|$
are uniform in $R$.

If $C_R(t)=m$ for a trajectory starting at $(x,a)$, its terminal
collision state is $y=T_R^mx$, its terminal age is
$v=t+a-S_{m,R}(x)$, and exactly
\begin{equation}\label{eq:exact-physical-cell-correction}
 0\le v<\tau_R(y),\qquad
 W_R(t)=K^{\rm c}_{m,R}(x)+d_R(y,v)-d_R(x,a).
\end{equation}
\end{lemma}
\begin{proof}
The segment in \eqref{eq:within-flight-cell} is contained in the fixed
ball of radius $47/100+\tau_+$. Only finitely many translates of the
bounded fundamental cell meet that ball; take their labels as
$\mathcal D$. Convexity makes the intersection of each cell with a
straight flight an interval. Flights lying along a cell edge form a
null set and may be assigned the fixed boundary convention. This
proves the partition and uniformity without differentiating a cell
crossing time.

The $j$th future collision occurs at $S_{j,R}(x)-a$. These times are
strictly increasing, so the terminal age condition is equivalent to
$C_R(t)=m$, apart from suspension-null endpoints. The terminal lifted
position is
$\mathcal L K^{\rm c}_{m,R}(x)+Rn_{\alpha_y}+v v(y)$.
Subtracting the initial cell label from its cell label gives
\eqref{eq:exact-physical-cell-correction}. In particular the last two
terms cannot be omitted at a singleton lattice label.
\end{proof}

For intervals $I,J\subset\R$ put
\begin{equation}\label{eq:flight-overlap-kernel}
 H_{I,J}(s)=\int_\R\one_I(a)\one_J(s+a)\,da
                         =(\one_{-I}*\one_J)(s).
\end{equation}
The value does not depend on choices at interval endpoints. Define
$p_{m,R}(k,t)=\mathbb P_R\{W_R(t)=k,C_R(t)=m\}$ for $t\ge0$ and
extend it by zero to $t<0$.

\begin{theorem}[Exact overlap formula and uniform physical-time regularity]
\label{thm:stationary-overlap-regularity}
For all $m\ge1$ the profiles $p_{m,R}(k,\cdot)$ have continuous,
compactly supported representatives given by
\begin{align}
 p_{m,R}(k,t)
 =\frac1{\bar\tau_R}\int_M\sum_{d,e\in\mathcal D}
 &\one_{\{K^{\rm c}_{m,R}(x)+e-d=k\}}\notag\\[-1mm]
 &\cdot H_{I_{R,d}(x),I_{R,e}(T_R^mx)}(t-S_{m,R}(x))\,d\nu(x).
 \label{eq:exact-overlap-probability}
\end{align}
There are constants $A_0,A_1,A_2,L_0<\infty$, independent of $m,R$,
for which
\begin{align}
 \sum_k\|p_{m,R}(k,\cdot)\|_1&\le A_0,&
 \sum_k\|\partial_t p_{m,R}(k,\cdot)\|_1&\le A_1,
                                      \label{eq:stationary-time-W1}\\
 \sum_k\|D_t^2p_{m,R}(k,\cdot)\|_{\TV}&\le A_2,&
 \sum_k|p_{m,R}(k,t)-p_{m,R}(k,s)|&\le L_0|t-s|.
                                      \label{eq:stationary-time-curvature}
\end{align}
One may take, writing $\bar\tau_- =\min_R\bar\tau_R$,
\[
 A_0=\tau_+^2/\bar\tau_-,\quad
 A_1=2J_*\tau_+/\bar\tau_-,\quad
 A_2=4J_*^2/\bar\tau_-,\quad L_0=2J_*/\bar\tau_-.
\]
Here $D_t^2p$ is a finite signed measure; an $L^1$ density for this
second derivative is neither needed nor asserted. The profiles are
probabilities at fixed observation times, not a probability density
of total mass one on $\Z^2\times\R$.
\end{theorem}
\begin{proof}
Insert \eqref{eq:exact-physical-cell-correction} into
\eqref{eq:stationary-endpoint-source} and integrate the initial age.
This is exactly \eqref{eq:exact-overlap-probability}. Since
$S_m\ge\tau_R(x)$, its integrand vanishes for $t<0$; it is supported
in $0\le t\le(m+1)\tau_+$. There are only finitely many lattice
labels for each $m$, by finite horizon.

For an interval, $D\one_I$ is the difference of its endpoint masses,
with total variation at most two. Consequently
\[
 \|H_{I,J}\|_1=|I||J|,\qquad
 \|H'_{I,J}\|_1\le2|I|,\qquad
 \|D^2H_{I,J}\|_{\TV}\le4.
\]
Apply these inequalities before integrating over $x$. Summing over
$k$ counts each pair $(d,e)$ exactly once. The interval partitions give
$\sum_d|I_{R,d}(x)|=\tau_R(x)$ and the analogous terminal identity.
They yield all three bounds $A_0,A_1,A_2$. Differentiation is justified
in the distributional sense by testing against compactly supported
smooth functions and using these finite absolute bounds.

For the Lipschitz assertion it is sharper to work on the source.
Denote the integrand event indicator at time $t$ by $a_k(t;x,a)$.
The unique initial interval and the final partition imply
\begin{align*}
 \sum_k|a_k(t;x,a)-a_k(s;x,a)|
 \le \sum_{d,e}\one_{I_{R,d}(x)}(a)
 \big|\one_{I_{R,e}(T^mx)}(t+a-S_m)
       -\one_{I_{R,e}(T^mx)}(s+a-S_m)\big|.
\end{align*}
Sum first over $d$, integrate in $a$, and enlarge its domain to the
real line. Translation of each final interval costs at most $2|t-s|$.
Integrating over $\nu$ proves the last bound with $L_0$. It also proves
continuity of the formula, including at $t=0$.
\end{proof}

\subsection{An absolutely convergent physical inversion}
Use $\widehat f(b)=\int e^{ibt}f(t)\,dt$. For $u\in\T^2$ and
$b\in\R$ define the single-flight amplitudes
\begin{align}
 A^-_R(x;u,b)&=\int_0^{\tau_R(x)}
           e^{-i(u\cdot d_R(x,a)+ba)}\,da,\notag\\
 A^+_R(x;u,b)&=\int_0^{\tau_R(x)}
           e^{ i(u\cdot d_R(x,a)+ba)}\,da,
                                    \label{eq:single-flight-amplitudes}\\
 \mathcal F_{m,R}(u,b)&=\int_M
 e^{i(u\cdot K^{\rm c}_{m,R}+bS_{m,R})}
 A^-_R(x;u,b)A^+_R(T_R^mx;u,b)\,d\nu(x).
                                    \label{eq:stationary-physical-transform}
\end{align}
The torus is represented by $[-\pi,\pi]^2$. For real frequencies the
two single-flight amplitudes are complex conjugates at the same state.

\begin{theorem}[Uniform pointwise far-frequency bound]
\label{thm:stationary-polynomial-band}
The probability of the original singleton physical event satisfies
\begin{equation}\label{eq:stationary-absolute-inversion}
 p_{m,R}(k,t)=\frac1{(2\pi)^3\bar\tau_R}
 \int_{[-\pi,\pi]^2}\int_\R
 e^{-i(u\cdot k+bt)}\mathcal F_{m,R}(u,b)\,db\,du.
\end{equation}
The integral is absolutely convergent. If $p^{\rm sharp}_{m,R,B}$
denotes the same inverse restricted to $|b|\le B$, then
\begin{equation}\label{eq:stationary-sharp-pointwise-tail}
 \sum_k\|p_{m,R}(k,\cdot)-p^{\rm sharp}_{m,R,B}(k,\cdot)\|_\infty
                          \le A_2/(\pi B),\qquad B>0.
\end{equation}
In particular $B_m=m^{P+3/2}$ gives error $O(m^{-P})$ after
multiplication by the physical local scale $m^{3/2}$.
\end{theorem}
\begin{proof}
Each amplitude is a sum of at most $J_*$ interval integrals, so
\begin{equation}\label{eq:flight-amplitude-decay}
 |A_R^\pm(x;u,b)|\le\min\{\tau_+,2J_*/|b|\}.
\end{equation}
The product in \eqref{eq:stationary-physical-transform} is therefore
bounded by $\min\{\tau_+^2,4J_*^2/b^2\}$, uniformly in $m,R,u$.
Fourier transformation of \eqref{eq:flight-overlap-kernel} shows that
$\mathcal F_{m,R}/\bar\tau_R$ is exactly the Fourier transform of the
finite family of time profiles. Fourier inversion and torus
orthogonality prove \eqref{eq:stationary-absolute-inversion}; continuity
from Theorem~\ref{thm:stationary-overlap-regularity} makes the identity
pointwise. No dynamical decay inside the finite band is invoked.

For the stronger summed bound, integration by parts in the sense of
distributions gives, for $b\ne0$,
\[
 \sum_k|\widehat p_{m,R}(k,b)|
 \le |b|^{-2}\sum_k\|D_t^2p_{m,R}(k,\cdot)\|_{\TV}
 \le A_2/b^2.
\]
Integrating over $|b|>B$ with the one-dimensional inverse factor
$1/(2\pi)$ proves \eqref{eq:stationary-sharp-pointwise-tail}.
The factor from the lattice inversion has already been accounted for
in each Fourier coefficient. This avoids a factor equal to the number
of spatial labels.
\end{proof}

\subsection{A positive likelihood on the same physical trajectory}
Let $k_B$ and $\eta$ be the positive kernel and its transform from
Lemma~\ref{lem:positive-kernel-moments}. For a set
$\mathcal K\subset\Z^2$ define, on the original stationary source,
\begin{equation}\label{eq:physical-time-positive-likelihood}
 A_{m,\mathcal K}(t)=\{C_R(t)=m,W_R(t)\in\mathcal K\},\qquad
 \lambda_{m,\mathcal K,B}(t;\omega)
   =\int_\R k_B(h)\one_{A_{m,\mathcal K}(t-h)}(\omega)\,dh.
\end{equation}
Negative-time events in this formula are zero. These are numerical
weights on one trajectory. They do not change its physical dynamics,
its age, its lattice convention or the target event.

\begin{theorem}[Polynomial-band reconstruction of the original event]
\label{thm:physical-time-selector-reconstruction}
For every $m\ge1$, $R$, $t\ge0$, $B>0$ and spatial set $\mathcal K$,
\begin{equation}\label{eq:physical-time-source-TV}
 0\le\lambda_{m,\mathcal K,B}\le1,\qquad
 \|\one_{A_{m,\mathcal K}(t)}-\lambda_{m,\mathcal K,B}(t)\|_{L^1(\mathbb P_R)}
                          \le L_0\kappa_1/B=:E_B^{\rm phys}.
\end{equation}
Consequently, for every bounded measurable trajectory statistic $W$,
\begin{equation}\label{eq:physical-time-all-selectors}
 |\mathbb E_R[W\one_{A_{m,\mathcal K}(t)}]
     -\mathbb E_R[W\lambda_{m,\mathcal K,B}(t)]|
                           \le\|W\|_\infty E_B^{\rm phys}.
\end{equation}
Uniformity permits $W$ to depend on $m,t,B$. It follows by source
contraction, not a spectral estimate for the weighted transform.

For $\mathcal K=\{k\}$ the second expectation has the finite formula
\begin{equation}\label{eq:physical-time-weighted-formula}
 \begin{aligned}
 \frac1{(2\pi)^3\bar\tau_R}
 \int_{[-\pi,\pi]^2}\int_{-B}^B
 &e^{-i(u\cdot k+bt)}\eta(b/B)\\
 {}\cdot\int_M &e^{i(u\cdot K^{\rm c}_{m,R}+bS_{m,R})}
 A^-_{R,W}(x;u,b)A^+_R(T_R^mx;u,b)\,d\nu(x)\,db\,du,
 \end{aligned}
\end{equation}
where
\[
 A^-_{R,W}(x;u,b)=\int_0^{\tau_R(x)}
             W(x,a)e^{-i(u\cdot d_R(x,a)+ba)}\,da.
\]
The statistic $W$ is held fixed while taking the observation-time
transform. A fixed bounded target-time interval $I$ may also be
integrated: its absolute weighted error is at most
$|I|\|W\|_\infty E_B^{\rm phys}$.
\end{theorem}
\begin{proof}
The source argument in the proof of
Theorem~\ref{thm:stationary-overlap-regularity}, before averaging the
event indicators, shows the stronger estimate
\[
 \sum_k\mathbb E_R|\one_{A_{m,\{k\}}(t)}-
                       \one_{A_{m,\{k\}}(s)}|\le L_0|t-s|
\]
for all real $t,s$, with the zero extension already established.
A spatial set can only reduce the corresponding bound. Average it
against $k_B(t-s)$ and use positivity and
$\int |h|k_B(h)\,dh=\kappa_1/B$. This proves
\eqref{eq:physical-time-source-TV}. Multiplication by $W$ and then
integration prove \eqref{eq:physical-time-all-selectors}. In particular
this is total variation of unnormalized measures on the original
stationary source, not merely a difference of scalar probabilities.

For fixed $(x,a)$ and $m$, the time event is supported in an interval
of length at most $\tau_+$. Its time Fourier transform is obtained by
the substitution $t=S_m+v-a$. Integrate $v$ over the final cell
intervals and $a$ over the initial ones. The factor $W(x,a)$ remains
inside the latter integral and gives $A^-_{R,W}$. Invert the positive
kernel and the two lattice coordinates to obtain
\eqref{eq:physical-time-weighted-formula}. The band is finite and all
age integrals and the statistic are bounded, so Fubini is legitimate.
Finally integrate the established error over $I$.
\end{proof}

\begin{corollary}[Denominator-sensitive physical posteriors]
\label{cor:direct-physical-posterior}
Let $A=A_{m,\mathcal K}(t)$, $p_A=\mathbb P_R(A)$ and
$p_{A,B}=\mathbb E_R\lambda_{m,\mathcal K,B}(t)$. If
$p_{A,B}>E_B^{\rm phys}$, then
\[
 p_A\ge p_{A,B}-E_B^{\rm phys}>0,\qquad
 \left\|\mathbb P_R(\cdot\mid A)-
          \frac{\lambda_{m,\mathcal K,B}(t)}{p_{A,B}}\mathbb P_R\right\|_{\TV}
                \le\frac{2E_B^{\rm phys}}{p_{A,B}-E_B^{\rm phys}}.
\]
If instead a lower bound for $p_A>E_B^{\rm phys}$ is known, the right
side may be replaced by $2E_B^{\rm phys}/p_A$. These assertions hold
also after multiplication by any fixed nonnegative selector at most
one, using its own denominator. For $p_A\ge c m^{-3/2}$ and
$B_m=m^{P+3/2}$, the relative error is $O(m^{-P})$.
\end{corollary}
\begin{proof}
Apply the normalization inequality in
Corollary~\ref{cor:positive-microscopic-posterior} to
\eqref{eq:physical-time-source-TV}. Multiplication by a selector
contracts the source $L^1$ error. The last assertion is division by
the explicitly stated denominator bound, not a proof of that bound.
\end{proof}

\paragraph{Which tail has been closed.}
Equations~\eqref{eq:stationary-sharp-pointwise-tail} and
\eqref{eq:physical-time-source-TV} use a polynomial bandwidth for the
\emph{original stationary observation} at a fixed collision count.
They require no removal of a small-mass source set, no long-itinerary
coarea derivative bound, and no comparison with an event at the preceding
section origin. The time profiles have uniform distributional
curvature because two physical flight ages have been integrated.
The return density in Theorem W has a different output and does not
inherit this estimate. In particular a bound for these profiles is not
a pointwise estimate for the four-coordinate return density.

\section{The Gaussian central term for an exact physical singleton}
\label{sec:stationary-gaussian-complement}
We now evaluate the central part of
\eqref{eq:stationary-absolute-inversion}. This calculation takes place
in the actual three frequency variables $(u_1,u_2,b)$ of a prescribed
physical collision count. It is not a trace of a four-dimensional
integrated estimate. Together with the preceding uniform time tail it
reduces the original stationary singleton probability to one finite,
signed middle-frequency integral at a polynomial bandwidth.

Define the bounded centered collision record and its covariance by
\begin{equation}\label{eq:physical-three-collision-record}
 f^{\rm c}_R=(\kappa_{R,1},\kappa_{R,2},\tau_R-\bar\tau_R)
       =P_Rh_R,\qquad
 P_R(w_1,w_2,a,b)=(w_1,w_2,b-\bar\tau_Ra),\qquad
 \Sigma_R=P_R\Gamma_RP_R^{\mathsf T}.
\end{equation}
The map $P_R$ has full row rank with uniformly bounded coefficients.
Thus $\Sigma_R$ is uniformly positive definite by Theorem E. Its
Green--Kubo formula and its smoothed version follow by applying $P_R$
to the collision formulas; no new correlation theorem is needed.

\subsection{A three-dimensional endpoint-weighted central estimate}
\begin{lemma}[Two roof weights at the collision endpoints]
\label{lem:two-roof-central-comparison}
For $m\ge2$, uniformly in $R$,
\begin{align}
 &\int_{|v|\le2m^{1/200}}\left|
 \int_M\tau_R(x)\tau_R(T_R^mx)
       e^{iv\cdot\sum_{j<m}f^{\rm c}_R(T_R^jx)/\sqrt m}\,d\nu(x)
       -\bar\tau_R^2e^{-v^{\mathsf T}\Sigma_Rv/2}
 \right|\,dv\notag\\
 &\hspace{24mm}\le C m^{-11/700}\sqrt{\log(2+m)}.
 \label{eq:two-roof-central-comparison}
\end{align}
The integration is with respect to three-dimensional Lebesgue measure.
\end{lemma}
\begin{proof}
Put $h_\delta=\mathsf S_\delta h_R$ and
$\tau_\delta=\mathsf S_\delta\tau_R$, and let
$z=P_R^{\mathsf T}v/\sqrt m$. The collision transfer convention gives
exactly the smoothed pairing
\[
 \ell_j\bigl(M_{\tau_\delta}\mathcal L_{R,\delta,z}^{m}
                                        (\tau_\delta\nu)\bigr).
\]
The initial vector and terminal functional each have norm
$O(\delta^{-2})$. The spectral splitting of
Lemma~\ref{lem:smooth-collision-expansion} therefore has complementary
error $O(\delta^{-4}\rho^m)$. Its principal amplitude at $z=0$ is
$\bar\tau_R^2$, because smoothing preserves the mean. The difference
of that amplitude from its value at zero is bounded by
$C\delta^{-6}|v|/\sqrt m$: two endpoint norms and one projection
difference have all been paid.

Write $\ell_\delta=1+|\log\delta|$. The same spectral lemma gives
\[
 m\log\lambda_{R,\delta}(z)
 =-\tfrac12v^{\mathsf T}P_R\Gamma_{R,\delta}P_R^{\mathsf T}v
                           +O(\delta^{-6}|v|^3/\sqrt m).
\]
When the displayed cubic remainder and $|z|\delta^{-2}$ are small,
positive semidefiniteness of the smoothed covariance bounds the
principal power by a fixed constant. The covariance replacement costs
$C|v|^2\delta\ell_\delta$.

Replacing either endpoint $\tau_\delta$ by $\tau_R$ costs $C\delta$,
by invariance, the $L^1$ smoothing error and the uniformly bounded
other factors. The product of the two smoothed roofs is bounded
pointwise independently of $m,\delta$. The small-BV variance bound
\eqref{eq:small-bv-variance}, applied to each component of
$P_R(h_R-h_\delta)$, consequently bounds the unsmoothing of the phase
by $C|v|\sqrt{\delta\ell_\delta}$. The resulting pointwise error is
\[
 C\left[\delta+|v|\sqrt{\delta\ell_\delta}
 +|v|^2\delta\ell_\delta
 +m^{-1/2}\delta^{-6}(|v|+|v|^3)+\delta^{-4}\rho^m\right].
\]
No long pullback of an endpoint weight is estimated in BV.

Take $\delta=m^{-1/14}/4$ and integrate on the three-dimensional
ball of radius $2m^{1/200}$. In the order just displayed, the five
polynomial decay exponents are
\[
 \frac{79}{1400},\qquad \frac{11}{700},\qquad
 \frac{13}{280},\qquad \frac9{175},\qquad \frac{29}{700}.
\]
The second term is the largest; the extra logarithm in the covariance
term is absorbed by its strict power margin. The spectral smallness
requirements follow from
\[
 \frac12-\frac1{200}-\frac17>0,\qquad
 \frac12-\frac3{200}-\frac37>0.
\]
The complementary term is exponentially small. Increasing $C$ for
finitely many initial $m$ proves \eqref{eq:two-roof-central-comparison}.
\end{proof}

Set $r_m=2m^{-99/200}$, so that $\sqrt m\,r_m=2m^{1/200}$.
In the neighborhood $|(u,b)|\le r_m$, the bounded cell labels and ages
in Lemma~\ref{lem:finite-flight-cell-partition} give
\begin{equation}\label{eq:flight-amplitude-central-approximation}
 |A_R^\pm(x;u,b)-\tau_R(x)|\le C|(u,b)|.
\end{equation}
This estimate uses the exact cell labels. It does not require the
amplitudes themselves to have a uniform BV norm at every frequency.
For $m$ sufficiently large, this ball is inside the chosen torus chart.
Define its inverse contribution by
\begin{equation}\label{eq:stationary-central-inverse}
 G^{\rm cen}_{m,R}(k,t)=\frac1{(2\pi)^3\bar\tau_R}
 \int_{|(u,b)|\le r_m}
    e^{-i(u\cdot k+bt)}\mathcal F_{m,R}(u,b)\,du\,db.
\end{equation}

\begin{theorem}[Evaluated Gaussian central term]
\label{thm:stationary-singleton-central-term}
Writing
\[
 \zeta_{m,R}(k,t)=\left(\frac{k}{\sqrt m},
                             \frac{t-m\bar\tau_R}{\sqrt m}\right),
\]
one has, uniformly in $R,k,t$ and all sufficiently large $m$,
\begin{equation}\label{eq:stationary-singleton-central-term}
 \left|m^{3/2}G^{\rm cen}_{m,R}(k,t)
             -\bar\tau_R g_{\Sigma_R}(\zeta_{m,R}(k,t))\right|
                \le C m^{-11/700}\sqrt{\log(2+m)}.
\end{equation}
Here $g_{\Sigma}$ is the normalized Gaussian density on $\R^3$.
\end{theorem}
\begin{proof}
Remove $e^{ibm\bar\tau_R}$ from
\eqref{eq:stationary-physical-transform}. Substitution
$(u,b)=v/\sqrt m$ in \eqref{eq:stationary-central-inverse} gives
an inverse on the actual three-dimensional ball. By
\eqref{eq:flight-amplitude-central-approximation}, replacement of
the two amplitudes by the endpoint roofs costs at most
\[
 C m^{-1/2}\int_{|v|\le2m^{1/200}}|v|\,dv
                         \le C m^{-12/25}
\]
after multiplication by $m^{3/2}$. Apply
Lemma~\ref{lem:two-roof-central-comparison} to the resulting transform.
The error in an inverse is bounded by its integrated transform error,
uniformly in the target. Uniform ellipticity gives a Gaussian tail
$Ce^{-c m^{1/100}}$ outside this ball. The full Gaussian inverse is
$g_{\Sigma_R}(\zeta)$ with factor $\bar\tau_R^2/\bar\tau_R$.
This proves the normalization and the asserted bound.
\end{proof}

\subsection{One finite signed complement for the original physical event}
For $B\ge1$ and sufficiently large $m$ put
\begin{equation}\label{eq:stationary-middle-complement}
 \mathfrak M_{m,R,B}(k,t)
 =\frac1{(2\pi)^3\bar\tau_R}
 \int_{\substack{u\in[-\pi,\pi]^2,\ |b|\le B\\ |(u,b)|>r_m}}
 e^{-i(u\cdot k+bt)}\mathcal F_{m,R}(u,b)\,db\,du.
\end{equation}
This is a signed inverse evaluated at the target, not an integral of
an absolute value. Its domain contains all noncentral spatial torus
frequencies and all remaining time frequencies below $B$.

\begin{theorem}[A pointwise Gaussian-plus-complement identity]
\label{thm:physical-singleton-reduction}
For every $B\ge1$ and all sufficiently large $m$,
\begin{align}
 &\sup_{R,k,t}\big|m^{3/2}p_{m,R}(k,t)
  -\bar\tau_Rg_{\Sigma_R}(\zeta_{m,R}(k,t))
  -m^{3/2}\mathfrak M_{m,R,B}(k,t)\big|\notag\\
 &\hspace{18mm}\le C\left(m^{-11/700}\sqrt{\log(2+m)}+\frac{m^{3/2}}B\right).
 \label{eq:physical-singleton-local-ledger}
\end{align}
There is no discarded-density term in this identity. For
$B_m=m^{P+3/2}$, its last term is $O(m^{-P})$.

More intrinsically, put
\[
 \xi_{t,m,R}(k)=\left(\frac{k}{\sqrt t},
                              \frac{m-t/\bar\tau_R}{\sqrt t}\right).
\]
For every fixed $M<\infty$, uniformly in $R$ and $k,m,t$ satisfying
$t\to\infty$ and $|\xi_{t,m,R}(k)|\le M$,
\begin{align}
 t^{3/2}p_{m,R}(k,t)
 &=g_{\mathcal V_R}(\xi_{t,m,R}(k))
       +t^{3/2}\mathfrak M_{m,R,B_m}(k,t)\notag\\
 &\quad+O_M\left(t^{-11/700}\sqrt{\log(2+t)}+t^{-P}\right).
 \label{eq:physical-singleton-intrinsic-ledger}
\end{align}
Thus the Gaussian denominator asymptotic for these original physical
singletons is equivalent to the vanishing of the displayed normalized
signed complement on each such central set.
\end{theorem}
\begin{proof}
Split the absolutely convergent integral in
\eqref{eq:stationary-absolute-inversion} into its central ball, the
finite middle region \eqref{eq:stationary-middle-complement}, and
$|b|>B$. Theorem~\ref{thm:stationary-singleton-central-term} evaluates
the first part and Theorem~\ref{thm:stationary-polynomial-band} bounds
the third. This gives \eqref{eq:physical-singleton-local-ledger} for
the actual event, pointwise in observation time. No essential supremum
of a discarded return density has entered the argument.

Let $D_R^{\rm clk}=\operatorname{diag}(1,1,-1/\bar\tau_R)$. The
physical covariance defined earlier satisfies exactly
\begin{equation}\label{eq:physical-covariance-jacobian}
 \mathcal V_R=\bar\tau_R^{-1}
              D_R^{\rm clk}\Sigma_R(D_R^{\rm clk})^{\mathsf T},\qquad
 \det\mathcal V_R=\det\Sigma_R/\bar\tau_R^5.
\end{equation}
Indeed $D_R^{\rm clk}P_R$ is the map $B_R$ of the physical-time
functional theorem. If $|\xi|\le M$, then
$m=t/\bar\tau_R+O_M(\sqrt t)$ uniformly. Set
$A_R=\sqrt{\bar\tau_R}\operatorname{diag}(1,1,-\bar\tau_R)$.
Its absolute determinant is $\bar\tau_R^{5/2}$ and
$A_R\mathcal V_RA_R^{\mathsf T}=\Sigma_R$. Also
\[
 \zeta_{m,R}=A_R\xi+O_M(t^{-1/2}),\qquad
 \bar\tau_R(t/m)^{3/2}=\bar\tau_R^{5/2}+O_M(t^{-1/2}).
\]
Uniform Gaussian derivative bounds on compact sets therefore show
\[
 \bar\tau_R(t/m)^{3/2}g_{\Sigma_R}(\zeta_{m,R})
                  =g_{\mathcal V_R}(\xi)+O_M(t^{-1/2}).
\]
Multiplying \eqref{eq:physical-singleton-local-ledger} by $(t/m)^{3/2}$
proves \eqref{eq:physical-singleton-intrinsic-ledger}. Since the
remaining explicit error tends to zero, subtraction proves both
directions of the asserted equivalence. A signed complement may
vanish even if its absolute transform integral is not small.
\end{proof}

\begin{corollary}[What a complement estimate would imply]
\label{cor:stationary-singleton-denominator-criterion}
If the normalized middle term in
\eqref{eq:physical-singleton-intrinsic-ledger} tends to zero uniformly
on a fixed central set, then there
\[
 \mathbb P_R\{W_R(t)=k,C_R(t)=m\}
           =t^{-3/2}g_{\mathcal V_R}(\xi_{t,m,R}(k))(1+o(1)),
\]
and this denominator is uniformly bounded below by $c_Mt^{-3/2}$.
For any $P>0$, Theorem~\ref{thm:physical-time-selector-reconstruction}
with $B_m=m^{P+3/2}$ then approximates its conditional initial law by
the positive finite-band likelihood with total variation error
$O_M(t^{-P})$. A selected event needs its own selected denominator;
the unselected Gaussian formula does not establish it.
\end{corollary}
\begin{proof}
Uniform ellipticity makes the Gaussian density uniformly positive on
the compact central set. Use
Theorem~\ref{thm:physical-singleton-reduction} and divide the source
error in Corollary~\ref{cor:direct-physical-posterior} by this lower
bound. The final sentence specifies the normalization required after
multiplication by a selector.
\end{proof}

\subsection{Relation to the return-density and suspension theories}
The cell-displacement local theorem of \cite{SV} and the suspension
mixing-local-limit framework of \cite[Definition 2.2 and Section 2.4]{DN}
are natural comparisons. In the latter framework a local theorem on
the base is an input to the suspension argument. We have not imported
that input for the moving-section four-coordinate record. The present
calculation instead keeps one prescribed physical collision count and
the complete cell corrections, integrates the true flight ages, and
obtains a uniform pointwise far-time-frequency estimate directly.
Its two single-flight amplitudes and their inverse-square product
are the mechanism behind the polynomial bandwidth. They do not by
themselves evaluate the middle-frequency collision correlation.

For reference the three norms and the completed steps can be recorded
without identifying their endpoints:
\begin{center}
\small
\renewcommand{\arraystretch}{1.15}
\begin{tabular}{@{}p{0.23\textwidth}p{0.37\textwidth}p{0.32\textwidth}@{}}
\hline
Object & Proved estimate & Remaining local task\\
\hline
Preceding-section return law & Theorem W: global mixed $L^1$ and
source-TV approximation; $B_n$ may be $\exp(O(n\log n))$ & The common
pointwise return correction and its finite complement.\\
Actual physical singleton & Theorem X: a pointwise $O(B^{-1})$ tail,
a Gaussian central term and a polynomial-band exact ledger & The
normalized signed middle term in
\eqref{eq:physical-singleton-intrinsic-ledger}.\\
Physical weighted event & Absolute source-TV error $E_B^{\rm phys}$
for every bounded selector & Its evaluated weighted formula and positive
selected denominator before relative conditioning.\\
\hline
\end{tabular}
\end{center}
The local physical remainder is now one explicitly specified integral,
with no discarded-density or event-replacement error. This is a direct
analysis of the original physical event, not a proof that it coincides
with a preceding-section event. The full four-coordinate return LLT
and the microscopic Gaussian denominator remain the stated targets;
neither is relabeled as a consequence of the new tail estimate.

\section{Conditioning, clocks, and geometric prediction}
\begin{proposition}[Central density conditioning]\label{prop:conditioning}
Suppose an unweighted and an $f$-weighted joint density satisfy
\[
 n^2p_n(z)=G_n(z)+r_n(z),\qquad
 n^2p_n^f(z)=a_fG_n(z)+r_n^f(z)
\]
on a common central window, with $G_n\ge g_0>0$ and $|r_n|,|r_n^f|\le\varepsilon_n<g_0$. Then their density ratio differs from $a_f$ by at most
\begin{equation}\label{eq:ratio}
          \frac{(1+|a_f|)\varepsilon_n}{g_0-\varepsilon_n}.
\end{equation}
The same bound holds after conditioning on fixed lattice coordinates and any positive-length roof interval contained in that window, including intervals whose lengths shrink with $n$.
\end{proposition}
\begin{proof}
Subtract $a_f$ times the first identity from the second and divide by $G_n+r_n\ge g_0-\varepsilon_n$. For an interval, integrate both identities first; the errors are at most its length times $\varepsilon_n$ and the denominator is at least its length times $g_0-\varepsilon_n$. The length cancels.
\end{proof}
At an exact roof value this is a density-defined version of a regular conditional expectation for almost every value with positive density. It is not conditioning on a positive-probability singleton. It supplies no estimate outside the common central window and no functional path convergence by itself. These are the precise limitations of this interface to stopped-path large deviations.

\begin{proposition}[A conditional physical-clock transfer]\label{prop:clock}
Let $(X_j,\tau_j)\in\R^d\times(0,\infty)$ be stationary dynamical observations with means $(a,\mu)$, $\mu>0$, and a joint functional CLT with covariance $\Sigma$. Assume the uniform law of large numbers for their clock and that initial and terminal residual records on bounded rescaled time intervals are $o_P(\sqrt t)$. Then the reward accumulated up to physical time $t$, centered by $(a/\mu)t$, has limiting covariance per unit time
\begin{equation}\label{eq:clock}
 \frac1\mu Q\Sigma Q^T,\qquad Q=(I_d,-a/\mu).
\end{equation}
For $X=(\kappa_1,\kappa_2,r)$ this transfers induced displacement and collision counts, not the number of experimental bits.
\end{proposition}
\begin{proof}
The clock law gives its inverse index $N(t)/t\to1/\mu$. The centered functional CLT evaluated at this inverse time and the identity $S_{N(t)}\tau=t+o_P(\sqrt t)$ give
\[
 S_{N(t)}X-(a/\mu)t
 = (S_{N(t)}X-N(t)a)-(a/\mu)(S_{N(t)}\tau-N(t)\mu)+o_P(\sqrt t).
\]
The continuous limit permits the random-time substitution. Projection by $Q$ and evaluation at time $1/\mu$ prove \eqref{eq:clock}. The assumed residual estimate includes unfinished induced returns, which are not bounded merely because individual flights are bounded.
\end{proof}

\begin{lemma}[Covariance continuity from a finite-time comparison]\label{lem:covariance}
Suppose a centered correlation matrix satisfies $\|C_j(R)\|\le Ce^{-cj}$ and
$\|C_j(R)-C_j(R')\|\le C\delta^\alpha e^{a j}$, where $\delta=|R-R'|$, $a,c,\alpha>0$. Then the absolutely convergent Green--Kubo covariance satisfies
\[
 \|\Sigma_R-\Sigma_{R'}\|\le C'\delta^{\alpha c/(a+c)}.
\]
If $a=0$, the bound is instead $C'\delta^\alpha(1+|\log\delta|)$.
\end{lemma}
\begin{proof}
Split the defining correlation series at $J=\lfloor\alpha\log(1/\delta)/(a+c)\rfloor$. For $a>0$, the finite sum is bounded by $C\delta^\alpha e^{aJ}$ and the remaining sum by $Ce^{-cJ}$. Both have the stated order. For $a=0$, sum the constant comparison bound over $J+1$ terms and use the exponential tail. Symmetrizing the correlations only changes the constant.
\end{proof}
The hypotheses of this lemma must concern the actual induced observables and their varying domains. Spectral continuity for a different observable norm is not a substitute for this comparison.

\begin{proposition}[Geometric error and the central prediction scale]\label{prop:prediction}
Assume \eqref{eq:LLT} with error $\epsilon_n$, uniformly elliptic covariances, and estimates $(\widehat m,\widehat\Sigma)$ satisfying $|\widehat m-m_R|+\|\widehat\Sigma-\Sigma_R\|\le\omega(\nu)$. If $\sqrt n\,\omega(\nu)\le1$, replacing the Gaussian parameters in \eqref{eq:LLT} changes its right side on a fixed central window by at most $C(1+\sqrt n)\omega(\nu)$. The total prediction error is bounded by this quantity plus $\epsilon_n$.
\end{proposition}
\begin{proof}
The argument of the Gaussian moves by $\sqrt n(\widehat m-m_R)$. On a slightly enlarged compact window the Gaussian and its derivatives in its argument and covariance are bounded uniformly under ellipticity. The mean value theorem gives the bound. Adding the LLT error finishes the proof.
\end{proof}

There is also a direct equilibrium conclusion that needs no local-limit hypothesis. The mean collision rate is
\[
 \lambda(R)=\frac{2R}{\sqrt3/2-\pi R^2}.
\]
Its derivative is $2/A_R+4\pi R^2/A_R^2$, where $A_R=\sqrt3/2-\pi R^2$. On $I$, $A_R>17/100$, using $\sqrt3>173/100$ and $\pi<22/7$, and direct substitution gives $0<\lambda'(R)<110$. Therefore a radius estimate with error at most $\nu$ gives collision-rate error at most $110\nu$. On the circular submodel the average support function identifies the radius, independently of translation. The finite geometric conclusion of \cite{A2G} therefore transfers to this rate with its geometric observation cost, without recovering the entire unknown launch law.


\section{A logarithmic unfinished-block bound under conditioning}
\label{sec:conditioned-clock}
The stationary law of an induced block is length-biased. We use this law
for the initial block and the exact visit intensity for subsequent blocks.
These two ingredients avoid applying a launch-law union bound to a
stationary initial state.

Write $\overline\tau_R^*=\bar\tau_R/c_*$ and represent the flow as the
suspension over $(Y_R^*,F_R^*,\nu_R^*)$ with roof $\tau_R^*$. Let $t_j$
be its two-sided visit times and $x_j$ their outgoing section states.
The minimum time between consecutive visits is at least $3/50$.

\begin{lemma}[Marked visit intensity]
\label{lem:marked-visit-intensity}
For every nonnegative measurable $f$ on $Y_R^*$ and every $t\ge0$,
\begin{equation}\label{eq:marked-visit-intensity}
 \mathbb E_{\mu_R}\sum_{t_j\in(0,t]} f(x_j)
       =\frac{t}{\overline\tau_R^*}\int f\dd\nu_R^*.
\end{equation}
\end{lemma}
\begin{proof}
First take bounded $f$. Stationarity makes the locally finite measure
$A\mapsto\mathbb E_{\mu_R}\sum_{t_j\in A}f(x_j)$ translation-invariant
on $\R$, hence equal to $i_f$ times Lebesgue measure. Local finiteness
follows from the positive lower roof bound. Choose $0<\delta<3/50$.
At a stationary time, being in the first $\delta$ units of an induced
roof is the same event as a visit in the preceding interval of length
$\delta$, apart from null endpoints. There is at most one such visit,
and its mark is the current base state. Integration in the suspension
therefore gives
\[
 \delta i_f=\frac1{\overline\tau_R^*}
       \int_{Y_R^*}\int_0^\delta f(x)\dd s\dd\nu_R^*(x)
       =\frac{\delta}{\overline\tau_R^*}\int f\dd\nu_R^*.
\]
This determines $i_f$ and proves the formula for bounded $f$.
Monotone convergence extends it to arbitrary nonnegative marks, including
the case of an infinite integral. No renewal independence is involved.
\end{proof}

Let $\mathcal M_R(t)$ be the maximum of $Q_R$ over all complete induced
blocks whose time intervals meet $[0,t]$, including the block containing
time zero. Endpoint conventions at a visit are immaterial for $\mu_R$.
The marked size of any unfinished piece is bounded by a fixed multiple
of $1+\mathcal M_R(t)$.

\begin{theorem}[Uniform exponential maximum and conditioned residual]
\label{thm:conditioned-maximal-block}
There are constants $a_2,C_2>0$ such that, for every $R\in I$, $t\ge0$
and $q\ge0$,
\begin{equation}\label{eq:exponential-maximal-block}
 \mu_R\{\mathcal M_R(t)>q\}\le C_2(1+t)e^{-a_2q}.
\end{equation}
Consequently, for any measurable event $B_{R,t}$ of probability at least
$d_t>0$,
\begin{equation}\label{eq:conditioned-maximal-block}
 \mu_R\{\mathcal M_R(t)>q\mid B_{R,t}\}
       \le\min\{1,C_2(1+t)e^{-a_2q}/d_t\}.
\end{equation}
In particular, if $d_t\ge d(1+t)^{-A}$ with $d>0$ and $A\ge0$, then
$\mathcal M_R(t)=O_P(\log(2+t))$ under these conditional laws, uniformly
in $R$ and in the events. The maximum marked unfinished return during
$[0,t]$ is then $o_P(\sqrt t)$ under the same laws. More generally the
last assertion holds if $\log((1+t)/d_t)=o(\sqrt t)$.
\end{theorem}
\begin{proof}
The initial containing block has law
$\tau_R^*\dd\nu_R^*/\overline\tau_R^*$. If the maximum exceeds $q$,
either that block has $Q_R>q$ or a visit in $(0,t]$ starts a block with
$Q_R>q$. The union bound and Lemma~\ref{lem:marked-visit-intensity} give
\begin{equation}\label{eq:exact-maximal-tail-budget}
 \mu_R\{\mathcal M_R(t)>q\}
 \le\frac1{\overline\tau_R^*}
       \left(\int\tau_R^*\one_{\{Q_R>q\}}\dd\nu_R^*
                       +t\,\nu_R^*(Q_R>q)\right).
\end{equation}
The denominator is bounded below on $I$. Since $\tau_R^*\le Q_R$,
the exponential moment in \eqref{eq:exponential-return} bounds the first
integral by $C e^{-a_1q/2}$ and the probability by $C e^{-a_1q}$.
For example, $Q\le (2/a_1)e^{a_1Q/2}$ and
$e^{a_1Q/2}\one_{Q>q}\le e^{a_1Q}e^{-a_1q/2}$ suffice for the
first bound. This proves \eqref{eq:exponential-maximal-block}.

Dividing the unconditional exceptional probability by $\mu_R(B_{R,t})$
proves \eqref{eq:conditioned-maximal-block}. With a polynomial lower
bound on $d_t$, setting $q=K\log(2+t)$ and taking
$a_2K>A+1$ makes the right side tend to zero. With the more general
hypothesis, set $q=\epsilon\sqrt t$ for any $\epsilon>0$; the logarithm
of the displayed upper bound tends to minus infinity. The deterministic
bound of each marked piece by $C(1+Q_R)$ proves the assertions for all
unfinished pieces at once. The event may depend on the whole orbit and
on an estimator from that orbit; no independence is assumed.
\end{proof}

For a stationary starting point $(x,s)$ put
$V_R(v)=\max\{j:T_{j,R}(x)\le s+v\}$. Define the two path records
\[
 Z_R(v)=(V_R(v),C_R(v),W_R(v)),\qquad
 Z_R^\circ(v)=(V_R(v),N_{V_R(v),R}(x),K_{V_R(v),R}(x)).
\]
The second record counts full blocks from the beginning of the initial
containing roof, not from the observation time. It is used solely as an
exact intermediate record for comparison.

\begin{corollary}[A conditional path-comparison budget]
\label{cor:conditional-path-budget}
There is $C_3$ such that, pathwise outside the collision singularities,
\begin{equation}\label{eq:deterministic-prefix-bound}
 \sup_{0\le v\le t}|Z_R(v)-Z_R^\circ(v)|
                         \le C_3(1+\mathcal M_R(t)).
\end{equation}
Let $\widehat Z_R$ and $\widehat Z_R^\circ$ be the corresponding paths
on $[0,1]$, divided by $\sqrt t$ and with any identical deterministic
centering. For $|F|\le1$ Lipschitz with constant $L_F$ for the uniform
path metric and $\mu_R(B_{R,t})\ge d_t>0$,
\begin{align}\label{eq:conditional-path-budget}
 \left|\mathbb E_{\mu_R}[F(\widehat Z_R)\mid B_{R,t}]
       -\mathbb E_{\mu_R}[F(\widehat Z_R^\circ)\mid B_{R,t}]\right|
 &\le\frac{C_3L_F(1+q)}{\sqrt t}
        +\frac{2C_2(1+t)e^{-a_2q}}{d_t}.
\end{align}
For $t\ge2$ choose
\begin{equation}\label{eq:conditional-log-budget}
 q=a_2^{-1}\log\left(\frac{\max\{1,2C_2\}(1+t)\sqrt t}{d_t}\right).
\end{equation}
The second term is at most $t^{-1/2}$; for polynomially bounded
conditioning probabilities the total error is
$O((1+L_F)\log t/\sqrt t)$.
\end{corollary}
\begin{proof}
The collision counts differ by the initial prefix and the final
incomplete block. Each is bounded by the full block's collision count.
The same holds for the lattice displacement after multiplying by the
uniform bound on a single-flight lattice step and adding a fixed
fundamental-cell boundary constant. The visit counts agree exactly by
construction. These estimates are simultaneous for all $v\in[0,t]$,
and prove \eqref{eq:deterministic-prefix-bound}. On
$\{\mathcal M_R(t)\le q\}$ the Lipschitz bound gives the first term in
\eqref{eq:conditional-path-budget}. On its complement the difference of
two values of $F$ is at most two. Apply
\eqref{eq:conditioned-maximal-block}. The stated choice of $q$ and
substitution prove the final assertions.
\end{proof}

\begin{remark}[Which conditional interface has been closed]
\label{rem:conditional-interface-scope}
The square-root unfinished-block requirement is now a theorem for the
actual mechanical family, even after any conditioning of polynomially
small positive probability and over the whole observation interval. A
functional central limit theorem for completed records is still a
separate premise in Proposition~\ref{prop:clock}. For a local limit, the
conditioning event on the two sides of \eqref{eq:conditional-path-budget}
is the same event. Replacing an exact lattice observation by the
completed-block observation changes that event and needs a further
boundary estimate. If two events $B,B'$ have probabilities at least $d$,
then for $|f|\le1$ direct subtraction of the ratios gives
\[
 |\mathbb E[f\mid B]-\mathbb E[f\mid B']|
             \le2\mu_R(B\mathbin\triangle B')/d.
\]
Thus an $O(\log t)$ error in an integer record is not discarded when
conditioning on one lattice value. The weighted raw-density estimates
of Proposition~\ref{prop:conditioning} remain the requisite local input.
A zero-probability singleton in the continuous coordinate is not an event
to which \eqref{eq:conditioned-maximal-block} applies.
\end{remark}

\section{Induced covariance and raw-summation criteria}
\label{sec:closure-contracts}
We record the exact analytical implications needed in the low-frequency
and outer-frequency parts of the original density theorem. These
criteria distinguish proved finite-record facts from the two infinite
sums that still require estimates. The collision matrix $\Gamma_R$ and
the actual Gaussian covariance $D_R$ have already been established in
Theorems~\ref{thm:collision-covariance} and
\ref{thm:actual-record-gaussian}. The following induced Green--Kubo
criterion is stronger: neither its infinite sum nor its periodic-evaluable
representative is presumed by those theorems. The direct argument of
Theorem~\ref{thm:joint-uniform-nondegeneracy} now proves positive
definiteness of $D_R$ without either of these additional premises;
the criterion below is retained for the stronger, absolutely convergent
induced Green--Kubo representation.

\begin{proposition}[A covariance closure criterion]
\label{prop:covariance-closure}
Put $g_R=G_R-\bar G_R$. Suppose that
\begin{equation}\label{eq:covariance-tail-contract}
 \lim_{m\to\infty}\sup_R\sum_{k>m}
       \left\|\int g_R\otimes(g_R\circ(F_R^*)^k)\dd\nu_R^*\right\|=0.
\end{equation}
Suppose also that zero asymptotic variance of $v\cdot g_R$ implies a
coboundary whose transfer function has a representative for which the
coboundary identity telescopes on the five real-rank periods of
Corollary~\ref{cor:real-period-rank}. Then the Green--Kubo matrices
\[
 \Sigma_R=\int g_R\otimes g_R\dd\nu_R^*
  +\sum_{k\ge1}\int\bigl(g_R\otimes(g_R\circ F^k)
                    +(g_R\circ F^k)\otimes g_R\bigr)\dd\nu_R^*
\]
are continuous and satisfy $cI_4\le\Sigma_R\le CI_4$ on $I$ for some
$c,C>0$.
\end{proposition}
\begin{proof}
Corollary~\ref{cor:all-moment-continuity} makes every fixed-lag matrix
continuous, including the zero-lag term. The uniformly summable tail in
\eqref{eq:covariance-tail-contract} therefore gives continuity of the
series. The usual expansion of the covariance of $\sum_{j<n}g_R\circ F^j$
is a finite identity; dividing it by $n$ and using absolute summability
identifies its limit with $\Sigma_R$. Thus $\Sigma_R$ is positive
semidefinite and its directional quadratic form is the asymptotic
variance.

If $v^T\Sigma_Rv=0$, the stipulated coboundary identity gives on each of
the five selected periods
\[
 v_1K_1+v_2K_2+v_3N+v_4T-h(v\cdot\bar G_R)=0.
\]
The augmented vector $(v_1,v_2,v_3,-v\cdot\bar G_R,v_4)$ annihilates
their five-dimensional record matrix, whose determinant is nonzero by
Corollary~\ref{cor:real-period-rank}. Hence $v=0$. Each matrix is
positive definite. Continuity and compactness of $I$ give uniform lower
and upper eigenvalue bounds. The periodic representative premise is
used explicitly here; a merely measurable identity is insufficient.
\end{proof}

\begin{lemma}[An absolute raw-residual summation criterion]
\label{lem:raw-residual-sum}
Let $\mathcal W$ be countable. For $w\in\mathcal W$ assign a lattice
label $k_w\in\Z^3$ and a function $q_w\in L^1(\R)$ whose second
distributional derivative is a finite signed measure. Suppose
\[
 A_0=\sum_w\|q_w\|_1<\infty,\qquad
 A_2=\sum_w\|D^2q_w\|_{\TV}<\infty.
\]
Then
$H(u,b)=\sum_w e^{iu\cdot k_w}\widehat q_w(b)$ belongs to
$L^1(\T^3\times\R)$ and, for $B\ge1$,
\begin{align}\label{eq:absolute-residual-sum}
 \int_{\T^3\times\R}|H|&\le2(2\pi)^3(A_0+A_2),\\
 \int_{\T^3\times\{|b|>B\}}|H|&\le2(2\pi)^3A_2/B.\nonumber
\end{align}
The same estimates hold uniformly for a parameter family whenever the
corresponding sums are uniformly bounded.
\end{lemma}
\begin{proof}
Distributional differentiation and the stated Fourier convention give
$-b^2\widehat q_w(b)=\widehat{D^2q_w}(b)$. Thus
\[
 |\widehat q_w(b)|\le\min\{\|q_w\|_1,
                    |b|^{-2}\|D^2q_w\|_{\TV}\}\quad(b\ne0).
\]
Integrate the first estimate on $|b|\le1$ and the second on $|b|>1$,
and then sum by Tonelli. The torus has volume $(2\pi)^3$.
The series itself converges uniformly as a transform of the absolutely
summable $L^1$ densities. Integrating only $|b|>B$ gives the second
bound. All lattice multiplicities are included in the two sums; no
bound on the number of words has been suppressed.
\end{proof}

To use this lemma, the residuals must be those of a decomposition of the
actual full measure, including singular itinerary boundaries. A jump left
in $q_w$ contributes a derivative of a point mass to $D^2q_w$ and does
not satisfy the hypothesis. The individual jump estimates of
Proposition~\ref{prop:general-edge} do not prove the sum $A_2<\infty$.
Nor does the scalar tail \eqref{eq:exponential-return} bound inverse
Jacobian factors or the variation of a density near a critical value.
If $A_2$ grows with $n$, its actual growth must be used in the frequency
splice rather than treated as a fixed constant.

There is nevertheless a useful exact truncation budget. For an insertion
$|w|\le M_0$, deleting the physical words with $N_{n,R}>L$ changes its
record measure, and its raw transform in uniform norm, by at most
\begin{equation}\label{eq:physical-count-truncation}
 M_0C_1e^{-a_1L/n}.
\end{equation}
Indeed total variation contracts under pushforward, and
\eqref{eq:block-exponential-moment} bounds the discarded mass.
On a frequency band of volume $V$ the integrated error is at most
$VM_0C_1e^{-a_1L/n}$. Thus the sufficient physical-count budget for a
prescribed integrated error $\varepsilon$ is
\[
 L\ge\frac{n}{a_1}\max\{0,\log(VM_0C_1/\varepsilon)\}.
\]
This controls a finite-band error, not a pointwise density error or an
infinite frequency band. It is compatible with retaining the original
raw datum and with the separate all-branch derivative contract above.

\section{The joint local-limit conclusion and its realization}
\label{sec:final-realization}
The endpoint of this article remains the parameter-uniform raw mixed
density local limit \eqref{eq:LLT} for the physical record
$(K_{n,R},N_{n,R},T_{n,R})$, with the weighted positive-denominator
interface of Proposition~\ref{prop:conditioning}. Neither the return
section nor the recorded observables are replaced by independent flights,
a smoothed observation, or an unrelated geometric problem.

The change from $Y_R$ to $Y_R^*$ is explicit, and the periodic records
are those of its actual first return. The joint periodic annihilator is
trivial. The adjacent-increment comparison yields a uniform positive
period discrepancy with logarithmic collision length, and
Theorem~\ref{thm:one-step-coercivity} gives its positive-measure
consequence for controlled regular phases. The phase-reconstruction and
Fredholm hypotheses of Theorem~\ref{thm:phase-resolvent-criterion} state
precisely how these estimates would give an induced resolvent bound.
Their verification on the actual induced strong/weak spaces is still
required. No division by an uncontrolled spectral modulus is implicit.

The collision-space spectral input has a different, now completed use:
Theorem~\ref{thm:exponential-return} proves uniform exponential moments
of the genuine return block. It strengthens finite-record stability to
all finite moments. The common-chart construction in
Proposition~\ref{prop:finite-common-patches} gives an explicit
initial-state remainder for the moving-domain density comparison.
The bounded compensation of Section~\ref{sec:collision-covariance}
then proves a uniformly summable collision Green--Kubo formula. The
exact stopped identity gives the actual Gaussian covariance
$D_R=c_*^{-1}\Gamma_R$, its continuity, and its $L^2$ cohomological
kernel. These conclusions no longer depend on an assumed induced
spectral expansion. The stronger induced-correlation criterion of
Proposition~\ref{prop:covariance-closure} remains available, but is not
used to obtain these Gaussian laws. The actual normalized second
moments now converge by Theorem~\ref{thm:marked-quadratic-moments},
and Theorem~\ref{thm:joint-uniform-nondegeneracy} proves uniform positive
definiteness. Its measurable phase argument removes the section term
before using symplectic area and finite physical covers; periodic
evaluation of the transfer function is not required.

At the physical-clock level, the first-order functional law and its
geometric rate comparison remain valid. The stronger maximum estimate
in Theorem~\ref{thm:conditioned-maximal-block} now controls all unfinished
returns on a growing interval, at logarithmic scale, even after
conditioning of polynomially small positive probability. Corollary
\ref{cor:conditional-path-budget} supplies the corresponding explicit
square-root path error. Thus the unfinished-block premise no longer
needs to be assumed in that range of conditioning probabilities.
Theorem~\ref{thm:actual-return-functional} now supplies the unconditioned
functional Gaussian theorem for the completed physical return record;
Corollary~\ref{cor:physical-functional-gaussian} gives unconditioned
physical-time fluctuations as well. Neither an unconditioned functional
limit nor a path comparison under one event justifies changing an exact
lattice conditioning event. That boundary comparison is a separate
weighted local-density problem, as recorded in
Remark~\ref{rem:conditional-interface-scope}.

For the raw inversion itself, the exact edge subtraction and the local
edge criterion retain every correction term. The all-branch residual
must include central critical words and singular itinerary boundaries.
Lemma~\ref{lem:raw-residual-sum} specifies an absolute derivative-norm
criterion for that sum; no uncontrolled word multiplicity is hidden in
it. The required uniform sum has not been established here. Its growth
constants, together with the actual induced operator estimates, must
satisfy the strict splice inequalities \eqref{eq:splice} or directly
prove \eqref{eq:local-tail}. The weighted insertion class needs the same
verification. A positive conditioning denominator then follows from the
raw LLT on central windows by Proposition~\ref{prop:conditioning}.

The resulting dependency chain is exact: the geometry and periodic
coercivity, the genuine exponential return tail, the common finite-record
charts, and the conditioned unfinished-block estimate are proved;
the continuous collision Green--Kubo formula, actual Gaussian laws,
functional limits, and measurable zero-variance cohomology are also
proved. The marked moment laws and full joint nondegeneracy are proved
as well. Quantitative induced phase reconstruction or a direct complete
complementary-frequency estimate, together with the global raw residual
and weighted bounds, remains necessary to instantiate the full density
theorem. No complete unconditional raw
LLT or full A3 density input is asserted before those premises are
verified. The collision spectral theory is used at its stated smooth-multiplier
scope for killing and for the compensated Gaussian limit. No spectral
theorem about the unbounded induced observables is imported.

\subsection*{Acknowledgment and reproducibility}
AI-assisted analysis contributed to this manuscript. The supplementary
source audit fixes the mathematical baseline, the referee report and
all imported source files; the response describes the revisions to each
numbered request. Finite computations check algebra and explicit orbit
margins and do not replace the analytical proofs. Independent human
specialist review is not claimed by this manuscript.

\appendix
\section{Uniform variation and the collision-space embeddings}
\label{app:bv-details}
We give the geometric and norm details used in
Section~\ref{sec:collision-covariance}. Variation is always taken in
initial collision coordinates. The statements below do not concern
variation of a density after a many-collision pushforward.

\begin{lemma}[A finite set of possible next centers]
\label{lem:finite-next-centers}
There is a finite subset $\mathcal D$ of the triangular lattice such
that every regular next collision, for every $R\in I$, is with a disk
centered at a member of $\mathcal D$. The cardinality and the set may
be chosen independently of the initial collision state.
\end{lemma}
\begin{proof}
Let $H$ be a uniform upper bound for the free flight, as supplied by the
finite-horizon geometry in Lemma~\ref{lem:collision-spectral-input}.
A flight from $q=Rn_\alpha$ to the disk centered at $d$ ends at
$q'=d+Rn'$ and has length at most $H$. Hence
$|d|\le |q|+|q'-q|+|q'-d|\le H+2R_{\max}$.
The intersection of the lattice with this fixed ball is finite. For
example the open disks of radius $1/2$ about its lattice points are
disjoint, and lie in a ball of radius $H+2R_{\max}+1/2$; area comparison
gives a bound of $4(H+2R_{\max}+1/2)^2$ for their number. Removing the
initial disk or inadmissible centers only decreases this bound.
\end{proof}

\begin{proposition}[Uniform semialgebraic slicing]
\label{prop:semialgebraic-slicing}+Let $\Lambda$ be a semialgebraic parameter set and
$Q=(a,b)\times(c,d)$ a bounded rectangle. Suppose
$u:\Lambda\times Q\to\R$ is semialgebraic and $|u|\le M$.
There is an integer $N$, depending on the family but not on
$\lambda\in\Lambda$, such that $u_\lambda\in\mathrm{BV}(Q)$ and
\begin{equation}\label{eq:slicing-variation-bound}
 |D_1u_\lambda|(Q)\le 4NM(d-c),\qquad
 |D_2u_\lambda|(Q)\le 4NM(b-a).
\end{equation}
The zero extension outside $Q$ has bounded variation as well, with an
additional bound $2M[(b-a)+(d-c)]$ for the sum of its two variations.
The assertions depend only on the almost-everywhere class of the
functions, not on chosen values on lower-dimensional sets.
\end{proposition}
\begin{proof}
Apply differentiable cell decomposition to the graph of $u$, ordering
$(\lambda,y)$ before $x$, and refine the resulting cells by the sign
of the derivative in $x$ wherever that derivative is defined. Uniform
finiteness and one-dimensional monotonicity give a single bound $N$
for the number of monotone intervals and exceptional points in every
fiber. These are the parameterized cell and monotonicity theorems of
\cite[Sections 2.1--2.2 and 6.2]{Coste}. On a monotone interval the
variation is at most $2M$. Jumps between intervals are also at most
$2M$. Isolated assigned values are immaterial for distributional
variation. Thus the variation of the almost-everywhere one-dimensional
fiber is bounded by $4NM$, after enlarging $N$ if necessary.

For $\phi\in C_c^1(Q)$, Fubini's theorem and the one-dimensional
integration-by-parts inequality now give
\[
 \left|\int_Q u_\lambda\,\partial_x\phi\dd x\dd y\right|
 \le 4NM(d-c)\|\phi\|_\infty.
\]
This bounded functional defines a finite distributional derivative
measure, proving the first estimate. Reorder the parameters as
$(\lambda,x)$ before $y$ and repeat the argument for the second.
Taking the larger of the two finite fiber bounds gives the stated
single $N$. Extending each fiber by zero adds at most two endpoint
jumps, each of size at most $M$. The same Fubini argument proves the
claimed zero-extension bound. In particular no bound on a derivative
in the parameter is needed.
\end{proof}

\begin{proposition}[Angular chart gluing for the physical record]
\label{prop:physical-chart-gluing}
The one-collision record $f_R$ has a uniformly bounded supremum norm
and flat-coordinate variation on $M$. Adding any fixed finite union
of moving rectangular section indicators with uniformly bounded
perimeters preserves these bounds. This proves the variation assertions
for $h_R$ and $s_R$ in Lemma~\ref{lem:physical-bv}.
\end{proposition}
\begin{proof}
Choose a fixed finite angular atlas on which
$t=\tan((\alpha-\alpha_0)/2)$ lies in a bounded interval, and choose a
smooth partition of unity with supports strictly inside these charts.
The chart maps and their inverses have bounded first derivatives.
The components of $n_\alpha$ and $t_\alpha$ are rational in $t$, and
$\sqrt{1-p^2}n_\alpha+pt_\alpha$ is semialgebraic in $(t,p)$.
Lemma~\ref{lem:finite-next-centers} reduces the next collision to
finitely many candidate roots. The conditions for a root to be real,
positive, and the first admissible entry, and all comparisons between
two such roots, are semialgebraic. Therefore its minimum and the finite
lattice label define bounded semialgebraic families in $(R,t,p)$.
Ties and singular trajectories can be assigned a fixed bounded value.
They do not change the almost-everywhere record.

Proposition~\ref{prop:semialgebraic-slicing} gives a uniform variation
bound on each chart. Multiplication by a fixed smooth partition
function bounds the variation by its supremum times the original
variation plus its derivative supremum times the original $L^1$ norm.
A bounded smooth change of the angular variable preserves a variation
bound with a constant depending only on the chart and its inverse.
For completeness, this follows for smooth functions by the chain rule
and change of variables, and for a bounded-variation function by
convolution on a slightly smaller chart followed by lower semicontinuity
of distributional variation. The cutoff has support in that smaller
chart. There is therefore no unaccounted jump on the chart boundary.
The sum of the finitely many localized functions gives the asserted
bound on the cylinder. At $p=\pm1$, the fiber proof already includes
finite one-sided traces; no inverse transformation from $p$ to the
angle of incidence is used.

A rectangle indicator has variation equal to its coordinate perimeter,
up to the fixed normalization of $\nu$. The variation of the indicator
of a finite union is bounded by the sum of these perimeters, including
rectangles represented across the angular seam. Its moving center does
not change that estimate. The bounded coefficients in
\eqref{eq:compensation-normalizations} finish the assertion for $h_R$
and $s_R$.
\end{proof}

\begin{lemma}[Smooth densities and test pairings in the local spaces]
\label{lem:density-norm-details}
On each local space $(\mathcal B_j,\mathcal B_{w,j})$ used in
Lemma~\ref{lem:collision-spectral-input}, smooth functions $a,g,b$
satisfy, uniformly for the radii in that neighborhood,
\begin{align}
 \|a\nu\|_{\mathcal B_j}+|a\nu|_{w,j}
       &\le C\|a\|_{C^1},\label{eq:smooth-density-embedding}\\
 \|M_gV\|_{\mathcal B_j}&\le C\|g\|_{C^2}\|V\|_{\mathcal B_j},
 &|\ell_j(M_bV)|&\le C\|b\|_{C^2}\|V\|_{\mathcal B_j}.
 \label{eq:smooth-multiplier-pairing}
\end{align}
The function norms are in the collision coordinates used to define
the local spaces. In particular, for the mean-preserving smoothing,
\begin{equation}\label{eq:smoothed-density-vector}
 \| (\mathsf S_\delta a)\nu\|_{\mathcal B_j}
       \le C\delta^{-2}\|a\|_\infty .
\end{equation}
Here $a$ is a density relative to $\nu$, not relative to unweighted
area in the incidence-angle coordinates.
\end{lemma}
\begin{proof}
We record the norm identification explicitly. In
\cite[Sections 3.2--3.3, (3.11)--(3.13)]{DZ}, a smooth function is
identified with its density times the invariant reference measure
$\mu=c\cos\varphi\dd r\dd\varphi$. Under the local boundary
reparametrization $r=R_0\alpha$, this probability is precisely $\nu$.
Thus the smooth function representing $a\nu$ is $a$; an additional
factor of $\cos\varphi$ is not inserted into that function.

For the weak norm, each stable-curve integral is bounded by
$C\|a\|_\infty$ because the stable curves have uniformly bounded
length and the test functions have bounded supremum. For the strong
stable norm, tests on a curve $W$ have supremum at most $|W|^{-\alpha}$.
The integral is therefore at most
$\|a\|_\infty |W|^{1-\alpha}$, uniformly bounded since $\alpha<1$.
For the strong unstable norm, take two matched graph curves at distance
at most $\epsilon$ in the source metric. Their unmatched parameter
intervals have length $O(\epsilon)$. On their common interval the graph
coordinates and arclength factors differ by $O(\epsilon)$, and the
matched tests differ in supremum by at most $\epsilon$.
The difference of the two integrals is consequently bounded by
$C\|a\|_{C^1}\epsilon$. Dividing by $\epsilon^\beta$ is harmless,
since the source exponent satisfies $\beta<1$. These three estimates
prove \eqref{eq:smooth-density-embedding}, with constants uniform in
the finite local cover and its fixed charts.

For a multiplier, stable and weak tests are simply multiplied by the
restriction of $g$ to their curves. In the matched-curve term, the two
restrictions differ in $C^q$ by at most
$C\|g\|_{C^2}\epsilon^{1-q}$, by interpolation of their supremum
and Lipschitz bounds. The admissible source exponents have
$\beta<1-q$. Splitting into matched tests and this difference gives
the first bound in \eqref{eq:smooth-multiplier-pairing}; it is the
multiplier estimate already used in
Lemma~\ref{lem:collision-spectral-input}. Applying the bounded mass
functional proves the second. Finally
\eqref{eq:bv-smoothing-bounds} bounds the smoothed $C^2$ norm by
$C\delta^{-2}\|a\|_\infty$. Composition with $p=\sin\varphi$
uses only bounded derivatives of this forward map. This proves
\eqref{eq:smoothed-density-vector}.
\end{proof}

The estimates above concern smooth functions on the local collision
spaces. They do not make the hard section projection a bounded
multiplier on those spaces, nor do they realize an unbounded induced
twist. The use of that projection in the compensation argument remains
entirely on the level of ordinary functions and smooth approximations.

\section{Chronological spectral products and coarse interpolation}
\label{app:functional-details}
This appendix supplies explicit product and interpolation estimates for
Section~\ref{sec:functional-gaussian}. It also separates the treatment
of short increments from the long-block spectral argument.

\begin{lemma}[Products of differently twisted collision powers]
\label{lem:chronological-product}
Fix an integer $q\ge1$. Use one local collision space, a common smoothing
scale $\delta$, and real vectors $z_1,\ldots,z_q$ with
$|z_j|\le c\delta^2$. Let $m_j\ge1$ and assume
\[
 \delta^{-6}\sum_{j=1}^q m_j|z_j|^3\le c.
\]
Write $L_j=\mathcal L_{R,\delta,z_j}$ and let $a$ be a complex initial
insertion with bounded supremum and variation norms. Then
\begin{align}
 &\left|\ell_j\bigl(L_q^{m_q}\cdots L_1^{m_1}
                   (\mathsf S_\delta a)\nu\bigr)
       -\left(\int a\dd\nu\right)
       \exp\left(-\tfrac12\sum_{j=1}^q
                      m_j z_j^{\mathsf T}\Gamma_{R,\delta}z_j\right)
        \right|\notag\\
 &\quad\le C_q\left\{
       \delta^{-4}\sum_{j=1}^q|z_j|
       +\delta^{-6}\sum_{j=1}^qm_j|z_j|^3
       +\delta^{-2}\sum_{j=1}^q\rho^{m_j}\right\}.
       \label{eq:chronological-product-bound}
\end{align}
The index on $\ell_j$ names the chosen local space and is fixed in the
formula. The powers are in physical chronological order: the earliest
increment acts first, on the right.
\end{lemma}
\begin{proof}
Write $L_j^{m_j}=\lambda_j^{m_j}\Pi_j+E_j$.
Lemma~\ref{lem:smooth-collision-expansion} gives
$\|\Pi_j\|\le C$, $\|E_j\|\le C\rho^{m_j}$ and
$\|\Pi_j-\Pi_R\|\le C\delta^{-2}|z_j|$.
The real-frequency quadratic terms in $m_j\log\lambda_j$ are
nonpositive. The assumption on the sum of the cubic remainders bounds
each principal scalar power and every product of these powers by a
fixed constant. Expand the product into its $2^q$ terms. Every term
with a complementary factor is bounded, after evaluation, by a
constant times $\delta^{-2}\rho^{m_j}$ for at least one of its
complementary indices. The factor $\delta^{-2}$ comes only from the
initial vector, through \eqref{eq:smoothed-density-vector}; the other
projectors and powers have uniform bounds. Summing the finitely many
terms gives the last term of \eqref{eq:chronological-product-bound}.

For the all-principal term, telescope the product of projectors about
$\Pi_R$, retaining their order. Since $\Pi_R^2=\Pi_R=\nu\ell_j$,
\[
 \|\Pi_q\cdots\Pi_1-\Pi_R\|
       \le C_q\delta^{-2}\sum_j|z_j|.
\]
Evaluation on the initial vector gives the first error in
\eqref{eq:chronological-product-bound}. Its zero-frequency amplitude
is $\int a\dd\nu$. The logarithmic expansion and
$|e^w-1|\le |w|e^{|w|}$ give the second error. The principal scalar
powers remain bounded during this last comparison. This proves the
claim without introducing a factor $\delta^{-2q}$.
\end{proof}

\begin{corollary}[Finite-dimensional scales, zero blocks, and short blocks]
\label{cor:finite-dimensional-details}
For fixed distinct times $0=t_0<t_1<\cdots<t_q$ and fixed vectors
$v_j$, put
$m_j=\lfloor nt_j\rfloor-\lfloor nt_{j-1}\rfloor$,
$z_j=v_j/\sqrt n$, and $\delta_n=\tfrac14n^{-1/32}$.
The spectral product error is bounded by
\[
 C_q\left\{\delta_n^{-4}n^{-1/2}
             +\delta_n^{-6}n^{-1/2}
             +\delta_n^{-2}\sum_j\rho^{m_j}\right\}=o(1).
\]
Repeated times give zero blocks, which may be omitted exactly. A fixed
number of additional deterministic increments of lengths $m_j=o(n)$
may instead be removed from a real characteristic exponent at a cost
$C\sum_j |v_j|\sqrt{m_j/n}$ for the original collision observable.
The estimates are uniform for initial probability densities with uniformly
bounded supremum and variation norms. The short-increment estimate uses
only the supremum bound.
\end{corollary}
\begin{proof}
Distinct fixed times make all $m_j$ positive multiples of $n$ for
large $n$. Substitute the scales in
\eqref{eq:chronological-product-bound}. Smoothing the initial density
costs $O(\delta_n)$; smoothing the covariance costs
$O(\delta_n\ell_{\delta_n})$ by
\eqref{eq:collision-covariance-limit}. The residual-path bound
\eqref{eq:residual-process-maximal} removes smoothing from the original
finite-dimensional process. Integer parts in the covariance time
coefficients cost $O(n^{-1})$. This gives precisely the independent
Gaussian increments in Theorem~\ref{thm:collision-functional}.

When an interval has length zero its operator is the identity, not a
principal spectral projector, so it is omitted before this expansion.
For a deterministic short interval, uniform collision correlation
summability bounds the second moment of its original centered sum by
$Cm_j$, from every deterministic starting index under a bounded
initial density. The difference of the two characteristic exponents,
with the other increments left at their original times, is therefore
bounded in $L^1$ by $C|v_j|\sqrt{m_j/n}$. The elementary exponential
difference bound proves the last assertion. This argument does not
use a complementary spectral power for a short block, and does not
move the other increments in physical time.
\end{proof}

\begin{lemma}[The complex bound before the fourth derivative]
\label{lem:complex-fourth-details}
For one real component $u$ of $h_{R,\delta}$, the stationary pairing
has, on a complex disk $|z|<c\delta^2$, a decomposition
\[
 \int e^{izS_{m,R}u}\dd\nu=e^{m\psi(z)}A(z)+E_m(z)
\]
with $|A(z)|+|\psi(z)|\le C$ and $|E_m(z)|\le C\rho^m$, uniformly
on a slightly smaller such disk. Consequently
$|E_m^{(4)}(0)|\le C\delta^{-8}\rho^m$.
\end{lemma}
\begin{proof}
The two fixed resolvent contours in
Lemma~\ref{lem:smooth-collision-expansion} give uniformly bounded
spectral projections and complementary powers throughout the complex
ball, not just on its real diameter. Both the vector $\nu$ and the
mass functional have bounded norms. Evaluating the complementary
power between them gives the displayed complex bound for $E_m$.
The eigenvalue lies in a fixed disk about one not meeting zero, so its
chosen logarithm and the principal amplitude are uniformly bounded.
Cauchy's estimate on a circle of radius proportional to $\delta^2$
now yields the fourth-derivative estimate. There is no need to bound
the principal power uniformly for complex $z$ and arbitrary $m$.
Its derivative is calculated by the product rule, as in
Lemma~\ref{lem:smooth-fourth-moment}.
\end{proof}

\begin{proposition}[Coarse polygonal interpolation]
\label{prop:coarse-interpolation}
Let $(u_k)$ be bounded real-vector random variables, $|u_k|\le M$.
Let $L\ge1$ be an integer and suppose, under the probability in question,
that every deterministic interval of $m\ge L$ indices has
\[
 \mathbb E\left|\sum_{k=r}^{r+m-1}u_k\right|^4\le C m^2.
\]
At times $jL/n$, interpolate the sums normalized by $\sqrt n$ to obtain
$Z_n$. Let $X_n$ be the polygonal interpolation using every index.
Then, with constants independent of $L,n$,
\begin{equation}\label{eq:coarse-interpolation-bound}
 \mathbb E|Z_n(t)-Z_n(s)|^4\le C'|t-s|^2,
 \qquad \|Z_n-X_n\|_\infty\le 2ML/\sqrt n.
\end{equation}
A final incomplete coarse interval is defined by extending one full
block beyond the observation horizon and restricting the interpolation.
\end{proposition}
\begin{proof}
Write $\Delta=L/n$. For a segment of duration $h\le\Delta$ inside
one block, linear interpolation multiplies the full-block increment by
$h/\Delta$. Its fourth moment is at most
$C\Delta^2(h/\Delta)^4\le Ch^2$. For an interval spanning two adjacent
blocks, apply this estimate separately to the durations $h_1,h_2$ and
use the $L^4$ triangle inequality. Since
$(\sqrt{h_1}+\sqrt{h_2})^4\le4(h_1+h_2)^2$, the same form holds.

For a general interval, split it into its first fractional block, its
complete-block middle interval, and its last fractional block. The
middle interval, when nonempty, contains at least $L$ indices, so the
hypothesis gives the required fourth moment for its total duration.
The $L^4$ triangle inequality and
$(\sqrt{h_1}+\sqrt{h_2}+\sqrt{h_3})^4
 \le9(h_1+h_2+h_3)^2$
prove the first bound with one fixed constant. Empty pieces are simply
omitted. Within one coarse block both interpolations remain within
$ML/\sqrt n$ of the sum at the left endpoint, proving the deterministic
second bound. The endpoint convention uses a full block, so these same
estimates apply on a fixed observation interval without a new constant.
\end{proof}

Applying this proposition with $u_k=h_{R,\delta_n}\circ T_R^k$ and
$L_n=\lceil\delta_n^{-8}\rceil$ gives the interpolation estimates in
Theorem~\ref{thm:collision-functional}. Indeed
\eqref{eq:smooth-fourth-moment} is bounded by $Cm^2$ for $m\ge L_n$.
Stationarity gives the interval estimates under $\nu$, and bounded
initial densities give them by domination. The fourth-moment constants
are uniform. Together with the residual maximal estimate this proves
tightness of the original process, not just of its coarse approximation.

\begin{lemma}[Continuity of the covariance square root]
\label{lem:covariance-square-root}
On every bounded set of positive semidefinite matrices, the map
$C\mapsto C^{1/2}$ is continuous in operator norm. Consequently
$R\mapsto\operatorname{Law}(\mathcal W_{\Gamma_R})$ is continuous on
continuous path space, including at singular matrices.
\end{lemma}
\begin{proof}
Uniformly approximate the scalar square-root function on a bounded
interval containing all the spectra by polynomials. The spectral theorem
bounds the matrix approximation by the same uniform scalar error, and
each polynomial in the matrix is norm-continuous. This proves the first
claim. On one probability space take a standard four-dimensional
Brownian motion $W$ and set $\mathcal W_C=C^{1/2}W$. On any compact time
interval its sample paths are bounded almost surely. Norm continuity
of the square root therefore gives almost-sure uniform convergence of
the coupled paths and hence convergence of their laws.
\end{proof}

\clearpage
\section{Locations of the collision norm definitions}
\label{app:collision-norm-map}
The following map fixes the source conventions in
Lemma~\ref{lem:density-norm-details}. Locations refer to the 45-page
preprint of \cite{DZ}; the last column records the deduction used here,
not an additional theorem attributed to that source.

\begin{center}
\renewcommand{\arraystretch}{1.2}
\begin{tabular}{>{\raggedright\arraybackslash}p{0.24\textwidth}>{\raggedright\arraybackslash}p{0.24\textwidth}>{\raggedright\arraybackslash}p{0.40\textwidth}}
Object & Source location & Use in this article\\ \hline
Relative density & Section 3.2; Section 3.3, footnote 7 & A smooth $a$
represents $a\nu$, without another cosine factor.\\
Weak norm & Equation (3.11) & Bounded stable-curve integrals.\\
Stable norm & Equation (3.12) & The test normalization gives
$\|a\|_\infty |W|^{1-\alpha}$.\\
Unstable norm & Equation (3.13) & Matched curves give
$C\|a\|_{C^1}\epsilon^{1-\beta}$.\\
Exponent restrictions & Section 3.3, before (3.12) &
$\beta\le p-q$, $p\le1/3$, hence $\beta<1-q$.\\
Mass functional & Lemma 3.4, with test $1$ &
$|\ell(V)|\le C|V|_w$.\\
Complementary powers & Theorem 2.1(3) &
Local uniform spectral power bound.
\end{tabular}
\end{center}

Smooth multiplication is proved in
Lemma~\ref{lem:collision-spectral-input} from these definitions.
The changing boundary length is handled by Remark 2.7(b) of the source.
The table concerns collision-space distributions. It neither embeds the
hard induced twist nor estimates second derivatives of raw densities.

\clearpage
\section{Dependencies and the remaining raw estimates}
\label{app:dependency-guide}
This guide distinguishes the completed input theorems, the two exact
residual constructions and the remaining pointwise local theorem.
The covariance is $D_R$ throughout; the physical projection is
$\mathcal V_R$. All finite-frequency reconstructions use the original
record, but are not Gaussian evaluations of their finite integrals.

\begin{center}
\small
\renewcommand{\arraystretch}{1.2}
\begin{tabular}{@{}p{0.22\textwidth}p{0.35\textwidth}p{0.36\textwidth}@{}}
\hline
Stage & Established result & Interface still required\\
\hline
Collision and return statistics & Gaussian and functional laws,
fixed moments, actual covariance and phase-defect results & Their
specified frequency and regularity budgets must be respected.\\
Prescribed-count Fourier bands & Theorems N--P: isotropic, coordinate
and shape-adaptive regions & The full fixed-count complement is not
covered by those regions or by averaged-orbit bounds.\\
Coherent raw correction & Theorem Q: exact and quantitative transport
between cutoffs; Theorem R: signed window control & The common
pointwise correction itself still needs a local bound.\\
Power--logarithm residual & Exact finite-record preparation and
$W^{2,1}$ residual & Its original long-time jet budget is not supplied
by the different geometric residual.\\
Geometric residual & Theorem V: all-label second-derivative budget
at linear collision cutoff; Proposition~\ref{prop:stationary-geometric-extraction}:
actual entrance-bias extension & Removed mass is small, but its density
is not known to be uniformly small on central windows.\\
Positive reconstruction & Theorem W: mixed $L^1$ approximation and
source-TV microscopic event approximation, uniform in bounded selectors
& Evaluate the remaining finite integral; a positive likelihood alone
does not provide a Gaussian denominator.\\
Physical conditioning & Theorems S--U: mesoscopic denominators,
relative comparison and Gaussian bridges & The original microscopic
physical-event comparison requires its own raw-scale denominator.\\
\hline
\end{tabular}
\end{center}

The exact pointwise responsibility is recorded in
\eqref{eq:geometric-common-raw-ledger}. The far-roof inverse and the
convolution of the discarded measure have quantitative bounds. Its
remaining signed terms are the discarded density and the finite
complementary inverse. Their joint local essential supremum, or an
equivalent complete raw inversion estimate, has not been derived from
the total-variation reconstruction. The geometric bandwidth satisfies
$\log B_n=O(n\log n)$; it is not merely of size $n\log n$.

For weights, three statements must not be interchanged. Finite-product
BV return marks admit the geometric extraction. The inherited central
Fourier theorems have their own insertion classes. The positive event
formula in Theorem W permits every bounded measurable trajectory
selector because its error is measured on the original source, not
because all such selectors are smooth multipliers. Its normalized
bound contains the actual selected denominator, or the explicit
finite-band lower-bound test of
Corollary~\ref{cor:positive-microscopic-posterior}.

The section law and the stationary entrance law are also different.
The latter uses the original roof bias and elapsed age in
\eqref{eq:equilibrium-source-for-inversion}. It is not replaced by the
bounded current-collision marginal. Finally, collision correlation
summability and induced Cesaro covariance convergence remain distinct.
The norm imports are in Appendix~\ref{app:collision-norm-map}; the new
boundary and positivity arguments are in
Sections~\ref{sec:source-boundary-calculus} and
\ref{sec:positive-microscopic-conditioning}.

\clearpage
\begin{thebibliography}{99}
\bibitem{A2G} Q. Qi, \emph{Scalar collision laws and recognition of periodic dispersing billiards}, A2 v43, primary article and technical supplement, theta-theory repository, 2026.
\bibitem{SV} D. Sz\'asz and T. Varj\'u, \emph{Local limit theorem for the Lorentz process and its recurrence in the plane}, Ergodic Theory Dynam. Systems \textbf{24} (2004), 257--278.
\bibitem{DN} D. Dolgopyat and P. N\'andori, \emph{On mixing and the local central limit theorem for hyperbolic flows}, Ergodic Theory Dynam. Systems \textbf{40} (2020), 142--174; arXiv:1710.08568v1, Definition 2.2 and Proposition 6.3.
\bibitem{DZ} M. F. Demers and H.-K. Zhang, \emph{A functional analytic approach to perturbations of the Lorentz gas}, Comm. Math. Phys. \textbf{324} (2013), 767--830; arXiv:1210.1261.
\bibitem{BDL} V. Baladi, M. F. Demers and C. Liverani, \emph{Exponential decay of correlations for finite horizon Sinai billiard flows}, Invent. Math. \textbf{211} (2018), 39--177; arXiv:1506.02836v3, including the correction note.
\bibitem{Young} L.-S. Young, \emph{Statistical properties of dynamical systems with some hyperbolicity}, Ann. of Math. (2) \textbf{147} (1998), 585--650; author preprint ESI 445 (1997), Section 8, Theorem 6.
\bibitem{BHN} H. Bruin, M. Holland and M. Nicol, \emph{Liv\v{s}ic regularity for Markov systems}, Ergodic Theory Dynam. Systems \textbf{25} (2005), 1739--1765; doi:10.1017/S0143385705000179; arXiv:math/0503690.
\bibitem{DZ11} M. F. Demers and H.-K. Zhang, \emph{Spectral analysis of the transfer operator for the Lorentz gas}, J. Mod. Dyn. \textbf{5} (2011), 665--709; doi:10.3934/jmd.2011.5.665.
\bibitem{Coste} M. Coste, \emph{An introduction to o-minimal geometry}, lecture notes, November 1999; Dottorato di Ricerca in Matematica, Universit\`a di Pisa, Istituti Editoriali e Poligrafici Internazionali, 2000.
\bibitem{GouezelRenewal} S. Gou\"ezel, \emph{Sharp polynomial estimates for the decay of correlations}, Israel J. Math. \textbf{139} (2004), 29--65; arXiv:math/0202147.
\bibitem{MN} I. Melbourne and M. Nicol, \emph{Almost sure invariance principle for nonuniformly hyperbolic systems}, Comm. Math. Phys. \textbf{260} (2005), 131--146; arXiv:math/0503693.
\bibitem{DPZ} M. F. Demers, F. P\`ene and H.-K. Zhang, \emph{Local limit theorem for randomly deforming billiards}, preprint, 2019, arXiv:1902.06850v1.
\bibitem{AFP}
L.~Ambrosio, N.~Fusco and D.~Pallara,
\emph{Functions of Bounded Variation and Free Discontinuity Problems},
Oxford Mathematical Monographs, Clarendon Press, Oxford, 2000.
\bibitem{BPR}
S.~Basu, R.~Pollack and M.-F.~Roy,
\emph{An asymptotically tight bound on the number of semi-algebraically
connected components of realizable sign conditions},
Combinatorica \textbf{29} (2009), 523--546;
arXiv:math/0603256v3, Section 3.2, equation (3.3).
\bibitem{CMInt} R. Cluckers and D. J. Miller,
\emph{Stability under integration of sums of products of real globally
subanalytic functions and their logarithms}, Duke Math. J. \textbf{156}
(2011), no. 2, 311--348; arXiv:0911.4373v1, Theorems 1.3 and 3.11.

\bibitem{DJS} H. D\"oring, S. Jansen and K. Schubert,
\emph{The method of cumulants for the normal approximation},
Probab. Surv. \textbf{19} (2022), 185--270; doi:10.1214/22-PS7.

\bibitem{CVExt} E. Carneiro and J. D. Vaaler,
\emph{Some extremal functions in Fourier analysis, II},
Trans. Amer. Math. Soc. \textbf{362} (2010), 5803--5843;
arXiv:0809.4050, Introduction.
\bibitem{SYZ} M. Stenlund, L.-S. Young and H. Zhang,
\emph{Dispersing billiards with moving scatterers},
Comm. Math. Phys. \textbf{322} (2013), 909--955;
arXiv:1210.0011v4, Section 3.1.
\end{thebibliography}

\end{document}
